\part*{Calculus II}

\iftoggle{bsc}{

\apexchapter{Applications of Integration}{chapter:app_of_int}

We begin this chapter with a reminder of a few key concepts from \autoref{chapter:integration}. Let $f$ be a continuous function on $[a,b]$ which is partitioned into $n$ equally spaced subintervals as 
\[a=x_0 < x_1 < \cdots < x_{n-1}<x_n=b.\]
Let $\Delta x=(b-a)/n$ denote the length of the  subintervals, and let $c_i$ be any $x$-value in the $i^{\text{th}}$ subinterval. \autoref{def:rie_sum} states that the sum
\[\sum_{i=1}^n f(c_i)\Delta x\]
is a \textit{Riemann Sum.} Riemann Sums are often used to approximate some quantity (area, volume, work, pressure, etc.). The \textit{approximation} becomes \textit{exact} by taking the limit 
\[\lim_{n\to\infty} \sum_{i=1}^n f(c_i)\Delta x.\]
\autoref{thm:riemannSum} connects limits of Riemann Sums to definite integrals:
\[\lim_{n\to\infty} \sum_{i=1}^n f(c_i)\Delta x = \int_a^b f(x)\ dx.\]
Finally, the Fundamental Theorem of Calculus states how definite integrals can be evaluated using antiderivatives. 

This chapter employs the following technique to a variety of applications. Suppose the value $Q$ of a quantity is to be calculated. We first approximate the value of $Q$ using a Riemann Sum, then find the exact value via a definite integral. We spell out this technique in the following Key Idea.

\begin{keyidea}[Application of Definite Integrals Strategy]\label{idea:app_of_defint}
Let a quantity be given whose value $Q$ is to be computed.\index{integration!general application technique}
\begin{enumerate}
\item	Divide the quantity into $n$ smaller ``subquantities'' of value $Q_i$.
\item	Identify a variable $x$ and function $f(x)$ such that each subquantity can be approximated with the product $f(c_i)\Delta x$, where $\Delta x$ represents a small change in $x$. Thus $Q_i \approx f(c_i)\Delta x$.
%% A sample approximation $f(c_i)\Delta x$ of $Q_i$ is called a \textit{differential element}.
\item	Recognize that $\ds Q\approx \sum_{i=1}^n Q_i = \sum_{i=1}^n f(c_i)\Delta x$, which is a Riemann Sum.
\item	Taking the appropriate limit gives $\ds Q = \int_a^b f(x)\ dx$
\end{enumerate}
\end{keyidea}

This Key Idea will make more sense after we have had a chance to use it several times. We begin with Area Between Curves.%, which we addressed briefly in \autoref{sec:FTC}.

\section{Area Between Curves}\label{sec:ABC}

We are often interested in knowing the area of a region. Forget momentarily that we addressed this already in \autoref{sec:FTC} and approach it instead using the technique described in \autoref{idea:app_of_defint}. 

\bgroup % leaving newcommand in the mtable causes errors
\newcommand{\f}[1]{.25*sin(deg(2*#1))+1.25}
\newcommand{\g}[1]{.25*(#1-2)^2+.2}
\mtable{Subdividing a region into vertical slices and approximating the areas with rectangles.}{fig:abcintro}{\small%
\begin{tikzpicture}[alt={(a) Blue upper curve f(x) and lower curve g(x); area between them shaded from a to b; axes x and y.}]
\begin{axis}[width=1.16\marginparwidth,tick label style={font=\scriptsize},
axis y line=middle,axis x line=middle,name=myplot,axis on top,xtick=\empty,
extra x ticks={.5,3},extra x tick labels={$a$,$b$},ytick=\empty,
ymin=-.2,ymax=2,xmin=-.5,xmax=3.5]
\addplot[smooth,draw={\coloronefill},fill={\coloronefill},area style,domain=.5:3]
 {.25*sin(deg(2*x))+1.25}\closedcycle;
\addplot[smooth,draw=white,fill=white,area style,domain=.5:3]
 {.25*(x-2)^2+.2}\closedcycle;
\addplot [smooth,thick, draw={\colorone},domain=-.25:3.25] {.25*sin(deg(2*x))+1.25} node [shift={(5pt,7pt)} ,black] {\scriptsize $f(x)$};
\addplot [smooth,thick, draw={\colorone},domain=-.25:3.25] {.25*(x-2)^2+.2}node [shift={(5pt,7pt)} ,black] {\scriptsize $g(x)$};
\end{axis}
\node [right] at (myplot.right of origin) {\scriptsize $x$};
\node [above] at (myplot.above origin) {\scriptsize $y$};
\end{tikzpicture}
\\ (a)
\\
\begin{tikzpicture}[alt={(b) Same curves with evenly spaced vertical slices across the shaded region; axes x and y.}]
\begin{axis}[width=1.16\marginparwidth,tick label style={font=\scriptsize},
axis y line=middle,axis x line=middle,name=myplot,axis on top,xtick=\empty,
extra x ticks={.5,3},extra x tick labels={$a$,$b$},ytick=\empty,
ymin=-.2,ymax=2,xmin=-.5,xmax=3.5]
\addplot[smooth,draw={\coloronefill},fill={\coloronefill},area style,domain=.5:3]
 {.25*sin(deg(2*x))+1.25}\closedcycle;
\addplot[smooth,draw=white,fill=white,area style,domain=.5:3]
 {.25*(x-2)^2+.2}\closedcycle;
\addplot [smooth,thick, draw={\colorone},domain=-.25:3.25] {.25*sin(deg(2*x))+1.25} node [shift={(5pt,7pt)} ,black] {\scriptsize $f(x)$};
\addplot [smooth,thick, draw={\colorone},domain=-.25:3.25] {.25*(x-2)^2+.2}node [shift={(5pt,7pt)} ,black] {\scriptsize $g(x)$};
\foreach\x in{.5,.75,...,3} {
 \edef\temp{\noexpand\draw[thick]
  (axis cs:\x,{.25*(\x-2)^2+.2})--(axis cs:\x,{.25*sin(deg(2*\x))+1.25});}\temp
}
\end{axis}
\node [right] at (myplot.right of origin) {\scriptsize $x$};
\node [above] at (myplot.above origin) {\scriptsize $y$};
\end{tikzpicture}
\\(b)
\\
\begin{tikzpicture}[alt={(c) Same curves; each slice shows a sample point and a rectangle whose height is f(ci) − g(ci); axes x and y.}]
\begin{axis}[width=1.16\marginparwidth,tick label style={font=\scriptsize},
axis y line=middle,axis x line=middle,name=myplot,axis on top,xtick=\empty,
extra x ticks={.5,3},extra x tick labels={$a$,$b$},ytick=\empty,
ymin=-.2,ymax=2,xmin=-.5,xmax=3.5]
\addplot[smooth,draw={\coloronefill},fill={\coloronefill},area style,domain=.5:3]
 {\f{x}}\closedcycle;
\addplot[smooth,draw=white,fill=white,area style,domain=.5:3]
 {\g{x}}\closedcycle;
% {.25*(x-2)^2+.2}\closedcycle;
\addplot [smooth,thick, draw={\colorone},domain=-.25:3.25] {\f{x}} node [shift={(5pt,7pt)} ,black] {\scriptsize $f(x)$};
\addplot [smooth,thick, draw={\colorone},domain=-.25:3.25] {\g{x}}node [shift={(5pt,7pt)} ,black] {\scriptsize $g(x)$};
\foreach\x in{.5,.75,...,2.75} {
 \edef\c{(\x+.25*(1+sin(deg(3*\x)))/2)}
 \edef\temp{
  \noexpand\draw[thick](axis cs:\x,{\f{\c}})rectangle(axis cs:{\x+.25},{\g{\c}});
  \noexpand\filldraw(axis cs:{\c},{\f{\c}})circle(1.3pt);
  \noexpand\filldraw(axis cs:{\c},{\g{\c}})circle(1.3pt);
 }\temp
}
\end{axis}
\node [right] at (myplot.right of origin) {\scriptsize $x$};
\node [above] at (myplot.above origin) {\scriptsize $y$};
\end{tikzpicture}
\\ (c)
}
\egroup

Let $Q$ be the area of a region bounded by continuous functions $f$ and $g$. If we break the region into many subregions, we have an obvious equation:
\begin{center}
Total Area = sum of the areas of the subregions.
\end{center}
The issue to address next is how to systematically break a region into subregions. A graph will help. Consider \autoref{fig:abcintro} (a) where a region between two curves is shaded. While there are many ways to break this into subregions, one particularly efficient way is to ``slice'' it vertically, as shown in \autoref{fig:abcintro} (b), into $n$ equally spaced slices.

We now approximate the area of a slice. Again, we have many options, but using a rectangle seems simplest. Picking any $x$-value $c_i$ in the $i^{\text{th}}$ slice, we set the height of the rectangle to be $f(c_i)-g(c_i)$, the difference of the corresponding $y$-values. The width of the rectangle is a small difference in $x$-values, which we represent with $\Delta x$. \autoref{fig:abcintro} (c) shows sample points $c_i$ chosen in each subinterval and appropriate rectangles drawn.
%(Each of these rectangles represents a differential element.)
Each slice has an area approximately equal to $\bigl(f(c_i)-g(c_i)\bigr)\Delta x$; hence, the total area is approximately the Riemann Sum\vspace{-.5\baselineskip}
\[Q \approx \sum_{i=1}^n \bigl(f(c_i)-g(c_i)\bigr)\Delta x.\]
Taking the limit as $n\to \infty$ gives the exact area as $\int_a^b \bigl(f(x)-g(x)\bigr)\dd x.$

\begin{theorem}[Area Between Curves]\label{thm:areabetweencurves}%
Let $f(x)$ and $g(x)$ be continuous functions defined on $[a,b]$ where $f(x)\geq g(x)$ for all $x$ in $[a,b]$. The area of the region bounded by the curves $y=f(x)$, $y=g(x)$ and the lines $x=a$ and $x=b$ is \index{integration!area between curves}
\[\int_a^b \bigl(f(x)-g(x)\bigr)\dd x.\]
\end{theorem}

Often, we do not know which function is greater (or they switch within the domain of integration).  If so, we can say that the area is $\int_a^b\abs{f(x)-g(x)}\dd x$, which may involve dividing the domain of integration into pieces.

\youtubeVideo{DRFyNHdVgUA}{Finding Areas Between Curves}

\begin{example}[Finding area enclosed by curves]\label{ex_abc1}%
Find the area of the region bounded by $f(x) = \sin x+2$, $g(x) = \frac12\cos (2x)-1$, $x=0$ and $x=4\pi$, as shown in \autoref{fig:abc1}.
\solution
\mtable{Graphing an enclosed region in \autoref{ex_abc1}.}{fig:abc1}{\begin{tikzpicture}[alt={Blue upper curve sin x + 2 and lower curve 0.5 cos 2x − 1; shaded area between curves from x 0 to 4π; axes x and y.}]
\begin{axis}[width=\marginparwidth,tick label style={font=\scriptsize},
axis y line=middle,axis x line=middle,name=myplot,axis on top,extra x ticks={12.57},
extra x tick labels={$4\pi$},ymin=-2,ymax=3.5,xmin=-.5,xmax=13.5]
\addplot[smooth,draw={\coloronefill},fill={\coloronefill},area style,domain=0:4*pi]
 {.5*cos(deg(2*x))-1.2}\closedcycle;
\addplot [smooth,draw={\coloronefill},fill={\coloronefill},area style,domain=0:4*pi]
 {sin(deg(x))+2} \closedcycle;
\addplot [smooth,thick, draw={\colorone},domain=0:4*pi,samples=40]
 {sin(deg(x))+2} node [shift={(-15pt,7pt)} ,black] {\scriptsize $f(x)$};
\addplot [smooth,thick, draw={\colorone},domain=0:4*pi,samples=40]
 {.5*cos(deg(2*x))-1.2}node [shift={(-25pt,-16pt)} ,black] {\scriptsize $g(x)$};
\end{axis}
\node [right] at (myplot.right of origin) {\scriptsize $x$};
\node [above] at (myplot.above origin) {\scriptsize $y$};
\end{tikzpicture}}
%
The graph verifies that the upper boundary of the region is given by $f$ and the lower bound is given by $g$. Therefore the area of the region is the value of the integral
\begin{align*} 
	\int_0^{4\pi} \bigl(f(x)- g(x)\bigr)\dd x
	& = \int_0^{4\pi} \Bigl(\sin x+2 - \bigl(\frac12\cos (2x)-1\bigr)\Bigr)\dd x \\
	&= -\cos x -\frac14\sin(2x)+3x\Big|_0^{4\pi}\\
	&=	12\pi %\approx 37.7
	\ \text{units}^2.
\end{align*}
\end{example}

% originally in FTC section
\begin{example}[Finding area between curves]\label{ex_ftc6}%
Find the area of the region enclosed by $y=x^2+x-5$ and $y=3x-2$.
\solution
It will help to sketch these two functions, as done in \autoref{fig:ftc6}. 
\mtable{Sketching the region enclosed by $y=x^2+x-5$ and $y=3x-2$ in \autoref{ex_ftc6}.}{fig:ftc6}{\begin{tikzpicture}[alt={Blue parabola y = x²+x−5 and blue line y=3x−2 intersect at −1 and 3; axes x and y.}]
\begin{axis}[width=\marginparwidth,tick label style={font=\scriptsize},
axis y line=middle,axis x line=middle,name=myplot,axis on top,xtick={-2,-1,1,2,3,4},
ytick={5,10,15},ymin=-10,ymax=16,xmin=-2.5,xmax=4.5]
\addplot [smooth,thick, draw={\colorone},domain=-2:4] {x^2+x-5}; 
\draw (axis cs: 2.5,15) node [black] {\scriptsize $y=x^2+x-5$};
\addplot [smooth,thick, draw={\colorone},domain=-2:4] {3*x-2}; 
\draw (axis cs: -1.5,-9) node [black] {\scriptsize $y=3x-2$};
\end{axis}
\node [right] at (myplot.right of origin) {\scriptsize $x$};
\node [above] at (myplot.above origin) {\scriptsize $y$};
\end{tikzpicture}}
The region whose area we seek is completely bounded by these two functions; they seem to intersect at $x=-1$ and $x=3$. To check, set $x^2+x-5=3x-2$ and solve for $x$:\vspace{-.5\baselineskip}
\begin{align*}
	x^2+x-5 &= 3x-2 \\
	(x^2+x-5) - (3x-2) &= 0\\
	x^2-2x-3 &= 0\\
	(x-3)(x+1) &= 0\\
	x&=-1,\ 3.
\end{align*}
Following \autoref{thm:areabetweencurves}, the area is 
\begin{align*}
	\int_{-1}^3\bigl(3x-2 -(x^2+x-5)\bigr)\dd x &= \int_{-1}^3 (-x^2+2x+3)\dd x \\
	&=\left.\left(-\frac13x^3+x^2+3x\right)\right|_{-1}^3 \\
	&=-\frac13(27)+9+9-\left(\frac13+1-3\right)\\
	&= 10\frac23 = 10.\overline{6}
\end{align*}
\end{example}

\mtable{Graphing a region enclosed by two functions in \autoref{ex_abc2}.}{fig:abc2}{\begin{tikzpicture}[alt={Blue cubic curve and blue line intersect at x = 1 and x = 4;; region between curves is shaded; axes x and y.}]
\begin{axis}[width=1.16\marginparwidth,tick label style={font=\scriptsize},
axis y line=middle,axis x line=middle,name=myplot,axis on top,xtick={1,2,3,4},
ymin=-5,ymax=3.5,xmin=-.5,xmax=4.7]
\addplot[smooth,draw={\coloronefill},fill={\coloronefill},area style,domain=1:2]
 {x^3-7*x^2+12*x-3}\closedcycle;
\addplot[smooth,draw=white,fill=white,area style,domain=1:2]
 {5-2*x}\closedcycle;
\addplot[smooth,draw={\coloronefill},fill={\coloronefill},area style,domain=2:2.5]
 {5-2*x}\closedcycle;
\addplot[smooth,draw=white,fill=white,area style,domain=2:2.2391]
 {x^3-7*x^2+12*x-3}\closedcycle;
\addplot[smooth,draw={\coloronefill},fill={\coloronefill},area style,domain=2.2391:4]
 {x^3-7*x^2+12*x-3}\closedcycle;
\addplot[smooth,draw=white,fill=white,area style,domain=2.5:4]
 {5-2*x}\closedcycle;
\addplot [smooth,thick, draw={\colorone},domain=-.2:4.5,samples=30] {x^3-7*x^2+12*x-3};
\addplot [smooth,thick, draw={\colorone},domain=0:4.5] {5-2*x};
\end{axis}
\node [right] at (myplot.right of origin) {\scriptsize $x$};
\node [above] at (myplot.above origin) {\scriptsize $y$};
\end{tikzpicture}}

\begin{example}[Finding total area enclosed by curves]\label{ex_abc2}%
Find the total area of the region enclosed by the functions $f(x) = -2x+5$ and $g(x) = x^3-7x^2+12x-3$ as shown in \autoref{fig:abc2}.
\solution
A quick calculation shows that $f=g$ at $x=1, 2$ and 4. One can proceed thoughtlessly by computing $\ds \int_1^4\bigl(f(x)-g(x)\bigr)\dd x$, but this ignores the fact that on $[1,2]$, $g(x)>f(x)$. (In fact, the thoughtless integration returns $-9/4$, hardly the expected value of an \emph{area}.) Thus we compute the total area by breaking the interval $[1,4]$ into two subintervals, $[1,2]$ and $[2,4]$ and using the proper integrand in each.
\begin{align*}
	\text{Total Area}
	&= \int_1^2 \bigl(g(x)-f(x)\bigr)\dd x + \int_2^4\bigl(f(x)-g(x)\bigr)\dd x\\
	&= \int_1^2 \bigl(x^3-7x^2+14x-8\bigr) \dd x
	+ \int_2^4\bigl(-x^3+7x^2-14x+8\bigr)\dd x\\
	&= \frac5{12} + \frac83
	 %\\&
	= \frac{37}{12} % \approx 3.083
	\ \text{units}^2.
\end{align*}
\end{example}

The previous example makes note that we are expecting area to be \emph{positive}. When first learning about the definite integral, we interpreted it as ``signed area under the curve,'' allowing for ``negative area.'' That doesn't apply here; area is to be positive.

The previous example also demonstrates that we often have to break a given region into subregions before applying \autoref{thm:areabetweencurves}. The following example shows another situation where this is applicable, along with an alternate view of applying the Theorem.

\mtable{Graphing a region for \autoref{ex_abc3}.}{fig:abc3}{\begin{tikzpicture}[alt={Blue region: upper curve y=2+√x on 0≤x≤1, then y=-(x-1)²+3 on 1≤x≤2; lower edge y=2; axes x and y.}]
\begin{axis}[width=1.16\marginparwidth,tick label style={font=\scriptsize},
axis y line=middle,axis x line=middle,name=myplot,axis on top,xtick={1,2},
ymin=-.1,ymax=3.5,xmin=-.1,xmax=2.5]
\addplot[draw={\coloronefill},fill={\coloronefill},area style,domain=0:1]{sqrt(x)+2}\closedcycle;
\addplot[draw={\coloronefill},fill={\coloronefill},area style,domain=1:2]{3-(x-1)^2}\closedcycle;
\draw[fill=white,draw=white](axis cs:0,0)rectangle(axis cs:2,2);
\addplot[draw={\colorone},domain=0:1]{sqrt(x)+2};
\addplot[draw={\colorone},domain=1:2]{3-(x-1)^2};
\draw[draw={\colorone}](axis cs:0,2)--(axis cs:2,2);
\draw (axis cs:.5,3.1) node {\scriptsize $y=\sqrt{x}+2$} (axis cs:1.8,3.2) node {\scriptsize $y=-(x-1)^2+3$};
\end{axis}
\node [right] at (myplot.right of origin) {\scriptsize $x$};
\node [above] at (myplot.above origin) {\scriptsize $y$};
\end{tikzpicture}}

\begin{example}[Finding area: integrating with respect to $y$]\label{ex_abc3}%
Find the area of the region enclosed by the functions $y=\sqrt{x}+2$, $y=-(x-1)^2+3$ and $y=2$, as shown in \autoref{fig:abc3}.
\solution
We give two approaches to this problem. In the first approach, we notice that the region's ``top'' is defined by two different curves. On $[0,1]$, the top function is $y=\sqrt{x}+2$; on $[1,2]$, the top function is $y=-(x-1)^2+3$. Thus we compute the area as the sum of two integrals:

\begin{align*}
	\text{Total Area}
	&= \int_0^1 \Bigl(\bigl(\sqrt{x}+2\bigr)-2\Bigr)\dd x + \int_1^2 \Bigl(\bigl(-(x-1)^2+3\bigr)-2\Bigr)\dd x \\
	&= 2/3 + 2/3
	%\\&
	=4/3.
\end{align*}

The second approach is clever and very useful in certain situations. We are used to viewing curves as functions of $x$; we input an $x$-value and a $y$-value is returned. Some curves can also be described as functions of $y$: input a $y$-value and an $x$-value is returned. We can rewrite the equations describing the boundary by solving for $x$:\vspace{-.3\baselineskip}
\begin{align*}
y=\sqrt{x}+2 & \quad\Rightarrow\quad x=(y-2)^2 \\
y=-(x-1)^2+3 & \quad\Rightarrow\quad x=\sqrt{3-y}+1.
\end{align*}

\mtable{The region used in \autoref{ex_abc3} with boundaries relabeled as functions of $y$.}{fig:abc3b}{\begin{tikzpicture}[alt={Horizontal view: shaded band between y=2 and y=3; left edge x=(y-2)², right edge x=1+√(3-y); a red horizontal slice shows Delta y; axes x and y.}]
\begin{axis}[width=1.16\marginparwidth,tick label style={font=\scriptsize},
axis y line=middle,axis x line=middle,name=myplot,axis on top,xtick={1,2},
ymin=-.1,ymax=3.5,xmin=-.1,xmax=2.5]
\addplot[draw={\coloronefill},fill={\coloronefill},area style,domain=0:1]{sqrt(x)+2}\closedcycle;
\addplot[draw={\coloronefill},fill={\coloronefill},area style,domain=1:2]{3-(x-1)^2}\closedcycle;
\draw[fill=white,draw=white](axis cs:0,0)rectangle(axis cs:2,2);
\addplot[draw={\colorone},domain=0:1]{sqrt(x)+2};
\addplot[draw={\colorone},domain=1:2]{3-(x-1)^2};
\draw[draw={\colorone}](axis cs:0,2)--(axis cs:2,2);
\draw (axis cs:.5,3.1) node {\scriptsize $x=(y-2)^2$} (axis cs:1.8,3.2) node {\scriptsize $x=\sqrt{3-y}+1$};
\filldraw [draw={\colortwo},fill={\colortwofill},thick] (axis cs: 0.0625,2.2) rectangle (axis cs:1.867,2.3);
\end{axis}
\node [right] at (myplot.right of origin) {\scriptsize $x$};
\node [above] at (myplot.above origin) {\scriptsize $y$};
\end{tikzpicture}}	

\autoref{fig:abc3b} shows the region with the boundaries relabeled. A
% differential element, a
horizontal rectangle is also pictured. The width of the rectangle is a small change in $y$: $\Delta y$. The height of the rectangle is a difference in $x$-values. The ``top'' $x$-value is the largest value, i.e., the rightmost. The ``bottom'' $x$-value is the smaller, i.e., the leftmost. Therefore the height of the rectangle is
\[\bigl(\sqrt{3-y}+1\bigr) - (y-2)^2.\]

The area is found by integrating the above function with respect to $y$ with the appropriate bounds. We determine these by considering the $y$-values the region occupies. It is bounded below by $y=2$, and bounded above by $y=3$. That is, both the ``top'' and ``bottom'' functions exist on the $y$ interval $[2,3]$. Thus
\begin{align*}
	\text{Total Area}
	&= \int_2^3 \bigl(\sqrt{3-y}+1 - (y-2)^2\bigr)\dd y \\
	&= \Bigl(-\frac23(3-y)^{3/2}+y-\frac13(y-2)^3\Bigr)\Big|_2^3 \\
	&= 4/3.
\end{align*}
The important thing to notice is that by integrating with respect to $y$ instead of $x$, we only had to do one integral and did not need to find the point at which to switch from one integration to another.
\end{example}

This calculus-based technique of finding area can be useful even with shapes that we normally think of as ``easy.'' \autoref{ex_abc4} computes the area of a triangle. While the formula ``$\frac12\times\text{base}\times\text{height}$'' is well known, in arbitrary triangles it can be nontrivial to compute the height. Calculus makes the problem simple.

\begin{example}[Finding the area of a triangle]\label{ex_abc4}%
Compute the area of the regions bounded by the lines
\mtable{Graphing a triangular region in \autoref{ex_abc4}.}{fig:abc4}{\begin{tikzpicture}[alt={Blue triangular region bounded by y = 3 − x, y = x + 1, and y = 5x − 15; vertices (1, 2), (3, 0), and (4, 5); axes x and y.}]
\begin{axis}[width=1.16\marginparwidth,axis equal,
	tick label style={font=\scriptsize},axis y line=middle,axis x line=middle,
	name=myplot,axis on top,%xtick={1,2,3},
	ymin=-.1,ymax=5.5,xmin=-.1,xmax=5.5]
\draw [draw={\colorone},thick,fill={\coloronefill}]
  (axis cs:3,0)
  -- node[below left,color=black] {$y=3-x$} (axis cs:1,2)
  -- node[above left,color=black] {$y=x+1$} (axis cs:4,5)
  -- node[pos=.7,below right,color=black] {$y=5x-15$} cycle;
\end{axis}
\node [right] at (myplot.right of origin) {\scriptsize $x$};
\node [above] at (myplot.above origin) {\scriptsize $y$};
\end{tikzpicture}}
$y=3-x$, $y=x+1$ and $y=5x-15$, as shown in \autoref{fig:abc4}.
\solution
Recognize that there are two ``bottom'' functions to this region, causing us to use two definite integrals.
\begin{align*}
	\text{Total Area}
	&= \int_1^3\bigl((x+1)-(3-x)\bigr)\dd x + \int_3^4\bigl((x+1)-(5x-15)\bigr)\dd x \\
%	&= \int_1^3 2x-2\dd x + \int_3^4 16-4x\dd x \\
%	&= \left.x^2-2x\right|_1^3 + \left.16x-2x^2\right|_3^4 \\
%	&= (3- -1) + (32-30) \\
	&= 4+2
	%\\&
	= 6.
\end{align*}
We can also approach this by converting each function into a function of $y$. This also requires 2 integrals, so there isn't really any advantage to doing so. We do it here for demonstration purposes.

The ``top'' function is always $x=\frac y5+3$ while there are two ``bottom'' functions: $x=3-y$ and $x=y-1$. Being mindful of the proper integration bounds, we have
\begin{align*}
	\text{Total Area}
	&= \int_0^2\left(\left(\frac y5+3\right) - (3-y)\right)\dd y
	+ \int_2^5\left(\left(\frac y5+3\right) - (y-1)\right)\dd y \\
%	&= \int_0^2\left(\frac{6y}5\right)\dd y+ \int_2^5\left(-\frac{4y}5+4\right)\dd y \\
%	&= \left.\frac{3y^2}5\right|_0^2 + \left[-\frac{2y^2}5+4y\right]_2^5 \\
%	&= \frac{12}5 + \left[10-\frac{32}5\right] \\
	&= \frac{12}5 + \frac{18}5
	% \\
%	&= \frac{30}5 \\
	%&
	= 6.
\end{align*}
Of course, the final answer is the same (and we see that integrating with respect to $x$ was probably easier, since it avoided fractions).
\end{example}

% move this to numerical integration?
%While we have focused on producing exact answers, we are also able to make approximations using the principle of \autoref{thm:areabetweencurves}. The integrand in the theorem is a distance (``top minus bottom''); integrating this distance function gives an area. By taking discrete measurements of distance, we can approximate an area using numerical integration techniques developed in \autoref{sec:numerical_integration}. The following example demonstrates this.
%
%\begin{example}[Numerically approximating area]\label{ex_abc5}%
%To approximate the area of a lake, shown in \autoref{fig:abc5}(a),  the ``length'' of the lake is measured at 200-foot increments as shown in \autoref{fig:abc5}(b), where the lengths are given in hundreds of feet. Approximate the area of the lake.
%\solution
%The measurements of length can be viewed as measuring ``top minus bottom'' of two functions. The exact answer is found by integrating $\ds \int_0^{12} \bigl(f(x)-g(x)\bigr)\dd x$, but of course we don't know the functions $f$ and $g$. Our discrete measurements instead allow us to approximate.
%
%\mtable{(a) A sketch of a lake, and (b) the lake with length measurements.}{fig:abc5}{\begin{tikzpicture}[alt={A blue lake seen from above.}]
%\begin{axis}[width=1.16\marginparwidth,tick label style={font=\scriptsize},
%axis y line=middle,axis x line=middle,name=myplot,axis on top,
%xtick={1,2,3,4,5,6,7,8,9,10,11,12},ytick={1,2,3,4,5,6,7,8},ymin=-.1,ymax=8.5,xmin=-.1,xmax=12.5]
%\addplot [draw={\colorone},thick,fill={\coloronefill},smooth] coordinates {(0,3.2)(0.5,3.417)(1.,3.637)(1.5,4.041)(2.,4.505)(2.5,5.)(3.,5.495)(3.5,5.959)(4.,6.363)(4.5,6.683)(5.,6.898)(5.5,6.995)(6.,6.968)(6.5,6.819)(7.,6.556)(7.5,6.197)(8.,5.763)(8.5,5.282)(9.,4.784)(9.5,4.298)(10.,3.857)(10.5,3.486)(11.,3.21)(11.5,2.8)(11.5,2)(11.,1.545)(10.5,1.314)(10.,1.102)(9.5,0.9154)(9.,0.7585)(8.5,0.6361)(8.,0.5514)(7.5,0.5069)(7.,0.5038)(6.5,0.5421)(6.,0.6208)(5.5,0.7378)(5.,0.8897)(4.5,1.072)(4.,1.281)(3.5,1.509)(3.,1.751)(2.5,2.)(2.,2.249)(1.5,2.491)(1.,2.719)(0.5,2.9102)(0,3.15)(0,3.2)};
%\end{axis}
%\end{tikzpicture}
%\\(a)\\
%\begin{tikzpicture}[alt={The same blue lake, now with width measurements from top to bottom at various points.}]
%\begin{axis}[width=1.16\marginparwidth,tick label style={font=\scriptsize},
%axis y line=middle,axis x line=middle,name=myplot,axis on top,
%xtick={1,2,3,4,5,6,7,8,9,10,11,12},ytick={1,2,3,4,5,6,7,8},
%ymin=-.1,ymax=8.5,xmin=-.1,xmax=12.5]
%\addplot [draw={\colorone},thick,fill={\coloronefill},smooth] coordinates {(0,3.2)(0.5,3.417)(1.,3.637)(1.5,4.041)(2.,4.505)(2.5,5.)(3.,5.495)(3.5,5.959)(4.,6.363)(4.5,6.683)(5.,6.898)(5.5,6.995)(6.,6.968)(6.5,6.819)(7.,6.556)(7.5,6.197)(8.,5.763)(8.5,5.282)(9.,4.784)(9.5,4.298)(10.,3.857)(10.5,3.486)(11.,3.21)(11.5,2.8)(11.5,2)(11.,1.545)(10.5,1.314)(10.,1.102)(9.5,0.9154)(9.,0.7585)(8.5,0.6361)(8.,0.5514)(7.5,0.5069)(7.,0.5038)(6.5,0.5421)(6.,0.6208)(5.5,0.7378)(5.,0.8897)(4.5,1.072)(4.,1.281)(3.5,1.509)(3.,1.751)(2.5,2.)(2.,2.249)(1.5,2.491)(1.,2.719)(0.5,2.9102)(0,3.15)(0,3.2)};
%\draw (axis cs:2,4.5) -- (axis cs:2,2.249) node [shift={(-3pt,0pt)},rotate=90,pos=.5] {\scriptsize 2.25}
%(axis cs:4,6.36) -- (axis cs:4,1.28) node [shift={(-3pt,0pt)},rotate=90,pos=.5] {\scriptsize 5.08}
%(axis cs:6,6.97) -- (axis cs:6,0.62) node [shift={(-3pt,0pt)},rotate=90,pos=.5] {\scriptsize 6.35}
%(axis cs:8,5.76) -- (axis cs:8,.55) node  [shift={(-3pt,0pt)},rotate=90,pos=.5] {\scriptsize 5.21}
%(axis cs:10,3.86) -- (axis cs:10,1.1) node [shift={(-3pt,0pt)},rotate=90,pos=.5] {\scriptsize 2.76};
%\end{axis}
%\node [right] at (myplot.right of origin) {\scriptsize $x$};
%\node [above] at (myplot.above origin) {\scriptsize $y$};
%\end{tikzpicture}
%\\ (b)}
%We have the following data points:
%\[(0,0),\ (2,2.25),\ (4,5.08),\ (6,6.35),\ (8,5.21),\ (10,2.76),\ (12,0).\]
%We also have that $\Delta x=\frac{b-a}{n} = 2$, so Simpson's Rule gives
%\begin{align*}
%	\text{Area}
%	&\approx \frac{2}{3}
%\Bigl(1\cdot0+4\cdot2.25+2\cdot5.08+4\cdot6.35+2\cdot5.21+4\cdot2.76+1\cdot0\Bigr)\\
%	&= 44.01\overline{3} \ \text{units}^2.
%\end{align*}
%
%Since the measurements are in hundreds of feet, units${}^2 = (100\ \text{ft})^2 = 10,000\ \text{ft}^2$, giving a total area of $440,133\ \text{ft}^2$. (Since we are approximating, we'd likely say the area was about $440,000\ \text{ft}^2$, which is a little more than 10 acres.)
%\end{example}

In the next section we apply \autoref{idea:app_of_defint} to finding the volumes of certain solids.

\printexercises{exercises/07-01-exercises}

\section[Volume by Cross-Sectional Area; Disk and Washer Methods]{Volume by Cross-Sectional Area;\\Disk and Washer Methods}\label{sec:disk}

The volume of a general right cylinder, as shown in \autoref{fig:cross1}, is 
\begin{center}
Area of the base $\times$ height.
\end{center}
\mtable{The volume of a general right cylinder is the product of its height and its base's area}{fig:cross1}{
\myincludeasythree{
3Droll=5.855410611764742,
3Dortho=0.00922776386141777,
3Dc2c=0.38511115312576294 0.8229191303253174 0.41772428154945374,
3Dcoo=9.545748710632324 -13.143834114074707 1.8314944505691528,
3Droo=150}{.5\marginparwidth}{Irregular right cylinder: matching top and bottom outlines, vertical side walls, base area A, height h.}{figures/figcross1_3D}}
We can use this fact as the building block in finding volumes of a variety of shapes.

Given an arbitrary solid, we can \emph{approximate} its volume by cutting it into $n$  thin slices. When the slices are thin, each slice can be approximated well by a general right cylinder. Thus the volume of each slice is approximately its cross-sectional area $\times$ thickness.
% (These slices are the differential elements.)

By orienting a solid along the $x$-axis, we can let $A(x_i)$ represent the cross-sectional area
of the $i\,^\text{th}$ slice, and let $\Delta x$ represent the thickness of the slices (the thickness is a small change in $x$). The total volume of the solid is approximately:
	\begin{align*} \text{Volume} &\approx \sum_{i=1}^n \Bigl[\text{Area}\ \times\ \text{thickness}\Bigr] \\
			&= \sum_{i=1}^n A(x_i)\Delta x.
	\end{align*}
	
Recognize that this is a Riemann Sum. By taking a limit (as the thickness of the slices goes to 0) we can find the volume exactly. 
\[\text{Volume}=\lim_{n\to \infty} \sum_{i=1}^n A(x_i)\Delta x\]
with $\Delta x=\frac{b-a}{n}$ and $x_i=a+i\Delta x$. We recognize this as a definite integral.

\begin{theorem}[Volume By Cross-Sectional Area]\label{thm:volume_by_cross_section}%
The volume $V$ of a solid, oriented along the $x$-axis with cross-sectional area $A(x)$ from $x=a$ to $x=b$, is \index{integration!volume!cross-sectional area}
\[V = \int_a^b A(x)\dd x.\]
\end{theorem}

\begin{example}[Finding the volume of a solid]\label{ex_disk0}%
Find the volume of a pyramid with a square base of side length 10 in and a height of 5 in.
\solution
There are many ways to ``orient'' the pyramid along the $x$-axis; \autoref{fig:disk0}(a) gives one such way, with the pointed top of the pyramid at the origin and the $x$-axis going through the center of the base.
\mtable{Orienting a pyramid along the $x$-axis (top) and cutting a slice in the pyramid (bottom) in \autoref{ex_disk0}.}{fig:disk0}{%
\myincludeasythree{
3Droll=130.3044775136707,
3Dortho=0.004594136495143175,
3Dc2c=0.4693154990673065 0.5382519960403442 0.7000198364257812,
3Dcoo=60.78499984741211 0.48087620735168457 -14.641631126403809,
3Droo=150.00000157638937}{\marginparwidth}{Pyramid on x axis, apex at origin, square base at x=5 of side length 10; cross section square side length 2x; axes x and y.}{figures/figcross_area1_3D}
\\(a)\\
\myincludeasythree{
3Droll=130.3044775136707,
3Dortho=0.004594136495143175,
3Dc2c=0.4693154990673065 0.5382519960403442 0.7000198364257812,
3Dcoo=60.78499984741211 0.48087620735168457 -14.641631126403809,
3Droo=150.00000157638937}{\marginparwidth}{Slice of pyramid at x showing square cross section of side length 2x and thickness Delta x, with dashed outline of full pyramid; axes x and y.}{figures/figcross_area1a_3D}
\\(b)}

Each cross section of the pyramid is a square.
%; this is a sample differential element.
To determine its area $A(x)$, we need to determine the side lengths of the square.

When $x=5$, the square has side length 10; when $x=0$, the square has side length 0. Since the edges of the pyramid are lines, it is easy to figure that each cross-sectional square has side length $2x$, giving $A(x) = (2x)^2=4x^2$.

If one were to cut a slice out of the pyramid at $x=3$, as shown in \autoref{fig:disk0}(b), one would have a shape with square bottom and top with sloped sides. If the slice were thin, both the bottom and top squares would have sides lengths of about 6, and thus the cross-sectional area of the bottom and top would be about 36in$^2$. Letting $\Delta x$ represent the thickness of the slice, the volume of this slice would then be about $36\Delta x$in$^3$. 

Cutting the pyramid into $n$ slices divides the total volume into $n$ equally-spaced smaller pieces, each with volume $(2x_i)^2\Delta x$, where $x_i$ is the approximate location of the slice along the $x$-axis and $\Delta x$ represents the thickness of each slice. One can approximate total volume of the pyramid by summing up the volumes of these slices:
\[\text{Volume } \approx \sum_{i=1}^n (2x_i)^2\Delta x.\]
Taking the limit as $n\to\infty$ gives the actual volume of the pyramid; recognizing this sum as a Riemann Sum allows us to find the exact answer using a definite integral, matching the definite integral given by \autoref{thm:volume_by_cross_section}.

We have 
\begin{align*}
	V
	&= \lim_{n\to\infty} \sum_{i=1}^n (2x_i)^2\Delta x\\
	&= \int_0^5 4x^2\dd x\\
	&= \frac43x^3\Big|_0^5 \\
	&=\frac{500}{3}\ %\text{in}^3 \approx 166.67
	\ \text{in}^3.
\end{align*}
We can check our work by consulting the general equation for the volume of a pyramid (see the back cover under ``Volume of A General Cone''): 
\begin{center}
$\frac13\times \text{area of base}\times \text{height}$.
\end{center}
Certainly, using this formula from geometry is faster than our new method, but the calculus-based method can be applied to much more than just cones.
\end{example}

An important special case of \autoref{thm:volume_by_cross_section} is when the solid is a \textbf{solid of revolution}, that is, when the solid is formed by rotating a shape around an axis.

Start with a function $y=f(x)$ from $x=a$ to $x=b$. Revolving this curve about a horizontal axis encloses a three-dimensional solid whose cross sections are disks (thin circles), perpendicular to the axis of rotation. Let $R(x)$ represent the radius of the cross-sectional disk at $x$; the area of this disk is $\pi [R(x)]^2$. Applying \autoref{thm:volume_by_cross_section} gives the Disk Method.

\begin{keyidea}[The Disk Method]\label{idea:disk_method}%
Let a solid be enclosed by revolving the curve $y=f(x)$ from $x=a$ to $x=b$ around a horizontal axis, and let $R(x)$ be the radius of the cross-sectional disk at $x$. The volume of the solid is
\index{integration!volume!Disk Method}\index{Disk Method}
\[V = \pi \int_a^b [R(x)]^2\dd x.\]
\end{keyidea}

\youtubeVideo{nZqOKc067Z8}{Longer Version --- Volumes using Disks/Washers}

\mtable[2.5in]{Sketching a solid in \autoref{ex_disk1}.}{fig:disk1}{%
  \begin{tikzpicture}[alt={Graph of y = 1/x from x = 1 to 2; region under curve shaded; vertical line at x shows radius R(x); axes x and y.}]
   \begin{axis}[width=\marginparwidth,
     tick label style={font=\scriptsize},axis y line=middle,
     axis x line=middle,name=myplot,axis on top,
     ymin=-.2,ymax=1.2,xmin=-.1,xmax=2.2]
    \draw[draw={\colorone},fill={\coloronefill}] plot[smooth,domain=1:2]
     (axis cs:\x,1/\x) |- (axis cs:1,0) --(axis cs:2,0)--(axis cs:1,0)-- cycle;
    \draw [thick,draw={\colortwo}] (axis cs: 1.5,0) --
      node[pos=.5,right,color=black]{\scriptsize$R(x)$} (axis cs: 1.5,.667);
   \end{axis}
   \node [right] at (myplot.right of origin) {\scriptsize $x$};
   \node [above] at (myplot.above origin) {\scriptsize $y$};
  \end{tikzpicture}
  \\ (a) \\
	\myincludeasythree{
3Droll=126.0060482236849,
3Dortho=0.004875946324318647,
3Dc2c=0.44462037086486816 0.39410829544067383 0.8043577671051025,
3Dcoo=67.21983337402344 3.5258262157440186 -29.566415786743164,
3Droo=150.00000160608386}{\marginparwidth}{Red disk at sample x with radius equal to the curve height 1/x; curve drawn; axes x and y.}{figures/figdisk1_3D}\\
	(b) \\
	\myincludeasythree{
3Droll=126.0060482236849,
3Dortho=0.004875946324318647,
3Dc2c=0.44462037086486816 0.39410829544067383 0.8043577671051025,
3Dcoo=67.21983337402344 3.5258262157440186 -29.566415786743164,
3Droo=150.00000160608386}{\marginparwidth}{3D view of solid formed by rotating shaded region under 1/x from 1 to 2 around x axis; sample disk shown; axes x and y.}{figures/figdisk1b_3D}\\
	(c)}

\begin{example}[Finding volume using the Disk Method]\label{ex_disk1}%
Find the volume of the solid formed by revolving about the $x$-axis the region bounded by the curves $y=1/x$, $x=1$, $x=2$ and the $x$-axis.
\solution
A sketch is always useful in helping us understand the problem. In \autoref{fig:disk1}(a) we have sketched the region we will be rotating.  In \autoref{fig:disk1}(b), the curve $y=1/x$ is sketched along with the sample slice, a disk, at $x$ with radius $R(x)=1/x$. In \autoref{fig:disk1}(c) the whole solid is pictured, along with the sample slice.

The volume of the sample slice shown in part (b) of the figure is approximately $\pi R(x_i)^2\Delta x$, where $R(x_i)$ is the radius of the disk shown and $\Delta x$ is the thickness of that slice. The radius $R(x_i)$ is the distance from the $x$-axis to the curve, hence $R(x_i) = 1/x_i$.

Slicing the solid into $n$ equally-spaced slices, we can approximate the total volume by adding up the approximate volume of each slice:
\[\text{Approximate volume } = \sum_{i=1}^n \pi \left(\frac1{x_i}\right)^2\Delta x.\]

Taking the limit of the above sum as $n\to\infty$ gives the actual volume; recognizing this sum as a Riemann sum allows us to evaluate the limit with a definite integral, which matches the formula given in \autoref{idea:disk_method}:

%Using \autoref{idea:disk_method}, we have 
\begin{align*}
	V &= \lim_{n\to\infty}\sum_{i=1}^n \pi \left(\frac1{x_i}\right)^2\Delta x\\
		&= \pi\int_1^2 \left(\frac1x\right)^2\dd x \\
		&= \pi\int_1^2 \frac1{x^2}\dd x \\
		&= \pi\left[-\frac1x\right]\Big|_1^2 \\
		&= \pi \left[-\frac12 - \left(-1\right)\right] \\
		&= \frac{\pi}{2}\ \text{units}^3.
\end{align*}
\end{example}

While \autoref{idea:disk_method} is given in terms of functions of $x$, the principle involved can be applied to functions of $y$ when the axis of rotation is vertical, not horizontal. We demonstrate this in the next example.

\mtable[-3in]{Sketching a solid in \autoref{ex_disk2}.}{fig:disk2}{%
  \begin{tikzpicture}[alt={Shaded region between y=0.5 and 1 bounded by y-axis and curve x=1/y; red line labeled R(y); axes x and y.}]
   \begin{axis}[width=1.16\marginparwidth,
     tick label style={font=\scriptsize},axis y line=middle,
     axis x line=middle,name=myplot,axis on top,
     ymin=-.2,ymax=1.2,xmin=-.1,xmax=2.2]
    \draw[draw={\colorone},fill={\coloronefill}] plot[smooth,domain=1:2]
     (axis cs:\x,1/\x) -| (axis cs:0,.5) --(axis cs:0,.5)--(axis cs:0,1)-- cycle;
    \draw [thick,draw={\colortwo}] (axis cs: 0,.75) --
     node[pos=.5,below,color=black]{\scriptsize$R(y)$} (axis cs: 1.333,.75);
   \end{axis}
   \node [right] at (myplot.right of origin) {\scriptsize $x$};
   \node [above] at (myplot.above origin) {\scriptsize $y$};
  \end{tikzpicture}
  \\ (a) \\
	\myincludeasythree{
3Droll=112.9954735907269,
3Dortho=0.003950905054807663,
3Dc2c=0.5512889623641968 0.2862306833267212 0.7836788296699524,
3Dcoo=-9.51923942565918 49.542564392089844 -9.694497108459473,
3Droo=150.00000339216956}{\marginparwidth}{Red disk at sample y around the y-axis, radius R(y)=1/y labeled; blue curve x=1/y drawn; axes x, y.}{figures/figdisk1a_3D}
\\(b)\\
	\myincludeasythree{
3Droll=112.9954735907269,
3Dortho=0.003950905054807663,
3Dc2c=0.5512889623641968 0.2862306833267212 0.7836788296699524,
3Dcoo=-9.51923942565918 49.542564392089844 -9.694497108459473,
3Droo=150.00000339216956}{\marginparwidth}{3D solid from rotating region between y-axis and curve x = 1/y around y-axis; blue translucent outer, red inner circle; axes x, y.}{figures/figdisk2a_3D}
\\(c)}

\begin{example}[Finding volume using the Disk Method]\label{ex_disk2}%
Find the volume of the solid formed by revolving about the $y$-axis the region bounded by the curves $y=1/x$, $y=1$, $y=0.5$, and the $y$-axis.
\solution
Since the axis of rotation is vertical, our perpendicular cross sections have thickness $\Delta y$ and radius $x=R(y)$. We need to convert the function into a function of $y$. % and convert the $x$-bounds to $y$-bounds.
Since $y=1/x$ defines the curve, we rewrite it as $x=1/y$.% The bound $x=1$ corresponds to the $y$-bound $y=1$, and the bound $x=2$ corresponds to the $y$-bound $y=1/2$. 

Thus we are rotating about the $y$-axis the region bounded by the curves $x=1/y$, $y=1/2$, $y=1$, and the $y$-axis to form a solid. The region of revolution is sketched in \autoref{fig:disk2}(a), the curve and sample sample disk are sketched in \autoref{fig:disk2}(b), and a full sketch of the solid is in \autoref{fig:disk2}(b). We integrate to find the volume:
\begin{align*}
	V &= \pi\int_{1/2}^1 \frac{1}{y^2}\dd y \\
	&= -\frac{\pi}y\Big|_{1/2}^1 \\
	&= \pi\ \text{units}^3.
\end{align*}
\end{example}

% todo create and work out an example where the same region is rotated about the x axis and y axis but the disk method (for both) shows that they have different volumes

% not the same region
%The previous two examples demonstrate how taking the same region and rotating it about two different axes will result in different solids and thus volumes.

We can also compute the volume of solids of revolution that have a hole in the center. The general principle is simple: compute the volume of the solid irrespective of the hole, then subtract the volume of the hole. If the outside radius of the solid is $R(x)$ and the inside radius (defining the hole) is $r(x)$, then the volume is 
\[
V = \pi\int_a^b [R(x)]^2 \dd x - \pi\int_a^b [r(x)]^2\dd x
= \pi\int_a^b \left([R(x)]^2-[r(x)]^2\right)\dd x.
\]

One can generate a solid of revolution with a hole in the middle by revolving a region about an axis. Consider \autoref{fig:washeridea}(a), where a region is sketched along with a dashed, horizontal axis of rotation. By rotating the region about the axis, a solid is formed. Each cross section of this solid will be a washer (a disk with a hole in the center) as sketched in \autoref{fig:washeridea}(b). The outside of the washer has radius $R(x)$, whereas the inside has radius $r(x)$. The entire solid is sketched in \autoref{fig:washeridea}(c).  This leads us to the Washer Method.

%\mtable{fig:washeridea}{Establishing the Washer Method.}{%
\noindent\begin{minipage}[t]{\linewidth}\noindent%
\captionsetup{type=figure}%
\flushinner{%
\tagpdfsetup{table/header-rows={2}}
\begin{tabular}{ c c c }
 \begin{tikzpicture}[alt={Blue curves R(x) (upper) and r(x) (lower) from x=1 to x=3; shaded region between them; axes x, y.}]
  \begin{axis}[width=1.16\marginparwidth,
     tick label style={font=\scriptsize},axis y line=middle,
     axis x line=middle,name=myplot,axis on top,
     ymin=-3.5,ymax=3.5,xmin=-.2,xmax=3.5]
    \draw[draw={\colorone},fill={\coloronefill}]
     (axis cs:1,2.5) parabola bend (axis cs:2,3) (axis cs:3,2.5) --
     (axis cs:3,2) parabola bend (axis cs:2,1) (axis cs:1,2) -- cycle;
    \draw [dashed] (axis cs:-.2,.5) -- (axis cs:3.5,.5);
    \draw [draw={\colortwo}] (axis cs:1.2,2.68) --
     node [pos=.6,left,color=black] {\scriptsize $R(x)$} (axis cs:1.2,.5);
    \draw [draw={\colortwo}] (axis cs: 2.8,1.64) --
     node [pos=.5,right,color=black] {\scriptsize $r(x)$} (axis cs: 2.8,.5);
    \filldraw (axis cs:1.2,2.68) circle (1pt) (axis cs:1.2,.5) circle (1pt)
     (axis cs: 2.8,1.64)circle (1pt)  (axis cs: 2.8,.5) circle (1pt)  ;
  \end{axis}
  \node [right] at (myplot.right of origin) {\scriptsize $x$};
  \node [above] at (myplot.above origin) {\scriptsize $y$};
 \end{tikzpicture}
%\\(a)\\
&
\myincludeasythree{
3Droll=96.94265756936434,
3Dortho=0.005309578496962786,
3Dc2c=0.547386884689331 0.07732175290584564 0.833299994468689,
3Dcoo=64.78053283691406 10.129053115844727 -47.566162109375,
3Droo=149.99999789136484}{\marginparwidth}{Blue curves R(x) and r(x) from a to b; red washer ellipse at sample x shows outer R(x) and inner r(x); axes x,y.}{figures/figwasher_idea_c_3D}
%\\(b)\\
&
\myincludeasythree{
3Droll=96.94265756936434,
3Dortho=0.005309578496962786,
3Dc2c=0.547386884689331 0.07732175290584564 0.833299994468689,
3Dcoo=64.78053283691406 10.129053115844727 -47.566162109375,
3Droo=149.99999789136484}{\marginparwidth}{3D hollow solid from rotating region between curves around x axis; blue outer, red inner shells; axes x and y.}{figures/figwasher_idea_b_3D}
\\
(a) & (b) &
(c)%}
\end{tabular}}
\caption{Establishing the Washer Method.}
\label{fig:washeridea}
\end{minipage}

\begin{keyidea}[The Washer Method]\label{idea:washermethod}%
Let a region bounded by $y=f(x)$, $y=g(x)$, $x=a$ and $x=b$ be rotated about a horizontal axis that does not intersect the region, forming a solid. Each cross section at $x$ will be a washer with outside radius $R(x)$ and inside radius $r(x)$. The volume of the solid is
\index{integration!volume!Washer Method}\index{Washer Method}
\[V = \pi\int_a^b \Bigl([R(x)]^2-[r(x)]^2\Bigr)\dd x.\]
\end{keyidea}

Even though we introduced it first, the Disk Method is just a special case of the Washer Method with an inside radius of $r(x)=0$.

%\mtable[-3\baselineskip]{Sketching the region, a sample slice, and solid in \autoref{ex_wash1}.}{fig:wash1}{%
\noindent\begin{minipage}[t]{\linewidth}\noindent%
\captionsetup{type=figure}%
\flushinner{%
\tagpdfsetup{table/header-rows={2}}
\begin{tabular}{ c c c }
 \begin{tikzpicture}[alt={Blue curves R(x) (upper) and r(x) (lower) from x=1 to 3; shaded region between them; axes x, y.}]
  \begin{axis}[width=\marginparwidth,
     tick label style={font=\scriptsize},axis y line=middle,
     axis x line=middle,name=myplot,axis on top,
     ymin=-.2,ymax=5.5,xmin=-.2,xmax=3.5]
    \draw[draw={\colorone},fill={\coloronefill}]
     (axis cs:1,1) parabola bend (axis cs:1,1) (axis cs:3,5) -- cycle;
    \draw [draw={\colortwo}] (axis cs:2.5,0) --
     node [pos=.4,right,color=black] {\scriptsize $R(x)$} (axis cs:2.5,4);
    \draw [draw={\colortwo}] (axis cs: 1.5,0) --
     node [pos=.5,right,color=black] {\scriptsize $r(x)$} (axis cs: 1.5,1.25);
  \end{axis}
  \node [right] at (myplot.right of origin) {\scriptsize $x$};
  \node [above] at (myplot.above origin) {\scriptsize $y$};
 \end{tikzpicture}
%\\(a)\\
&
\myincludeasythree{
3Droll=97.32968340849395,
3Dortho=0.004881054162979126,
3Dc2c=0.43117064237594604 0.05905330181121826 0.9003357887268066,
3Dcoo=64.78053283691406 10.12905216217041 -47.566158294677734,
3Droo=149.99999769784193}{\marginparwidth}{Washer slice at x=2: red disk with hole, blue outer radius to upper curve and red inner radius to lower curve; axes x, y.}{figures/figwash1c_3D}
%\\(b)\\
&
\myincludeasythree{
3Droll=97.32968340849395,
3Dortho=0.004881054162979126,
3Dc2c=0.43117064237594604 0.05905330181121826 0.9003357887268066,
3Dcoo=64.78053283691406 10.12905216217041 -47.566158294677734,
3Droo=149.99999769784193}{\marginparwidth}{3D hollow solid from rotating the region between curves around x axis; blue and red translucent surfaces; axes x, y.}{figures/figwash1b_3D}
%\\(c)}
\\(a)&(b)&(c)
\end{tabular}}
\caption{Sketching the region, a sample slice, and solid in \autoref{ex_wash1}.}
\label{fig:wash1}
\end{minipage}

\begin{example}[Finding volume with the Washer Method]\label{ex_wash1}%
Find the volume of the solid formed by rotating the region bounded by $y=x^2-2x+2$ and $y=2x-1$ about the $x$-axis.
\solution
A sketch of the region will help, as given in \autoref{fig:wash1}(a). Rotating about the $x$-axis will produce cross sections in the shape of washers, as shown in \autoref{fig:wash1}(b); the complete solid is shown in part (c). The outside radius of this washer is $R(x) = 2x+1$; the inside radius is $r(x) = x^2-2x+2$. As the region is bounded from $x=1$ to $x=3$, we integrate as follows to compute the volume.
\begin{align*}
V &= \pi\int_1^3 \Bigl((2x-1)^2-(x^2-2x+2)^2\Bigr)\dd x \\
		&= \pi\int_1^3 \bigl(-x^4+4x^3-4x^2+4x-3\bigr)\dd x \\
		&= \pi\Bigl[-\frac{1}{5}x^5+x^4-\frac43x^3+2x^2-3x\Bigr]\Big|_1^3 \\
		&=\frac{104}{15}\pi %\approx 21.78
		\ \text{units}^3.
\end{align*}
\end{example}

When rotating about a vertical axis, the outside and inside radius functions must be functions of $y$.

%\mtable{Sketching the region, a sample slice, and the solid in \autoref{ex_wash_3}.}{fig:ex_sq_and_rt}{%
\noindent\begin{minipage}[t]{\linewidth}\noindent%
\captionsetup{type=figure}%
\flushinner{%
\tagpdfsetup{table/header-rows={2}}
\begin{tabular}{ c c c }
 \begin{tikzpicture}[alt={Blue curves y=x² (lower) and y=2x (upper) from x=0 to 2; shaded region between curves; red horizontal sample lines at r(x) and R(x); axes x,y.}]
  \begin{axis}[width=\marginparwidth,tick label style={font=\scriptsize},
    axis equal,axis y line=middle,axis x line=middle,name=myplot,axis on top,
    xtick=\empty,extra x ticks={1},ytick={1},ymin=-.2,ymax=1.2,xmin=-.1,xmax=1.2]
   \addplot [draw={\coloronefill},fill={\coloronefill},domain=0:1] {sqrt(x)};
   \addplot [draw={\coloronefill},fill={\coloronefill},domain=0:1] {x^2};
   \addplot [draw={\colorone},smooth,thick,domain=0:1] {x^2};
   \addplot [draw={\colorone},smooth,thick,domain=0:1,samples=100] {sqrt x};
   \draw [thick,draw={\colortwo}] (axis cs: 0,.81) --
    node[pos=.4,above,color=black]{\scriptsize $R(x)$}(axis cs:.9,.81);
   \draw [thick,draw={\colortwo}] (axis cs: 0,.5) --
    node[pos=.5,above,color=black]{\scriptsize $r(x)$} (axis cs:.25,.5);
  \end{axis}
  \node [right] at (myplot.right of origin) {\scriptsize $x$};
  \node [above] at (myplot.above origin) {\scriptsize $y$};
 \end{tikzpicture}
% \\ (a) \\
&
\myincludeasythree{
3Droll=126.0060482236849,
3Dortho=0.004875946324318647,
3Dc2c=0.44462037086486816 0.39410829544067383 0.8043577671051025,
3Dcoo=3.631758213043213 62.91783142089844 -23.517351150512695,
3Droo=150}{\marginparwidth}{Red horizontal washer slice showing outer radius to y = 2x and inner radius to y = x^2; blue rotated surface; axes x,y.}{figures/figsq_rt_3D}
% \\ (b) \\
&
\myincludeasythree{
3Droll=126.0060482236849,
3Dortho=0.004875946324318647,
3Dc2c=0.44462037086486816 0.39410829544067383 0.8043577671051025,
3Dcoo=3.631758213043213 62.91783142089844 -23.517351150512695,
3Droo=150}{\marginparwidth}{3D hollow solid formed by rotating region between y = x^2 and y = 2x about the y-axis; blue outer and red inner shells; axes x,y.}{figures/figsq_rt_b_3D}
\\(a)&(b)&(c)
\end{tabular}}
% \\ (c)}
\caption{Sketching the region, a sample slice, and the solid in \autoref{ex_wash_3}.}
\label{fig:ex_sq_and_rt}
\end{minipage}

\begin{example}[Finding volume with the Washer Method]\label{ex_wash_3}%
Find the volume of the solid formed by rotating the region bounded by $y=x^2$ and $x=y^2$ about the $y$-axis.
\solution
In \autoref{fig:ex_sq_and_rt} we have a sketch of the region (a), a sample slice (b), and the solid (c). Rotating about the $y$-axis will produce cross sections in the shape of washers, as shown in (b); the complete solid is shown in  part (c). Since the axis of rotation is vertical, each radius is a function of $y$. The outside radius of this washer is $R(y)=\sqrt y$ and the inside radius is $r(y)=y^2$. As the region is bounded from $y=0$ to $y=1$, we integrate as follows to compute the volume. 
\begin{align*}
V&=\pi \int_0^1 \left((\sqrt y)^2-(y^2)^2\right) \dd y\\
&=\pi \int_0^1 y-y^4 \dd y\\
&=\pi\left.\left[\frac{1}{2} y^2-\frac{1}{5} y^5\right] \right|_0^1\\
&=\frac{3\pi}{10} \text{units}^3.
\end{align*}
\end{example}

%\mtable{Sketching the region, a sample slice, and the solid in \autoref{ex_wash2}.}{fig:wash2}{%
\noindent\begin{minipage}[t]{\linewidth}\noindent%
\captionsetup{type=figure}%
\flushinner{%
\tagpdfsetup{table/header-rows={2}}
\begin{tabular}{ c c c }
 \begin{tikzpicture}[alt={Blue curves R(x) upper and r(x) lower from x=a to b; shaded region between red lines; axes x, y.}]
  \begin{axis}[width=1.16\marginparwidth,axis equal,
    tick label style={font=\scriptsize},
    axis y line=middle,axis x line=middle,axis lines=center,name=myplot,
    xtick={1,2},ztick={1,2,3},ymin=-.2,ymax=3.2,xmin=-.2,xmax=2.2]
    \draw[draw={\colorone},fill={\coloronefill}] (axis cs:1,1) -- (axis cs:2,1) --(axis cs:2,3)-- cycle;
    \draw [thick,draw={\colortwo}] (axis cs: 0,2.6) --
      node[pos=.4,below,color=black]{\scriptsize$r(x)$} (axis cs: 1.8,2.6);
    \draw [thick,draw={\colortwo}] (axis cs: 0,1.2) --
      node[pos=.3,above,color=black]{\scriptsize$R(x)$} (axis cs: 2,1.2);
   \end{axis}
   \node [right] at (myplot.right of origin) {\scriptsize $x$};
   \node [above] at (myplot.above origin) {\scriptsize $y$};
  \end{tikzpicture}
%  \\(a) \\
&
\myincludeasythree{
3Droll=123.42078064170005,
3Dortho=0.00488104997202754,
3Dc2c=0.41855040192604065 0.31316813826560974 0.8524911999702454,
3Dcoo=-7.115118026733398 66.85215759277344 -16.372848510742188,
3Droo=149.99999819286032}{\marginparwidth}{Blue shaded region between the two curves; red washer slice at sample x shows outer and inner radii; axes x, y.}{figures/figwash2b_3D}
%\\(b) \\
&
\myincludeasythree{
3Droll=123.42078064170005,
3Dortho=0.00488104997202754,
3Dc2c=0.41855040192604065 0.31316813826560974 0.8524911999702454,
3Dcoo=-7.115118026733398 66.85215759277344 -16.372848510742188,
3Droo=149.99999819286032}{\marginparwidth}{3D hollow solid: blue outer shell and red inner shell from rotating region between two curves; axes x, y.}
{figures/figwash2c_3D}
\\(a)&(b)&(c)
\end{tabular}}
%\\(c)}
\caption{Sketching the region, a sample slice, and the solid in \autoref{ex_wash2}.}
\label{fig:wash2}
\end{minipage}

\begin{example}[Finding volume with the Washer Method]\label{ex_wash2}%
Find the volume of the solid formed by rotating the triangular region with vertices at $(1,1)$, $(2,1)$ and $(2,3)$ about the $y$-axis.
\solution
The triangular region is sketched in \autoref{fig:wash2}(a); the sample slice is sketched in (b) and the full solid is drawn in (c). They help us establish the outside and inside radii. Since the axis of rotation is vertical, each radius is a function of $y$. 

The outside radius $R(y)$ is formed by the line connecting $(2,1)$ and $(2,3)$; it is a constant function, as regardless of the $y$-value the distance from the line to the axis of rotation is 2. Thus $R(y)=2$.

The inside radius is formed by the line connecting $(1,1)$ and $(2,3)$. The equation of this line is $y=2x-1$, but we need to refer to it as a function of $y$. Solving for $x$ gives $r(y) = \frac12(y+1)$. 

We integrate over the $y$-bounds of $y=1$ to $y=3$. Thus the volume is
\begin{align*}
V 	&=	\pi\int_1^3\Bigl(2^2 - \bigl(\frac12(y+1)\bigr)^2\Bigr)\dd y \\
		&=	\pi\int_1^3\Bigl(-\frac14y^2-\frac12y+\frac{15}4\Bigr)\dd y \\
		&= 	\pi\Bigl[-\frac1{12}y^3-\frac14y^2+\frac{15}4y\Bigr]\Big|_1^3\\
		&= \frac{10}3\pi %\approx 10.47
		\ \text{units}^3.
\end{align*}
\end{example}

In the previous examples, the axis of rotation has either been the $x$ or $y$ axis. We will now consider a problem where the axis of rotation is some other horizontal line.

\mtable{Sketching the solid in \autoref{ex_wash_4}.}{fig:wash4}{%
 \begin{tikzpicture}[alt={Blue region between two curves from x=0 to x=2; red vertical lines mark R(x) (lower) and r(x) (higher); axes x and y.}]
  \begin{axis}[width=\marginparwidth,axis equal,
     tick label style={font=\scriptsize},axis y line=middle,
     axis x line=middle,name=myplot,axis on top,
     ymin=-.5,ymax=2.5,xmin=-.2,xmax=1.2]
   \addplot [draw={\coloronefill},fill={\coloronefill},domain=0:1] {sqrt(x)};
   \addplot [draw={\coloronefill},fill={\coloronefill},domain=0:1] {x};
   \addplot [draw={\colorone},smooth,thick,domain=0:1] {x};
   \addplot [draw={\colorone},smooth,thick,domain=0:1,samples=100] {sqrt x};
   \draw [dashed] (axis cs:-.2,2) -- (axis cs:2.2,2);
   \draw [draw={\colortwo},thick] (axis cs:.81,2) --
    node[pos=.5,right,color=black]{\scriptsize$r(x)$} (axis cs:.81,.9);
   \draw [draw={\colortwo},thick] (axis cs:.3,2) --
    node[pos=.5,right,color=black]{\scriptsize$R(x)$} (axis cs:.3,.3);
  \end{axis}
  \node [right] at (myplot.right of origin) {\scriptsize $x$};
  \node [above] at (myplot.above origin) {\scriptsize $y$};
 \end{tikzpicture}
 \\ (a) \\
\myincludeasythree{
3Droll=110.40940687660915,
3Dortho=0.004087134264409542,
3Dc2c=0.5237183570861816 0.21099673211574554 0.8253480792045593,
3Dcoo=57.86643981933594 53.05048751831055 -36.72541809082031,
3Droo=150}{\marginparwidth}{Blue region between two curves; dashed slanted line marks axis of rotation; red washer slice at sample x shows inner and outer radii; axes x, y.}{figures/figwash4_3D}
\\ (b) \\
\myincludeasythree{
3Droll=110.40940687660915,
3Dortho=0.004087134264409542,
3Dc2c=0.5237183570861816 0.21099673211574554 0.8253480792045593,
3Dcoo=57.86643981933594 53.05048751831055 -36.72541809082031,
3Droo=150}{\marginparwidth}{3D solid from rotating region between curves about dashed slanted axis; blue outer shell and red inner core; axes x, y.}{figures/figwash4b_3D}
\\ (c)}

\begin{example}[Finding volume with the Washer Method]\label{ex_wash_4}%
Find the volume of the solid formed by rotating the region bounded by $y=\sqrt x$ and $y=x$ about $y=2$.
\solution
\autoref{fig:wash4} shows the region we are rotating (a), a sample slice (b) and the full solid (c). The axis of rotation is horizontal so the radii must be functions of $x$. The radii is the distance from the axis of rotation to the curve so the outside radius of this washer is $R(x)=2-x$ and the inside radius is $r(x)=2-\sqrt x$. The region is bounded from  $x=0$ to $x=1$, thus the volume is
\begin{align*}
V&=\pi \int_0^1 \left( (2-x)^2-(2-\sqrt x)^2\right) \dd x\\
&=\pi \int_0^1 (4-4x+x^2)-(4-4\sqrt x+x)\dd x\\
&=\pi \int_0^1 x^2 -5x+4\sqrt x \dd x\\
&=\pi \left.\left[\frac13 x^3-\frac52 x^2+\frac83 x^{3/2}\right]\right|_0^1\\
&=\frac{\pi}{2} \text{units}^3.
\end{align*}
\end{example}

This section introduced a new application of the definite integral. Our default view of the definite integral is that it gives ``the area under the curve.'' However, we can establish definite integrals that represent other quantities; in this section, we computed volume.

The ultimate goal of this section is not to compute volumes of solids. That can be useful, but what is more useful is the understanding of this basic principle of integral calculus, outlined in \autoref{idea:app_of_defint}: to find the exact value of some quantity, 
\begin{itemize}
	\item we start with an approximation (in this section, slice the solid and approximate the volume of each slice), 
	\item then make the approximation better by refining our original approximation (i.e., use more slices), 
	\item	then use limits to establish a definite integral which gives the exact value.
\end{itemize}

We practice this principle in the next section where we find volumes by slicing solids in a different way.

\printexercises{exercises/07-02-exercises}

\section{The Shell Method}\label{sec:shell_method}

Often a given problem can be solved in more than one way. A particular method may be chosen out of convenience, personal preference, or perhaps necessity. Ultimately, it is good to have options.

The previous section introduced the Disk and Washer Methods, which computed the volume of solids of revolution by integrating the cross-sectional area of the solid. This section develops another method of computing volume, the \textbf{Shell Method.} Instead of slicing the solid perpendicular to the axis of rotation creating cross-sections, we now slice it parallel to the axis of rotation, creating ``shells.''

\mtable{Introducing the Shell Method.}{fig:shell_intro}{%
  \begin{tikzpicture}[alt={Blue curve y=1+1/x²; pale blue region under curve; red vertical line at sample x; axes x, y.}]
   \begin{axis}[width=1.16\marginparwidth,
     tick label style={font=\scriptsize},axis y line=middle,
     axis x line=middle,name=myplot,axis on top,
     ymin=-.2,ymax=1.2,xmin=-.2,xmax=1.2]
    \draw[draw={\colorone},fill={\coloronefill}] plot[smooth,domain=0:1]
     (axis cs:\x,{1/(\x*\x+1)}) |-(axis cs:1,0)
      --(axis cs:1,0)--(axis cs:0,0)
      -- cycle;
    \draw [thick,draw={\colortwo}] (axis cs: .75,0) -- (axis cs:.75,.64);
    \draw (axis cs:.8,.85) node[above] {\scriptsize$y=\frac1{1+x^2}$};
   \end{axis}
   \node [right] at (myplot.right of origin) {\scriptsize $x$};
   \node [above] at (myplot.above origin) {\scriptsize $y$};
  \end{tikzpicture}
\\(a)\smallskip\\
\myincludeasythree{
3Droll=134.9923368693739,
3Dortho=0.00484071671962738,
3Dc2c=0.3376615047454834 0.4032796025276184 0.8504999876022339,
3Dcoo=3.760643243789673 66.28701782226562 -12.64914608001709,
3Droo=149.99997892897957}{\marginparwidth}{Red cylindrical shell at sample x: height to curve y = 1 + 1/x^2 and thickness delta x; axes x, y.}{figures/figshell_intro_d_3D}
\\(b)\smallskip\\
\myincludeasythree{
3Droll=134.9923368693739,
3Dortho=0.00484071671962738,
3Dc2c=0.3376615047454834 0.4032796025276184 0.8504999876022339,
3Dcoo=3.760643243789673 66.28701782226562 -12.64914608001709,
3Droo=149.99997892897957}{\marginparwidth}{3D translucent solid from rotating region under y = 1 + 1/x^2; shells visible; axes x, y.}{figures/figshell_intro_a_3D}
\\(c)}

Consider \autoref{fig:shell_intro}, where the region shown in (a) is rotated around the $y$-axis forming the solid shown in (c). A small slice of the region is drawn in (a), \emph{parallel} to the axis of rotation. When the region is rotated, this thin slice forms a \textbf{cylindrical shell}, as pictured in part (b) of the figure. The previous section approximated a solid with lots of thin disks (or washers); we now approximate a solid with many thin cylindrical shells.

To compute the volume of one shell, first consider the paper label on a soup can with radius $r$ and height $h$. What is the area of this label? A simple way of determining this is to cut the label and lay it out flat, forming a rectangle with height $h$ and length $2\pi r$. Thus the area is $A = 2\pi rh$; see \autoref{fig:soupcan}(a).

Do a similar process with a cylindrical shell, with height $h$, thickness $\Delta x$, and approximate radius $r$. Cutting the shell and laying it flat forms a rectangular solid with length $2\pi r$, height $h$ and depth $\Delta x$. Thus the volume is $V \approx 2\pi rh\Delta x$; see \autoref{fig:soupcan}(b). (We say ``approximately'' since our radius was an approximation.)

By breaking the solid into $n$ cylindrical shells, we can approximate the volume of the solid as
\[V \approx \sum_{i=1}^n 2\pi r_ih_i\Delta x,\]
where $r_i$, $h_i$ and $\Delta x$ are the radius, height and thickness of the $i^{\text{th}}$ shell, respectively. 

This is a Riemann Sum. Taking a limit as the thickness of the shells approaches 0 leads to a definite integral.
\begin{align*}
V&=\lim_{n\to \infty} \sum_{i=1}^n 2\pi r_i h_i \Delta x\\
&=2\pi \int_a^b r(x)h(x)\dd x
\end{align*}

\begin{figure}[!hbt]
\flushinner{%
\tagpdfsetup{table/header-rows={2}}
\begin{tabular}{cc}
% todo Tim turn these into asymptote files
\begin{tikzpicture}[alt={Blue cylinder of height h and radius r cut along a vertical line; arrow to rectangle width 2 π r, height h labelled area = 2 π r h},scale=.67]
 \draw [draw={\colorone},thick,left color=\colorone!60,right color=\colorone!20]
  (0,0) ellipse [x radius=1.5cm, y radius=1cm];
\begin{scope}[xscale=1.5]
 \draw [draw={\colorone},thick,left color=\colorone!60,right color=\colorone!20]
 (-1,0) -- (-1,-3.5) arc (180:360:1) -- (1,0) arc (360:180:1);
 \draw [draw={\colorone},fill=white](-1,-.2) -- node [pos=.5,left,color=black] {\scriptsize $h$} (-1,-3.3) arc (180:360:1) -- (1,-.2) arc (360:180:1);
\end{scope}
 \draw [draw={\colorone},dashed] (1.06,-.91) -- node [pos=.5, above,rotate=90,color=black] {\scriptsize cut here} (1.06,-4.01);
 \draw (0,0) -- node [pos=.5,above] {\scriptsize $r$} (1.5,0);
 \draw [->,>=latex,ultra thick] (2,-2)--(3,-2);
\draw[draw={\colorone}] (3.5,-.5) -- node [pos=.5,above,color=black] {\scriptsize $2\pi r$} (8.5,-.5) -- node [pos=.5,right,color=black] {\scriptsize $h$} (8.5,-3.6)--(3.5,-3.6) -- cycle;
\draw (6,-2.05) node {\scriptsize $A = 2\pi r h$};
\end{tikzpicture}
&
\begin{tikzpicture}[alt={Blue cylindrical shell height h, radius r and thickness delta x cut and flattened; rectangle width 2 π r, height h, thickness ∆x, labelled V ≈ 2 π r h ∆x},scale=.67]
\draw [draw={\colorone},thick,left color=\colorone!60,right color=\colorone!20]
 (0,0) ellipse [x radius=1.5cm,y radius=1cm];
\begin{scope}[xscale=1.5]
\draw [draw={\colorone},thick,left color=\colorone!60,right color=\colorone!20]
 (-1,0) -- (-1,-3.5)  arc (180:360:1) -- (1,0) arc (360:180:1);
\end{scope}
\draw [draw={\colorone},thick,left color=\colorone!20,right color=\colorone!60]
 (0,0) ellipse [x radius=1.2cm,y radius=.8cm];
\draw (-1.5,-1.75) node [left,color=black] {\scriptsize $h$};
\draw [draw={\colorone},dashed] (.8475,-.565) --(1.065,-.71) -- node [pos=.5, above,rotate=90,color=black] {\scriptsize cut here} (1.065,-4.2);
\draw (0,0) -- node [pos=.5,above] {\scriptsize $r$} (1.35,0);
\draw [->,>=stealth](-1.8,.3) node [above] {\scriptsize $\Delta x$}-- (-1.3,0);
\draw [->,>=latex,ultra thick] (2,-2)--(3,-2);
\draw [draw={\colorone},thick,left color=\colorone!60,right color=\colorone!20]
 (3.5,-.3) -- (3.7,-.1) --node [pos=.5,above,color=black] {\scriptsize $2\pi r$} (8.7,-.1) -- (8.5,-.3);
\draw [draw={\colorone},thick,left color=\colorone!60,right color=\colorone!20]
 (8.5,-.3) -- (8.7,-.1) --node [pos=.5,right,color=black] {\scriptsize $h$} (8.7,-3.6) -- (8.5,-3.8);
\draw [draw={\colorone},thick,left color=\colorone!60,right color=\colorone!20]
 (3.5,-.3) --  (8.5,-.3) --  (8.5,-3.8)--(3.5,-3.8) -- cycle;
\draw [>=stealth,->] (4,0) node [above] {\scriptsize $\Delta x$} --(3.8,-.2);
\draw (6,-2.05) node {\scriptsize $V\approx 2\pi r h\Delta x$};
\end{tikzpicture}
\\
(a) & (b)
\end{tabular}}
\caption{Determining the volume of a thin cylindrical shell.%
\label{fig:soupcan}}
% otherwise, `Package nag Warning: \label in float, but not after \caption'
\end{figure}

\bigskip

\begin{keyidea}[The Shell Method]\label{idea:shell_method}%
Let a solid be formed by revolving a region $R$, bounded by $x=a$ and $x=b$, around a vertical axis. Let $r(x)$ represent the distance from the axis of rotation to $x$ (i.e., the radius of a sample shell) and let $h(x)$ represent the height of the solid at $x$ (i.e., the height of the shell). The volume of the solid is 
\index{integration!volume!Shell Method}\index{Shell Method}
\[V = 2\pi\int_a^b r(x)h(x)\dd x.\]
\end{keyidea}

\textbf{Special Cases:}
	\begin{enumerate}
	\item		When the region $R$ is bounded above by $y=f(x)$ and below by $y=g(x)$, then $h(x) = f(x)-g(x)$.
	\item		When the axis of rotation is the $y$-axis (i.e., $x=0$) then $r(x) = x$.
	\end{enumerate}

\youtubeVideo{V6nTsxumjgU}{Volumes of Revolution --- Cylindrical Shells}

Let's practice using the Shell Method.

\begin{example}[Finding volume using the Shell Method]\label{ex_shell1}%
Find the volume of the solid formed by rotating the region bounded by $y=0$, $y=1/(1+x^2)$, $x=0$ and $x=1$ about the $y$-axis.
%
\mtable{Graphing a region in \autoref{ex_shell1}.}{fig:shell1}{\begin{tikzpicture}[alt={Blue curve y=1+1/x² on 0.5<x<3; pale blue region under curve; red vertical line at sample x; axes x, y.}]
\begin{axis}[width=1.16\marginparwidth,tick label style={font=\scriptsize},
axis y line=middle,axis x line=middle,name=myplot,axis on top,xtick={1},
extra x ticks={.4},extra x tick labels={$x$},ytick={1},
ymin=-.2,ymax=1.25,xmin=-.15,xmax=1.1,]
\addplot [draw={\coloronefill},thick,smooth,domain=0:1,fill={\coloronefill}] plot {(1/(1+x^2))} -- (axis cs:1,0) -- (axis cs:0,0) -- cycle;
\addplot [draw={\colorone},thick,smooth,domain=0:1] plot {(1/(1+x^2))};
\draw [thick,draw={\colortwo}] (axis cs:.4,0) -- (axis cs:.4,.86);
\draw (axis cs:.4,.45) node [left] {$\scriptstyle h(x)\left\{\rule{0pt}{33pt}\right.$}
      (axis cs:.2,-.12) node {\scriptsize $\underbrace{\makebox[35pt]{}}_{r(x)}$}
      (axis cs: .7,1) node {\scriptsize $\ds y=\frac{1}{1+x^2}$};
\end{axis}
\node [right] at (myplot.right of origin) {\scriptsize $x$};
\node [above] at (myplot.above origin) {\scriptsize $y$};
\end{tikzpicture}}
%
\solution
This is the region used to introduce the Shell Method in \autoref{fig:shell_intro}, but is sketched again in \autoref{fig:shell1} for closer reference. A line is drawn in the region parallel to the axis of rotation representing a shell that will be carved out as the region is rotated about the $y$-axis.

The distance this line is from the axis of rotation determines $r(x)$; as the distance from $x$ to the $y$-axis is $x$, we have $r(x)=x$. The height of this line determines $h(x)$; the top of the line is at $y=1/(1+x^2)$, whereas the bottom of the line is at $y=0$. Thus $h(x) = 1/(1+x^2)-0 = 1/(1+x^2)$. The region is bounded from $x=0$ to $x=1$, so the volume is 
\begin{align*}
	V
	&= 2\pi\int_0^1 \frac{x}{1+x^2}\dd x.
%\end{align*}
\intertext{This requires substitution. Let \intertextmath{u=1+x^2}, so \intertextmath{\dd u = 2x\dd x}. We also change the bounds: \intertextmath{u(0) = 1} and \intertextmath{u(1) = 2}. Thus we have:}
%\begin{align*}
	&= \pi\int_1^2 \frac{1}{u}\dd u \\
	&= \pi\ln u\Big|_1^2\\
	&= \pi\ln 2 %\approx 2.178
	\ \text{units}^3.
\end{align*}
Note: in order to find this volume using the Disk Method, two integrals would be needed to account for the regions above and below $y=1/2$.
\end{example}

With the Shell Method, nothing special needs to be accounted for to compute the volume of a solid that has a hole in the middle, as demonstrated next.

\mtable{Graphing a region in \autoref{ex_shell2}.}{fig:shell2}{%
\begin{tikzpicture}[alt={Blue line y=2x+1; triangle under curve 0<x<1 shaded; red slice shows h(x); horizontal bracket r(x) at base, axes x,y.}]
\begin{axis}[width=1.16\marginparwidth,tick label style={font=\scriptsize},
axis y line=middle,axis x line=middle,name=myplot,axis on top,xtick={1,2,3},
extra x ticks={.4},extra x tick labels={$x$},ytick={1,2,3},
ymin=-.2,ymax=3.25,xmin=-.15,xmax=3.2,]
\draw [draw={\colorone},thick,smooth,,fill={\coloronefill}] (axis cs:0,1) -- node [black,rotate=57,pos=.5,above]  {\scriptsize $y=2x+1$} (axis cs:1,3) --  (axis cs: 1,1) -- cycle;
\draw [dashed] (axis cs:3,0) -- (axis cs:3,3.2);
\draw [thick,draw={\colortwo}] (axis cs: .4,1) -- (axis cs:.4,1.8);
\draw(axis cs:1,1.4)node [right] {$\scriptstyle\left.\rule{0pt}{14pt}\right\} h(x).$}
     (axis cs:1.7,.7) node { $\underbrace{\makebox[100pt]{}}_{r(x)}$};
\end{axis}
\node [right] at (myplot.right of origin) {\scriptsize $x$};
\node [above] at (myplot.above origin) {\scriptsize $y$};
\end{tikzpicture}
\\(a)\smallskip\\
\myincludeasythree{
3Droll=138.26595091591264,
3Dortho=0.0048407199792563915,
3Dc2c=0.16831636428833008 0.19470007717609406 0.966313362121582,
3Dcoo=74.36175537109375 71.21556854248047 -40.660457611083984,
3Droo=149.9999886313645}{\marginparwidth}{Red cylindrical shell from rotating the blue triangle under y=2x+1 where 0<x<1 around the y-axis; axes x, y.}{figures/figshell2b_3D}
\\(b)\smallskip\\
\myincludeasythree{
3Droll=138.26595091591264,
3Dortho=0.0048407199792563915,
3Dc2c=0.16831636428833008 0.19470007717609406 0.966313362121582,
3Dcoo=74.36175537109375 71.21556854248047 -40.660457611083984,
3Droo=149.9999886313645}{\marginparwidth}{Blue 3D translucent solid from rotating the triangle under y=2x+1 on 0<x<1 about the y-axis; axes x, y.}{figures/figshell2c_3D}
\\(c)}

\begin{example}[Finding volume using the Shell Method]\label{ex_shell2}%
Find the volume of the solid formed by rotating the triangular region determined by the points $(0,1)$, $(1,1)$ and $(1,3)$ about the line $x=3$.
\solution
The region is sketched in \autoref{fig:shell2}(a) along with a line within the region parallel to the axis of rotation. In part (b) of the figure, we see a sample shell, and in part (c) the whole solid is shown.

The height of the sample shell is the distance from $y=1$ to $y=2x+1$, the line that connects the points $(0,1)$ and $(1,3)$. Thus $h(x) = 2x+1-1 = 2x$. The radius of the sample shell is the distance from $x$ to $x=3$; that is, it is $r(x)=3-x$. The $x$-bounds of the region are $x=0$ to $x=1$, giving
\begin{align*}
	V
	&= 2\pi\int_0^1 (3-x)(2x)\dd x \\
	&= 2\pi\int_0^1 \bigl(6x-2x^2\bigr)\dd x \\
	&= 2\pi\left.\left(3x^2-\frac23x^3\right)\right|_0^1\\
	&= \frac{14}{3}\pi%\approx 14.66
	\ \text{units}^3.
\end{align*}
\end{example}

When revolving a region around a horizontal axis, we must consider the radius and height functions in terms of $y$, not $x$.

\mtable{Graphing a region in \autoref{ex_shell3}.}{fig:shell3}{%
\begin{tikzpicture}[alt={Blue triangular region bounded by x=0, y=1, and line y=2x+1; red horizontal slice shows h(y); bracket marks r(y); axes x,y.}]
\begin{axis}[width=1.16\marginparwidth,tick label style={font=\scriptsize},
axis y line=middle,axis x line=middle,name=myplot,axis on top,xtick={1,2,3},
ytick={1,2,3},extra y ticks={1.8},extra y tick labels={$y$},
ymin=-.2,ymax=3.25,xmin=-.15,xmax=1.5,]
\draw [draw={\colorone},thick,smooth,,fill={\coloronefill}] (axis cs:0,1) -- node [black,rotate=41,pos=.5,above]  {\scriptsize $x=\frac12y-\frac12$} (axis cs:1,3) --  (axis cs: 1,1) -- cycle;
\draw [thick,draw={\colortwo}] (axis cs: .4,1.8) -- node [black,pos=.5,below] {$\underbrace{\makebox[42pt]{}}_{h(y)}$} (axis cs:1,1.8);
\draw (axis cs:1,.9) node [right] {$\scriptstyle\left.\rule{0pt}{29pt}\right\} r(y)$};
\end{axis}
\node [right] at (myplot.right of origin) {\scriptsize $x$};
\node [above] at (myplot.above origin) {\scriptsize $y$};
\end{tikzpicture}
\\(a) \\
\myincludeasythree{
3Droll=97.25039174241927,
3Dortho=0.0048407199792563915,
3Dc2c=0.34401965141296387 0.04830114170908928 0.9377192854881287,
3Dcoo=71.46640014648438 9.081968307495117 -11.759183883666992,
3Droo=149.99998743176678}{\marginparwidth}{Red shell from rotating blue triangle under y = 2x + 1 for 0<x<1 about a slanted axis; axes x, y.}
{figures/figshell3b_3D}
\\(b)\\
\myincludeasythree{
3Droll=97.25039174241927,
3Dortho=0.0048407199792563915,
3Dc2c=0.34401965141296387 0.04830114170908928 0.9377192854881287,
3Dcoo=71.46640014648438 9.081968307495117 -11.759183883666992,
3Droo=149.99998743176678}{\marginparwidth}{3D translucent solid rotating the triangle with corners (0,1),(1,3),(3,1) about the x-axis; axes x, y.}{figures/figshell3c_3D}
\\(c)}

\begin{example}[Finding volume using the Shell Method]\label{ex_shell3}%
Find the volume of the solid formed by rotating the region given in \autoref{ex_shell2} about the $x$-axis.
\solution
The region is sketched in \autoref{fig:shell3}(a). In part (b) of the figure the sample shell is drawn, and the solid is sketched in (c). (Note that the triangular region looks ``short and wide'' here, whereas in the previous example the same region looked ``tall and narrow.'' This is because the bounds on the graphs are different.)

The height of the sample shell is an $x$-distance, between $x=\frac12y-\frac12$ and $x=1$. Thus $h(y) = 1-(\frac12y-\frac12) = -\frac12y+\frac32.$ The radius is the distance from $y$ to the $x$-axis, so $r(y) =y$. The $y$ bounds of the region are $y=1$ and $y=3$, leading to the integral
\begin{align*}
V &= 2\pi\int_1^3\left[y\left(-\frac12y+\frac32\right)\right]\dd y \\
	&= 2\pi\int_1^3\left[-\frac12y^2+\frac32y\right]\dd y \\
	&= 2\pi\left.\left[-\frac16y^3+\frac34y^2\right]\right|_1^3 \\
	&= 2\pi\left[\frac94-\frac7{12}\right]\\
	&=	\frac{10}{3}\pi %\approx 10.472
	\ \text{units}^3.
\end{align*}
\end{example}

The following example shows how there are times when it does not matter which method you choose to evaluate the volume of a solid. In \autoref{ex_wash_4} we found the volume of the solid formed by rotating the region bounded by $y=\sqrt x$ and $y=x$ about $y=2$. We will now demonstrate how to find the volume with the shell method. Note that your answer should be the same whichever method you choose.

\begin{example}[Using the shell method instead of the washer method]\label{ex_shell_wash_eq}%
Find the volume of the solid formed by rotating the region bounded by $y=\sqrt x$ and $y=x$ about $y=2$ using the Shell Method.
\solution
Since our shells are parallel to the axis of rotation, we must consider the radius and height functions in terms of $y$. The radius of a sample shell will be $r(y)=2-y$ and the height of a sample shell will be $h(y)=y-y^2$. The $y$ bounds for the region will be $y=0$ to $y=1$ resulting in the integral
\begin{align*}
V&=2\pi \int_0^1 (2-y)(y-y^2)\dd y\\
&=2\pi\int_0^1 y^3-3y^2+2y\dd y\\
&=\frac{\pi}{2} \text{units}^3.
\end{align*}
\end{example}

At the beginning of this section it was stated that ``it is good to have options.'' The next example finds the volume of a solid rather easily with the Shell Method, but using the Washer Method would be quite a chore.

\mtable{Graphing a region in \autoref{ex_shell4}.}{fig:shell4}{%
\begin{tikzpicture}[alt={Blue curve y=1+1/x² on 0.5<x<3; shaded beneath; red line labeled h(x); bracket marks r(x); axes x and y.}]
\begin{axis}[width=\marginparwidth,tick label style={font=\scriptsize},
			axis y line=middle,axis x line=middle,name=myplot,axis on top,
			xtick=\empty,extra x ticks={1,2,3},ytick={1,2},
			ymin=-.2,ymax=2.5,xmin=-.1,xmax=3.5]
\addplot [draw={\coloronefill},fill={\coloronefill},domain=0:2] {3*x-x^2};
\addplot [draw={\colorone},smooth,thick,domain=0:3] {3*x-x^2};
\addplot[draw={\colorone},smooth,thick,domain=0:3] {x};
\draw [thick,draw={\colortwo}] (axis cs: 1,1) -- (axis cs:1,2);
\draw (axis cs: .9,1.5) node[right]{$\scriptstyle\left.\rule{0pt}{20pt}\right\} h(x)$};
\draw (axis cs:.5,1.1) node [below] {$\underbrace{\makebox[30pt]{}}_{r(x)}$} ;
\end{axis}
\node [right] at (myplot.right of origin) {\scriptsize $x$};
\node [above] at (myplot.above origin) {\scriptsize $y$};
\end{tikzpicture}
\\(a)\smallskip\\
\myincludeasythree{
3Droll=120.03842998357472,
3Dortho=0.0036980914883315568,
3Dc2c=0.36060017347335815 0.2289222925901413 0.9041913747787476,
3Dcoo=22.334463119506836 36.29957962036133 -19.963319778442383,
3Droo=150}{\marginparwidth}{Red cylindrical shell 1<y<2 around y-axis, formed by rotating a sample strip of the blue curve region; axes x, y.}{figures/figshellparab_3D}
\\(b)\smallskip\\
\myincludeasythree{
3Droll=120.03842998357472,
3Dortho=0.0036980914883315568,
3Dc2c=0.36060017347335815 0.2289222925901413 0.9041913747787476,
3Dcoo=22.334463119506836 36.29957962036133 -19.963319778442383,
3Droo=150}{\marginparwidth}{Blue funnel-shaped solid: region between y=x and y=3x−x^2 for 0<y<3 rotated about y-axis; axes x, y.}{figures/figshellparab_b_3D}
\\(c)}

\begin{example}[Finding volume using the Shell Method]\label{ex_shell4}%
Find the volume of the solid formed by rotating the region bounded by $y=3x-x^2$ and $y=x$ about the $y$-axis.
\solution
The region, a sample shell, and the resulting solid are shown in \autoref{fig:shell4}. The radius of a sample shell is $r(x)=x$; the height of a sample shell is $h(x)=(3x-x^2)-x=2x-x^2$. The $x$ bounds on the region are $x=0$ to $x=2$ leading to the integral
\begin{align*}
V&=2\pi \int_0^2 x(2x-x^2) \dd x\\
&=2\pi \int_0^2 2x^2-x^3 \dd x\\
&=2\pi \left.\left[ \frac23 x^3-\frac14 x^4 \right] \right|_0^2\\
&=\frac{8\pi}{3} \text{units}^3.
\end{align*}

Note that in order to use the Washer Method, we would need to solve $y=3x-x^2$ for $x$, requiring us to complete the square. We must evaluate two integrals as we have two different sample slices. The volume can be computed as 
\begin{align*}
V&=\pi\int_0^2\left(y^2-\left(\frac32-\sqrt{\frac94 -y}\right)^2\right)\dd y \\
&\qquad+\pi\int_2^{9/4}\left(\left(\frac32+\sqrt{\frac94 -y}\right)^2
-\left(\frac32-\sqrt{\frac94 -y}\right)^2\right)\dd y
% this next expression is incorrect, if we want to go back to including it
%\\&=\pi\int_0^2\left(y-\frac32 +\sqrt{\frac94 -y}\right)^2 \dd y + 2\pi\int_2^{9/4} \left(2\sqrt{\frac94 -y}\right)^2\dd y
\end{align*}
While this integral is not impossible to solve, using the Shell Method gave us a significantly easier way to compute the volume.
\end{example}

We finish with an example where we can use either method to find the volume.

\begin{example}[Finding volume using both methods%the Washer Method and the Shell Method% otherwise, Method goes to the next line
]\label{ex_shell_wash}%
Find the volume of the solid formed by rotating the region bounded by the curves $y=x^2+3$, $y=2x+1$, $x=0$, and $x=1$ about the y-axis using both the washer method and the shell method.
\solution
We'll start with the shell method, since that turns out to be easier.  The region, a sample shell, and the resulting solid are shown in \autoref{shell_wash_fig} parts a, b, and e respectively.  The volume is
\[
V=2\pi\int_0^1 x((x^2+3)-(2x+1))\dd x
=\frac{7\pi}6\text{units}^3.
\]
%\begin{align*}
%V&=2\pi\int_0^1 x((x^2+3)-(2x+1))\dd x\\
%&=2\pi\int_0^1 x^3-2x^2+2x\dd x\\
%%&=2\pi\left[\frac14-\frac23+1\right]\\
%&=\frac{7\pi}6\text{units}^3.
%\end{align*}

With the washer method, we need to integrate with respect to $y$ because we are rotating around a vertical axis.  We also need to divide the region in two because the washers will run into different boundaries at different heights.  We have indicated the region and the two different types of washers in \autoref{shell_wash_fig} parts c and d.  The volume is
\begin{align*}
V&=\pi\int_1^3\left(\frac{y-1}2\right)^2\dd y
+\pi\int_3^4\left(1^2-\left(\sqrt{y-3}\right)^2\right)\dd y\\
%&=\pi\int_1^3\frac{y^2}4-\frac y2+\frac14\dd y+\pi\int_3^4 4-y\dd y\\
%&=\pi\left[\frac{y^3}{12}-\frac{y^2}4+\frac y4\right]_1^3+\pi\left[4y-\frac{y^2}2\right]_3^4\\
%&=\pi\left[\frac{27}{12}-\frac94+\frac34\right]
%-\pi\left[\frac1{12}-\frac14+\frac14\right]
%+\pi\left[16-8\right]-\pi\left[12-\frac92\right]\\
%&=\pi\frac34-\pi\frac1{12}+8\pi-12\pi+\pi\frac92\\
&=\pi\frac23+\pi\frac12
%\\&
=\frac{7\pi}6\text{units}^3.
\end{align*}
We are reassured to find the same answer using either method.
\end{example}

\noindent\begin{minipage}{\linewidth}\tagstructbegin{tag=Sect}\noindent%
\captionsetup{type=figure}%
\flushinner{%
%\captionsetup{type=figure}
\tagpdfsetup{table/header-rows={2,4}}
\begin{tabular}{ c c c }
\begin{tikzpicture}[alt={Blue triangle with vertices (0,1), (1,1), (1,4); red vertical slice shows h(y); bracket to y-axis marks r(y); axes x,y.}]
\begin{axis}[width=\marginparwidth,tick label style={font=\scriptsize},
			axis y line=middle,axis x line=middle,name=myplot,axis on top,
			xtick=\empty,extra x ticks={1,2},ytick={1,2,3,4},
			ymin=-.2,ymax=4.5,xmin=-.1,xmax=1.5]%[,axis equal]
\draw[\colorone,thick](axis cs:0,1)--(axis cs:1,3)
%node [color=black,left,pos=.5] {$\underbrace{\makebox[30pt]{}}_{r(x)}$}
--(axis cs:1,4);
\node at(axis cs:.25,2){$\underbrace{\makebox[30pt]{}}_{r(x)}$};
\addplot[draw={\colorone},smooth,thick,domain=0:1] {x^2+3};
\draw[thick,draw={\colortwo}] (axis cs: .5,2) -- (axis cs:.5,3.25)%;
node[color=black,right,pos=.5]{$\scriptstyle\left.\rule{0pt}{12pt}\right\} h(x)$};
\end{axis}
\node [right] at (myplot.right of origin) {\scriptsize $x$};
\node [above] at (myplot.above origin) {\scriptsize $y$};
\end{tikzpicture}
&
\myincludeasythree{
3Droll=115.86380470248459,
3Dortho=0.003982423339039087,
3Dc2c=0.35041913390159607 0.18568213284015656 0.9180024862289429,
3Dcoo=5.665955543518066 74.65992736816406 -21.359724044799805,
3Droo=150}{\marginparwidth}{Red cylindrical shell 1<y<3 from rotating the sample slice of the blue triangle about the y-axis; axes x,y.}{figures/shellwash}
&
\begingroup
 \newcommand{\localimage}{%
  \myincludeasythree{
   3Droll=124.28706719451111,
   3Dortho=0.003981335088610649,
   3Dc2c=0.3404673635959625 0.24756169319152832 0.9070805907249451,
   3Dcoo=-5.416632652282715 90.77053833007812 -22.686864852905273,
   3Droo=150}{\marginparwidth}{Blue triangle; red horizontal slice at y=3 highlights inner r(x) and outer R(x); axes x,y.}{figures/shellwash_b}%%
 }
 % todo Tim
 % see https://github.com/brucemiller/LaTeXML/issues/1709
 \ifbool{latexml}{
  \localimage
 }{
  \tagpdfsetup{table/multirow={2}}
  \multirow{12.6}{*}{\localimage}
  % see https://tex.stackexchange.com/a/364935/107497
 }
\endgroup
\smallskip\\(a)&(b)\smallskip\\
\begin{tikzpicture}[alt={Two red washers at y=2 and 3 illustrate cross-sections when triangle rotates about the x-axis; axes x,y.}]
\begin{axis}[width=\marginparwidth,tick label style={font=\scriptsize},
			axis y line=middle,axis x line=middle,name=myplot,axis on top,
			xtick=\empty,extra x ticks={1,2},ytick={1,2,3,4},
			ymin=-.2,ymax=4.5,xmin=-.1,xmax=1.5]%[,axis equal]
\draw[\colorone,thick](axis cs:0,1)--(axis cs:1,3)
%node [color=black,left,pos=.5] {$\underbrace{\makebox[25pt]{}}_{r(x)}$}
--(axis cs:1,4);
\addplot[draw={\colorone},smooth,thick,domain=0:1] {x^2+3};
\draw[thick,draw={\colortwo}] (axis cs:0,3.64)--(axis cs: .8,3.64)
node[color=black,above,pos=.5]{\scriptsize$r(x)$};
\draw[thick,draw={\colortwo}] (axis cs: 0,2) -- (axis cs:.5,2)%;
node[color=black,above,pos=.5]{\scriptsize$R(x)$};
\end{axis}
\node [right] at (myplot.right of origin) {\scriptsize $x$};
\node [above] at (myplot.above origin) {\scriptsize $y$};
\end{tikzpicture}
&
\myincludeasythree{
3Droll=120.0383890024397,
3Dortho=0.004286274313926697,
3Dc2c=0.36060017347335815 0.2289222776889801 0.9041913747787476,
3Dcoo=9.441178321838379 68.38707733154297 -22.945234298706055,
3Droo=150}{\marginparwidth}{Blue translucent solid formed by rotating the triangle about the x-axis; inner cavity visible; axes x,y.}{figures/shellwash_c}
\smallskip\\(c)&(d)&(e)
\end{tabular}%
}% closing flushinner
\captionof{figure}{Graphing the region in \autoref{ex_shell_wash}.}
\label{shell_wash_fig}\tagstructend
\end{minipage}%
\bigskip

%We end this section with a table summarizing the usage of the Washer and Shell Methods.
%
%\begin{keyidea}[Summary of the Washer and Shell Methods]\label{idea:shell_and_washer}%
%Let a region $R$ be given with $x$-bounds $x=a$ and $x=b$ and $y$-bounds $y=c$ and $y=d$.\smallskip\\
%\tagpdfsetup{table/header-rows={1}}
%\begin{tabular}{ c c c c }
% 		& Washer Method & & Shell Method \\
% Horizontal Axis  & $\ds \pi\int_a^b \bigl(R(x)^2-r(x)^2\bigr)\dd x$ & & $\ds 2\pi\int_c^d r(y)h(y)\dd y$ \\ \\
% Vertical Axis &  $\ds\pi \int_c^d\bigl(R(y)^2-r(y)^2\bigr)\dd y$ & & $\ds 2\pi\int_a^b r(x)h(x)\dd x$
% \end{tabular}
%\index{integration!volume!Washer Method}\index{Washer Method}\index{integration!volume!Shell Method}\index{Shell Method}
%\end{keyidea}

As in the previous section, the real goal of this section is not to be able to compute volumes of certain solids. Rather, it is to be able to solve a problem by first approximating, then using limits to refine the approximation to give the exact value. In this section, we approximate the volume of a solid by cutting it into thin cylindrical shells. By summing up the volumes of each shell, we get an approximation of the volume. By taking a limit as the number of equally spaced shells goes to infinity, our summation can be evaluated as a definite integral, giving the exact value.

%We use this same principle again in the next section, where we find the length of curves in the plane.

\printexercises{exercises/07-03-exercises}

\section{Work}\label{sec:work}

\emph{Work} is the scientific term used to describe the action of a force which moves an object. When a constant force $F$ is applied to move an object a distance $d$, the amount of work performed is $W=F\cdot d$. 

The SI unit of force is the newton, (kg$\cdot$m/s$^2$), and the SI unit of distance is a meter (m). The fundamental unit of work is one newton-meter, or a joule (J). That is, applying a force of one newton for one meter performs one joule of work. In Imperial units (as used in the United States), force is measured in pounds (lb) and distance is measured in feet (ft), hence work is measured in ft-lb. 

\mnote{\textbf{Note:} \emph{Mass} and \emph{weight} are closely related, yet different, concepts. The mass $m$ of an object is a quantitative measure of that object's resistance to acceleration. The weight $w$  of an object is a measurement of the force applied to the object by the acceleration of gravity $g$.\bigskip\\
Since the two measurements are proportional, $w=m\cdot g$, they are often used interchangeably in everyday conversation. When computing work, one must be careful to note which is being referred to. When mass is given, it must be multiplied by the acceleration of gravity to reference the related force.}

When force is constant, the measurement of work is straightforward. For instance, lifting a 200 lb object 5 ft performs $200\cdot 5 = 1000$ ft-lb of work. 

What if the force applied is variable? For instance, imagine a climber pulling a 200 ft rope up a vertical face. The rope becomes lighter as more is pulled in, requiring less force and hence the climber performs less work.

In general, let $F(x)$ be a force function on an interval $[a,b]$. We want to measure the amount of work done applying the force $F$ from $x=a$ to $x=b$. We can approximate the amount of work being done by partitioning $[a,b]$ into subintervals $a=x_0<x_1 <\dots <x_n=b$ and assuming that $F$ is constant on each subinterval. Let $c_i$ be a value in the $i\,^{\text{th}}$ subinterval $[x_i,x_{i+1}]$. Then the work done on this interval is approximately $W_i\approx F(c_i)\cdot(x_{i+1}-x_i) = F(c_i)\Delta x_i$, a constant force $\times$ the distance over which it is applied. The total work is 
\[ W = \sum_{i=1}^n W_i \approx \sum_{i=1}^n F(c_i)\Delta x_i.\]
This, of course, is a Riemann sum. Taking a limit as the subinterval lengths go to zero give an exact value of work which can be evaluated through a definite integral.

\begin{keyidea}[Work]\label{idea:work}%
Let $F(x)$ be a continuous function on $[a,b]$ describing the amount of force being applied to an object in the direction of travel from distance $x=a$ to distance $x=b$. The total work $W$ done on $[a,b]$ is
\index{integration!work}\index{work}
\[W = \int_a^b F(x)\dd x.\]
\end{keyidea}

\youtubeVideo{2pbInn9PkHQ}{Finding Work using Calculus --- The Cable/Rope Problem}

\begin{example}[Computing work performed: applying variable force]\label{ex_work1}%
A 60 m climbing rope is hanging over the side of a tall cliff. How much work is performed in pulling the rope up to the top, where the rope has a mass of 66 g/m? 
%How much work is performed pulling a 60 m climbing rope up a cliff face, where the rope has a mass of 66 g/m?
\solution
We need to create a force function $F(x)$ for $0\le x\le60$. To do so, we must first decide what $x$ is measuring: is it the length of the rope still hanging or is it the amount of rope pulled in? As long as we are consistent, either approach is fine. We adopt for this example the convention that $x$ is the amount of rope pulled in. This seems to match intuition better; pulling up the first 10 meters of rope involves $x=0$ to $x=10$ instead of $x=60$ to $x=50$. 

As $x$ is the amount of rope pulled in, the amount of rope still hanging is $60-x$. This length of rope has a mass of 66 g/m, or $0.066$ kg/m. The the mass of the rope still hanging is $0.066(60-x)$ kg; multiplying this mass by the acceleration of gravity, 9.8 m/s$^2$, gives our variable force function
\[F(x) = (9.8)(0.066)(60-x) = 0.6468(60-x).\]

Thus the total work performed in pulling up the rope is 
\[W = \int_0^{60} 0.6468(60-x)\dd x = 1,164.24\ \text{J}.\]

By comparison, consider the work done in lifting the entire rope 60 meters. The rope weights $60\times 0.066 \times 9.8 = 38.808$ N, so the work applying this force for 60 meters is $60\times 38.808 = 2,328.48$ J. This is exactly twice the work calculated before (and we leave it to the reader to understand why.)
\end{example}

\begin{example}[Computing work performed: applying variable force]\label{ex_work1_5}%
Consider again pulling a 60 m rope up a cliff face, where the rope has a mass of 66 g/m. At what point is exactly half the work performed?
\solution
From \autoref{ex_work1} we know the total work performed is $1,164.24$ J. We want to find a height $h$ such that the work in pulling the rope from a height of $x=0$ to a height of $x=h$ is 582.12, half the total work. Thus we want to solve for $h$ in the equation
\[\int_0^h 0.6468(60-x)\dd x = 582.12.\]

We see that
\begin{align*}
	\int_0^h 0.6468(60-x)\dd x &= 582.12 \\
	\left(38.808x-0.3234x^2\right)\Big|_0^h &=582.12 \\
	38.808h-0.3234h^2 &=582.12 \\
	-0.3234h^2+38.808h-582.12 &=0 \qquad\qquad\text{\small(Apply the Quadratic Formula)} \\
%\intertext{Apply the Quadratic Formula.}
	h&\approx17.57 \ \text{and}\ 102.43
\end{align*}
As the rope is only 60m long, the only sensible answer is $h=17.57$. Thus about half the work is done pulling up the first 17.57m the other half of the work is done pulling up the remaining 42.43m.
\end{example}

\mnote{\textbf{Note:} In \autoref{ex_work1_5}, we find that half of the work performed in pulling up a 60 m rope is done in the last 42.43 m. Why is it not coincidental that $60/\sqrt{2}=42.43$?}

\begin{example}[Computing work performed: applying variable force]\label{ex_work2}%
A box of 100 lb of sand is being pulled up at a uniform rate a distance of 50 ft over 1 minute. The sand is leaking from the box at a rate of 1 lb/s. The box itself weighs 5 lb and is pulled by a rope weighing .2 lb/ft. 
\begin{enumerate}
	\item	How much work is done lifting just the rope?
	\item	How much work is done lifting just the box and sand?
	\item	What is the total amount of work performed?
\end{enumerate}
\solution
\begin{enumerate}
	\item	We start by forming the force function $F_r(x)$ for the rope (where the subscript denotes we are considering the rope). As in the previous example, let $x$ denote the amount of rope, in feet, pulled in. (This is the same as saying $x$ denotes the height of the box.) The weight of the rope with $x$ feet pulled in is $F_r(x) = 0.2(50-x) = 10-0.2x$. (Note that we do not have to include the acceleration of gravity here, for the \emph{weight} of the rope per foot is given, not its \emph{mass} per meter as before.) The work performed lifting the rope is 
	\[W_r = \int_0^{50} (10-0.2x)\dd x = 250\ \text{ft-lb}.\]
	
	\item	The sand is leaving the box at a rate of 1 lb/s. As the vertical trip is to take one minute, we know that 60 lb will have left when the box reaches its final height of 50 ft. Again letting $x$ represent the height of the box, we have two points on the line that describes the weight of the sand: when $x=0$, the sand weight is 100 lb, producing the point $(0,100)$; when $x=50$, the sand in the box weighs 40 lb, producing the point $(50,40)$. The slope of this line is $\frac{100-40}{0-50} = -1.2$, giving the equation of the weight of the sand at height $x$ as $w(x) = -1.2x+100$. The box itself weighs a constant 5 lb, so the total force function is $F_b(x) = -1.2x+105$. Integrating from $x=0$ to $x=50$ gives the work performed in lifting box and sand:
	\[W_b = \int_0^{50} (-1.2x+105)\dd x = 3750\ \text{ft-lb.}\]
	
	\item	The total work is the sum of $W_r$ and $W_b$: $250+3750=4000$ ft-lb. We can also arrive at this via integration:
	\begin{align*}
	W
	&= \int_0^{50} (F_r(x)+F_b(x))\dd x \\
	&= \int_0^{50} (10-0.2x-1.2x+105)\dd x \\
	&= \int_0^{50} (-1.4x+115) \dd x \\
	&= 4000 \ \text{ft-lb.}
	\end{align*}	
\end{enumerate}
\end{example}

\subsection{Hooke's Law and Springs}

Hooke's Law states that the force required to compress or stretch a spring $x$ units from its natural length is proportional to $x$; that is, this force is $F(x) = kx$ for some constant $k$. For example, if a force of 1 N stretches a given spring 2 cm, then a force of 5 N will stretch the spring 10 cm. Converting the distances to meters, we have that stretching this spring 0.02 m requires a force of $F(0.02) = k(0.02) = 1$ N, hence $k = 1/0.02 = 50$ N/m. 
\index{Hooke's Law}

\begin{example}[Computing work performed: stretching a spring]\label{ex_spring1}%
A force of 20 lb stretches a spring from a natural length of 7 inches to a length of 12 inches. How much work was performed in stretching the spring to this length?
\solution
In many ways, we are not at all concerned with the actual length of the spring, only with the amount of its change. Hence, we do not care that 20 lb of force stretches the spring to a length of 12 inches, but rather that a force of 20 lb stretches the spring by 5 in. This is illustrated in \autoref{fig:spring1}; we only measure the change in the spring's length, not the overall length of the spring.

\begin{center}\tagstructbegin{tag=Sect}
\begin{tikzpicture}[alt={Two views: spring with block; first block at x≈1 on 0<x<6 scale, second after stretch at x≈6; arrow F right.}]
 \begin{scope}[xscale=.2,yscale=.5,shift={(.5,0)}]
  \draw [thick,draw={\colorone}] (-1,.5)--(-.5,.5)--(0,0);
  \foreach \x in {0,1,...,5} {
   \begin{scope}[shift={(\x*2,0)}]
    \draw [thick,draw={\colorone}] (0,0)--(1,1)--(2,0);
   \end{scope}
  }
  \begin{scope}[shift={(12,0)}]
   \draw [thick,draw={\colorone}] (0,0) --(1,1) --(1.5,.5)--(2,.5);
  \end{scope}
 \end{scope}
 \draw [shift={(.9,0)},thick,{\colorone},left color={\colorone!70},right color={\coloronefill}]
 (2,0) rectangle (2.5,.5);
 \draw [shift={(.9,0)},->,>=stealth]
  (2.75,.15)-= node [pos=.5,above] {$F$} (3.25,.15);
 \begin{scope}[shift={(0,-.25)}]
  \draw [thick] (-.5,0)--(5.5,0);
  \foreach \x in {0,...,6} {
   \draw (\x*.4+2.8,.1)--(\x*.4+2.8,-.1) node [below] {\scriptsize $\x$};
  }
 \end{scope}
 \begin{scope}[shift={(0,-1.75)}]
  \begin{scope}[xscale=.343,yscale=.5,shift={(.5,0)}]
   \draw [thick,draw={\colorone}] (-1,.5)--(-.5,.5)--(0,0);
   \foreach \x in {0,1,...,5} {
    \begin{scope}[shift={(\x*2,0)}]
     \draw [thick,draw={\colorone}] (0,0)--(1,1)--(2,0);
    \end{scope}
   }
   \begin{scope}[shift={(12,0)}]
    \draw [thick,draw={\colorone}] (0,0) --(1,1) --(1.5,.5)--(2,.5);
   \end{scope}
  \end{scope}
  \begin{scope}[shift={(0,-.25)}]
   \draw [thick] (-.5,0)--(5.5,0);
   \foreach \x in {0,...,6} {
    \draw (\x*.4+2.8,.1)--(\x*.4+2.8,-.1) node [below] {\scriptsize $\x$};
   }
  \end{scope}
  \draw [shift={(2.9,0)},thick,{\colorone},left color={\colorone!70},right color={\coloronefill}]
  (2,0) rectangle (2.5,.5);
 \end{scope}             
\end{tikzpicture}
\captionsetup{type=figure}
\caption{Illustrating the important aspects of stretching a spring in computing work in \autoref{ex_spring1}.}
\label{fig:spring1}\tagstructend
\end{center}

Converting the units of length to feet, we have
\[F(5/12) = (5/12)k = 20\ \text{lb}.\]
Thus $k = 48$ lb/ft and $F(x) = 48x$. 

We compute the total work performed by integrating $F(x)$ from $x=0$ to $x=5/12$:
\begin{align*}
	W
	&= 	\int_0^{5/12} 48x \dd x \\
	&=	24x^2\Big|_0^{5/12} \\
	&=  25/6 %\approx 4.1667
	\ \text{ft-lb.}
\end{align*}
\end{example}

\subsection{Pumping Fluids}

Another useful example of the application of integration to compute work comes in the pumping of fluids, often illustrated in the context of emptying a storage tank by pumping the fluid out the top. This situation is different than our previous examples for the forces involved are constant. After all, the force required to move one cubic foot of water (about 62.4 lb) is the same regardless of its location in the tank. What is variable is the distance that cubic foot of water has to travel; water closer to the top travels less distance than water at the bottom, producing less work.

\mtable{Weight and Mass densities}{fig:weight_density}{%
\tagpdfsetup{table/header-rows={1}}
\begin{tabular}{lll}
Fluid & lb/ft$^3$ & kg/m$^3$ \\ \midrule
Gasoline & \phantom{0}45.93 & \phantom{00}737.22\\
Methanol & \phantom{0}49.3 & \phantom{00}791.3\\
Fuel Oil & \phantom{0}55.46 & \phantom{00}890.13 \\
Water & \phantom{0}62.4 & \phantom{0}1000 \\
Milk, whole & \phantom{0}63.6 & \phantom{0}1020 \\
Milk, nonfat & \phantom{0}65.4 & \phantom{0}1050 \\
Concrete & 150 & \phantom{0}2400 \\
Iodine & 307 & \phantom{0}4927 \\
Mercury & 844 & 13546
\end{tabular}}

We demonstrate how to compute the total work done in pumping a fluid out of the top of a tank in the next two examples.

\begin{example}[Computing work performed: pumping fluids]\label{ex_pump1}%
A cylindrical storage tank with a radius of 10 ft and a height of 30 ft is filled with water, which weighs approximately 62.4 lb/ft$^3$. Compute the amount of work performed by pumping the water up to a point 5 feet above the top of the tank.
\solution
We will refer often to \autoref{fig:pump1a} which illustrates the salient aspects of this problem.

\mtable{Illustrating a water tank in order to compute the work required to empty it in \autoref{ex_pump1}.}{fig:pump1a}{\begin{tikzpicture}[alt={Cylinder 10 ft diameter; blue water slice between y_i<y<y_{i+1}; arrow marks lift 35−y_i; y-axis 0<y<35.},x=.1cm,y=.125cm,>=stealth]
\draw [->] (0,-2) -- (0,37) node [above] {\scriptsize $y$};
\draw (-1,0) node [left] {\scriptsize 0} -- (1,0)
      (-1,30) node [left] {\scriptsize 30} -- (1,30)
      (-1,35) node [left] {\scriptsize 35} -- (1,35);
\draw (2,12)--(4,12)  (2,35)--(4,35)
      (3,12) -- node [pos=.5,fill=white,rotate=90] {\scriptsize $35-y_i$} (3,35);
\begin{scope}[xscale=4,shift={(4.5,0)}]         
\draw (0,30) -- node [pos=.5,above] {\scriptsize $10$} (2.9,29.22);             
\draw [thick] (3,30) --(3,0) arc (0:-180:3) -- (-3,30);
\draw [thick](0,30) circle (3);
\draw [thick,dashed] (3,0) arc (0:180:3);
\foreach \y/\x in {12/{$y_i$},16/{$y_{i+1}$}} {
	\draw (3,\y) node [right] {\scriptsize\x} arc (0:-180:3);
    \draw [dashed] (3,\y) arc (0:180:3);
}
\draw (5.1,14) node {\scriptsize $\left.\rule{0pt}{.3cm}\right\}\Delta y_i$};
\draw [thick,left color={\colorone},right color={\coloronefill}]
 (0,16) circle (3);
\draw [thick,left color={\colorone},right color={\coloronefill}]
 (-3,16) -- (-3,12)  arc (180:360:3) -- (3,16) arc (360:180:3);
\end{scope}
\end{tikzpicture}}
% todo Tim make the water tank figure into 3d?

We start as we often do: we partition an interval into subintervals. We orient our tank vertically since this makes intuitive sense with the base of the tank at $y=0$. Hence the top of the water is at $y=30$, meaning we are interested in subdividing the $y$-interval $[0,30]$ into $n$ subintervals as 
\[0 = y_0 < y_1 < \dots < y_n = 30.\]
Consider the work $W_i$ of pumping only the water residing in the $i^{\text{th}}$ subinterval, illustrated in \autoref{fig:pump1a}. The force required to move this water is equal to its weight which we calculate as volume $\times $ density. The volume of water in this subinterval is 
$V_i = 10^2\pi \Delta y_i$; its density is $62.4$ lb/ft$^3$. Thus the required force is $6240\pi\Delta y_i$ lb.

We approximate the distance the force is applied by using any $y$-value contained in the $i^{\text{th}}$ subinterval; for simplicity, we arbitrarily use $y_i$ for now (it will not matter later on). The water will be pumped to a point 5 feet above the top of the tank, that is, to the height of $y=35$ ft. Thus the distance the water at height $y_i$ travels is $35-y_i$ ft. 

In all, the approximate work $W_i$ performed in moving the water in the $i^{\text{th}}$ subinterval to a point 5 feet above the tank is 
\[W_i \approx 6240\pi\Delta y_i(35-y_i).\]
To approximate the total work performed in pumping out all the water from the tank, we sum all the work $W_i$ performed in pumping the water from each of the $n$ subintervals of $[0,30]$:
\[W \approx \sum_{i=1}^n W_i = \sum_{i=1}^n 6240\pi\Delta y_i(35-y_i).\]
This is a Riemann sum. Taking the limit as the subinterval length goes to 0 gives 
\begin{align*}
W 	&=	\int_0^{30} 6240\pi(35-y)\dd y \\
		&=  6240\pi\left(35y-\frac{y^2}2\right)\Big|_0^{30} \\
		&= 	11,762,123 \ \text{ft-lb}\\
		&\approx  1.176 \times 10^7 \ \text{ft-lb}.
\end{align*}
\end{example}

\mtable{A simplified illustration for computing work.}{fig:pump1b}{\begin{tikzpicture}[alt={Upside-down cone apex y=-10, top radius 2 at y=0; blue slice at depth y, volume V(y)=π(2+y/5)² dy; axes x,y.},x=.1cm,y=.125cm,>=stealth,scale=.9]
\draw [->] (0,-2) -- (0,37) node [above] {\scriptsize $y$};
\draw (-1,0) node [left] {\scriptsize 0} -- (1,0)
      (-1,30) node [left] {\scriptsize 30} -- (1,30)
      (-1,35) node [left] {\scriptsize 35} -- (1,35)
      (-1,14) node [left] {\scriptsize $y$} -- (1,14);
\draw (2,14)--(4,14)  (2,35)--(4,35)
      (3,14) -- node [pos=.5,fill=white,rotate=90] {\scriptsize $35-y$} (3,35);
\begin{scope}[xscale=4,shift={(4.5,0)}]         
\draw (0,30) -- node [pos=.5,above] {\scriptsize $10$} (2.9,29.22);                             
\draw [thick](3,30) --(3,0) arc (0:-180:3) -- (-3,30);
\draw [thick] (0,30) circle (3);
\draw [thick,dashed] (3,0) arc (0:180:3);
\draw [thick,draw={\colorone},left color={\colorone},right color={\coloronefill}]
 (0,14) circle (3);
\draw [->] (5.25,16) node [above] {\scriptsize $V(y)=100\pi\dd y$} -- (2,14);
\end{scope}
\end{tikzpicture}}
% todo Tim make the simplified water tank into 3d?

We can ``streamline'' the above process a bit as we may now recognize what the important features of the problem are. \autoref{fig:pump1b} shows the tank from \autoref{ex_pump1} without the $i^{\text{th}}$ subinterval identified. Instead, we just draw a sample slice. This helps establish the height a small amount of water must travel along with the force required to move it (where the force is volume $\times$ density).

We demonstrate the concepts again in the next examples.

\begin{example}[Computing work performed: pumping fluids]\label{ex_pump2}%
A conical water tank has its top at ground level and its base 10 feet below ground. The radius of the cone at ground level is 2 ft. It is filled with water weighing 62.4 lb/ft$^3$ and is to be emptied by pumping the water to a spigot 3 feet above ground level. Find the total amount of work performed in emptying the tank.
\solution
The conical tank is sketched in \autoref{fig:pump2}. We can orient the tank in a variety of ways; we could let $y=0$ represent the base of the tank and $y=10$ represent the top of the tank, but we choose to keep the convention of the wording given in the problem and let $y=0$ represent ground level and hence $y=-10$ represents the bottom of the tank. The actual ``height'' of the water does not matter; rather, we are concerned with the distance the water travels. 

\mtable{A graph of the conical water tank in \autoref{ex_pump2}.}{fig:pump2}{\begin{tikzpicture}[alt={Upside-down cone apex y=-10, top radius 2 at y=0; blue slice at depth y, volume V(y)=π(y/5+2)² dy; axes x,y.},x=.1cm,y=.125cm,>=stealth]
\draw [->] (0,-2) -- (0,37) node [above] {\scriptsize $y$};
\draw (-1,0) node [left] {\scriptsize $-10$} -- (1,0)
      (-1,30) node [left] {\scriptsize 0} -- (1,30)
      (-1,35) node [left] {\scriptsize 3} -- (1,35)
      (-1,14) node [left] {\scriptsize $y$} -- (1,14);
\draw (2,14)--(4,14)  (2,35)--(4,35)
      (3,14) -- node [pos=.5,fill=white,rotate=90] {\scriptsize $3-y$} (3,35);
\begin{scope}[xscale=4,shift={(4.5,0)}]         
\draw (0,30) -- node [pos=.5,above] {\scriptsize $2$} (2.9,29.22);                              
\draw [thick](3,30) --(0,0) -- (-3,30);
\draw [thick] (0,30) circle (3);
\draw [thick,draw={\colorone},left color={\colorone},right color={\coloronefill}]
 (0,14) circle (1.4);
\draw [->] (3.4,6) node [below] {\scriptsize $V(y)=\pi(\frac y5+2)^2\dd y$} -- (0,14);
\end{scope}
\end{tikzpicture}}
% todo Tim make the water cone into 3d?

The figure also sketches %a differential element,
a cross-sectional circle. The radius of this circle is variable, depending on $y$. When $y=-10$, the circle has radius 0; when $y=0$, the circle has radius 2. These two points, $(-10,0)$ and $(0,2)$, allow us to find the equation of the line that gives the radius of the cross-sectional circle, which is $r(y) = y/5+2$. Hence the volume of water at this height is $V(y)=\pi(y/5+2)^2\dd y$, where $\dd y$ represents a very small height of the slice. The force required to move the water at height $y$ is $F(y) = 62.4\times V(y)$.

The distance the water at height $y$ travels is given by $h(y)=3-y$. Thus the total work done in pumping the water from the tank is 
\begin{align*}
	W
	&= \int_{-10}^0 62.4\pi(y/5+2)^2(3-y)\dd y\\
	&= 62.4\pi\int_{-10}^0\left(-\frac1{25}y^3-\frac{17}{25}y^2-\frac85y+12\right)\dd y\\
	&= 62.2\pi\cdot\frac{220}{3} %\approx 14,376
	\text{ ft-lb.}
\end{align*}
\end{example}

\begin{example}[Computing work performed: pumping fluids]\label{ex_pump3}%
A rectangular swimming pool is 20 ft wide and has a 3 ft ``shallow end'' and a 6 ft ``deep end.'' It is to have its water pumped out to a point 2 ft above the current top of the water. The cross-sectional dimensions of the water in the pool are given in \autoref{fig:pump3}(a). (Note that the ``20 ft wide'' is into the picture; the pool is 25 ft long.) Compute the amount of work performed in draining the pool.
\solution
For the purposes of this problem we choose to set $y=0$ to represent the bottom of the pool, meaning the top of the water is at $y=6$.
\mtable{The cross-section of a swimming pool filled with water in \autoref{ex_pump3} and two sample slices.}{fig:pump3}{%
 \begin{tikzpicture}[alt={Pool 25 ft: depth 6 ft to 3 ft via slope (10,0)-(15,3); rotated view 0<x<15,0<y<8 shows slices 0<y<3,3<y<6.},x=.14cm,y=.2cm,>=stealth]
  \draw [fill={\coloronefill}] (0,0)
   -- node [below] {\scriptsize 10 ft.} (10,0)
   -- node [below,yshift=-3ex] {(a)} (15,3)
   -- node [below] {\scriptsize 10 ft.} (25,3)
   -- node [right] {\scriptsize 3 ft.} (25,6)
   -- node [above] {\scriptsize 25 ft} (0,6)
   -- node [left] {\scriptsize 6 ft.} (0,0);
  \begin{scope}[yshift=-1.3in]
   \draw[>=stealth,->] (-2,-1) -- (-2,10) node [above] {\scriptsize $y$};
   \foreach \y/\x in {0/0,1.5/$y$,3/3,6/6,8/8}
   {
    \draw (-2.5,\y) node [left] {\scriptsize $\x$} -- (-1.5,\y);
   }
   \draw [thick] (0,0) --  (10,0) -- node [below,yshift=-7ex] {(b)} (15,3) --  (25,3)--  (25,6)-- (0,6) -- (0,0);
   \draw [thick,draw={\colorone}] (0,5) -- (25,5)
    (0,1.5) -- (12.5,1.5);
   \draw [fill=black] (10,0) circle (1.5pt) node [shift={(15pt,-2pt)}] {\scriptsize $(10,0)$}
    (15,3) circle (1.5pt) node [shift={(15pt,-6pt)}] {\scriptsize $(15,3)$};
   \draw[>=stealth,->] (-2,-2) -- (20,-2) node [right] {\scriptsize $x$};
   \foreach \y in {0,10,15}
   {
    \draw (\y,-2.5) node [below] {\scriptsize $\y$} -- (\y,-1.5);
   }
  \end{scope}
 \end{tikzpicture}}
\autoref{fig:pump3}(b) shows the pool oriented with this $y$-axis, along with 2 sample slices as the pool must be split into two different regions. 

The top region lies in the $y$-interval of $[3,6]$, where the length of the sample slice is $25$ ft as shown. As the pool is 20 ft wide, this sample slice of water has a volume of $V(y) = 20\cdot25\cdot\dd y$.  The water is to be pumped to a height of $y=8$, so the height function is $h(y) = 8-y$. The work done in pumping this top region of water is 
\[W_t = 62.4\int_3^6 500(8-y)\dd y = 327,600 \text{ ft-lb}.\]

The bottom region lies in the $y$-interval of $[0,3]$; we need to compute the length of the sample slice in this interval.

One end of the sample slice is at $x=0$ and the other is along the line segment joining the points $(10,0)$ and $(15,3)$. The equation of this line is $y= 3(x-10)/5$; as we will be integrating with respect to $y$, we rewrite this equation as $x=5y/3+10$. So the length of the sample slice is a difference of $x$-values: $x=0$ and $x=5y/3+10$, giving a length of $x=5y/3+10$. 

Again, as the pool is 20 ft wide, this slice of water has a volume of $V(y) = 20\cdot(5y/3+10)\cdot\dd y$; the height function is the same as before at $h(y)=8-y$. The work performed in emptying this part of the pool is
\[W_b = 62.4\int_0^3 20(5y/3+10)(8-y)\dd y = 299,520\ \text{ft-lb}.\]
The total work in emptying the pool is 
\[W = W_b+W_t = 327,600+299,520 = 627,120\ \text{ft-lb}.\]
Notice how the emptying of the bottom of the pool performs almost as much work as emptying the top. The top portion travels a shorter distance but has more water. In the end, this extra water produces more work.
\end{example}

The next section introduces one final application of the definite integral, the calculation of fluid force on a plate.

\printexercises{exercises/07-05-exercises}

\section{Fluid Forces}\label{sec:fluid_force}

In the unfortunate situation of a car driving into a body of water, the conventional wisdom is that the water pressure on the doors will quickly be so great that they will be effectively unopenable. (Survival techniques suggest immediately opening the door, rolling down or breaking the window, or waiting until the water fills up the interior at which point the pressure is equalized and the door will open. See Mythbusters episode \#72 to watch Adam Savage test these options.)

How can this be true? How much force does it take to open the door of a submerged car? In this section we will find the answer to this question by examining the forces exerted by fluids.

We start with \textbf{pressure}, which is related to \textbf{force} by the following equations:
\[
\text{Pressure} = \frac{\text{Force}}{\text{Area}}
\quad  \Leftrightarrow \quad
\text{Force} = \text{Pressure}\times\text{Area}.
\]
In the context of fluids, we have the following definition.

\begin{definition}[Fluid Pressure]\label{def:fluid_pressure}%
Let $w$ be the weight-density of a fluid. The \textbf{pressure} $p$ exerted on an object at depth $d$ in the fluid is $p = w\cdot d.$\index{fluid pressure/force}\index{integration!fluid force}
\end{definition}

We use this definition to find the \textbf{force} exerted on a horizontal sheet by considering the sheet's area.

\youtubeVideo{dNYbyKQQmQ4}{Calculating the Fluid Force on a Semi-Circle}

\mtable{A cylindrical tank in \autoref{ex_fluid1}.}{fig:fluid1a}{\begin{tikzpicture}[alt={Vertical cylinder, radius 2 ft; fluid fills full 10 ft height (0<y<10); dashed base circle and height arrow.},x=.3cm,y=.3cm,>=stealth]
\begin{scope}[xscale=2]
 \draw[thick,left color={\coloronefill},right color={\coloronefill!20}]
 (-2,10) -- (-2,0) arc (180:360:2) -- (2,10);
 \draw[thick,left color={\coloronefill!30},right color={\coloronefill}]
 (0,10)circle(2);
\draw [thick] (0,12) circle (2);
\draw [thick](-2,12)--(-2,10) (2,12)--(2,10);
\draw [dashed,thick] (-2,0) arc (180:0:2);
\draw (0,0) -- node [above,pos=.5] {\scriptsize 2 ft} (2,0);
\end{scope}
\begin{scope}[shift={(5,0)}]
\draw (-.5,0) -- (.5,0) (-.5,10) -- (.5,10)
 (0,0) -- node [pos=.5,fill=white,rotate=90] {\scriptsize 10 ft} (0,10);
\end{scope}
\end{tikzpicture}}
% todo Tim convert the cylindrical tank into a 3d image?

\begin{example}[Computing fluid force]\label{ex_fluid1}%
\mbox{}\\[-2\baselineskip]%\parbox[t]{\linewidth}{%
\begin{enumerate}
	\item	A cylindrical storage tank has a radius of 2 ft and holds 10 ft of a fluid with a weight-density of 50 lb/ft$^3$. (See \autoref{fig:fluid1a}.) What is the force exerted on the base of the cylinder by the fluid?
	\item	A rectangular tank whose base is a 5 ft square has a circular hatch at the bottom with a radius of 2 ft. The tank holds 10 ft of a fluid with a weight-density of 50 lb/ft$^3$. (See \autoref{fig:fluid1b}.) What is the force exerted on the hatch by the fluid?
%
\mtable{A rectangular tank in \autoref{ex_fluid1}.}{fig:fluid1b}{\begin{tikzpicture}[alt={5 ft × 5 ft square tank, 0<y<10 ft fluid depth; dashed ellipse (r = 2 ft) on base shows hatch below fluid column.},x=.25cm,y=.25cm,>=stealth]
 \begin{scope}[xscale=2]
  \begin{scope}[rotate=-30]
   \draw [thick](-2.5,0) node (A) {}
                (2.5,0) node (B) {} 
                (2.5,5) node (C) {};
   \draw [dashed,thick] (A.center) -- (-2.5,5) node (D) {} -- (C.center);
   \draw (0,2.5) node (O) {}
         (-1,.77) node (OC) {};
  \end{scope}
  \begin{scope}[shift={(0,10)}]
   \begin{scope}[rotate=-30,]
    \draw (-2.5,0) node (AAA) {} 
          (2.5,0) node (BBB) {}
          (2.5,5) node (CCC) {}
          (-2.5,5) node (DDD) {};
   \end{scope}
  \end{scope}
  \begin{scope}[shift={(0,12)}]
   \begin{scope}[rotate=-30,]
    \draw [thick] (-2.5,0) node (AA) {} 
                  (2.5,0) node (BB) {} 
                  (2.5,5) node (CC) {}
                  (-2.5,5) node (DD) {};
   \end{scope}
  \end{scope}
  \draw[thick,left color={\coloronefill},right color={\coloronefill!20}]
  (AAA.center) -- (A.center)-- node [pos=.5,rotate=-15,below] {\scriptsize 5 ft} (B.center)
   --node [pos=.5,rotate=45,below] {\scriptsize 5 ft}(C.center)--(CCC.center);
  \draw[thick,right color={\coloronefill},left color={\coloronefill!20}]
   (AAA.center) -- (BBB.center) -- (CCC.center)--(DDD.center) -- cycle;
  \draw [thick] (AA.center) -- (BB.center) -- (CC.center)--(DD.center) -- cycle
   (AA.center) -- (AAA.center)
   (CC.center) -- (CCC.center)
   (B.center) -- (BB.center)
   (DD.center) --(DDD.center);
  \draw [thick,dashed] (DDD.center)--(D.center)--(C.center)
   (D.center) -- (A.center)
   (O.center) circle (2)
   (O.center) -- node [pos=.5,above,rotate=15] {\scriptsize 2 ft} (OC.center);
  \begin{scope}[shift={(5.5,3)}]
   \draw (-.5,0) -- (.5,0)
         (-.5,10) -- (.5,10)
         (0,0) -- node [pos=.5,fill=white,rotate=90] {\scriptsize 10 ft} (0,10);
  \end{scope}
 \end{scope}
\end{tikzpicture}}
% todo Tim simplify this tikz? maybe convert it to 3d anyway?
%
\end{enumerate}%}
\vspace{0pt}
\solution
%
\begin{enumerate}
	\item	Using \autoref{def:fluid_pressure}, we calculate that the pressure exerted on the cylinder's base is $w\cdot d = 50 \text{ lb/ft}^3\times 10\text{ ft} = 500$ lb/ft$^2$. The area of the base is $\pi\cdot 2^2 = 4\pi$ ft$^2$. So the force exerted by the fluid is 
	\[F = 500\times 4\pi = 6283\text{ lb}.\]
Note that we effectively just computed the \emph{weight} of the fluid in the tank.

	\item	The dimensions of the tank in this problem are irrelevant. All we are concerned with are the dimensions of the hatch and the depth of the fluid. Since the dimensions of the hatch are the same as the base of the tank in the previous part of this example, as is the depth, we see that the fluid force is the same. That is, $F = 6283$ lb. 
	
	A key concept to understand here is that we are effectively measuring the weight of a 10 ft column of water above the hatch. The size of the tank holding the fluid does not matter.
\end{enumerate}
\end{example}

The previous example demonstrates that computing the force exerted on a horizontally oriented plate is relatively easy to compute. What about a vertically oriented plate? For instance, suppose we have a circular porthole located on the side of a submarine. How do we compute the fluid force exerted on it?

Pascal's Principle states  that the pressure exerted by a fluid at a depth is equal in all directions. Thus the pressure on any portion of a plate that is 1 ft below the surface of water is the same no matter how the plate is oriented. (Thus a hollow cube submerged at a great depth will not simply be ``crushed'' from above, but the sides will also crumple in. The fluid will exert force on \emph{all} sides of the cube.)

\mtable{A thin, vertically oriented plate submerged in a fluid with weight-density $w$.}{fig:fluid_intro1}{\begin{tikzpicture}[alt={Submerged trapezoid plate; wavy surface above. Red slice at depth $d_i$, width $\ell(c_i)$, thickness delta y_i.},x=1cm,y=1cm,>=stealth]
 \begin{scope}
  \draw [thick] (0,0) -- (1,2) -- (2,2) -- (4,0)--cycle;
  \draw [draw={\colortwo},fill={\colortwofill},thick] (.5,.9) rectangle (3,1.1);
  \draw (3,1) node [right] {\scriptsize $\left. \rule{0pt}{.2cm}\right\}\Delta y_i$};
  \draw (.5,.5) -- (.5,.7)
        (3,.5) -- (3,.7)
        (.5,.6) -- node [pos=.5,below,] {\scriptsize $\ell(c_i)$} (3,.6);
  \foreach \x in {0,1,2,3} {
   \begin{scope}[shift={(\x*1,2.5)}]               
    \draw [draw={\colorone},thick] (-.5,.25) parabola bend (0,0) (.5,.25);
   \end{scope}
  }
  \draw (.2,1) -- (.4,1)
        (.2,2.5) -- (.4,2.5)
        (.3,1) -- node [pos=.5,left] {\scriptsize $d_i$} (.3,2.5);
 \end{scope}
\end{tikzpicture}}

So consider a vertically oriented plate as shown in \autoref{fig:fluid_intro1} submerged in a fluid with weight-density $w$. What is the total fluid force exerted on this plate? We find this force by first approximating the force on small horizontal strips.

Let the top of the plate be at depth $b$ and let the bottom be at depth $a$. (For now we assume that surface of the fluid is at depth 0, so if the bottom of the plate is 3 ft under the surface, we have $a=-3$. We will come back to this later.) We partition the interval $[a,b]$ into $n$ subintervals 
\[a=y_0<y_1<\dotso<y_n=b,\]
with the $i\,^\text{th}$ subinterval having length $\Delta y_i$. The force $F_i$ exerted on the plate in the $i\,^\text{th}$ subinterval is $F_i = \text{Pressure}\times \text{Area}.$

The pressure is depth $\times w$. We approximate the depth of this thin strip by choosing any value $d_i$ in $[y_i,y_{i+1}]$; the depth is approximately $-d_i$. (Our convention has $d_i$ being a negative number, so $-d_i$ is positive.) For convenience, we let $d_i$ be an endpoint of the subinterval; we let $d_i = y_i$. 

The area of the thin strip is approximately length $\times$\ width. The width is $\Delta y_i$. The length is a function of some $y$-value $c_i$ in the $i\,^\text{th}$ subinterval. We state the length is $\ell(c_i)$. Thus 
\begin{align*}
	F_i	&= \text{Pressure} \times \text{Area} \\
		&=	-y_i\cdot w \times \ell(c_i)\cdot\Delta y_i.
\end{align*}
To approximate the total force, we add up the approximate forces on each of the $n$ thin strips:
\[
F = \sum_{i=1}^n F_i \approx \sum_{i=1}^n -w\cdot y_i\cdot\ell(c_i)\cdot\Delta y_i.
\]
This is, of course, another Riemann Sum. We can find the exact force by taking a limit as the subinterval lengths go to 0; we evaluate this limit with a definite integral.

\begin{keyidea}[Fluid Force on a Vertically Oriented Plate]\label{idea:fluid_force}%
Let a vertically oriented plate be submerged in a fluid with weight-density $w$ where the top of the plate is at $y=b$ and the bottom is at $y=a$. Let $\ell(y)$ be the length of the plate at $y$.
\index{fluid pressure/force}\index{integration!fluid force}
\begin{enumerate}
	\item	If $y=0$ corresponds to the surface of the fluid, then the force exerted on the plate by the fluid is
	\[F=\int_a^b w\cdot(-y)\cdot\ell(y)\dd y.\]
	\item	In general, let $d(y)$ represent the distance between the surface of the fluid and the plate at $y$. Then the force exerted on the plate by the fluid is 
	\[F=\int_a^b w\cdot d(y)\cdot\ell(y)\dd y.\]
\end{enumerate}
\end{keyidea}

\begin{example}[Finding fluid force]\label{ex_fluid2}%
Consider a thin plate in the shape of an
isosceles triangle as shown in \autoref{fig:fluid2a} submerged in water with a
%
\mtable{A thin plate in the shape of an isosceles triangle in \autoref{ex_fluid2}.}{fig:fluid2a}{\begin{tikzpicture}[alt={An upside-down isosceles triangle; top edge spans 4ft, overall height is 4ft.},x=.4cm,y=.4cm,>=stealth]
\draw[thick] (0,0) -- (2,4)
 -- node [pos=.5,above] {\scriptsize 4 ft} (-2,4) -- cycle;
\draw (2.2,0) -- (2.6,0)
      (2.2,4) -- (2.6,4)
      (2.4,0) -- node [pos=.5,rotate=-90,above] {\scriptsize 4 ft} (2.4,4);
\end{tikzpicture}}
%
weight-density of 62.4 lb/ft$^3$. If the bottom of the plate is 10 ft below the surface of the water, what is the total fluid force exerted on this plate?
\solution
We approach this problem in two different ways to illustrate the different ways \autoref{idea:fluid_force} can be implemented. First we will let $y=0$ represent the surface of the water, then we will consider an alternate convention.

	\mtable{Sketching the triangular plate in \autoref{ex_fluid2} with the convention that the water level is at $y=0$.}{fig:fluid2b}{\begin{tikzpicture}[alt={Axes shown; water at y=0. Inverted plate joins (-2,-6) to (2,-6) to apex (0,-10). Red sample strip lies −8<y<−6.}]  
\begin{axis}[width=\marginparwidth, axis y line=middle,axis x line=middle,
tick label style={font=\scriptsize},ymin=-11,ymax=2,xmin=-3,xmax=5,axis on top,
xtick={-2,-1,1,2},ytick={-10,-8,-4,-2},height=1.625\marginparwidth,
xlabel={\scriptsize$x$},ylabel={\scriptsize$y$},xlabel style=right,ylabel style=above]
 \draw[thick] (axis cs:0,-10) -- (axis cs:2,-6) --  (axis cs:-2,-6) -- cycle;
 \draw [fill=black] (axis cs:2,-6) circle (1pt) node [above] {\scriptsize $(2,-6)$}
  (axis cs:-2,-6)circle (1pt) node [above] {\scriptsize $(-2,-6)$}
  (axis cs:0,-10)circle (1pt);
 \draw [draw={\colortwo},thick]  (axis cs:-1.5,-7) -- (axis cs:1.5,-7) node [right,black] {\scriptsize $y$};
 \foreach \x in {-2,-1,0,1,2} {
  \edef\temp{\noexpand\draw[draw={\colorone},thick]
  (axis cs:{\x-.5},.25) parabola bend (axis cs:\x,0) (axis cs:{\x+.5},.25);}\temp
 }
 \draw(axis cs:2,1) node {\scriptsize water line};
 \draw(axis cs:4.1,-7) -- (axis cs:4.5,-7)
      (axis cs:4.1,0) -- (axis cs:4.5,0)
      (axis cs:4.3,-7) -- node [pos=.5,rotate=-90,above] {\scriptsize $d(y) = 10-y$} (axis cs:4.3,0);
\end{axis}
\end{tikzpicture}}

\begin{enumerate}
	\item	We let $y=0$ represent the surface of the water; therefore the bottom of the plate is at $y=-10$. We center the triangle on the $y$-axis as shown in \autoref{fig:fluid2b}. The depth of the plate at $y$ is $-y$ as indicated by the Key Idea. We now consider the length of the plate at $y$.

	We need to find equations of the left and right edges of the plate. The right hand side is a line that connects the points $(0,-10)$ and $(2,-6)$: that line has equation $x=\frac12(y+10)$. (Find the equation in the familiar $y=mx+b$ format and solve for $x$.) Likewise, the left hand side is described by the line $x=-\frac12(y+10)$. The total length is the distance between these two lines: $\ell(y)=\frac12(y+10) - (-\frac12(y+10)) = y+10.$
	
The total fluid force is then:
\begin{align*}
	F 	&=	\int_{-10}^{-6} 62.4(-y)(y+10)\dd y \\
		&=	62.4\cdot \frac{176}{3} \approx	3660.8\text{ lb}.
\end{align*}	

	\item	Sometimes it seems easier to orient the thin plate nearer the origin. For instance, consider the convention that the bottom of the triangular plate is at $(0,0)$, as shown in \autoref{fig:fluid2c}. The equations of the left and right hand sides are easy to find. They are $y=2x$ and $y=-2x$, respectively, which we rewrite as $x=y/2$ and $x=-y/2$. Thus the length function is $\ell(y) =y/2-(-y/2) = y$. 

\mtable{Sketching the triangular plate in \autoref{ex_fluid2} with the convention that the base of the triangle is at $(0,0)$.}{fig:fluid2c}{\begin{tikzpicture}[alt={Water surface at y=10. Inverted triangle apex (0,0), base (−2,4)–(2,4). Red slice at y, thickness dy. Bracket labels depth 10−y.}]
\begin{axis}[width=\marginparwidth, axis y line=middle,axis x line=middle,
tick label style={font=\scriptsize},ymin=-1,ymax=12,xmin=-3,xmax=5,axis on top,
xtick={-2,-1,1,2},ytick={2,6,8,10},height=1.625\marginparwidth, % 13/8
xlabel={\scriptsize$x$},ylabel={\scriptsize$y$},xlabel style=right,ylabel style=above]
 \draw[thick] (axis cs:0,0) -- (axis cs:2,4) --  (axis cs:-2,4) -- cycle;
 \draw [fill=black] (axis cs:2,4) circle (1pt) node [above] {\scriptsize $(2,4)$}
  (axis cs:-2,4)circle (1pt) node [above] {\scriptsize $(-2,4)$}
  (axis cs:0,0)circle (1pt);
 \draw [draw={\colortwo},thick]  (axis cs:-1.5,3) -- (axis cs:1.5,3) node [right,black] {\scriptsize $y$};
 \foreach \x in {-2,-1,0,1,2} {
  \edef\temp{\noexpand\draw[draw={\colorone},thick]
  (axis cs:{\x-.5},10.25) parabola bend (axis cs:\x,10) (axis cs:{\x+.5},10.25);}\temp
 }
 \draw(axis cs:2,11) node {\scriptsize water line};
 \draw(axis cs:4.1,3) -- (axis cs:4.5,3)
      (axis cs:4.1,10) -- (axis cs:4.5,10)
      (axis cs:4.3,3) -- node [pos=.5,rotate=-90,above] {\scriptsize $d(y) = 10-y$} (axis cs:4.3,10);
\end{axis}
\end{tikzpicture}}

	As the surface of the water is 10 ft above the base of the plate, we have that the surface of the water is at $y=10$. Thus the depth function is the distance between $y=10$ and $y$; $d(y) = 10-y$. We compute the total fluid force as:
	\begin{align*}
	F	&=\int_0^4 62.4(10-y)(y)\dd y \\
		&\approx 3660.8\text{ lb}.
	\end{align*}
\end{enumerate}
The correct answer is, of course, independent of the placement of the plate in the coordinate plane as long as we are consistent.
\end{example}

\begin{example}[Finding fluid force]\label{ex_fluid3}%
Find the total fluid force on a car door submerged up to the bottom of its window in water, where the car door is a rectangle 40'' long and 27'' high (based on the dimensions of a 2005 Fiat Grande Punto).
\solution
The car door, as a rectangle, is drawn in \autoref{fig:fluid3}. Its length is $10/3$ ft and its height is 2.25 ft. We adopt the convention that the top of the door is at the surface of the water, both of which are at $y=0$. Using the weight-density of water of 62.4 lb/ft$^3$, we have the total force as
%
\mtable{Sketching a submerged car door in \autoref{ex_fluid3}.}{fig:fluid3}{\begin{tikzpicture}[alt={Water at y=0. Door is rectangle 0<x<10/3, −2.25<y<0. Red horizontal slice at depth y, width 10/3, thickness dy.}
,x=.5cm,y=.5cm,>=stealth]
 \draw[thick] (0,0) -- (3.33,0) --  (3.33,-2.25) -- (0,-2.25) --cycle;
 \draw [fill=black]      (3.33,0) circle (1pt)
  node [above right] {\scriptsize $(3.\overline{3},0)$}
  (3.33,-2.25)circle (1pt) node [below right] {\scriptsize $(3.\overline{3},-2.25)$}
  (0,-2.25)circle (1pt) node [below left] {\scriptsize $(0,-2.25)$};
 \draw [fill=black]      (0,0) circle (1pt) node [above left] {\scriptsize $(0,0)$};
 \draw [draw={\colortwo},thick]  (0,-1.5) -- (3.33,-1.5)node[right,black]{\scriptsize$y$};
 \draw [->] (0,-3) -- (0,1) node [above] {\scriptsize $y$};
 \draw [->] (-1,0) -- (6,0) node [right] {\scriptsize $x$};
 \foreach \x in {-1,0,1,2,3,4} {
  \begin{scope}[shift={(\x*1,0)}]         
   \draw [draw={\colorone},thick] (-.5,.25) parabola bend (0,0) (.5,.25);
  \end{scope}
 }
\end{tikzpicture}}
%
\begin{align*}
	F &=	\int_{-2.25}^0 62.4(-y)10/3\dd y \\
	&= 	\int_{-2.25}^0 -208y\dd y\\
	&= -104y^2\Big|_{-2.25}^0 \\
	&=	526.5 \text{ lb.}
\end{align*}
Most adults would find it very difficult to apply over 500 lb of force to a car door while seated inside, making the door effectively impossible to open. This  is  counter-intuitive as most assume that the door would be relatively easy to open. The truth is that it is not, hence the survival tips mentioned at the beginning of this section.
\end{example}

\mtable{Measuring the fluid force on an underwater porthole in \autoref{ex_fluid4}.}{fig:fluid4}{\begin{tikzpicture}[alt={Circle radius 1.5 ft at origin; waterline y=50 ft; depth d(y)=50−y; red horizontal slice at y, width ℓ(y) dy, spans −1.5<x<1.5.},x=.6cm,y=.6cm,>=stealth]
 \draw[thick] (0,0) circle (1.5);
 \draw [draw={\colortwo},thick] (-1.3,.75) -- (1.3,.75) node[right,black]{\scriptsize$y$};
 \draw [->] (0,-2.2) -- (0,3)--(-.25,3.25) -- (0,3.5) -- (.25,3.75) -- (0,4)
  --(0,6) node [above] {\scriptsize $y$};
 \draw [->] (-2.5,0) -- (2.5,0) node [right] {\scriptsize $x$};
 \foreach \x in {-2,-1,1,2} {
  \draw (\x,0.2) -- (\x,-.2) node [below] {\scriptsize $\x$};
 }
 \foreach \x in {-2,-1,1,2} {
  \draw (.2,\x) -- (-.2,\x) node [left] {\scriptsize $\x$};
 }
 \draw (.2,5) -- (-.2,5) node [left] {\scriptsize $50$};
 \foreach \x in {-1.5,-0.5,0.5,1.5} {
  \draw[draw={\colorone},thick,shift={(\x,5)}](-.5,.25)parabola bend(0,0) (.5,.25);
 }
 \draw (2,6) node {\scriptsize water line}
       (2,-2.5) node {\scriptsize \emph{not to scale}};
 \draw (2.1,.75) -- (2.5,.75)
       (2.1,5) -- (2.5,5)
       (2.3,.75) -- node [pos=.5,rotate=-90,above] {\scriptsize $d(y)=50-y$}(2.3,5);
\end{tikzpicture}}

\begin{example}[Finding fluid force]\label{ex_fluid4}%
An underwater observation tower is being built with circular viewing portholes enabling visitors to see underwater life. Each vertically oriented porthole is to have a 3 ft diameter whose center is to be located 50 ft underwater. Find the total fluid force exerted on each porthole. Also, compute the fluid force on a horizontally oriented porthole that is under 50 ft of water.
\solution
We place the center of the porthole at the origin, meaning the surface of the water is at $y=50$ and the depth function will be $d(y)=50-y$; see \autoref{fig:fluid4} 

The equation of a circle with a radius of 1.5 is $x^2+y^2=2.25$; solving for $x$ we have $x=\pm \sqrt{2.25-y^2}$, where the positive square root corresponds to the right side of the circle and the negative square root corresponds to the left side of the circle. Thus the length function at depth $y$ is $\ell(y) = 2\sqrt{2.25-y^2}$. Integrating on $[-1.5,1.5]$ we have:
\begin{align*}
	F
	&= 62.4\int_{-1.5}^{1.5} 2(50-y)\sqrt{2.25-y^2}\dd y \\
	&= 62.4\int_{-1.5}^{1.5} \bigl(100\sqrt{2.25-y^2} - 2y\sqrt{2.25-y^2}\bigr)\dd y \\
	&= 6240\int_{-1.5}^{1.5} \bigl(\sqrt{2.25-y^2}\bigr)\dd y - 62.4\int_{-1.5}^{1.5} \bigl(2y\sqrt{2.25-y^2}\bigr)\dd y.
\end{align*}			
The second integral above can be evaluated using Substitution. Let $u=2.25-y^2$ with $\dd u = -2y\dd y$. The new bounds are: $u(-1.5)=0$ and $u(1.5)=0$; the new integral will integrate from $u=0$ to $u=0$, hence the integral is 0.

The first integral above finds the area of half a circle of radius 1.5, thus the first integral evaluates to $6240\cdot\pi\cdot1.5^2/2 = 22,054$. Thus the total fluid force on a vertically oriented porthole is $22,054$ lb.

Finding the force on a horizontally oriented porthole is more straightforward:
\[
F = \text{Pressure}\times\text{Area} = 62.4\cdot50\times \pi\cdot1.5^2
 = 22,054\text{ lb}.
\]
That these two forces are equal is not coincidental; it turns out that the fluid force applied to a vertically oriented circle whose center is at depth $d$ is the same as force applied to a horizontally oriented circle at depth $d$.
\end{example}

We end this chapter with a reminder of the true skills meant to be developed here. We are not truly concerned with an ability to find fluid forces or the volumes of solids of revolution. Work done by a variable force is important, though measuring the work done in pulling a rope up a cliff is probably not.

What we are actually concerned with is the ability to solve certain problems by first approximating the solution, then refining the approximation, then recognizing if/when this refining process results in a definite integral through a limit. Knowing the formulas found inside the special boxes within this chapter is beneficial as it helps solve problems found in the exercises, and other mathematical skills are strengthened by properly applying these formulas. However, more importantly, understand how each of these formulas was constructed. Each is the result of a summation of approximations; each summation was a Riemann sum, allowing us to take a limit and find the exact answer through a definite integral. %\bigskip

%The next chapter addresses an entirely different topic: sequences and series. In short, a sequence is a list of numbers, where a series is the summation of a list of numbers. These seemingly-simple ideas lead to very powerful mathematics.

\printexercises{exercises/07-06-exercises}


}{}

\iftoggle{isEarlyTrans}{}{%
\apexchapter{Inverse Functions and L'Hôpital's Rule}{chapter:diff_conc}
% todo write a Inverse Functions prerequisite section reviewing logs

% todo write a better chapter intro for Inverse Functions
This chapter completes our differentiation toolkit.  The first and most important tool will be how to differentiate inverse functions.  We'll be able to use this to differentiate exponential and logarithmic functions, which we stated in \autoref{thm:deriv_common} but did not prove.
\section{Inverse Functions}\label{sec:inv_funcs}

We say that two functions $f$ and $g$ are \emph{inverses} if $g(f(x))=x$ for all $x$ in the domain of $f$ and $f(g(x))=x$ for all $x$ in the domain of $g$. A function can only have an inverse if it is one-to-one, i.e. if we never have $f(x_1)=f(x_2)$ for different elements $x_1$ and $x_2$ of the domain. This is equivalent to saying that the graph of the function passes the horizontal line test.
%Functions that are not one-to-one may sometimes have inverses on part of their domains.
%
%If $f$ and $g$ are inverses, the domain of $g$ will be the range of $f$ and the range of $g$ will be the domain of $f$. The graphs of $f$ and $g$ will be reflections of each other across the line $y=x$ since $y=f(x)$ if and only if $x=g(y)$ (since the point $(y,x)$ is on the graph of $g$ whenever $(x,y)$ is on the graph of $f$.)
The inverse of $f$ is denoted $f^{-1}$, which should not be confused with the function $1/f(x)$.\index{one-to-one}

\begin{keyidea}[Inverse Functions]\label{ki_inv_funcs}%
For a one-to–one function $f$,
\begin{itemize}
\item The domain of $f^{-1}$ is the range of $f$; the range of $f^{-1}$ is the domain of $f$.
\item $f^{-1}(f(x))=x$ for all $x$ in the domain of $f$.
\item $f(f^{-1}(x))=x$ for all $x$ in the domain of $f^{-1}$.
\item The graph of $y=f^{-1}(x)$ is the reflection across $y=x$ of the graph of $y=f(x)$.
\item $y=f^{-1}(x)$ if and only if $f(y)=x$ and $y$ is in the domain of $f$.
\end{itemize}
\end{keyidea}

\youtubeVideo{BmjbDINGZGg}{Finding the Inverse of a Function or Showing One Does not Exist, Ex 3}

To determine whether or not $f$ and $g$ are inverses for each other, we check to see whether or not $g(f(x))=x$ for all $x$ in the domain of $f$,and $f(g(x))=x$ for all $x$ in the domain of $g$.

\mtable{A function $f$ along with its inverse $f\primeskip^{-1}$. (Note how it does not matter which function we refer to as $f$; the other is $f\primeskip^{-1}$.)}{fig_inverse1}{\begin{tikzpicture}[alt={Blue curve f and red inverse reflect over dashed y=x; blue passes (-0.5,0.375) to (1,1.5); red passes (0.375,-0.5) to (1.5,1).}]
 \begin{axis}[width=1.16\marginparwidth,axis equal,tick label style={font=\scriptsize},
   minor x tick num=1,minor y tick num=1,axis y line=middle,axis x line=middle,
   ymin=-1.2,ymax=1.7,xmin=-1.2,xmax=1.7,name=myplot]
  \addplot[draw={\colorone},smooth,thick,domain=-1:1.05] (x,x^3+.5);
  \addplot[draw={\colortwo},smooth,thick,domain=-1:1.05] (x^3+.5,x);
  \addplot[dashed,thin] {x};
  \filldraw[black] (axis cs:-.5,.375)
   node [shift={(-10pt,10pt)}] {\tiny $(-0.5,0.375)$} circle (1pt);
  \filldraw[black] (axis cs:.375,-.5)
   node [shift={(10pt,-10pt)}] {\tiny $(0.375,-0.5)$} circle (1pt);
  \filldraw[black] (axis cs:1,1.5) node [left] {\tiny $(1,1.5)$} circle (1pt);
  \filldraw[black] (axis cs:1.5,1) node [below] {\tiny $(1.5,1)$} circle (1pt);
 \end{axis}
 \node[right] at (myplot.right of origin) {\scriptsize $x$};
 \node[above] at (myplot.above origin) {\scriptsize $y$};
\end{tikzpicture}}

\begin{example}[Verifying Inverses]\label{ex_inv_determine}%
Determine whether or not the following pairs of functions are inverses:
\begin{enumerate}
\item $f(x)=3x+1$; $g(x)=\dfrac{x-1}3$
\item $f(x)=x^3+1$; $g(x)=x^{1/3}-1$
\end{enumerate}
\solution
\begin{enumerate}
\item To check the composition we plug $f(x)$ in for $x$ in the definition of $g$ as follows:\vspace{-.3\baselineskip}
\[ g(f(x))=\frac{f(x)-1}3=\frac{(3x+1)-1}3=\frac{3x}3=x\]
So $g(f(x))=x$ for all $x$ in the domain of $f$. Likewise, you can check that $f(g(x))=x$ for all $x$ in the domain of $g$, so $f$ and $g$ are inverses.
\item If we try to proceed as before, we find that:
\[g(f(x))=(f(x))^{1/3}-1=(x^3+1)^{1/3}-1\]
This doesn't seem to be the same as the identity function $x$. To verify this, we find a number $a$ in the domain of $f$ and show that $g(f(a))\neq a$ for that value. Let's try $x=1$. Since $f(1)=1^3+1=2$, we find that $g(f(1))=g(2)=2^{1/3}-1\approx 0.26$. Since $g(f(1))\neq 1$, these functions are not inverses.
\end{enumerate}
\end{example}

%To find the inverse $f^{-1}$ of a given function $y=f(x)$, we follow these steps (when possible):
%\begin{enumerate}
%\item In the equation defining $f$, solve for $x$.
%\item Interchange $x$ and $y$.
%\item The equation now defines $y=f^{-1}(x)$.
%\end{enumerate}
%
%\begin{example}[Finding Inverses]\label{ex_inv_find}%
%Find the inverses of the following functions.
%\[\text{1.}\quad f(x)=5x+3\qquad\qquad\text{2.}\quad g(x)=x^3+1.\]
%\solution\begin{enumerate}
%\item Start with the equation $y=5x+3$ and solve for $x$:
%\begin{align*}
%y&=5x+3\\
%y-3&=5x\\
%\frac{y-3}5 &=x
%\end{align*}
%Next interchange $x$ and $y$ to get $y=\dfrac{x-3}5$. The inverse of $f$ is:
%\[ f^{-1}(x)=\frac{x-3}5 \]
%\item Proceeding as before, we get:
%\begin{align*}
%y&=x^3+1\\
%y-1&=x^3\\
%(y-1)^{1/3}&=x
%\end{align*}
%Interchanging $x$ and $y$ yields $y=(x-1)^{1/3}$, so
%\[g^{-1}(x)=(x-1)^{1/3}.\]
%\end{enumerate}
%\end{example}


\subsection{Functions that are not one-to-one.}

\mtable{The function $f(x)=x^2$ is not one-to-one.}{fig_parab_oto}{\begin{tikzpicture}[alt={Blue parabola y = x² spanning −2 < x < 2; vertex at (0,0); symmetric points (−2,4) and (2,4) show not one-to-one.}]
 \begin{axis}[width=\marginparwidth,
   tick label style={font=\scriptsize},axis y line=middle,axis x line=middle,
   ymin=-1,ymax=5,xmin=-2.5,xmax=2.5,xtick={-2,2},ytick={},name=myplot]
  \addplot [draw={\colorone},smooth,thick,domain=-2.5:2.5] (x,{x^2});
  \draw (axis cs:-2,4)node[above right]{\scriptsize$(-2,4)$}
      --(axis cs:2,4)node[above left]{\scriptsize$(2,4)$};
 \end{axis}
 \node [right] at (myplot.right of origin) {\scriptsize $x$};
 \node [above] at (myplot.above origin) {\scriptsize $y$};
\end{tikzpicture}}
Unfortunately, not every function we would like to find an inverse for is one-to-one. For example, the function $f(x)=x^2$ is not one-to-one because $f(-2)=f(2)=4$. If $f^{-1}$ is an inverse for $f$, then $f^{-1}(f(-2))=-2$ implies that $f^{-1}(4)=-2$. On the other hand, $f^{-1}(f(2))=2$, so $f^{-1}(4)=2$. We cannot have it both ways if $f^{-1}$ is a function, so no such inverse exists. We can find a partial solution to this dilemma by restricting the domain of $f$. There are many possible choices, but traditionally we restrict the domain to the interval $[0,\infty)$. The function $f^{-1}(x)=\sqrt x$ is now an inverse for this restricted version of $f$.


\subsection{The inverse sine function}

We consider the function $f(x)=\sin x$, which is not one-to-one. A piece of the graph of $f$ is in \autoref{fig_sin_arc}(a). In order to find an appropriate restriction of the domain of $f$, we look for consecutive critical points where $f$ takes on its minimum and maximum values. In this case, we use the interval $[-\pi/2,\pi/2]$. We define the inverse of $f$ on this restricted range by $y=\sin^{-1}x$ if and only if $\sin y=x$ and $-\pi/2\leq y\leq \pi/2$. The graph is a reflection of the graph of $g$ across the line $y=x$, as seen in \autoref{fig_sin_arc}(b).

\noindent\begin{minipage}[t]{\linewidth}\noindent%
\captionsetup{type=figure}%
\centering
\tagpdfsetup{table/header-rows={2}}
\begin{tabular}{cc}
\begin{tikzpicture}[alt={Blue sine curve on −π < x < π; crosses origin, peaks y = 1 at x = π/2, trough y = −1 at x = −π/2.},baseline={(current bounding box.center)}]
 \begin{axis}[x=.7cm, y=.7cm,
   tick label style={font=\scriptsize},axis y line=middle,axis x line=middle,
   ymin=-1.6,ymax=1.6,xmin=-3.3,xmax=3.3,xtick={-3.14,-1.57,1.57,3.14},
   xticklabels={$-\pi$,$-\frac\pi2$,$\frac\pi2$,$\pi$},name=myplot]
  \addplot [draw={\colorone},smooth,thick,domain=-3.14:3.14] (x,{sin(deg(x))});
 \end{axis}
 \node[anchor=base] at (myplot.origin) {};
 \node [right] at (myplot.right of origin) {\scriptsize $x$};
 \node [above] at (myplot.above origin) {\scriptsize $y$};
\end{tikzpicture}
&
\begin{tikzpicture}[alt={Blue one-to-one arc of sin x (−π/2 < x < π/2) and red inverse sine mirrored across dashed y = x line.},baseline={(current bounding box.center)}]
 \begin{axis}[x=.7cm,y=.7cm,
   tick label style={font=\scriptsize},axis y line=middle,axis x line=middle,
   ymin=-2,ymax=2,xmin=-2,xmax=2,name=myplot,
   xtick={-1.57,-1,1,1.57},xticklabels={$-\frac\pi2$,$-1$,$1$,$\frac\pi2$},
   ytick={-1.57,-1,1,1.57},yticklabels={$-\frac\pi2$,$-1$,$1$,$\frac\pi2$}]
  \addplot [draw={\colorone},smooth,thick,domain=-1.57:1.57] (x,{sin(deg(x))})
   node[pos=.8,below right]{$\sin x$};
  \addplot [draw={\colortwo},smooth,thick,domain=-1.57:1.57] ({sin(deg(x))},x)
   node[pos=.95,above]{$\sin^{-1}x$};
  \addplot[dashed,thin] {x};
 \end{axis}
 \node[anchor=base] at (myplot.origin) {};
 \node [right] at (myplot.right of origin) {\scriptsize $x$};
 \node [above] at (myplot.above origin) {\scriptsize $y$};
\end{tikzpicture}
\\ (a) & (b)
\end{tabular}
\caption{(a) A portion of $y=\sin x$. (b) A one-to-one portion of $y=\sin x$ along with $y=\sin^{-1}x$.}
\label{fig_sin_arc}
\end{minipage}

\subsection{The inverse tangent function}

Next we consider the function $f(x)=\tan x$, which is also not one-to-one. A piece of the graph of $f$ is given in \autoref{fig_tan_arc}(a).  In order to find an interval on which the function is one-to-one and on which the function takes on all values in the range, we use an interval between consecutive vertical asymptotes. Traditionally, the interval $(-\pi/2,\pi/2)$ is chosen. Note that we choose the open interval in this case because the function $f$ is not defined at the endpoints. So we define $y=\tan^{-1} x$ if and only if $\tan y=x$ and $-\pi/2< y<\pi/2$. The graph of $y=\tan^{-1} x$ is shown in \autoref{fig_tan_arc}(b). Also note that the vertical asymptotes of the original function are reflected to become horizontal asymptotes of the inverse function.

\noindent\begin{minipage}[t]{\linewidth}\noindent%
\captionsetup{type=figure}%
\centering
\tagpdfsetup{table/header-rows={2}}
\begin{tabular}{cc}
\begin{tikzpicture}[alt={Blue tan curves on −3π/2 < x < −π/2, −π/2 < x < π/2, π/2 < x < 3π/2 with vertical asymptotes and y-axis center.},baseline={(current bounding box.center)}]
 \begin{axis}[x=.5cm,y=.5cm,
   tick label style={font=\scriptsize},axis y line=middle,axis x line=middle,
   ymin=-3,ymax=3,xmin=-5,xmax=5,xtick={-4.71,-3.14,-1.57,1.57,3.14,4.71},
   xticklabels={$-\frac{3\pi}2$,$-\pi$,$-\frac\pi2$,$\frac\pi2$,$\pi$,$\frac{3\pi}2$},
   name=myplot]
  \addplot [draw={\colorone},smooth,thick,domain=-1.5:1.5] (x,{tan(deg(x))});
  \addplot [draw={\colorone},smooth,thick,domain=-1.5:1.5] (x+3.14,{tan(deg(x))});
  \addplot [draw={\colorone},smooth,thick,domain=-1.5:1.5] (x-3.14,{tan(deg(x))});
 \end{axis}
 \node [right] at (myplot.right of origin) {\scriptsize $x$};
 \node [above] at (myplot.above origin) {\scriptsize $y$};
\end{tikzpicture}
&
\begin{tikzpicture}[alt={Blue tan arc on −π/2 < x < π/2 and red tan^(-1) x reflected across dashed line y = x, showing inverse relation.},baseline={(current bounding box.center)}]
 \begin{axis}[x=.5cm,y=.5cm,
   tick label style={font=\scriptsize},axis y line=middle,axis x line=middle,
   ymin=-3,ymax=3,xmin=-3,xmax=3,name=myplot,
   xtick={-1.57,1.57},xticklabels={$-\frac\pi2$,$\frac\pi2$},
   ytick={-1.57,1.57},yticklabels={$-\frac\pi2$,$\frac\pi2$}]
  \addplot [draw={\colorone},smooth,thick,domain=-1.3:1.3] (x,{tan(deg(x))})
   node[pos=.8,above right]{$\tan x$};
  \addplot [draw={\colortwo},smooth,thick,domain=-1.3:1.3] ({tan(deg(x))},x)
   node[pos=.6,below right,yshift=1.3ex]{$\tan^{-1}x$};
  \addplot[dashed,thin] {x};
 \end{axis}
 \node [right] at (myplot.right of origin) {\scriptsize $x$};
 \node [above] at (myplot.above origin) {\scriptsize $y$};
\end{tikzpicture}
\\ (a) & (b)
\end{tabular}
\caption{(a) A portion of $y=\tan x$. (b) A one-to-one portion of $y=\tan x$ along with $y=\tan^{-1}$.}
\label{fig_tan_arc}
\end{minipage}

The other inverse trigonometric functions are defined in a similar fashion. The resulting domains and ranges are summarized in \autoref{fig:domain_trig}.\bigskip

\noindent\begin{minipage}[t]{\linewidth}\noindent%
\captionsetup{type=figure}%
\small%
\flushinner{%
\tagpdfsetup{table/header-rows={1}}
\begin{tabular}{ c c c @{\hspace{2em}} c c c }
Function & \parbox[b]{40pt}{\centering Restricted Domain} & Range &\parbox[b]{40pt}{\centering Inverse Function} & Domain & Range\\ \cmidrule(r{2em}){1-3}\cmidrule(l{-1em}){4-6}
$\sin x$ & $[-\pi/2,\pi/2]$ & $[-1,1]$&$\sin^{-1}x$ & $[-1,1]$ & $[-\pi/2,\pi/2]$ \\
$\cos x$ & $[0,\pi]$ & $[-1,1]$&$\cos^{-1}x$ & $[-1,1]$ & $[0,\pi]$ \\
$\tan x$ & $(-\pi/2,\pi/2)$ & $(-\infty,\infty)$ & $\tan^{-1}x$ & $(-\infty,\infty)$ & $(-\pi/2,\pi/2)$	\\
$\csc x$ & $[-\pi/2,0)\cup (0, \pi/2]$ & $(-\infty,-1]\cup [1,\infty)$&$\csc^{-1} x$ & $(-\infty,-1]\cup [1,\infty)$ & $[-\pi/2,0)\cup (0, \pi/2]$  \\
$\sec x$ & $[0,\pi/2)\cup (\pi/2,\pi]$ & $(-\infty,-1]\cup [1,\infty)$&$\sec^{-1}x$ & $(-\infty,-1]\cup [1,\infty)$ & $[0,\pi/2)\cup (\pi/2,\pi]$ \\
$\cot x$ & $(0,\pi)$ & $(-\infty,\infty)$&$\cot^{-1}x$ &  $ (-\infty,\infty)$ & $(0,\pi)$	
\end{tabular}}
\caption{Domains and ranges of the trigonometric and inverse trigonometric functions.}
\label{fig:domain_trig}
\end{minipage}

\begin{example}[Evaluating Inverse Trigonometric Functions]\label{eg_inv_trig}%
Find exact values for the following:
\mnote{Sometimes, $\arcsin$ is used instead of $\sin^{-1}$.  Similar ``arc'' functions are used for the other inverse trigonometric functions as well.}
\\
\begin{minipage}{.5\linewidth}
\begin{enumerate}
 \item $\tan^{-1}(1)$
 \item $\cos(\sin^{-1}(\sqrt3/2))$
\end{enumerate}
\end{minipage}%
\begin{minipage}{.5\linewidth}
\begin{enumerate}
 \setcounter{enumi}{2}
 \item $\sin^{-1}(\sin(7\pi/6))$
 \item $\tan(\cos^{-1}(11/15))$
\end{enumerate}
\end{minipage}
\solution
\begin{enumerate}
\item $\tan^{-1}(1)=\pi/4$
\item $\cos(\sin^{-1}(\sqrt3/2))=\cos(\pi/3)=1/2$
\item Since $7\pi/6$ is not in the range of the inverse sine function, we should be careful with this one.
\[\sin^{-1}(\sin(7\pi/6))=\sin^{-1}(-1/2)=-\pi/6.\]
\item Since we don't know the value of $\cos^{-1}(11/15)$, we let $\theta$ stand for this value. We know that $\theta$ is an angle between $0$ and $\pi$ and that $\cos(\theta)=11/15$.
In \autoref{fig_inv_trig}, we use this information to construct a right triangle with angle $\theta$, where the adjacent side over the hypotenuse must equal 11/15. Applying the Pythagorean Theorem we find that
\mtable{A right triangle for the situation in \autoref{eg_inv_trig} (4).}{fig_inv_trig}{%
\begin{tikzpicture}[alt={Right triangle: base 11, hypotenuse 15, unknown vertical side y, angle θ at left base.},x=1ex,y=1ex]
 \draw(0,0)
  --node[pos=.1,above right]{$\theta$}node[pos=.6,below]{11}(11,0)
  --node[pos=.5,right]{$y$}(11,15)
  --node[pos=.5,above left]{15}cycle;
\end{tikzpicture}}
\[y=\sqrt{15^2-11^2}=\sqrt{104}=2\sqrt{26}.\]
Finally, we have:
\[\tan(\cos^{-1}(11/15))=\tan(\theta)=\frac{2\sqrt{26}}{11}.\]
\end{enumerate}
\end{example}

\printexercises{exercises/07-07-exercises}

% todo write an example of finding the inverse of a function

% todo write more exercises for 7.1

\section{Derivatives of Inverse Functions}\label{sec:deriv_inverse_function}

%Recall that a function $y=f(x)$ is said to be \emph{one-to-one} if it passes the horizontal line test; that is, for two different $x$ values $x_1$ and $x_2$, we do \emph{not} have $f(x_1)=f(x_2)$. In some cases the domain of $f$ must be restricted so that it is one-to-one. For instance, consider $f(x)=x^2$. Clearly, $f(-1)= f(1)$, so $f$ is not one-to-one on its regular domain, but by restricting $f$ to $(0,\infty)$, $f$ is one-to-one.\index{derivative!inverse function}
%
%Now recall that one-to-one functions have \emph{inverses}. That is, if $f$ is one-to-one, it has an inverse function, denoted by $f\primeskip^{-1}$, such that if $f(a)=b$, then $f\primeskip^{-1}(b) = a$. The domain of $f\primeskip^{-1}$ is the range of $f$, and vice-versa. For ease of notation, we set $g=f\primeskip^{-1}$ and treat $g$ as a function of $x$.
%
%Since $f(a)=b$ implies $g(b)=a$, when we compose $f$ and $g$ we get a nice result: \[f\bigl(g(b)\bigr) = f(a) = b.\] In general, $f\bigl(g(x)\bigr) =x$ and $g\bigl(f(x)\bigr) = x$. This gives us a convenient way to check if two functions are inverses of each other: compose them and if the result is $x$, then they are inverses (on the appropriate domains.)
%
%When the point $(a,b)$ lies on the graph of $f$, the point $(b,a)$ lies on the graph of $g$. This leads us to discover that the graph of $g$ is the reflection of $f$ across the line $y=x$. In \autoref{fig:inverse1} we see a function graphed along with its inverse. See how the point $(1,1.5)$ lies on one graph, whereas $(1.5,1)$ lies on the other. Because of this relationship, whatever we know about $f$ can quickly be transferred into knowledge about $g$.
%
%\mtable[-1in]{A function $f$ along with its inverse $f\primeskip^{-1}$. (Note how it does not matter which function we refer to as $f$; the other is $f\primeskip^{-1}$.)}{fig:inverse1}{\begin{tikzpicture}[alt={Blue curve f and red inverse reflect over dashed y=x; blue passes (-0.5,0.375)→(1,1.5); red passes (0.375,-0.5)→(1.5,1).}]
%\begin{axis}[width=1.16\marginparwidth,tick label style={font=\scriptsize},
%minor x tick num=1,minor y tick num=1,axis y line=middle,axis x line=middle,
%ymin=-1.2,ymax=1.7,xmin=-1.2,xmax=1.7,name=myplot]
%\addplot [draw={\colorone},smooth,thick,domain=-1:1.05] (x,x^3+.5);
%\addplot [draw={\colortwo},smooth,thick,domain=-1:1.05] (x^3+.5,x);
%\addplot [dashed,thin] {x};
%\filldraw [black] (axis cs:-.5,.375) node [above left] {\tiny $(-0.5,0.375)$} circle (1pt);
%\filldraw [black] (axis cs:.375,-.5) node [below right] {\tiny $(0.375,-0.5)$} circle (1pt);
%\filldraw [black] (axis cs:1,1.5) node [left] {\tiny $(1,1.5)$} circle (1pt);
%\filldraw [black] (axis cs:1.5,1) node [below] {\tiny $(1.5,1)$} circle (1pt);
%\end{axis}
%\node [right] at (myplot.right of origin) {\scriptsize $x$};
%\node [above] at (myplot.above origin) {\scriptsize $y$};
%\end{tikzpicture}}
%
%For example, consider \autoref{fig:inverse2} where the tangent line to $f$ at the point $(a,b)$ is drawn. That line has slope $\fp(a)$. Through reflection across $y=x$, we can see that the tangent line to $g$ at the point $(b,a)$ should have slope $\ds \frac{1}{\fp(a)}$. This then tells us that $\ds g\primeskip'(b) = \frac{1}{\fp(a)}.$
%
%\mtable{Corresponding tangent lines drawn to $f$ and $f\primeskip^{-1}$.}{fig:inverse2}{\begin{tikzpicture}[alt={Blue f and red f^(-1) reflected across dashed y=x; tangents at (a,b) on f and (b,a) on f^(-1); axes x, y.}]
%\begin{axis}[width=1.16\marginparwidth,tick label style={font=\scriptsize},
%minor x tick num=1,minor y tick num=1,axis y line=middle,axis x line=middle,
%ymin=-1.2,ymax=1.7,xmin=-1.2,xmax=1.7,name=myplot]
%\addplot [draw={\colorone},smooth,thick,domain=-1:1.05] (x,x^3+.5);
%\addplot [draw={\colortwo},smooth,thick,domain=-1:1.05] (x^3+.5,x);
%\addplot [dashed,thin] {x};
%\addplot [domain=-.9:-.1] {.75*(x+.5)+.375};
%\addplot [domain=.1:.7] {4/3*(x-.375)-.5};
%\filldraw [black] (axis cs:-.5,.375) node [above left] {\tiny $(a,b)$} circle (1pt);
%\filldraw [black] (axis cs:.375,-.5) node [below right] {\tiny $(b,a)$} circle (1pt);
%\end{axis}
%\node [right] at (myplot.right of origin) {\scriptsize $x$};
%\node [above] at (myplot.above origin) {\scriptsize $y$};
%\end{tikzpicture}}
%
%Consider:
%\begin{center} % this had some ``rule'' commands as well
%\tagpdfsetup{table/header-rows={1}}
%	\begin{tabular}{ c c c }
%	Information about $f$ & & Information about $g=f\primeskip^{-1}$ \\ \hline
%	\parbox{100pt}{\centering $(-0.5,0.375)$ lies on $f$} & \h skip 40pt & \parbox{100pt}{\centering $(0.375,-0.5)$ lies on $g$}\\
%	\parbox{100pt}{\centering Slope of tangent line to $f$ at $x=-0.5$ is $3/4$} & & \parbox{100pt}{\centering Slope of tangent line to $g$ at $x=0.375$ is $4/3$} \\
%	$\fp(-0.5) = 3/4$ & & $g\primeskip'(0.375) = 4/3$
%	\end{tabular}
%\end{center}
%
%We have discovered a relationship between $\fp$ and $g\primeskip'$ in a mostly graphical way. We can realize this relationship analytically as well. Let $y = g(x)$, where again $g = f\primeskip^{-1}$. We want to find $\ds y\primeskip'$. Since $y = g(x)$, we know that $f(y) = x$. Using the Chain Rule and Implicit Differentiation, take the derivative of both sides of this last equality.
%		\begin{align*}
%			\frac{\dd}{\dd x}\Bigl(f(y)\Bigr) &= \frac{\dd}{\dd x}\Bigl(x\Bigr) \\
%			\fp(y)\cdot y\primeskip' &= 1\\
%			y\primeskip' &= \frac{1}{\fp(y)}\\
%			y\primeskip' &= \frac{1}{\fp(g(x))}
%		\end{align*}
%		
%This leads us to the following theorem.

In this section we will figure out how to differentiate the inverse of a function. To do so, we recall that if $f$ and $g$ are inverses, then $f(g(x))=x$ for all $x$ in the domain of $f$. Differentiating and simplifying yields:
\begin{align*}
f(g(x))&=x\\
\fp(g(x))g\primeskip'(x)&=1\\
g\primeskip'(x)&=\frac 1{\fp(g(x))} \quad\text{assuming $\fp(x)$ is nonzero}
\end{align*}
Note that the derivation above assumes that the function $g$ is differentiable. It is possible to prove that $g$ must be differentiable if $\fp$ is nonzero, but the proof is beyond the scope of this text. However, assuming this fact we have shown the following:

\begin{theorem}[Derivatives of Inverse Functions]\label{thm:deriv_inverse_functions}%
Let $f$ be differentiable and one-to-one on an open interval $I$, where $\fp(x) \neq 0$ for all $x$ in $I$, let $J$ be the range of $f$ on $I$, let $g$ be the inverse function of $f$, and let $f(a) = b$ for some $a$ in $I$. Then $g$ is a differentiable function on $J$, and in particular,
\begin{align*}
 \left(f\primeskip^{-1}\right)'(b) &=g\primeskip'(b) = \frac{1}{\fp(a)} \\
 \left(f\primeskip^{-1}\right)'(x) &=g\primeskip'(x) = \frac{1}{\fp(g(x))}
\end{align*}
\end{theorem}

The results of \autoref{thm:deriv_inverse_functions} are not trivial; the notation may seem confusing at first. Careful consideration, along with examples, should earn understanding.

\youtubeVideo{RKfGMX0pn2k}{Derivative of an Inverse Function, Ex 2}

In the next example we apply \autoref{thm:deriv_inverse_functions} to the arcsine function.

\begin{example}[Finding the derivative of an inverse trigonometric function]\label{ex_deriv_arcsin}%
Let $y = \sin^{-1} x$. Find $y\primeskip'$ using \autoref{thm:deriv_inverse_functions}.
\solution
Adopting our previously defined notation, let $g(x) = \sin^{-1} x$ and $f(x) = \sin x$. Thus $\fp(x) = \cos x$. Applying \autoref{thm:deriv_inverse_functions}, we have 
\begin{align*}
	g\primeskip'(x) &= \frac{1}{\fp(g(x))} \\
	&= \frac{1}{\cos(\sin^{-1}x)}.
\end{align*}
			
\mtable{A right triangle defined by $y=\sin ^{-1}(x/1)$ with the length of the third leg found using the Pythagorean Theorem.}{fig:inverse3}{\begin{tikzpicture}[alt={Right triangle: hypotenuse 1 sloping up right; vertical leg length x; base √(1-x²); left angle y (arcsin x).},x=.5\marginparwidth,y=.5\marginparwidth,>=stealth]
\draw (0,0) node [shift={(15pt,5pt)}] {\scriptsize$y$}
 -- (1,0) node [pos=.5,below] {\scriptsize $\sqrt{1-x^2}$}
 -- (1,.75) node [pos=.5,right] {\scriptsize $x$}
 -- (0,0) node [pos=.5,above] {\scriptsize $1$};
\end{tikzpicture}}

This last expression is not immediately illuminating. Drawing a figure will help, as shown in \autoref{fig:inverse3}. Recall that the sine function can be viewed as taking in an angle and returning a ratio of sides of a right triangle, specifically, the ratio ``opposite over hypotenuse.'' This means that the arcsine function takes as input a ratio of sides and returns an angle. The equation $y=\sin^{-1} x$ can be rewritten as $y=\sin^{-1}(x/1)$; that is, consider a right triangle where the hypotenuse has length 1 and the side opposite of the angle with measure $y$ has length $x$. This means the final side has length $\sqrt{1-x^2}$, using the Pythagorean Theorem.

Therefore $\cos (\sin^{-1} x) = \cos y = \sqrt{1-x^2}/1 = \sqrt{1-x^2}$, resulting in
\[\frac \dd{\dd x}\bigl(\sin^{-1}x\bigr)=g\primeskip'(x)=\frac1{\sqrt{1-x^2}}.\]
\end{example}

Remember that the input $x$ of the arcsine function is a ratio of a side of a right triangle to its hypotenuse; the absolute value of this ratio will be less than 1. Therefore $1-x^2$ will be positive.

\mtable{Graphs of $y=\sin x$ and $y=\sin^{-1}x$ along with corresponding tangent lines.}{fig:inverse4}{\begin{tikzpicture}[alt={Top: blue y = sin x on −π/2<x<π/2 with black tangent at (π/3, √3/2). Bottom: red y = sin^(-1) x with tangent at (√3/2, π/3).}]
\begin{axis}[width=1.16\marginparwidth,tick label style={font=\scriptsize},
minor x tick num=1,minor y tick num=1,axis y line=middle,axis x line=middle,
ymin=-1.7,ymax=1.7,xmin=-1.7,xmax=1.7,xtick={-1.57,-.79,.79,1.57},
xticklabels={$-\frac{\pi}{2}$,$-\frac{\pi}{4}$,$\frac{\pi}{4}$,$\frac{\pi}{2}$},
name=myplot,axis equal]
\addplot [draw={\colorone},smooth,thick,domain=-1.57:1.57] (x,{sin(deg(x))});
\draw (axis cs:-1.1,-1.1) node {\scriptsize $y=\sin x$};
\filldraw [black] (axis cs:1.05,.866) node [above left] {\scriptsize $(\frac{\pi}{3},\frac{\sqrt{3}}{2})$} circle (1pt);
\addplot [draw=black,thick,smooth,domain=.5:1.5] {.5*(x-1.05)+.866};
\end{axis}
\node [right] at (myplot.right of origin) {\scriptsize $x$};
\node [above] at (myplot.above origin) {\scriptsize $y$};
\begin{scope}[shift={(0,-120pt)}]
\begin{axis}[width=1.16\marginparwidth,tick label style={font=\scriptsize},
minor x tick num=1,minor y tick num=1,axis y line=middle,axis x line=middle,
ymin=-1.7,ymax=1.7,xmin=-1.7,xmax=1.7,ytick={-1.57,-.79,.79,1.57},
yticklabels={$-\frac{\pi}{2}$,$-\frac{\pi}{4}$,$\frac{\pi}{4}$,$\frac{\pi}{2}$},
name=myplot,axis equal]
\addplot [draw={\colortwo},smooth,thick,domain=-1.57:1.57,] ({sin(deg(x))},x);
\draw (axis cs:-1.2,-.6) node {\scriptsize $y=\sin^{-1} x$};
\filldraw [black] (axis cs:.866,1.05) node [above left] {\scriptsize $(\frac{\sqrt{3}}{2},\frac{\pi}{3})$} circle (1pt);
\addplot [draw=black,thick,smooth,domain=.5:1.5] {2*(x-.866)+1.05};
\end{axis}
\end{scope}
\end{tikzpicture}}

In order to make $y=\sin x$ one-to-one, we restrict its domain to $[-\pi/2,\pi/2]$; on this domain, the range is $[-1,1]$. Therefore the domain of $y=\sin^{-1}x$ is $[-1,1]$ and the range is $[-\pi/2,\pi/2]$. When $x=\pm 1$, note how the derivative of the arcsine function is undefined; this corresponds to the fact that as $x\to \pm1$, the tangent lines to arcsine approach vertical lines with undefined slopes.

In \autoref{fig:inverse4} we see $f(x) = \sin x$ and $f\primeskip^{-1}(x)=\sin^{-1} x$ graphed on their respective domains. The line tangent to $\sin x$ at the point $(\pi/3, \sqrt{3}/2)$ has slope $\cos \pi/3 = 1/2$. The slope of the corresponding point on $\sin^{-1}x$, the point $(\sqrt{3}/2,\pi/3)$, is
\[\frac{1}{\sqrt{1-(\sqrt{3}/2)^2}} = \frac{1}{\sqrt{1-3/4}} = \frac{1}{\sqrt{1/4}} = \frac{1}{1/2}=2,\]
verifying \autoref{thm:deriv_inverse_functions} yet again: at corresponding points, a function and its inverse have reciprocal slopes.\bigskip

Using similar techniques, we can find the derivatives of all the inverse trigonometric functions after first restricting their domains according to \autoref{fig:domain_trig} to allow them to be invertible.

{\tcbset{grow to right by=3em}
% 2 em is sufficient, but 3em shortens the paragraph by a line
\begin{theorem}[Derivatives of Inverse Trigonometric Functions]\label{thm:deriv_inverse_trig}%
The inverse trigonometric functions are differentiable on all open sets contained in their domains (as listed in \autoref{fig:domain_trig}) and their derivatives are as follows:\index{derivative!inverse trig.}
\begin{multicols}{2}
	\begin{enumerate}\lxAddClass{columns2}
		\item	$\ds\frac \dd{\dd x}\bigl(\sin^{-1}x\bigr)= \frac1{\sqrt{1-x^2}}$ 
		\item	$\ds\frac \dd{\dd x}\bigl(\sec^{-1}x\bigr)= \frac1{\abs{x}\sqrt{x^2-1}}$
		\item	$\ds\frac \dd{\dd x}\bigl(\tan^{-1}x\bigr)= \frac1{1+x^2}$
		\item	$\ds\frac \dd{\dd x}\bigl(\cos^{-1}x\bigr)=-\frac1{\sqrt{1-x^2}}$ 
		\item	$\ds\frac \dd{\dd x}\bigl(\csc^{-1}x\bigr)=-\frac1{\abs{x}\sqrt{x^2-1}}$
		\item	$\ds\frac \dd{\dd x}\bigl(\cot^{-1}x\bigr)=-\frac1{1+x^2}$
	\end{enumerate}
\end{multicols}
\end{theorem}}

Note how the last three derivatives are merely the negatives of the first three, respectively. Because of this, the first three are used almost exclusively throughout this text.

%In \autoref{sec:basic_diff_rules}, we stated without proof or explanation that $\ds \frac{\dd}{\dd x}\bigl(\ln x\bigr) = \frac1x$. We can justify that now using \autoref{thm:deriv_inverse_functions}, as shown in the example.
%
%\begin{example}[Finding the derivative of $y=\ln x$]\label{ex_deriv_lnx}%
%Use \autoref{thm:deriv_inverse_functions} to compute $\ds \frac{\dd}{\dd x}\bigl(\ln x\bigr)$.
%\solution
%View $y= \ln x$ as the inverse of $y = e^x$. Therefore, using our standard notation, let $f(x) = e^x$ and $g(x) = \ln x$. We wish to find $g\primeskip'(x)$. \autoref{thm:deriv_inverse_functions} gives:
%\begin{align*}
%	g\primeskip'(x)
%	&= \frac1{\fp(g(x))} \\
%	&= \frac1{e^{\ln x}} \\
%	&= \frac1x.
%\end{align*}
%\end{example}

\begin{example}[Finding derivatives of inverse functions]\label{eg_inv_derivs}%
Find the derivatives of the following functions:
\[
 \text{1.}\quad f(x)=\cos^{-1}(x^2)\qquad
 \text{2.}\quad g(x)=\frac{\sin^{-1}x}{\sqrt{1-x^2}}\qquad
 \text{3.}\quad f(x)=\sin^{-1}(\cos x)
\]
\solution
\begin{enumerate}
\item We use \autoref{thm:deriv_inverse_trig} and the Chain Rule to find:
\[\fp(x)=-\frac{1}{\sqrt{1-(x^2)^2}}(2x)=- \frac{2x}{\sqrt{1-x^4}}\]
\item We use \autoref{thm:deriv_inverse_trig} and the Quotient Rule to compute: 
\begin{align*}
 g\primeskip'(x)
 &=\frac{\left(\frac 1{\sqrt{1-x^2}}\right)\sqrt{1-x^2} - (\sin^{-1} x)\left( \frac 1{2\sqrt{1-x^2}} (-2x)\right)}{\left(\sqrt{1-x^2}\right)^2}\\
 &=\frac{\sqrt{1-x^2}+x\sin^{-1}x}{\left(\sqrt{1-x^2}\right)^3}
\end{align*}
\item We apply \autoref{thm:deriv_inverse_trig} and the Chain Rule again to compute:
\begin{align*}
 \fp(x)&=\frac{1}{\sqrt{1-\cos^2x}}(-\sin x)\\
 &=\frac{-\sin x}{\sqrt{\sin^2x}}\\
 &=\frac{-\sin x}{\abs{\sin x}}.
\end{align*}
\end{enumerate}
\end{example}

\autoref{thm:deriv_inverse_trig} allows us to integrate some functions that we could not integrate before. For example,
\[\int\frac{\dd x}{\sqrt{1-x^2}}=\sin^{-1}x+C.\]
Combining these formulas with $u$-substitution yields the following:

{\tcbset{grow to right by=1em}
\begin{theorem}[Integrals Involving Inverse Trigonometric Functions]\label{thm:int_inverse_trig}%
Let $a>0$.  Then
\begin{enumerate}
	\item	$\ds\int\frac1{a^2+x^2}\dd x=\frac1a\tan^{-1}\left(\frac xa\right) + C$
	\item	$\ds\int\frac1{\sqrt{a^2-x^2}}\dd x=\sin^{-1}\left(\frac xa\right)+C$
	\item	$\ds\int\frac1{x\sqrt{x^2-a^2}}\dd x=\frac1a\sec^{-1}\left(\frac{\abs x}a\right)+C$
\end{enumerate}
\end{theorem}}

We will look at the second part of this theorem. The other parts are similar and are left as exercises.

First we note that the integrand involves the number $a^2$, but does not explicitly involve $a$. We make the assumption that $a>0$ in order to simplify what follows. We can rewrite the integral as follows:
\[\int\frac{\dd x}{\sqrt{a^2-x^2}}=\int\frac{\dd x}{\sqrt{a^2(1-(x/a)^2)}}=\int\frac{\dd x}{a\sqrt{1-(x/a)^2}}\]
We next use the substitution $u=x/a$ and $\dd u=\dd x/a$ to find: 
\begin{align*}
\int\frac{\dd x}{a\sqrt{1-(x/a)^2}}
&=\int\frac{a}{a\sqrt{1-u^2}}\dd u\\
&=\int \frac{\dd u}{\sqrt{1-u^2}}\\
&=\sin^{-1}u+C\\
&=\sin^{-1}(x/a)+C
\end{align*}

We conclude this section with several examples.

\begin{example}[Finding antiderivatives involving inverse functions]\label{eg_inv_deriv_harder}%
Find the following integrals.
\[
 \text{1.}\quad\int\frac{\dd x}{100+x^2}\qquad
 \text{2.}\quad\int\frac{\sin^{-1} x}{\sqrt{1-x^2}}\dd x\qquad
 \text{3.}\quad\int\frac{\dd x}{x^2+2x+5}
\]
\solution
\begin{enumerate}
\item $\ds\int\frac{\dd x}{100+x^2}=\int\frac{\dd x}{10^2+x^2}=\frac1{10}\tan^{-1}(x/10)+C$
\item We use the substitution $u=\sin^{-1}x$ and $\dd u=\frac{\dd x}{\sqrt{1-x^2}}$ to find:
\[\int\frac{\sin^{-1}x}{\sqrt{1-x^2}}\dd x=\int u\dd u=\frac12 u^2+C=\frac 12\left(\sin^{-1}x\right)^2+C\]
\item This does not immediately look like one of the forms in \autoref{thm:int_inverse_trig}, but we can complete the square in the denominator to see that
\[\int\frac{\dd x}{x^2+2x+5} =\int\frac{\dd x}{(x^2+2x+1)+4}=\int\frac{\dd x}{4+(x+1)^2}\]
We now use the substitution $u=x+1$ and $\dd u=\dd x$ to find:
\begin{multline*}
 \int\frac{\dd x}{4+(x+1)^2} =\int\frac{\dd u}{4+u^2}\\
 =\frac12 \tan^{-1}(u/2)+C =\frac 12\tan^{-1}\left(\frac{x+1}2\right)+C.
\end{multline*}
\end{enumerate}
\end{example}

% moved to end of implicit differentiation.  but do we want to give it again?
%In this chapter we have defined the derivative, given rules to facilitate its computation, and given the derivatives of a number of standard functions. We restate the most important of these in the following theorem, intended to be a reference for further work.
%
%\begin{theorem}[Glossary of Derivatives of Elementary Functions]\label{thm:deriv_glossary}%
%Let $u$ and $v$ be differentiable functions, and let $c$ and $n$ be real numbers, $n\neq 0$.
%\begin{enumext*}[columns=2, widest=10]
%	\item		$\frac{\dd}{\dd x}\bigl(cu\bigr) = cu'$
%	\item		$\frac{\dd}{\dd x}\bigl(u\pm v\bigr) = u'\pm v'$
%	\item		$\frac{\dd}{\dd x}\bigl(u\cdot v\bigr) = uv'+u'v$
%	\item		$\frac{\dd}{\dd x}\bigl(\frac uv\bigr) = \frac{u'v-uv'}{v^2}$
%	\item		$\frac{\dd}{\dd x}\bigl(u(v)\bigr) = u'(v)v'$
%	\item		$\frac{\dd}{\dd x}\bigl(x^n\bigr) = nx^{n-1}$
%	\item		$\frac{\dd}{\dd x}\bigl(c\bigr) = 0$
%	\item		$\frac{\dd}{\dd x}\bigl(x\bigr) = 1$
%	\item		$\frac{\dd}{\dd x}\bigl(\ln x\bigr) = \frac{1}{x}$
%	\item		$\frac{\dd}{\dd x}\bigl(e^x\bigr) = e^x$
%	\item		$\frac{\dd}{\dd x}\bigl(\sin x\bigr) = \cos x$
%	\item		$\frac{\dd}{\dd x}\bigl(\cos x\bigr) = -\sin x$
%	\item		$\frac{\dd}{\dd x}\bigl(\tan x\bigr) = \sec^2x$
%	\item		$\frac{\dd}{\dd x}\bigl(\cot x\bigr) = -\csc^2x$
%	\item		$\frac{\dd}{\dd x}\bigl(\sec x\bigr) = \sec x\tan x$
%	\item		$\frac{\dd}{\dd x}\bigl(\csc x\bigr) = -\csc x\cot x$
%	\item		$\frac{\dd}{\dd x}\bigl(a^x\bigr) = \ln a\cdot a^x$
%	\item		$\frac{\dd}{\dd x}\bigl(\log_a x\bigr) = \frac{1}{\ln a}\cdot\frac{1}{x}$
%	\item		$\frac{\dd}{\dd x}\bigl(\sin^{-1}x\bigr) = \frac{1}{\sqrt{1-x^2}}$
%	\item		$\frac{\dd}{\dd x}\bigl(\csc^{-1}x\bigr) = -\frac{1}{\abs{x}\sqrt{x^2-1}}$
%	\item		$\frac{\dd}{\dd x}\bigl(\tan^{-1}x\bigr) = \frac{1}{1+x^2}$
%	\item		$\frac{\dd}{\dd x}\bigl(\cos^{-1}x\bigr) = -\frac{1}{\sqrt{1-x^2}}$
%	\item		$\frac{\dd}{\dd x}\bigl(\sec^{-1}x\bigr) = \frac{1}{\abs{x}\sqrt{x^2-1}}$
%	\item		$\frac{\dd}{\dd x}\bigl(\cot^{-1}x\bigr) = -\frac{1}{1+x^2}$
%\end{enumext*}
%\end{theorem}

\printexercises{exercises/02-07-exercises}

\section{Exponential and Logarithmic Functions}\label{sec:exp_log}

In this section we will define general exponential and logarithmic functions and find their derivatives. 

\subsection{General exponential functions}

\mtable{The function $2^x$ for rational values of $x$.}{fig:two_to_x}{%
\begin{tikzpicture}[alt={Blue dotted curve y=2^x: passes (0,1), trend 0<y<3 for −1<x<1.5, climbs fast rightward, approaches y=0 leftward.}]
\begin{axis}[width=\marginparwidth, axis y line=middle,axis x line=middle,
tick label style={font=\scriptsize},
ymin=-1,ymax=5,xmin=-4,xmax=4,name=myplot,xtick={1},ytick={1,2}]
\addplot[draw={\colorone},densely dotted,domain=-4:4]{2^x};
\end{axis}
\node [right] at (myplot.right of origin) {\scriptsize $x$};
\node [above] at (myplot.above origin) {\scriptsize $y$};
\end{tikzpicture}}
Consider first the function $f(x)=2^x$. If $x$ is rational, then we know how to compute $2^x$. What do we mean by $2^\pi$ though? We compute this by first looking at $2^r$ for rational numbers $r$ that are very close to $\pi$, then finding a limit. In our case we might compute $2^3$, $2^{3.1}$, $2^{3.14}$, etc. We then define $2^\pi$ to be the limit of these numbers. Note that this is actually a different kind of limit than we have dealt with before since we only consider rational numbers close to $\pi$, not all real numbers close to $\pi$. We will see one way to make this more precise in \autoref{chapter:sequences_series}. Graphically, we can plot the values of $2^x$ for $x$ rational and get something like the dotted curve in \autoref{fig:two_to_x}. In order to define the remaining values, we are ``connecting the dots'' in a way that makes the function continuous.

\mtable{The functions $2^x$ and $2^{-x}$.}{fig:two_to_minus_x}{%
\begin{tikzpicture}[alt={Blue y=2^x increasing and red y=2^−x decreasing; meet at (0,1);  −1.2<x<1.2, 0.3<y<3; mirror images about y-axis.}]
\begin{axis}[width=\marginparwidth, axis y line=middle,axis x line=middle,
tick label style={font=\scriptsize},
ymin=-1,ymax=5,xmin=-4,xmax=4,name=myplot,xtick={-1,1},ytick={.5,1,2}]
\addplot[draw={\colorone},domain=-4:4]{2^x};
\addplot[draw={\colortwo},domain=-4:4]{2^(-x)};
\end{axis}
\node [right] at (myplot.right of origin) {\scriptsize $x$};
\node [above] at (myplot.above origin) {\scriptsize $y$};
\end{tikzpicture}}
It follows from continuity and the properties of limits that exponential functions will satisfy the familiar properties of exponents (see \autoref{sec:deriv_prereqs}).
This implies that
\[\left(\frac12\right)^x =(2^{-1})^x=2^{-x},\]
so the graph of $g(x)=(1/2)^x$ is the reflection of $f$ across the $y$-axis, as in \autoref{fig:two_to_minus_x}.

We can go through the same process as above for any base $a>0$, though we are not usually interested in the constant function $1^x$.

\begin{keyidea}[Properties of Exponential Functions]\label{ki_exp_func_props}%
For $a>0$ and $a\neq1$ the exponential function $f(x)=a^x$ satisfies:
\begin{multicols}{2}
\begin{enumerate}\lxAddClass{columns2}
 \item $a^0=1$
 \item $\ds\lim_{x\to\infty}a^x=\begin{cases}\infty&a>1\\0&a<1\end{cases}$
 \item $a^x>0$ for all $x$
 \item $\ds\lim_{x\to-\infty}a^x=\begin{cases}0&a>1\\\infty&a<1\end{cases}$
\end{enumerate}
\end{multicols}
\end{keyidea}

\subsection{Derivatives of exponential functions}

Suppose $f(x)=a^x$ for some $a>0$. We can use the rules of exponents to find the derivative of $f$:
\begin{align*}
	\fp(x)
	&=\lim_{h\to 0}\frac{f(x+h)-f(x)}{h}\\
	&=\lim_{h\to 0}\frac{a^{x+h}-a^x}{h}\\
	&=\lim_{h\to 0}\frac{a^xa^h-a^x}{h}\\
	&=\lim_{h\to 0}\frac{a^x(a^h-1)}{h}\\
	&=a^x\lim_{h\to 0}\frac{a^h-1}{h} \qquad\text{(since $a^x$ does not depend on $h$)}\\
\end{align*}
So we know that $\fp(x)=\ds a^x \lim_{h\to 0}\frac{a^h-1}{h}$, but can we say anything about that remaining limit? First we note that
\[\fp(0)=\lim_{h\to 0}\frac  {a^{0+h}-a^0}{h}=\lim_{h\to 0}\frac{a^h-1}{h},\]
so we have $\fp(x)=a^x\fp(0)$. The actual value of the limit $\ds\lim_{h\to 0}\frac{a^h-1}{h}$ depends on the base $a$, but it can be proved that it does exist. We will figure out just what this limit is later, but for now we note that the easiest differentiation formulas come from using a base $a$ that makes $\ds\lim_{h\to 0}\frac{a^h-1}{h}=1$. This base is the
% irrational ?
number $e\approx 2.71828$ and the exponential function $e^x$ is called the natural exponential function. This leads to the following result.

\begin{theorem}[Derivative of Exponential Functions]\label{thm_baby_exp_deriv}%
For any base $a>0$, the exponential function $f(x)=a^x$ has derivative $\fp(x)=a^x\fp(0)$. The natural exponential function $g(x)=e^x$ has derivative $g\primeskip'(x)=e^x$.
\end{theorem}

\youtubeVideo{U3PyUcEd7IU}{Derivatives of Exponential Functions}

%\begin{example}[Finding Exponential Derivatives]\label{ex_exp_derivs}%
%Find derivatives of the following functions.
%\[\text{1. }f(x)=e^{x^2}\qquad\text{2. }g(t)=t e^t\]
%\solution
%\begin{enumerate}
%\item We note that this is a composition of functions where the \emph{inside} function is $x^2$. Hence we can apply the Chain Rule to see that: \[\fp(x)=\left(e^{x^2}\right) (2x)\]
%\item In this case we apply the Product Rule for derivatives to find: \[g\primeskip'(t)=e^t+te^t=e^t(t+1).\]
%\end{enumerate}
%\end{example}

\subsection{General logarithmic functions}

% todo Tim move this to 7.0 once that's written
Before reviewing general logarithmic functions, we'll first remind ourselves of the laws of logarithms.

\begin{keyidea}[Properties of Logarithms]\label{ki_log_laws}%
For $a,x,y>0$ and $a\ne1$, we have
\begin{enumext*}[columns=2]
\item$\log_a(xy)=\log_ax+\log_ay$
\item$\log_a\dfrac xy=\log_ax-\log_ay$\vspace{-.5ex}
\item$\log_xy=\dfrac{\log_ay}{\log_ax}$, when $x\ne1$
\item$\log_ax^y=y\log_ax$
\item$\log_a1=0$
\item$\log_aa=1$
\end{enumext*}
\end{keyidea}

\mtable{The functions $y=a^x$ and $y=\log_a x$ for $a>1$.}{fig_exp_log}{%
\begin{tikzpicture}[alt={Blue y=a^x and red y=log\_a x (a>1) reflected across dashed y=x; meet at (1,1); axes ticked at 1 and e on both scales.}]
\begin{axis}[width=\marginparwidth,name=myplot,axis equal,
tick label style={font=\scriptsize},
axis y line=middle,axis x line=middle,ymin=-2,ymax=4,xmin=-2,xmax=4,
xtick={1,2.718},xticklabels={1,$e$},
ytick={1,2.718},yticklabels={1,$e$}]
\addplot[draw={\colorone},domain=-4:1.4]{e^x};
\addplot[draw={\colortwo},domain=.1:4]{ln(x)};
\addplot[dashed,domain=-4:4]{x};
\end{axis}
\node [right] at (myplot.right of origin) {\scriptsize $x$};
\node [above] at (myplot.above origin) {\scriptsize $y$};
\end{tikzpicture}}
Let us consider the function $f(x)=a^x$ where $a\neq1$. We know that $\fp(x)=\fp(0)a^x$, where $\fp(0)$ is a constant that depends on the base $a$. Since $a^x>0$ for all $x$, this implies that $\fp(x)$ is either always positive or always negative, depending on the sign of $\fp(0)$. This in turn implies that $f$ is strictly monotonic, so $f$ is one-to-one. We can now say that $f$ has an inverse. We call this inverse the logarithm with base $a$, denoted $f^{-1}(x)=\log_ax$. When $a=e$, this is the natural logarithm function $\ln x$. So we can say that $y=\log_a x$ if and only if $a^y=x$. Since the range of the exponential function is the set of positive real numbers, the domain of the logarithm function is also the set of positive real numbers. Reflecting the graph of $y=a^x$ across the line $y=x$ we find that (for $a>1$) the graph of the logarithm looks like \autoref{fig_exp_log}.

\begin{keyidea}[Properties of Logarithmic Functions]\label{ki_log_func_props}%
For $a>0$ and $a\ne1$ the logarithmic function $f(x)=\log_a x$ satisfies:
\begin{enumerate}
\item The domain of $f(x)=log_a x$ is $(0,\infty)$ and the range is $(-\infty,\infty)$.
\item $y=\log_a x$ if and only if $a^y=x$.
\item $\ds\lim_{x\to\infty}\log_a x=
\begin{cases}\infty&\text{ if $a>1$}\\ -\infty&\text{ if $a<1$}\end{cases}$
\item $\ds\lim_{x\to0^+}\log_a x=\begin{cases}-\infty&\text{ if $a>1$}\\\infty&\text{ if $a<1$}\end{cases}$
\end{enumerate}
\end{keyidea}

Using the inverse of the natural exponential function, we can determine what the value of $\fp(0)$ is in the formula $(a^x)'=\fp(0)a^x$. To do so, we note that $a=e^{\ln a}$ since the exponential and logarithm functions are inverses. Hence we can write:
\[a^x=\left(e^{\ln a}\right)^x=e^{x\ln a}\]
Now since $\ln a$ is a constant, we can use the Chain Rule to see that:
\[\frac d{\dd x} a^x=\frac d{\dd x} e^{x\ln a} =e^{x\ln a}(\ln a) =a^x\ln a\]
Comparing this to our previous result, we can restate our theorem:

\begin{theorem}[Derivative of Exponential Functions]\label{thm_exp_deriv}%
For any base $a>0$, the exponential function $f(x)=a^x$ has derivative $\fp(x)=a^x\ln a$. The natural exponential function $g(x)=e^x$ has derivative $g\primeskip'(x)=e^x$.
\end{theorem}

\subsection{Change of base}

In the previous computation, we found it convenient to rewrite the general exponential function in terms of the natural exponential function. A related formula allows us to rewrite the general logarithmic function in terms of the natural logarithm.  To see how this works, suppose that $y=\log_ax$, then we have:
\begin{align*}
a^y&=x \\
\ln(a^y)&=\ln x\\
y\ln a&=\ln x\\
y&=\frac{\ln x}{\ln a}\\
\log_a x&=\frac{\ln x}{\ln a}.
\end{align*}
This change of base formula allows us to use facts about the natural logarithm to derive facts about the general logarithm.

\subsection{Derivatives of logarithmic functions}

Since the natural logarithm function is the inverse of the natural exponential function, we can use the formula $(f^{-1}(x))'=\dfrac1{\fp(f^{-1}(x))}$ to find the derivative of $y=\ln x$. We know that $\dfrac\dd{\dd x}e^x=e^x$, so we get:
\[\frac\dd{\dd x}\ln x=\frac1{e^y}=\frac1{e^{\ln x}}=\frac1x.\]
Now we
%In \autoref{sec:deriv_inverse_function}, \autoref{ex_deriv_lnx}, we found that $\frac{\dd}{\dd x}\ln x=\frac1x$. We
can apply the change of base formula to find the derivative of a general logarithmic function:
\[\frac{\dd}{\dd x}\log_ax=\frac{\dd}{\dd x}\left(\frac{\ln x}{\ln a}\right) =\frac 1{\ln a}\left(\frac{\dd}{\dd x}\ln x\right)=\frac 1{x\ln a}.\]

\begin{example}[Finding Derivatives of Logs and Exponentials]\label{ex_find_d_exp_log}%
Find derivatives of the following functions.
\[
 \text{1.}\quad f(x)=x3^{4x-7}\qquad
 \text{2.}\quad g(x)=2^{x^2}\qquad
 \text{3.}\quad h(x)=\frac x{\log_5x}
\]
\solution
\begin{enumerate}
\item We apply both the Product and Chain Rules:
\[\fp(x)=3^{4x-7}+x\left(3^{4x-7}\ln 3\right)(4)=(1+4x\ln 3)3^{4x-7}\]
\item We apply the Chain Rule:
\[g\primeskip'(x)=2^{x^2}(\ln 2)(2x)=2^{x^2+1}x\ln2.\]
\item Applying the Quotient Rule:
\[h\primeskip'(x)=\frac{\log_5x-x\left(\frac1{x\ln 5}\right)}{(\log_5x)^2}=\frac{(\log_5x)(\ln 5)-1}{(\log_5x)^2\ln 5}\]
\end{enumerate}
\end{example}

\begin{example}[The Derivative of the Natural Log]\label{ex_d_ln_abs_x}%
Find the derivative of the function $y=\ln\abs x$.
\solution
We can rewrite our function as
\[y=\begin{cases} \ln x & \text{if $x>0$}\\ \ln(-x) & \text{if $x<0$}\\ \end{cases}\]
Applying the Chain Rule, we see that $\frac{\dd y}{\dd x}=\frac 1x$ for $x>0$, and $\frac{\dd y}{\dd x}=\frac{-1}{-x}=\frac 1x$ for $x<0$. Hence we have
\[\frac{\dd}{\dd x}\ln\abs x=\frac 1x \quad\text{ for $x\neq0$.}\]
\end{example}

\subsection{Antiderivatives}

Combining these new results, we arrive at the following theorem:

\begin{theorem}[Derivatives and Antiderivatives of Exponentials and Logarithms]\label{thm_int_exp_log}%
Given a base $a>0$ and $a\neq 1$, the following hold:\\[-1.5\baselineskip]
\begin{multicols}{2}
\begin{enumerate}\lxAddClass{columns2}
\item $\dfrac\dd{\dd x}e^x=e^x$
\item $\dfrac\dd{\dd x}a^x=a^x\ln a$
\item $\dfrac\dd{\dd x}\ln x=\dfrac1x$
\item $\dfrac\dd{\dd x}\log_a x=\dfrac1{x\ln a}$
\item $\ds\int e^x\dd x=e^x+C$
\item $\ds\int a^x\dd x=\frac{a^x}{\ln a}+C$
\item $\ds\int \frac{\dd x}{x}=\ln\abs x+C$
\item[]
\end{enumerate}
\end{multicols}
\end{theorem}

\begin{example}[Finding Antiderivatives]\label{ex_exp_log_anti}%
Find the following antiderivatives.
\[
 \text{1. }\int 3^x\dd x\qquad
 \text{2. }\int x^2 e^{x^3}\dd x\qquad
 \text{3. }\int \frac{x\dd x}{x^2+1}
\]
\solution
\begin{enumerate}
\item Applying our theorem,
\[\int 3^x\dd x=\frac{3^x}{\ln 3}+C\]
\item We use the substitution $u=x^3$, $\dd u=3x^2\dd x$:
\begin{align*}
\int x^2e^{x^3}\dd x &=\frac13 \int e^u\dd u\\
&=\frac 13 e^u+C\\
&=\frac 13 e^{x^3}+C\\
\end{align*}
\item Using the substitution $u=x^2+1$, $\dd u=2x\dd x$:
\begin{align*}
\int \frac{x\dd x}{x^2+1}&=\frac 12 \int \frac{\dd u}{u}\\
&=\frac 12 \ln\abs u+C\\
&=\frac 12 \ln\abs{x^2+1}+C\\
&=\frac 12 \ln(x^2+1)+C
\end{align*}
\end{enumerate}
\end{example}

Note that we do not yet have an antiderivative for the function $f(x)=\ln x$. We remedy this in \autoref{sec:IBP} with \autoref{ex_ibp5}.
% although see Example 2.4.4

%with the following example.
%
%example: Compute $\int \ln x\dd x$.
%
%\solution While this does not look like a product of functions, we find integration by parts useful. In particular we use $u=\ln x$, $\dd u=\dd x/x$, $\dd v=\dd x$, and $v=x$:
%\begin{align*}
%\int \ln x\dd x &=x\ln x-\int (x)\left(\frac{\dd x}x\right)\\
%&=x\ln x-\int \dd x\\
%&=x\ln x-x+C\\
%\end{align*}

\subsection{Logarithmic Differentiation}

\mtable{A plot of $y=x^x$.}{fig:logdiffa}{\begin{tikzpicture}[alt={Blue y=x^x for 0<x<2; open dot at (0,1); min near x=0.37,y=0.69; rises through (1,1) to about (2,4).}]
\begin{axis}[width=1.16\marginparwidth,tick label style={font=\scriptsize},
axis y line=middle,axis x line=middle,name=myplot,
ymin=-.1,ymax=4.2,xmin=-.1,xmax=2.2]
\addplot[draw={\colorone},smooth,thick,domain=.01:2] {exp(x*ln(x))};
\draw [fill=white,draw=black,thick] (axis cs:0,1) circle (1.5pt);
\end{axis}
\node [right] at (myplot.right of origin) {\scriptsize $x$};
\node [above] at (myplot.above origin) {\scriptsize $y$};
\end{tikzpicture}}

Consider the function $y=x^x$; it is graphed in \autoref{fig:logdiffa}. It is well-defined for $x>0$ and we might be interested in finding equations of lines tangent and normal to its graph. How do we take its derivative?\index{logarithmic differentiation}

The function is not a power function: it has a ``power'' of $x$, not a constant. It is not an exponential function: it has a ``base'' of $x$, not a constant. 

A differentiation technique known as \emph{logarithmic differentiation} becomes useful here. The basic principle is this: take the natural log of both sides of an equation $y=f(x)$, then use implicit differentiation to find $y\primeskip'$. We demonstrate this in the following example.

\begin{example}[Using Logarithmic Differentiation]\label{ex_implicit10}%
Given $y=x^x$, use logarithmic differentiation to find $y\primeskip'$.
\solution
As suggested above, we start by taking the natural log of both sides then applying implicit differentiation.
\begin{align*}
y &= x^x \\
\ln (y) &= \ln (x^x) && \text{\small (apply logarithm rule)}\\
\ln (y) &= x\ln x && \text{\small (now use implicit differentiation)}\\
\frac{\dd}{\dd x}\Bigl(\ln (y)\Bigr) &= \frac{\dd}{\dd x}\Bigl(x\ln x\Bigr) \\
\frac{y\primeskip'}{y} &= \ln x + x\cdot\frac1x\\
\frac{y\primeskip'}{y} &= \ln x + 1\\
y\primeskip'
&= y\bigl(\ln x+1\bigr) && \text{\small (substitute $y=x^x$)}\\
y\primeskip' &= x^x\bigl(\ln x+1\bigr).
\end{align*} 

\mtable{A graph of $y=x^x$ and its tangent line at $x=1.5$.}{fig:implicit10}{\begin{tikzpicture}[alt={Blue y=x^x with red tangent drawn at x=1.5; tangent touches curve near (1.5, 3.2) and slopes upward rightward.}]
\begin{axis}[width=1.16\marginparwidth,tick label style={font=\scriptsize},
axis y line=middle,axis x line=middle,name=myplot,
ymin=-.1,ymax=4.2,xmin=-.1,xmax=2.2]
\addplot[draw={\colorone},smooth,thick,domain=.01:2] {exp(x*ln(x))};
\addplot [draw={\colortwo},thick,domain=.9:2] {2.582*(x-1.5)+1.837};
\draw [fill=white,draw=black,thick] (axis cs:0,1) circle (1.5pt);
\end{axis}
\node [right] at (myplot.right of origin) {\scriptsize $x$};
\node [above] at (myplot.above origin) {\scriptsize $y$};
\end{tikzpicture}}

To ``test'' our answer, let's use it to find the equation of the tangent line at $x=1.5$. The point on the graph our tangent line must pass through is $(1.5, 1.5^{1.5}) \approx (1.5, 1.837)$. Using the equation for $y\primeskip'$, we find the slope as
\[y\primeskip' = 1.5^{1.5}\bigl(\ln 1.5+1\bigr) \approx 1.837(1.405) \approx 2.582.\]
Thus the equation of the tangent line is $y = 2.582(x-1.5)+1.837$. \autoref{fig:implicit10} graphs $y=x^x$ along with this tangent line.
\end{example}

\printexercises{exercises/07-exp-log-exercises}

% todo write another logarithmic differentiation example

\section{Hyperbolic Functions}\label{sec:hyperbolic}

The \textbf{hyperbolic functions} are functions that have many applications to mathematics, physics, and engineering. Among many other applications, they are used to describe the formation of satellite rings around planets, to describe the shape of a rope hanging from two points, and have application to the theory of special relativity. This section defines the hyperbolic functions and describes many of their properties, especially their usefulness to calculus.

\mtable{Using trigonometric functions to define points on a circle and hyperbolic functions to define points on a hyperbola.}{fig:hfcircle}{\begin{tikzpicture}[alt={The unit circle, with a shaded sector in the first quadrant of angle θ and area θ/2.}]
 \begin{axis}[height=\marginparwidth,width=\marginparwidth,
   tick label style={font=\scriptsize},
   axis y line=middle,axis x line=middle,name=myplot,axis on top,
   xtick={-1,1},ytick={-1,1},ymin=-1.1,ymax=1.1,xmin=-1.1,xmax=1.1]
  \addplot [draw={\coloronefill},fill={\coloronefill},domain=0:45] ({cos(x)},{sin(x)})
   -- (axis cs:0,0)--cycle;
  \addplot [draw={\colorone},domain=0:360,thick,smooth,samples=40] ({cos(x)},{sin(x)});
  \filldraw (axis cs:.707,.707) circle (1pt) node [left]%[shift={(4pt,11pt)}]
   {\scriptsize ($\cos \theta$,$\sin \theta$)};
  \draw (axis cs:.6,.25) node {\scriptsize $A=\dfrac{\theta}{2}$};
  \draw (axis cs:-.5,.2) node {\scriptsize $x^2+y^2=1$};
 \end{axis}
% \draw[draw={\colorone}](0,0)circle(1);
 \node [right] at (myplot.right of origin) {\scriptsize $x$};
 \node [above] at (myplot.above origin) {\scriptsize $y$};
\end{tikzpicture}
\vspace{10pt}
\begin{tikzpicture}[alt={The unit hyperbola.  The region in the first quadrant left of the right half and to the right of the segment connecting the origin to (coshθ,sinhθ) has been shaded.}]
 \begin{axis}[height=\marginparwidth,width=\marginparwidth,
   tick label style={font=\scriptsize},
   axis y line=middle,axis x line=middle,name=myplot,axis on top,
   ymin=-3.1,ymax=3.1,xmin=-3.1,xmax=3.1]
  \addplot [draw={\coloronefill},fill={\coloronefill},domain=0:1.6] ({cosh(x)},{sinh(x)})
   -- (axis cs:0,0)--cycle;
  \addplot [draw={\colorone},domain=-2:2,thick,smooth] ({cosh(x)},{sinh(x)});
  \addplot [draw={\colorone},domain=-2:2,thick,smooth] ({-cosh(x)},{sinh(x)});
  \filldraw (axis cs:2.577,2.376) circle (1pt) node [left]
   {\scriptsize ($\cosh \theta$,$\sinh \theta$)};
  \draw (axis cs:2,.6) node {\scriptsize $A=\dfrac{\theta}{2}$};
  \draw (axis cs:-1.75,2.75) node {\scriptsize $x^2-y^2=1$};
 \end{axis}
 \node [right] at (myplot.right of origin) {\scriptsize $x$};
 \node [above] at (myplot.above origin) {\scriptsize $y$};
\end{tikzpicture}}

These functions are sometimes referred to as the ``hyperbolic trigonometric functions'' as there are many connections between them and the standard trigonometric functions. \autoref{fig:hfcircle} demonstrates one such connection. Just as cosine and sine are used to define points on the circle defined by $x^2+y^2=1$, the functions \textbf{hyperbolic cosine} and \textbf{hyperbolic sine} are used to define points on the hyperbola $x^2-y^2=1$.

We begin with their definitions.

\begin{definition}[Hyperbolic Functions]\label{def:hyperbolic_functions}%
\mbox{}\\[-2\baselineskip]\parbox[t]{\linewidth}{%
\begin{multicols}{2}\index{hyperbolic function!definition}%
\begin{enumerate}\lxAddClass{columns2}
\item		$\ds \cosh x = \frac{e^x+e^{-x}}2$
\item		$\ds \sinh x = \frac{e^x-e^{-x}}2$
\item		$\ds \tanh x = \frac{\sinh x}{\cosh x}$
\item		$\ds \sech x = \frac{1}{\cosh x}$
\item		$\ds \csch x = \frac{1}{\sinh x}$
\item		$\ds \coth x = \frac{\cosh x}{\sinh x}$
\end{enumerate}
\end{multicols}}
\end{definition}

The hyperbolic functions are graphed in \autoref{fig:hyperbolic}. In the graphs of $\cosh x$ and $\sinh x$, graphs of $e^x/2$ and $e^{-x}/2$ are included with dashed lines. As $x$ gets ``large,'' $\cosh x$ and $\sinh x$ each act like $e^x/2$; when $x$ is a large negative number, $\cosh x$ acts like $e^{-x}/2$ whereas $\sinh x$ acts like $-e^{-x}/2$.

\mnote{\textbf{Pronunciation Note:} \\
``cosh'' rhymes with ``gosh,'' \\
``sinh'' rhymes with ``pinch,'' and\\
``tanh'' rhymes with ``ranch,'' \\
%``sech'' rhymes with ``fetch,'' and \\
%``coth'' rhymes with ``moth.''
}

Notice the domains of $\tanh x$ and $\sech x$ are $(-\infty,\infty)$, whereas both $\coth x$ and $\csch x$ have vertical asymptotes at $x=0$. Also note the ranges of these functions, especially $\tanh x$: as $x\to\infty$, both $\sinh x$ and $\cosh x$ approach $e^x/2$, hence $\tanh x$ approaches $1$.

\noindent
\begin{minipage}{\linewidth}\tagstructbegin{tag=Sect}
\centering
\begin{tikzpicture}[alt={cosh is in blue, and makes a U shape above the x axis.  For positive x, cosh is well approximated by exp(x/2).  For negative x, cosh is well approximated by exp(-x/2).}]
\begin{axis}[width=.55\textwidth,tick label style={font=\scriptsize},
axis y line=middle,axis x line=middle,name=myplot,axis on top,ymin=-10.9,ymax=10.9,
xmin=-3.5,xmax=3.5,scaled ticks=false]
\addplot [draw={\colorone},thick,smooth,domain=-3:3] {cosh(x)};
\draw (axis cs:2,-6) node {\scriptsize $f(x)=\cosh x$};
\addplot [draw={\colortwo},smooth,thick,dashed,domain=-3:3] {exp(x)/2};
\addplot [draw={\colortwo},smooth,thick,dashed,domain=-3:3] {exp(-x)/2};
\end{axis}
\node [right] at (myplot.right of origin) {\scriptsize $x$};
\node [above] at (myplot.above origin) {\scriptsize $y$};
\end{tikzpicture}%
\qquad
\begin{tikzpicture}[alt={sinh is in blue, and looks vaguely like y=x*x*x.  For positive x, sinh is well approximated by exp(x/2).  For negative x, sinh is well approximated by -exp(-x/2).}]
\begin{axis}[width=.55\textwidth,tick label style={font=\scriptsize},
axis y line=middle,axis x line=middle,name=myplot,axis on top,ymin=-10.9,ymax=10.9,
xmin=-3.5,xmax=3.5,scaled ticks=false]
\addplot [draw={\colorone},thick,smooth,domain=-3:3] {sinh(x)};
\draw (axis cs:2,-6) node {\scriptsize $f(x)=\sinh x$};
\addplot [draw={\colortwo},smooth,thick,dashed,domain=-3:3] {exp(x)/2};
\addplot [draw={\colortwo},smooth,thick,dashed,domain=-3:3] {exp(-x)/2};
\end{axis}
\node [right] at (myplot.right of origin) {\scriptsize $x$};
\node [above] at (myplot.above origin) {\scriptsize $y$};
\end{tikzpicture}
\\[20pt]
\begin{tikzpicture}[alt={tahn is in blue, has a horizontal asymptote of y=±1 as x approaches ±∞, and behaves like y=x near the origin.  coth is in red, and looks like 1+1/x for positive x and -1+1/x for negative x.}]
\begin{axis}[width=.55\textwidth,tick label style={font=\scriptsize},
axis y line=middle,axis x line=middle,name=myplot,axis on top,ytick={-2,2},
ymin=-3.5,ymax=3.5,xmin=-3.5,xmax=3.5,scaled ticks=false]
\addplot [draw={\colorone},thick,smooth,domain=-3:3] {tanh(x)};
\draw (axis cs:1.7,-1.5) node {\scriptsize $f(x)=\tanh x$};
\draw (axis cs:2.2,2) node {\scriptsize $f(x)=\coth x$};
\draw [->,>=stealth] (axis cs:1,-1) -- (axis cs:.6,.45);
\addplot [draw={\colortwo},smooth,thick,domain=-3:-.1] {1/tanh(x)};
\addplot [draw={\colortwo},smooth,thick,domain=.1:3] {1/tanh(x)};
\draw [loosely dashed] (axis cs:-3,1) -- (axis cs:3,1);
\draw [loosely dashed] (axis cs:-3,-1) -- (axis cs:3,-1);
\end{axis}
\node [right] at (myplot.right of origin) {\scriptsize $x$};
\node [above] at (myplot.above origin) {\scriptsize $y$};
\end{tikzpicture}%
\qquad
\begin{tikzpicture}[alt={sech is in blue, has a horizontal asymptote of y=0 as x approaches ±∞, and behaves like 1-x*x/2 near the origin.  csch is in red, and looks like y=1/x.}]
\begin{axis}[width=.55\textwidth,tick label style={font=\scriptsize},
axis y line=middle,axis x line=middle,name=myplot,axis on top,ytick={-3,-2,-1,1,2,3},
ymin=-3.5,ymax=3.5,xmin=-3.5,xmax=3.5,scaled ticks=false]
\addplot [draw={\colorone},thick,smooth,domain=-3:3] {1/cosh(x)};
\draw (axis cs:-2,1.5) node {\scriptsize $f(x)=\sech x$};
\draw (axis cs:2.2,2) node {\scriptsize $f(x)=\csch x$};
\addplot [draw={\colortwo},smooth,thick,domain=-3:-.1] {1/sinh(x)};
\addplot [draw={\colortwo},smooth,thick,domain=.1:3] {1/sinh(x)};
\end{axis}
\node [right] at (myplot.right of origin) {\scriptsize $x$};
\node [above] at (myplot.above origin) {\scriptsize $y$};
\end{tikzpicture}
\captionsetup{type=figure}
\caption{Graphs of the hyperbolic functions.}
\label{fig:hyperbolic}\tagstructend
\end{minipage}

\youtubeVideo{G1C1Z5aTZSQ}{Hyperbolic Functions --- The Basics}

%It is no coincidence that these functions share a name similar to the trigonometric functions. 
The following example explores some of the properties of these functions that bear remarkable resemblance to the properties of their trigonometric counterparts.

\begin{example}[Exploring properties of hyperbolic functions]\label{ex_hf1}%
Use \autoref{def:hyperbolic_functions} to rewrite the following expressions.
\begin{multicols}{2}
\begin{enumerate}\lxAddClass{columns2}
\item		$\cosh^2 x-\sinh^2x$
\item		$\tanh^2 x+\sech^2 x$
\item		$2\cosh x\sinh x$
\item		$\frac{\dd}{\dd x}\bigl(\cosh x\bigr)$
\item		$\frac{\dd}{\dd x}\bigl(\sinh x\bigr)$
\item		$\frac{\dd}{\dd x}\bigl(\tanh x\bigr)$
\end{enumerate}
\end{multicols}
\solution
\begin{enumerate}
\item \hfill$\begin{aligned}[t]
 \cosh^2x-\sinh^2x
 &= \left(\frac{e^x+e^{-x}}2\right)^2 -\left(\frac{e^x-e^{-x}}2\right)^2\\
 &= \frac{e^{2x}+2e^xe^{-x} + e^{-2x}}4 - \frac{e^{2x}-2e^xe^{-x} + e^{-2x}}4\\
 &= \frac44=1.
\end{aligned}$\hfill\null\\
So $\cosh^2 x-\sinh^2x=1$.

\item \hfill$\begin{aligned}[t]
 \tanh^2 x+\sech^2 x
 &=\frac{\sinh^2x}{\cosh^2 x} + \frac{1}{\cosh^2 x} \\
 &= \frac{\sinh^2x+1}{\cosh^2 x}\qquad \text{\small Now use identity from \#1.}\\
 &= \frac{\cosh^2 x}{\cosh^2 x} = 1.
\end{aligned}$\hfill\null\\
So $\tanh^2 x+\sech^2 x=1$.

\item \hfill$\begin{aligned}[t]
	2\cosh x\sinh x
	&= 2\left(\frac{e^x+e^{-x}}2\right)\left(\frac{e^x-e^{-x}}2\right) \\
	&= 2 \cdot\frac{e^{2x} - e^{-2x}}4\\
	&= \frac{e^{2x} - e^{-2x}}2 = \sinh (2x).\\
\end{aligned}$\hfill\null\\
Thus $2\cosh x\sinh x = \sinh (2x)$.

\item \hfill$\begin{aligned}[t]
	\frac{\dd}{\dd x}\bigl(\cosh x\bigr)
	&= \frac{\dd}{\dd x}\left(\frac{e^x+e^{-x}}2\right) \\
	&= \frac{e^x-e^{-x}}2\\
	&= \sinh x.
\end{aligned}$\hfill\null\\
So $\frac{\dd}{\dd x}\bigl(\cosh x\bigr) = \sinh x.$
	
\item \hfill$\begin{aligned}[t]
	\frac{\dd}{\dd x}\bigl(\sinh x\bigr)
	&= \frac{\dd}{\dd x}\left(\frac{e^x-e^{-x}}2\right) \\
	&= \frac{e^x+e^{-x}}2\\
	&= \cosh x.
\end{aligned}$\hfill\null\\
So $\frac{\dd}{\dd x}\bigl(\sinh x\bigr) = \cosh x.$

\item \hfill$\begin{aligned}[t]
	\frac{\dd}{\dd x}\bigl(\tanh x\bigr)
	&= \frac{\dd}{\dd x}\left(\frac{\sinh x}{\cosh x}\right) \\
	&= \frac{\cosh x \cosh x - \sinh x \sinh x}{\cosh^2 x}\\
	&= \frac{1}{\cosh^2 x}\\
	&= \sech^2 x.
\end{aligned}$\hfill\null\\
So $\frac{\dd}{\dd x}\bigl(\tanh x\bigr) = \sech^2 x.$
\end{enumerate}
\end{example}

The following Key Idea summarizes many of the important identities relating to hyperbolic functions. Each can be verified by referring back to \autoref{def:hyperbolic_functions}.

{
\tcbset{grow to right by=16em}
\begin{keyidea}[Useful Hyperbolic Function Properties]\label{idea:hyperbolic_identities}%
\lxAddClass{columns3}
\begin{minipage}[t]{.33\linewidth}
\textbf{Basic Identities}\par
\index{hyperbolic function!identities}\index{hyperbolic function!derivatives}\index{hyperbolic function!integrals}\index{derivative!hyperbolic funct.}\index{integration!hyperbolic funct.}%
\begin{enumerate}
\item $\cosh^2x-\sinh^2x=1$
\item	$\tanh^2x+\sech^2x=1$
\item	$\coth^2x-\csch^2x = 1$
\item	$\cosh 2x=\cosh^2x+\sinh^2x$
\item	$\sinh 2x = 2\sinh x\cosh x$
\item	$\ds\cosh^2x = \frac{\cosh 2x+1}{2}$
\item $\ds \sinh^2x=\frac{\cosh 2x-1}{2}$
\end{enumerate}
\end{minipage}%
\begin{minipage}[t]{.33\linewidth}
\textbf{Derivatives}
\begin{enumerate}
\item $\frac{\dd}{\dd x}\bigl(\cosh x\bigr) = \sinh x$
\item $\frac{\dd}{\dd x}\bigl(\sinh x\bigr) = \cosh x$
\item $\frac{\dd}{\dd x}\bigl(\tanh x\bigr) = \sech^2 x$
\item $\frac{\dd}{\dd x}\bigl(\sech x\bigr) = -\sech x\tanh x$
\item $\frac{\dd}{\dd x}\bigl(\csch x\bigr) = -\csch x\coth x$
\item $\frac{\dd}{\dd x}\bigl(\coth x\bigr) = -\csch^2x$
\end{enumerate}
\end{minipage}%
\begin{minipage}[t]{.33\linewidth}
\textbf{Integrals}
\begin{enumerate}
\item $\ds\int \cosh x\dd x = \sinh x+C$
\item $\ds\int \sinh x\dd x = \cosh x+C$
\item $\ds\int \tanh x\dd x = \ln(\cosh x) +C$
\item $\ds\int \coth x\dd x = \ln\abs{\sinh x}+C$
\end{enumerate}
\end{minipage}
\end{keyidea}
}

We practice using \autoref{idea:hyperbolic_identities}.

\begin{example}[Derivatives and integrals of hyperbolic functions]\label{ex_hf2}%
Evaluate the following derivatives and integrals.\\
\begin{minipage}[t]{.5\linewidth}
\begin{enumerate}
\item		$\ds\frac{\dd}{\dd x}\bigl(\cosh 2x\bigr)$
\item		$\ds\int \sech^2(7t-3)\dd t$
\end{enumerate}
\end{minipage}%
\begin{minipage}[t]{.5\linewidth}
\begin{enumerate}\addtocounter{enumi}{2}
\item		$\ds \int_0^{\ln 2} \cosh x\dd x$
\end{enumerate}
\end{minipage}
\solution
\begin{enumerate}
\item		Using the Chain Rule directly, we have $\frac{\dd}{\dd x} \bigl(\cosh 2x\bigr) = 2\sinh 2x$.

Just to demonstrate that it works, let's also use the Basic Identity found in \autoref{idea:hyperbolic_identities}: $\cosh 2x = \cosh^2x+\sinh^2x$.
\begin{align*}
\frac{\dd}{\dd x}\bigl(\cosh 2x\bigr) = \frac{\dd}{\dd x}\bigl(\cosh^2x+\sinh^2x\bigr)
&= 2\cosh x\sinh x+ 2\sinh x\cosh x\\
&= 4\cosh x\sinh x.
\end{align*}
Using another Basic Identity, we can see that $4\cosh x\sinh x = 2\sinh 2x$. We get the same answer either way.

\item	  We employ substitution, with $u = 7t-3$ and $\dd u = 7\dd t$. Applying
\autoref{idea:hyperbolic_identities} we have:
\[\int \sech^2 (7t-3)\dd t = \frac17\tanh (7t-3) + C.\]

\item \[\int_0^{\ln 2} \cosh x\dd x = \sinh x\Big|_0^{\ln 2} = \sinh (\ln 2) - \sinh 0 = \sinh(\ln 2).\]

We can simplify this last expression as $\sinh x$ is based on exponentials:
\[\sinh(\ln 2) = \frac{e^{\ln 2}-e^{-\ln 2}}2 = \frac{2-1/2}{2} = \frac34.\]
\end{enumerate}
\end{example}

\subsection{Inverse Hyperbolic Functions}

Just as the inverse trigonometric functions are useful in certain integrations, the inverse hyperbolic functions are useful with others. \autoref{fig:hfinverse2} shows the restrictions on the domains to make each function one-to-one and the resulting domains and ranges of their inverse functions. Their graphs are shown in \autoref{fig:hfinverse1}.\index{hyperbolic function!inverse}

Because the hyperbolic functions are defined in terms of exponential functions, their inverses can be expressed in terms of logarithms as shown in \autoref{idea:hyperbolic_log}. It is often more convenient to use $\sinh^{-1}x$ than $\ln\bigl(x+\sqrt{x^2+1}\bigr)$, especially when one is working on theory and does not need to compute actual values. On the other hand, when computations are needed, technology is often helpful but many hand-held calculators lack a \emph{convenient} $\sinh^{-1}x$ button. (Often it can be accessed under a menu system, but not conveniently.) In such a situation, the logarithmic representation is useful. The reader is not encouraged to memorize these, but rather know they exist and know how to use them when needed.

\noindent\begin{minipage}[t]{\linewidth}\tagstructbegin{tag=Sect}\noindent%
\captionsetup{type=figure}%
\flushinner{%
\small
\tagpdfsetup{table/header-rows={1}}
\begin{tabular}{ l c c @{\hspace{2em}} c c c }
Function & Domain & Range & Function & Domain & Range \\ \cmidrule(r{2em}){1-3} \cmidrule(l{-1em}){4-6}
$\cosh x$ & $[0,\infty)$ & $[1,\infty)$ &
 $\cosh^{-1} x$ & $[1,\infty)$ & $[0,\infty)$ \\
$\sinh x$ & $(-\infty,\infty)$ & $(-\infty,\infty)$ &
 $\sinh^{-1} x$ & $(-\infty,\infty)$ & $(-\infty,\infty)$\\
$\tanh x$ & $(-\infty,\infty)$ & $(-1,1)$ &
 $\tanh^{-1} x$ & $(-1,1)$ & $(-\infty,\infty)$\\
$\sech x$ & $[0,\infty)$ & $(0,1]$ & $\sech^{-1} x$ & $(0,1]$ & $[0,\infty)$\\
$\csch x$ & $(-\infty,0) \cup (0,\infty)$ & $(-\infty,0) \cup (0,\infty)$ &
 $\csch^{-1} x$ & $(-\infty,0) \cup (0,\infty)$ & $(-\infty,0) \cup (0,\infty)$\\
$\coth x$ & $(-\infty,0) \cup (0,\infty)$ & $(-\infty,-1) \cup (1,\infty)$ &
 $\coth^{-1} x$ & $(-\infty,-1) \cup (1,\infty)$ & $(-\infty,0) \cup (0,\infty)$
\end{tabular}}
\caption{Domains and ranges of the hyperbolic and inverse hyperbolic functions.}
\label{fig:hfinverse2}\tagstructend
\end{minipage}

\noindent\begin{minipage}[t]{\linewidth}\tagstructbegin{tag=Sect}\noindent%
\captionsetup{type=figure}%
%\flushinner{%
\addtolength{\tabcolsep}{6pt}
\tagpdfsetup{table/tagging=presentation}
\begin{tabular}{cc}
\begin{tikzpicture}[alt={cosh for positive x is in blue.  Its inverse is in red, and has been reflected across the dashed line y=x.}]
\begin{axis}[width=1.16\marginparwidth,tick label style={font=\scriptsize},
axis y line=middle,axis x line=middle,name=myplot,axis on top,axis equal,
ymin=-.9,ymax=10.9,xmin=-.9,xmax=10.9,scaled ticks=false]
\addplot [draw={\colorone},thick,smooth,domain=0:3] {cosh(x)};
\draw (axis cs:8,4) node {\scriptsize $y=\cosh^{-1} x$};
\draw (axis cs:5.6,10) node {\scriptsize $y=\cosh x$};
\addplot [draw={\colortwo},smooth,thick,domain=0:3] ({cosh(x)},x);
\addplot [dashed,domain=-.5:10] {x};
\end{axis}
\node [right] at (myplot.right of origin) {\scriptsize $x$};
\node [above] at (myplot.above origin) {\scriptsize $y$};
\end{tikzpicture}
&
\begin{tikzpicture}[alt={sinh is in blue.  Its inverse is in red, and has been reflected across the dashed line y=x.}]
\begin{axis}[width=1.16\marginparwidth,tick label style={font=\scriptsize},
axis y line=middle,axis x line=middle,name=myplot,axis on top,axis equal,
ymin=-10.9,ymax=10.9,xmin=-10.9,xmax=10.9,scaled ticks=false]
\addplot [draw={\colorone},thick,smooth,domain=-3:3] {sinh(x)};
\draw (axis cs:-6,7) node {\scriptsize $y=\sinh x$};
\draw (axis cs:6,-5) node {\scriptsize $y=\sinh^{-1} x$};
\draw[->,>=stealth] (axis cs:-1,7) -- (axis cs:2,7);
\draw[->,>=stealth] (axis cs:7,-3) -- (axis cs:7,2);
\addplot [draw={\colortwo},smooth,thick,domain=-3:3] ({sinh(x)},x);
\addplot [dashed,domain=-10:10] {x};
\end{axis}
\node [right] at (myplot.right of origin) {\scriptsize $x$};
\node [above] at (myplot.above origin) {\scriptsize $y$};
\end{tikzpicture}
\\[15pt]
\begin{tikzpicture}[alt={Inverse tanh is in blue, and looks like (trigonometric) tangent.  Inverse coth is in red, and looks like 1/(x-1) for x>1, and 1/(x+1) for x<-1.}]
\begin{axis}[width=1.16\marginparwidth,tick label style={font=\scriptsize},
axis y line=middle,axis x line=middle,name=myplot,axis on top,xtick={-2,2},
ytick={-2,2},ymin=-3.2,ymax=3.2,xmin=-3.2,xmax=3.2,scaled ticks=false,axis equal]
\addplot [draw={\colortwo},thick,smooth,domain=.348:3] ({1/tanh(x)},x);
\addplot [draw={\colortwo},thick,smooth,domain=-3:-.348] ({1/tanh(x)},x);
\draw (axis cs:2.2,1.5) node {\tiny $y=\coth^{-1} x$};
\draw (axis cs:2.2,-1.5) node {\tiny $y=\tanh^{-1} x$};
\draw [->,>=stealth] (axis cs:1.8,-1.3) -- (axis cs:.2,-.2);
\draw [loosely dashed] (axis cs:-1,-3)--(axis cs:-1,3);
\draw [loosely dashed] (axis cs:1,-3)--(axis cs:1,3);
\addplot [draw={\colorone},smooth,thick,domain=-3.8:3.8] ({tanh(x)},x);
\end{axis}
\node [right] at (myplot.right of origin) {\scriptsize $x$};
\node [above] at (myplot.above origin) {\scriptsize $y$};
\end{tikzpicture}
&
\begin{tikzpicture}[alt={Inverse sech is in blue, starting with a vertical asymptote as x approaches 0 from the right, and going to the point (1,0).  Inverse csch is in red, and looks like 1/x.}]
\begin{axis}[width=1.16\marginparwidth,tick label style={font=\scriptsize},
axis y line=middle,axis x line=middle,name=myplot,axis on top,
xtick={-3,-2,-1,1,2,3},ytick={-3,-2,-1,1,2,3},
ymin=-3.5,ymax=3.5,xmin=-3.5,xmax=3.5,scaled ticks=false,axis equal]
\draw (axis cs:1.5,-1.5) node {\tiny $y=\sech^{-1} x$};
\draw (axis cs:-2.2,-1.6) node {\tiny $y=\csch^{-1} x$};
\draw [->,>=latex] (axis cs:1,-1.2) -- (axis cs:1,-.2);
\addplot [draw={\colortwo},smooth,thick,domain=-3:-.3275] ({1/sinh(x)},x);
\addplot [draw={\colortwo},smooth,thick,domain=.3275:3] ({1/sinh(x)},x);
\addplot [draw={\colorone},thick,smooth,domain=0:3] ({1/cosh(x)},x);
\end{axis}
\node [right] at (myplot.right of origin) {\scriptsize $x$};
\node [above] at (myplot.above origin) {\scriptsize $y$};
\end{tikzpicture}
\end{tabular}
%}
\caption{Graphs of the hyperbolic functions and their inverses.}\label{fig:hfinverse1}\tagstructend
\end{minipage}

Now let's consider the inverses of the hyperbolic functions. We begin with the function $f(x)=\sinh x$. Since $\fp(x)=\cosh x>0$ for all real $x$, $f$ is increasing and must be one-to-one.
% todo Tim once we have an example of finding an inverse in text/07_Inverse_Functions.tex
% We proceed as in \autoref{sec:inv_funcs}:
\begin{align*}\allowdisplaybreaks
y&=\frac{e^x-e^{-x}}2\\
2y&=e^x-e^{-x} \qquad\text{(now multiply by $e^x$)}\\
2ye^x&=e^{2x}-1 \qquad\text{(a quadratic form )}\\
\left(e^x\right)^2-2ye^x-1&=0 \qquad\text{(use the quadratic formula)}\\
e^x&=\frac{2y\pm\sqrt{4y^2+4}}2\\
e^x&=y\pm\sqrt{y^2+1} \qquad\text{(use the fact that $e^x>0$)}\\
e^x&=y+\sqrt{y^2+1}\\
x&=\ln(y+\sqrt{y^2+1})\\
\end{align*}
Finally, interchange the variable to find that
\[\sinh^{-1} x=\ln(x+\sqrt{x^2+1}).\]
In a similar manner we find that the inverses of the other hyperbolic functions are given by:

{
\tcbset{grow to right by=5em} % 5 sufficient
\begin{keyidea}[Logarithmic definitions of Inverse Hyperbolic Functions]\label{idea:hyperbolic_log}%
\mbox{}\\[-2\baselineskip]%
\index{hyperbolic function!inverse!logarithmic def.}%
\begin{multicols}{2}
\begin{enumerate}\lxAddClass{columns2}
\item $\ds\cosh^{-1}x=\ln\bigl(x+\sqrt{x^2-1}\bigr)$;\\\null\qquad $x\geq1$
\item $\ds\tanh^{-1}x=\frac12\ln\left(\frac{1+x}{1-x}\right)$;\\\null\qquad$\abs x<1$
\item $\ds\sech^{-1}x=\ln\left(\frac{1+\sqrt{1-x^2}}x\right)$;\\\null\qquad$0<x\leq1$
\item $\ds\sinh^{-1}x=\ln\bigl(x+\sqrt{x^2+1}\bigr)$\\\mbox{}
\item $\ds\coth^{-1}x=\frac12\ln\left(\frac{x+1}{x-1}\right)$;\\\null\qquad$\abs x>1$
\item $\ds\csch^{-1}x=\ln\left(\frac1x+\frac{\sqrt{1+x^2}}{\abs x}\right)$;\\\null\qquad $x\neq0$
\end{enumerate}
\end{multicols}
\end{keyidea}
}

The following Key Ideas give the derivatives and integrals relating to the inverse hyperbolic functions. In \autoref{idea:hyperbolic_inverse_integrals}, both the inverse hyperbolic and logarithmic function representations of the antiderivative are given, based on \autoref{idea:hyperbolic_log}. Again, these latter functions are often more useful than the former.
%Note how inverse hyperbolic functions can be used to solve integrals we used Trigonometric Substitution to solve in \autoref{sec:trig_sub}.



%\mtable{Logarithmic definitions of the inverse hyperbolic functions.}{fig:hfinverse5}{%
%\begin{align*}
%\cosh^{-1}x&=\ln\bigl(x+\sqrt{x^2-1}\bigr);\ x\geq1\\
%\sinh^{-1}x &= \ln\bigl(x+\sqrt{x^2+1}\bigr)\\
%\tanh^{-1}x &= \frac12\ln\left(\frac{1+x}{1-x}\right);\ \abs x<1\\
%\sech^{-1}x &= \ln\left(\frac{1+\sqrt{1-x^2}}x\right);\ 0<x\leq1\\
%\csch^{-1}x &= \ln\left(\frac1x+\frac{\sqrt{1+x^2}}{\abs x}\right);\ x\neq0\\
%\coth^{-1}x &= \frac12\ln\left(\frac{x+1}{x-1}\right);\ \abs x>1
%\end{align*}
%}

{
\tcbset{grow to right by=3em}
\begin{keyidea}[Derivatives Involving Inverse Hyperbolic Functions]\label{idea:hyperbolic_inverse_derivatives}%
\index{hyperbolic function!inverse!derivative}\index{derivative!inverse hyper.}
\mbox{}\\[-2\baselineskip]\parbox[t]{\linewidth}{%
\begin{multicols}{2}
\begin{enumerate}\lxAddClass{columns2}
\item $\ds\frac\dd{\dd x}\bigl(\cosh^{-1}x\bigr)=\frac1{\sqrt{x^2-1}}$;\\\null\qquad$x>1$
\item $\ds\frac\dd{\dd x}\bigl(\sinh^{-1}x\bigr)=\frac1{\sqrt{x^2+1}}$\\\vphantom{$x\ne0$}
\item $\ds\frac\dd{\dd x}\bigl(\tanh^{-1}x\bigr)=\frac1{1-x^2}$;\\\null\qquad$\abs x<1$
\item $\ds\frac\dd{\dd x}\bigl(\sech^{-1}x\bigr)=\frac{-1}{x\sqrt{1-x^2}}$;\\\null\qquad$0<x<1$
\item $\ds\frac\dd{\dd x}\bigl(\csch^{-1}x\bigr)=\frac{-1}{\abs x\sqrt{1+x^2}}$;\\\null\qquad$x\neq0$
\item $\ds\frac\dd{\dd x}\bigl(\coth^{-1}x\bigr)=\frac1{1-x^2}$;\\\null\qquad$\abs x>1$
\end{enumerate}
\end{multicols}}
\end{keyidea}}

{
\tcbset{grow to right by=8.1em}
% 8.1em sufficient
\begin{keyidea}[Integrals Involving Inverse Hyperbolic Functions]\label{idea:hyperbolic_inverse_integrals}%
\index{integration!inverse hyper.}\index{hyperbolic function!inverse!integration}%
\begin{anywhereenum}
\begin{align*}
&\item \int\frac1{\sqrt{x^2-a^2}}\dd x && =\cosh^{-1}\left(\frac xa\right)+C;~0<a<x && =\ln\left(x+\sqrt{x^2-a^2}\right)+C \\
&\item \int\frac1{\sqrt{x^2+a^2}}\dd x && =\sinh^{-1}\left(\frac xa\right)+C;~a>0 && =\ln\left(x+\sqrt{x^2+a^2}\right)+C \\
&\item \int\frac1{a^2-x^2}\dd x && ={\begin{cases}
\frac1a\tanh^{-1}\left(\frac xa\right)+C & \abs{x}<\abs{a} \\
\frac1a\coth^{-1}\left(\frac xa\right)+C & \abs{a}<\abs{x}
\end{cases}} && =\frac1{2a}\ln\abs{\frac{a+x}{a-x}}+C \\
&\item \int\frac1{x\sqrt{a^2-x^2}}\dd x\!\!\! &&
=-\frac1a\sech^{-1}\frac{\abs x}a+C;~0<\abs x<a && =\frac1a\ln\abs{\frac x{a+\sqrt{a^2-x^2}}}+C \\
% $\ds{}=-\frac1a\sech^{-1}\left(\frac xa\right)+C$ & $\ds{}=\frac1a\ln\left(\frac{x}{a+\sqrt{a^2-x^2}}\right)+C$ \\[\defaultaddspace]
%&\null\hfill{$\scriptstyle0<x<a$}\\[\defaultaddspace]
&\item \int\frac1{x\sqrt{x^2+a^2}}\dd x\!\!\! && =-\frac1a\csch^{-1}\frac{\abs x}a+C;~x\neq 0,\ a>0\!\!\! && =
\frac1a\ln\abs{\frac x{a+\sqrt{a^2+x^2}}}
%-\frac1a\ln\left(\frac ax+\frac{\sqrt{x^2+a^2}}{a\abs x}\right)
% sqrt(1+(x/a)^2)/|x/a| = sqrt(a^2+x^2)/|ax|
+C
\end{align*}
\end{anywhereenum}
\end{keyidea}}

We practice using the derivative and integral formulas in the following example.

\begin{example}[Derivatives and integrals involving inverse hyperbolic functions]\label{ex_hf3}%
Evaluate the following.
\begin{multicols}{2}
\begin{enumerate}\lxAddClass{columns2}
\item	$\ds \frac{\dd}{\dd x}\left[\cosh^{-1}\left(\frac{3x-2}{5}\right)\right]$
\item	$\ds \int\frac{1}{x^2-1}\dd x$
\item	$\ds \int \frac{1}{\sqrt{9x^2+10}}\dd x$
\item[]
\end{enumerate}
\end{multicols}
\solution
\begin{enumerate}
\item	Applying \autoref{idea:hyperbolic_inverse_derivatives} with the Chain Rule gives:
		\[\frac{\dd}{\dd x}\left[\cosh^{-1}\left(\frac{3x-2}5\right)\right] = \frac{1}{\sqrt{\left(\frac{3x-2}5\right)^2-1}}\cdot\frac35.\]

\item		Multiplying the numerator and denominator by $(-1)$ gives 
a %: $\ds \int \frac{1}{x^2-1}\dd x = \int \frac{-1}{1-x^2}\dd x$. The
second integral can be solved with a direct application of item \#3 from \autoref{idea:hyperbolic_inverse_integrals}, with $a=1$. Thus
\begin{align}
\int \frac{1}{x^2-1}\dd x &= -\int \frac{1}{1-x^2}\dd x \notag \\
		&= \begin{cases}-\tanh^{-1}\left(x\right)+C & x^2<1 \\
-\coth^{-1}\left(x\right)+C & 1<x^2 \end{cases} \notag\\
     &=-\frac12\ln\abs{\frac{x+1}{x-1}}+C\notag\\
     &=\frac12\ln\abs{\frac{x-1}{x+1}}+C.\label{eq:hf3}
     \end{align}

%We should note that this exact problem was solved at the beginning of \autoref{sec:partial_fraction}. In that example the answer was given as $\frac12\ln\abs{x-1}-\frac12\ln\abs{x+1}+C.$ Note that this is equivalent to the answer given in \autoeqref{eq:hf3}, as $\ln(a/b) = \ln a - \ln b$.

%The key to linking the two seemingly different answers together is \autoref{fig:hfinverse5}, where the logarithmic definitions of the inverse hyperbolic functions are given. Note that the definitions of $\tanh^{-1}x$ and $\coth^{-1}x$ are very similar; the conditions placed on $\abs x$ ensure that the argument of $\ln$ is always positive. Thus one could say 
%\[\frac12\ln\abs{\frac{x+1}{x-1}} = \begin{cases}\tanh^{-1}x+C & \abs x<1 \\ \\
%\coth^{-1}x+C & \abs x>1 \end{cases}.\]
%
%We reconcile the two answers by returning to \autoeqref{eq:hf3} and continuing:
%\begin{align*}
%\int \frac{1}{x^2-1}\dd x &= \int \frac{-1}{1-x^2}\dd x \\
%			&= \begin{cases}-\frac1a\tanh^{-1}\left(\frac xa\right)+C & x^2<a^2 \\ \\
%-\frac1a\coth^{-1}\left(\frac xa\right)+C & a^2<x^2 \end{cases}\\
%			&= -\frac12\ln\abs{\frac{x+1}{x-1}}+C \\
%			&= -\frac12\ln\abs{x+1} + \frac12\ln\abs{x-1} +C,
%\end{align*}
%matching the answer previously obtained.

\item		This requires a substitution, then item \#2 of \autoref{idea:hyperbolic_inverse_integrals} can be applied.

Let $u = 3x$, hence $\dd u = 3\dd x$. We have 
\begin{align*}
	\int \frac{1}{\sqrt{9x^2+10}}\dd x
	&= \frac13\int\frac{1}{\sqrt{u^2+10}}\dd u.
	\intertext{Note \intertextmath{a^2=10}, hence \intertextmath{a = \sqrt{10}}. Now apply the integral rule.}
	 &= \frac13 \sinh^{-1}\left(\frac{3x}{\sqrt{10}}\right) + C \\
	 &= \frac13 \ln \abs{3x+\sqrt{9x^2+10}}+C.
\end{align*}
\end{enumerate}
\end{example}

This section covers a lot of ground. New functions were introduced, along with some of their fundamental identities, their derivatives and antiderivatives, their inverses, and the derivatives and antiderivatives of these inverses. Four Key Ideas were presented, each including quite a bit of information.

Do not view this section as containing a source of information to be memorized, but rather as a reference for future problem solving. \autoref{idea:hyperbolic_inverse_integrals} contains perhaps the most useful information. Know the integration forms it helps evaluate and understand how to use the inverse hyperbolic answer and the logarithmic answer.

The next section takes a brief break from demonstrating new integration techniques. It instead demonstrates a technique of evaluating limits that return indeterminate forms. This technique will be useful in \autoref{sec:improper_integration}, where limits will arise in the evaluation of certain definite integrals.


\printexercises{exercises/06-05-exercises}

\section{L'Hôpital's Rule}\label{sec:lhopitals_rule}

This section is concerned with a technique for evaluating certain limits that will be useful in later chapters.

Our treatment of limits exposed us to ``0/0'', an indeterminate form. If both $\ds \lim_{x\to c}f(x)$ and $\ds \lim_{x\to c} g(x)$ are zero, we do not conclude that $\ds \lim_{x\to c} f(x)/g(x)$ is $0/0$; rather, we use $0/0$ as notation to describe the fact that both the numerator and denominator approach 0. The expression 0/0 has no numeric value; other work must be done to evaluate the limit.

Other indeterminate forms exist; they are: $\infty/\infty$, $0\cdot\infty$, $\infty-\infty$, $0^0$, $1^\infty$ and $\infty^0$. Just as ``0/0'' does not mean ``divide 0 by 0,'' the expression ``$\infty/\infty$'' does not mean ``divide infinity by infinity.'' Instead, it means ``a quantity is growing without bound and is being divided by another quantity that is growing without bound.'' We cannot determine from such a statement what value, if any, results in the limit. Likewise, ``$0\cdot \infty$'' does not mean ``multiply zero by infinity.'' Instead, it means ``one quantity is shrinking to zero, and is being multiplied by a quantity that is growing without bound.'' We cannot determine from such a description what the result of such a limit will be.

This section introduces L'Hôpital's Rule, a method of resolving limits that produce the indeterminate forms 0/0 and $\infty/\infty$. We'll also show how algebraic manipulation can be used to convert other indeterminate expressions into one of these two forms so that our new rule can be applied.

\begin{theorem}[L'Hôpital's Rule, Part 1]\label{thm:LHR_1}%
Let $f$ and $g$ be differentiable functions on an open interval $I$ containing $a$.
\begin{enumerate}
\item If $\ds\lim_{x\to a}f(x)=0$, $\ds\lim_{x\to a}g(x)=0$, and $g\primeskip'(x)\neq 0$ except possibly at $x=a$, then \[\lim_{x\to a}\frac{f(x)}{g(x)}=\lim_{x\to a} \frac{\fp(x)}{g\primeskip'(x)},\]
assuming that the limit on the right exists.
\item If  $\ds\lim_{x\to a}f(x)=\pm\infty$ and $\ds\lim_{x\to a}g(x)=\pm\infty$, then \[\lim_{x\to a}\frac{f(x)}{g(x)}=\lim_{x\to a} \frac{\fp(x)}{g\primeskip'(x)},\]
assuming that the limit on the right exists.
\index{LHopitals Rule@L'Hôpital's Rule}
\end{enumerate}
\end{theorem}

A similar statement holds if we just look at the one sided limits $\ds\lim_{x\to a^-}$ and $\ds\lim_{x\to a^+}$.

\begin{theorem}[L'Hôpital's Rule, Part 2]\label{thm:LHR_2}%
Let $f$ and $g$ be differentiable functions on the open interval $(c,\infty)$ for some value $c$ and $g\primeskip'(x)\neq0$ on $(c,\infty)$.
\begin{enumerate}
\item If $\ds\lim_{x\to\infty}f(x)=0$ and $\ds\lim_{x\to \infty}g(x)=0$, then
\[\lim_{x\to\infty}\frac{f(x)}{g(x)}=\lim_{x\to\infty}\frac{\fp(x)}{g\primeskip'(x)},\]
assuming that the limit on the right exists.
\item If  $\ds\lim_{x\to \infty}f(x)=\pm\infty$ and $\ds\lim_{x\to \infty}g(x)=\pm\infty$, then
\[\lim_{x\to\infty}\frac{f(x)}{g(x)}=\lim_{x\to\infty}\frac{\fp(x)}{g\primeskip'(x)},\]
assuming that the limit on the right exists.
\end{enumerate}
Similar statements can be made where $x$ approaches $-\infty$.
\index{LHopitals Rule@L'Hôpital's Rule}
\end{theorem}

We demonstrate the use of L'Hôpital's Rule in the following examples; we will often use ``LHR'' as an abbreviation of ``L'Hôpital's Rule.''

\begin{example}[Using L'Hôpital's Rule]\label{ex_lhr1}%
Evaluate the following limits, using L'Hôpital's Rule as needed.\\
\begin{minipage}[t]{.5\textwidth}
\begin{enumerate}
\item		$\ds \lim_{x\to0}\frac{\sin x}x$
\item		$\ds \lim_{x\to 1}\frac{\sqrt{x+3}-2}{1-x}$
\item		$\ds \lim_{x\to0}\frac{x^2}{1-\cos x}$
\end{enumerate}
\end{minipage}%
\begin{minipage}[t]{.5\textwidth}
\begin{enumerate}\addtocounter{enumi}{3}
	\item	$\ds\lim_{x\to-3}\frac{x^3+27}{x^2+9}$
	\item	$\ds\lim_{x\to\infty}\frac{3x^2-100x+2}{4x^2+5x-1000}$
	\item	$\ds\lim_{x\to\infty}\frac{e^x}{x^3}$
\end{enumerate}
\end{minipage}
\solution
\begin{enumerate}
	\item	This has the indeterminate form $0/0$. We proved this limit is 1 in \autoref{ex_limit_sinx_prove} using the Squeeze Theorem. Here we use L'Hôpital's Rule to show its power.
\[\lim_{x\to0}\frac{\sin x}x \LHequals \lim_{x\to0} \frac{\cos x}{1}=1.\]
While this seems easier than using the Squeeze Theorem to find this limit, we note that applying L'Hôpital's Rule here requires us to know the derivative of $\sin x$. We originally encountered this limit when we were trying to find that derivative.

	\item	This has the indeterminate form $0/0$.
\[\lim_{x\to 1}\frac{\sqrt{x+3}-2}{1-x} 	 \LHequals \lim_{x \to 1} \frac{\frac12(x+3)^{-1/2}}{-1} =-\frac 14.\]

	\item	This has the indeterminate form $0/0$.
\[\lim_{x\to 0}\frac{x^2}{1-\cos x}  \LHequals  \lim_{x\to 0} \frac{2x}{\sin x}.\]
This latter limit also evaluates to the $0/0$ indeterminate form. To evaluate it, we apply L'Hôpital's Rule again.
\[
 \lim_{x\to 0} \frac{2x}{\sin x}
 \LHequals \lim_{x\to 0} \frac{2}{\cos x} = 2 .
\]
Thus $\ds \lim_{x\to0}\frac{x^2}{1-\cos x}=2.$

	\item \hfill$\ds\lim_{x\to-3}\frac{x^3+27}{x^2+9}
	 =\frac0{18}=0$\hfill\null\smallskip\\
We cannot use L'Hôpital's Rule in this case because the original limit does not return an indeterminate form, so L'Hôpital's Rule does not apply. In fact, the inappropriate use of L'Hôpital's Rule here would result in the incorrect limit $-\frac92$.

	\item	We can evaluate this limit already using \autoref{thm:lim_rational_fn_at_infty}; the answer is 3/4. We apply L'Hôpital's Rule to demonstrate its applicability.
\[
 \lim_{x\to\infty} \frac{3x^2-100x+2}{4x^2+5x-1000}
 \LHequals \lim_{x\to\infty} \frac{6x-100}{8x+5}
 \LHequals \lim_{x\to\infty} \frac68 = \frac34.
\]

	\item	$\ds\lim_{x\to \infty}\frac{e^x}{x^3} \LHequals \lim_{x\to\infty} \frac{e^x}{3x^2} \LHequals \lim_{x\to\infty} \frac{e^x}{6x} \LHequals \lim_{x\to\infty} \frac{e^x}{6} = \infty.$

Recall that this means that the limit does not exist; as $x$ approaches $\infty$, the expression $e^x/x^3$ grows without bound. We can infer from this that $e^x$ grows ``faster'' than $x^3$; as $x$ gets large, $e^x$ is far larger than $x^3$. (This has important implications in computing when considering efficiency of algorithms.)
\end{enumerate}
\end{example}

\subsection{Indeterminate Forms \texorpdfstring{$0\cdot\infty$ and $\infty-\infty$}{0·∞ and ∞-∞}}

L'Hôpital's Rule can only be applied to ratios of functions. When faced with an indeterminate form such as $0\cdot\infty$ or $\infty-\infty$, we can sometimes apply algebra to rewrite the limit so that L'Hôpital's Rule can be applied. We demonstrate the general idea in the next example.
\index{limit!indeterminate form}\index{indeterminate form}

\youtubeVideo{kEnwac_9lyg}{L'Hôpital's Rule --- Indeterminate Powers}

\begin{example}[Applying L'Hôpital's Rule to other indeterminate forms]\label{ex_LHR3}%
Evaluate the following limits.\\
\begin{minipage}[t]{.5\textwidth}
\begin{enumerate}
\item		$\ds \lim_{x\to0^+} x\cdot e^{1/x}$
\item		$\ds \lim_{x\to0^-} x\cdot e^{1/x}$
\end{enumerate}
\end{minipage}%
\begin{minipage}[t]{.5\textwidth}
\begin{enumerate}\addtocounter{enumi}{2}
\item		$\ds \lim_{x\to\infty} \left(\ln(x+1)-\ln x\right)$
\item		$\ds \lim_{x\to\infty} \left(x^2-e^x\right)$
\end{enumerate}
\end{minipage}
\solution
\begin{enumerate}
	\item	As $x\to 0^+$, note that $x\to 0$ and $e^{1/x}\to \infty$. Thus we have the indeterminate form $0\cdot\infty$. We rewrite the expression $x\cdot e^{1/x}$ as $\ds\frac{e^{1/x}}{1/x}$; now, as $x\to 0^+$, we get the indeterminate form $\infty/\infty$ to which L'Hôpital's Rule can be applied. 
\[
\lim_{x\to0^+} x\cdot e^{1/x} = \lim_{x\to 0^+} \frac{e^{1/x}}{1/x} \LHequals \lim_{x\to 0^+}\frac{(-1/x^2)e^{1/x}}{-1/x^2} =\lim_{x\to 0^+}e^{1/x} =\infty.
\]

Interpretation: $e^{1/x}$ grows ``faster'' than $x$ shrinks to zero, meaning their product grows without bound.

	\item	As $x\to 0^-$, note that $x\to 0$ and $e^{1/x}\to e^{-\infty}\to 0$. The the limit evaluates to $0\cdot 0$ which is not an indeterminate form. We conclude then that
	\[\lim_{x\to 0^-}x\cdot e^{1/x} = 0.\]

	\item	This limit initially evaluates to the indeterminate form $\infty-\infty$. By applying a logarithmic rule, we can rewrite the limit as 
\[
\lim_{x\to\infty}\left(\ln(x+1)-\ln x\right) = \lim_{x\to \infty} \ln \left(\frac{x+1}x\right).
\]

As $x\to \infty$, the argument of the natural logarithm approaches $\infty/\infty$, to which we can apply L'Hôpital's Rule.
\[\lim_{x\to\infty} \frac{x+1}x \LHequals \lim_{x\to\infty}\frac11=1.
\]

Since $x\to\infty$ implies $\ds\frac{x+1}x\to 1$, it follows that 
\[x\to\infty \quad \text{ implies }\quad \ln\left(\frac{x+1}x\right)\to\ln 1=0.\]

Thus
\[
 \lim_{x\to\infty} \left(\ln(x+1)-\ln x\right)
 = \lim_{x\to \infty} \ln \left(\frac{x+1}x\right)=0.
\]
Interpretation: since this limit evaluates to 0, it means that for large $x$, there is essentially no difference between $\ln (x+1)$ and $\ln x$; their difference is essentially 0.

	\item	The limit $\ds \lim_{x\to\infty} \left(x^2-e^x\right)$ initially returns the indeterminate form $\infty-\infty$. We can rewrite the expression by factoring out $x^2$; $\ds x^2-e^x = x^2\left(1-\frac{e^x}{x^2}\right).$ We need to evaluate how $e^x/x^2$ behaves as $x\to\infty$:
\[
\lim_{x\to\infty}\frac{e^x}{x^2} \LHequals \lim_{x\to\infty} \frac{e^x}{2x}
\LHequals \lim_{x\to\infty} \frac{e^x}{2} = \infty.
\]

Thus $\lim_{x\to\infty}x^2(1-e^x/x^2)$ evaluates to $\infty\cdot(-\infty)$, which is not an indeterminate form; rather, $\infty\cdot(-\infty)$ evaluates to $-\infty$. We conclude that 
$\ds \lim_{x\to\infty} \left(x^2-e^x\right) = -\infty.$

Interpretation: as $x$ gets large, the difference between $x^2$ and $e^x$ grows very large.
\end{enumerate}
\end{example}

\subsection{Indeterminate Forms\ \ \texorpdfstring{$0^0$, $1^\infty$, and $\infty^0$}{0\^{}0, 1\^{}∞, and ∞\^{}0}}

When faced with a limit that returns one of the indeterminate forms $0^0$, $1^\infty$, or $\infty^0$, it is often useful to use the natural logarithm to convert to an indeterminate form we already know how to find the limit of, then use the natural exponential function find the original limit. This is possible because the natural logarithm and natural exponential functions are inverses and because they are both continuous. The following Key Idea expresses the concept, which is followed by an example that demonstrates its use.

\begin{keyidea}[Evaluating Limits Involving Indeterminate Forms $0^0$, $1^\infty$ and $\infty^0$]\label{idea:LHR_power}%
If $\ds \lim_{x\to c} \ln\bigl(f(x)\bigr) = L$,\quad then 
$\ds \lim_{x\to c} f(x) = \lim_{x\to c} e^{\ln(f(x))} = e\,^L.$ \index{limit!indeterminate form}\index{indeterminate form}
\end{keyidea}

\begin{example}[Using L'Hôpital's Rule with indeterminate forms involving exponents]\label{ex_LHR4}%
Evaluate the following limits.
\[
 \text{1.}\quad\lim_{x\to\infty} \left(1+\frac1x\right)^x \qquad\qquad
 \text{2.}\quad\lim_{x\to0^+} x^x.
\]
\solution
\begin{enumerate}
\item		This is equivalent to a special limit given in \autoref{thm:special_limits}; these limits have important applications in mathematics and finance. Note that the exponent approaches $\infty$ while the base approaches 1, leading to the indeterminate form $1^\infty$. Let $f(x) = (1+1/x)^x$; the problem asks to evaluate $\ds\lim_{x\to\infty}f(x)$. Let's first evaluate $\ds \lim_{x\to\infty}\ln\bigl(f(x)\bigr)$.
\begin{align*}
\lim_{x\to\infty}\ln\bigl(f(x)\bigr)
			&= \lim_{x\to\infty} \ln \left(1+\frac1x\right)^x \\
			&= \lim_{x\to\infty} x\ln\left(1+\frac1x\right)\\
			&= \lim_{x\to\infty} \frac{\ln\left(1+\frac1x\right)}{1/x}\\
			\intertext{This produces the indeterminate form 0/0, so we apply L'Hôpital's Rule.}
			&=	\lim_{x\to\infty} \frac{\frac{1}{1+1/x}\cdot(-1/x^2)}{(-1/x^2)} \\
			&= \lim_{x\to\infty}\frac{1}{1+1/x}\\
			&= 1.
\end{align*}
Thus $\ds\lim_{x\to\infty} \ln \bigl(f(x)\bigr) = 1.$ We return to the original limit and apply \autoref{idea:LHR_power}.
\[\lim_{x\to\infty}\left(1+\frac1x\right)^x = \lim_{x\to\infty} f(x) =  \lim_{x\to\infty}e^{\ln (f(x))} = e^1 = e.\]
This is another way to determine the value of the number $e$.

\item		This limit leads to the indeterminate form $0^0$. Let $f(x) = x^x$ and consider first $\ds\lim_{x\to0^+} \ln\bigl(f(x)\bigr)$. 
%
\mtable{A graph of $f(x)=x^x$ supporting the fact that as $x\to 0^+$, $f(x)\to 1$.}{fig:LHR4}{\begin{tikzpicture}[alt={A graph starting near (0,1), moving downward as it goes right, and then moving up to finish near (2,4).}]
\begin{axis}[width=1.16\marginparwidth,tick label style={font=\scriptsize},
axis y line=middle,axis x line=middle,name=myplot,axis on top,ytick={1,2,3,4},
ymin=-.4,ymax=4.5,xmin=-.1,xmax=2.2,scaled ticks=false]
\addplot[draw={\colorone},thick,smooth,domain=.01:2]{exp(x*ln(x))};
\draw (axis cs:1,1) node [below right] {\scriptsize $f(x)=x^x$};
\end{axis}
\node [right] at (myplot.right of origin) {\scriptsize $x$};
\node [above] at (myplot.above origin) {\scriptsize $y$};
\end{tikzpicture}}%
%
\begin{align*}
\lim_{x\to0^+} \ln\bigl(f(x)\bigr) &= \lim_{x\to0^+} \ln\left(x^x\right) \\
			&= \lim_{x\to0^+} x\ln x \\
			\intertext{This produces the indeterminate form \intertextmath{0(-\infty)}, so we rewrite it in order to apply L'Hôpital's Rule.}
			&= \lim_{x\to0^+} \frac{\ln x}{1/x}.\\
			\intertext{This produces the indeterminate form \intertextmath{-\infty/\infty} so we apply L'Hôpital's Rule.}
			&=	\lim_{x\to0^+} \frac{1/x}{-1/x^2} \\
			&= \lim_{x\to0^+} -x \\
			&= 0.
\end{align*}%
Thus $\ds\lim_{x\to0^+} \ln\bigl(f(x)\bigr) =0$. We return to the original limit and apply \autoref{idea:LHR_power}.
\[
\lim_{x\to0^+} x^x = \lim_{x\to0^+} f(x) = \lim_{x\to0^+} e^{\ln(f(x))} = e^0 = 1.
\]
This result is supported by the graph of $f(x)=x^x$ given in \autoref{fig:LHR4}.
\end{enumerate}
\end{example}

% todo Tim do we want to add a hierarchy of function growth to the end of LH section?

%Our brief revisit of limits will be rewarded in the next section where we consider \emph{improper integration.} So far, we have only considered definite integrals where the bounds are finite numbers, such as $\ds \int_0^1 f(x)\dd x$. Improper integration considers integrals where one, or both, of the bounds are ``infinity.'' Such integrals have many uses and applications, in addition to generating ideas that are enlightening.

\printexercises{exercises/06-06-exercises}

}


\apexchapter{Techniques of Integration}{chapter:anti_tech}

\autoref{chapter:integration} introduced the antiderivative and connected it to signed areas under a curve through the Fundamental Theorem of Calculus. The chapter after explored more applications of definite integrals than just area. As evaluating definite integrals will become even more important, we will want to find antiderivatives of a variety of functions.

This chapter is devoted to exploring techniques of antidifferentiation. While not every function has an antiderivative in terms of elementary functions,
% (a concept introduced in the section on Numerical Integration),
we can still find antiderivatives of a wide variety of functions.
\input{text/06_Int_By_Parts}

% this section gets 06_01_ex_27 ?

\section{Trigonometric Integrals}\label{sec:trigint}

Trigonometric functions are useful for describing periodic behavior. This section describes several techniques for finding antiderivatives of certain combinations of trigonometric functions.

\subsection{Integrals of the form \texorpdfstring{$\ds \int \sin^m x\cos^n x\dd x$}{∫(sin x)\^{}m (cos x)\^{}n dx}}

In learning the technique of Substitution, we saw the integral $\int \sin x\cos x\dd x$ in \autoref{ex_sub10}. The integration was not difficult, and one could easily evaluate the indefinite integral by letting $u=\sin x$ or by letting $u = \cos x$. This integral is easy since the power of both sine and cosine is 1.

We generalize this and consider integrals of the form $\int \sin^mx\cos^nx\dd x$, where $m,n$ are nonnegative integers. Our strategy for evaluating these integrals is to use the identity $\cos^2x+\sin^2x=1$ to convert high powers of one trigonometric function into the other, leaving a single sine or cosine term in the integrand. We summarize the general technique in the following Key Idea.

\youtubeVideo{zyg9k1je7Fg}{Trigonometric Integrals --- Part 2 of 6}


{%
\tcbset{grow to right by=4em} % 4 sufficient
\begin{keyidea}[Integrals Involving Powers of Sine and Cosine]\label{idea:trig_int_1}%
Consider $\ds \int \sin^mx\cos^nx\dd x$, where $m,n$ are nonnegative integers.\index{integration!of trig.\ powers}
\begin{enumerate}
	\item	If $m$ is odd, then $m=2k+1$ for some integer $k$. Rewrite \small
		\[
		\sin^mx = \sin^{2k+1}x = \sin^{2k}x\sin x = (\sin^2x)^k\sin x = (1-\cos^2x)^k\sin x.
		\]
		\normalsize Then \small
		\[
		\int \sin^mx\cos^nx\dd x = \int (1-\cos^2x)^k\sin x\cos^nx\dd x = -\int (1-u^2)^ku^n\dd u,
		\]
		\normalsize where $u = \cos x$ and $\dd u = -\sin x\dd x$. 
	\item	If $n$ is odd, then using substitutions similar to that outlined above we have \small
		\[\int \sin^mx\cos^nx\dd x = \int u^m(1-u^2)^k\dd u,\]
		\normalsize where $u = \sin x$ and $\dd u = \cos x\dd x$.
	\item	If both $m$ and $n$ are even, use the half-angle identities \small
		\[
		\cos^2x = \frac{1+\cos (2x)}{2} \quad \text{and}\quad \sin^2x = \frac{1-\cos(2x)}2
		\]
		\normalsize to reduce the degree of the integrand. Expand the result and apply the principles of this Key Idea again.
	\end{enumerate}
\end{keyidea}%
}

We practice applying \autoref{idea:trig_int_1} in the next examples.

\begin{example}[Integrating powers of sine and cosine]\label{ex_trigint1}%
Evaluate $\ds\int\sin^5x\cos^8x\dd x$.
\solution
The power of the sine factor is odd, so we rewrite $\sin^5x$ as
\[\sin^5x = \sin^4x\sin x = (\sin^2x)^2\sin x = (1-\cos^2x)^2\sin x.\]

Our integral is now $\ds \int (1-\cos^2x)^2\cos^8x\sin x\dd x$. Let $u = \cos x$, hence $\dd u = -\sin x\dd x$. Making the substitution and expanding the integrand gives
\begin{align*}
 \int (1-\cos^2x)^2\cos^8x\sin x\dd x
 &= -\int (1-u^2)^2u^8\dd u \\
 &= -\int \bigl(1-2u^2+u^4\bigr)u^8\dd u \\
 &= -\int \bigl(u^8-2u^{10}+u^{12}\bigr)\dd u \\
 &= -\frac19u^9 + \frac2{11}u^{11} - \frac1{13}u^{13} + C \\
 &=-\frac19\cos^9 x + \frac2{11}\cos^{11} x - \frac1{13}\cos^{13} x + C.
\end{align*}
\end{example}

\begin{example}[Integrating powers of sine and cosine]\label{ex_trigint2}%
Evaluate $\ds \int\sin^5x\cos^9x\dd x$.
\solution
Because the powers of both the sine and cosine factors are odd, we can apply the techniques of \autoref{idea:trig_int_1} to either power.
We choose to work with the power of the sine factor since that has a smaller exponent.
% We choose to work with the power of the cosine factor since the previous example used the sine factor's power.

We rewrite $\sin^5x$ as
\begin{align*}
 \sin^5x&=\sin^4x\sin x\\
 &=(1-\cos^2x)^2\sin x\\
 &=(1-2\cos^2x+\cos^4x)\sin x.
\end{align*}
This lets us rewrite the integral as
\[
\int\sin^5x\cos^9x\dd x=\int\bigl(1-2\cos^2x+\cos^4x\bigr)\sin x\cos^9x\dd x.
\]

Substituting and integrating with $u=\cos x$ and $\dd u=-\sin x\dd x$, we have
\begin{align*}
\int\bigl(1-2\cos^2x+\cos^4x\bigr)&\sin x\cos^9x\dd x\\
&=-\int\bigl(1-2u^2+u^4\bigr)u^9\dd u\\
&=-\int u^9-2u^{11}+u^{13}\dd u\\
&=-\frac1{10}u^{10}+\frac16u^{12}-\frac1{14}u^{14}+C\\
&=-\frac1{10}\cos^{10}x+\frac16\cos^{12}x-\frac1{14}\cos^{14}x+C.
\end{align*}

Instead, another approach would be to rewrite $\cos^9x$ as
\begin{align*} \cos^9 x &= \cos^8x\cos x \\
				&= (\cos^2x)^4\cos x \\
				&= (1-\sin^2x)^4\cos x \\
				&= (1-4\sin^2x+6\sin^4x-4\sin^6x+\sin^8x)\cos x.
\end{align*}

We rewrite the integral as 
\[\int\sin^5x\cos^9x\dd x = \int\bigl(\sin^5x\bigr)\bigl(1-4\sin^2x+6\sin^4x-4\sin^6x+\sin^8x\bigr)\cos x\dd x.\]

Now substitute and integrate, using $u = \sin x $ and $\dd u = \cos x\dd x$.
\begin{align*}
 \int & \bigl(\sin^5x\bigr)\bigl(1-4\sin^2x+6\sin^4x-4\sin^6x+\sin^8x\bigr)\cos x\dd x \\
 &=\int u^5(1-4u^2+6u^4-4u^6+u^8)\dd u \\
 &= \int\bigl(u^5-4u^7+6u^9-4u^{11}+u^{13}\bigr)\dd u \\
 &= \frac16u^6-\frac12u^8+\frac35u^{10}-\frac13u^{12}+\frac{1}{14}u^{14}+C\\
 &= \frac16\sin^6 x-\frac12\sin^8 x+\frac35\sin^{10} x-\frac13\sin^{12} x+\frac{1}{14}\sin^{14} x+C.
\end{align*}
%
\end{example}

\textbf{Technology Note:} The work we are doing here can be a bit tedious, but the skills developed (problem solving, algebraic manipulation, etc.) are important. Nowadays problems of this sort are often solved using a computer algebra system. The powerful program \textit{Mathematica}\textsuperscript{\textregistered} integrates $\int \sin^5x\cos^9x\dd x$ as
{\small
\begin{multline*}
 f(x)=\\
 -\frac{45 \cos (2 x)}{16384}-\frac{5 \cos (4 x)}{8192}+\frac{19 \cos (6
   x)}{49152}+\frac{\cos (8 x)}{4096}-\frac{\cos (10 x)}{81920}-\frac{\cos (12
   x)}{24576}-\frac{\cos (14 x)}{114688},
\end{multline*}}

\mtable{A plot of $f(x)$ and $g(x)$ from \autoref{ex_trigint2} and the Technology Note.}{fig:trigint2}{\begin{tikzpicture}[alt={g(x) starts flat near the origin, then moves up to near (1,0.0045), is flat until near (2,0,0045), then moves down to finish flat near (π,0).  f(x) is similar, but slightly below g(x).}]
\begin{axis}[width=\marginparwidth,tick label style={font=\scriptsize},
axis y line=middle,axis x line=middle,name=myplot,axis on top,
ytick={-.002,.002,.004},yticklabels={$-0.002$,$0.002$,$0.004$},
ymin=-.003,ymax=0.005,xmin=-.1,xmax=3.15,scaled ticks=false]
% should we actually do this instead of coordinates?
\addplot [draw={\colortwo},thick,smooth] coordinates {(0,-0.0027879) (0.15708,-0.0027856) (0.31416,-0.0026798) (0.47124,-0.0020322) (0.62832,-0.00060997) (0.7854,0.00089518)(0.94248,0.0017233) (1.0996,0.0019485) (1.2566,0.0019734) (1.4137,0.0019741) (1.5708,0.0019741) (1.7279,0.0019741) (1.885,0.0019734) (2.042,0.0019485) (2.1991,0.0017233) (2.3562,0.00089518) (2.5133,-0.00060997) (2.6704,-0.0020322) (2.8274,-0.0026798) (2.9845,-0.0027856) (3.1416,-0.0027879)};
\draw (axis cs:2.6,0.004) node {\scriptsize $g(x)$};
\addplot [draw={\colorone},thick,smooth] coordinates {(0,0) (0.15708,0) (0.31416,0.00010807) (0.47124,0.0007557) (0.62832,0.0021779) (0.7854,0.003683) (0.94248,0.0045111)
(1.0996,0.0047364) (1.2566,0.0047612) (1.4137,0.0047619)
(1.5708,0.0047619) (1.7279,0.0047619) (1.885,0.0047612) (2.042,0.0047364) (2.1991,0.0045111) (2.3562,0.003683) (2.5133,0.0021779) (2.6704,0.0007557) (2.8274,0.00010807) (2.9845,0) (3.1416,0)};
\draw (axis cs:2.4,-0.002) node {\scriptsize $f(x)$};
\end{axis}
\node [right] at (myplot.right of origin) {\scriptsize $x$};
\node [above] at (myplot.above origin) {\scriptsize $y$};
\end{tikzpicture}}
\noindent
which clearly has a different form than our second answer in \autoref{ex_trigint2}, which is
\[
 g(x)=\frac16\sin^6 x-\frac12\sin^8 x+\frac35\sin^{10} x-\frac13\sin^{12} x+\frac{1}{14}\sin^{14} x.
\]
\autoref{fig:trigint2} shows a graph of $f$ and $g$; they are clearly not equal, but they differ \emph{only by a constant}: $g(x) = f(x) + C$ for some constant $C$. We have two different antiderivatives of the same function, meaning both answers are correct.\bigskip

%This is a common trigonometric integral.\bigskip

\begin{example}[Integrating powers of sine and cosine]\label{ex_sub8}%
Evaluate $\ds \int \sin^2x\dd x$.
\solution
The power of sine is even so we employ a half-angle identity, algebra and a u- substitution as follows:
\begin{align*}
	\int \sin^2x\dd x
	&= \int \frac{1-\cos(2x)}2\dd x \\
	&= \frac12\int 1-\cos(2x)\dd x \\
	&= \frac12\left(x - \frac12\sin(2x)\right)+C \\
	&= \frac12x - \frac14\sin(2x) + C.
\end{align*}
\end{example}

\begin{example}[Integrating powers of sine and cosine]\label{ex_trigint3}%
Evaluate $\ds\int\cos^4x\sin^2x\dd x$.
\solution
The powers of sine and cosine are both even, so we employ the half-angle formulas and algebra as follows.
\begin{align*}
	\int \cos^4x\sin^2x\dd x
	&= \int\left(\frac{1+\cos(2x)}{2}\right)^2\left(\frac{1-\cos(2x)}2\right)\dd x \\
	&= \int\frac{1+2\cos(2x)+\cos^2(2x)}4\cdot\frac{1-\cos(2x)}2\dd x\\
	&=	\int \frac18\bigl(1+\cos(2x)-\cos^2(2x)-\cos^3(2x)\bigr)\dd x
\end{align*}
The $\cos(2x)$ term is easy to integrate.
%, especially with \autoref{idea:linearsub}.
The $\cos^2(2x)$ term is another trigonometric integral with an even power, requiring the half-angle formula again. The $\cos^3(2x)$ term is a cosine function with an odd power, requiring a substitution as done before. We integrate each in turn below.

\begin{gather*}
\int\cos(2x)\dd x = \frac12\sin(2x)+C.\\
\int\cos^2(2x)\dd x = \int \frac{1+\cos(4x)}2\dd x
= \frac12\bigl(x+\frac14\sin(4x)\bigr)+C.
\end{gather*}

Finally, we rewrite $\cos^3(2x)$ as
\[\cos^3(2x) = \cos^2(2x)\cos(2x) = \bigl(1-\sin^2(2x)\bigr)\cos(2x).\]
Letting $u=\sin(2x)$, we have $\dd u = 2\cos(2x)\dd x$, hence
\begin{align*}
\int \cos^3(2x)\dd x &= \int\bigl(1-\sin^2(2x)\bigr)\cos(2x)\dd x\\
							&= \int \frac12(1-u^2)\dd u\\
							&= \frac12\Bigl(u-\frac13u^3\Bigr)+C\\
							&=	\frac12\Bigl(\sin(2x)-\frac13\sin^3(2x)\Bigr)+C
\end{align*}

Putting all the pieces together, we have
\begin{align*}
	\int &\cos^4x\sin^2x\dd x \\
	&=\int \frac18\bigl(1+\cos(2x)-\cos^2(2x)-\cos^3(2x)\bigr)\dd x \\
	&= \frac18
	\Bigl[x+\frac12\sin(2x)-\frac12\bigl(x+\frac14\sin(4x)\bigr)
	-\frac12\Bigl(\sin(2x)-\frac13\sin^3(2x)\Bigr)\Bigr]
	+C \\
	&=\frac18\Bigl[\frac12x-\frac18\sin(4x)+\frac16\sin^3(2x)\Bigr]+C.
\end{align*}
\end{example}

The process above was a bit long and tedious, but being able to work a problem such as this from start to finish is important.

\subsection{Integrals of the form \texorpdfstring{$\ds\int\tan^mx\sec^nx\dd x$}{∫(tan x)\^{}m (sec x)\^{}n dx}}

When evaluating integrals of the form $\int \sin^mx\cos^nx\dd x$, the Pythagorean Theorem allowed us to convert even powers of sine into even powers of cosine, and vice versa. If, for instance, the power of sine was odd, we pulled out one $\sin x$ and converted the remaining even power of $\sin x$ into a function using powers of $\cos x$, leading to an easy substitution.

The same basic strategy applies to integrals of the form $\int \tan^mx\sec^n x\dd x$, albeit a bit more nuanced. The following three facts will prove useful:
\begin{itemize}
\item $\frac{\dd}{\dd x}(\tan x) = \sec^2x$, 
\item $\frac{\dd}{\dd x}(\sec x) = \sec x\tan x$ , and 
\item	$1+\tan^2x = \sec^2x$ (the Pythagorean Theorem).
\end{itemize}

If the integrand can be manipulated to separate a $\sec^2x$ term with the remaining secant power even, or if a $\sec x\tan x$ term can be separated with the remaining $\tan x$ power even, the Pythagorean Theorem can be employed, leading to a simple substitution. This strategy is outlined in the following Key Idea.

{
\tcbset{grow to right by=13em}
\begin{keyidea}[Integrals Involving Powers of Tangent and Secant]\label{idea:trig_int_2}%
Consider $\ds\int\tan^mx\sec^nx\dd x$, where $m$ and $n$ are nonnegative integers.\index{integration!of trig. powers}
\begin{enumerate}
\item		If $n$ is even, then $n=2k$ for some integer $k$. Rewrite $\sec^nx$ as 
\[\sec^nx = \sec^{2k}x = \sec^{2k-2}x\sec^2x = (1+\tan^2x)^{k-1}\sec^2x.\]
Then
\[
\int\tan^mx\sec^nx\dd x=\int\tan^mx(1+\tan^2x)^{k-1}\sec^2x\dd x = \int u^m(1+u^2)^{k-1}\dd u,
\]
where $u = \tan x$ and $\dd u = \sec^2x\dd x$.

\item		If $m$ is odd and $n>0$, then $m=2k+1$ for some integer $k$. Rewrite $\tan^mx\sec^nx$ as
\[
\tan^mx\sec^nx = \tan^{2k+1}x\sec^nx = \tan^{2k}x\sec^{n-1}x\sec x\tan x = (\sec^2x-1)^k\sec^{n-1}x\sec x\tan x.
\]
Then
\[
\int\tan^mx\sec^nx\dd x=\int(\sec^2x-1)^k\sec^{n-1}x\sec x\tan x\dd x = \int(u^2-1)^ku^{n-1}\dd u,
\]
where $u = \sec x$ and $\dd u = \sec x\tan x\dd x$.

\item If $n$ is odd and $m$ is even, then $m=2k$ for some integer $k$. Convert $\tan^mx $ to $(\sec^2x-1)^k$. Expand the new integrand and use Integration By Parts, with $\dd v = \sec^2x\dd x$.

\item		If $m$ is even and $n=0$, rewrite $\tan^mx$ as
\[
\tan^mx = \tan^{m-2}x\tan^2x = \tan^{m-2}x(\sec^2x-1) = \tan^{m-2}\sec^2x-\tan^{m-2}x.
\]
So
\[
\int\tan^mx\dd x = \underbrace{\int\tan^{m-2}x\sec^2x\dd x}_{\text{\small apply rule \#1}}\quad - \underbrace{\int\tan^{m-2}x\dd x}_{\text{\small apply rule \#4 again}}.
\]

\end{enumerate}
\end{keyidea}
}

The techniques described in items 1 and 2 of \autoref{idea:trig_int_2} are relatively straightforward, but the techniques in items 3 and 4 can be rather tedious. A few examples will help with these methods.

\begin{example}[Integrating powers of tangent and secant]\label{ex_trigint5}%
Evaluate $\ds\int \tan^2x\sec^6x\dd x$.
\solution
Since the power of secant is even, we use rule \#1 from \autoref{idea:trig_int_2} and pull out a $\sec^2x$ in the integrand. We convert the remaining powers of secant into powers of tangent.
\begin{align*}
\int \tan^2x\sec^6x\dd x &= \int\tan^2x\sec^4x\sec^2x\dd x \\
		&= \int \tan^2x\bigl(1+\tan^2x\bigr)^2\sec^2x\dd x \\
\intertext{Now substitute, with $u=\tan x$, with $\dd u = \sec^2x\dd x$.}
		&=\int u^2\bigl(1+u^2\bigr)^2\dd u\\
\intertext{We leave the integration and subsequent substitution to the reader. The final answer is}
		&=\frac13\tan^3x+\frac25\tan^5x+\frac17\tan^7x+C.
\end{align*}
\end{example}

We derived integrals for tangent and secant in \autoref{sec:substitution} and will regularly use them when evaluating integrals of the form $\tan^m x \sec^n x \dd x$.  As a reminder:
\begin{align*}
 \int\tan x\dd x &=\ln\abs{\sec x}+C \\
 \int\sec x\dd x &=\ln\abs{\sec x+\tan x}+C
\end{align*}

\begin{example}[Integrating powers of tangent and secant]\label{ex_trigint6}%
Evaluate $\ds\int \sec^3x\dd x$.
\solution
We apply rule \#3 from \autoref{idea:trig_int_2} as the power of secant is odd and the power of tangent is even (0 is an even number). We use Integration by Parts; the rule suggests letting $\dd v = \sec^2x\dd x$, meaning that $u = \sec x$. \\
\noindent\begin{minipage}[t]{\linewidth}\noindent%
\captionsetup{type=figure}%
\[
\begin{aligned}
u&= \sec x & \dd v&=\sec^2 x\dd x\\
\dd u&= \text{?} & v&=\text{?}
\end{aligned}
\qquad\Rightarrow\qquad
\begin{aligned}
u&= \sec x & \dd v&=\sec^2 x\dd x\\
\dd u&= \sec x\tan x\dd x & v&=\tan x
\end{aligned}
\]
\caption{Setting up Integration by Parts.}\label{fig:trigint1}
\end{minipage}

Employing Integration by Parts, we have
\begin{align*}
\int \sec^3x\dd x
 	&=	\int \underbrace{\sec x}_u\cdot\underbrace{\sec^2 x\dd x}_{\dd v}\\
	&=	\sec x\tan x - \int \sec x\tan^2x\dd x. \\
\intertext{This new integral also requires applying rule \#3 of \autoref{idea:trig_int_2}:}
	&= \sec x\tan x - \int \sec x \bigl(\sec^2 x-1\bigr)\dd x\\
	&=	\sec x\tan x - \int \sec^3x\dd x + \int \sec x\dd x \\
	&= \sec x\tan x -\int \sec^3x\dd x + \ln\abs{\sec x+\tan x}
\end{align*}
%
\mnote{\textbf{Note:} Remember that in \autoref{ex_sub7}, we found that $\int\sec x\dd x=\ln\abs{\sec x+\tan x}+C$}
%
In previous applications of Integration by Parts, we have seen where the original integral has reappeared in our work. We resolve this by adding $\int \sec^3x\dd x$ to both sides, giving:
\begin{align*}
2\int \sec^3x\dd x &= \sec x\tan x + \ln\abs{\sec x+\tan x} \\
\int \sec^3x\dd x &= \frac12\Bigl(\sec x\tan x + \ln\abs{\sec x+\tan x}\Bigr)+C.
\end{align*}
\end{example}

We give one more example.

\begin{example}[Integrating powers of tangent and secant]\label{ex_trigint7}%
Evaluate $\ds\int\tan^6x\dd x$.
\solution
We employ rule \#4 of \autoref{idea:trig_int_2}. 
\begin{align*}
	\int \tan^6x\dd x
	&= \int \tan^4x\tan^2x\dd x \\
	&= \int\tan^4x\bigl(\sec^2x-1\bigr)\dd x\\
	&= \int\tan^4x\sec^2x\dd x - \int\tan^4x\dd x \\
\intertext{We integrate the first integral with substitution, $u=\tan x$ and $\dd u=\sec^2x\dd x$; and the second by employing rule \#4 again.}
	&= \int u^4\dd u-\int\tan^2 x\tan^2 x\dd x \\
	&=	\frac15\tan^5x-\int\tan^2x\bigl(\sec^2x-1\bigr)\dd x \\
	&= \frac15\tan^5x -\int\tan^2x\sec^2x\dd x + \int\tan^2x\dd x\\
\intertext{Again, use substitution for the first integral and rule \#4 for the second.}
	&= \frac15\tan^5x-\frac13\tan^3x+\int\bigl(\sec^2x-1\bigr)\dd x \\
	&=	 \frac15\tan^5x-\frac13\tan^3x+\tan x - x+C.
\end{align*}
\end{example}

\subsection{Integrals of the form \texorpdfstring{$\ds\int\cot^mx\csc^nx\dd x$}{∫(cot x)\^{}m (csc x)\^{}n dx}}

Not surprisingly, evaluating integrals of the form $\int\cot^mx\csc^nx\dd x$ is similar to evaluating $\int\tan^mx\sec^nx\dd x$. The guidelines from \autoref{idea:trig_int_2} and the following three facts will be useful:
\begin{align*}
 \frac{\dd}{\dd x}(\cot x) &= -\csc^2x \\
 \frac{\dd}{\dd x}(\csc x) &= -\csc x\cot x,\qquad\text{and} \\
 \csc^2 x &= \cot^2x+1
\end{align*}

\begin{example}[Integrating powers of cotangent and cosecant]\label{ex_int_cot_csc}%
Evaluate $\ds\int\cot^2x\csc^4x\dd x$
\solution
Since the power of cosecant is even we will let $u=\cot x$ and save a $\csc^2x$ for the resulting $\dd u=-\csc^2x\dd x$.
\begin{align*}
 \int\cot^2x\csc^4x\dd x
 &=\int\cot^2x\csc^2x\csc^2x\dd x \\
 &=\int\cot^2x(1+\cot^2x)\csc^2x\dd x \\
 &=-\int u^2(1+u^2)\dd u.
\end{align*}
The integration and substitution required to finish this example are similar to that of previous examples in this section. The result is
\[-\frac13\cot^3x-\frac15\cot^5x+C.\]
\end{example}

\subsection{Integrals of the form \texorpdfstring{$\ds \int\sin(mx)\sin(nx)\dd x,$ $\ds\int \cos(mx)\cos(nx)\dd x$, and $\ds\int \sin(mx)\cos(nx)\dd x$.}{∫sin(mx)sin(nx)dx, ∫cos(mx)cos(nx)dx, and ∫sin(mx)cos(nx)dx}}

Functions that contain products of sines and cosines of differing periods are important in many applications including the analysis of sound waves. Integrals of the form 
\[
\int\sin(mx)\sin(nx)\dd x,\quad \int \cos(mx)\cos(nx)\dd x \quad \text{and}\quad\int \sin(mx)\cos(nx)\dd x
\]
are best approached by first applying the Product to Sum Formulas of Trigonometry found in the back cover of this text, namely
\begin{align*}
\sin(mx)\sin(nx) &= \frac12\Bigl[\cos\bigl((m-n)x\bigr)-\cos\bigl((m+n)x\bigr)\Bigr] \\
\cos(mx)\cos(nx) &= \frac12\Bigl[\cos\bigl((m-n)x\bigr)+\cos\bigl((m+n)x\bigr)\Bigr] \\
\sin(mx)\cos(nx) &=	\frac12\Bigl[\sin\bigl((m-n)x\bigl)+\sin\bigl((m+n)x\bigr)\Bigr]
\end{align*}

\begin{example}[Integrating products of $\sin(mx)$ and $\cos(nx)$]\label{ex_trigint4}%
Evaluate $\ds\int\sin(5x)\cos(2x)\dd x$.
\solution
The application of the formula and subsequent integration are straightforward:
\begin{align*}
	\int\sin(5x)\cos(2x)\dd x
	&= \int \frac12\Bigl[\sin(3x)+\sin(7x)\Bigr]\dd x \\
	&= -\frac16\cos(3x) - \frac1{14}\cos(7x) + C.
\end{align*}
\end{example}

\subsection{Integrating other combinations of trigonometric functions}

Combinations of trigonometric functions that we have not discussed in this chapter are evaluated by applying algebra, trigonometric identities and other integration strategies to create an equivalent integrand that we can evaluate. To evaluate ``crazy'' combinations, those not readily manipulated into a familiar form, one should use integral tables. A table of ``common crazy'' combinations can be found at the end of this text.

These latter examples were admittedly long, with repeated applications of the same rule. Try to not be overwhelmed by the length of the problem, but rather admire how robust this solution method is. A trigonometric function of a high power can be systematically reduced to trigonometric functions of lower powers until all antiderivatives can be computed. 

The next section introduces an integration technique known as Trigonometric Substitution, a clever combination of Substitution and the Pythagorean Theorem.

\printexercises{exercises/06-03-exercises}

% the following was taken from u-substitution

%\subsection{Substitution and Inverse Trigonometric Functions}
%
%%In \autoref{sec:deriv_inverse_function} 
%When studying derivatives of inverse functions, we learned that
%\[\frac{\dd}{\dd x}\bigl(\tan^{-1}x\bigr) = \frac{1}{1+x^2}.\]
%Applying the Chain Rule to this is not difficult; for instance,
%\[\frac{\dd}{\dd x}\bigl(\tan^{-1}5x\bigr) = \frac{5}{1+25x^2}.\]
%We now explore how Substitution can be used to ``undo'' certain derivatives that are the result of the Chain Rule applied to Inverse Trigonometric functions. We begin with an example.
%
%\begin{example}[Integrating by substitution: inverse trigonometric functions]\label{ex_subst14}%
%Evaluate $\ds \int \frac{1}{25+x^2}\dd x$.
%\solution
%The integrand looks similar to the derivative of the arctangent function. Note:
%\begin{align*}
%\frac{1}{25+x^2} &= \frac{1}{25(1+\frac{x^2}{25})}\\
%							&= \frac{1}{25(1+\left(\frac{x}{5}\right)^2)} \\
%							&= \frac{1}{25}\frac{1}{1+\left(\frac{x}{5}\right)^2}\ .
%\end{align*}
%Thus
%\[\int\frac{1}{25+x^2}\dd x = \frac{1}{25}\int \frac{1}{1+\left(\frac{x}{5}\right)^2}\dd x.\]
%This can be integrated using Substitution. Set $u = x/5$, hence $du = \dd x/5$ or $\dd x=5du$. Thus
%\begin{align*}
%	\int\frac{1}{25+x^2}\dd x
%	&= \frac{1}{25}\int \frac{1}{1+\left(\frac{x}{5}\right)^2}\dd x \\
%	&= \frac15\int \frac{1}{1+u^2}\dd u \\
%	&= \frac15\tan^{-1}u + C \\
%	&= \frac15\tan^{-1}\left(\frac x5\right)+C
%\end{align*}
%\end{example}
%
%\autoref{ex_subst14} demonstrates a general technique that can be applied to other integrands that result in inverse trigonometric functions. The results are summarized here.
%
%% moved to 7.2
%
%Let's practice using \autoref{thm:int_inverse_trig}.\bigskip
%
%\begin{example}[Integrating by substitution: inverse trigonometric functions]\label{ex_subst15}%
%Evaluate the given indefinite integrals.
%\[\int \frac{1}{9+x^2}\dd x,\quad \int \frac{1}{x\sqrt{x^2-\frac{1}{100}}}\dd x\quad \text{ and }\quad  \int \frac{1}{\sqrt{5-x^2}}\dd x.\]
%\solution
%Each can be answered using a straightforward application of \autoref{thm:int_inverse_trig}.
%\begin{gather*}
%\int \frac{1}{9+x^2}\dd x = \frac13\tan^{-1} \frac x3 + C,\text{ as }a=3.
%\int\frac1{x\sqrt{x^2-\frac1{100}}}\dd x=10\sec^{-1}10x+C,\text{ as }a=\frac1{10}.
%\int \frac{1}{\sqrt{5-x^2}} = \sin^{-1}\frac{x}{\sqrt{5}}+C,\text{ as }a = \sqrt{5}.
%\end{example}
%
%Most applications of \autoref{thm:int_inverse_trig} are not as straightforward. The next examples show some common integrals that can still be approached with this theorem.
%
%\begin{example}[Integrating by substitution: completing the square]\label{ex_subst16}%
%Evaluate $\ds \int\frac{1}{x^2-4x+13}\dd x$.
%\solution
%Initially, this integral seems to have nothing in common with the integrals in \autoref{thm:int_inverse_trig}. As it lacks a square root, it almost certainly is not related to arcsine or arcsecant. It is, however, related to the arctangent function.
%
%We see this by \emph{completing the square} in the denominator. We give a brief reminder of the process here. 
%
%Start with a quadratic with a leading coefficient of 1. It will have the form of $x^2 + bx + c$. Take 1/2 of $b$, square it, and add/subtract it back into the expression. I.e., 
%\begin{align*}
%	x^2+bx+ c
%	&= \underbrace{x^2 + bx + \frac{b^2}4}_{(x+b/2)^2} - \frac{b^2}4 + c\\
%	&= \left(x+\frac b2\right)^2 + c-\frac{b^2}4
%\end{align*}
%In our example, we take half of $-4$ and square it, getting $4$. We add/subtract it into the denominator as follows:
%
%\begin{align*}
%	\frac{1}{x^2-4x+13}
%	&= \frac{1}{\underbrace{x^2-4x+4}_{(x-2)^2}-4+13}\\
%	&=\frac{1}{(x-2)^2 + 9}
%\end{align*}
%We can now integrate this using the arctangent rule. Technically, we need to substitute first with $u=x-2$, but by now we can do this easily. Thus we have 
%\[
%\int \frac{1}{x^2-4x+13}\dd x
%= \int \frac{1}{(x-2)^2+9}\dd x
%= \frac13\tan^{-1}\frac{x-2}{3}+C.
%\]
%\end{example}
%
%\begin{example}[Integrals requiring multiple methods]\label{ex_subst17}%
%Evaluate $\ds \int \frac{4-x}{\sqrt{16-x^2}}\dd x$.
%\solution
%This integral requires two different methods to evaluate it. We get to those methods by splitting up the integral: 
%\[
%\int \frac{4-x}{\sqrt{16-x^2}}\dd x
%= \int \frac{4}{\sqrt{16-x^2}}\dd x - \int \frac{x}{\sqrt{16-x^2}}\dd x.
%\]
%The first integral is handled using a straightforward application of \autoref{thm:int_inverse_trig}; the second integral is handled by substitution, with $u = 16-x^2$. We handle each separately.
%\[\int \frac{4}{\sqrt{16-x^2}}\dd x = 4\sin^{-1}\frac{x}{4} + C.\]
%For $\ds \int\frac{x}{\sqrt{16-x^2}}\dd x$, we set $u = 16-x^2$, so $du = -2x\dd x$ and $x\dd x = -du/2$. We have
%\begin{align*}
%	\int\frac{x}{\sqrt{16-x^2}}\dd x
%	&= \int\frac{-du/2}{\sqrt{u}}\\
%	&= -\frac12\int \frac{1}{\sqrt{u}}\dd u \\
%	&= - \sqrt{u} + C\\
%	&= -\sqrt{16-x^2} + C.
%\end{align*}
%Combining these together, we have 
%\[\int \frac{4-x}{\sqrt{16-x^2}}\dd x = 4\sin^{-1}\frac x4 + \sqrt{16-x^2}+C.\]
%\end{example}


\section{Trigonometric Substitution}\label{sec:trig_sub}

In \autoref{sec:def_int} we defined the definite integral as the ``signed area under the curve.'' In that section we had not yet learned the Fundamental Theorem of Calculus, so we evaluated special definite integrals which described nice, geometric shapes. For instance, we were able to evaluate
\begin{equation}
\int_{-3}^3\sqrt{9-x^2}\dd x = \frac{9\pi}{2}\label{eq:trigsub1}%
\end{equation}
 as we recognized that $f(x) = \sqrt{9-x^2}$ described the upper half of a circle with radius 3. 

We have since learned a number of integration techniques, including Substitution and Integration by Parts, yet we are still unable to evaluate the above integral without resorting to a geometric interpretation. This section introduces Trigonometric Substitution, a method of integration that fills this gap in our integration skill. This technique works on the same principle as Substitution as found in \autoref{sec:substitution}, though it can feel ``backward.'' In \autoref{sec:substitution}, we set $u=f(x)$, for some function $f$, and replaced $f(x)$ with $u$. In this section, we will set $x=f(\theta)$, where $f$ is a trigonometric function, then replace $x$ with $f(\theta)$. 

\youtubeVideo{yW6Odu0YHL0}{Trigonometric Substitution --- Example 3 / Part 1}

We start by demonstrating this method in evaluating the integral in \autoeqref{eq:trigsub1}. After the example, we will generalize the method and give more examples.

\begin{example}[Using Trigonometric Substitution]\label{ex_trigsub1}%
Evaluate $\ds \int_{-3}^3\sqrt{9-x^2}\dd x$.
\solution
We begin by noting that $9\sin^2\theta + 9\cos^2\theta = 9$, and hence $9\cos^2\theta = 9-9\sin^2\theta$. If we let $x=3\sin\theta$, then $9-x^2 = 9-9\sin^2\theta = 9\cos^2\theta$. 

Setting $x=3\sin \theta$ gives  $\dd x = 3\cos\theta\dd\theta$. We are almost ready to substitute. We also change our bounds of integration. The bound $x=-3$ corresponds to $\theta = -\pi/2$ (for when $\theta = -\pi/2$, $x=3\sin \theta = -3$). Likewise, the bound of $x=3$ is replaced by the bound $\theta = \pi/2$. Thus
\begin{align*}
	\int_{-3}^3\sqrt{9-x^2}\dd x
	&= \int_{-\pi/2}^{\pi/2} \sqrt{9-9\sin^2\theta} (3\cos\theta)\dd\theta \\
	&= \int_{-\pi/2}^{\pi/2} 3\sqrt{9\cos^2\theta} \cos\theta\dd\theta \\
	&=\int_{-\pi/2}^{\pi/2} 3\abs{3\cos \theta} \cos\theta\dd\theta.
	\intertext{On \intertextmath{[-\pi/2,\pi/2]}, \intertextmath{\cos \theta} is always positive, so we can drop the absolute value bars, then employ a half-angle formula:}
	&= \int_{-\pi/2}^{\pi/2} 9\cos^2 \theta\dd\theta\\
	&= \int_{-\pi/2}^{\pi/2} \frac{9}{2}\bigl(1+\cos(2\theta)\bigr)\dd\theta\\
	& = \frac92 \left.\left(\theta +\frac12\sin(2\theta)\right)\right|_{-\pi/2}^{\pi/2}= \frac92\pi.
\end{align*}
This matches our answer from before.
\end{example}

We now describe in detail Trigonometric Substitution. This method excels when dealing with integrands that contain $\sqrt{a^2-x^2}$, $\sqrt{x^2-a^2}$ and $\sqrt{x^2+a^2}$. The following Key Idea outlines the procedure for each case, followed by more examples.
% Each right triangle acts as a reference to help us understand the relationships between $x$ and $\theta$.

\begin{keyidea}[Trigonometric Substitution]\label{idea:trigsub}%
% todo Tim check that the label is good
\mbox{}\\[-2\baselineskip]\index{integration!trig. subst.}\begin{enumerate}%[label=(\alph*)]
	\item\label{trigsubsin} For integrands containing $\sqrt{a^2-x^2}$:\smallskip\\
		Let $x=a\sin\theta$, \quad for $-\pi/2\leq \theta\leq \pi/2$ and $a>0$. \smallskip\\
	On this interval, $\cos\theta\geq 0$, so	$\sqrt{a^2-x^2} = a\cos\theta$
		
	\item\label{trigsubtan} For integrands containing $\sqrt{x^2+a^2}$:\smallskip\\
		Let $x=a\tan\theta$, \quad for $-\pi/2 < \theta < \pi/2$ and $a>0$. \smallskip\\
	On this interval, $\sec\theta> 0$, so $\sqrt{x^2+a^2} = a\sec\theta$
		
	\item\label{trigsubsec} For integrands containing $\sqrt{x^2-a^2}$:\smallskip\\
		Let $x=a\sec\theta$, \quad restricting our work to where $x\geq a>0$,\\
		so $x/a\geq 1$, and $0\leq\theta<\pi/2$. \smallskip\\
	On this interval, $\tan\theta\geq 0$, so	$\sqrt{x^2-a^2} = a\tan\theta$
\end{enumerate}
\end{keyidea}

\begin{example}[Using Trigonometric Substitution]\label{ex_trigsub3}%
Evaluate $\ds \int \frac{1}{\sqrt{5+x^2}}\dd x$.
\solution
Using \autoref{idea:trigsub}\ref*{trigsubtan}, we recognize $a=\sqrt{5}$ and  set $x= \sqrt{5}\tan \theta$. This makes $\dd x = \sqrt{5}\sec^2\theta\dd\theta$. We will use the fact that $\sqrt{5+x^2} = \sqrt{5+5\tan^2\theta} = \sqrt{5\sec^2\theta} = \sqrt{5}\sec\theta.$ Substituting, we have:
\begin{align*}
\int \frac{1}{\sqrt{5+x^2}}\dd x &= \int \frac{1}{\sqrt{5+5\tan^2\theta}}\sqrt{5}\sec^2\theta\dd\theta \\
			&= \int \frac{\sqrt{5}\sec^2\theta}{\sqrt{5}\sec\theta} \dd\theta\\
			&= \int \sec\theta\dd\theta\\
			&= \ln\abs{\sec\theta+\tan\theta}+C.
\end{align*}
While the integration steps are over, we are not yet done. The original problem was stated in terms of $x$, whereas our answer is given in terms of $\theta$. We must convert back to $x$.

\mtable{A reference triangle for \autoref{ex_trigsub3}}{fig:tan_tri}{\begin{tikzpicture}[alt={Right triangle: base √5, vertical leg x, hypotenuse √(x²+5); right angle at foot of x, θ at left vertex.}]
	\draw [very thick] (0,0) -- node[below,pos=.5] {$\sqrt5$} (3,0)
	 -- node [right,pos=.5] {$x$} (3,2)
	 -- node [pos=.5,above,sloped] {$\sqrt{x^2+5}$} cycle;
    \draw [thick] (2.7,0) -- (2.7,.3) -- (3,.3);
	\draw (.75,.25) node {$\theta$};
\end{tikzpicture}}
The lengths of the sides of the reference triangle in \autoref{fig:tan_tri} are determined by the Pythagorean Theorem. With $x=\sqrt{5}\tan\theta$, we have 
\[\tan \theta = \frac x{\sqrt{5}}\quad \text{and}\quad \sec\theta = \frac{\sqrt{x^2+5}}{\sqrt{5}}.\]
This gives\vspace{-.3\baselineskip}
\begin{align*}
	\int\frac1{\sqrt{5+x^2}}\dd x
	&= \ln\abs{\sec\theta+\tan\theta}+C \\
	&= \ln\abs{\frac{\sqrt{x^2+5}}{\sqrt5}+ \frac x{\sqrt5}}+C.
\end{align*}
We can leave this answer as is, or we can use a logarithmic identity to simplify it. Note:\vspace{-.3\baselineskip}
\begin{align*}
	\ln\abs{\frac{\sqrt{x^2+5}}{\sqrt5}+ \frac x{\sqrt5}}+C
	&= \ln\abs{\frac1{\sqrt5}\left(\sqrt{x^2+5}+ x\right)}+C \\
	&= \ln\abs{\frac1{\sqrt5}} + \ln\abs{\sqrt{x^2+5}+ x}+C\\
	&=	\ln\abs{\sqrt{x^2+5}+ x}+C,
\end{align*}
where the $\ln\bigl(1/\sqrt5\bigr)$ term is absorbed into the constant $C$. (In \autoref{sec:hyperbolic} we learned another way of approaching this problem.)
\end{example}

\begin{example}[Using Trigonometric Substitution]\label{ex_trigsub2}%
Evaluate $\ds \int \sqrt{4x^2-1}\dd x$.
\solution
We start by rewriting the integrand so that it has the form $\sqrt{x^2-a^2}$ for some value of $a$:
\begin{align*}
\sqrt{4x^2-1} &= \sqrt{4\left(x^2-\frac14\right)}\\
		&= 2\sqrt{x^2-\left(\frac12\right)^2}.
\end{align*}
So we have $a=1/2$, and following \autoref{idea:trigsub}\ref*{trigsubsec}, we set $x= \frac12\sec\theta$, and hence $\dd x = \frac12\sec\theta\tan\theta\dd\theta$. %The Key Idea also shows that $\sqrt{x^2-1/2^2} = \frac12\tan\theta$. 
We now rewrite the integral with these substitutions:\vspace{-.3\baselineskip}
\begin{align*}
	\int \sqrt{4x^2-1}\dd x &= \int 2\sqrt{x^2-\left(\frac12\right)^2}\dd x\\
	&= \int 2\sqrt{\frac14\sec^2\theta - \frac14}
	\left(\frac12\sec\theta\tan\theta\right)\dd\theta\\
	&=\int \sqrt{\frac14(\sec^2\theta-1)}\Bigl(\sec\theta\tan\theta\Bigr)\dd\theta\\
	&=\int\sqrt{\frac14\tan^2\theta}\Bigl(\sec\theta\tan\theta\Bigr)\dd\theta\\
	&=\int \frac12\tan^2\theta\sec\theta\dd\theta\\
	&=\frac12\int \Bigl(\sec^2\theta-1\Bigr)\sec\theta\dd\theta\\
	&=\frac12\int \bigl(\sec^3\theta - \sec\theta\bigr)\dd\theta.
\end{align*}
We integrated $\sec^3\theta$ in \autoref{ex_trigint6}, finding its antiderivatives to be
\[
\int\sec^3\theta\dd\theta
=\frac12\Bigl(\sec\theta\tan\theta+\ln\abs{\sec\theta+\tan\theta}\Bigr)+C.
\]
Thus
\begin{align*}
\int &\sqrt{4x^2-1}\dd x\\
	&=\frac12\int \bigl(\sec^3\theta-\sec\theta\bigr)\dd\theta\\
	&= \frac12\left(\frac12\Bigl(\sec\theta\tan\theta
	+\ln\abs{\sec\theta+\tan \theta}\Bigr)-\ln\abs{\sec\theta+\tan\theta}\right) + C\\
	%\end{align*}
	%\begin{align*}
	&= \frac14\left(\sec\theta\tan\theta -\ln\abs{\sec\theta+\tan\theta}\right)+C.
\end{align*}
We are not yet done. Our original integral is given in terms of $x$, whereas our final answer, as given, is in terms of $\theta$. We need to rewrite our answer in terms of $x$. With $a=1/2$, and $x=\frac12\sec\theta$, we use the Pythagorean Theorem to determine the lengths of the sides of the reference triangle in \autoref{fig:sec_tri}.
\mtable{A reference triangle for \autoref{ex_trigsub2}}{fig:sec_tri}{\begin{tikzpicture}[alt={Right triangle: base 1/2, vertical leg √(x²-1/4), hypotenuse x; right angle at right, θ at left.}]
	\draw [very thick] (0,0) -- node[below,pos=.5] {$1/2$} (3,0)
	 -- node [right,pos=.5] {$\sqrt{x^2-1/4}$} (3,2)
	 -- node [pos=.5,above] {$x$} cycle;
    \draw [thick] (2.7,0) -- (2.7,.3) -- (3,.3);
	\draw (.75,.25) node {$\theta$};
\end{tikzpicture}}
\[\tan \theta = \frac{\sqrt{x^2-\frac14}}{\frac12} = 2\sqrt{x^2-\frac14}\qquad \text{and}\qquad\sec\theta = 2x.\]
Therefore,
\begin{align*}
	\int \sqrt{4x^2-1}\dd x
	& =\frac14\Bigl(\sec\theta\tan\theta -\ln\abs{\sec\theta+\tan\theta}\Bigr)+C \\
	&=\frac14\Bigl(2x\cdot 2\sqrt{x^2-\frac14} - \ln\abs{2x + 2\sqrt{x^2-\frac14}}\Bigr)+C\\
	&= \frac14\Bigl(4x\sqrt{x^2-\frac14}-\ln\abs{2x + 2\sqrt{x^2-\frac14}}\Bigr)+C \\
	& = \frac14\Bigl(2x\sqrt{4x^2-1} - \ln\abs{2x + \sqrt{4x^2-1}}\Bigr)+C.
\end{align*}
\end{example}

\begin{example}[Using Trigonometric Substitution]\label{ex_trigsub4}%
Evaluate $\ds \int \frac{\sqrt{4-x^2}}{x^2}\dd x$.
\solution
We use \autoref{idea:trigsub}\ref*{trigsubsin} with $a=2$, $x=2\sin \theta$, $\dd x = 2\cos \theta d\theta$ and hence $\sqrt{4-x^2} = 2\cos\theta$. This gives
\begin{align*}
\int \frac{\sqrt{4-x^2}}{x^2}\dd x &= \int \frac{2\cos\theta}{4\sin^2\theta}(2\cos\theta)\dd\theta\\
		&= \int \cot^2\theta\dd\theta\\
		&=	\int (\csc^2\theta -1)\dd\theta\\
		&= -\cot\theta -\theta + C.
\end{align*}
We need to rewrite our answer in terms of $x$. Using the Pythagorean Theorem we determine the lengths of the sides of the reference triangle in \autoref{fig:sin_tri}. We have $\cot\theta = \sqrt{4-x^2}/x$ and $\theta = \sin^{-1}(x/2)$. Thus
\[\int \frac{\sqrt{4-x^2}}{x^2}\dd x = -\frac{\sqrt{4-x^2}}x-\sin^{-1}\left(\frac x2\right) + C.\]
\end{example}

\mtable[yshift=8\baselineskip]{A reference triangle for \autoref{ex_trigsub4}}{fig:sin_tri}{\begin{tikzpicture}[alt={Right triangle: base √(4-x²), vertical leg x, hypotenuse 2; right angle at right, θ at left.}]
	\draw [very thick] (0,0) -- node[below,pos=.5] {$\sqrt{4-x^2}$} (3,0)
	 -- node [right,pos=.5] {$x$} (3,2) -- node [pos=.5,above] {$2$} cycle;
    \draw [thick] (2.7,0) -- (2.7,.3) -- (3,.3);
	\draw (.75,.25) node {$\theta$};
\end{tikzpicture}}%

Trigonometric Substitution can be applied in many situations, even those not of the form $\sqrt{a^2-x^2}$, $\sqrt{x^2-a^2}$ or $\sqrt{x^2+a^2}$. In the following example, we apply it to an integral we already know how to handle.

\begin{example}[Using Trigonometric Substitution]\label{ex_trigsub5}%
Evaluate $\ds \int\frac1{x^2+1}\dd x$.
\solution
We know the answer already as $\tan^{-1}x+C$. We apply Trigonometric Substitution here to show that we get the same answer without inherently relying on knowledge of the derivative of the arctangent function.

Using \autoref{idea:trigsub}\ref*{trigsubtan}, let $x=\tan\theta$, $\dd x=\sec^2\theta\dd\theta$ and note that $x^2+1 = \tan^2\theta+1 = \sec^2\theta$. Thus
\begin{align*}
\int \frac1{x^2+1}\dd x &= \int \frac{1}{\sec^2\theta}\sec^2\theta\dd\theta \\
			&= \int 1\dd\theta\\
			&= \theta + C.
\end{align*}
Since $x=\tan \theta$, $\theta = \tan^{-1}x$, and we conclude that $\ds \int\frac1{x^2+1}\dd x = \tan^{-1}x+C.$
\end{example}

The next example is similar to the previous one in that it does not involve a square-root. It shows how several techniques and identities can be combined to obtain a solution.

\begin{example}[Using Trigonometric Substitution]\label{ex_trigsub7}%
Evaluate $\ds\int\frac1{(x^2+6x+10)^2}\dd x$.
\solution
We start by completing the square, then make the substitution $u=x+3$, followed by the trigonometric substitution of $u=\tan\theta$:
\mnote{\textbf{Note:} Remember the sine and cosine double angle identities:
\begin{align*}
 \sin2\theta &= 2\sin\theta\cos\theta\\
 \cos2\theta &= \cos^2\theta-\sin^2\theta \\
 &= 2\cos^2\theta-1 \\
 &= 1-2\sin^2\theta
\end{align*}
They are often needed for writing your final answer in terms of $x$.}
\begin{align}
\int \frac1{(x^2+6x+10)^2}\dd x &=\int \frac1{\bigl((x+3)^2+1\bigr)^2}\dd x = \int \frac1{(u^2+1)^2}\dd u. \notag
\intertext{Now make the substitution \intertextmath{u=\tan\theta}, \intertextmath{\dd u=\sec^2\theta\dd\theta}:}
   &=	\int \frac1{(\tan^2\theta+1)^2}\sec^2\theta\dd\theta\notag\\
	&= \int\frac 1{(\sec^2\theta)^2}\sec^2\theta\dd\theta\notag\\
	&= \int \cos^2\theta\dd\theta.\notag
	\intertext{Applying a half-angle formula, we have}
	&= \int \left(\frac12 +\frac12\cos(2\theta)\right)\dd\theta \notag\\
	&= \frac12\theta + \frac14\sin(2\theta) + C.\label{eq:extrigsub7}
\end{align}
We need to return to the variable $x$. As $u=\tan\theta$, $\theta = \tan^{-1}u$. Using the identity $\sin(2\theta) = 2\sin\theta\cos\theta$ and using a reference triangle, %found in \autoref{idea:trigsub}(b),
we have 
\[\frac14\sin(2\theta) = \frac12\frac u{\sqrt{u^2+1}}\cdot\frac 1{\sqrt{u^2+1}} = \frac12\frac u{u^2+1}.\]
Finally, we return to $x$ with the substitution $u=x+3$. We start with the expression in \autoeqref{eq:extrigsub7}:
\begin{align*}
	\frac12\theta + \frac14\sin(2\theta) + C
	&= \frac12\tan^{-1}u + \frac12\frac{u}{u^2+1}+C\\
	&= \frac12\tan^{-1}(x+3) + \frac{x+3}{2(x^2+6x+10)}+C.
\end{align*}
Stating our final result in one line,
\[
\int\frac1{(x^2+6x+10)^2}\dd x=\frac12\tan^{-1}(x+3)+\frac{x+3}{2(x^2+6x+10)}+C.
\]
\end{example}

Our last example returns us to definite integrals, as seen in our first example. Given a definite integral that can be evaluated using Trigonometric Substitution, we could first evaluate the corresponding indefinite integral (by changing from an integral in terms of $x$ to one in terms of $\theta$, then converting back to $x$) and then evaluate using the original bounds. It is much more straightforward, though, to change the bounds as we substitute.

\begin{example}[Definite integration and Trigonometric Substitution]\label{ex_trigsub6}%
Evaluate $\ds\int_0^5\frac{x^2}{\sqrt{x^2+25}}\dd x$.
\solution
Using \autoref{idea:trigsub}\ref*{trigsubtan}, we set $x=5\tan\theta$, $\dd x = 5\sec^2\theta\dd\theta$, and note that $\sqrt{x^2+25} = 5\sec\theta$. As we substitute, we change the bounds of integration.

The lower bound of the original integral is $x=0$. As $x=5\tan\theta$, we solve for $\theta$ and find $\theta = \tan^{-1}(x/5)$. Thus the new lower bound is $\theta = \tan^{-1}(0) = 0$. The original upper bound is $x=5$, thus the new upper bound is $\theta = \tan^{-1}(5/5) = \pi/4$. 

Thus we have 
\begin{align*}
\int_0^5\frac{x^2}{\sqrt{x^2+25}}\dd x &= \int_0^{\pi/4} \frac{25\tan^2\theta}{5\sec\theta}5\sec^2\theta\dd\theta\\
		&= 25\int_0^{\pi/4} \tan^2\theta\sec\theta\dd\theta.
\end{align*}
We encountered this indefinite integral in \autoref{ex_trigsub2} where we found 
\[\int \tan^2\theta\sec\theta \dd\theta = \frac12\bigl(\sec\theta\tan\theta-\ln\abs{\sec\theta+\tan\theta}\bigr).\]
So
\begin{align*}
	25\int_0^{\pi/4} \tan^2\theta\sec\theta\dd\theta
	&= \left.\frac{25}2\left(\sec\theta\tan\theta-\ln\abs{\sec\theta+\tan\theta}\right)\right|_0^{\pi/4}\\
	&= \frac{25}2\bigl(\sqrt2-\ln(\sqrt2+1)\bigr)
%	\\
%	&\approx 6.661
	.
\end{align*}
\end{example}

%The following equalities are very useful when evaluating integrals using Trigonometric Substitution.
%
%\begin{keyidea}[Useful Equalities with Trigonometric Substitution]\label{idea:useful_trigsub}%
%\begin{enumerate}
%	\item	$\sin(2\theta) = 2\sin\theta\cos\theta$
%	\item	$\cos(2\theta) = \cos^2\theta - \sin^2\theta = 2\cos^2\theta-1 = 1-2\sin^2\theta$
%	\item $\ds \int \sec^3\theta\dd\theta = \frac12\Bigl(\sec \theta\tan \theta + \ln\abs{\sec \theta+\tan \theta}\Bigr)+C$
%	\item	$\ds \int \cos^2\theta\dd\theta = \int \frac12\bigl(1+\cos(2\theta)\bigr)\dd\theta = \frac12\bigl(\theta+\sin\theta\cos\theta\bigr)+C.$
%\end{enumerate}
%\end{keyidea}

The next section introduces Partial Fraction Decomposition, which is an algebraic technique that turns ``complicated'' fractions into sums of ``simpler'' fractions, making integration easier.

\printexercises{exercises/06-08-exercises}

% this was orphaned from u-substitution

%\begin{example}[Integration by substitution: simplifying first]\label{ex_sub9}%
%Evaluate $\ds\int \frac{x^3+4x^2+8x+5}{x^2+2x+1}\dd x$.
%\solution
%One may try to start by setting $u$ equal to either the numerator or denominator; in each instance, the result is not workable. 
%
%When dealing with rational functions (i.e., quotients made up of polynomial functions), it is an almost universal rule that everything works better when the degree of the numerator is less than the degree of the denominator. Hence we use polynomial division.
%
%We skip the specifics of the steps, but note that when $x^2+2x+1$ is divided into $x^3+4x^2+8x+5$, it goes in $x+2$ times with a remainder of $3x+3$. Thus 
%\[\frac{x^3+4x^2+8x+5}{x^2+2x+1} = x+2 + \frac{3x+3}{x^2+2x+1}.\]
%Integrating $x+2$ is simple. The fraction can be integrated by setting $u = x^2+2x+1$, giving $\dd u = (2x+2)\dd x$. This is very similar to the numerator. Note that $\dd u/2 = (x+1)\dd x$ and then consider the following:
%\begin{align*} % this had ``rule'' commands scattered throughout
%\int \frac{x^3+4x^2+8x+5}{x^2+2x+1}\dd x
% & = \int \left(x+2 + \frac{3x+3}{x^2+2x+1}\right)\dd x \\
% &= \int (x+2)\dd x + \int \frac{3(x+1)}{x^2+2x+1}\dd x \\
% & = \frac12x^2+2x+C_1 + \int \frac{3}{u}\frac{\dd u}{2} \\
% &= \frac12x^2+2x+C_1 + \frac32\ln\abs u + C_2 \\
% &= \frac12x^2+2x+\frac32\ln\abs{x^2+2x+1} + C.
%\end{align*}
%In some ways, we ``lucked out'' in that after dividing, substitution was able to be done. In later sections we'll develop techniques for handling rational functions where substitution is not directly feasible.
%\end{example}



\section{Partial Fraction Decomposition}\label{sec:partial_fraction}

In this section we investigate the antiderivatives of rational functions. Recall that rational functions are functions of the form $f(x)= \frac{p(x)}{q(x)}$, where $p(x)$ and $q(x)$ are polynomials and $q(x)\neq 0$. Such functions arise in many contexts, one of which is the solving of certain fundamental differential equations.

We begin with an example that demonstrates the motivation behind this section. Consider the integral $\ds\int \frac{1}{x^2-1}\dd x$. We do not have a simple formula for this (if the denominator were $x^2+1$, we would recognize the antiderivative as being the arctangent function). It can be solved using Trigonometric Substitution, but note how the integral is easy to evaluate once we realize:

%This integral is not difficult to evaluate once one realizes the following fact: 
\[\frac{1}{x^2-1} = \frac{1/2}{x-1} - \frac{1/2}{x+1}.\]
Thus 
\begin{align*}
\int\frac{1}{x^2-1}\dd x &= \int\frac{1/2}{x-1}\dd x - \int\frac{1/2}{x+1}\dd x \\
			&= \frac12\ln\abs{x-1}-\frac12\ln\abs{x+1}+C.
\end{align*}

This section teaches how to \emph{decompose}
\[\frac1{x^2-1}\quad  \text{into}\quad  \frac{1/2}{x-1}-\frac{1/2}{x+1}.\]

We start with a rational function $f(x)=\dfrac{p(x)}{q(x)}$, where $p$ and $q$ do not have any common factors. We first consider the degree of $p$ and $q$. 
\begin{itemize}
	\item If the $\deg(p)\ge\deg(q)$ then we use polynomial long division to divide $q$ into $p$ to determine a remainder $r(x)$ where $\deg(r)<\deg(q)$. We then write $f(x) =s(x)+\dfrac{r(x)}{q(x)}$ and apply partial fraction decomposition to $\dfrac{r(x)}{q(x)}$.
	\item If the $\deg(p)<\deg(q)$ we can apply partial fraction decomposition to $\dfrac{p(x)}{q(x)}$ without additional work.
\end{itemize}

Partial fraction decomposition is based on an algebraic theorem that guarantees that any polynomial, and hence $q$, can use real numbers to factor into the product of linear and irreducible quadratic factors.
\mnote{An \emph{irreducible quadratic} is one that cannot factor into linear terms with real coefficients.}\index{irreducible quadratic}
The following Key Idea states how to decompose a rational function into a sum of rational functions whose denominators are all of lower degree than $q$.

{
\tcbset{grow to right by=6em}
\begin{keyidea}[Partial Fraction Decomposition]\label{idea:partial_fraction}%
Let $\dfrac{p(x)}{q(x)}$ be a rational function, where $\deg(p)<\deg(q)$.
\begin{enumerate}
	\item \textbf {Factor} $\mathbf{q(x):}$ Write $q(x)$ as the product of its linear and irreducible quadratic factors of the form $(ax+b)^m$ and $(ax^2+bx+c)^n$ where $m$ and $n$ are the highest powers of each factor that divide $q$.
	\begin{itemize}
		\item \textbf{Linear Terms:} For each linear factor of $q(x)$ the decomposition of $\dfrac{p(x)}{q(x)}$ will contain the following terms:
		\[\frac{A_1}{(ax+b)}+\frac {A_2}{(ax+b)^2}+\dotsb\frac{A_m}{(ax+b)^m}\]
		\item \textbf{Irreducible Quadratic Terms:}  For each irreducible quadratic factor of $q(x)$ the decomposition of $\dfrac {p(x)}{q(x)}$ will contain the following terms:
		\[
		 \dfrac{B_1x+C_1}{(ax^2+bx+c)}+\frac{B_2x+C_2}{(ax^2+bx+c)^2}
		 +\dotsb\frac{B_nx+C_n}{(ax^2+bx+c)^n}
		\]
	\end{itemize}
	\item \textbf{Finding the Coefficients $\mathbf{A_i}$, $\mathbf{B_i}$, and $\mathbf{C_i}$:}
	\begin{itemize}
		\item Set $\dfrac{p(x)}{q(x)}$ equal to the sum of its linear and irreducible quadratic terms.
		\[
		 \frac{p(x)}{q(x)}
		 =\frac{A_1}{(ax+b)}+\dotsb\frac{A_m}{(ax+b)^m}
		 +\frac{B_1x+C_1}{(ax^2+bx+c)}+\dotsb\frac{B_nx+C_n}{(ax^2+bx+c)^n}
		\]
		\item Multiply this equation by the factored form of $q(x)$ and simplify to clear the denominators. 
		\item Solve for the coefficients $A_i, B_i,$ and $C_i$ by
		\begin{enumerate}
			\item multiplying out the remaining terms and collecting like powers of $x$, equating the resulting coefficients and solving the resulting system of linear equations, \textbf{or}
			\item substituting in values for $x$ that eliminate terms so the simplified equation can be solved for a coefficient.
		\end{enumerate}
	\end{itemize}
\end{enumerate}
\end{keyidea}
}

\youtubeVideo{6qVgHWxdlZ0}{Integration Using method of Partial Fractions}

The following examples will demonstrate how to put this Key Idea into practice. In \autoref{ex_pf1}, we focus on the setting up the decomposition of a rational function.

\begin{example}[Decomposing into partial fractions]\label{ex_pf1}%
Decompose $\ds f(x)=\frac{1}{(x+5)(x-2)^3(x^2+x+2)(x^2+x+7)^2}$ without solving for the resulting coefficients.
\solution
The denominator is already factored, as both $x^2+x+2$ and $x^2+x+7$ are irreducible quadratics. We need to decompose $f(x)$ properly. Since $(x+5)$ is a linear factor that divides the denominator, there will be a
\[\frac{A}{x+5}\]
term in the decomposition.

As $(x-2)^3$ divides the denominator, we will have the following terms in the decomposition:
\[\frac{B}{x-2},\quad \frac{C}{(x-2)^2}\quad \text{and}\quad \frac{D}{(x-2)^3}.\]

The $x^2+x+2$ term in the denominator results in a $\ds\frac{Ex+F}{x^2+x+2}$ term.

Finally, the $(x^2+x+7)^2$ term results in the terms
\[\frac{Gx+H}{x^2+x+7}\quad \text{and}\quad \frac{Ix+J}{(x^2+x+7)^2}.\]
All together, we have
\begin{multline*}
	\frac{1}{(x+5)(x-2)^3(x^2+x+2)(x^2+x+7)^2}= \\
	\frac{A}{x+5} + \frac{B}{x-2}+ \frac{C}{(x-2)^2}+\frac{D}{(x-2)^3}+\\
	\frac{Ex+F}{x^2+x+2}+\frac{Gx+H}{x^2+x+7}+\frac{Ix+J}{(x^2+x+7)^2}
\end{multline*}
Solving for the coefficients $A$, $B$, \dots, $J$ would be a bit tedious but not ``hard.''  In the next example we demonstrate solving for the coefficients using both methods given in \autoref{idea:partial_fraction}.
\end{example}

%(-37 x-39)/(140800 (x^2+x+2))+(67804 x+21113)/(520524225 (x^2+x+7))+(89 x-32)/(296595 (x^2+x+7)^2)+665617/(5015768576 (x-2))-1/(5501034 (x+5))-1119/(6889792 (x-2)^2)+1/(9464 (x-2)^3)

\begin{example}[Decomposing into partial fractions]\label{ex_pf2}%
Perform the partial fraction decomposition of $\ds \frac{1}{x^2-1}$.
\solution
The denominator can be written as the product of two linear factors: $x^2-1 = (x-1)(x+1)$. Thus 
\begin{equation}\label{eq:decomp2}%
 \frac{1}{x^2-1} = \frac{A}{x-1} + \frac{B}{x+1}.
\end{equation}
Using the method described in \autoref{idea:partial_fraction} 2(a) to solve for $A$ and $B$, first multiply through by $x^2-1 = (x-1)(x+1)$:
\begin{align}\label{eq:pf2}%
	1
	&= \frac{A(x-1)(x+1)}{x-1}+\frac{B(x-1)(x+1)}{x+1} \notag\\
	&= A(x+1) + B(x-1)\\
	&= Ax+A + Bx-B \notag\\
	&= (A+B)x + (A-B)\qquad\text{collect like terms.}\notag
\end{align}
The next step is key. %Note the equality we have:
%\[1 = (A+B)x+(A-B).\]
For clarity's sake, rewrite the equality we have as
\[0x+1 = (A+B)x+(A-B).\]
On the left, the coefficient of the $x$ term is 0; on the right, it is $(A+B)$. Since both sides are equal for all values of $x$, we must have that $0=A+B$. Likewise, on the left, we have a constant term of 1; on the right, the constant term is $(A-B)$. Therefore we have $1=A-B$.

We have two linear equations with two unknowns. This one is easy to solve by hand, leading to 
\[
 \begin{aligned}A+B&=0\\A-B&=1\end{aligned}
 \qquad\Rightarrow\qquad
 \begin{aligned}A&=1/2\\B&=-1/2.\end{aligned}
\]
Thus
\[\frac{1}{x^2-1}=\frac{1/2}{x-1}-\frac{1/2}{x+1}.\]

Before solving for $A$ and $B$ using the method described in \autoref{idea:partial_fraction} 2(b), we note that Equations \eqref{eq:decomp2} and \eqref{eq:pf2} are not equivalent. Only the second equation holds for all values of $x$, including $x=-1$ and $x=1$, by continuity of polynomials. Thus, we can choose values for $x$ that eliminate terms in the polynomial to solve for $A$ and $B$.
\[1=A(x+1) + B(x-1).\]
If we choose $x=-1$,
\begin{align*}
 1&=A(0) + B(-2) \\
 B&=-\frac{1}{2}.
\end{align*}
Next choose $x=1$:
\begin{align*}
 1&=A(2) + B(0) \\
 A&=\frac{1}{2}.
\end{align*}
Resulting in the same decomposition as above.
\end{example}

In \autoref{ex_pf3}, we solve for the decomposition coefficients using the system of linear equations (method 2a). The margin note explains how to solve using substitution (method 2b).

\begin{example}[Integrating using partial fractions]\label{ex_pf3}%
Use partial fraction decomposition to integrate $\ds\int\frac{1}{(x-1)(x+2)^2}\dd x$.
\solution
We decompose the integrand as follows, as described by \autoref{idea:partial_fraction}:
\begin{equation}\label{eq:decomp3}%
 \frac{1}{(x-1)(x+2)^2} = \frac{A}{x-1} + \frac{B}{x+2} + \frac{C}{(x+2)^2}.
\end{equation}
To solve for $A$, $B$ and $C$, we multiply both sides by $(x-1)(x+2)^2$ and collect like terms:
%
\mnote{\textbf{Note:} Equations \eqref{eq:decomp3} and \eqref{eq:pf3} are not equivalent for $x=1$ and $x=-2$. However, due to the continuity of polynomials we can let $x=1$ to simplify the right hand side to $A(1+2)^2=9A$. Since the left hand side is still $1$, we have $1=9A$, so that $A=1/9$.\bigskip\\
Likewise,when $x=-2$; this leads to the equation $1=-3C$. Thus $C = -1/3$.\bigskip\\
Knowing $A$ and $C$, we can find the value of $B$ by choosing yet another value of $x$, such as $x=0$, and solving for $B$.}
%
\begin{align}
	1
	&= A(x+2)^2 + B(x-1)(x+2) + C(x-1)\label{eq:pf3}\\
	&= Ax^2+4Ax+4A + Bx^2 + Bx-2B + Cx-C \notag \\
	&= (A+B)x^2 + (4A+B+C)x + (4A-2B-C)\notag
\end{align}
We have
\[0x^2+0x+ 1 = (A+B)x^2 + (4A+B+C)x + (4A-2B-C)\]
leading to the equations 
\[A+B = 0, \quad 4A+B+C = 0 \quad \text{and} \quad 4A-2B-C = 1.\]
These three equations of three unknowns lead to a unique solution:
\[A = 1/9,\quad B = -1/9 \quad \text{and} \quad C = -1/3.\]
Thus 
\[
\int\frac{1}{(x-1)(x+2)^2}\dd x = \int \frac{1/9}{x-1}\dd x + \int \frac{-1/9}{x+2}\dd x + \int \frac{-1/3}{(x+2)^2}\dd x.
\]

Each can be integrated with a simple substitution with $u=x-1$ or $u=x+2$.
% (or by directly applying \autoref{idea:linearsub} as the denominators are linear functions).
The end result is
\[\int\frac{1}{(x-1)(x+2)^2}\dd x = \frac19\ln\abs{x-1}-\frac19\ln\abs{x+2} +\frac1{3(x+2)}+C.\]
\end{example}

\begin{example}[Integrating using partial fractions]\label{ex_pf4}%
Use partial fraction decomposition to integrate $\ds \int \frac{x^3}{(x-5)(x+3)}\dd x$.
\solution
\autoref{idea:partial_fraction} presumes that the degree of the numerator is less than the degree of the denominator. Since this is not the case here, we begin by using polynomial division to reduce the degree of the numerator. We omit the steps, but encourage the reader to verify that
\[\frac{x^3}{(x-5)(x+3)} = x+2+\frac{19x+30}{(x-5)(x+3)}.\]
Using \autoref{idea:partial_fraction}, we can rewrite the new rational function as:
\[\frac{19x+30}{(x-5)(x+3)} = \frac{A}{x-5} + \frac{B}{x+3}\]
for appropriate values of $A$ and $B$. Clearing denominators, we have 
\[19x+30=A(x+3)+B(x-5).\]
As in the previous examples we choose values of $x$ to eliminate terms in the polynomial.  If we choose $x=-3$,
\begin{align*}
 19(-3)+30&=A(0) + B(-8) \\
 B&= \frac{27}{8}.
\end{align*}
Next choose $x=5$:
\begin{align*}
 19(5)+30&=A(8) + B(0)\\
 A&= \frac{125}{8}.
\end{align*}
We can now integrate:
\begin{align*}
	\int \frac{x^3}{(x-5)(x+3)}\dd x
	&= \int\left(x+2+\frac{125/8}{x-5}+\frac{27/8}{x+3}\right)\dd x \\
	&= \frac{x^2}2+2x+\frac{125}{8}\ln\abs{x-5}+\frac{27}8\ln\abs{x+3}+C.
\end{align*}
\end{example}

Before the next example we remind the reader of a rational integrand evaluated by trigonometric substitution:
\[\int\frac1{x^2+a^2}\dd x=\frac1a\tan^{-1}\left(\frac xa\right) + C.\]

\begin{example}[Integrating using partial fractions]\label{ex_pf5}%
Use partial fraction decomposition to evaluate $\ds \int\frac{7x^2+31x+54}{(x+1)(x^2+6x+11)}\dd x$.
\solution
The degree of the numerator is less than the degree of the denominator so we begin by applying \autoref{idea:partial_fraction}. We have:
\begin{align*}
\frac{7x^2+31x+54}{(x+1)(x^2+6x+11)} &= \frac{A}{x+1} + \frac{Bx+C}{x^2+6x+11}. \\
\intertext{Now clear the denominators.}
7x^2+31x+54 &= A(x^2+6x+11) + (Bx+C)(x+1).
\end{align*}
Again, we choose values of $x$ to eliminate terms in the polynomial.  If we choose $x=-1$,
\begin{align*}
 30&=6A + (-B+C)(0)\\
 A&= 5.
\end{align*}

Although none of the other terms can be zeroed out, we continue by letting $A=5$ and substituting helpful values of $x$. 
Choosing $x=0$, we notice
\begin{align*}
 54&= 55 +C \\
 C&= -1.
\end{align*}
Finally, choose $x=1$ (any value other than $-1$ and $0$ can be used, $1$ is easy to work with)
\begin{align*}
 92&=90 + (B-1)(2)\\
 B&= 2.
\end{align*}
Thus
\[
 \int\frac{7x^2+31x+54}{(x+1)(x^2+6x+11)}\dd x
 = \int\left(\frac{5}{x+1} + \frac{2x-1}{x^2+6x+11}\right)\dd x.
\]

The first term of this new integrand is easy to evaluate; it leads to a $5\ln\abs{x+1}$ term. The second term is not hard, but takes several steps and uses substitution techniques.

The integrand $\dfrac{2x-1}{x^2+6x+11}$ has a quadratic in the denominator and a linear term in the numerator. This leads us to try substitution. Let $u=x^2+6x+11$, so $du=(2x+6)\dd x$. The numerator is $2x-1$, not $2x+6$, but we can get a $2x+6$ term in the numerator by adding 0 in the form of ``$7-7$.''
\begin{align*}
	\frac{2x-1}{x^2+6x+11} &= \frac{2x-1+7-7}{x^2+6x+11} \\
	&= \frac{2x+6}{x^2+6x+11} - \frac{7}{x^2+6x+11}.
\end{align*}
We can now integrate the first term with substitution, yielding $\ln\abs{x^2+6x+11}$. The final term can be integrated using arctangent. First, complete the square in the denominator:
\[\frac{7}{x^2+6x+11} = \frac{7}{(x+3)^2+2}.\]
An antiderivative of the latter term can be found using \autoref{idea:trigsub} and substitution:
\[
 \int \frac7{x^2+6x+11}\dd x=\frac7{\sqrt2}\tan^{-1}\left(\frac{x+3}{\sqrt2}\right)+C.
\]

Let's start at the beginning and put all of the steps together.
\begin{align*}
	\int&\frac{7x^2+31x+54}{(x+1)(x^2+6x+11)}\dd x \\
	&= \int\left(\frac{5}{x+1} + \frac{2x-1}{x^2+6x+11}\right)\dd x \\
	&= \int\frac{5}{x+1}\dd x  + \int\frac{2x+6}{x^2+6x+11}\dd x -\int\frac{7}{x^2+6x+11}\dd x \\
	&= 5\ln\abs{x+1}+\ln\abs{x^2+6x+11}-\frac7{\sqrt2}\tan^{-1}\left(\frac{x+3}{\sqrt2}\right)+C.
\end{align*}
As with many other problems in calculus, it is important to remember that one is not expected to ``see'' the final answer immediately after seeing the problem. Rather, given the initial problem, we break it down into smaller problems that are easier to solve. The final answer is a combination of the answers of the smaller problems.
\end{example}

Partial Fraction Decomposition is an important tool when dealing with rational functions. Note that at its heart, it is a technique of algebra, not calculus, as we are rewriting a fraction in a new form. Regardless, it is very useful in the realm of calculus as it lets us evaluate a certain set of ``complicated'' integrals.  The next section will require the reader to determine an appropriate method for evaluating a variety of integrals.

\printexercises{exercises/06-04-exercises}

\section{Integration Strategies}\label{sec:int_techniques}

We've now seen a fair number of different integration techniques and so we should probably pause at this point to talk a little bit about a strategy to use for determining the correct technique to use when faced with an integral.

There are a couple of points that need to be made about this strategy. First, it isn't a hard and fast set of rules for determining the method that should be used. It is really nothing more than a general set of guidelines that will help us to identify techniques that may work. Some integrals can be done in more than one way and so depending on the path you take through the strategy you may end up with a different technique than someone else who also went through this strategy.

Second, while the strategy is presented as a way to identify the technique that could be used on an integral keep in mind that, for many integrals, it can also automatically exclude certain techniques as well. When going through the strategy keep two lists in mind. The first list is integration techniques that simply won't work and the second list is techniques that look like they might work. After going through the strategy, if the second list has only one entry then that is the technique to use. If on the other hand, there is more than one possible technique to use we will have to decide on which is liable to be the best for us to use. Unfortunately there is no way to teach which technique is the best as that usually depends upon the person and which technique they find to be the easiest.

Third, don't forget that many integrals can be evaluated in multiple ways and so more than one technique may be used on it. This has already been mentioned in each of the previous points, but is important enough to warrant a separate mention. Sometimes one technique will be significantly easier than the others and so don't just stop at the first technique that appears to work. Always identify all possible techniques and then go back and determine which you feel will be the easiest for you to use.

Next, it's entirely possible that you will need to use more than one method to completely evaluate an integral. For instance a substitution may lead to using integration by parts or partial fractions integral.

\clearpage

\begin{keyidea}[Guidelines for Choosing an Integration Strategy]\label{ki_int_strat}%
\mbox{}\\[-2\baselineskip]\parbox[t]{\linewidth}{%
\begin{enumerate}
	\item Simplify the integrand, if possible.
	\item See if a ``simple'' substitution will work.
	\item Identify the type of integral.
	\item Relate the integral to an integral we already know how to do.
	\item Try multiple techniques.
	\item Try again.
\end{enumerate}}
\end{keyidea}

Let's expand on the ideas of the previous Key Idea.
\begin{enumerate}
\item \textbf{Simplify the integrand, if possible.} This step is very important in the integration process. Many integrals can be taken from very difficult to very easy with a little simplification or manipulation. Don't forget basic trigonometric and algebraic identities as these can often be used to simplify the integral.

We used this idea when we were looking at integrals involving trigonometric functions. For example consider the integral
\[\int \cos^2 x\dd x.\]
The integral can't be done as it is, however by recalling the identity,
\[\cos^2 x = \frac{1}{2}(1 + \cos 2x)\]
the integral becomes very easy to do.

Note that this example also shows that simplification does not necessarily mean that we'll write the integrand in a ``simpler'' form. It only means that we'll write the integrand in a form that we can deal with and this is often longer and/or ``messier'' than the original integral.

\item \textbf{See if a ``simple'' substitution will work.} Look to see if a simple substitution can be used instead of the often more complicated methods from this chapter. For example consider both of the following integrals.
\[\int \frac{x}{x^2-1}\dd x \qquad \int x\sqrt{x^2 -1}\dd x\]
The first integral can be done with the method of partial fractions and the second could be done with a trigonometric substitution.

However, both could also be evaluated using the substitution $u=x^2 -1$ and the work involved in the substitution would be significantly less than the work involved in either partial fractions or trigonometric substitution.

So, always look for quick, simple substitutions before moving on to the more complicated techniques of this chapter.

\item \textbf{Identify the type of integral.} Note that any integral may fall into more than 
one of these types. Because of this fact it's usually best to go all the way through the list and identify all possible types since one may be easier than the other and it's entirely possible that the easier type is listed lower in the list.

\begin{enumerate}
\item Is the integrand a rational expression (i.e. is the integrand a polynomial divided by a polynomial)? If so then partial fractions (\autoref{sec:partial_fraction}) may work on the integral.

\item Is the integrand a polynomial times a trigonometric function, exponential, or logarithm? If so, then integration by parts (\autoref{sec:IBP}) may work.

\item Is the integrand a product of sines and cosines, secants and tangents, or cosecants and cotangents? If so, then the topics from \autoref{sec:trigint} may work. Likewise, don't forget that some quotients involving these functions can also be done using these techniques.

\item Does the integrand involve $\sqrt {b^2x^2 + a^2}, \sqrt {b^2x^2 - a^2}, \text{or} \sqrt { a^2-b^2x^2}$? If so, then a trigonometric substitution (\autoref{sec:trig_sub}) might work nicely.

\item Does the integrand have roots other than those listed above in it? If so then the substitution $u=\sqrt[n]{g(x)}$ might work.

\item Does the integrand have a quadratic in it? If so then completing the square on the quadratic might put it into a form that we can deal with. 
\end{enumerate}

\item \textbf{Relate the integral to an integral we already know how to do.} In other words, can we use a substitution or manipulation to write the integrand into a  form that does fit into the forms we've looked at previously in this chapter. A typical example is the following integral.
\[\int \cos x\sqrt{1 + \sin^2 x}\dd x\]
This integral doesn't obviously fit into any of the forms we looked at in this chapter. However, with the substitution $u = \sin x$ we can reduce the integral to the form
\[\int \sqrt{1+u^2}\dd x\]
which is a trigonometric substitution problem.

\item \textbf{Try multiple techniques.} In this step we need to ask ourselves if it is possible that we'll need to use multiple techniques. The example in the previous part is a good example. Using a substitution didn't allow us to actually do the integral. All it did was put the integral into a form that we could use a different technique on.

Don't ever get locked into the idea that an integral will only require one step to completely evaluate it. Many will require more than one step.

\item \textbf{Try again.} If everything that you've tried to this point doesn't work then go back through the process again. This time try a technique that you didn't use the first time around.
\end{enumerate}

As noted above, this strategy is not a hard and fast set of rules. It is only intended to guide you through the process of best determining how to do any given integral. Note as well that the only place Calculus II actually arises is the third step. Steps 1, 2, and 4 involve nothing more than manipulation of the integrand either through direct manipulation of the integrand or by using a substitution. The last two steps are simply ideas to think about in going through this strategy.

Many students go through this process and concentrate almost exclusively on Step 3 (after all this is Calculus II, so it's easy to see why they might do that...) to the exclusion of the other steps. One very large consequence of that exclusion is that often a simple manipulation or substitution is overlooked that could make the integral very easy to do.

Before moving on to the next section we will work a couple of examples illustrating a couple of not so obvious simplifications/manipulations and a not so obvious substitution.

\begin{example}[Strategies of Integration]\label{ex_int_strat_tan_sec}%
Evaluate the integral
\[\int \frac{\tan x}{\sec^4 x}\dd x\]
\solution
This integral almost falls into the form given in 3c. It is a quotient of tangent and secant and we know that sometimes we can use the same methods for products of tangents and secants on quotients.

The process from \autoref{sec:trigint} tells us that if we have even powers of secant to save two of them and convert the rest to tangents. That won't work here. We can save two secants, but they would be in the denominator and they won't do us any good here. Remember that the point of saving them is so they could be there for the substitution $u = \tan x$. That requires them to be in the numerator. So, that won't work. We need to find another solution method.

There are in fact two solution methods to this integral depending on how you want to go about it.

\textbf{Solution 1} \indent In this solution method we could just convert everything to sines and cosines and see if that gives us an integral we can deal with.
\begin{align*}
	\int \frac{\tan x}{\sec^4 x}\dd x
	&= \int \frac{\sin x}{\cos x} \cos^4 x\dd x\\
	&= \int \sin x \cos^3 x\dd x \qquad\text{substitute $u = \cos x$}\\
	&= -\int u^3\dd u \\
	&=-\frac{1}{4} \cos^4 x +C
\end{align*}

Note that just converting to sines and cosines won't always work and if it does it won't always work this nicely. Often there will be a lot more work that would need to be done to complete the integral.\bigskip

\textbf{Solution 2} \indent This solution method goes back to dealing with secants and tangents. Let's notice that if we had a secant in the numerator we could just use $u = \sec x$ as a substitution and it would be a fairly quick and simple substitution to use. We don't have a secant in the numerator. However we could very easily get a secant in the numerator by multiplying the numerator and denominator by secant (i.e. we multiply the integrand by ``1''). 
\begin{align*}
	\int \frac{\tan x}{\sec^4 x}\dd x
	&= \int \frac{\tan x \sec x}{\sec^5 x}\dd x \qquad\text{substitute $u = \sec x$}\\
	&= \int \frac{1}{u^5}\dd u \\
	&=-\frac{1}{4} \frac{1}{\sec^4 x}+C\\
	&=-\frac{1}{4} \cos^4 x +C
\end{align*}
\end{example}

In the previous example we saw two ``simplifications'' that allowed us to evaluate the integral. The first was using identities to rewrite the integral into terms we could deal with and the second involved multiplying the numerator and denominator by something to again put the integral into terms we could deal with.

Using identities to rewrite an integral is an important ``simplification'' and we should not forget about it. Integrals can often be greatly simplified or at least put into a form that can be dealt with by using an identity.

The second ``simplification'' is not used as often, but does show up on occasion so again, it's best to remember it. In fact, let's take another look at an example in which multiplying the integrand by ``1'' will allow us to evaluate an integral.

\begin{example}[Strategy for Integration]\label{ex_int_strat_sin}%
Evaluate the integral
\[\int \frac{1}{1+\sin x}\dd x\]
\solution
This is an integral which if we just concentrate on the third step we won't get anywhere. This integral doesn't appear to be any of the kinds of integrals that we worked on in this chapter. We can evaluate the integral however, if we do the following,
\begin{align*}
	\int \frac{1}{1+\sin x}\dd x
	&= \int \frac{1}{1+\sin x} \frac{1-\sin x}{1-\sin x}\dd x \\
	&= \int \frac{1-\sin x}{1-\sin^2 x}\dd x 
\end{align*}

This does not appear to have done anything for us. However, if we now remember the first ``simplification'' we looked at above we will notice that we can use an identity to rewrite the denominator. Once we do that we can further manipulate the integrand into something we can evaluate.
\begin{align*}
	\int \frac{1}{1+\sin x}\dd x
	&= \int \frac{1-\sin x}{\cos^2 x}\dd x \\
	&= \int \frac{1}{\cos^2 x} - \frac{\sin x}{\cos x}\frac{1}{\cos x}\dd x \\
	&= \int \sec^2 x - \tan x \sec x\dd x \\
	&= \tan x - \sec x +C
\end{align*}
\end{example}

So, we've just seen once again that multiplying by a helpful form of ``1'' can put the integral into a form we can integrate. Notice as well that this example also showed that ``simplifications'' do not necessarily put an integral into a simpler form. They only put the integrand into a form that is easier to integrate.

Let's now take a quick look at an example of a substitution that is not so obvious.

\begin{example}[Strategy for Integration]\label{ex_int_strat_cos_sqrt}%
Evaluate the integral
\[\int \cos \sqrt x\dd x\]
\solution
We introduced this integral by saying that the substitution was not so obvious. However, this is really an integral that falls into the form given by 3e in \autoref{ki_int_strat}. Many people miss that form and so don't think about it. So, let's try the following substitution.
\[u = \sqrt x \qquad x=u^2 \qquad\dd x=2u\dd u\]
With this substitution the integral becomes,
\[\int \cos \sqrt x\dd x = 2\int u \cos u\dd u\]
This is now an integration by parts. Remember that often we will need to use more than one technique to completely do the integral. This is a fairly simple integration by parts problem so we'll leave the remainder of the details for you to check.
\[\int \cos \sqrt x\dd x =2(\cos \sqrt x + \sqrt x \sin \sqrt x) + C.\]
\end{example}

\bigskip
It will be possible to integrate every integral assigned in this class, but it is important to note that there are integrals that just can't be evaluated. We should also note that after we look at series in \autoref{chapter:sequences_series} we will be able to write down a series representation of many of these types of integrals.

% todo write exercises using reduction formulas

\printexercises{exercises/06-09-exercises}

\section{Improper Integration}\label{sec:improper_integration}

We begin this section by considering the following definite integrals:
\begin{itemize}
\item	$\ds \int_0^{100}\frac1{1+x^2}\dd x \approx 1.5608,$
\item	$\ds \int_0^{1000}\frac1{1+x^2}\dd x \approx 1.5698,$
\item	$\ds \int_0^{10,000}\frac1{1+x^2}\dd x \approx 1.5707.$
\end{itemize}

Notice how the integrand is $1/(1+x^2)$ in each integral (which is sketched in \autoref{fig:improper1}). As the upper bound gets larger, one would expect the ``area under the curve'' would also grow. While the definite integrals do increase in value as the upper bound grows, they are not  increasing by much. In fact, consider:
\[\int_0^b \frac{1}{1+x^2}\dd x = \tan^{-1}x\Big|_0^b = \tan^{-1}b-\tan^{-1}0 = \tan^{-1}b.\]
As $b\rightarrow \infty$, $\tan^{-1}b \rightarrow \pi/2.$ Therefore it seems that as the upper bound $b$ grows, the value of the definite integral $\ds \int_0^b\frac{1}{1+x^2}\dd x$ approaches $\pi/2\approx 1.5708$. This should strike the reader as being a bit amazing: even though the curve extends ``to infinity,'' it has a finite amount of area underneath it.

\mtable{Graphing $\ds f(x)=\frac{1}{1+x^2}$.}{fig:improper1}{\begin{tikzpicture}[alt={Blue curve f(x)=1/(1+x²) on 0 < x < 10, falling from y=1 at x=0 toward 0; area under curve to x axis is shaded.}]
\begin{axis}[width=1.16\marginparwidth,tick label style={font=\scriptsize},
axis y line=middle,axis x line=middle,name=myplot,axis on top,
ymin=-.1,ymax=1.1,xmin=-1,xmax=11]
\addplot[draw={\colorone},fill={\coloronefill},domain=0:11,area style,thick]
 {1/(1+x*x)} |-(axis cs:0,0);
\end{axis}
\node [right] at (myplot.right of origin) {\scriptsize $x$};
\node [above] at (myplot.above origin) {\scriptsize $y$};
\end{tikzpicture}}

When we defined the definite integral $\ds\int_a^b f(x)\dd x$, we made two stipulations:
	\begin{enumerate}
	\item		The interval over which we integrated, $[a,b]$, was a finite interval, and
	\item		The function $f(x)$ was continuous on $[a,b]$ (ensuring that the range of $f$ was finite).
	\end{enumerate}
	
In this section we consider integrals where one or both of the above conditions do not hold. Such integrals are called \textbf{improper integrals.}

\clearpage
% the page breaks here regardless, but without the clearpage,
% it gets justified and looks bad
% clearpage causes the page to be raggedbottom
% we can consider
% https://tex.stackexchange.com/q/226716/107497
% to make just this page raggedbottom and not need the clearpage

\subsection{Improper Integrals with Infinite Bounds}

\begin{definition}[Improper Integrals with Infinite Bounds]\label{def:imp_int1}%
\index{integration!improper}\index{improper integration}\index{convergence!of improper int.}\index{divergence!of improper int.}
\mbox{}\\[-2\baselineskip]\parbox[t]{\linewidth}{%
\begin{enumerate}
\item		Let $f$ be a continuous function on $[a,\infty)$. For $t \geq a$ let
\[\int_a^\infty f(x)\dd x = \lim_{t\to\infty}\int_a^t f(x)\dd x.\]

\item		Let $f$ be a continuous function on $(-\infty,b]$. For $t \leq b$ let
\[\int_{-\infty}^b f(x)\dd x = \lim_{t\to-\infty}\int_t^b f(x)\dd x.\]

\item		Let $f$ be a continuous function on $(-\infty,\infty)$. For any real number $c$ (which one doesn't matter), let
\[
\int_{-\infty}^\infty f(x)\dd x
= \lim_{a\to-\infty}\int_a^c f(x)\dd x\ +\ \lim_{b\to\infty}\int_c^b f(x)\dd x.
\]
\end{enumerate}}
\end{definition}

An improper integral is said to \textbf{converge} if its corresponding limit exists and is finite; otherwise, it \textbf{diverges}. The improper integral in part 3 converges if and only if both of its limits exist.

\youtubeVideo{f6cGotvktxs}{Improper Integral --- Infinity in Upper and Lower Limits}

\begin{example}[Evaluating improper integrals]\label{ex_impint1}%
Evaluate the following improper integrals.\\
\begin{minipage}[t]{.5\textwidth}
\begin{enumerate}
\item		$\ds\int_1^\infty \frac1{x^2}\dd x$
\item		$\ds\int_1^\infty \frac1x\dd x$
\end{enumerate}
\end{minipage}%
\begin{minipage}[t]{.5\textwidth}
\begin{enumerate}\addtocounter{enumi}{2}
\item		$\ds\int_{-\infty}^0 e^x\dd x$
\item		$\ds\int_{-\infty}^\infty \frac1{1+x^2}\dd x$
\end{enumerate}
\end{minipage}%
\solution
\begin{enumerate}
	\item
	%
\mtable{A graph of $f(x) = \frac{1}{x^2}$ in\\\autoref{ex_impint1} part 1.}{fig:impint1a}{\begin{tikzpicture}[alt={Blue curve f(x)=1/x² on 1 < x < 10; rapid drop from y=1 at x=1 then tails to 0; region under curve and above x axis is shaded.}]
\begin{axis}[width=1.16\marginparwidth,tick label style={font=\scriptsize},
axis y line=middle,axis x line=middle,name=myplot,axis on top,xtick={1,5,10},
ymin=-.1,ymax=1.1,xmin=-1,xmax=11]
\addplot [draw={\colorone},fill={\coloronefill},area style,domain=1:10.5,thick]
 {1/(x*x)} |- (axis cs:1,0);
\draw (axis cs:5,.75) node {\scriptsize $\ds f(x)=\frac{1}{x^2}$};
\end{axis}
\node [right] at (myplot.right of origin) {\scriptsize $x$};
\node [above] at (myplot.above origin) {\scriptsize $y$};
\end{tikzpicture}}
%
	\hfill$\begin{aligned}[t]
		\int_1^\infty \frac{1}{x^2}\dd x
		& = \lim_{t\to\infty} \int_1^t\frac1{x^2}\dd x \\
		& = \lim_{t\to\infty} \left.\frac{-1}{x}\right|_1^t \\ 
		%&= \lim_{t\to\infty} \frac{-1}{x}\Big|_1^t \\
		&= \lim_{t\to\infty} \frac{-1}{t} + 1\\
		&= 1.
	\end{aligned}$\hfill\null\\
	A graph of the area defined by this integral is given in \autoref{fig:impint1a}.

	\item
%
\mtable{A graph of $f(x) = \frac{1}{x}$ in\\\autoref{ex_impint1} part 2.}{fig:impint1b}{\begin{tikzpicture}[alt={Blue curve f(x)=1/x on 1 < x < 10; decreases from y=1 at x=1 to y=0.1 at x=10; shaded area under curve and above x axis.}]
\begin{axis}[width=1.16\marginparwidth,tick label style={font=\scriptsize},
axis y line=middle,axis x line=middle,name=myplot,axis on top,xtick={1,5,10},
ymin=-.1,ymax=1.1,xmin=-1,xmax=11]
\addplot [draw={\colorone},fill={\coloronefill},area style,domain=1:10.5,thick]
 {1/x} |- (axis cs:1,0);
\draw (axis cs:5,.75) node {\scriptsize $\ds f(x)=\frac{1}{x}$};
\end{axis}
\node [right] at (myplot.right of origin) {\scriptsize $x$};
\node [above] at (myplot.above origin) {\scriptsize $y$};
\end{tikzpicture}}
%
	\hfill$\begin{aligned}[t]
		\int_1^\infty \frac1x\dd x
		& = \lim_{t\to\infty}\int_1^t\frac1x\dd x \\
		&= \lim_{t\to\infty} \ln \abs{x}\Bigr|_1^t \\
		&= \lim_{t\to\infty} \ln (t)\\
		&= \infty.
	\end{aligned}$\hfill\null\\
	The limit does not exist, hence the improper integral $\ds\int_1^\infty\frac1x\dd x$ diverges. Compare the graphs in Figures \ref{fig:impint1a} and \ref{fig:impint1b}; notice how the values of $f(x) = 1/x$ are noticeably larger than those of $f(x)=1/x^2$. This difference is enough to cause the improper integral to diverge.

	\item
%
\mtable{A graph of $f(x) = e^x$ in\\\autoref{ex_impint1} part 3.}{fig:impint1c}{\begin{tikzpicture}[alt={Blue curve f(x)=e^x on –10 < x < –1; climbs steeply from near 0 to y 1; region right of x=–10, left of x=–1 and above x axis is shaded.}]
\begin{axis}[width=1.16\marginparwidth,tick label style={font=\scriptsize},
axis y line=middle,axis x line=middle,name=myplot,axis on top,xtick={-1,-5,-10},
ytick={1},ymin=-.1,ymax=1.1,xmin=-11,xmax=1]
\addplot [draw={\colorone},fill={\coloronefill},area style,domain=-10.5:0,thick]
 {exp(x)} |- (axis cs:0,0);
\draw (axis cs:-5,.75) node {\scriptsize $\ds f(x)=e^x$};
\end{axis}
\node [right] at (myplot.right of origin) {\scriptsize $x$};
\node [above] at (myplot.above origin) {\scriptsize $y$};
\end{tikzpicture}}
%
	\hfill$\begin{aligned}[t]
		\int_{-\infty}^0 e^x\dd x
		&= \lim_{t\to-\infty} \int_t^0 e^x\dd x \\
		&=  \lim_{t\to-\infty} e^x\Big|_t^0 \\
		&= \lim_{t\to-\infty} \bigl(e^0-e^t\bigr) \\
		&= 1.
	\end{aligned}$\hfill\null\\
	A graph of the area defined by this integral is given in \autoref{fig:impint1c}.

	\item	We will need to break this into two improper integrals and choose a value of $c$ as in part 3 of \autoref{def:imp_int1}. Any value of $c$ is fine; we choose $c=0$.
%
\mtable{A graph of $f(x) = \frac{1}{1+x^2}$ in \autoref{ex_impint1} part 4.}{fig:impint1d}{\begin{tikzpicture}[alt={Blue graph of f(x)=1/(1+x²), symmetric, peak y=1 at x=0, then falls to y=0 by x=10; area below curve to x-axis shaded.}]
\begin{axis}[width=1.16\marginparwidth,tick label style={font=\scriptsize},
axis y line=middle,axis x line=middle,name=myplot,axis on top,ytick={1},
ymin=-.1,ymax=1.1,xmin=-11,xmax=11]
\addplot [draw={\colorone},fill={\coloronefill},area style,domain=-10.5:10.5,thick,samples=100]
 {1/(1+x*x)} |- (axis cs:0,0);
\draw (axis cs:6,.75) node {\scriptsize $\ds f(x)=\frac{1}{1+x^2}$};
\end{axis}
\node [right] at (myplot.right of origin) {\scriptsize $x$};
\node [above] at (myplot.above origin) {\scriptsize $y$};
\end{tikzpicture}}
%
\begin{align*}
	\int_{-\infty}^\infty \frac1{1+x^2}\dd x
	&= \lim_{t\to-\infty} \int_t^0\frac{1}{1+x^2}\dd x + \lim_{t\to\infty} \int_0^t\frac{1}{1+x^2}\dd x \\
	&= \lim_{t\to-\infty} \tan^{-1}x\Big|_t^0 + \lim_{t\to\infty} \tan^{-1}x\Big|_0^t\\
	&= \lim_{t\to-\infty} \left(\tan^{-1}0-\tan^{-1}t\right) + \lim_{t\to\infty} \left(\tan^{-1}t-\tan^{-1}0\right)\\		
	&= \left(0-\frac{-\pi}2\right) + \left(\frac{\pi}2-0\right).\\
	&= \pi.
\end{align*}
A graph of the area defined by this integral is given in \autoref{fig:impint1d}.
\end{enumerate}
\end{example}

\autoref{sec:lhopitals_rule} introduced L'Hôpital's Rule, a method of evaluating limits that return indeterminate forms. It is not uncommon for the limits resulting from improper integrals to need this rule as demonstrated next.

\begin{example}[Improper integration and L'Hôpital's Rule]\label{ex_impint2}%
Evaluate the improper integral $\ds \int_1^\infty \frac{\ln x}{x^2}\dd x$.
\solution
This integral will require the use of Integration by Parts. Let $u = \ln x$ and $dv = 1/x^2\dd x$. Then
%
\mtable{A graph of $f(x) = \frac{\ln x}{x^2}$ in \autoref{ex_impint2}.}{fig:impint2}{\begin{tikzpicture}[alt={Blue curve f(x)=(ln x)/ x² on 1<x<10; rises to small hump near x=1.3 then decays toward 0; region under curve is filled.}]
\begin{axis}[width=1.16\marginparwidth,tick label style={font=\scriptsize},
axis y line=middle,axis x line=middle,name=myplot,axis on top,xtick={1,5,10},
ymin=-.1,ymax=.5,xmin=-1,xmax=11]
\addplot [draw={\colorone},fill={\coloronefill},area style,domain=1:10.5,thick,samples=100]
 {ln(x)/(x*x)} |- (axis cs:1,0);
\draw (axis cs:5,.3) node {\scriptsize $\ds f(x)=\frac{\ln x}{x^2}$};
\end{axis}
\node [right] at (myplot.right of origin) {\scriptsize $x$};
\node [above] at (myplot.above origin) {\scriptsize $y$};
\end{tikzpicture}}
%
\begin{align*}
	\int_1^\infty\frac{\ln x}{x^2}\dd x
	&= \lim_{t\to\infty}\int_1^t\frac{\ln x}{x^2}\dd x \\
	&=  \lim_{t\to\infty}\left(-\frac{\ln x}{x}\Big|_1^t +\int_1^t \frac{1}{x^2}\dd x \right)\\
	&=  \lim_{t\to\infty} \left.\left(-\frac{\ln x}{x} -\frac1x\right)\right|_1^t\\
	&=	\lim_{t\to\infty} \left(-\frac{\ln t}{t}-\frac1t - \left(-\ln 1-1\right)\right).
\end{align*}
The $1/t$ goes to 0, and $\ln 1=0$, leaving
% $\ds \lim_{t\to\infty} -\frac{\ln t}t + 1$. We need to evaluate
$\ds \lim_{t\to\infty} \frac{\ln t}{t}$ with L'Hôpital's Rule. We have:
\[\lim_{t\to\infty}\frac{\ln t}t \LHequals \lim_{t\to\infty} \frac{1/t}{1} = 0.\]
Thus the improper integral evaluates as:
\[\int_1^\infty\frac{\ln x}{x^2}\dd x = 1.\]
\end{example}

\subsection{Improper Integrals with Infinite Range}

We have just considered definite integrals where the interval of integration was infinite. We now consider another type of improper integration, where the range of the integrand is infinite.

\begin{definition}[Improper Integration with Infinite Range]\label{def:imp_int2}%
Let $f(x)$ be a continuous function on $[a,b]$ except at $c$, $a\leq c\leq b$, where $x=c$ is a vertical asymptote of $f$. Define\index{integration!improper}\index{improper integration}
\[\int_a^b f(x)\dd x = \lim_{t\to c^-}\int_a^t f(x)\dd x + \lim_{t\to c^+}\int_t^b f(x)\dd x.\]
%The integral converges if both limits exist and diverges otherwise.
\end{definition}

Note that $c$ can be one of the endpoints ($a$ or $b$). In that case, there is only one limit to consider as part of the definition.

\begin{example}[Improper integration of functions with infinite range]\label{ex_impint3}%
Evaluate the following improper integrals:
\[
 \text{1. }\int_0^1\frac1{\sqrt{x}}\dd x\qquad\qquad
 \text{2. }\int_{-1}^1\frac{1}{x^2}\dd x.
\]
\solution
\begin{enumerate}
\item		A graph of $f(x) = 1/\sqrt{x}$ is given in \autoref{fig:impint3}.
%
\mtable{A graph of $f(x)=\frac{1}{\sqrt{x}}$ in \autoref{ex_impint3}.}{fig:impint3}{\begin{tikzpicture}[alt={Blue curve f(x)=1/√x on 0<x<1; shoots to y 10 near x=0, decreases to y=1 at x=1; vertical strip at x=1 and area beneath are shaded.}]
\begin{axis}[width=1.16\marginparwidth,tick label style={font=\scriptsize},
axis y line=middle,axis x line=middle,name=myplot,axis on top,
ymin=-.1,ymax=11,xmin=-.1,xmax=1.1]
\addplot [draw={\colorone},fill={\coloronefill},area style,thick,domain=.01:1] {1/sqrt(x)} |- (axis cs:.01,0);
\draw (axis cs:.5,5) node {\scriptsize $\ds f(x)=\frac{1}{\sqrt{x}}$};
\end{axis}
\node [right] at (myplot.right of origin) {\scriptsize $x$};
\node [above] at (myplot.above origin) {\scriptsize $y$};
\end{tikzpicture}}
%
Notice that $f$ has a vertical asymptote at $x=0$. In some sense, we are trying to compute the area of a region that has no ``top.'' Could this have a finite value? 
\begin{align*}
	\int_0^1 \frac{1}{\sqrt{x}}\dd x
	&= \lim_{t\to0^+}\int_t^1 \frac1{\sqrt{x}}\dd x \\
	&= \lim_{t\to0^+} 2\sqrt{x}\Big|_t^1 \\
	&= \lim_{t\to0^+} 2\left(\sqrt{1}-\sqrt{t}\right)\\
	&=	2.
\end{align*}
It turns out that the region does have a finite area even though it has no upper bound (strange things can occur in mathematics when considering the infinite).

\item		The function $f(x) = 1/x^2$ has a vertical asymptote at $x=0$, as shown in \autoref{fig:impint3b}, so this integral is an improper integral. Let's eschew using limits for a moment and proceed without recognizing the improper nature of the integral. This leads to:
\begin{align*}
\int_{-1}^1\frac1{x^2}\dd x &= -\frac1x\Big|_{-1}^1\\
			&= -1 - (1)\\
			&=-2.
\end{align*}
%
\mtable{A graph of $f(x)=\frac{1}{x^2}$ in \autoref{ex_impint3}.}{fig:impint3b}{\begin{tikzpicture}[alt={Blue graph f(x)=1/x² on −1<x<1 with a tall central spike (clipped at y=10) and symmetric flanks; full region under curve shaded.}]
\begin{axis}[width=1.16\marginparwidth,tick label style={font=\scriptsize},
axis y line=middle,axis x line=middle,name=myplot,axis on top,
ymin=-.1,ymax=11,xmin=-1.1,xmax=1.1]
\addplot [draw={\colorone},fill={\coloronefill},smooth,thick,domain=-1:-.2]
 {1/(x*x)} |-(axis cs:-1,0);
\draw[draw={\coloronefill},fill={\coloronefill}](axis cs:-.25,0)rectangle(axis cs:.25,11);
\addplot [draw={\coloronefill},fill={\coloronefill},smooth,thick,domain=.2:1]
 {1/(x*x)} |-(axis cs:.2,0);
% for some reason, the previous puts a \colorone bar on the right
\addplot [draw={\colorone},smooth,thick,domain=.2:1] {1/(x*x)};
\draw (axis cs:.7,7) node {\scriptsize $\ds f(x)=\frac{1}{x^2}$};
\end{axis}
\node [right] at (myplot.right of origin) {\scriptsize $x$};
\node [above] at (myplot.above origin) {\scriptsize $y$};
\end{tikzpicture}}
%
Clearly the area in question is above the $x$-axis, yet the area is supposedly negative. In this example we noted the discontinuity of the integrand on $[-1,1]$ (its improper nature) but continued anyway to apply the Fundamental Theorem of Calculus. Violating the hypothesis of the FTC led us to an incorrect area of $-2$. If we now evaluate the integral using \autoref{def:imp_int2} we will see that the area is unbounded.
\begin{align*}
	\int_{-1}^1\frac1{x^2}\dd x
	&= \lim_{t\to0^-}\int_{-1}^t \frac1{x^2}\dd x + \lim_{t\to0^+}\int_t^1\frac1{x^2}\dd x \\
	&= \lim_{t\to0^-}-\frac1x\Big|_{-1}^t + \lim_{t\to0^+}-\frac1x\Big|_t^1\\
	&= \lim_{t\to0^-}\left(-\frac1t-1\right) + \lim_{t\to0^+}\left(-1+\frac1t\right).
\end{align*}
Neither limit converges hence the original improper integral diverges. The nonsensical answer we obtained by ignoring the improper nature of the integral is just that: nonsensical.
\end{enumerate}
\end{example}

\subsection{Understanding Convergence and Divergence}

Oftentimes we are interested in knowing simply whether or not an improper integral converges, and not necessarily the value of a convergent integral. We provide here several tools that help determine the convergence or divergence of improper integrals without integrating.

Our first tool is knowing the behavior of functions of the form $\dfrac1{x\primeskip^p}$.

\begin{example}[Improper integration of $1/x^p$]\label{ex_impint4}%
Determine the values of $p$ for which $\ds \int_1^\infty \frac1{x^p}\dd x$ converges.
\solution
We begin by integrating and then evaluating the limit.
\begin{align*}
	\int_1^\infty \frac1{x\primeskip^p}\dd x
	&= \lim_{t\to\infty}\int_1^t\frac1{x\primeskip^p}\dd x\\
	&= \lim_{t\to\infty}\int_1^t x^{-p}\dd x \qquad \text{\small (assume $p\neq 1$)}\\
	&= \lim_{t\to\infty} \frac{1}{-p+1}x^{-p+1}\Big|_1^t\\
	&= \lim_{t\to\infty} \frac{1}{1-p}\bigl(t^{1-p}-1^{1-p}\bigr).\\
\end{align*}
%
\mtable{Plotting functions of the form $1/x\,^p$ in \autoref{ex_impint4}.}{fig:impint4}{\begin{tikzpicture}[alt={Three decreasing curves on log axes: red f(x)=1/x^p (gentlest), dashed 1/x, blue f(x)=1/x^q (steepest); p<1<q.}]
\begin{axis}[width=1.16\marginparwidth,tick label style={font=\scriptsize},
axis y line=middle,axis x line=middle,name=myplot,axis on top,xtick={1},
ytick=\empty,ymin=-.1,ymax=6,xmin=-.1,xmax=2.1]
\addplot [dashed,thick,smooth,domain=.1:2] {1/x};
\addplot [draw={\colorone},thick,smooth,domain=.1:2] {1/x^1.5};
\addplot [draw={\colortwo},thick,smooth,domain=.05:2] {1/sqrt(x)};
\draw (axis cs:.7,5) node {\scriptsize $\ds f(x)=\frac{1}{x\,^q}$};
\draw (axis cs:.35,.8) node {\scriptsize $\ds f(x)=\frac{1}{x\,^p}$};
\draw (axis cs:1.5,3) node {\scriptsize $p<1<q$};
\end{axis}
\node [right] at (myplot.right of origin) {\scriptsize $x$};
\node [above] at (myplot.above origin) {\scriptsize $y$};
\end{tikzpicture}}%
%
When does this limit converge --- i.e., when is this limit \emph{not} $\infty$? This limit converges precisely when the power of $t$ is less than 0: when $1-p<0 \Rightarrow 1<p$. 

Our analysis shows that if $p>1$, then $\ds\int_1^\infty \frac1{x\primeskip^p}\dd x $ converges. When $p<1$ the improper integral diverges; we showed in \autoref{ex_impint1} that when $p=1$ the integral also diverges. 

\autoref{fig:impint4} graphs $y=1/x$ with a dashed line, along with graphs of $y=1/x^p$, $p<1$, and $y=1/x^q$, $q>1$. Somehow the dashed line forms a dividing line between convergence and divergence. %A function of the form $1/x^q$ will be under the dashed line on $[1,\infty)$ when $q>1$. Even if $q$ is ``very close'' to 1, the difference will be enough to force convergence.
\end{example}

The result of \autoref{ex_impint4} provides an important tool in determining the convergence of other integrals. A similar result is proved in the exercises about improper integrals of the form $\ds \int_0^1\frac1{x\primeskip^p}\dd x$. These results are summarized in the following Key Idea.

{
\tcbset{grow to right by=8em}
\begin{keyidea}[Convergence of Improper Integrals $\ds \int_1^\infty\frac1{x\primeskip^p}\dd x$ and $\ds \int_0^1\frac1{x\primeskip^p}\dd x$.]\label{idea:impint1}%
\index{convergence!of improper int.}\index{divergence!of improper int.}
\mbox{}\\[-1.5\baselineskip]\parbox[t]{\linewidth}{%
\begin{enumerate}
\item		The improper integral $\ds \int_1^\infty\frac1{x\primeskip^p}\dd x$ converges when $p>1$ and diverges when $p\leq 1.$
\item		The improper integral $\ds \int_0^1\frac1{x\primeskip^p}\dd x$ converges when $p<1$ and diverges when $p\geq 1.$
\end{enumerate}}
\end{keyidea}}

A basic technique in determining convergence of improper integrals is to compare an integrand whose convergence is unknown to an integrand whose convergence is known. We often use integrands of the form $1/x\primeskip^p$ in comparisons as their convergence on certain intervals is known. This is described in the following theorem.

\mnote{\textbf{Note:} We used the upper and lower bound of ``1'' in \autoref{idea:impint1} for convenience. It can be replaced by any $a$ where $a>0$. }

\begin{theorem}[Direct Comparison Test for Improper Integrals]\label{thm:impint_comparison}%
Let $f$ and $g$ be continuous on $[a,\infty)$ where $0\leq f(x)\leq g(x)$ for all $x$ in $[a,\infty)$.
\index{integration!improper}\index{convergence!Direct Comparison Test!for integration}\index{divergence!Direct Comparison Test!for integration}\index{Direct Comparison Test!for integration}\index{convergence!of improper int.}\index{divergence!of improper int.}
	\begin{enumerate}
	\item		If $\ds \int_a^\infty g(x)\dd x$ converges, then $\ds \int_a^\infty f(x)\dd x$ converges.
	\item		If $\ds \int_a^\infty f(x)\dd x$ diverges, then $\ds \int_a^\infty g(x)\dd x$ diverges.
	\end{enumerate}

%\item		Let $f$ and $g$ be continuous functions on $[a,b]$ except at $x=c$, where each has a vertical asymptote, and $0\leq f(x)\leq g(x)$ for all $x$ in $[a,b]$, $x\neq c$.  
%	\begin{itemize}
%	\item		If $\ds \int_a^b g(x)\dd x$ converges, then $\ds \int_a^b f(x)\dd x$ converges.
%	\item		If $\ds \int_a^b f(x)\dd x$ diverges, then $\ds \int_a^b g(x)\dd x$ diverges.
%	\end{itemize}
%\end{itemize}
\end{theorem}

\begin{example}[Determining convergence of improper integrals]\label{ex_impint5}%
Determine the convergence of the following improper integrals.
\[
 \text{1. }\int_1^\infty e^{-x^2}\dd x\qquad\qquad
 \text{2. }\int_3^\infty \frac{1}{\sqrt{x^2-x}}\dd x
\]
\solution
\begin{enumerate}
\item		The function $f(x) = e^{-x^2}$ does not have an antiderivative expressible in terms of elementary functions, so we cannot integrate directly. It is comparable to $g(x)=1/x^2$, and as demonstrated in \autoref{fig:impint5}, $e^{-x^2} < 1/x^2$ on $[1,\infty)$. We know from \autoref{idea:impint1} that $\ds \int_1^\infty \frac{1}{x^2}\dd x$ converges, hence $\ds\int_1^\infty e^{-x^2}\dd x$ also converges.

\mtable{Graphs of $f(x) = e^{-x^2}$ and $f(x)= 1/x^2$ in \autoref{ex_impint5}.}{fig:impint5}{\begin{tikzpicture}[alt={Red 1/x² and blue e^(–x²) on 1<x<4; both fall toward the axis, red stays above blue after x=1.3.}]
\begin{axis}[width=\marginparwidth,tick label style={font=\scriptsize},
axis y line=middle,axis x line=middle,name=myplot,axis on top,
ymin=-.1,ymax=1.1,xmin=-.1,xmax=4.1]
\addplot [draw={\colorone},thick,smooth,domain=1:4] {exp(-x*x)};
\addplot [draw={\colortwo},thick,smooth,domain=1:4] {1/(x*x)};
\draw (axis cs:.7,.45) node {\scriptsize $\ds f(x)=e^{-x^2}$};
\draw (axis cs:1.9,.8) node {\scriptsize $\ds f(x)=\frac{1}{x^2}$};
\end{axis}
\node [right] at (myplot.right of origin) {\scriptsize $x$};
\node [above] at (myplot.above origin) {\scriptsize $y$};
\end{tikzpicture}}

\item		Note that for large values of $x$, $\ds \frac{1}{\sqrt{x^2-x}} \approx \frac{1}{\sqrt{x^2}} =\frac{1}{x}$. We know from \autoref{idea:impint1} and the subsequent note that  $\ds \int_3^\infty \frac1x\dd x$ diverges, so we seek to compare the original integrand to $1/x$.

It is easy to see that when $x>0$, we have $x = \sqrt{x^2} > \sqrt{x^2-x}$. Taking reciprocals reverses the inequality, giving
\[\frac1x < \frac1{\sqrt{x^2-x}}.\]

Using \autoref{thm:impint_comparison}, we conclude that since $\ds\int_3^\infty\frac1x\dd x$ diverges, then $\ds\int_3^\infty\frac1{\sqrt{x^2-x}}\dd x$ diverges as well. \autoref{fig:impint5b} illustrates this.

\mtable{Graphs of $f(x) = 1/\sqrt{x^2-x}$ and $f(x)= 1/x$ in \autoref{ex_impint5}.}{fig:impint5b}{\begin{tikzpicture}[alt={Blue 1 / √(x² – x) and red 1 / x on 2<x<6; curves descend, blue slightly above red throughout.}]
\begin{axis}[width=\marginparwidth,tick label style={font=\scriptsize},
axis y line=middle,axis x line=middle,name=myplot,axis on top,
ymin=-.1,ymax=.5,xmin=-.1,xmax=6.2]
\addplot [draw={\colorone},thick,smooth,domain=3:6] {1/sqrt(x*x-x)};
\addplot [draw={\colortwo},thick,smooth,domain=3:6] {1/x};
\draw (axis cs:4,.45) node {\scriptsize $\ds f(x)=\frac{1}{\sqrt{x^2-x}}$};
\draw (axis cs:2.5,.2) node {\scriptsize $\ds f(x)=\frac{1}{x}$};
\end{axis}
\node [right] at (myplot.right of origin) {\scriptsize $x$};
\node [above] at (myplot.above origin) {\scriptsize $y$};
\end{tikzpicture}}
\end{enumerate}
\end{example}

Being able to compare ``unknown'' integrals to ``known'' integrals is very useful in determining convergence. However, some of our examples were a little ``too nice.'' For instance, it was convenient that $\ds \frac{1}x < \frac{1}{\sqrt{x^2-x}}$, but what if the ``$-x$'' were replaced with a ``$+2x+5$''? That is, what can we say about the convergence of $\ds \int_3^\infty\frac{1}{\sqrt{x^2+2x+5}}\dd x$? We have $\ds \frac{1}{x} > \frac1{\sqrt{x^2+2x+5}}$, so we cannot use \autoref{thm:impint_comparison}.

In cases like this (and many more) it is useful to employ the following theorem.

\begin{theorem}[Limit Comparison Test for Improper Integrals]\label{thm:impint_limit}%
Let $f$ and $g$ be continuous functions on $[a,\infty)$ where $f(x)>0$ and $g(x)>0$ for all $x$. If\vspace{-.5\baselineskip}
\[\lim_{x\to\infty} \frac{f(x)}{g(x)} = L,\qquad 0<L<\infty,\]
	then\vspace{-.5\baselineskip}
	\[\int_a^\infty f(x)\dd x \quad \text{and} \quad \int_a^\infty g(x)\dd x\]
	either both converge or both diverge.%
	\index{integration!improper}\index{convergence!Limit Comparison Test!for integration}\index{divergence!Limit Comparison Test!for integration}\index{Limit Comparison Test!for integration}\index{convergence!of improper int.}\index{divergence!of improper int.}
\end{theorem}

\begin{example}[Determining convergence of improper integrals]\label{ex_impint6}%
Determine the convergence of $\ds \int_3^{\infty} \frac{1}{\sqrt{x^2+2x+5}}\dd x$.
\solution
As $x$ gets large, the square root of a quadratic function will begin to behave much like $y=x$. So we compare \small$\ds\frac{1}{\sqrt{x^2+2x+5}}$\normalsize\ to \small$\ds\frac1x$\normalsize\ with the Limit Comparison Test:
\[
\lim_{x\to\infty} \frac{1/\sqrt{x^2+2x+5}}{1/x}
= \lim_{x\to\infty}\frac{x}{\sqrt{x^2+2x+5}}.
\]

The immediate evaluation of this limit returns $\infty/\infty$, an indeterminate form. Using L'Hôpital's Rule seems appropriate, but in this situation, it does not lead to useful results. (We encourage the reader to employ L'Hôpital's Rule at least once to verify this.)

The trouble is the square root function.
%To get rid of it, we employ the following fact: If $\ds \lim_{x\to c} f(x) = L$, then $\ds\lim_{x\to c} f(x)^2 = L^2.$ (This is true when either $c$ or $L$ is $\infty$.) So we consider now the limit
%\[\lim_{x\to\infty} \frac{x^2}{x^2+2x+5}.\]
%This converges to 1, meaning the original limit also converged to 1. As $x$ gets very large, the function $\frac1{\sqrt{x^2+2x+5}}$ looks very much like $\frac1x$.
We determine the limit by using a technique we learned in \autoref{key:rat_lim_at_inf}:
%
\mtable{Graphing $f(x)=\frac{1}{\sqrt{x^2+2x+5}}$ and $f(x)=\frac1x$ in \autoref{ex_impint6}.}{fig:impint6}{\begin{tikzpicture}[alt={Blue 1 / √(x²+2x+5) and red 1/x on 3<x<20; red starts above then meets blue near x=10 before both approach zero.}]
\begin{axis}[width=\marginparwidth,tick label style={font=\scriptsize},
axis y line=middle,axis x line=middle,name=myplot,axis on top,
ymin=-.1,ymax=.35,xmin=-.1,xmax=21]
\addplot [draw={\colortwo},thick,smooth,domain=3:20] {1/sqrt(x*x+2*x+5)};
\addplot [draw={\colorone},thick,smooth,domain=3:20] {1/x};
\draw (axis cs:10,.3) node {\scriptsize $\ds f(x)=\frac1x$};
\draw (axis cs:4,.09) node {\scriptsize $\ds f(x)=\frac1{\sqrt{x^2+2x+5}}$};
\end{axis}
\node [right] at (myplot.right of origin) {\scriptsize $x$};
\node [above] at (myplot.above origin) {\scriptsize $y$};
\end{tikzpicture}}
%
\[
 \lim_{x\to\infty}\frac x{\sqrt{x^2+2x+5}}
 =\lim_{x\to\infty}\frac{\frac xx}{\sqrt{\frac{x^2+2x+5}{x^2}}}
 =\lim_{x\to\infty}\frac1{\sqrt{1+\frac2x+\frac5{x^2}}}=1
\]
Since we know that $\ds\int_3^{\infty} \tfrac1x\dd x$ diverges, by the Limit Comparison Test we know that $\ds\int_3^\infty\tfrac1{\sqrt{x^2+2x+5}}\dd x$ also diverges. \autoref{fig:impint6} graphs $f(x)=1/\sqrt{x^2+2x+5}$ and $f(x)=1/x$, illustrating that as $x$ gets large, the functions become indistinguishable.
\end{example}

Both the Direct and Limit Comparison Tests were given in terms of integrals over an infinite interval. There are versions that apply to improper integrals with an infinite range, but as they are a bit wordy and a little more difficult to employ, they are omitted from this text.\bigskip

This chapter has explored many integration techniques. We learned
% Substitution, which reverses the Chain Rule of differentiation, as well as
Integration by Parts, which reverses the Product Rule of differentiation. We also learned specialized techniques for handling trigonometric and rational functions. All techniques effectively have this goal in common: rewrite the integrand in a new way so that the integration step is easier to see and implement.

As stated before, integration is, in general, hard. It is easy to write a function whose antiderivative is impossible to write in terms of elementary functions, and even when a function does have an antiderivative expressible by elementary functions, it may be really hard to discover what it is. The powerful computer algebra system \textit{Mathematica}\textsuperscript{\textregistered} has approximately 1,000 pages of code dedicated to integration. 

Do not let this difficulty discourage you. There is great value in learning integration techniques, as they allow one to manipulate an integral in ways that can illuminate a concept for greater understanding. There is also great value in understanding the need for good numerical techniques: the Trapezoidal and Simpson's Rules are just the beginning of powerful techniques for approximating the value of integration.\bigskip

%The next chapter stresses the uses of integration. We generally do not find antiderivatives for antiderivative's sake, but rather because they provide the solution to some type of problem. The following chapter introduces us to several different problems whose solution is provided by integration.

\printexercises{exercises/06-07-exercises}

\section{Numerical Integration}\label{sec:numerical_integration}

The Fundamental Theorem of Calculus gives a concrete technique for finding the exact value of a definite integral. That technique is based on computing antiderivatives. Despite the power of this theorem, there are still situations where we must \emph{approximate} the value of the definite integral instead of finding its exact value. The first situation we explore is where we \emph{cannot} compute an antiderivative of the integrand. The second case is when we actually do not know the integrand, but only its value when evaluated at certain points.\index{integration!numerical}\index{numerical integration} %While we handle both situations in the same way, we address them separately here.
%
%\subsection{An Antiderivative That Cannot Be Computed}
\bigskip

%\autoref{sec:antider} introduced antiderivatives and the indefinite integral; given a function, its indefinite integral is a set of \emph{functions}. In \autoref{sec:def_int}, we learned about definite integrals: given a function and two bounds, the definite integral computes the ``area under the curve,'' i.e., a \emph{number}. The Fundamental Theorem of Calculus states that these two concepts are related: we use antiderivatives to evaluate definite integrals. Definite integrals have immense importance in mathematics, science and engineering. We will explore some of these applications in \autoref{SEVEN}. 
%
%This method of evaluating definite integrals has one significant shortcoming: not all functions have an antiderivative expressible in terms of elementary functions.  (Elementary functions are combinations of polynomials, $n^{\text{th}}$ root, rational, exponential, logarithmic and trigonometric functions.). To be clear, we are not claiming that the antiderivatives are just \emph{hard} to find; rather, we are saying that we \emph{cannot} write them out using functions that we normally use. 
An \textbf{elementary function}\index{elementary function} is any function that is a combination of polynomials, $n^{\text{th}}$ roots, rational, exponential, logarithmic and trigonometric functions and their inverses. We can compute the derivative of any elementary function, but there are many elementary functions of which we cannot compute an antiderivative. For example, the following functions do not have antiderivatives that we can express with elementary functions:
\[e^{-x^2}, \quad \sin(x^3)\quad \text{and} \quad \frac{\sin x}{x}.\]

The simplest way to refer to the antiderivatives of $e^{-x^2}$ is to simply write $\ds\int e^{-x^2}\dd x$. 

\mtable{Graphically representing three definite integrals that cannot be evaluated using antiderivatives.}{fig:numerical1}{\begin{tikzpicture}[alt={Blue curve y = e^(–x² on 0 ≤ x ≤ 1, starting at (0, 1) then falling smoothly to about (1, 0.4); region beneath is shaded.}]
\begin{axis}[width=1.16\marginparwidth,tick label style={font=\scriptsize},
axis y line=middle,axis x line=middle,name=myplot,axis on top,
ymin=-.2,ymax=1.2,xmin=-.2,xmax=1.1]
\addplot[draw={\colorone},fill={\coloronefill},area style,domain=0:1]
 {exp(-x*x)}\closedcycle;
\draw (axis cs:.75,1) node {\scriptsize $y=e^{-x^2}$};
\end{axis}
\node [right] at (myplot.right of origin) {\scriptsize $x$};
\node [above] at (myplot.above origin) {\scriptsize $y$};
\end{tikzpicture}
\bigskip\\
\begin{tikzpicture}[alt={Blue curve y=sin(x³); small negative arch left of the y-axis, large positive hump centered near x=1; shaded between curve and axis.}]
\begin{axis}[width=1.16\marginparwidth,tick label style={font=\scriptsize},
axis y line=middle,axis x line=middle,name=myplot,axis on top,
ymin=-.6,ymax=1.2,xmin=-1,xmax=1.7]
\addplot[draw={\colorone},fill={\coloronefill},area style,domain=-pi/4:pi/2]
 {sin(deg(x^3))}\closedcycle;
\draw (axis cs:.6,.9) node {\scriptsize $y=\sin(x^3)$};
\end{axis}
\node [right] at (myplot.right of origin) {\scriptsize $x$};
\node [above] at (myplot.above origin) {\scriptsize $y$};
\end{tikzpicture}
\bigskip\\
\begin{tikzpicture}[alt={Blue curve y = sin x / x on 0 ≤ x ≤ 15; initial tall peak at x=3, alternating smaller oscillations that decay toward the x-axis; entire area is shaded.}]
\begin{axis}[width=1.16\marginparwidth,tick label style={font=\scriptsize},
axis y line=middle,axis x line=middle,name=myplot,axis on top,xtick={5,10,15},
ymin=-.3,ymax=1.1,xmin=-.9,xmax=15.9]
\addplot[draw={\colorone},fill={\coloronefill},area style,domain=.5:4*pi]
 {sin(deg(x))/x}\closedcycle;
\draw (axis cs:10,.7) node {\scriptsize $\displaystyle y=\frac{\sin x}{x}$};
\end{axis}
\node [right] at (myplot.right of origin) {\scriptsize $x$};
\node [above] at (myplot.above origin) {\scriptsize $y$};
\end{tikzpicture}}

%How, then, can we solve problems involving definite integrals of such functions? We \emph{approximate.} 
This section outlines three common methods of approximating the value of definite integrals. We describe each as a systematic method of approximating area under a curve. By approximating this area accurately, we find an accurate approximation of the corresponding definite integral.

We will apply the methods we learn in this section to the following definite integrals:\vspace{-.5\baselineskip}
\[
 \int_0^1 e^{-x^2}\dd x, \quad
 \int_{-\frac{\pi}{4}}^{\frac{\pi}{2}} \sin(x^3)\dd x,
 \quad \text{and} \quad
 \int_{0.5}^{4\pi} \frac{\sin(x)}{x}\dd x,
\]
% todo Tim maybe use \pi/4 to \pi/2 in (b) instead?
as pictured in \autoref{fig:numerical1}.

\subsection{The Left and Right Hand Rule Methods}

In \autoref{sec:riemann} we addressed the problem of evaluating definite integrals by approximating the area under the curve using rectangles. We revisit those ideas here before introducing other methods of approximating definite integrals. \index{numerical integration!Left/Right Hand Rule}\index{Right Hand Rule}\index{Left Hand Rule}\index{integration!numerical!Left/Right Hand Rule}

We start with a review of notation. Let $f$ be a continuous function on the interval $[a,b]$. We wish to approximate $\ds \int_a^b f(x)\dd x$. We partition $[a,b]$ into $n$ equally spaced subintervals, each of length $\ds\Delta x = \frac{b-a}{n}$. The endpoints of these subintervals are labeled as
\[x_0=a,\ x_1 = a+\Delta x,\ x_2 = a+ 2\Delta x,\ \dotsc,\ x_i = a+i\Delta x,\ \dotsc,\ x_n = b.\]

%\autoref{idea:riemann}
\autoref{sec:riemann} showed that to use the Left Hand Rule we use the summation $\ds \sum_{i=1}^n f(x_{i-1})\Delta x$ and to use the Right Hand Rule we use $\ds \sum_{i=1}^n f(x_i)\Delta x$. We review the use of these rules in the context of examples.

\begin{example}[Approximating definite integrals with rectangles]\label{ex_num1}%
Approximate $\ds \int_0^1e^{-x^2}\dd x$ using the Left and Right Hand Rules with 5 equally spaced subintervals.
\solution
We begin by partitioning the interval $[0,1]$ into 5 equally spaced intervals. We have $\Delta x = \frac{1-0}5 = 1/5=0.2$, so
\[x_0=0,\ x_1=0.2,\ x_2=0.4,\ x_3=0.6,\ x_4=0.8,\text{ and }x_5=1.\]

\mtable{Approximating $\ds\int_0^1e^{-x^2}\dd x$ in \autoref{ex_num1} using (top) the left hand rule and (bottom) the right hand rule.}{fig:num1}{\begin{tikzpicture}[alt={Five uniform-width red rectangles from x = 0 to 1 whose left tops touch the decreasing curve e^(–x²); rectangles overshoot the curve, giving an over-estimate of area.}]
\begin{axis}[width=1.16\marginparwidth,tick label style={font=\scriptsize},
axis y line=middle,axis x line=middle,name=myplot,axis on top,xtick={.2,.4,.6,.8,1},
ymin=-.2,ymax=1.2,xmin=-.2,xmax=1.1]
\addplot[draw={\colorone},fill={\coloronefill},area style,domain=0:1]
 {exp(-x*x)}\closedcycle;
\foreach \x in {0,0.2,...,.8} {
 \edef\y{exp(-\x*\x)}
 \edef\temp{\noexpand\draw[thick,draw={\colortwo}]
  (axis cs:\x,0)rectangle(axis cs:{\x+.2},{\y});
 }\temp
}
\draw (axis cs:.75,1) node {\scriptsize $y=e^{-x^2}$};
\end{axis}
\node [right] at (myplot.right of origin) {\scriptsize $x$};
\node [above] at (myplot.above origin) {\scriptsize $y$};
\end{tikzpicture}
\bigskip\\
\begin{tikzpicture}[alt={five rectangles shifted so their right tops meet e^(–x²); rectangles lie fully below the curve, giving an under-estimate.}]
\begin{axis}[width=1.16\marginparwidth,tick label style={font=\scriptsize},
axis y line=middle,axis x line=middle,name=myplot,axis on top,xtick={.2,.4,.6,.8,1},
ymin=-.2,ymax=1.2,xmin=-.2,xmax=1.1]
\addplot[draw={\colorone},fill={\coloronefill},area style,domain=0:1]
 {exp(-x*x)}\closedcycle;
\foreach \x in {0.2,0.4,...,1} {
 \edef\y{exp(-\x*\x)}
 \edef\temp{\noexpand\draw[thick,draw={\colortwo}]
  (axis cs:\x,0)rectangle(axis cs:{\x-.2},{\y});
 }\temp
}
\draw (axis cs:.75,1) node {\scriptsize $y=e^{-x^2}$};
\end{axis}
\node [right] at (myplot.right of origin) {\scriptsize $x$};
\node [above] at (myplot.above origin) {\scriptsize $y$};
\end{tikzpicture}}

Using the Left Hand Rule, we have:
\begin{align*}
	\sum_{i=1}^n f(x_{i-1})\Delta x
	&= \bigl(f(x_0)+f(x_1)+f(x_2) + f(x_3) + f(x_4)\bigr)\Delta x \\
	&= \bigl(f(0) + f(0.2) + f(0.4) + f(0.6) + f(0.8)\bigr)\Delta x \\
	&\approx \bigl(1+0.961 + 0.852 + 0.698 + 0.527\bigr)(0.2)\\
	&\approx 0.808.
\end{align*}

Using the Right Hand Rule, we have:
\begin{align*}
	\sum_{i=1}^n f(x_i)\Delta x
	&= \bigl(f(x_1)+f(x_2) + f(x_3) + f(x_4) + f(x_5)\bigr)\Delta x \\
	&= \bigl(f(0.2) + f(0.4) + f(0.6) + f(0.8)+f(1)\bigr)\Delta x \\
	&\approx \bigl(0.961 +0.852 + 0.698 + 0.527 + 0.368\bigr)(0.2)\\
	&\approx 0.681.
\end{align*}

\autoref{fig:num1} shows the rectangles used in each method to approximate the definite integral. These graphs show that in this particular case, the Left Hand Rule is an over approximation and the Right Hand Rule is an under approximation. To get a better approximation, we could use more rectangles, as we did in \autoref{sec:riemann}. We could also average the Left and Right Hand Rule results together, giving
\[\frac{0.808 + 0.681}{2} = 0.7445.\]
The actual answer, accurate to 4 places after the decimal, is 0.7468, showing our average is a good approximation.
\end{example}

\mtable{Table of values used to approximate $\ds\int_{-\frac{\pi}4}^{\frac{\pi}2}\sin(x^3)\dd x$ in \autoref{ex_num2}.}{fig:num2a}{%
	\tagpdfsetup{table/header-rows={1}}
	\begin{tabular}{ l r l l }\lxBeginTableHead
	$x_i$ & Exact & Approx. & $\sin(x_i^3)$ \\\lxEndTableHead\midrule
	$x_0$ & $-\pi/4\phantom{0}$ & $-0.785$ & $-0.466$ \\
	$x_1$ & $-7\pi/40$ & $-0.550$ & $-0.165$ \\
	$x_2$ & $-\pi/10$ & $-0.314$ & $-0.031$ \\
	$x_3$ & $-\pi/40$ & $-0.0785$ & $\phantom{-}0$ \\
	$x_4$ & $\pi/20$ & $\phantom{-}0.157$ & $\phantom{-}0.004$ \\
	$x_5$ & $\pi/8\phantom{0}$ & $\phantom{-}0.393$ & $\phantom{-}0.061$ \\
	$x_6$ & $\pi/5\phantom{0}$ & $\phantom{-}0.628$ & $\phantom{-}0.246$ \\
	$x_7$ & $11\pi/40$ & $\phantom{-}0.864$ & $\phantom{-}0.601$ \\
	$x_8$ & $7\pi/20$ & $\phantom{-}1.10$ & $\phantom{-}0.971$ \\
	$x_9$ & $17\pi/40$ & $\phantom{-}1.34$ & $\phantom{-}0.690$ \\
	$x_{10}$ & $\pi/2\phantom{0}$ & $\phantom{-}1.57$ & $-0.670$
\end{tabular}}

\begin{example}[Approximating definite integrals with rectangles]\label{ex_num2}%
Approximate $\ds\int_{-\frac{\pi}4}^{\frac{\pi}2} \sin (x^3)\dd x$ using the Left and Right Hand Rules with 10 equally spaced subintervals.
\solution
We begin by finding $\Delta x$:
\[\frac{b-a}{n} = \frac{\pi/2 - (-\pi/4)}{10} = \frac{3\pi}{40}\approx 0.236.\]
It is useful to write out the endpoints of the subintervals in a table; in \autoref{fig:num2a}, we give the exact values of the endpoints, their decimal approximations, and decimal approximations of $\sin(x^3)$ evaluated at these points. 

\mtable{Approximating\\ $\ds\int_{-\frac{\pi}4}^{\frac{\pi}2}\sin(x^3)\dd x$ in \autoref{ex_num2} using (top) the left hand rule and (bottom) the right hand rule.}{fig:num2b}{%
\begin{tikzpicture}[alt={Blue curve sin(x³) with seven red rectangles whose left tops hit the curve; three negative region rectangles from left of the y-axis, others rise into the positive hump; shaded bars illustrate left-hand rule.}]
\begin{axis}[width=1.16\marginparwidth,tick label style={font=\scriptsize},
axis y line=middle,axis x line=middle,name=myplot,axis on top,
ymin=-.7,ymax=1.2,xmin=-1,xmax=1.7]
\addplot[draw={\colorone},fill={\coloronefill},area style,domain=-pi/4:pi/2]
 {sin(deg(x^3))}\closedcycle;
\foreach \x / \y in %
                {-0.785/ -0.466, -0.550/-0.165, -0.314/ -0.0310,%
                 -0.0785/-0.000484, 0.157/ 0.00388, 0.393/ 0.0605,
                 0.628/ 0.246, 0.864/ 0.601, 1.10/ 0.971, 1.34/ 0.690} {
 \addplot [thick,draw={\colortwo}] coordinates
  {(\x+.2356,0) (\x+.2356,\y) (\x,\y) (\x,0) (\x+.2356,0)};
}
\draw (axis cs:.6,.9) node {\scriptsize $y=\sin(x^3)$};
\end{axis}
\node [right] at (myplot.right of origin) {\scriptsize $x$};
\node [above] at (myplot.above origin) {\scriptsize $y$};
\end{tikzpicture}
\\
\begin{tikzpicture}[alt={Blue curve sin(x³) with seven red rectangles aligned to right endpoints; bars now sit under the hump and above the troughs, showing right-hand rule approximation.}]
\begin{axis}[width=1.16\marginparwidth,tick label style={font=\scriptsize},
axis y line=middle,axis x line=middle,name=myplot,axis on top,
ymin=-.7,ymax=1.2,xmin=-1,xmax=1.7]
\addplot[draw={\colorone},fill={\coloronefill},area style,domain=-pi/4:pi/2]
 {sin(deg(x^3))}\closedcycle;
\foreach \x / \y in %
               {-0.550/-0.165, -0.314/ -0.0310,%
                -0.0785/-0.000484, 0.157/ 0.00388, 0.393/ 0.0605,
                0.628/ 0.246, 0.864/ 0.601, 1.10/ 0.971, 1.34/ 0.690, 1.57/ -0.670} {
  \addplot [thick,draw={\colortwo}] coordinates
    {(\x-.2356,0) (\x-.2356,\y) (\x,\y) (\x,0) (\x-.2356,0)};
}
\draw (axis cs:.6,1.1) node {\scriptsize $y=\sin(x^3)$};
\end{axis}
\node [right] at (myplot.right of origin) {\scriptsize $x$};
\node [above] at (myplot.above origin) {\scriptsize $y$};
\end{tikzpicture}}

Once this table is created, it is straightforward to approximate the definite integral using the Left and Right Hand Rules. (Note: the table itself is easy to create, especially with a standard spreadsheet program on a computer. The last two columns are all that are needed.) The Left Hand Rule sums the first 10 values of $\sin(x_i^3)$ and multiplies the sum by $\Delta x$; the Right Hand Rule sums the last 10 values of $\sin(x_i^3)$ and multiplies by $\Delta x$. Therefore we have:
\begin{align*}
	\text{Left Hand Rule: }\int_{-\frac{\pi}4}^{\frac{\pi}2}\sin(x^3)\dd x
	& \approx (1.91)(0.236) = 0.451. \\
	\text{Right Hand Rule: }\int_{-\frac{\pi}4}^{\frac{\pi}2}\sin(x^3)\dd x
	&\approx (1.71)(0.236) = 0.404.
\end{align*}

The average of the Left and Right Hand Rules is 0.4275.  The actual answer, accurate to 3 places after the decimal, is 0.460. Our approximations were once again fairly good. The rectangles used in each approximation are shown in \autoref{fig:num2b}. It is clear from the graphs that using more rectangles (and hence, narrower rectangles) should result in a more accurate approximation.
\end{example}

\subsection{The Trapezoidal Rule}

In \autoref{ex_num1} we approximated the value of $\ds \int_0^1 e^{-x^2}\dd x$ with 5 rectangles of equal width. \autoref{fig:num1} showed the rectangles used in the Left and Right Hand Rules. These graphs clearly show that rectangles do not match the shape of the graph all that well, and that accurate approximations will only come by using lots of rectangles. \index{Trapezoidal Rule}\index{numerical integration!Trapezoidal Rule}\index{integration!numerical!Trapezoidal Rule}

\mtable{Approximating $\int_0^1 e^{-x^2}\dd x$ using 5 trapezoids of equal widths.}{fig:num3a}{\begin{tikzpicture}[alt={Blue y = e^(–x²) on 0 ≤ x ≤ 1; five equal-width red trapezoids under curve, shaded region below.}]
\begin{axis}[width=1.16\marginparwidth,tick label style={font=\scriptsize},
axis y line=middle,axis x line=middle,name=myplot,axis on top,xtick={.2,.4,...,1},
ymin=-.2,ymax=1.2,xmin=-.2,xmax=1.1]
\addplot[draw={\colorone},fill={\coloronefill},area style,domain=0:1]
 {exp(-x*x)}\closedcycle;
\foreach\x in{0,.2,...,.8}{
  \edef\y{{exp(-\x*\x)}}
  \edef\z{{exp(-(\x+.2)*(\x+.2))}}
  \edef\temp{\noexpand\addplot[thick,draw={\colortwo}] coordinates
   {(\x,0) (\x,\y) ({\x+.2},\z) ({\x+.2},0) (\x,0)};
 }\temp
}
\draw (axis cs:.75,1) node {\scriptsize $y=e^{-x^2}$};
\end{axis}
\node [right] at (myplot.right of origin) {\scriptsize $x$};
\node [above] at (myplot.above origin) {\scriptsize $y$};
\end{tikzpicture}}

Instead of using rectangles to approximate the area, we can instead use \emph{trapezoids.} In \autoref{fig:num3a}, we show the region under $f(x) = e^{-x^2}$ on $[0,1]$ approximated with 5 trapezoids of equal width; the top ``corners'' of each trapezoid lie on the graph of $f(x)$. It is clear from this figure that these trapezoids more accurately approximate the area under $f$ and hence should give a better approximation of $\int_0^1 e^{-x^2}\dd x$. (In fact, these trapezoids seem to give a \emph{great} approximation of the area.)

\youtubeVideo{8z6JRFvjkpc}{The Trapezoid Rule for Approximating Integrals}

\mtable{The area of a trapezoid is $\frac{a+b}2h$.}{fig:trapezoid}{\begin{tikzpicture}[alt={Right trapezoid sketch: parallel sides a and b, vertical height h}]
\draw(0,0)--node[pos=.5,left]{\small $a$}(0,1)--(1,1.5)
 --node[pos=.5,right]{\small$b$}(1,0)--node[pos=.5,below]{\small$h$}(0,0);
\draw(0,.1)--(.1,.1)--(.1,0);
\draw(.9,0)--(.9,.1)--(1,.1);
%\draw(2.5,.75)node{Area = };
\end{tikzpicture}}

The formula for the area of a trapezoid is given in \autoref{fig:trapezoid}. We approximate $\int_0^1 e^{-x^2}\dd x$ with these trapezoids in the following example.

\begin{example}[Approximating definite integrals using trapezoids]\label{ex_num3}%
Use 5 trapezoids of equal width to approximate $\ds \int_0^1e^{-x^2}\dd x$.
\solution
To compute the areas of the 5 trapezoids in \autoref{fig:num3a}, it will again be useful to create a table of values as shown in \autoref{fig:num3b}.

\mtable{A table of values of $e^{-x^2}$.}{fig:num3b}{%
\tagpdfsetup{table/header-rows={1}}
\begin{tabular}{ll}\lxBeginTableHead
$x_i$ & $e^{-x_i^2}$ \\\lxEndTableHead\midrule
0 & 1\\
0.2 & 0.961 \\
0.4 & 0.852 \\
0.6 & 0.698 \\
0.8 & 0.527 \\
1 & 0.368
\end{tabular}}

The leftmost trapezoid has legs of length 1 and 0.961 and a height of 0.2. Thus, by our formula, the area of the leftmost trapezoid is:
\[\frac{1+0.961}{2}(0.2) = 0.1961.\]
Moving right, the next trapezoid has legs of length 0.961 and 0.852 and a height of 0.2. Thus its area is:
\[\frac{0.961+0.852}2(0.2) = 0.1813.\]

The sum of the areas of all 5 trapezoids is:
\begin{align*}
\frac{1+0.961}{2}(0.2) + \frac{0.961+0.852}2(0.2)+\frac{0.852+0.698}2(0.2)& \\
+\frac{0.698+0.527}2(0.2)+\frac{0.527+0.368}2(0.2)&= 0.7445.
\end{align*}
We approximate $\ds \int_0^1 e^{-x^2}\dd x \approx 0.7445$.
\end{example}

There are many things to observe in this example. Note how each term in the final summation was multiplied by both 1/2 and by $\Delta x = 0.2$. We can factor these coefficients out, leaving a more concise summation as:
\begin{multline*}
\frac12(0.2)\Bigl[(1+0.961) + (0.961+0.852) +\\
 (0.852+0.698) + ( 0.698+ 0.527) +(0.527 + 0.368)\Bigr].
\end{multline*}
Now notice that all numbers except for the first and the last are added twice. Therefore we can write the summation even more concisely as
\[\frac{0.2}{2}\Bigl[1 + 2(0.961+0.852+0.698+0.527) + 0.368\Bigr].\]

This is the heart of the \textbf{Trapezoidal Rule}, where a definite integral $\ds \int_a^b f(x)\dd x$ is approximated by using trapezoids of equal widths to approximate the corresponding area under $f$. Using $n$ equally spaced subintervals with endpoints $x_0$, $x_1$, \dots, $x_n$, we again have $\ds \Delta x = \frac{b-a}n$. Thus:
\begin{align*}
	\int_a^b f(x)\dd x
	& \approx \sum_{i=1}^n \frac{f(x_{i-1})+f(x_i)}2\Delta x \\
	& = \frac{\Delta x}2 \sum_{i=1}^n \bigl(f(x_{i-1})+f(x_i)\bigr)\\
	& = \frac{\Delta x}2\left[f(x_0)+ 2\sum_{i=1}^{n-1} f(x_i) + f(x_n)\right].
\end{align*}

\begin{example}[Using the Trapezoidal Rule]\label{ex_num4}%
Revisit \autoref{ex_num2} and approximate $\ds\int_{-\frac{\pi}{4}}^{\frac{\pi}{2}} \sin (x^3)\dd x$ using the Trapezoidal Rule and 10 equally spaced subintervals.
\solution
We refer back to \autoref{fig:num2a} for the table of values of $\sin(x^3)$. Recall that $\Delta x = 3\pi/40 \approx 0.236$. Thus we have:\small
\begin{align*}
	\int_{-\frac{\pi}4}^{\frac{\pi}2} & \sin (x^3)\dd x \\
	&\approx \frac{0.236}{2}\Bigl[-0.466 + 2\Bigl(-0.165+(-0.031)+\dotsb+%0.971+
0.69\Bigr)+(-0.67)\Bigr]\\
	&= 0.4275.
\end{align*}
\end{example}

Notice how ``quickly'' the Trapezoidal Rule can be implemented once the table of values is created. This is true for all the methods explored in this section; the real work is creating a table of $x_i$ and $f(x_i)$ values. Once this is completed, approximating the definite integral is not difficult. Again, using technology is wise. Spreadsheets can make quick work of these computations and make using lots of subintervals easy. 

Also notice the  approximations the Trapezoidal Rule gives. It is  the average of the approximations given by the Left and Right Hand Rules! This effectively renders the Left and Right Hand Rules obsolete. They are useful when first learning about definite integrals, but if a real approximation is needed, one is generally better off using the Trapezoidal Rule instead of either the Left or Right Hand Rule.

We will also show that the Trapezoidal Rule makes using the Midpoint Rule obsolete as well.  With much more work, it will turn out that the Midpoint Rule has only a marginal gain in accuracy.  But we will include it in our results for the sake of completeness.\bigskip

How can we improve on the Trapezoidal Rule, apart from using more and more trapezoids? The answer is clear once we look back and consider what we have \emph{really} done so far. The Left Hand Rule is not \emph{really} about using rectangles to approximate area. Instead, it approximates a function $f$ with constant functions on small subintervals and then computes the definite integral of these constant functions. The Trapezoidal Rule is really approximating a function $f$ with a linear function on a small subinterval, then computes the definite integral of this linear function. In both of these cases the definite integrals are easy to compute in geometric terms.

So we have a progression: we start by approximating $f$ with a constant function and then with a linear function. What is next? A quadratic function. By approximating the curve of a function with lots of parabolas, we generally get an even better approximation of the definite integral. We call this process \textbf{Simpson's Rule}, named after Thomas Simpson (1710--1761), even though others had used this rule as much as 100 years prior.

\subsection{Simpson's Rule}

Given one point, we can create a constant function that goes through that point. Given two points, we can create a linear function that goes through those points. Given three points, we can create a quadratic function that goes through those three points (given that no two have the same $x$-value).\index{Simpson's Rule}\index{numerical integration!Simpson's Rule}\index{integration!numerical!Simpson's Rule}

Consider three points $(x_1,y_1)$, $(x_2,y_2)$ and $(x_3,y_3)$ whose $x$-values are equally spaced and $x_1<x_2<x_3$. Let $f$ be the quadratic function that goes through these three points. An exercise will ask you to show that 
\begin{equation}
\int_{x_1}^{x_3} f(x)\dd x = \frac{x_3-x_1}{6}\bigl(y_1+4y_2+y_3\bigr).\label{eq:simpsons}
\end{equation}

\mtable{A graph of a function $f$ and a parabola that approximates it well on $[1,3]$.}{fig:numsimpsons}{\begin{tikzpicture}[alt={Blue curve f from x = 1 to 3; red dashed parabola fits through endpoints and midpoint.}]
\begin{axis}[width=1.16\marginparwidth,tick label style={font=\scriptsize},
axis y line=middle,axis x line=middle,name=myplot,axis on top,xtick={1,2,3},
ytick={1,2,3},ymin=-.2,ymax=4,xmin=.5,xmax=3.5]
\addplot [draw={\colorone},domain=.9:3.1,thick] {11-12*x+4.5*x^2-.5*x^3};
\addplot [draw={\colortwo},dashed,domain=.9:3.1,thick] {8-6.5*x+1.5*x^2};
\filldraw (axis cs:1,3) circle (1pt);
\filldraw (axis cs:2,1) circle (1pt);
\filldraw (axis cs:3,2) circle (1pt);
\end{axis}
\node [right] at (myplot.right of origin) {\scriptsize $x$};
\node [above] at (myplot.above origin) {\scriptsize $y$};
\end{tikzpicture}}

Consider \autoref{fig:numsimpsons}. A function $f$ goes through the 3 points shown and the parabola $g$ that also goes through those points is graphed with a dashed line. Using our equation from above, we know exactly that
\[\int_1^3 g(x)\dd x = \frac{3-1}{6}\bigl(3+4(1)+2\bigr)= 3.\]
Since $g$ is a good approximation for $f$ on $[1,3]$, we can state that
\[\int_1^3 f(x)\dd x \approx 3.\]

Notice how the interval $[1,3]$ was split into two subintervals as we needed 3 points. Because of this, whenever we use Simpson's Rule, we need to break the interval into an even number of subintervals. 

In general, to approximate $\ds \int_a^b f(x)\dd x$ using Simpson's Rule, subdivide $[a,b]$ into $n$ subintervals, where $n$ is even and each subinterval has width $\Delta x = (b-a)/n$. We approximate $f$ with $n/2$ parabolic curves, using \autoeqref{eq:simpsons} to compute the area under these parabolas. Adding up these areas gives the formula:
\begin{multline*}
\int_a^b f(x)\dd x\approx\\
\frac{\Delta x}3
\left[f(x_0)+4f(x_1)+2f(x_2)+4f(x_3)+\dotsb+2f(x_{n-2})+4f(x_{n-1})+f(x_n)\right].
\end{multline*}
Note how the coefficients of the terms in the summation have the pattern 1, 4, 2, 4, 2, 4, \dots, 2, 4, 1.

Let's demonstrate Simpson's Rule with a concrete example.

\mtable{A table of values to approximate $\int_0^1e^{-x^2}\dd x$ in \autoref{ex_num5}, along with a graph of the function.}{fig:num5a}{%
\tagpdfsetup{table/header-rows={1}}
\begin{tabular}{ll}\lxBeginTableHead
$x_i$ & $e^{-x_i^2}$ \\\lxEndTableHead\midrule
0 & 1 \\
0.25 & 0.939 \\
0.5 & 0.779 \\
0.75 & 0.570 \\
1 & 0.368
\end{tabular}
\\ (a) \\
\begin{tikzpicture}[alt={Piecewise-linear red join of those points beneath blue y = e^(–x²); region shaded.}]
\begin{axis}[width=1.16\marginparwidth,tick label style={font=\scriptsize},
axis y line=middle,axis x line=middle,name=myplot,axis on top,xtick={.25,.5,.75,1},
ymin=-.2,ymax=1.2,xmin=-.2,xmax=1.1]
\addplot[draw={\colorone},fill={\coloronefill},area style,domain=0:1]
 {exp(-x*x)}\closedcycle;
\addplot [draw={\colortwo},domain=0:.5,thick] {1-0.042*x-0.8*x^2};
\addplot [draw={\colortwo},domain=.5:1,thick] {1.218-0.907*x+0.0569*x^2};
\addplot [thick,draw={\colortwo}] coordinates {(0,1) (0,0) (.5,0) (.5,.779)};
\addplot [thick,draw={\colortwo}] coordinates {(.5,0) (1,0) (1,.368)};
\draw (axis cs:.75,1) node {\scriptsize $y=e^{-x^2}$};
\filldraw (axis cs:0,1) circle (1pt)
          (axis cs:.25,.939) circle (1pt)
          (axis cs:.5,.779) circle (1pt)
          (axis cs:.75,.57) circle (1pt)
          (axis cs:1,.368) circle (1pt);
\end{axis}
\node [right] at (myplot.right of origin) {\scriptsize $x$};
\node [above] at (myplot.above origin) {\scriptsize $y$};
\end{tikzpicture}
\\ (b)}

\begin{example}[Using Simpson's Rule]\label{ex_num5}%
Approximate $\ds\int_0^1 e^{-x^2}\dd x$ using Simpson's Rule and 4 equally spaced subintervals.
\solution
We begin by making a table of values as we have in the past, as shown in \autoref{fig:num5a}(a). Simpson's Rule states that
\[
 \int_0^1e^{-x^2}\dd x
 \approx \frac{0.25}{3}\Bigl[1+4(0.939)+2(0.779)+4(0.570) + 0.368\Bigr]
 = 0.7468\overline{3}.
\]

Recall in \autoref{ex_num1} we stated that the correct answer, accurate to 4 places after the decimal, was 0.7468. Our approximation with Simpson's Rule, with 4 subintervals, is better than our approximation with the Trapezoidal Rule using 5.

\autoref{fig:num5a}(b) shows $f(x) = e^{-x^2}$ along with its approximating parabolas, demonstrating how good our approximation is. The approximating curves are nearly indistinguishable from the actual function.
\end{example}

\mtable{A table of values to approximate $\int_{-\frac{\pi}4}^{\frac{\pi}2}\sin(x^3)\dd x$ in \autoref{ex_num6}, along with a graph of the function.}{fig:num6a}{%
\tagpdfsetup{table/header-rows={1}}
\begin{tabular}{ll}\lxBeginTableHead
  $x_i$ & $\sin(x_i^3)$ \\\lxEndTableHead\midrule
  $-0.785$ & $-0.466$ \\
  $-0.550$ & $-0.165$   \\
  $-0.314$ & $-0.031$ \\
  $-0.0785$ & $\phantom{-}0$   \\
  $\phantom{-}0.157$ & $\phantom{-}0.004$ \\
  $\phantom{-}0.393$ & $\phantom{-}0.061$ \\
  $\phantom{-}0.628$ & $\phantom{-}0.246$ \\
  $\phantom{-}0.864$ & $\phantom{-}0.601$ \\
  $\phantom{-}1.10$ & $\phantom{-}0.971$ \\
  $\phantom{-}1.34$ & $\phantom{-}0.690$ \\
  $\phantom{-}1.57$ & $-0.670$ \\
\end{tabular}
\\ (a) \\
\begin{tikzpicture}[alt={Red polygon joins sample points under blue y=sin x³ on –π/4 < x < π/2; shaded lobe.}]
\begin{axis}[width=1.16\marginparwidth,tick label style={font=\scriptsize},
axis y line=middle,axis x line=middle,name=myplot,axis on top,
ymin=-.7,ymax=1.2,xmin=-1,xmax=1.7]
\addplot[draw={\colorone},fill={\coloronefill},area style,domain=-pi/4:pi/2]
 {sin(deg(x^3))}\closedcycle;
\addplot [draw={\colortwo},domain=-.785:-.314,thick] {-.11-.721*x-1.49*x^2};
\addplot [thick,draw={\colortwo}] coordinates {(-.785,-.466) (-.785,0) (-.314,0) (-.314,-.03)};
\addplot [draw={\colortwo},domain=-.314:.157,thick] {.004+.037*x-.236*x^2};
\addplot [thick,draw={\colortwo}] coordinates {(-.314,0) (.157,0) (.157,.004)};
\addplot [draw={\colortwo},domain=.157:.628,thick] {.037-.395*x+1.156*x^2};
\addplot [thick,draw={\colortwo}] coordinates {(.157,0) (.628,0) (.628,.246)};
\addplot [draw={\colortwo},domain=.628:1.1,thick] {-.632+1.32*x+.13*x^2};
\addplot [thick,draw={\colortwo}] coordinates {(.628,0) (1.1,0) (1.1,.971)};
\addplot [draw={\colortwo},domain=1.1:1.57,thick] {-11.98+22.46*x-9.72*x^2};
\addplot [thick,draw={\colortwo}] coordinates {(1.1,0) (1.57,0) (1.57,-.67)};
\filldraw (axis cs:-.785,-.466 ) circle (1pt)
          (axis cs:-.55,-.165 ) circle (1pt)
          (axis cs:-.314,-.031 ) circle (1pt)
          (axis cs:-.0785,0 ) circle (1pt)
          (axis cs:.157,.004 ) circle (1pt)
          (axis cs:.393,.061 ) circle (1pt)
          (axis cs:.628,.246 ) circle (1pt)
          (axis cs:.864,.601 ) circle (1pt)
          (axis cs:1.1,.971 ) circle (1pt)
          (axis cs:1.34,.69 ) circle (1pt)
          (axis cs:1.57,-.67 ) circle (1pt);
\draw (axis cs:.6,.9) node {\scriptsize $y=\sin(x^3)$};
\end{axis}
\node [right] at (myplot.right of origin) {\scriptsize $x$};
\node [above] at (myplot.above origin) {\scriptsize $y$};
\end{tikzpicture}
\\ (b)}%

\begin{example}[Using Simpson's Rule]\label{ex_num6}%
Approximate $\ds\int_{-\frac{\pi}4}^{\frac{\pi}2} \sin (x^3)\dd x$ using Simpson's Rule and 10 equally spaced intervals.
\solution
\autoref{fig:num6a}(a) shows the table of values that we used in the past for this problem, shown here again for convenience. Again, $\Delta x = (\pi/2+\pi/4)/10 \approx 0.236$.

Simpson's Rule states that
{\small\begin{align*}
	\int_{-\frac{\pi}4}^{\frac{\pi}2} \sin (x^3)\dd x
	&\approx \frac{0.236}3\Bigl[(-0.466)+4(-0.165)+2(-0.031) + \dotsb \\
	&\qquad\qquad\qquad \dotsb + 2(0.971) + 4(0.69) + (-0.67)\Bigr]\\
	&= 0.4701
\end{align*}}

Recall that the actual value, accurate to 3 decimal places, is 0.460. Our approximation is within $0.01$ of the correct value. The graph in \autoref{fig:num6a}(b) shows how closely the parabolas match the shape of the graph.
\end{example}

%\noindent\begin{minipage}[t]{\linewidth}\tagstructbegin{tag=Sect}\noindent%
%\captionsetup{type=figure}%
%
%\caption{Approximating $\int_{-\frac{\pi}4}^{\frac{\pi}2}\sin(x^3)\dd x$ in \autoref{ex_num6} with Simpson's Rule and 10 equally spaced intervals.}
%\label{fig:num6b}\tagstructend
%\end{minipage}

\subsection{Summary and Error Analysis}

We summarize the key concepts of this section thus far in the following Key Idea.

{
\tcbset{grow to right by=8em}
\begin{keyidea}[Numerical Integration]\label{idea:numerical}%
Let $f$ be a continuous function on $[a,b]$, let $n$ be a positive integer, and let $\ds\Delta x = \frac{b-a}{n}$.
Set $x_0=a$, $x_1=a+\Delta x$, \dots, $x_i = a+i\Delta x$, $x_n=b$.
Then $\ds\int_a^b f(x)\dd x$ can be approximated by:\\
\index{integration!numerical!Left Hand Rule}%
\index{integration!numerical!Right Hand Rule}%
\index{integration!numerical!Midpoint Rule}%
\index{integration!numerical!Trapezoidal Rule}%
\index{integration!numerical!Simpson's Rule}%
\index{Left/Right Hand Rule}\index{Midpoint Rule}%
\index{Trapezoidal Rule}\index{Simpson's Rule}%
\index{numerical integration!Left Hand Rule}%
\index{numerical integration!Right Hand Rule}%
\index{numerical integration!Midpoint Rule}%
\index{numerical integration!Trapezoidal Rule}%
\index{numerical integration!Simpson's Rule}%
\tagpdfsetup{table/header-columns={1},table/header-rows={}}
\begin{tabular}{ll}
Left Hand Rule: &
%$\ds\int_a^b f(x)\dd x \approx 
$\Delta x\Bigl[f(x_0) + f(x_1) + \dotsb + f(x_{n-1})\Bigr]$.\\\addlinespace
Right Hand Rule: &
%$\ds\int_a^b f(x)\dd x \approx 
$\Delta x\Bigl[f(x_1) + f(x_2) + \dotsb + f(x_n)\Bigr]$.\\\addlinespace
Midpoint Rule: &
%$\ds\int_a^b f(x)\dd x \approx 
$\Delta x\Bigl[f\bigl(\frac{x_0+x_1}2\bigr) + f\bigl(\frac{x_1+x_2}2\bigr) + \dotsb + f\bigl(\frac{x_{n-1}+x_n}2\bigr)\Bigr]$.\\\addlinespace
Trapezoidal Rule: &
%$\ds\int_a^b f(x)\dd x \approx 
$\frac{\Delta x}2\Bigl[f(x_0) + 2f(x_1) + 2f(x_2) +\dotsb + 2f(x_{n-1})+ f(x_n)\Bigr]$.\\\addlinespace
Simpson's Rule: &
%$\ds\int_a^b f(x)\dd x \approx 
$\frac{\Delta x}3\Bigl[f(x_0) + 4f(x_1) + 2f(x_2) +\dotsb + 4f(x_{n-1})+ f(x_n)\Bigr]$ {\small ($n$ even)}.
\end{tabular}
\end{keyidea}
}

In our examples, we approximated the value of a definite integral using a given method then compared it to the ``right'' answer. This should have raised several questions in the reader's mind, such as:
\begin{enumerate}
	\item	How was the ``right'' answer computed?
	\item	If the right answer can be found, what is the point of approximating?
	\item	If there is value to approximating, how are we supposed to know if the approximation is any good?
\end{enumerate}

These are good questions, and their answers are educational. In the examples, \emph{the} right answer was never computed. Rather, an approximation accurate to a certain number of places after the decimal was given. In \autoref{ex_num1}, we do not know the \emph{exact} answer, but we know it starts with 0.7468. These more accurate approximations were computed using numerical integration but with more precision (i.e., more subintervals and the help of a computer). 

Since the exact answer cannot be found, approximation still has its place. How are we to tell if the approximation is any good?

``Trial and error'' provides one way. %(No, this does not refer to seeing whether or not the bridge collapses.) 
Using technology, make an approximation with, say, 10, 100, and 200 subintervals. This likely will not take much time at all, and a trend should emerge. If a trend does not emerge, try using yet more subintervals. Keep in mind that trial and error is never foolproof; you might stumble upon a problem in which a trend will not emerge.

A second method is to use Error Analysis. While the details are beyond the scope of this text, there are some formulas that give \emph{bounds} for how good your approximation will be. For instance, the formula might state that the approximation is within 0.1 of the correct answer. If the approximation is 1.58, then one knows that the correct answer is between 1.48 and 1.68. By using lots of subintervals, one can get an approximation as accurate as one likes. \autoref{thm:numerical_error} states what these bounds are.

% can't use: L, R, M, T, S, n, a, b
\begin{theorem}[Error Bounds in Numerical Integration]\label{thm:numerical_error}%
Suppose that $K_m$ is an upper bound on $\abs{f^{(m)}(x)}$ on $[a,b]$.  Then a bound for the error of the numerical method of integration is given by:
\index{integration!numerical!Left Hand Rule}
\index{integration!numerical!Right Hand Rule}
\index{integration!numerical!Midpoint Rule}
\index{integration!numerical!Trapezoidal Rule}
\index{integration!numerical!Simpson's Rule}
\index{Left/Right Hand Rule!error bounds}
\index{Midpoint Rule!error bounds}
\index{Trapezoidal Rule!error bounds}
\index{Simpson's Rule!error bounds}
\index{numerical integration!Left Hand Rule!error bounds}
\index{numerical integration!Right Hand Rule!error bounds}
\index{numerical integration!Midpoint Rule!error bounds}
\index{numerical integration!Trapezoidal Rule!error bounds}
\index{numerical integration!Simpson's Rule!error bounds}
\begin{center}
\tagpdfsetup{table/header-rows={1}}
\begin{tabular}{lc}
Method & Error Bound \\\midrule
Left/Right Hand Rule & $\dfrac{K_1(b-a)^2}{2n}$ \\\addlinespace
Midpoint Rule & $\dfrac{K_2(b-a)^3}{24n^2}$ \\\addlinespace
Trapezoidal Rule & $\dfrac{K_2(b-a)^3}{12n^2}$ \\\addlinespace
Simpson's Rule & $\dfrac{K_4(b-a)^5}{180n^4}$
\end{tabular}
\end{center}
\end{theorem}

There are some key things to note about this theorem.
\begin{enumerate}
	\item	The larger the interval, the larger the error. This should make sense intuitively.
	\item	The error shrinks as more subintervals are used (i.e., as $n$ gets larger).
	\item	When $n$ doubles, the Left and Right Hand Rules double in accuracy, the Midpoint and Trapezoidal Rules quadruple in accuracy, and Simpson's Rule is 16 times more accurate.
	\item	The error in Simpson's Rule has a term relating to the 4$^{\text{th}}$ derivative of $f$. Consider a cubic polynomial: its $4^{\text{th}}$ derivative is 0. Therefore, the error in approximating the definite integral of a cubic polynomial with Simpson's Rule is 0 --- Simpson's Rule computes the exact answer!
\end{enumerate}

We revisit Examples \ref{ex_num3} and \ref{ex_num5} and compute the error bounds using \autoref{thm:numerical_error} in the following example.

\begin{example}[Computing error bounds]\label{ex_num7}%
Find the error bounds when approximating $\ds \int_0^1 e^{-x^2}\dd x$ using the Trapezoidal Rule and 5 subintervals, and using Simpson's Rule with 4 subintervals.
\solution
\textbf{Trapezoidal Rule with $n=5$:}

We start by computing the 2\textsuperscript{nd} derivative of $f(x) = e^{-x^2}$:
%
\mtable{Graphing $\fpp(x)$ in \autoref{ex_num7} to help establish error bounds.}{fig:num7a}{\begin{tikzpicture}[alt={Blue y = e^(-x²)(4x²–2) on 0 < x < 1; curve starts near –2, rises through 0 around 0.7, ends about 1.7; axes x,y.}]
\begin{axis}[width=1.16\marginparwidth,tick label style={font=\scriptsize},
axis y line=middle,axis x line=middle,name=myplot,axis on top,
ymin=-2.1,ymax=.7,xmin=-.1,xmax=1.1]
\addplot[draw={\colorone},smooth,thick,domain=0:1] {exp(-x*x)*(4*x*x-2)};
\draw (axis cs:.8,-1.5) node {\scriptsize $y=e^{-x^2}(4x^2-2)$};
\end{axis}
\node [right] at (myplot.right of origin) {\scriptsize $x$};
\node [above] at (myplot.above origin) {\scriptsize $y$};
\end{tikzpicture}}
%
\[\fpp(x) = e^{-x^2}(4x^2-2).\]
\autoref{fig:num7a} shows a graph of $\fpp(x)$ on $[0,1]$. It is clear that the largest value of $\fpp$, in absolute value, is 2. Thus we let $M=2$ and apply the error formula from \autoref{thm:numerical_error}.
\[E_T = \frac{(1-0)^3}{12\cdot 5^2}\cdot 2 = 0.00\overline{6}.\]

Our error estimation formula states that our approximation of 0.7445 found in \autoref{ex_num3} is within 0.0067 of the correct answer, hence we know that
\[0.7445-0.0067 = .7378 \leq \int_0^1e^{-x^2}\dd x \leq 0.7512 = 0.7445 + 0.0067.\]
We had earlier computed the exact answer, correct to 4 decimal places, to be 0.7468, affirming the validity of \autoref{thm:numerical_error}.\bigskip\\
\textbf{Simpson's Rule with $n=4$:}

We start by computing the 4\textsuperscript{th} derivative of $f(x) = e^{-x^2}$:
%
\mtable{Graphing $f\,^{(4)}(x)$ in \autoref{ex_num7} to help establish error bounds.}{fig:num7b}{\begin{tikzpicture}[alt={Blue y = e^(-x²)(16x⁴–48x²+12) on 0 < x < 1; peak (0, 12), crosses 0 near 0.8, falls to –8; axes x,y.}]
\begin{axis}[width=1.16\marginparwidth,tick label style={font=\scriptsize},
axis y line=middle,axis x line=middle,name=myplot,axis on top,minor y tick num=4,
ymin=-8,ymax=13.5,xmin=-.1,xmax=1.1]
\addplot[draw={\colorone},smooth,thick,domain=0:1] {exp(-x*x)*(16*x^4-48*x*x+12)};
\draw (axis cs:.7,11) node {\scriptsize $y=e^{-x^2}(16x^4-48x^2+12)$};
\end{axis}
\node [right] at (myplot.right of origin) {\scriptsize $x$};
\node [above] at (myplot.above origin) {\scriptsize $y$};
\end{tikzpicture}}
%
\[f\,^{(4)}(x) = e^{-x^2}(16x^4-48x^2+12).\]
\autoref{fig:num7b} shows a graph of $f\,^{(4)}(x)$ on $[0,1]$. It is clear that the largest value of $f\,^{(4)}$, in absolute value, is 12. Thus we let $M=12$ and apply the error formula from \autoref{thm:numerical_error}.
\[E_s = \frac{(1-0)^5}{180\cdot 4^4}\cdot 12 = 0.00026.\]

Our error estimation formula states that our approximation of $0.7468\overline{3}$ found in \autoref{ex_num5} is within 0.00026 of the correct answer, hence we know that
\[0.74683-0.00026 = .74657 \leq \int_0^1e^{-x^2}\dd x \leq 0.74709 = 0.74683 + 0.00026.\]
Once again we affirm the validity of \autoref{thm:numerical_error}.
\end{example}

We have seen that, for $\ds\int_0^1 e^{-x^2}\dd x$, Simpson's Rule with 4 subintervals is far more accurate than the Trapezoidal Rule with 5 subintervals.  We now find how many intervals we would need to match that accuracy.

\begin{example}[Finding a number subintervals]\label{ex_num_find_num_sub_int}%
Find the number of subintervals necessary to estimate $\ds\int_0^1 e^{-x^2}\dd x$ to within $0.00026$ using the Trapezoidal Rule.
\solution
We can again use that $\fpp(x)$ is bounded by 2, so that
\[E_T=\frac{(1-0)^3}{12\cdot n^2}\cdot 2=\frac1{6n^2}.\]
In order for this to be at most $0.00026$, we need to have
\[n\ge\frac1{\sqrt{6\cdot0.00026}}\approx25.3.\]
Therefore, we will need at least 26 subintervals in order to have as much accuracy as Simpson's Rule with 4 subintervals.
\end{example}

At the beginning of this section we mentioned two main situations where numerical integration was desirable. We have considered the case where an antiderivative of the integrand cannot be computed. We now investigate the situation where the integrand is not known. This is, in fact, the most widely used application of Numerical Integration methods. ``Most of the time'' we observe behavior but do not know ``the'' function that describes it. We instead collect data about the behavior and make approximations based off of this data. We demonstrate this in an example.

\begin{example}[Approximating distance traveled]\label{ex_num8}%
One of the authors drove his daughter home from school while she recorded their speed every 30 seconds. The data is given in \autoref{fig:num8}.
%
\mtable{Speed data collected at 30 second intervals for \autoref{ex_num8}.}{fig:num8}{%
\tagpdfsetup{table/header-rows={1}}
\begin{tabular}{rr}\lxBeginTableHead
Time & Speed\\(min)&(mph)
 \\\lxEndTableHead\midrule
0\phantom{.5} & 0 \\
 0.5 & 25 \\ 1\phantom{.5} & 22 \\
 1.5 & 19 \\ 2\phantom{.5} & 39 \\
 2.5 & 0 \\ 3\phantom{.5} & 43\\
 3.5 & 59 \\ 4\phantom{.5} & 54 \\
 4.5 & 51 \\ 5\phantom{.5} & 43 \\
 5.5 & 35 \\ 6\phantom{.5} & 40 \\
 6.5 & 43 \\ 7\phantom{.5} & 30 \\
 7.5 & 0 \\ 8\phantom{.5} & 0 \\
 8.5 & 28 \\ 9\phantom{.5} & 40 \\
 9.5 & 42 \\ 10\phantom{.5} & 40 \\
 10.5 & 39 \\ 11\phantom{.5} & 40 \\
 11.5 & 23 \\ 12\phantom{.5} & 0 \\ 
\end{tabular}}
%
Approximate the distance they traveled.
\solution
Recall that by integrating a speed function we get distance traveled. We have information about $v(t)$; we will use Simpson's Rule to approximate $\ds \int_a^b v(t)\dd t.$ 

The most difficult aspect of this problem is converting the given data into the form we need it to be in. The speed is measured in miles per hour, whereas the time is measured in 30 second increments.

We need to compute $\Delta x = (b-a)/n$. Clearly, $n=24$. What are $a$ and $b$? Since we start at time $t=0$, we have that $a=0$. The final recorded time came after 24 periods of 30 seconds, which is 12 minutes or 1/5 of an hour. Thus we have
\[\Delta x = \frac{b-a}{n} = \frac{1/5-0}{24} = \frac1{120}; \quad \frac{\Delta x}{3} = \frac{1}{360}.\]

Thus the distance traveled is approximately:
\begin{align*}
\int_0^{0.2}v(t)\dd t &\approx \frac{1}{360}\Bigl[f(x_1)+4f(x_2) + 2f(x_3) + \dots + 4f(x_n)+f(x_{n+1})\Bigr]\\
					&= \frac{1}{360}\Bigl[0+4\cdot25+2\cdot 22 + \dots + 2\cdot40+4\cdot 23 + 0\Bigr]\\
					&\approx 6.2167 \ \text{miles.}
\end{align*}
We approximate the author drove 6.2 miles. (Because we are sure the reader wants to know, the author's odometer recorded the distance as about 6.05 miles.)
\end{example}

%We started this chapter learning about antiderivatives and indefinite integrals. We then seemed to change focus by looking at areas between the graph of a function and the $x$-axis. We defined these areas as the definite integral of the function, using a notation very similar to the notation of the indefinite integral. The Fundamental Theorem of Calculus tied these two seemingly separate concepts together: we can find areas under a curve, i.e., we can evaluate a definite integral, using antiderivatives. 
%
%We ended the chapter by noting that antiderivatives are sometimes more than difficult to find: they are impossible. Therefore we developed numerical techniques that gave us good approximations of definite integrals.
%
%We used the definite integral to compute areas, and also to compute displacements and distances traveled. There is far more we can do than that. In \autoref{chapter:app_of_int} we'll see more applications of the definite integral. Before that, in \autoref{chapter:anti_tech} we'll learn advanced techniques of integration, analogous to learning rules like the Product, Quotient and Chain Rules of differentiation.

\printexercises{exercises/05-05-exercises}

\iftoggle{isEarlyTrans}{%
 \section{Hyperbolic Functions}\label{sec:hyperbolic}

The \textbf{hyperbolic functions} are functions that have many applications to mathematics, physics, and engineering. Among many other applications, they are used to describe the formation of satellite rings around planets, to describe the shape of a rope hanging from two points, and have application to the theory of special relativity. This section defines the hyperbolic functions and describes many of their properties, especially their usefulness to calculus.

\mtable{Using trigonometric functions to define points on a circle and hyperbolic functions to define points on a hyperbola.}{fig:hfcircle}{\begin{tikzpicture}[alt={The unit circle, with a shaded sector in the first quadrant of angle θ and area θ/2.}]
 \begin{axis}[height=\marginparwidth,width=\marginparwidth,
   tick label style={font=\scriptsize},
   axis y line=middle,axis x line=middle,name=myplot,axis on top,
   xtick={-1,1},ytick={-1,1},ymin=-1.1,ymax=1.1,xmin=-1.1,xmax=1.1]
  \addplot [draw={\coloronefill},fill={\coloronefill},domain=0:45] ({cos(x)},{sin(x)})
   -- (axis cs:0,0)--cycle;
  \addplot [draw={\colorone},domain=0:360,thick,smooth,samples=40] ({cos(x)},{sin(x)});
  \filldraw (axis cs:.707,.707) circle (1pt) node [left]%[shift={(4pt,11pt)}]
   {\scriptsize ($\cos \theta$,$\sin \theta$)};
  \draw (axis cs:.6,.25) node {\scriptsize $A=\dfrac{\theta}{2}$};
  \draw (axis cs:-.5,.2) node {\scriptsize $x^2+y^2=1$};
 \end{axis}
% \draw[draw={\colorone}](0,0)circle(1);
 \node [right] at (myplot.right of origin) {\scriptsize $x$};
 \node [above] at (myplot.above origin) {\scriptsize $y$};
\end{tikzpicture}
\vspace{10pt}
\begin{tikzpicture}[alt={The unit hyperbola.  The region in the first quadrant left of the right half and to the right of the segment connecting the origin to (coshθ,sinhθ) has been shaded.}]
 \begin{axis}[height=\marginparwidth,width=\marginparwidth,
   tick label style={font=\scriptsize},
   axis y line=middle,axis x line=middle,name=myplot,axis on top,
   ymin=-3.1,ymax=3.1,xmin=-3.1,xmax=3.1]
  \addplot [draw={\coloronefill},fill={\coloronefill},domain=0:1.6] ({cosh(x)},{sinh(x)})
   -- (axis cs:0,0)--cycle;
  \addplot [draw={\colorone},domain=-2:2,thick,smooth] ({cosh(x)},{sinh(x)});
  \addplot [draw={\colorone},domain=-2:2,thick,smooth] ({-cosh(x)},{sinh(x)});
  \filldraw (axis cs:2.577,2.376) circle (1pt) node [left]
   {\scriptsize ($\cosh \theta$,$\sinh \theta$)};
  \draw (axis cs:2,.6) node {\scriptsize $A=\dfrac{\theta}{2}$};
  \draw (axis cs:-1.75,2.75) node {\scriptsize $x^2-y^2=1$};
 \end{axis}
 \node [right] at (myplot.right of origin) {\scriptsize $x$};
 \node [above] at (myplot.above origin) {\scriptsize $y$};
\end{tikzpicture}}

These functions are sometimes referred to as the ``hyperbolic trigonometric functions'' as there are many connections between them and the standard trigonometric functions. \autoref{fig:hfcircle} demonstrates one such connection. Just as cosine and sine are used to define points on the circle defined by $x^2+y^2=1$, the functions \textbf{hyperbolic cosine} and \textbf{hyperbolic sine} are used to define points on the hyperbola $x^2-y^2=1$.

We begin with their definitions.

\begin{definition}[Hyperbolic Functions]\label{def:hyperbolic_functions}%
\mbox{}\\[-2\baselineskip]\parbox[t]{\linewidth}{%
\begin{multicols}{2}\index{hyperbolic function!definition}%
\begin{enumerate}\lxAddClass{columns2}
\item		$\ds \cosh x = \frac{e^x+e^{-x}}2$
\item		$\ds \sinh x = \frac{e^x-e^{-x}}2$
\item		$\ds \tanh x = \frac{\sinh x}{\cosh x}$
\item		$\ds \sech x = \frac{1}{\cosh x}$
\item		$\ds \csch x = \frac{1}{\sinh x}$
\item		$\ds \coth x = \frac{\cosh x}{\sinh x}$
\end{enumerate}
\end{multicols}}
\end{definition}

The hyperbolic functions are graphed in \autoref{fig:hyperbolic}. In the graphs of $\cosh x$ and $\sinh x$, graphs of $e^x/2$ and $e^{-x}/2$ are included with dashed lines. As $x$ gets ``large,'' $\cosh x$ and $\sinh x$ each act like $e^x/2$; when $x$ is a large negative number, $\cosh x$ acts like $e^{-x}/2$ whereas $\sinh x$ acts like $-e^{-x}/2$.

\mnote{\textbf{Pronunciation Note:} \\
``cosh'' rhymes with ``gosh,'' \\
``sinh'' rhymes with ``pinch,'' and\\
``tanh'' rhymes with ``ranch,'' \\
%``sech'' rhymes with ``fetch,'' and \\
%``coth'' rhymes with ``moth.''
}

Notice the domains of $\tanh x$ and $\sech x$ are $(-\infty,\infty)$, whereas both $\coth x$ and $\csch x$ have vertical asymptotes at $x=0$. Also note the ranges of these functions, especially $\tanh x$: as $x\to\infty$, both $\sinh x$ and $\cosh x$ approach $e^x/2$, hence $\tanh x$ approaches $1$.

\noindent
\begin{minipage}{\linewidth}\tagstructbegin{tag=Sect}
\centering
\begin{tikzpicture}[alt={cosh is in blue, and makes a U shape above the x axis.  For positive x, cosh is well approximated by exp(x/2).  For negative x, cosh is well approximated by exp(-x/2).}]
\begin{axis}[width=.55\textwidth,tick label style={font=\scriptsize},
axis y line=middle,axis x line=middle,name=myplot,axis on top,ymin=-10.9,ymax=10.9,
xmin=-3.5,xmax=3.5,scaled ticks=false]
\addplot [draw={\colorone},thick,smooth,domain=-3:3] {cosh(x)};
\draw (axis cs:2,-6) node {\scriptsize $f(x)=\cosh x$};
\addplot [draw={\colortwo},smooth,thick,dashed,domain=-3:3] {exp(x)/2};
\addplot [draw={\colortwo},smooth,thick,dashed,domain=-3:3] {exp(-x)/2};
\end{axis}
\node [right] at (myplot.right of origin) {\scriptsize $x$};
\node [above] at (myplot.above origin) {\scriptsize $y$};
\end{tikzpicture}%
\qquad
\begin{tikzpicture}[alt={sinh is in blue, and looks vaguely like y=x*x*x.  For positive x, sinh is well approximated by exp(x/2).  For negative x, sinh is well approximated by -exp(-x/2).}]
\begin{axis}[width=.55\textwidth,tick label style={font=\scriptsize},
axis y line=middle,axis x line=middle,name=myplot,axis on top,ymin=-10.9,ymax=10.9,
xmin=-3.5,xmax=3.5,scaled ticks=false]
\addplot [draw={\colorone},thick,smooth,domain=-3:3] {sinh(x)};
\draw (axis cs:2,-6) node {\scriptsize $f(x)=\sinh x$};
\addplot [draw={\colortwo},smooth,thick,dashed,domain=-3:3] {exp(x)/2};
\addplot [draw={\colortwo},smooth,thick,dashed,domain=-3:3] {exp(-x)/2};
\end{axis}
\node [right] at (myplot.right of origin) {\scriptsize $x$};
\node [above] at (myplot.above origin) {\scriptsize $y$};
\end{tikzpicture}
\\[20pt]
\begin{tikzpicture}[alt={tahn is in blue, has a horizontal asymptote of y=±1 as x approaches ±∞, and behaves like y=x near the origin.  coth is in red, and looks like 1+1/x for positive x and -1+1/x for negative x.}]
\begin{axis}[width=.55\textwidth,tick label style={font=\scriptsize},
axis y line=middle,axis x line=middle,name=myplot,axis on top,ytick={-2,2},
ymin=-3.5,ymax=3.5,xmin=-3.5,xmax=3.5,scaled ticks=false]
\addplot [draw={\colorone},thick,smooth,domain=-3:3] {tanh(x)};
\draw (axis cs:1.7,-1.5) node {\scriptsize $f(x)=\tanh x$};
\draw (axis cs:2.2,2) node {\scriptsize $f(x)=\coth x$};
\draw [->,>=stealth] (axis cs:1,-1) -- (axis cs:.6,.45);
\addplot [draw={\colortwo},smooth,thick,domain=-3:-.1] {1/tanh(x)};
\addplot [draw={\colortwo},smooth,thick,domain=.1:3] {1/tanh(x)};
\draw [loosely dashed] (axis cs:-3,1) -- (axis cs:3,1);
\draw [loosely dashed] (axis cs:-3,-1) -- (axis cs:3,-1);
\end{axis}
\node [right] at (myplot.right of origin) {\scriptsize $x$};
\node [above] at (myplot.above origin) {\scriptsize $y$};
\end{tikzpicture}%
\qquad
\begin{tikzpicture}[alt={sech is in blue, has a horizontal asymptote of y=0 as x approaches ±∞, and behaves like 1-x*x/2 near the origin.  csch is in red, and looks like y=1/x.}]
\begin{axis}[width=.55\textwidth,tick label style={font=\scriptsize},
axis y line=middle,axis x line=middle,name=myplot,axis on top,ytick={-3,-2,-1,1,2,3},
ymin=-3.5,ymax=3.5,xmin=-3.5,xmax=3.5,scaled ticks=false]
\addplot [draw={\colorone},thick,smooth,domain=-3:3] {1/cosh(x)};
\draw (axis cs:-2,1.5) node {\scriptsize $f(x)=\sech x$};
\draw (axis cs:2.2,2) node {\scriptsize $f(x)=\csch x$};
\addplot [draw={\colortwo},smooth,thick,domain=-3:-.1] {1/sinh(x)};
\addplot [draw={\colortwo},smooth,thick,domain=.1:3] {1/sinh(x)};
\end{axis}
\node [right] at (myplot.right of origin) {\scriptsize $x$};
\node [above] at (myplot.above origin) {\scriptsize $y$};
\end{tikzpicture}
\captionsetup{type=figure}
\caption{Graphs of the hyperbolic functions.}
\label{fig:hyperbolic}\tagstructend
\end{minipage}

\youtubeVideo{G1C1Z5aTZSQ}{Hyperbolic Functions --- The Basics}

%It is no coincidence that these functions share a name similar to the trigonometric functions. 
The following example explores some of the properties of these functions that bear remarkable resemblance to the properties of their trigonometric counterparts.

\begin{example}[Exploring properties of hyperbolic functions]\label{ex_hf1}%
Use \autoref{def:hyperbolic_functions} to rewrite the following expressions.
\begin{multicols}{2}
\begin{enumerate}\lxAddClass{columns2}
\item		$\cosh^2 x-\sinh^2x$
\item		$\tanh^2 x+\sech^2 x$
\item		$2\cosh x\sinh x$
\item		$\frac{\dd}{\dd x}\bigl(\cosh x\bigr)$
\item		$\frac{\dd}{\dd x}\bigl(\sinh x\bigr)$
\item		$\frac{\dd}{\dd x}\bigl(\tanh x\bigr)$
\end{enumerate}
\end{multicols}
\solution
\begin{enumerate}
\item \hfill$\begin{aligned}[t]
 \cosh^2x-\sinh^2x
 &= \left(\frac{e^x+e^{-x}}2\right)^2 -\left(\frac{e^x-e^{-x}}2\right)^2\\
 &= \frac{e^{2x}+2e^xe^{-x} + e^{-2x}}4 - \frac{e^{2x}-2e^xe^{-x} + e^{-2x}}4\\
 &= \frac44=1.
\end{aligned}$\hfill\null\\
So $\cosh^2 x-\sinh^2x=1$.

\item \hfill$\begin{aligned}[t]
 \tanh^2 x+\sech^2 x
 &=\frac{\sinh^2x}{\cosh^2 x} + \frac{1}{\cosh^2 x} \\
 &= \frac{\sinh^2x+1}{\cosh^2 x}\qquad \text{\small Now use identity from \#1.}\\
 &= \frac{\cosh^2 x}{\cosh^2 x} = 1.
\end{aligned}$\hfill\null\\
So $\tanh^2 x+\sech^2 x=1$.

\item \hfill$\begin{aligned}[t]
	2\cosh x\sinh x
	&= 2\left(\frac{e^x+e^{-x}}2\right)\left(\frac{e^x-e^{-x}}2\right) \\
	&= 2 \cdot\frac{e^{2x} - e^{-2x}}4\\
	&= \frac{e^{2x} - e^{-2x}}2 = \sinh (2x).\\
\end{aligned}$\hfill\null\\
Thus $2\cosh x\sinh x = \sinh (2x)$.

\item \hfill$\begin{aligned}[t]
	\frac{\dd}{\dd x}\bigl(\cosh x\bigr)
	&= \frac{\dd}{\dd x}\left(\frac{e^x+e^{-x}}2\right) \\
	&= \frac{e^x-e^{-x}}2\\
	&= \sinh x.
\end{aligned}$\hfill\null\\
So $\frac{\dd}{\dd x}\bigl(\cosh x\bigr) = \sinh x.$
	
\item \hfill$\begin{aligned}[t]
	\frac{\dd}{\dd x}\bigl(\sinh x\bigr)
	&= \frac{\dd}{\dd x}\left(\frac{e^x-e^{-x}}2\right) \\
	&= \frac{e^x+e^{-x}}2\\
	&= \cosh x.
\end{aligned}$\hfill\null\\
So $\frac{\dd}{\dd x}\bigl(\sinh x\bigr) = \cosh x.$

\item \hfill$\begin{aligned}[t]
	\frac{\dd}{\dd x}\bigl(\tanh x\bigr)
	&= \frac{\dd}{\dd x}\left(\frac{\sinh x}{\cosh x}\right) \\
	&= \frac{\cosh x \cosh x - \sinh x \sinh x}{\cosh^2 x}\\
	&= \frac{1}{\cosh^2 x}\\
	&= \sech^2 x.
\end{aligned}$\hfill\null\\
So $\frac{\dd}{\dd x}\bigl(\tanh x\bigr) = \sech^2 x.$
\end{enumerate}
\end{example}

The following Key Idea summarizes many of the important identities relating to hyperbolic functions. Each can be verified by referring back to \autoref{def:hyperbolic_functions}.

{
\tcbset{grow to right by=16em}
\begin{keyidea}[Useful Hyperbolic Function Properties]\label{idea:hyperbolic_identities}%
\lxAddClass{columns3}
\begin{minipage}[t]{.33\linewidth}
\textbf{Basic Identities}\par
\index{hyperbolic function!identities}\index{hyperbolic function!derivatives}\index{hyperbolic function!integrals}\index{derivative!hyperbolic funct.}\index{integration!hyperbolic funct.}%
\begin{enumerate}
\item $\cosh^2x-\sinh^2x=1$
\item	$\tanh^2x+\sech^2x=1$
\item	$\coth^2x-\csch^2x = 1$
\item	$\cosh 2x=\cosh^2x+\sinh^2x$
\item	$\sinh 2x = 2\sinh x\cosh x$
\item	$\ds\cosh^2x = \frac{\cosh 2x+1}{2}$
\item $\ds \sinh^2x=\frac{\cosh 2x-1}{2}$
\end{enumerate}
\end{minipage}%
\begin{minipage}[t]{.33\linewidth}
\textbf{Derivatives}
\begin{enumerate}
\item $\frac{\dd}{\dd x}\bigl(\cosh x\bigr) = \sinh x$
\item $\frac{\dd}{\dd x}\bigl(\sinh x\bigr) = \cosh x$
\item $\frac{\dd}{\dd x}\bigl(\tanh x\bigr) = \sech^2 x$
\item $\frac{\dd}{\dd x}\bigl(\sech x\bigr) = -\sech x\tanh x$
\item $\frac{\dd}{\dd x}\bigl(\csch x\bigr) = -\csch x\coth x$
\item $\frac{\dd}{\dd x}\bigl(\coth x\bigr) = -\csch^2x$
\end{enumerate}
\end{minipage}%
\begin{minipage}[t]{.33\linewidth}
\textbf{Integrals}
\begin{enumerate}
\item $\ds\int \cosh x\dd x = \sinh x+C$
\item $\ds\int \sinh x\dd x = \cosh x+C$
\item $\ds\int \tanh x\dd x = \ln(\cosh x) +C$
\item $\ds\int \coth x\dd x = \ln\abs{\sinh x}+C$
\end{enumerate}
\end{minipage}
\end{keyidea}
}

We practice using \autoref{idea:hyperbolic_identities}.

\begin{example}[Derivatives and integrals of hyperbolic functions]\label{ex_hf2}%
Evaluate the following derivatives and integrals.\\
\begin{minipage}[t]{.5\linewidth}
\begin{enumerate}
\item		$\ds\frac{\dd}{\dd x}\bigl(\cosh 2x\bigr)$
\item		$\ds\int \sech^2(7t-3)\dd t$
\end{enumerate}
\end{minipage}%
\begin{minipage}[t]{.5\linewidth}
\begin{enumerate}\addtocounter{enumi}{2}
\item		$\ds \int_0^{\ln 2} \cosh x\dd x$
\end{enumerate}
\end{minipage}
\solution
\begin{enumerate}
\item		Using the Chain Rule directly, we have $\frac{\dd}{\dd x} \bigl(\cosh 2x\bigr) = 2\sinh 2x$.

Just to demonstrate that it works, let's also use the Basic Identity found in \autoref{idea:hyperbolic_identities}: $\cosh 2x = \cosh^2x+\sinh^2x$.
\begin{align*}
\frac{\dd}{\dd x}\bigl(\cosh 2x\bigr) = \frac{\dd}{\dd x}\bigl(\cosh^2x+\sinh^2x\bigr)
&= 2\cosh x\sinh x+ 2\sinh x\cosh x\\
&= 4\cosh x\sinh x.
\end{align*}
Using another Basic Identity, we can see that $4\cosh x\sinh x = 2\sinh 2x$. We get the same answer either way.

\item	  We employ substitution, with $u = 7t-3$ and $\dd u = 7\dd t$. Applying
\autoref{idea:hyperbolic_identities} we have:
\[\int \sech^2 (7t-3)\dd t = \frac17\tanh (7t-3) + C.\]

\item \[\int_0^{\ln 2} \cosh x\dd x = \sinh x\Big|_0^{\ln 2} = \sinh (\ln 2) - \sinh 0 = \sinh(\ln 2).\]

We can simplify this last expression as $\sinh x$ is based on exponentials:
\[\sinh(\ln 2) = \frac{e^{\ln 2}-e^{-\ln 2}}2 = \frac{2-1/2}{2} = \frac34.\]
\end{enumerate}
\end{example}

\subsection{Inverse Hyperbolic Functions}

Just as the inverse trigonometric functions are useful in certain integrations, the inverse hyperbolic functions are useful with others. \autoref{fig:hfinverse2} shows the restrictions on the domains to make each function one-to-one and the resulting domains and ranges of their inverse functions. Their graphs are shown in \autoref{fig:hfinverse1}.\index{hyperbolic function!inverse}

Because the hyperbolic functions are defined in terms of exponential functions, their inverses can be expressed in terms of logarithms as shown in \autoref{idea:hyperbolic_log}. It is often more convenient to use $\sinh^{-1}x$ than $\ln\bigl(x+\sqrt{x^2+1}\bigr)$, especially when one is working on theory and does not need to compute actual values. On the other hand, when computations are needed, technology is often helpful but many hand-held calculators lack a \emph{convenient} $\sinh^{-1}x$ button. (Often it can be accessed under a menu system, but not conveniently.) In such a situation, the logarithmic representation is useful. The reader is not encouraged to memorize these, but rather know they exist and know how to use them when needed.

\noindent\begin{minipage}[t]{\linewidth}\tagstructbegin{tag=Sect}\noindent%
\captionsetup{type=figure}%
\flushinner{%
\small
\tagpdfsetup{table/header-rows={1}}
\begin{tabular}{ l c c @{\hspace{2em}} c c c }
Function & Domain & Range & Function & Domain & Range \\ \cmidrule(r{2em}){1-3} \cmidrule(l{-1em}){4-6}
$\cosh x$ & $[0,\infty)$ & $[1,\infty)$ &
 $\cosh^{-1} x$ & $[1,\infty)$ & $[0,\infty)$ \\
$\sinh x$ & $(-\infty,\infty)$ & $(-\infty,\infty)$ &
 $\sinh^{-1} x$ & $(-\infty,\infty)$ & $(-\infty,\infty)$\\
$\tanh x$ & $(-\infty,\infty)$ & $(-1,1)$ &
 $\tanh^{-1} x$ & $(-1,1)$ & $(-\infty,\infty)$\\
$\sech x$ & $[0,\infty)$ & $(0,1]$ & $\sech^{-1} x$ & $(0,1]$ & $[0,\infty)$\\
$\csch x$ & $(-\infty,0) \cup (0,\infty)$ & $(-\infty,0) \cup (0,\infty)$ &
 $\csch^{-1} x$ & $(-\infty,0) \cup (0,\infty)$ & $(-\infty,0) \cup (0,\infty)$\\
$\coth x$ & $(-\infty,0) \cup (0,\infty)$ & $(-\infty,-1) \cup (1,\infty)$ &
 $\coth^{-1} x$ & $(-\infty,-1) \cup (1,\infty)$ & $(-\infty,0) \cup (0,\infty)$
\end{tabular}}
\caption{Domains and ranges of the hyperbolic and inverse hyperbolic functions.}
\label{fig:hfinverse2}\tagstructend
\end{minipage}

\noindent\begin{minipage}[t]{\linewidth}\tagstructbegin{tag=Sect}\noindent%
\captionsetup{type=figure}%
%\flushinner{%
\addtolength{\tabcolsep}{6pt}
\tagpdfsetup{table/tagging=presentation}
\begin{tabular}{cc}
\begin{tikzpicture}[alt={cosh for positive x is in blue.  Its inverse is in red, and has been reflected across the dashed line y=x.}]
\begin{axis}[width=1.16\marginparwidth,tick label style={font=\scriptsize},
axis y line=middle,axis x line=middle,name=myplot,axis on top,axis equal,
ymin=-.9,ymax=10.9,xmin=-.9,xmax=10.9,scaled ticks=false]
\addplot [draw={\colorone},thick,smooth,domain=0:3] {cosh(x)};
\draw (axis cs:8,4) node {\scriptsize $y=\cosh^{-1} x$};
\draw (axis cs:5.6,10) node {\scriptsize $y=\cosh x$};
\addplot [draw={\colortwo},smooth,thick,domain=0:3] ({cosh(x)},x);
\addplot [dashed,domain=-.5:10] {x};
\end{axis}
\node [right] at (myplot.right of origin) {\scriptsize $x$};
\node [above] at (myplot.above origin) {\scriptsize $y$};
\end{tikzpicture}
&
\begin{tikzpicture}[alt={sinh is in blue.  Its inverse is in red, and has been reflected across the dashed line y=x.}]
\begin{axis}[width=1.16\marginparwidth,tick label style={font=\scriptsize},
axis y line=middle,axis x line=middle,name=myplot,axis on top,axis equal,
ymin=-10.9,ymax=10.9,xmin=-10.9,xmax=10.9,scaled ticks=false]
\addplot [draw={\colorone},thick,smooth,domain=-3:3] {sinh(x)};
\draw (axis cs:-6,7) node {\scriptsize $y=\sinh x$};
\draw (axis cs:6,-5) node {\scriptsize $y=\sinh^{-1} x$};
\draw[->,>=stealth] (axis cs:-1,7) -- (axis cs:2,7);
\draw[->,>=stealth] (axis cs:7,-3) -- (axis cs:7,2);
\addplot [draw={\colortwo},smooth,thick,domain=-3:3] ({sinh(x)},x);
\addplot [dashed,domain=-10:10] {x};
\end{axis}
\node [right] at (myplot.right of origin) {\scriptsize $x$};
\node [above] at (myplot.above origin) {\scriptsize $y$};
\end{tikzpicture}
\\[15pt]
\begin{tikzpicture}[alt={Inverse tanh is in blue, and looks like (trigonometric) tangent.  Inverse coth is in red, and looks like 1/(x-1) for x>1, and 1/(x+1) for x<-1.}]
\begin{axis}[width=1.16\marginparwidth,tick label style={font=\scriptsize},
axis y line=middle,axis x line=middle,name=myplot,axis on top,xtick={-2,2},
ytick={-2,2},ymin=-3.2,ymax=3.2,xmin=-3.2,xmax=3.2,scaled ticks=false,axis equal]
\addplot [draw={\colortwo},thick,smooth,domain=.348:3] ({1/tanh(x)},x);
\addplot [draw={\colortwo},thick,smooth,domain=-3:-.348] ({1/tanh(x)},x);
\draw (axis cs:2.2,1.5) node {\tiny $y=\coth^{-1} x$};
\draw (axis cs:2.2,-1.5) node {\tiny $y=\tanh^{-1} x$};
\draw [->,>=stealth] (axis cs:1.8,-1.3) -- (axis cs:.2,-.2);
\draw [loosely dashed] (axis cs:-1,-3)--(axis cs:-1,3);
\draw [loosely dashed] (axis cs:1,-3)--(axis cs:1,3);
\addplot [draw={\colorone},smooth,thick,domain=-3.8:3.8] ({tanh(x)},x);
\end{axis}
\node [right] at (myplot.right of origin) {\scriptsize $x$};
\node [above] at (myplot.above origin) {\scriptsize $y$};
\end{tikzpicture}
&
\begin{tikzpicture}[alt={Inverse sech is in blue, starting with a vertical asymptote as x approaches 0 from the right, and going to the point (1,0).  Inverse csch is in red, and looks like 1/x.}]
\begin{axis}[width=1.16\marginparwidth,tick label style={font=\scriptsize},
axis y line=middle,axis x line=middle,name=myplot,axis on top,
xtick={-3,-2,-1,1,2,3},ytick={-3,-2,-1,1,2,3},
ymin=-3.5,ymax=3.5,xmin=-3.5,xmax=3.5,scaled ticks=false,axis equal]
\draw (axis cs:1.5,-1.5) node {\tiny $y=\sech^{-1} x$};
\draw (axis cs:-2.2,-1.6) node {\tiny $y=\csch^{-1} x$};
\draw [->,>=latex] (axis cs:1,-1.2) -- (axis cs:1,-.2);
\addplot [draw={\colortwo},smooth,thick,domain=-3:-.3275] ({1/sinh(x)},x);
\addplot [draw={\colortwo},smooth,thick,domain=.3275:3] ({1/sinh(x)},x);
\addplot [draw={\colorone},thick,smooth,domain=0:3] ({1/cosh(x)},x);
\end{axis}
\node [right] at (myplot.right of origin) {\scriptsize $x$};
\node [above] at (myplot.above origin) {\scriptsize $y$};
\end{tikzpicture}
\end{tabular}
%}
\caption{Graphs of the hyperbolic functions and their inverses.}\label{fig:hfinverse1}\tagstructend
\end{minipage}

Now let's consider the inverses of the hyperbolic functions. We begin with the function $f(x)=\sinh x$. Since $\fp(x)=\cosh x>0$ for all real $x$, $f$ is increasing and must be one-to-one.
% todo Tim once we have an example of finding an inverse in text/07_Inverse_Functions.tex
% We proceed as in \autoref{sec:inv_funcs}:
\begin{align*}\allowdisplaybreaks
y&=\frac{e^x-e^{-x}}2\\
2y&=e^x-e^{-x} \qquad\text{(now multiply by $e^x$)}\\
2ye^x&=e^{2x}-1 \qquad\text{(a quadratic form )}\\
\left(e^x\right)^2-2ye^x-1&=0 \qquad\text{(use the quadratic formula)}\\
e^x&=\frac{2y\pm\sqrt{4y^2+4}}2\\
e^x&=y\pm\sqrt{y^2+1} \qquad\text{(use the fact that $e^x>0$)}\\
e^x&=y+\sqrt{y^2+1}\\
x&=\ln(y+\sqrt{y^2+1})\\
\end{align*}
Finally, interchange the variable to find that
\[\sinh^{-1} x=\ln(x+\sqrt{x^2+1}).\]
In a similar manner we find that the inverses of the other hyperbolic functions are given by:

{
\tcbset{grow to right by=5em} % 5 sufficient
\begin{keyidea}[Logarithmic definitions of Inverse Hyperbolic Functions]\label{idea:hyperbolic_log}%
\mbox{}\\[-2\baselineskip]%
\index{hyperbolic function!inverse!logarithmic def.}%
\begin{multicols}{2}
\begin{enumerate}\lxAddClass{columns2}
\item $\ds\cosh^{-1}x=\ln\bigl(x+\sqrt{x^2-1}\bigr)$;\\\null\qquad $x\geq1$
\item $\ds\tanh^{-1}x=\frac12\ln\left(\frac{1+x}{1-x}\right)$;\\\null\qquad$\abs x<1$
\item $\ds\sech^{-1}x=\ln\left(\frac{1+\sqrt{1-x^2}}x\right)$;\\\null\qquad$0<x\leq1$
\item $\ds\sinh^{-1}x=\ln\bigl(x+\sqrt{x^2+1}\bigr)$\\\mbox{}
\item $\ds\coth^{-1}x=\frac12\ln\left(\frac{x+1}{x-1}\right)$;\\\null\qquad$\abs x>1$
\item $\ds\csch^{-1}x=\ln\left(\frac1x+\frac{\sqrt{1+x^2}}{\abs x}\right)$;\\\null\qquad $x\neq0$
\end{enumerate}
\end{multicols}
\end{keyidea}
}

The following Key Ideas give the derivatives and integrals relating to the inverse hyperbolic functions. In \autoref{idea:hyperbolic_inverse_integrals}, both the inverse hyperbolic and logarithmic function representations of the antiderivative are given, based on \autoref{idea:hyperbolic_log}. Again, these latter functions are often more useful than the former.
%Note how inverse hyperbolic functions can be used to solve integrals we used Trigonometric Substitution to solve in \autoref{sec:trig_sub}.



%\mtable{Logarithmic definitions of the inverse hyperbolic functions.}{fig:hfinverse5}{%
%\begin{align*}
%\cosh^{-1}x&=\ln\bigl(x+\sqrt{x^2-1}\bigr);\ x\geq1\\
%\sinh^{-1}x &= \ln\bigl(x+\sqrt{x^2+1}\bigr)\\
%\tanh^{-1}x &= \frac12\ln\left(\frac{1+x}{1-x}\right);\ \abs x<1\\
%\sech^{-1}x &= \ln\left(\frac{1+\sqrt{1-x^2}}x\right);\ 0<x\leq1\\
%\csch^{-1}x &= \ln\left(\frac1x+\frac{\sqrt{1+x^2}}{\abs x}\right);\ x\neq0\\
%\coth^{-1}x &= \frac12\ln\left(\frac{x+1}{x-1}\right);\ \abs x>1
%\end{align*}
%}

{
\tcbset{grow to right by=3em}
\begin{keyidea}[Derivatives Involving Inverse Hyperbolic Functions]\label{idea:hyperbolic_inverse_derivatives}%
\index{hyperbolic function!inverse!derivative}\index{derivative!inverse hyper.}
\mbox{}\\[-2\baselineskip]\parbox[t]{\linewidth}{%
\begin{multicols}{2}
\begin{enumerate}\lxAddClass{columns2}
\item $\ds\frac\dd{\dd x}\bigl(\cosh^{-1}x\bigr)=\frac1{\sqrt{x^2-1}}$;\\\null\qquad$x>1$
\item $\ds\frac\dd{\dd x}\bigl(\sinh^{-1}x\bigr)=\frac1{\sqrt{x^2+1}}$\\\vphantom{$x\ne0$}
\item $\ds\frac\dd{\dd x}\bigl(\tanh^{-1}x\bigr)=\frac1{1-x^2}$;\\\null\qquad$\abs x<1$
\item $\ds\frac\dd{\dd x}\bigl(\sech^{-1}x\bigr)=\frac{-1}{x\sqrt{1-x^2}}$;\\\null\qquad$0<x<1$
\item $\ds\frac\dd{\dd x}\bigl(\csch^{-1}x\bigr)=\frac{-1}{\abs x\sqrt{1+x^2}}$;\\\null\qquad$x\neq0$
\item $\ds\frac\dd{\dd x}\bigl(\coth^{-1}x\bigr)=\frac1{1-x^2}$;\\\null\qquad$\abs x>1$
\end{enumerate}
\end{multicols}}
\end{keyidea}}

{
\tcbset{grow to right by=8.1em}
% 8.1em sufficient
\begin{keyidea}[Integrals Involving Inverse Hyperbolic Functions]\label{idea:hyperbolic_inverse_integrals}%
\index{integration!inverse hyper.}\index{hyperbolic function!inverse!integration}%
\begin{anywhereenum}
\begin{align*}
&\item \int\frac1{\sqrt{x^2-a^2}}\dd x && =\cosh^{-1}\left(\frac xa\right)+C;~0<a<x && =\ln\left(x+\sqrt{x^2-a^2}\right)+C \\
&\item \int\frac1{\sqrt{x^2+a^2}}\dd x && =\sinh^{-1}\left(\frac xa\right)+C;~a>0 && =\ln\left(x+\sqrt{x^2+a^2}\right)+C \\
&\item \int\frac1{a^2-x^2}\dd x && ={\begin{cases}
\frac1a\tanh^{-1}\left(\frac xa\right)+C & \abs{x}<\abs{a} \\
\frac1a\coth^{-1}\left(\frac xa\right)+C & \abs{a}<\abs{x}
\end{cases}} && =\frac1{2a}\ln\abs{\frac{a+x}{a-x}}+C \\
&\item \int\frac1{x\sqrt{a^2-x^2}}\dd x\!\!\! &&
=-\frac1a\sech^{-1}\frac{\abs x}a+C;~0<\abs x<a && =\frac1a\ln\abs{\frac x{a+\sqrt{a^2-x^2}}}+C \\
% $\ds{}=-\frac1a\sech^{-1}\left(\frac xa\right)+C$ & $\ds{}=\frac1a\ln\left(\frac{x}{a+\sqrt{a^2-x^2}}\right)+C$ \\[\defaultaddspace]
%&\null\hfill{$\scriptstyle0<x<a$}\\[\defaultaddspace]
&\item \int\frac1{x\sqrt{x^2+a^2}}\dd x\!\!\! && =-\frac1a\csch^{-1}\frac{\abs x}a+C;~x\neq 0,\ a>0\!\!\! && =
\frac1a\ln\abs{\frac x{a+\sqrt{a^2+x^2}}}
%-\frac1a\ln\left(\frac ax+\frac{\sqrt{x^2+a^2}}{a\abs x}\right)
% sqrt(1+(x/a)^2)/|x/a| = sqrt(a^2+x^2)/|ax|
+C
\end{align*}
\end{anywhereenum}
\end{keyidea}}

We practice using the derivative and integral formulas in the following example.

\begin{example}[Derivatives and integrals involving inverse hyperbolic functions]\label{ex_hf3}%
Evaluate the following.
\begin{multicols}{2}
\begin{enumerate}\lxAddClass{columns2}
\item	$\ds \frac{\dd}{\dd x}\left[\cosh^{-1}\left(\frac{3x-2}{5}\right)\right]$
\item	$\ds \int\frac{1}{x^2-1}\dd x$
\item	$\ds \int \frac{1}{\sqrt{9x^2+10}}\dd x$
\item[]
\end{enumerate}
\end{multicols}
\solution
\begin{enumerate}
\item	Applying \autoref{idea:hyperbolic_inverse_derivatives} with the Chain Rule gives:
		\[\frac{\dd}{\dd x}\left[\cosh^{-1}\left(\frac{3x-2}5\right)\right] = \frac{1}{\sqrt{\left(\frac{3x-2}5\right)^2-1}}\cdot\frac35.\]

\item		Multiplying the numerator and denominator by $(-1)$ gives 
a %: $\ds \int \frac{1}{x^2-1}\dd x = \int \frac{-1}{1-x^2}\dd x$. The
second integral can be solved with a direct application of item \#3 from \autoref{idea:hyperbolic_inverse_integrals}, with $a=1$. Thus
\begin{align}
\int \frac{1}{x^2-1}\dd x &= -\int \frac{1}{1-x^2}\dd x \notag \\
		&= \begin{cases}-\tanh^{-1}\left(x\right)+C & x^2<1 \\
-\coth^{-1}\left(x\right)+C & 1<x^2 \end{cases} \notag\\
     &=-\frac12\ln\abs{\frac{x+1}{x-1}}+C\notag\\
     &=\frac12\ln\abs{\frac{x-1}{x+1}}+C.\label{eq:hf3}
     \end{align}

%We should note that this exact problem was solved at the beginning of \autoref{sec:partial_fraction}. In that example the answer was given as $\frac12\ln\abs{x-1}-\frac12\ln\abs{x+1}+C.$ Note that this is equivalent to the answer given in \autoeqref{eq:hf3}, as $\ln(a/b) = \ln a - \ln b$.

%The key to linking the two seemingly different answers together is \autoref{fig:hfinverse5}, where the logarithmic definitions of the inverse hyperbolic functions are given. Note that the definitions of $\tanh^{-1}x$ and $\coth^{-1}x$ are very similar; the conditions placed on $\abs x$ ensure that the argument of $\ln$ is always positive. Thus one could say 
%\[\frac12\ln\abs{\frac{x+1}{x-1}} = \begin{cases}\tanh^{-1}x+C & \abs x<1 \\ \\
%\coth^{-1}x+C & \abs x>1 \end{cases}.\]
%
%We reconcile the two answers by returning to \autoeqref{eq:hf3} and continuing:
%\begin{align*}
%\int \frac{1}{x^2-1}\dd x &= \int \frac{-1}{1-x^2}\dd x \\
%			&= \begin{cases}-\frac1a\tanh^{-1}\left(\frac xa\right)+C & x^2<a^2 \\ \\
%-\frac1a\coth^{-1}\left(\frac xa\right)+C & a^2<x^2 \end{cases}\\
%			&= -\frac12\ln\abs{\frac{x+1}{x-1}}+C \\
%			&= -\frac12\ln\abs{x+1} + \frac12\ln\abs{x-1} +C,
%\end{align*}
%matching the answer previously obtained.

\item		This requires a substitution, then item \#2 of \autoref{idea:hyperbolic_inverse_integrals} can be applied.

Let $u = 3x$, hence $\dd u = 3\dd x$. We have 
\begin{align*}
	\int \frac{1}{\sqrt{9x^2+10}}\dd x
	&= \frac13\int\frac{1}{\sqrt{u^2+10}}\dd u.
	\intertext{Note \intertextmath{a^2=10}, hence \intertextmath{a = \sqrt{10}}. Now apply the integral rule.}
	 &= \frac13 \sinh^{-1}\left(\frac{3x}{\sqrt{10}}\right) + C \\
	 &= \frac13 \ln \abs{3x+\sqrt{9x^2+10}}+C.
\end{align*}
\end{enumerate}
\end{example}

This section covers a lot of ground. New functions were introduced, along with some of their fundamental identities, their derivatives and antiderivatives, their inverses, and the derivatives and antiderivatives of these inverses. Four Key Ideas were presented, each including quite a bit of information.

Do not view this section as containing a source of information to be memorized, but rather as a reference for future problem solving. \autoref{idea:hyperbolic_inverse_integrals} contains perhaps the most useful information. Know the integration forms it helps evaluate and understand how to use the inverse hyperbolic answer and the logarithmic answer.

The next section takes a brief break from demonstrating new integration techniques. It instead demonstrates a technique of evaluating limits that return indeterminate forms. This technique will be useful in \autoref{sec:improper_integration}, where limits will arise in the evaluation of certain definite integrals.


\printexercises{exercises/06-05-exercises}

 \section{L'Hôpital's Rule}\label{sec:lhopitals_rule}

This section is concerned with a technique for evaluating certain limits that will be useful in later chapters.

Our treatment of limits exposed us to ``0/0'', an indeterminate form. If both $\ds \lim_{x\to c}f(x)$ and $\ds \lim_{x\to c} g(x)$ are zero, we do not conclude that $\ds \lim_{x\to c} f(x)/g(x)$ is $0/0$; rather, we use $0/0$ as notation to describe the fact that both the numerator and denominator approach 0. The expression 0/0 has no numeric value; other work must be done to evaluate the limit.

Other indeterminate forms exist; they are: $\infty/\infty$, $0\cdot\infty$, $\infty-\infty$, $0^0$, $1^\infty$ and $\infty^0$. Just as ``0/0'' does not mean ``divide 0 by 0,'' the expression ``$\infty/\infty$'' does not mean ``divide infinity by infinity.'' Instead, it means ``a quantity is growing without bound and is being divided by another quantity that is growing without bound.'' We cannot determine from such a statement what value, if any, results in the limit. Likewise, ``$0\cdot \infty$'' does not mean ``multiply zero by infinity.'' Instead, it means ``one quantity is shrinking to zero, and is being multiplied by a quantity that is growing without bound.'' We cannot determine from such a description what the result of such a limit will be.

This section introduces L'Hôpital's Rule, a method of resolving limits that produce the indeterminate forms 0/0 and $\infty/\infty$. We'll also show how algebraic manipulation can be used to convert other indeterminate expressions into one of these two forms so that our new rule can be applied.

\begin{theorem}[L'Hôpital's Rule, Part 1]\label{thm:LHR_1}%
Let $f$ and $g$ be differentiable functions on an open interval $I$ containing $a$.
\begin{enumerate}
\item If $\ds\lim_{x\to a}f(x)=0$, $\ds\lim_{x\to a}g(x)=0$, and $g\primeskip'(x)\neq 0$ except possibly at $x=a$, then \[\lim_{x\to a}\frac{f(x)}{g(x)}=\lim_{x\to a} \frac{\fp(x)}{g\primeskip'(x)},\]
assuming that the limit on the right exists.
\item If  $\ds\lim_{x\to a}f(x)=\pm\infty$ and $\ds\lim_{x\to a}g(x)=\pm\infty$, then \[\lim_{x\to a}\frac{f(x)}{g(x)}=\lim_{x\to a} \frac{\fp(x)}{g\primeskip'(x)},\]
assuming that the limit on the right exists.
\index{LHopitals Rule@L'Hôpital's Rule}
\end{enumerate}
\end{theorem}

A similar statement holds if we just look at the one sided limits $\ds\lim_{x\to a^-}$ and $\ds\lim_{x\to a^+}$.

\begin{theorem}[L'Hôpital's Rule, Part 2]\label{thm:LHR_2}%
Let $f$ and $g$ be differentiable functions on the open interval $(c,\infty)$ for some value $c$ and $g\primeskip'(x)\neq0$ on $(c,\infty)$.
\begin{enumerate}
\item If $\ds\lim_{x\to\infty}f(x)=0$ and $\ds\lim_{x\to \infty}g(x)=0$, then
\[\lim_{x\to\infty}\frac{f(x)}{g(x)}=\lim_{x\to\infty}\frac{\fp(x)}{g\primeskip'(x)},\]
assuming that the limit on the right exists.
\item If  $\ds\lim_{x\to \infty}f(x)=\pm\infty$ and $\ds\lim_{x\to \infty}g(x)=\pm\infty$, then
\[\lim_{x\to\infty}\frac{f(x)}{g(x)}=\lim_{x\to\infty}\frac{\fp(x)}{g\primeskip'(x)},\]
assuming that the limit on the right exists.
\end{enumerate}
Similar statements can be made where $x$ approaches $-\infty$.
\index{LHopitals Rule@L'Hôpital's Rule}
\end{theorem}

We demonstrate the use of L'Hôpital's Rule in the following examples; we will often use ``LHR'' as an abbreviation of ``L'Hôpital's Rule.''

\begin{example}[Using L'Hôpital's Rule]\label{ex_lhr1}%
Evaluate the following limits, using L'Hôpital's Rule as needed.\\
\begin{minipage}[t]{.5\textwidth}
\begin{enumerate}
\item		$\ds \lim_{x\to0}\frac{\sin x}x$
\item		$\ds \lim_{x\to 1}\frac{\sqrt{x+3}-2}{1-x}$
\item		$\ds \lim_{x\to0}\frac{x^2}{1-\cos x}$
\end{enumerate}
\end{minipage}%
\begin{minipage}[t]{.5\textwidth}
\begin{enumerate}\addtocounter{enumi}{3}
	\item	$\ds\lim_{x\to-3}\frac{x^3+27}{x^2+9}$
	\item	$\ds\lim_{x\to\infty}\frac{3x^2-100x+2}{4x^2+5x-1000}$
	\item	$\ds\lim_{x\to\infty}\frac{e^x}{x^3}$
\end{enumerate}
\end{minipage}
\solution
\begin{enumerate}
	\item	This has the indeterminate form $0/0$. We proved this limit is 1 in \autoref{ex_limit_sinx_prove} using the Squeeze Theorem. Here we use L'Hôpital's Rule to show its power.
\[\lim_{x\to0}\frac{\sin x}x \LHequals \lim_{x\to0} \frac{\cos x}{1}=1.\]
While this seems easier than using the Squeeze Theorem to find this limit, we note that applying L'Hôpital's Rule here requires us to know the derivative of $\sin x$. We originally encountered this limit when we were trying to find that derivative.

	\item	This has the indeterminate form $0/0$.
\[\lim_{x\to 1}\frac{\sqrt{x+3}-2}{1-x} 	 \LHequals \lim_{x \to 1} \frac{\frac12(x+3)^{-1/2}}{-1} =-\frac 14.\]

	\item	This has the indeterminate form $0/0$.
\[\lim_{x\to 0}\frac{x^2}{1-\cos x}  \LHequals  \lim_{x\to 0} \frac{2x}{\sin x}.\]
This latter limit also evaluates to the $0/0$ indeterminate form. To evaluate it, we apply L'Hôpital's Rule again.
\[
 \lim_{x\to 0} \frac{2x}{\sin x}
 \LHequals \lim_{x\to 0} \frac{2}{\cos x} = 2 .
\]
Thus $\ds \lim_{x\to0}\frac{x^2}{1-\cos x}=2.$

	\item \hfill$\ds\lim_{x\to-3}\frac{x^3+27}{x^2+9}
	 =\frac0{18}=0$\hfill\null\smallskip\\
We cannot use L'Hôpital's Rule in this case because the original limit does not return an indeterminate form, so L'Hôpital's Rule does not apply. In fact, the inappropriate use of L'Hôpital's Rule here would result in the incorrect limit $-\frac92$.

	\item	We can evaluate this limit already using \autoref{thm:lim_rational_fn_at_infty}; the answer is 3/4. We apply L'Hôpital's Rule to demonstrate its applicability.
\[
 \lim_{x\to\infty} \frac{3x^2-100x+2}{4x^2+5x-1000}
 \LHequals \lim_{x\to\infty} \frac{6x-100}{8x+5}
 \LHequals \lim_{x\to\infty} \frac68 = \frac34.
\]

	\item	$\ds\lim_{x\to \infty}\frac{e^x}{x^3} \LHequals \lim_{x\to\infty} \frac{e^x}{3x^2} \LHequals \lim_{x\to\infty} \frac{e^x}{6x} \LHequals \lim_{x\to\infty} \frac{e^x}{6} = \infty.$

Recall that this means that the limit does not exist; as $x$ approaches $\infty$, the expression $e^x/x^3$ grows without bound. We can infer from this that $e^x$ grows ``faster'' than $x^3$; as $x$ gets large, $e^x$ is far larger than $x^3$. (This has important implications in computing when considering efficiency of algorithms.)
\end{enumerate}
\end{example}

\subsection{Indeterminate Forms \texorpdfstring{$0\cdot\infty$ and $\infty-\infty$}{0·∞ and ∞-∞}}

L'Hôpital's Rule can only be applied to ratios of functions. When faced with an indeterminate form such as $0\cdot\infty$ or $\infty-\infty$, we can sometimes apply algebra to rewrite the limit so that L'Hôpital's Rule can be applied. We demonstrate the general idea in the next example.
\index{limit!indeterminate form}\index{indeterminate form}

\youtubeVideo{kEnwac_9lyg}{L'Hôpital's Rule --- Indeterminate Powers}

\begin{example}[Applying L'Hôpital's Rule to other indeterminate forms]\label{ex_LHR3}%
Evaluate the following limits.\\
\begin{minipage}[t]{.5\textwidth}
\begin{enumerate}
\item		$\ds \lim_{x\to0^+} x\cdot e^{1/x}$
\item		$\ds \lim_{x\to0^-} x\cdot e^{1/x}$
\end{enumerate}
\end{minipage}%
\begin{minipage}[t]{.5\textwidth}
\begin{enumerate}\addtocounter{enumi}{2}
\item		$\ds \lim_{x\to\infty} \left(\ln(x+1)-\ln x\right)$
\item		$\ds \lim_{x\to\infty} \left(x^2-e^x\right)$
\end{enumerate}
\end{minipage}
\solution
\begin{enumerate}
	\item	As $x\to 0^+$, note that $x\to 0$ and $e^{1/x}\to \infty$. Thus we have the indeterminate form $0\cdot\infty$. We rewrite the expression $x\cdot e^{1/x}$ as $\ds\frac{e^{1/x}}{1/x}$; now, as $x\to 0^+$, we get the indeterminate form $\infty/\infty$ to which L'Hôpital's Rule can be applied. 
\[
\lim_{x\to0^+} x\cdot e^{1/x} = \lim_{x\to 0^+} \frac{e^{1/x}}{1/x} \LHequals \lim_{x\to 0^+}\frac{(-1/x^2)e^{1/x}}{-1/x^2} =\lim_{x\to 0^+}e^{1/x} =\infty.
\]

Interpretation: $e^{1/x}$ grows ``faster'' than $x$ shrinks to zero, meaning their product grows without bound.

	\item	As $x\to 0^-$, note that $x\to 0$ and $e^{1/x}\to e^{-\infty}\to 0$. The the limit evaluates to $0\cdot 0$ which is not an indeterminate form. We conclude then that
	\[\lim_{x\to 0^-}x\cdot e^{1/x} = 0.\]

	\item	This limit initially evaluates to the indeterminate form $\infty-\infty$. By applying a logarithmic rule, we can rewrite the limit as 
\[
\lim_{x\to\infty}\left(\ln(x+1)-\ln x\right) = \lim_{x\to \infty} \ln \left(\frac{x+1}x\right).
\]

As $x\to \infty$, the argument of the natural logarithm approaches $\infty/\infty$, to which we can apply L'Hôpital's Rule.
\[\lim_{x\to\infty} \frac{x+1}x \LHequals \lim_{x\to\infty}\frac11=1.
\]

Since $x\to\infty$ implies $\ds\frac{x+1}x\to 1$, it follows that 
\[x\to\infty \quad \text{ implies }\quad \ln\left(\frac{x+1}x\right)\to\ln 1=0.\]

Thus
\[
 \lim_{x\to\infty} \left(\ln(x+1)-\ln x\right)
 = \lim_{x\to \infty} \ln \left(\frac{x+1}x\right)=0.
\]
Interpretation: since this limit evaluates to 0, it means that for large $x$, there is essentially no difference between $\ln (x+1)$ and $\ln x$; their difference is essentially 0.

	\item	The limit $\ds \lim_{x\to\infty} \left(x^2-e^x\right)$ initially returns the indeterminate form $\infty-\infty$. We can rewrite the expression by factoring out $x^2$; $\ds x^2-e^x = x^2\left(1-\frac{e^x}{x^2}\right).$ We need to evaluate how $e^x/x^2$ behaves as $x\to\infty$:
\[
\lim_{x\to\infty}\frac{e^x}{x^2} \LHequals \lim_{x\to\infty} \frac{e^x}{2x}
\LHequals \lim_{x\to\infty} \frac{e^x}{2} = \infty.
\]

Thus $\lim_{x\to\infty}x^2(1-e^x/x^2)$ evaluates to $\infty\cdot(-\infty)$, which is not an indeterminate form; rather, $\infty\cdot(-\infty)$ evaluates to $-\infty$. We conclude that 
$\ds \lim_{x\to\infty} \left(x^2-e^x\right) = -\infty.$

Interpretation: as $x$ gets large, the difference between $x^2$ and $e^x$ grows very large.
\end{enumerate}
\end{example}

\subsection{Indeterminate Forms\ \ \texorpdfstring{$0^0$, $1^\infty$, and $\infty^0$}{0\^{}0, 1\^{}∞, and ∞\^{}0}}

When faced with a limit that returns one of the indeterminate forms $0^0$, $1^\infty$, or $\infty^0$, it is often useful to use the natural logarithm to convert to an indeterminate form we already know how to find the limit of, then use the natural exponential function find the original limit. This is possible because the natural logarithm and natural exponential functions are inverses and because they are both continuous. The following Key Idea expresses the concept, which is followed by an example that demonstrates its use.

\begin{keyidea}[Evaluating Limits Involving Indeterminate Forms $0^0$, $1^\infty$ and $\infty^0$]\label{idea:LHR_power}%
If $\ds \lim_{x\to c} \ln\bigl(f(x)\bigr) = L$,\quad then 
$\ds \lim_{x\to c} f(x) = \lim_{x\to c} e^{\ln(f(x))} = e\,^L.$ \index{limit!indeterminate form}\index{indeterminate form}
\end{keyidea}

\begin{example}[Using L'Hôpital's Rule with indeterminate forms involving exponents]\label{ex_LHR4}%
Evaluate the following limits.
\[
 \text{1.}\quad\lim_{x\to\infty} \left(1+\frac1x\right)^x \qquad\qquad
 \text{2.}\quad\lim_{x\to0^+} x^x.
\]
\solution
\begin{enumerate}
\item		This is equivalent to a special limit given in \autoref{thm:special_limits}; these limits have important applications in mathematics and finance. Note that the exponent approaches $\infty$ while the base approaches 1, leading to the indeterminate form $1^\infty$. Let $f(x) = (1+1/x)^x$; the problem asks to evaluate $\ds\lim_{x\to\infty}f(x)$. Let's first evaluate $\ds \lim_{x\to\infty}\ln\bigl(f(x)\bigr)$.
\begin{align*}
\lim_{x\to\infty}\ln\bigl(f(x)\bigr)
			&= \lim_{x\to\infty} \ln \left(1+\frac1x\right)^x \\
			&= \lim_{x\to\infty} x\ln\left(1+\frac1x\right)\\
			&= \lim_{x\to\infty} \frac{\ln\left(1+\frac1x\right)}{1/x}\\
			\intertext{This produces the indeterminate form 0/0, so we apply L'Hôpital's Rule.}
			&=	\lim_{x\to\infty} \frac{\frac{1}{1+1/x}\cdot(-1/x^2)}{(-1/x^2)} \\
			&= \lim_{x\to\infty}\frac{1}{1+1/x}\\
			&= 1.
\end{align*}
Thus $\ds\lim_{x\to\infty} \ln \bigl(f(x)\bigr) = 1.$ We return to the original limit and apply \autoref{idea:LHR_power}.
\[\lim_{x\to\infty}\left(1+\frac1x\right)^x = \lim_{x\to\infty} f(x) =  \lim_{x\to\infty}e^{\ln (f(x))} = e^1 = e.\]
This is another way to determine the value of the number $e$.

\item		This limit leads to the indeterminate form $0^0$. Let $f(x) = x^x$ and consider first $\ds\lim_{x\to0^+} \ln\bigl(f(x)\bigr)$. 
%
\mtable{A graph of $f(x)=x^x$ supporting the fact that as $x\to 0^+$, $f(x)\to 1$.}{fig:LHR4}{\begin{tikzpicture}[alt={A graph starting near (0,1), moving downward as it goes right, and then moving up to finish near (2,4).}]
\begin{axis}[width=1.16\marginparwidth,tick label style={font=\scriptsize},
axis y line=middle,axis x line=middle,name=myplot,axis on top,ytick={1,2,3,4},
ymin=-.4,ymax=4.5,xmin=-.1,xmax=2.2,scaled ticks=false]
\addplot[draw={\colorone},thick,smooth,domain=.01:2]{exp(x*ln(x))};
\draw (axis cs:1,1) node [below right] {\scriptsize $f(x)=x^x$};
\end{axis}
\node [right] at (myplot.right of origin) {\scriptsize $x$};
\node [above] at (myplot.above origin) {\scriptsize $y$};
\end{tikzpicture}}%
%
\begin{align*}
\lim_{x\to0^+} \ln\bigl(f(x)\bigr) &= \lim_{x\to0^+} \ln\left(x^x\right) \\
			&= \lim_{x\to0^+} x\ln x \\
			\intertext{This produces the indeterminate form \intertextmath{0(-\infty)}, so we rewrite it in order to apply L'Hôpital's Rule.}
			&= \lim_{x\to0^+} \frac{\ln x}{1/x}.\\
			\intertext{This produces the indeterminate form \intertextmath{-\infty/\infty} so we apply L'Hôpital's Rule.}
			&=	\lim_{x\to0^+} \frac{1/x}{-1/x^2} \\
			&= \lim_{x\to0^+} -x \\
			&= 0.
\end{align*}%
Thus $\ds\lim_{x\to0^+} \ln\bigl(f(x)\bigr) =0$. We return to the original limit and apply \autoref{idea:LHR_power}.
\[
\lim_{x\to0^+} x^x = \lim_{x\to0^+} f(x) = \lim_{x\to0^+} e^{\ln(f(x))} = e^0 = 1.
\]
This result is supported by the graph of $f(x)=x^x$ given in \autoref{fig:LHR4}.
\end{enumerate}
\end{example}

% todo Tim do we want to add a hierarchy of function growth to the end of LH section?

%Our brief revisit of limits will be rewarded in the next section where we consider \emph{improper integration.} So far, we have only considered definite integrals where the bounds are finite numbers, such as $\ds \int_0^1 f(x)\dd x$. Improper integration considers integrals where one, or both, of the bounds are ``infinity.'' Such integrals have many uses and applications, in addition to generating ideas that are enlightening.

\printexercises{exercises/06-06-exercises}
%
}{}



\apexchapter{Sequences and Series}{chapter:sequences_series}

This chapter introduces \textbf{sequences} and \textbf{series}, important mathematical constructions that are useful when solving a large variety of mathematical problems. The content of this chapter is considerably different from the content of the chapters before it. While the material we learn here definitely falls under the scope of ``calculus,'' we will make very little use of derivatives or integrals. Limits are extremely important, though, especially limits that involve infinity. 

One of the problems addressed by this chapter is this: suppose we know information about a function and its derivatives at a point, such as  $f(1) = 3$, $\fp(1) = 1$, $\fp'(1) = -2$, $\fp''(1) = 7$, and so on. What can I say about $f(x)$ itself? Is there any reasonable approximation of the value of $f(2)$? The topic of Taylor Series addresses this problem, and allows us to make excellent approximations of functions when limited knowledge of the function is available.


\section{Sequences}\label{sec:sequences}

We commonly refer to a set of events that occur one after the other as a \emph{sequence} of events. In mathematics, we use the word \emph{sequence} to refer to an ordered set of numbers, i.e., a set of numbers that ``occur one after the other.''

For instance, the numbers $2, 4, 6, 8, \dotsc$, form a sequence. The order is important; the first number is 2, the second is 4, etc. It seems natural to seek a formula that describes a given sequence, and often this can be done. For instance, the sequence above could be described by the function $a(n) = 2n$, for the values of $n = 1, 2, \dotsc$ (it could also be described by $n^4-10 n^3+35 n^2-48n+24$, to give one of infinitely many other options). To find the 10$^\text{th}$ term in the sequence, we would compute $a(10)$. This leads us to the following, formal definition of a sequence.

\begin{definition}[Sequence]\label{def:sequence}%
%
\mnote{\textbf{Notation:} We use $\mathbb{N}$ to describe the set of natural numbers, that is, the integers $1, 2, 3, \dotsc$}
%
A \textbf{sequence} is a function $a(n)$ whose domain is $\mathbb{N}$. The \textbf{range} of a sequence is the set of all distinct values of $a(n)$.
\index{sequences!definition}\bigskip

The \textbf{terms} of a sequence are the values $a(1), a(2), \dotsc$, which are usually denoted with subscripts as $a_1, a_2, \dotsc$.\bigskip

A sequence $a(n)$ is often denoted as $\{a_n\}$.
\end{definition}

\youtubeVideo{9K1xx6wfN-U}{Sequences --- Examples showing convergence or divergence}

\mnote[-1.2in]{\textbf{Factorial:} The expression $3!$ refers to the number $3\cdot2\cdot1 = 6$.
\index{factorial}\bigskip\\
In general, $n! = n\cdot (n-1)\cdot(n-2)\dotsm 2\cdot1$, where $n$ is a natural number.\bigskip\\
We define $0! = 1$. While this does not immediately make sense, it makes many mathematical formulas work properly.}

\mtable{Plotting sequences in \autoref{ex_seq1}.}{fig:seq1b}{%
\begin{tikzpicture}[alt={Blue dots for aₙ=3^n/n^! at 1≤n≤4; values about 3, 4.5, 4.5, 3.4; discrete plot, axes n,y.}]
\begin{axis}[width=\marginparwidth,tick label style={font=\scriptsize},
axis y line=middle,axis x line=middle,name=myplot,axis on top,xtick={1,2,3,4},
ytick={1,2,3,4,5},ymin=-.1,ymax=5.5,xmin=-.1,xmax=4.5]
\addplot [only marks,draw={\colorone},fill={\colorone},mark size={1.75pt}] coordinates {(1,3)(2,4.5)(3,4.5)(4,3.375)};
\draw (axis cs:2.5,1.5) node {\scriptsize $\ds a_n = \frac{3^n}{n!}$};
\end{axis}
\node [right] at (myplot.right of origin) {\scriptsize $n$};
\node [above] at (myplot.above origin) {\scriptsize $y$};
\end{tikzpicture}
\\(a)\\
\begin{tikzpicture}[alt={Blue dots for aₙ=4+(–1)^n at 1≤n≤4; points alternate 3 and 5; discrete plot, axes n,y.}]
\begin{axis}[width=\marginparwidth,tick label style={font=\scriptsize},
axis y line=middle,axis x line=middle,name=myplot,axis on top,xtick={1,2,3,4},
ytick={1,2,3,4,5},ymin=-.1,ymax=5.5,xmin=-.1,xmax=4.5]
\addplot [only marks,draw={\colorone},fill={\colorone},mark size={1.75pt},domain=1:4] {4+(-1)^x};
\draw (axis cs:2,1) node {\scriptsize $\ds a_n = 4+(-1)^n$};
\end{axis}
\node [right] at (myplot.right of origin) {\scriptsize $n$};
\node [above] at (myplot.above origin) {\scriptsize $y$};
\end{tikzpicture}
\\(b)\\
\begin{tikzpicture}[alt={{Blue dots for aₙ=(–1)^{(n+1)/2}/n² at 1≤n≤5; alternating sign, magnitudes shrink to 0; axes n,y.}}]
\begin{axis}[width=\marginparwidth,tick label style={font=\scriptsize},
axis y line=middle,axis x line=middle,name=myplot,axis on top,xtick={1,2,3,4,5},
ytick={-1},extra y ticks={.5,.25},extra y tick labels={$1/2$,$1/4$},
ymin=-1.1,ymax=0.6,xmin=-.1,xmax=5.5]
\addplot [only marks,draw={\colorone},fill={\colorone},mark size={1.75pt}] coordinates{(1,-1)(2,-.25)(3,.111)(4,0.0625)(5,-.04)};
\draw (axis cs:3,-.5) node {\scriptsize $\ds a_n = \frac{(-1)^{n(n+1)/2}}{n^2}$};
\end{axis}
\node [right] at (myplot.right of origin) {\scriptsize $n$};
\node [above] at (myplot.above origin) {\scriptsize $y$};
\end{tikzpicture}
\\(c)}

\begin{example}[Listing terms of a sequence]\label{ex_seq1}%
List the first four terms of the following sequences.
\[
 \text{1. }\{a_n\} = \left\{\frac{3^n}{n!}\right\}\qquad
 \text{2. }\{a_n\} = \{4+(-1)^n\}\qquad
 \text{3. }\{a_n\} = \left\{\frac{(-1)^{n(n+1)/2}}{n^2}\right\}
\]
\solution
\begin{enumerate}
\item		$\ds a_1=\frac{3^1}{1!} = 3$;\qquad	$\ds a_2= \frac{3^2}{2!} = \frac92$;\qquad $\ds a_3 = \frac{3^3}{3!} = \frac92$; \qquad $\ds a_4 = \frac{3^4}{4!} = \frac{27}8$

We can plot the terms of a sequence with a scatter plot. The ``$x$''-axis is used for the values of $n$, and the values of the terms are plotted on the $y$-axis. To visualize this sequence, see \autoref{fig:seq1b}(a).

\item		$a_1= 4+(-1)^1 = 3$;\qquad $a_2 = 4+(-1)^2 = 5$; 

\noindent $a_3=4+(-1)^3 = 3$; \qquad $a_4 = 4+(-1)^4 = 5$.\\
Note that the range of this sequence is finite, consisting of only the values 3 and 5. This sequence is plotted in \autoref{fig:seq1b}(b).

\item		$\ds a_1= \frac{(-1)^{1(2)/2}}{1^2} = -1$; \qquad $\ds a_2 = \frac{(-1)^{2(3)/2}}{2^2} =-\frac14$

\noindent $\ds a_3 = \frac{(-1)^{3(4)/2}}{3^2} = \frac19$ \qquad $\ds a_4 = \frac{(-1)^{4(5)/2}}{4^2} = \frac1{16}$; 

\noindent $\ds a_5 = \frac{(-1)^{5(6)/2}}{5^2}=-\frac1{25}$.

\noindent We gave one extra term to begin to show the pattern of signs is ``$-$, $-$, $+$, $+$, $-$, $-$, \dots'', due to the fact that the exponent of $-1$ is a special quadratic. This sequence is plotted in \autoref{fig:seq1b}(c).
\end{enumerate}
\end{example}

\begin{example}[Determining a formula for a sequence]\label{ex_seq2}%
Find the $n^\text{th}$ term of the following sequences, i.e., find a function that describes each of the given sequences.
\begin{multicols}{2}
\begin{enumerate}\lxAddClass{columns2}
\item		$2, 5, 8, 11, 14, \dotsc$
\item		$2, -5, 10, -17, 26, -37, \dotsc$
\item		$1, 1, 2, 6, 24, 120, 720, \dotsc$
\item		$\ds \frac52, \frac52, \frac{15}8, \frac54, \frac{25}{32}, \dotsc$
\end{enumerate}
\end{multicols}
\solution
We should first note that there is never exactly one function that describes a finite set of numbers as a sequence. There are many sequences that start with 2, then 5, as our first example does. We are looking for a simple formula that describes the terms given, knowing there is possibly more than one answer.
\begin{enumerate}
\item		Note how each term is 3 more than the previous one. This implies a linear function would be appropriate: $a(n) = a_n = 3n + b$ for some appropriate value of $b$. As we want $a_1=2$, we set $b=-1$. Thus $a_n = 3n-1$.

\item		First notice how the sign changes from term to term. This is most commonly accomplished by multiplying the terms by either $(-1)^n$ or $(-1)^{n+1}$. Using $(-1)^n$ multiplies the odd terms by $(-1)$; using $(-1)^{n+1}$ multiplies the even terms by $(-1)$. As this sequence has negative even terms, we will multiply by $(-1)^{n+1}$. 

After this, we might feel a bit stuck as to how to proceed. At this point, we are just looking for a pattern of some sort: what do the numbers 2, 5, 10, 17, etc., have in common? There are many correct answers, but the one that we'll use here is that each is one more than a perfect square. That is, $2=1^2+1$, $5=2^2+1$, $10=3^2+1$, etc. Thus our formula is $a_n= (-1)^{n+1}(n^2+1)$.

\item		One who is familiar with the factorial function will readily recognize these numbers. They are $0!$, $1!$, $2!$, $3!$, etc. Since our sequences start with $n=1$, we cannot write $a_n = n!$, for this misses the $0!$ term. Instead, we shift by 1, and write $a_n = (n-1)!$.

\item		This one may appear difficult, especially as the first two terms are the same, but a little ``sleuthing'' will help. Notice how the terms in the numerator are always multiples of 5, and the terms in the denominator are always powers of 2. Does something as simple as $a_n = \frac{5n}{2^n}$ work?

When $n=1$, we see that we indeed get $5/2$ as desired. When $n=2$, we get $10/4 = 5/2$. Further checking shows that this formula indeed matches the other terms of the sequence.
\end{enumerate}
\end{example}

A common mathematical endeavor is to create a new mathematical object (for instance, a sequence) and then apply previously known mathematics to the new object. We do so here. The fundamental concept of calculus is the limit, so we will investigate what it means to find the limit of a sequence.

%{\tcbset{grow to right by=8em}}
\begin{definition}[Limit of a Sequence, Convergent, Divergent]\label{def:seq_limit}%
Let $\{a_n\}$ be a sequence and let $L$ be a real number. Given any $\epsilon>0$, if an $m$ can be found such that $\abs{a_n-L}<\epsilon$ for all $n>m$, then we say the \textbf{limit of $\{a_n\}$, as $n$ approaches infinity, is $L$}, denoted \[\lim_{n\to\infty}a_n = L.\]

If $\ds\lim_{n\to\infty} a_n$ exists, we say the sequence \textbf{converges}; otherwise, the sequence \textbf{diverges}.\index{limit!of sequence}\index{sequences!limit}\index{convergence!of sequence}\index{divergence!of sequence}\index{sequences!convergent}\index{sequences!divergent}
\end{definition}

This definition states, informally, that if the limit of a sequence is $L$, then if you go far enough out along the sequence, all subsequent terms will be \emph{really close} to $L$. Of course, the terms ``far enough'' and ``really close'' are subjective terms, but hopefully the intent is clear.

This definition is reminiscent of the $\epsilon$-$\delta$ proofs of \autoref{chapter:limits}. In that chapter we developed other tools to evaluate limits apart from the formal definition; we do so here as well.

%{\tcbset{grow to right by=3em}}
\begin{theorem}[Limit of a Sequence]\label{thm:seq_limit}%
Let $\{a_n\}$ be a sequence and let $f(x)$ be a function whose domain contains the positive real numbers where $f(n) = a_n$ for all $n$ in $\mathbb{N}$.\bigskip

If $\ds \lim_{x\to\infty} f(x) = L$, then $\ds\lim_{n\to\infty} a_n = L$.
%\begin{enumerate}
%\item		If $\ds \lim_{x\to\infty} f(x) = L$, then $\ds\lim_{n\to\infty} a_n = L$.
%\item		If $\ds \lim_{x\to\infty} f(x)$ does not exist, then $\{a_n\}$ diverges.
%\end{enumerate}
\end{theorem}

\autoref{thm:seq_limit} allows us, in certain cases, to apply the tools developed in \autoref{chapter:limits} to limits of sequences (and even if the theorem doesn't apply, we can still use ideas from that chapter to prove similar theorems for sequences). Note two things \emph{not} stated by the theorem:
	\begin{enumerate}
		\item If $\ds \lim_{x\to\infty}f(x)$ does not exist, we cannot conclude that $\ds\lim_{n\to\infty} a_n$ does not exist. It may, or may not, exist. For instance, we can define a sequence $\{a_n\} = \{\cos(2\pi n)\}$. Let $f(x) = \cos (2\pi x)$. Since the cosine function oscillates over the real numbers, the limit $\ds \lim_{x\to\infty}f(x)$ does not exist. 
		
		However, for every positive integer $n$, $\cos(2\pi n) = 1$, so $\ds \lim_{n\to\infty} a_n = 1$.
		
		%For every positive integer $n$, $\cos(2\pi n) = 1$, so $\ds \lim_{n\to\infty} a_n = 1$. 
		%
		%It is natural to set $f(x) = \cos (2\pi x)$. Since the cosine function oscillates over the real numbers, the limit $\ds \lim_{x\to\infty}f(x)$ does not exist.
		\item	If we cannot find a function $f(x)$ whose domain contains the positive real numbers where $f(n) = a_n$ for all $n$ in $\mathbb{N}$, we cannot conclude $\ds\lim_{n\to\infty} a_n$ does not exist. It may, or may not, exist.
	\end{enumerate}

%When we considered limits before, the domain of the function was an interval of real numbers. Now, as we consider limits, the domain is restricted to $\mathbb{N}$, the natural numbers. \autoref{thm:seq_limit} states that if we can extend a function whose domain is $\mathbb{N}$ to a

%\autoref{thm:seq_limit} states that this restriction of the domain does not affect the outcome of the limit and whatever tools we developed in \autoref{chapter:limits} to evaluate limits can be applied here as well.\\

%Considering again \autoref{ex_seq3}, we can now state when $\{a_n\} = \{\frac1n\}$, $\ds \lim_{n\to\infty} a_n = 0$.\\

\begin{example}[Determining convergence/divergence of a sequence]\label{ex_seq4}%
Determine the convergence or divergence of the following sequences.
%
\mtable{Scatter plots of the sequences in \autoref{ex_seq4}.}{fig:seq4}{%
 \begin{tikzpicture}[alt={Blue dots for aₙ=(3n²–2n+1)/(n²–1000) at 1≤n≤100; ; values descend then level near 3; axes n,y}]
  \begin{axis}[width=\marginparwidth,tick label style={font=\scriptsize},
    axis y line=middle,axis x line=middle,name=myplot,axis on top,%
    xtick={20,40,60,80,100},ymin=-11,ymax=11,xmin=-.1,xmax=110]
   \addplot [only marks,draw={\colorone},fill={\colorone},mark size={.75pt},domain=1:100,samples=21]
    {(3*x^2 - 2*x + 1)/(x^2 - 1000)};
   \draw (axis cs:60,-8) node {\scriptsize $a_n=\dfrac{3n^2-2n+1}{n^2-1000}$};
  \end{axis}
  \node [right] at (myplot.right of origin) {\scriptsize $n$};
  \node [above] at (myplot.above origin) {\scriptsize $y$};
 \end{tikzpicture}
\\(a)\\
\begin{tikzpicture}[alt={Blue dots for aₙ=cos n at 1≤n≤100; scatter dense between –1 and 1 showing no trend; axes n,y.}]
\begin{axis}[width=\marginparwidth,tick label style={font=\scriptsize},
axis y line=middle,axis x line=middle,name=myplot,axis on top,
xtick={20,40,60,80,100},ymin=-1.1,ymax=1.1,xmin=-.1,xmax=110]
\addplot [only marks,draw={\colorone},fill={\colorone},mark size={.75pt},domain=1:100,samples=100]
 {cos(deg(x))};
\end{axis}
\node [right] at (myplot.right of origin) {\scriptsize $n$};
\node [above] at (myplot.above origin) {\scriptsize $y$};
\node [] at (myplot.north) {\scriptsize $\ds a_n = \cos n$};
\end{tikzpicture}
\\(b)\\
\begin{tikzpicture}[alt={Blue dots for aₙ=(–1)^n/n at 1≤n≤20; alternating sign terms decay toward 0; axes n,y.}]
\begin{axis}[width=\marginparwidth,tick label style={font=\scriptsize},
axis y line=middle,axis x line=middle,name=myplot,axis on top,
ymin=-1.1,ymax=1.1,xmin=-.1,xmax=22]
\addplot [only marks,draw={\colorone},fill={\colorone},mark size={.75pt},domain=1:20,samples=20] {((-1)^x)/x};
\draw (axis cs:10,-.7) node {\scriptsize $\ds a_n = \frac{(-1)^n}{n}$};
\end{axis}
\node [right] at (myplot.right of origin) {\scriptsize $n$};
\node [above] at (myplot.above origin) {\scriptsize $y$};
\end{tikzpicture}
\\(c)}
%
\[
 \text{1. }\{a_n\} = \left\{\frac{3n^2-2n+1}{n^2-1000}\right\}\qquad
 \text{2. }\{a_n\} = \{\cos n \}\qquad
 \text{3. }\{a_n\} = \left\{\frac{(-1)^n}{n}\right\}
\]
\solution
\begin{enumerate}
\item		Using \autoref{thm:lim_rational_fn_at_infty}, we can state that $\ds\lim_{x\to\infty} \frac{3x^2-2x+1}{x^2-1000} = 3$. (We could have also directly applied L'Hôpital's Rule.) Thus the sequence $\{a_n\}$ converges, and its limit is 3. A scatter plot of every 5 values of $a_n$ is given in \autoref{fig:seq4} (a). The values of $a_n$ vary widely near $n=30$, ranging from about $-73$ to $125$, but as $n$ grows, the values approach 3.

\item		The limit $\ds\lim_{x\to\infty}\cos x$ does not exist, as $\cos x$ oscillates (and takes on every value in $[-1,1]$ infinitely many times). This means that we cannot apply \autoref{thm:seq_limit}. 

The fact that the cosine function oscillates strongly hints that $\cos n$, when $n$ is restricted to $\mathbb{N}$, will also oscillate. \autoref{fig:seq4} (b), where the sequence is plotted, shows that this is true. Because only discrete values of cosine are plotted, it does not bear strong resemblance to the familiar cosine wave.

Based on the graph, we suspect that $\ds \lim_{n\to\infty} a_n$ does not exist, but we have not decisively proven it yet.

%conclude that the sequence $\{\cos n\}$ diverges. (And in this particular case, since the domain is restricted to $\mathbb{N}$, no value of $\cos n$ is repeated!) This sequence is plotted in \autoref{fig:seq4} (b); because only discrete values of cosine are plotted, it does not bear strong resemblance to the familiar cosine wave.

\item		We cannot actually apply \autoref{thm:seq_limit} here, as the function $f(x) = (-1)^x/x$ is not well defined. (What does $(-1)^{\sqrt{2}}$ mean? In actuality, there is an answer, but it involves \emph{complex analysis}, beyond the scope of this text.) So for now we say that we cannot determine the limit. (But we will be able to very soon.) By looking at the plot in \autoref{fig:seq4} (c), we would like to conclude that the sequence converges to 0. That is true, but at this point we are unable to decisively say so.
\end{enumerate}
\end{example}

It seems that  %$\ds \left\{\frac{(-1)^n}{n}\right\}$ 
$\{(-1)^n/n\}$ converges to 0 but we lack the formal tool to prove it. The following theorem gives us that tool.

\begin{theorem}[Absolute Value Theorem]\label{thm:abs_val_seq}%
Let $\{a_n\}$ be a sequence. If $\ds \lim_{n\to\infty}\abs{a_n}= 0$, then $\ds \lim_{n\to\infty} a_n = 0$\index{Absolute Value Theorem}\index{limit!Absolute Value Theorem}\index{sequence!Absolute Value Theorem}
\end{theorem}

% this proof is exercise 45, but copying the proof is not the worst learning outcome
\begin{proof}
We know $-\abs{a_n}\leq a_n\leq\abs{a_n}$ and $\ds \lim_{n\to \infty} (-\abs{a_n})=-\lim_{n\to\infty}\abs{a_n}=0$. Thus by the Squeeze Theorem, $\ds \lim_{n\to\infty} a_n =0$.
\end{proof}

\begin{example}[Determining the convergence / divergence of a sequence]\label{ex_seq5}%
Determine the convergence or divergence of the following sequences.
\[
 \text{1. }\{a_n\} = \left\{\frac{(-1)^n}{n}\right\}\qquad
 \text{2. }\{a_n\} = \left\{\frac{(-1)^n(n+1)}{n}\right\}
\]
\solution
\begin{enumerate}
\item		This appeared in \autoref{ex_seq4}. We want to apply \autoref{thm:abs_val_seq}, so consider the limit of $\{\abs{a_n}\}$:
\begin{align*}
\lim_{n\to\infty}\abs{a_n}&= \lim_{n\to\infty} \abs{\frac{(-1)^n}{n}} \\
					&= \lim_{n\to\infty} \frac1n \\
					&= 0.
\end{align*}
Since this limit is 0, we can apply \autoref{thm:abs_val_seq} and state that $\ds\lim_{n\to\infty} a_n=0$.

\item Because of the alternating nature of this sequence (i.e., every other term is multiplied by $-1$), we cannot simply look at the limit $\ds \lim_{x\to\infty} \frac{(-1)^x(x+1)}{x}$. We can try to apply the techniques of \autoref{thm:abs_val_seq}:
\begin{align*}
	\lim_{n\to\infty}\abs{a_n}
	&= \lim_{n\to\infty} \abs{\frac{(-1)^n(n+1)}{n}} \\
	&= \lim_{n\to\infty} \frac{n+1}{n}\\
	&= 1.
\end{align*}
%
\mtable{A plot of a sequence in \autoref{ex_seq5}, part 2.}{fig:seq5}{\begin{tikzpicture}[alt={Blue dots for aₙ = (−1)^n·(n + 1)/n with 1 < n < 20; +ve drop from 2 to 1, -ve rise −2 to −1; alternating toward 1 on y-axis.}]
\begin{axis}[width=\marginparwidth,tick label style={font=\scriptsize},
axis y line=middle,axis x line=middle,name=myplot,axis on top,ytick={-1,-2,1,2},
ymin=-2.5,ymax=2.5,xmin=-.1,xmax=22]
\addplot [only marks,draw={\colorone},fill={\colorone},mark size={1.5pt},domain=1:20,samples=20]
 {((-1)^x*(x+1))/x};
\draw (axis cs:15,1.9) node {\scriptsize $\ds a_n = \frac{(-1)^n(n+1)}{n}$};
\end{axis}
\node [right] at (myplot.right of origin) {\scriptsize $n$};
\node [above] at (myplot.above origin) {\scriptsize $y$};
\end{tikzpicture}}%
%
We have concluded that when we ignore the sign, the sequence approaches 1. This means we cannot apply \autoref{thm:abs_val_seq}; it states the the limit must be 0 in order to conclude anything.

Since we know that the signs of the terms alternate \emph{and} we know that the limit of $\abs{a_n}$ is 1, we know that as $n$ approaches infinity, the terms will alternate between values close to 1 and $-1$, meaning the sequence diverges. A plot of this sequence is given in \autoref{fig:seq5}.
\end{enumerate}
\end{example}

We continue our study of the limits of sequences by considering some of the properties of these limits.

\begin{theorem}[Properties of the Limits of Sequences]\label{thm:seq_properties}%
Let $\{a_n\}$ and $\{b_n\}$ be sequences such that $\ds \lim_{n\to\infty} a_n = L$ and $\ds \lim_{n\to\infty} b_n = K$, where $L$ and $K$ are real numbers, and let $c$ be a real number.\index{sequences!limit properties}
\begin{multicols}{2}
\begin{enumerate}\lxAddClass{columns2}
\item		$\ds \lim_{n\to\infty} (a_n\pm b_n) = L\pm K$
\item		$\ds \lim_{n\to\infty} (a_n\cdot b_n) = L\cdot K$
\item		$\ds \lim_{n\to\infty} (a_n/b_n) = L/K$, $K\neq 0$
\item		$\ds \lim_{n\to\infty} (c\cdot a_n) = c\cdot L$
\end{enumerate}
\end{multicols}
\end{theorem}

\begin{example}[Applying properties of limits of sequences]\label{ex_seq6}%
Let the following limits be given:
\begin{itemize}
\item	 	$\ds \lim_{n\to\infty} a_n = 0$;
\item		$\ds \lim_{n\to\infty} b_n = e$; and
\item	  $\ds \lim_{n\to\infty} c_n = 5$.
\end{itemize}
Evaluate the following limits.
\[
 \text{1.}\quad\lim_{n\to\infty} (a_n+b_n)\qquad
 \text{2.}\quad\lim_{n\to\infty} (b_n\cdot c_n)\qquad
 \text{3.}\quad\lim_{n\to\infty} (1000\cdot a_n)
\]
\solution
We will use \autoref{thm:seq_properties} to answer each of these.
\begin{enumerate} 
\item		Since $\ds \lim_{n\to\infty} a_n = 0$ and $\ds \lim_{n\to\infty} b_n = e$, we conclude that $\ds \lim_{n\to\infty} (a_n+b_n) = 0+e = e.$ So even though we are adding something to each term of the sequence $b_n$, we are adding something so small that the final limit is the same as before.

\item		Since $\ds \lim_{n\to\infty} b_n = e$ and $\ds \lim_{n\to\infty} c_n = 5$, we conclude that $\ds \lim_{n\to\infty} (b_n\cdot c_n) = e\cdot 5 = 5e.$

\item		Since $\ds \lim_{n\to\infty} a_n = 0$, we have $\ds \lim_{n\to\infty} 1000a_n =1000\cdot 0 = 0$. It does not matter that we multiply each term by 1000; the sequence still approaches 0. (It just takes longer to get close to 0.)
\end{enumerate}
\end{example}

% todo write a transition paragraph to geometric sequences

\begin{definition}[Geometric Sequence]\label{def:geom_seq}%
For a constant $r$, the sequence $\{r^n\}$ is known as a \textbf{geometric sequence}.\index{geometric sequence}
\end{definition}

\begin{theorem}[Convergence of Geometric Sequences]\label{thm:geom_seq}%
The sequence $\{r^n\}$ is convergent if $-1<r\leq 1$ and divergent for all other values of $r$.  Furthermore,
\[
\lim_{n\to \infty} r^n=\begin{cases} 
0&  -1<r<1\\
1& r=1
\end{cases}\]
\end{theorem}

\begin{proof}
We can see from \autoref{ki_exp_func_props} and by letting $a=r$ that
\[
\lim_{n\to \infty} r^n =
\begin{cases}
\infty &  r>1\\
0 & 0<r<1.
\end{cases}
\]
We also know that $\ds  \lim_{x\to \infty} 1^n=1$ and $\ds  \lim_{x\to \infty} 0^n=0$. If $-1<r<0$, we know $0<\abs r<1$ so $\ds \lim_{x\to \infty}\abs{r^n}=\lim_{x\to \infty}\abs r^n=0$ and thus by \autoref{thm:abs_val_seq},$\ds \lim_{x\to \infty} r^n=0$. If $r\le-1$, $\ds \lim_{x\to \infty} r^n$ does not exist. Therefore, the sequence $\{ r^n\}$ is convergent if $-1<r\leq 1$ and divergent for all other values of $r$.
\end{proof}

% todo add a geometric sequence example

% todo use more than just color to distinguish a_n from S_n here throughout chapter

There is more to learn about sequences than just their limits. We will also study their range and the relationships terms have with the terms that follow. We start with some definitions describing properties of the range.

\begin{definition}[Bounded and Unbounded Sequences]\label{def:bounded}%
A sequence $\{a_n\}$ is said to be \textbf{bounded} if there exist real numbers $m$ and $M$ such that $m < a_n < M$ for all $n$ in $\mathbb{N}$.\bigskip

A sequence $\{a_n\}$ is said to be \textbf{unbounded} if it is not bounded.\bigskip

A sequence $\{a_n\}$ is said to be \textbf{bounded above} if there exists an $M$ such that $a_n < M$ for all $n$ in $\mathbb{N}$; it is \textbf{bounded below} if there exists an $m$ such that $m<a_n$ for all $n$ in $\mathbb{N}$.
\index{sequences!boundedness}\index{bounded sequence}\index{unbounded sequence}
\end{definition}

It follows from this definition that an unbounded sequence may be bounded above or bounded below; a sequence that is both bounded above and below is simply a bounded sequence.

\begin{example}[Determining boundedness of sequences]\label{ex_seq3}%
Determine the boundedness of the following sequences.
\[
 \text{1.}\quad\{a_n\} = \left\{\frac1n\right\}\qquad\qquad
 \text{2.}\quad\{a_n\} = \{2^n\}
\]
%
\mtable{A plot of $\{a_n\} = \{1/n\}$ and $\{a_n\} = \{2^n\}$ from \autoref{ex_seq3}.}{fig:seq3}{%
\begin{tikzpicture}[alt={Blue dots: aₙ=1/n for 1≤n≤10; values decline from 1 to near 0 along the n-axis.}]
\begin{axis}[width=\marginparwidth,tick label style={font=\scriptsize},
axis y line=middle,axis x line=middle,name=myplot,axis on top,xtick={1,...,10},
ytick={1},extra y ticks={.5,.25,.1},extra y tick labels={$1/2$,$1/4$,$1/10$},
ymin=-.1,ymax=1.1,xmin=-.1,xmax=11]
\addplot [only marks,draw={\colorone},fill={\colorone},mark size={1.75pt},domain=1:10,samples=10]  {1/x};
\draw (axis cs:6,.75) node {\scriptsize $\ds a_n = \frac{1}{n}$};
\end{axis}
\node [right] at (myplot.right of origin) {\scriptsize $n$};
\node [above] at (myplot.above origin) {\scriptsize $y$};
\end{tikzpicture}
\\(a)\\[10pt]
\begin{tikzpicture}[alt={Blue dots: aₙ=2^n for 1≤n≤9; points soar from 2 up to about 256 on the y-axis.}]
\begin{axis}[width=\marginparwidth,tick label style={font=\scriptsize},
axis y line=middle,axis x line=middle,name=myplot,axis on top,
ymin=-.1,ymax=275,xmin=-.1,xmax=9]
\addplot [only marks,draw={\colorone},fill={\colorone},mark size={1.75pt},domain=1:8,samples=8]  {2^x};
\draw (axis cs:4,150) node {\scriptsize $\ds a_n = 2^{n}$};
\end{axis}
\node [right] at (myplot.right of origin) {\scriptsize $n$};
\node [above] at (myplot.above origin) {\scriptsize $y$};
\end{tikzpicture}
\\(b)}
\solution
\begin{enumerate}
\item	The terms of this sequence are always positive but are decreasing, so we have $0<a_n<2$ for all $n$. Thus this sequence is bounded. \autoref{fig:seq3}(a) illustrates this.

\item	The terms of this sequence obviously grow without bound. However, it is also true that these terms are all positive, meaning $0<a_n$. Thus we can say the sequence is unbounded, but also bounded below. \autoref{fig:seq3}(b) illustrates this.
\end{enumerate}
\end{example}

The previous example produces some interesting concepts. First, we can recognize that the sequence $\ds\left\{1/n\right\}$ converges to 0. This says, informally, that ``most'' of the terms of the sequence are ``really close'' to 0. This implies that the sequence is bounded, using the following logic. First, ``most'' terms are near 0, so we could find some sort of bound on these terms (using \autoref{def:seq_limit}, the bound is $\epsilon$). That leaves a ``few'' terms that are not near 0 (i.e., a \emph{finite} number of terms). A finite list of numbers is always bounded. 

This logic suggests that if a sequence converges, it must be bounded. This is indeed true, as stated by the following theorem.

\begin{theorem}[Convergent Sequences are Bounded]\label{thm:converge_bounded}%
Let $\ds \left\{a_n\right\}$ be a convergent sequence. Then $\{a_n\}$ is bounded.
\index{bounded sequence!convergence}\index{convergence!of sequence}\index{sequences!convergent}
\end{theorem}

\mnote[\baselineskip]{\textbf{Note:} Keep in mind what \autoref{thm:converge_bounded} does \emph{not} say. It does not say that bounded sequences must converge, nor does it say that if a sequence does not converge, it is not bounded.}

In \autoref{ex_LHR4} part 1, we found that $\ds \lim_{x\to \infty} (1+1/x)^x=e$. If we consider the sequence $\ds \{b_n\}=\{(1+1/n)^n\}$, we see that $\ds \lim_{n\to \infty}b_n=e$. Even though it may be difficult to intuitively grasp the behavior of this sequence, we know immediately that it is bounded.

Another interesting concept to come out of \autoref{ex_seq3} again involves the sequence $\{1/n\}$. We stated, without proof, that the terms of the sequence were decreasing. That is, that $a_{n+1} < a_n$ for all $n$. (This is easy to show. Clearly $n < n+1$. Taking reciprocals flips the inequality: $1/n > 1/(n+1)$. This is the same as $a_n > a_{n+1}$.) Sequences that either steadily increase or decrease are important, so we give this property a name.

\begin{definition}[Monotonic Sequences]\label{def:monotonic}%
\mbox{}\\[-2\baselineskip]\index{sequences!monotonic}\index{monotonic sequence}\begin{enumerate}
\item		A sequence $\{a_n\}$ is \textbf{monotonically increasing} if $a_n \leq a_{n+1}$ for all $n$, i.e.,
\[a_1 \leq a_2 \leq a_3 \leq \dotsb \leq a_n \leq a_{n+1} \dotsb\]
 \item	A sequence $\{a_n\}$ is \textbf{monotonically decreasing} if $a_n \geq a_{n+1}$ for all $n$, i.e.,
\[a_1 \geq a_2 \geq a_3 \geq \dotsb\geq a_n \geq a_{n+1} \dotsb\]
 \item	A sequence is \textbf{monotonic} if it is monotonically increasing or monotonically decreasing.
 \end{enumerate}
\end{definition}

\mnote{\textbf{Note:} It is sometimes useful to call a monotonically increasing sequence \emph{strictly increasing} if $a_n < a_{n+1}$ for all $n$; i.e, we remove the possibility that subsequent terms are equal.\bigskip\\
A similar statement holds for \emph{strictly decreasing.}}

\begin{example}[Determining monotonicity]\label{ex_seq7}%
Determine the monotonicity of the following sequences.
\begin{multicols}{2}
\begin{enumerate}\lxAddClass{columns2}
\item $\ds \{a_n\} = \left\{\frac{n+1}n\right\}$
\item	$\ds \{a_n\} = \left\{\frac{n^2+1}{n+1}\right\}$	
\item $\ds \{a_n\} = \left\{\frac{n^2-9}{n^2-10n+26}\right\}$
\item	$\ds \{a_n\} = \left\{\frac{n^2}{n!}\right\}$	
\end{enumerate}
\end{multicols}
\solution
In each of the following, we will examine $a_{n+1}-a_n$. If $a_{n+1}-a_n \ge0$, we conclude that $a_n\le a_{n+1}$ and hence the sequence is increasing. If $a_{n+1}-a_n\le0$, we conclude that $a_n\ge a_{n+1}$ and the sequence is decreasing. Of course, a sequence need not be monotonic and perhaps neither of the above will apply.

We also give a scatter plot of each sequence. These are useful as they suggest a pattern of monotonicity, but analytic work should be done to confirm a graphical trend.

\begin{enumerate}
\item \hfill$\begin{aligned}[t]
a_{n+1}-a_n &= \frac{n+2}{n+1} - \frac{n+1}{n} \\		
	&= \frac{(n+2)(n)-(n+1)^2}{(n+1)n} \\
	&=	\frac{-1}{n(n+1)} \\
	&<0 \quad\text{ for all $n$.}
\end{aligned}$\hfill\null\\
%
\mtable{A plot of $\{a_n\}=\{\frac{n+1}n\}$ in \autoref{ex_seq7}(a).}{fig:seq7a}{\begin{tikzpicture}[alt={Blue dots: aₙ=(n+1)/n for 1≤n≤11; sequence slips from 2 toward 1.}]
\begin{axis}[width=1.16\marginparwidth,tick label style={font=\scriptsize},
axis y line=middle,axis x line=middle,name=myplot,axis on top,ytick={-1,-2,1,2},
ymin=-.1,ymax=2.5,xmin=-1,xmax=11]
\addplot [only marks,draw={\colorone},fill={\colorone},mark size={1.5pt},domain=1:10,samples=10] {(x+1)/x};
\draw (axis cs:6,1.9) node {\scriptsize $\ds a_n = \frac{n+1}{n}$};
\end{axis}
\node [right] at (myplot.right of origin) {\scriptsize $n$};
\node [above] at (myplot.above origin) {\scriptsize $y$};
\end{tikzpicture}}
%
Since $a_{n+1}-a_n<0$ for all $n$, we conclude that the sequence is decreasing.

\item
%
\mtable{A plot of $\{a_n\}=\{\frac{n^2+1}{n+1}\}$ in \autoref{ex_seq7}(b).}{fig:seq7b}{\begin{tikzpicture}[alt={Blue dots: aₙ=(n²+1)/(n+1) for 1≤n≤11; roughly linear rise, y climbs 2 towards 10.}]
\begin{axis}[width=1.16\marginparwidth,tick label style={font=\scriptsize},
axis y line=middle,axis x line=middle,name=myplot,axis on top,
ymin=-1,ymax=11,xmin=-1,xmax=11]
\addplot [only marks,draw={\colorone},fill={\colorone},mark size={1.5pt},domain=1:10,samples=10]
 {(x^2+1)/(x+1)};
\draw (axis cs:7.5,1.9) node {\scriptsize $\ds a_n = \frac{n^2+1}{n+1}$};
\end{axis}
\node [right] at (myplot.right of origin) {\scriptsize $n$};
\node [above] at (myplot.above origin) {\scriptsize $y$};
\end{tikzpicture}}
%
\hfill$\begin{aligned}[t]
	a_{n+1}-a_n &= \frac{(n+1)^2+1}{n+2} - \frac{n^2+1}{n+1} \\		
	&= \frac{\bigl((n+1)^2+1\bigr)(n+1)- (n^2+1)(n+2)}{(n+1)(n+2)}\\
	&=	\frac{n(n+3)}{(n+1)(n+2)} \\
	&> 0 \quad \text{ for all $n$.}
\end{aligned}$\hfill\null

Since $a_{n+1}-a_n>0$ for all $n$, we conclude the sequence is increasing.

\item%
%
\mtable{A plot of $\{a_n\}=\{\frac{n^2-9}{n^2-10n+26}\}$ in \autoref{ex_seq7}(c).}{fig:seq7c}{\begin{tikzpicture}[alt={Blue dots: aₙ=(n²−9)/(n²−10n+26) for 1<n<11; peaks near n=3 then tapers.}]
\begin{axis}[width=\marginparwidth,tick label style={font=\scriptsize},
axis y line=middle,axis x line=middle,name=myplot,axis on top,
ymin=-1,ymax=16,xmin=-1,xmax=11]
\addplot [only marks,draw={\colorone},fill={\colorone},mark size={1.5pt},domain=1:10,samples=10]
 {(x^2-9)/(x^2-10*x+26)};
\draw (axis cs:3.4,11) node {\scriptsize $\ds a_n = \frac{n^2-9}{n^2-10n+26}$};
\end{axis}
\node [right] at (myplot.right of origin) {\scriptsize $n$};
\node [above] at (myplot.above origin) {\scriptsize $y$};
\end{tikzpicture}}%
%
	We can clearly see in \autoref{fig:seq7c}, where the sequence is plotted, that it is not monotonic. However, it does seem that after the first 4 terms it is decreasing. To understand why, perform the same analysis as done before:
\begin{align*}
	a_{n+1}-a_n &= \frac{(n+1)^2-9}{(n+1)^2-10(n+1)+26} - \frac{n^2-9}{n^2-10n+26} \\	
	&= \frac{n^2+2n-8}{n^2-8n+17}-\frac{n^2-9}{n^2-10n+26}\\
	&= \frac{(n^2+2n-8)(n^2-10n+26)-(n^2-9)(n^2-8n+17)}{(n^2-8n+17)(n^2-10n+26)}\\
	&= \frac{-10n^2+60n-55}{(n^2-8n+17)(n^2-10n+26)}.
\end{align*}

We want to know when this is greater than, or less than, 0. The denominator is always positive, therefore we are only concerned with the numerator. Using the quadratic formula, we can determine that $-10n^2+60n-55=0$ when $n\approx 1.13, 4.87$. So for $n<1.13$, the sequence is decreasing. Since we are only dealing with the natural numbers, this means that $a_1 > a_2$.

Between $1.13$ and $4.87$, i.e., for $n=2$, 3 and 4, we have that $a_{n+1}>a_n$ and the sequence is increasing. (That is, when $n=2$, 3 and 4, the numerator $-10n^2+60n-55$ from the fraction above is $>0$.)

When $n> 4.87$, i.e, for $n\geq 5$, we have that $-10n^2+60n-55<0$, hence $a_{n+1}-a_n<0$, so the sequence is decreasing.

In short, the sequence is simply not monotonic. However, it is useful to note that for $n\geq 5$, the sequence is monotonically decreasing. 

\item
%
\mtable{A plot of $\{a_n\}=\{n^2/n!\}$ in \autoref{ex_seq7}(d).}{fig:seq7d}{\begin{tikzpicture}[alt={Blue dots: aₙ=n²/n^! for 1≤n≤11; rapid drop from 2 to near 0.}]
\begin{axis}[width=1.16\marginparwidth,tick label style={font=\scriptsize},
axis y line=middle,axis x line=middle,name=myplot,axis on top,
ymin=-.1,ymax=2.1,xmin=-1,xmax=11]
\addplot [only marks,draw={\colorone},fill={\colorone},mark size={1.5pt},domain=1:10,samples=10]
 {(x^2)/(factorial(x))};
\draw (axis cs:7,1.25) node {\scriptsize $\ds a_n = \frac{n^2}{n!}$};
\end{axis}
\node [right] at (myplot.right of origin) {\scriptsize $n$};
\node [above] at (myplot.above origin) {\scriptsize $y$};
\end{tikzpicture}}
%
	Again, the plot in \autoref{fig:seq7d} shows that the sequence is not monotonic, but it suggests that it is monotonically decreasing after the first term. Instead of looking at $a_{n+1}-a_n$, this time we'll look at $a_n/a_{n+1}$:
\begin{align*}
	\frac{a_n}{a_{n+1}} &= \frac{n^2}{n!}\frac{(n+1)!}{(n+1)^2} \\
	&= \frac{n^2}{n+1} \\
	&= n-1+\frac1{n+1}
\end{align*}
When $n=1$, the above expression is $<1$; for $n\geq 2$, the above expression is $>1$. Thus this sequence is not monotonic, but it is monotonically decreasing after the first term.
\end{enumerate}
\end{example}

Knowing that a sequence is monotonic can be useful. In particular, if we know that a sequence is bounded and monotonic, we can conclude it converges. Consider, for example, a sequence that is monotonically decreasing and is bounded below. We know the sequence is always getting smaller, but that there is a bound to how small it can become. This is enough to prove that the sequence will converge, as stated in the following theorem.

\begin{theorem}[Bounded Monotonic Sequences are Convergent]\label{thm:monotonic_converge}%
Let $\{a_n\}$ be a bounded, monotonic sequence. Then $\{a_n\}$ converges; i.e., $\ds \lim_{n \to\infty}a_n$ exists.
\index{sequences!convergent}\index{convergence!of monotonic sequences}
\end{theorem}

Consider once again the sequence $\{a_n\} = \{1/n\}$. It is easy to show it is monotonically decreasing and that it is always positive (i.e., bounded below by 0). Therefore we can conclude by \autoref{thm:monotonic_converge} that the sequence converges. We already knew this by other means, but in the following section this theorem will become very useful.

Convergence of a sequence does not depend on the first $N$ terms of a sequence.  For example, we could adapt the sequence of the previous paragraph to be
\[
 1,\ 10,\ 100,\ 1000,\ \frac15,\ \frac16,\ \frac17,\ \frac18,\ \frac19,\ \frac1{10},
 \dotsc
\]
Because we only changed three of the first 4 terms, we have not affected whether the sequence converges or diverges.

Sequences are a great source of mathematical inquiry. The On-Line Encyclopedia of Integer Sequences (\url{http://oeis.org}) contains thousands of sequences and their formulae. (As of this writing, there are 348,000 sequences in the database.) Perusing this database quickly demonstrates that a single sequence can represent several different ``real life'' phenomena. 

Interesting as this is, our interest actually lies elsewhere. We are more interested in the \emph{sum} of a sequence. That is, given a sequence $\{a_n\}$, we are very interested in $a_1+a_2+a_3+\dotsb$. Of course, one might immediately counter with ``Doesn't this just add up to `infinity'?'' Many times, yes, but there are many important cases where the answer is no. This is the topic of \emph{series}, which we begin to investigate in the next section.

\printexercises{exercises/08-01-exercises}

\section{Infinite Series}\label{sec:series}

Given the sequence $\{a_n\} = \{1/2^n\} = 1/2,\ 1/4,\ 1/8,\ \dotsc$, consider the following sums:\vspace{-.5\baselineskip}
\[
\begin{array}{ c c c c c }
a_1				&=& 1/2					&=& 1/2\\
a_1+a_2			&=& 1/2+1/4				&=& 3/4\\
a_1+a_2+a_3		&=& 1/2+1/4+1/8			&=& 7/8\\
a_1+a_2+a_3+a_4	&=& 1/2+1/4+1/8+1/16		&=& 15/16
\end{array}
\]
Later, we will be able to show that\vspace{-.5\baselineskip}
\[a_1+a_2+a_3+\dotsb+a_n = \frac{2^n-1}{2^n} = 1-\frac{1}{2^n}.\]
Let $S_n$ be the sum of the first $n$ terms of the sequence $\{1/2^n\}$. From the above, we see that $S_1=1/2$, $S_2 = 3/4$, and that $S_n = 1-1/2^n$. 

Now consider the following limit: $\ds \lim_{n\to\infty}S_n = \lim_{n\to\infty}\bigl(1-1/2^n\bigr) = 1$. This limit can be interpreted as saying something amazing: \emph{the sum of \emph{all} the terms of the sequence $\{1/2^n\}$ is 1.} 


This example illustrates some interesting concepts that we explore in this section. We begin this exploration with some definitions.

{\tcbset{grow to right by=5em}
\begin{definition}[Infinite Series, $n^{\text{th}}$ Partial Sums, Convergence, Divergence]\label{def:series}%
Let $\{a_n\}$ be a sequence.
\index{series!definition}\index{series!partial sums}\index{series!convergent}\index{series!divergent}\index{convergence!of series}\index{divergence!of series}
\begin{enumerate}
\item		The sum $\ds \sum_{n=1}^\infty a_n$ is an \textbf{infinite series} (or, simply \textbf{series}).
\item		Let $\ds S_n = \sum_{i=1}^n a_i$\,; the sequence $\{S_n\}$ is the sequence of \textbf{$n^{\text{th}}$ partial sums} of $\{a_n\}$.\vspace{-.5\baselineskip}
\item		If the sequence $\{S_n\}$ converges to $L$, we say the series $\ds \sum_{n=1}^\infty a_n$ \textbf{converges} to $L$, and we write $\ds \sum_{n=1}^\infty a_n = L$.\vspace{-.5\baselineskip}
\item		If the sequence $\{S_n\}$ diverges, the series $\ds \sum_{n=1}^\infty a_n$ \textbf{diverges}.
\end{enumerate}
\end{definition}}

Using our new terminology, we can state that the series $\ds \sum_{n=1}^\infty 1/2^n$\vspace{-.5\baselineskip} converges, and $\ds \sum_{n=1}^\infty 1/2^n = 1.$

\youtubeVideo{cyoiIBs7kIg}{Finding a Formula for a Partial Sum of a Telescoping Series}

We will explore a variety of series in this section. We start with two series that diverge, showing how we might discern divergence.

\begin{example}[Showing series diverge]\label{ex_series1}%
\mbox{}\\[-2\baselineskip]\parbox[t]{\linewidth}{%
\begin{enumerate}
\item		Let $\{a_n\} = \{n^2\}$. Show $\ds \sum_{n=1}^\infty a_n$ diverges.
\item		Let $\{b_n\} = \{(-1)^{n+1}\}$. Show $\ds \sum_{n=1}^\infty b_n$ diverges.
\end{enumerate}}\vspace{0pt}
\solution
\begin{enumerate}
\item	Consider $S_n$, the $n^{\text{th}}$ partial sum.
%
\mtable{Scatter plots relating to the series of \autoref{ex_series1} part 1.}{fig:series1a}{\begin{tikzpicture}[alt={Blue aₙ rise gently to 100 while red Sₙ surge to almost 340; scatter for 1≤n≤10, y-axis 0 < y < 350.}]
\begin{axis}[width=1.16\marginparwidth,tick label style={font=\scriptsize},
axis y line=middle,axis x line=middle,name=myplot,axis on top,
xmin=-1,xmax=11,ymin=-10,legend columns=-1,legend style={at={(.5,-.1)},anchor=north}]%ymax=2.1,%
\addplot [only marks,draw={\colorone},fill={\colorone},mark size={1.5pt},domain=1:10,samples=10] {x^2};
\addlegendentry{\scriptsize$a_n$\quad\mbox{}}
\addplot [only marks,draw={\colortwo},fill={\colortwo},mark size={1.5pt},domain=1:10,samples=10] {x*(x+1)*(2*x+1)/6};
\addlegendentry{\scriptsize$S_n$}
\end{axis}
\node [right] at (myplot.right of origin) {\scriptsize $n$};
\node [above] at (myplot.above origin) {\scriptsize $y$};
\end{tikzpicture}}
%
\begin{align*}
	S_n &= a_1+a_2+a_3+\dotsb+a_n \\		
	&= 1^2+2^2+3^2\dotsb+ n^2 \\
	&= \frac{n(n+1)(2n+1)}{6}.
	&\iflatexml\qquad\fi&\text{by \autoref{thm:summation}}
\end{align*}
Since $\ds \lim_{n\to\infty}S_n = \infty$, we conclude that the series $\ds \sum_{n=1}^\infty n^2$ diverges. It is instructive to write $\ds \sum_{n=1}^\infty n^2=\infty$ for this tells us \emph{how} the series diverges: it grows without bound.

A scatter plot of the sequences $\{a_n\}$ and $\{S_n\}$ is given in \autoref{fig:series1a}. The terms of $\{a_n\}$ are growing, so the terms of the partial sums $\{S_n\}$ are growing even faster, illustrating that the series diverges.

\item	The sequence $\{b_n\}$ starts with $1, -1, 1, -1, \dotsc$. Consider some of the partial sums $S_n$ of $\{b_n\}$:
\begin{align*}
S_1 &= 1\\
S_2 &= 0\\
S_3 &= 1\\
S_4 &= 0
\end{align*}
This pattern repeats; we find that
$S_n = \begin{cases}
1  & n\ \text{ is odd}\\
0  & n\ \text{ is even.}
\end{cases}$ \smallskip
As $\{S_n\}$ oscillates, repeating $1, 0, 1, 0, \dotsc$, we conclude that $\ds\lim_{n\to\infty}S_n$ does not exist, hence $\ds\sum_{n=1}^\infty (-1)^{n+1}$ diverges.		

\mtable{Scatter plots relating to the series of \autoref{ex_series1} part 2.}{fig:series1b}{\begin{tikzpicture}[alt={Blue bₙ dots alternate ±1; red Sₙ alternate 0 and 1; scatter for 1≤n≤10, y-axis –1 < y < 1.}]
\begin{axis}[width=\marginparwidth,tick label style={font=\scriptsize},
axis y line=middle,axis x line=middle,name=myplot,axis on top,
ymin=-1.1,ymax=1.1,xmin=-1,xmax=11,legend columns=-1,legend style={at={(.5,-.1)},anchor=north}]
\addplot [only marks,draw={\colorone},fill={\colorone},mark size={2.5pt},domain=1:10,samples=10] {-(-1)^x};
\addlegendentry{\scriptsize$b_n$\quad\mbox{}}
\addplot [only marks,draw={\colortwo},fill={\colortwo},mark size={1.5pt},domain=1:10,samples=10] {(1-(-1)^x)/2};
\addlegendentry{\scriptsize$S_n$}
\end{axis}
\node [right] at (myplot.right of origin) {\scriptsize $n$};
\node [above] at (myplot.above origin) {\scriptsize $y$};
\end{tikzpicture}}

A scatter plot of the sequence $\{b_n\}$ and the partial sums $\{S_n\}$ is given in \autoref{fig:series1b}. When $n$ is odd, $b_n = S_n$ so the marks for $b_n$ are drawn oversized to show they coincide.
\end{enumerate}
\end{example}

While it is important to recognize when a series diverges, we are generally more interested in the series that  converge. In this section we will demonstrate a few general techniques for determining convergence; later sections will delve deeper into this topic.

\subsection{Geometric Series}

One important type of series is a \emph{geometric series}.

\begin{definition}[Geometric Series]\label{def:geom_series}%
A \textbf{geometric series} is a series of the form 
\[\sum_{n=0}^\infty ar^n = a+ar+ar^2+ar^3+\dotsb+ar^n+\dotsb\]
Note that the index starts at $n=0$. If the index starts at $n=1$ we have $\ds \sum_{n=1}^\infty ar^{n-1}$.%
\index{series!geometric}\index{geometric series}
\end{definition}

We started this section with a geometric series, although we dropped the first term of $1$. One reason geometric series are important is that they have nice convergence properties.

\begin{theorem}[Convergence of Geometric Series]\label{thm:geom_series}%
Consider the geometric series $\ds \sum_{n=0}^\infty ar^n$ where $a\neq0$.\vspace{-.5\baselineskip}
\index{series!geometric}\index{geometric series}\index{convergence!of geometric series}\index{divergence!of geometric series}
\begin{enumerate}
\item	If $r\neq1$, the $n^{\text{th}}$ partial sum is: $\ds S_n = \sum_{k=0}^{n-1}ar^k=\frac{a(1-r\,^n)}{1-r}$.
\item	The series converges if, and only if, $\abs r<1$. When $\abs r<1$, 
\[\sum_{n=0}^\infty ar^n = \frac{a}{1-r}.\]
\end{enumerate}
\end{theorem}

\begin{proof}
If $r=1$, then $S_n=a+a+a+\dotsb+a=na$. Since $\ds\lim_{n\to \infty} S_n=\pm \infty$, the geometric series diverges.\\
If $r\neq 1$, we have
\[S_n=a+ar+ar^2+\dotsb+ar^{n-1}.\]
Multiply each term by $r$ and we have 
\[rS_n=ar+ar^2+ar^3\dotsb+ar^n.\]
Subtract these two equations and solve for $S_n$.
\begin{align*}
S_n-rS_n &=a-ar^n \\
S_n &=\frac{a(1-r^n)}{1-r}\\
\end{align*}
From \autoref{thm:geom_seq}, we know that if $-1<r<1$, then $\ds \lim_{n\to \infty} r^n=0$ so
\[
\lim_{n\to \infty} S_n=\lim_{n\to \infty}=\frac{a(1-r^n)}{1-r}
=\frac{a}{1-r}- \frac{a}{1-r}\lim_{n\to \infty}r^n=\frac{a}{1-r}.
\]
So when $\abs r<1$ the geometric series converges and its sum is $\ds \frac{a}{1-r}$.

If either $r\leq -1$ or $r>1$, the sequence $\{r^n\}$ is divergent by \autoref{thm:geom_seq}. Thus $\ds \lim_{n\to \infty} S_n$ does not exist, so the geometric series diverges if $r\leq -1$ or \mbox{$r>1$.} % prevent a line break
\end{proof}

According to \autoref{thm:geom_series}, the series 
\[
\ds\sum_{n=0}^\infty \frac{1}{2^n}
=\sum_{n=0}^\infty \left(\frac 12\right)^n= 1+\frac12+\frac14+\dotsb
\]
converges as $r=1/2$, and $\ds \sum_{n=0}^\infty \frac{1}{2^n} = \frac{1}{1-1/2} = 2.$ This concurs with our introductory example; while there we got a sum of 1, we skipped the first term of 1.

\mtable{Scatter plots relating to the series in \autoref{ex_series2}.}{fig:series2}{%
\begin{tikzpicture}[alt={Blue aₙ decline toward 0; red Sₙ creep up to almost 2; scatter 2≤n≤10, y-axis –1 < y < 2.5.}]
\begin{axis}[width=\marginparwidth,tick label style={font=\scriptsize},
axis y line=middle,axis x line=middle,name=myplot,axis on top,xtick={2,4,6,8,10},
ymin=-.5,ymax=2.5,xmin=-1,xmax=11,legend columns=-1,legend style={at={(.5,-.1)},anchor=north}]
\addplot [only marks,draw={\colorone},fill={\colorone},mark size={1.5pt},domain=2:10,samples=9] {(.75)^x};
\addlegendentry{\scriptsize$a_n$\quad\mbox{}}
\addplot [only marks,draw={\colortwo},fill={\colortwo},mark size={1.5pt},domain=2:10,samples=9] {2.25-3*(.75)^x};
\addlegendentry{\scriptsize$S_n$}
\end{axis}
\node [right] at (myplot.right of origin) {\scriptsize $n$};
\node [above] at (myplot.above origin) {\scriptsize $y$};
\end{tikzpicture}
\smallskip\\(a)\bigskip\\
\begin{tikzpicture}[alt={Blue aₙ positive towards 0; red Sₙ flatten at almost 0.6; scatter 2≤n≤10, y-axis –0.5 < y < 1.}]
\begin{axis}[width=\marginparwidth,tick label style={font=\scriptsize},
axis y line=middle,axis x line=middle,name=myplot,axis on top,xtick={2,4,6,8,10},
ymin=-1.1,ymax=1.1,xmin=-1,xmax=11,legend columns=-1,legend style={at={(.5,0)},anchor=south}]
\addplot [only marks,draw={\colorone},fill={\colorone},mark size={1.5pt},domain=0:10,samples=11] {(-.5)^(x)};
\addlegendentry{\scriptsize$a_n$\quad\mbox{}}
\addplot [only marks,draw={\colortwo},fill={\colortwo},mark size={1.5pt},domain=0:10,samples=11] {(2+(-.5)^(x))/3};
\addlegendentry{\scriptsize$S_n$}
\end{axis}
\node [right] at (myplot.right of origin) {\scriptsize $n$};
\node [above] at (myplot.above origin) {\scriptsize $y$};
\end{tikzpicture}
\smallskip\\(b)\bigskip\\
\begin{tikzpicture}[alt={Blue aₙ explode past 800; red Sₙ climb past 1000; scatter 1≤n≤6, y-axis 0 < y < 1 100.}]
\begin{axis}[width=\marginparwidth,tick label style={font=\scriptsize},
axis y line=middle,axis x line=middle,name=myplot,axis on top,
xmin=-1,xmax=7,ymin=-50,legend columns=-1,legend style={at={(.5,-.1)},anchor=north}]%ymax=1.1,%
\addplot [only marks,draw={\colorone},fill={\colorone},mark size={1.5pt},domain=0:6,samples=7] {3^(x)};
\addlegendentry{\scriptsize$a_n$\quad\mbox{}}
\addplot [only marks,draw={\colortwo},fill={\colortwo},mark size={1.5pt},domain=0:6,samples=7] {(3^(x+1)-1)/2};
\addlegendentry{\scriptsize$S_n$}
\end{axis}
\node [right] at (myplot.right of origin) {\scriptsize $n$};
\node [above] at (myplot.above origin) {\scriptsize $y$};
\end{tikzpicture}
\smallskip\\(c)}

\begin{example}[Exploring geometric series]\label{ex_series2}%
Check the convergence of the following series. If the series converges, find its sum.
\[
 \text{1. }\sum_{n=2}^\infty \left(\frac34\right)^n\qquad
 \text{2. }\sum_{n=0}^\infty \left(\frac{-1}{2}\right)^n\qquad
 \text{3. }\sum_{n=0}^\infty 3^n
\]
\solution
\begin{enumerate}
\item		Since $r=3/4<1$, this series converges. By \autoref{thm:geom_series}, we have that
\[\sum_{n=0}^\infty \left(\frac34\right)^n = \frac{1}{1-3/4} = 4.\]
However, note the subscript of the summation in the given series: we are to start with $n=2$. Therefore we subtract off the first two terms, giving:
\[\sum_{n=2}^\infty \left(\frac34\right)^n = 4 - 1 - \frac34 = \frac94.\]
This is illustrated in \autoref{fig:series2}(a).

\item	Since $\abs r = 1/2 < 1$, this series converges, and by \autoref{thm:geom_series},
\[\sum_{n=0}^\infty \left(\frac{-1}{2}\right)^n = \frac{1}{1-(-1/2)} = \frac23.\]
The partial sums of this series are plotted in \autoref{fig:series2}(b). Note how the partial sums are not purely increasing as some of the terms of the sequence $\{(-1/2)^n\}$ are negative.

\item		Since $r>1$, the series diverges. (This makes ``common sense''; we expect the sum
\[1+3+9+27 + 81+243+\dotsb\]
to diverge.) This is illustrated in \autoref{fig:series2}(c).
\end{enumerate}
\end{example}

Later sections will provide tests by which we can determine whether or not a given series converges. This, in general, is much easier than determining \emph{what} a given series converges to. There are many cases, though, where the sum can be determined.

\begin{example}[Telescoping series]\label{ex_series3}%
Evaluate the sum~~$\ds \sum_{n=1}^\infty \left(\frac1n-\frac1{n+1}\right)$.
\index{series!telescoping}\index{telescoping series}
\solution
It will help to write down some of the first few partial sums of this series.
\begin{align*}
S_1 &=	\frac11-\frac12 & & = 1-\frac12\\
S_2 &=	\left(\frac11-\frac12\right) + \left(\frac12-\frac13\right) & & = 1-\frac13\\
S_3 &=	\left(\frac11-\frac12\right) + \left(\frac12-\frac13\right)+\left(\frac13-\frac14\right) & &= 1-\frac14\\
S_4 &=	\left(\frac11-\frac12\right) + \left(\frac12-\frac13\right)+\left(\frac13-\frac14\right) +\left(\frac14-\frac15\right)& &= 1-\frac15
\end{align*}
Note how most of the terms in each partial sum subtract out. In general, we see that $S_n = 1-\frac{1}{n+1}$. This means that the sequence $\{S_n\}$ converges,  as
\[\lim_{n\to\infty}S_n=\smallskip\lim_{n\to\infty}\left(1-\tfrac1{n+1}\right)=1,\]
and so we conclude that $\ds\sum_{n=1}^\infty \left(\tfrac1n-\tfrac1{n+1}\right) = 1$. Partial sums of the series are plotted in \autoref{fig:series3}.
\end{example}

\mtable[yshift=10\baselineskip]{Scatter plots relating to the series of \autoref{ex_series3}.}{fig:series3}{\begin{tikzpicture}[alt={Blue aₙ decay to 0; red Sₙ increase toward 0.9; scatter 1≤n≤10, y-axis 0 < y < 1.}]
\begin{axis}[width=1.16\marginparwidth,tick label style={font=\scriptsize},
axis y line=middle,axis x line=middle,name=myplot,axis on top,xtick={2,4,6,8,10},
ymin=-.1,ymax=1.1,xmin=-1,xmax=11,legend columns=-1,legend style={at={(.5,-.1)},anchor=north}]
\addplot [only marks,draw={\colorone},fill={\colorone},mark size={1.5pt},domain=1:10,samples=10]
 {1/x-1/(x+1)};
\addlegendentry{\scriptsize$a_n$\quad\mbox{}}
\addplot [only marks,draw={\colortwo},fill={\colortwo},mark size={1.5pt},domain=1:10,samples=10] {1-1/x};
\addlegendentry{\scriptsize$S_n$}
\end{axis}
\node [right] at (myplot.right of origin) {\scriptsize $n$};
\node [above] at (myplot.above origin) {\scriptsize $y$};
\end{tikzpicture}}%

The series in \autoref{ex_series3} is an example of a \textbf{telescoping series}. Informally, a telescoping series is one in which the partial sums reduce to just a fixed number of terms. The partial sum $S_n$ did not contain $n$ terms, but rather just two: 1 and $1/(n+1)$.\index{series!telescoping}\index{telescoping series}

When possible, seek a way to write an explicit formula for the $n^{\text{th}}$ partial sum $S_n$. This makes evaluating the limit $\ds\lim_{n\to\infty} S_n$ much more approachable. We do so in the next example.

%\noindent\textbf{Note on notation:} Most of the series we encounter will start with $n=1$. For ease of notation, we will often write $\sum a_n$ instead of writing $\ds\sum_{n=1}^\infty a_n$.\\

% the next mtable needs to be outside the {example} for tagging, before the example because the example is 1.4 pages long,
% but after the page break so that it appears next to the example
\newpage

\mtable[yshift=-0.5\textheight]{Scatter plots relating to the series of \autoref{ex_series4} part 1.}{fig:series4a}{\begin{tikzpicture}[alt={Scatter-blue aₙ drops toward 0 while red Sₙ rises and levels near 1.3; data for 1≤n≤10.}]
\begin{axis}[width=1.16\marginparwidth,tick label style={font=\scriptsize},
axis y line=middle,axis x line=middle,name=myplot,axis on top,xtick={2,4,6,8,10},
ymin=-.1,ymax=1.6,xmin=-1,xmax=11,legend columns=-1,legend style={at={(.5,-.1)},anchor=north}]
\addplot [only marks,draw={\colorone},fill={\colorone},mark size={1.5pt},domain=1:10,samples=10] {2/(x^2+2*x)};
\addlegendentry{\scriptsize$a_n$\quad\mbox{}}
\addplot [only marks,draw={\colortwo},fill={\colortwo},mark size={1.5pt},domain=1:10,samples=10] {1.5-1/(x+1)-1/(x+2)};
\addlegendentry{\scriptsize$S_n$}
\end{axis}
\node [right] at (myplot.right of origin) {\scriptsize $n$};
\node [above] at (myplot.above origin) {\scriptsize $y$};
\end{tikzpicture}}\vspace{-\topsep}\vspace{-\partopsep}% this shouldn't be needed.
% todo bugfix see https://tex.stackexchange.com/q/750678/107497

\begin{example}[Evaluating series]\label{ex_series4}%
Evaluate each of the following infinite series.
\[
 \text{1.}\quad\sum_{n=1}^\infty \frac{2}{n^2+2n}\qquad\qquad
 \text{2.}\quad\sum_{n=1}^\infty \ln\left(\frac{n+1}{n}\right)
\]
\solution
\begin{enumerate}
\item		We can decompose the fraction $2/(n^2+2n)$ as
\[\frac2{n^2+2n} = \frac1n-\frac1{n+2}.\]
(See \autoref{sec:partial_fraction}, Partial Fraction Decomposition, to recall how  this is done, if necessary.)

Expressing the terms of $\{S_n\}$ is now more instructive:%
{%
\small%
\begin{align*}
S_1 &= 1-\tfrac13 &&= 1-\tfrac13\\
S_2 &= (1-\tfrac13) + (\tfrac12-\tfrac14) &&= 1+\tfrac12-\tfrac13-\tfrac14\\
S_3 &= (1-\tfrac13) + (\tfrac12-\tfrac14) + (\tfrac13-\tfrac15)
&&= 1+\tfrac12-\tfrac14-\tfrac15\\
S_4 &= (1-\tfrac13) + (\tfrac12-\tfrac14) + (\tfrac13-\tfrac15) + (\tfrac14-\tfrac16)
&&= 1+\tfrac12-\tfrac15-\tfrac16\\
S_5 &= (1-\tfrac13) + (\tfrac12-\tfrac14) + (\tfrac13-\tfrac15) + (\tfrac14-\tfrac16) + (\tfrac15-\tfrac17)\mkern-18mu
&&= 1+\tfrac12-\tfrac16-\tfrac17
\end{align*}}%
We again have a telescoping series. In each partial sum, most of the terms pair up to add to zero and we obtain the formula $\ds S_n = 1+\frac12-\frac1{n+1}-\frac1{n+2}.$ Taking limits allows us to determine the convergence of the series:
\[
\lim_{n\to\infty}S_n
= \lim_{n\to\infty} \left(1+\frac12-\frac1{n+1}-\frac1{n+2}\right) = \frac32,
\]
so $\ds\sum_{n=1}^\infty \frac1{n^2+2n} = \frac32$.
This is illustrated in \autoref{fig:series4a}.

\item	We begin by writing the first few partial sums of the series:
\begin{align*}
S_1 &= \ln\left(2\right) \\
S_2 &= \ln\left(2\right)+\ln\left(\frac32\right) \\
S_3 &= \ln\left(2\right)+\ln\left(\frac32\right)+\ln\left(\frac43\right) \\
S_4 &= \ln\left(2\right)+\ln\left(\frac32\right)+\ln\left(\frac43\right)
+\ln\left(\frac54\right) 
\end{align*}
At first, this does not seem helpful, but recall the logarithmic identity: $\ln x+\ln y = \ln (xy).$ Applying this to $S_4$ gives:
\[
S_4 = \ln\left(2\right)+\ln\left(\frac32\right)+\ln\left(\frac43\right)
+\ln\left(\frac54\right)
=\ln\left(\frac21\cdot\frac32\cdot\frac43\cdot\frac54\right)
=\ln\left(5\right).
\]
We must generalize this for $S_n$.
\begin{multline*}
S_n=\ln\left(2\right)+\ln\left(\frac32\right)+\dots +\ln \left(\frac{n+1}{n}\right)\\
=\ln\left(\frac21\cdot\frac32 \dots  \frac{n}{n-1}\cdot \frac{n+1}{n}\right)
=\ln(n+1)
\end{multline*}

We can conclude that $\{S_n\} = \bigl\{\ln (n+1)\bigr\}$. This sequence  does not converge, as $\ds \lim_{n\to\infty}S_n=\infty$. Therefore  $\ds\sum_{n=1}^\infty  \ln\left(\frac{n+1}{n}\right)=\infty$; the series diverges. Note in \autoref{fig:series4b} how the sequence of partial sums grows slowly; after 100 terms, it is not yet over 5. Graphically we may be fooled into thinking the series converges, but our analysis above shows that it does not.
\end{enumerate}
\end{example}

\mtable[yshift=10\baselineskip]{Scatter plots relating to the series of \autoref{ex_series4} part 2.}{fig:series4b}{\begin{tikzpicture}[alt={Scatter-blue aₙ stays near 0, red Sₙ increases slowly to about 4.5 over 10≤n≤100.}]
\begin{axis}[width=1.16\marginparwidth,tick label style={font=\scriptsize},
axis y line=middle,axis x line=middle,name=myplot,axis on top,
ymin=-.5,ymax=5,xmin=-10,xmax=110,legend columns=-1,legend style={at={(.5,-.1)},anchor=north}]
\addplot [only marks,draw={\colorone},fill={\colorone},mark size={1.5pt},domain=10:100,samples=10] {ln((x+1)/x)};
\addlegendentry{\scriptsize$a_n$\quad\mbox{}}
\addplot [only marks,draw={\colortwo},fill={\colortwo},mark size={1.5pt},domain=10:100,samples=10] {ln(x+1)};
\addlegendentry{\scriptsize$S_n$}
\end{axis}
\node [right] at (myplot.right of origin) {\scriptsize $n$};
\node [above] at (myplot.above origin) {\scriptsize $y$};
\end{tikzpicture}}

We are learning about a new mathematical object, the series. As done before, we apply ``old'' mathematics to this new topic.

\begin{theorem}[Properties of Infinite Series]\label{thm:series_prop}%
Suppose that \quad$\ds \sum_{n=1}^\infty a_n$\quad and \quad $\ds\sum_{n=1}^\infty b_n$ are convergent series, and that \quad$\ds \sum_{n=1}^\infty a_n = L$,\quad  $\ds\sum_{n=1}^\infty b_n = K$, and $c$ is a constant.
\index{series!properties}\index{Sum/Difference Rule!of series}
\index{Constant Multiple Rule!of series}
\begin{enumerate}
\item  Constant Multiple Rule: $\ds\sum_{n=1}^\infty c\cdot a_n = c\cdot\sum_{n=1}^\infty a_n = c\cdot L.$
\item		Sum/Difference Rule: $\ds\sum_{n=1}^\infty \bigl(a_n\pm b_n\bigr) = \sum_{n=1}^\infty a_n \pm \sum_{n=1}^\infty b_n = L \pm K.$
\end{enumerate}
\end{theorem}

Before using this theorem, we will consider the harmonic series $\ds\sum_{n=1}^\infty\frac1n$.

\begin{example}[Divergence of the Harmonic Series]\label{ex_harm_series}%
Show that the harmonic series $\ds \sum_{n=1}^\infty\frac1n$ diverges.
\solution
We will use a proof by contradiction here. Suppose the harmonic series converges to $S$. That is
\[S=1+\frac12+ \frac13 +\frac14+\frac15+\frac16+\frac17+\frac18+\dotsb\]
We then have
\begin{align*}
S &\geq 1+\frac12+\frac14+\frac14+\frac16+\frac16+\frac18+\frac18+\dotsb\\
&=1+\frac12+\frac12\phantom{+\frac14}+\frac13\phantom{+\frac16}
+\frac14\phantom{+\frac18}+\dotsb\\
&=\frac12+S
\end{align*}
This gives us $S\geq \frac12+S$ which can never be true, thus our assumption that the harmonic series converges must be false. Therefore, the harmonic series diverges.
\end{example}

%Before using this theorem, we provide a few ``famous'' series.
%
%{\tcbset{grow to right by=2em}}
%\begin{keyidea}[Important Series]\label{idea:famous_series}%
%\begin{enumerate}
%\item	\parbox{90pt}{$\ds\sum_{n=0}^\infty \frac1{n!} = e$. } (Note that the index starts with $n=0$.)
%\item	$\ds\sum_{n=1}^\infty \frac1{n^2} = \frac{\pi^2}{6}$.
%\item	$\ds\sum_{n=1}^\infty \frac{(-1)^{n+1}}{n^2} = \frac{\pi^2}{12}$.
%\item	$\ds\sum_{n=0}^\infty \frac{(-1)^{n}}{2n+1} = \frac{\pi}{4}$.
%\item	\parbox{90pt}{$\ds\sum_{n=1}^\infty \frac{1}{n} $ \quad diverges.} (This is called the \emph{Harmonic Series}.)\index{Harmonic Series}
%\item	\parbox{90pt}{$\ds\sum_{n=1}^\infty \frac{(-1)^{n+1}}{n} = \ln 2$.} (This is called the \emph{Alternating Harmonic Series}.)\index{Alternating Harmonic Series}
%\end{enumerate}
%\end{keyidea}
%
%\begin{example}[Evaluating series]\label{ex_series5}%
%Evaluate the given series.
%
%\[
%\text{1.}\quad\sum_{n=1}^\infty \frac{(-1)^{n+1}\bigl(n^2-n\bigr)}{n^3}\qquad
%\text{2.}\quad\sum_{n=1}^\infty \frac{1000}{n!}\qquad
%\text{3.}\quad\frac1{16}+\frac1{25}+\frac1{36}+\frac1{49}+\dotsb
%\]
%\solution
%\begin{enumerate}
%\item	We start by using algebra to break the series apart:
%\begin{align*}
%\sum_{n=1}^\infty \frac{(-1)^{n+1}\bigl(n^2-n\bigr)}{n^3} &= \sum_{n=1}^\infty\left(\frac{(-1)^{n+1}n^2}{n^3}-\frac{(-1)^{n+1}n}{n^3}\right) \\
%						&= \sum_{n=1}^\infty\frac{(-1)^{n+1}}{n}-\sum_{n=1}^\infty\frac{(-1)^{n+1}}{n^2} \\
%						&= \ln(2) - \frac{\pi^2}{12}	\approx	-0.1293.
%\end{align*}
%
%This is illustrated in \autoref{fig:series5}(a).
%\mtable{Scatter plots relating to the series of \autoref{ex_series5} part 1.}{fig:series5a}{\begin{tikzpicture}[alt={Blue aₙ alternate signs but decrease in magnitude.  Red Sₙ are negative and settle down around -1.7.}]
%\begin{axis}[width=1.16\marginparwidth,tick label style={font=\scriptsize},
%axis y line=middle,axis x line=middle,name=myplot,axis on top,
%xtick={10,20,30,40,50},ymin=-.3,ymax=.3,xmin=-1,xmax=55,legend columns=-1,legend style={at={(.5,0)},anchor=south}]
%\addplot [only marks,draw={\colorone},fill={\colorone},mark size={1.5pt},domain=5:50,samples=10] {((-1)^(x+1)*(x^2-x))/(x^3)};
%\addlegendentry{\scriptsize$a_n$\quad\mbox{}}
%\addplot [only marks,draw={\colortwo},fill={\colortwo},mark size={1.5pt},domain=1:10,samples=10] coordinates {(5,-0.05528)(10,-0.1723)(15,-0.09917)(20,-0.1525)(25,-0.1105) (30,-0.1452)(35,-0.1156)(40,-0.1414)(45,-0.1186)(50,-0.139)};
%\addlegendentry{\scriptsize$S_n$}
%\end{axis}
%\node [right] at (myplot.right of origin) {\scriptsize $n$};
%\node [above] at (myplot.above origin) {\scriptsize $y$};
%\end{tikzpicture}}
%
%\item		This looks very similar to the series that involves $e$ in \autoref{idea:famous_series}. Note, however, that the series given in this example starts with $n=1$ and not $n=0$. The first term of the series in the Key Idea is $1/0! = 1$, so we will subtract this from our result below:
%\begin{align*}
%		\sum_{n=1}^\infty \frac{1000}{n!} &= 1000\cdot\sum_{n=1}^\infty \frac{1}{n!} \\
%							&= 1000\cdot (e-1) \approx  1718.28.
%\end{align*}
%This is illustrated in \autoref{fig:series5}(b). The graph shows how this particular series converges very rapidly.
%\mtable{Scatter plots relating to the series of \autoref{ex_series5} part 2.}{fig:series5b}{\begin{tikzpicture}[alt={Blue aₙ decrease from 1000 quickly to near 0.  Red Sₙ increase to around 1700.}]
%\begin{axis}[width=1.16\marginparwidth,tick label style={font=\scriptsize},
%axis y line=middle,axis x line=middle,name=myplot,axis on top,
%ymin=-.3,ymax=2100,xmin=-1,xmax=11,legend columns=-1,legend style={at={(.5,-.15)},anchor=north}]
%\addplot [only marks,draw={\colorone},fill={\colorone},mark size={1.5pt},domain=1:10,samples=10]coordinates {(1,1000.)(2,500.)(3,166.7)(4,41.67)(5,8.333)(6,1.389)(7,0.1984)(8,0.0248)(9,0.002756)(10,0.0002756)};
%\addlegendentry{\scriptsize$a_n$\quad\mbox{}}
%\addplot [only marks,draw={\colortwo},fill={\colortwo},mark size={1.5pt},domain=1:10,samples=10]coordinates {(1,1000.)(2,1500.)(3,1667.)(4,1708.)(5,1717.)(6,1718.)(7,1718.)(8,1718.)(9,1718.)(10,1718.)};
%\addlegendentry{\scriptsize$S_n$}
%\end{axis}
%\node [right] at (myplot.right of origin) {\scriptsize $n$};
%\node [above] at (myplot.above origin) {\scriptsize $y$};
%\end{tikzpicture}}
%\mtable{Scatter plots relating to the series in \autoref{ex_series5}.}{fig:series5}{%
%\begin{tikzpicture}[alt={Every five steps, blue aₙ alternate signs decreasing in magnitude toward 0.  Red Sₙ settle to around -0.15.  5≤n≤50.}]
%\begin{axis}[width=1.16\marginparwidth,tick label style={font=\scriptsize},
%axis y line=middle,axis x line=middle,name=myplot,axis on top,
%xtick={10,20,30,40,50},ymin=-.3,ymax=.3,xmin=-1,xmax=55,legend columns=-1,legend style={at={(.5,0)},anchor=south}]
%\addplot [only marks,draw={\colorone},fill={\colorone},mark size={1.5pt},domain=5:50,samples=10] {((-1)^(x+1)*(x^2-x))/(x^3)};
%\addlegendentry{\scriptsize$a_n$\quad\mbox{}}
%\addplot [only marks,draw={\colortwo},fill={\colortwo},mark size={1.5pt},domain=1:10,samples=10]coordinates {(5,-0.05528)(10,-0.1723)(15,-0.09917)(20,-0.1525)(25,-0.1105)(30,-0.1452)(35,-0.1156)(40,-0.1414)(45,-0.1186)(50,-0.139)};
%\addlegendentry{\scriptsize$S_n$}
%\end{axis}
%\node [right] at (myplot.right of origin) {\scriptsize $n$};
%\node [above] at (myplot.above origin) {\scriptsize $y$};
%\end{tikzpicture}
%\\[10pt](a)\\[15pt]
%\begin{tikzpicture}[alt={Blue aₙ decreasing from 1000 rapidly to nearly 0.  Red Sₙ increasing from 1000 to around 1700.}]
%\begin{axis}[width=1.16\marginparwidth,tick label style={font=\scriptsize},
%axis y line=middle,axis x line=middle,name=myplot,axis on top,
%ymin=-.3,ymax=2100,xmin=-1,xmax=11,legend columns=-1,legend style={at={(.5,-.15)},anchor=north}]
%\addplot [only marks,draw={\colorone},fill={\colorone},mark size={1.5pt},domain=1:10,samples=10]coordinates {(1,1000.)(2,500.)(3,166.7)(4,41.67)(5,8.333)(6,1.389)(7,0.1984)(8,0.0248)(9,0.002756)(10,0.0002756)};
%\addlegendentry{\scriptsize$a_n$\quad\mbox{}}
%\addplot [only marks,draw={\colortwo},fill={\colortwo},mark size={1.5pt},domain=1:10,samples=10]coordinates {(1,1000.)(2,1500.)(3,1667.)(4,1708.)(5,1717.)(6,1718.)(7,1718.)(8,1718.)(9,1718.)(10,1718.)};
%\addlegendentry{\scriptsize$S_n$}
%\end{axis}
%\node [right] at (myplot.right of origin) {\scriptsize $n$};
%\node [above] at (myplot.above origin) {\scriptsize $y$};
%\end{tikzpicture}}
%\\[10pt](b)
%
%\item		The denominators in each term are perfect squares; we are adding $\ds \sum_{n=4}^\infty \frac{1}{n^2}$ (note  we start with $n=4$, not $n=1$). This series will converge. Using the formula from \autoref{idea:famous_series}, we have the following:
%\begin{align*}
%\sum_{n=1}^\infty \frac1{n^2} &= \sum_{n=1}^3 \frac1{n^2} +\sum_{n=4}^\infty \frac1{n^2} \\
%\sum_{n=1}^\infty \frac1{n^2} - \sum_{n=1}^3 \frac1{n^2} &=\sum_{n=4}^\infty \frac1{n^2} \\
%\frac{\pi^2}{6} - \left(\frac11+\frac14+\frac19\right) &= \sum_{n=4}^\infty \frac1{n^2} \\
%\frac{\pi^2}{6} - \frac{49}{36} &= \sum_{n=4}^\infty \frac1{n^2} \\
%0.2838&\approx \sum_{n=4}^\infty \frac1{n^2} 
%\end{align*}
%\end{enumerate}
%\end{example}

It may take a while before one is comfortable with this statement, whose truth lies at the heart of the study of infinite series: \emph{it is possible that the sum of an infinite list of nonzero numbers is finite.} We have seen this repeatedly in this section, yet it still may ``take some getting used to.''

As one contemplates the behavior of series, a few facts become clear. 
\begin{enumerate}
\item		In order to add an infinite list of nonzero numbers and get a finite result, ``most'' of those numbers must be ``very near'' 0. 
\item		If a series diverges, it means that the sum of an infinite list of numbers is not finite (it may approach $\pm \infty$ or it may oscillate), and:
\begin{enumerate}
	\item	The series will still diverge if the first term is removed.
	\item	The series will still diverge if the first 10 terms are removed.
	\item	The series will still diverge if the first $1{,}000{,}000$ terms are removed.
	\item	The series will still diverge if any finite number of terms from anywhere in the series are removed.
\end{enumerate}
\end{enumerate}

These concepts are very important and lie at the heart of the next two theorems.

\begin{theorem}[Convergence of Sequence]\label{thm:series_conv}%
If the series $\ds \sum_{n=1}^\infty a_n$ converges, then $\ds\lim_{n\to\infty}a_n=0$.
\end{theorem}

\begin{proof}
Let $S_n=a_1+a_2+\dots+a_n$. We have 
\begin{align*}
S_n&=a_1+a_2+\dots+a_{n-1}+a_n\\
S_n&=S_{n-1}+a_n\\
a_n&=S_n-S_{n-1}
\end{align*}
Since  $\ds \sum_{n\to\infty}a_n$ converges, the sequence $\{ S_n\}$ converges.  Let $\ds \lim_{n\to\infty} S_n=S$. As $n\to \infty$, $n-1$ also goes to $\infty$, so $\ds \lim_{n\to\infty} S_{n-1}=S$. We now have
\begin{align*}
\lim_{n\to\infty} a_n&= \lim_{n\to\infty}(S_n-S_{n-1})\\
&= \lim_{n\to\infty}S_n - \lim_{n\to\infty} S_{n-1}\\
&=S-S=0\qedhere
\end{align*}
\end{proof}

\begin{theorem}[Test for Divergence]\label{thm:series_nth_term}%
If $ \ds \lim_{n\to\infty} a_n$ does not exist or $\ds  \lim_{n\to\infty}a_n\neq0$, then the series $\ds \sum_{n=1}^\infty a_n$ diverges.
\end{theorem}

The Test for Divergence follows from \autoref{thm:series_conv}. If the series does not diverge, it must converge and therefore $\ds  \lim_{n\to\infty}a_n=0$.

Note that the two statements in Theorems \ref{thm:series_conv} and \ref{thm:series_nth_term} are really the same. In order to converge, the terms of the sequence must approach 0; if they do not, the series will not converge. 

Looking back, we can apply this theorem to the series in \autoref{ex_series1}. In that example, we had $\{a_n\} = \{n^2\}$ and $\{b_n\} = \{(-1)^{n+1}\}$.
\[\lim_{n\to\infty} a_n=\lim_{n\to\infty} n^2=\infty\]
and
\[\lim_{n\to\infty} b_n=\lim_{n\to\infty}(-1)^{n+1}\text{ which does not exist.}\]
Thus by the Test for Divergence, both series will diverge.

\textbf{Important!} This theorem \emph{does not state} that if $\ds \lim_{n\to\infty} a_n = 0$ then $\ds \sum_{n=1}^\infty  a_n $ converges. The standard example of this is the Harmonic Series, as given in \autoref{ex_harm_series}. The Harmonic Sequence, $\{1/n\}$, converges to 0; the Harmonic Series, $\ds \sum_{n=1}^\infty 1/n$, diverges.

\begin{theorem}[Infinite Nature of Series]\label{thm:series_behavior}%
The convergence or divergence of a series remains unchanged by the insertion or deletion of any finite number of terms. That is:
\begin{enumerate}
	\item	A divergent series will remain divergent with the insertion or deletion of any finite number of terms.
	\item	A convergent series will remain convergent with the insertion or deletion of any finite number of terms. (Of course, the \emph{sum} will likely change.)
\end{enumerate}
\end{theorem}

In other words, when we are only interested in the convergence or divergence of a series, it is safe to ignore the first few billion terms.

\begin{example}[Removing Terms from the Harmonic Series]\label{ex_trunc_harm}%
Consider once more the Harmonic Series $\ds\sum_{n=1}^\infty\frac1n$\vspace{-.8\baselineskip} which diverges; that is, the partial sums $S_N=\ds\sum_{n=1}^N\frac1n$ grow (very, very slowly) without bound. One might think that by removing the ``large'' terms of the sequence that perhaps the series will converge. This is simply not the case. For instance, the sum of the first 10 million terms of the Harmonic Series is about 16.7. Removing the first 10 million terms from the Harmonic Series changes the partial sums,  effectively subtracting 16.7 from the sum. However, a sequence that is growing without bound will still grow without bound when 16.7 is subtracted from it. 

The equation below illustrates this. Even though we have subtracted off the first 10 million terms, this only subtracts a constant off of an expression that is still growing to infinity. Therefore, the modified series is still growing to infinity.
\begin{multline*}
 \sum_{n=10,000,001}^\infty\frac1n
 =\lim_{N\to\infty}\sum_{n=10,000,001}^N\frac1n
 =\lim_{N\to\infty}\sum_{n=1}^N\frac1n-\sum_{n=1}^{10,000,001}\frac1n \\
 =\lim_{N\to\infty}\sum_{n=1}^N\frac1n-16.7
 =\infty.
\end{multline*}
\end{example}

This section introduced us to series and defined a few special types of series whose convergence properties are well known. We know when a geometric series converges or diverges. Most series that we encounter are not one of these types, but we are still interested in knowing whether or not they converge. The next three sections introduce tests that help us determine whether or not a given series converges.

\printexercises{exercises/08-02-exercises}

\section{The Integral Test}\label{sec:int_tests}

Knowing whether or not a series converges is very important, especially when we discuss Power Series in \autoref{sec:power_series}. \autoref{thm:geom_series} gives criteria for when Geometric series converge and \autoref{thm:series_nth_term} gives a quick test to determine if a series diverges. There are many important series whose convergence cannot be determined by these theorems, though, so we introduce a set of tests that allow us to handle a broad range of series. We start with the Integral Test.


\subsection{Integral Test}

We stated in \autoref{sec:sequences} that a sequence $\{a_n\}$ is a function $a(n)$ whose domain is $\mathbb{N}$, the set of natural numbers. If we can extend $a(n)$ to have the domain of $\mathbb{R}$, the real numbers, and it is both positive and decreasing on $[1,\infty)$, then the convergence of $\ds \sum_{n=1}^\infty a_n$ is the same as $\ds\int_1^\infty a(x)\dd x$. 

\begin{theorem}[Integral Test]\label{thm:integral_test}%
Let a sequence $\{a_n\}$ be defined by $a_n=a(n)$, where $a(n)$ is continuous, positive, and decreasing on $[1,\infty)$. Then $\ds \sum_{n=1}^\infty a_n$ converges, if, and only if, $\ds\int_1^\infty a(x)\dd x$ converges. In other words:
\begin{enumerate}
\item If $\ds \int_1^\infty a(x)\dd x$ is convergent, then $\ds \sum_{n=1}^\infty a_n$ is convergent.
\item If $\ds \int_1^\infty a(x)\dd x$ is divergent, then $\ds \sum_{n=1}^\infty a_n$ is divergent.
\end{enumerate}
\index{series!Integral Test}\index{Integral Test}\index{convergence!Integral Test}\index{divergence!Integral Test}
\end{theorem}

\mnote{\textbf{Note:} \autoref{thm:integral_test} does not state that the integral and the summation have the same value.}

Note that it is not necessary to start the series or the integral at $n=1$. We may use any interval $[n,\infty)$ on which $a(n)$ is continuous, positive and decreasing. Also the sequence $\{a_n\}$ does not have to be strictly decreasing. It must be \emph{ultimately decreasing} which means it is decreasing for all $n$ larger than some number $N$.\bigskip



We can demonstrate the truth of the Integral Test with two simple graphs. In \autoref{fig:integral_test}(a), the height of each rectangle is $a(n)=a_n$ for $n=1,2,\dotsc$, and clearly the rectangles enclose more area than the area under $y=a(x)$. Therefore we can conclude that
\begin{equation}
\int_1^\infty a(x)\dd x < \sum_{n=1}^\infty a_n.\label{eq:integral_testa}
\end{equation}
In \autoref{fig:integral_test}(b), we draw rectangles under $y=a(x)$ with the Right-Hand rule, starting with $n=2$. This time, the area of the rectangles is less than the area under $y=a(x)$, so $\ds\sum_{n=2}^\infty a_n < \int_1^\infty a(x)\dd x$. Note how this summation starts with $n=2$; adding $a_1$ to both sides lets us rewrite the summation starting with $n=1$:
\begin{equation}
\sum_{n=1}^\infty a_n < a_1 +\int_1^\infty a(x)\dd x.\label{eq:integral_testb}
\end{equation} 

\mtable[yshift=5\baselineskip]{Illustrating the truth of the Integral Test.}{fig:integral_test}{%
\begin{tikzpicture}[alt={Left-end rectangles (red) over a decreasing y=a(x) curve; overestimate area on 1 < x < }]
\begin{axis}[width=1.16\marginparwidth,tick label style={font=\scriptsize},
axis y line=middle,axis x line=middle,name=myplot,xtick={1,2,3,4,5},
ytick={1,2},ymin=-.1,ymax=2.5,xmin=-.5,xmax=5.6]
\addplot [very thick,draw={\colorone}, smooth,domain=-.5:5.6,samples=20] ({x},{2/(x+1)});
\draw [very thick,draw={\colortwo}]
	(axis cs: 1,0)--(axis cs:1,1)--(axis cs: 2,1)--(axis cs: 2,0)
	(axis cs: 2,0)--(axis cs:2,.667)--(axis cs: 3,.667)--(axis cs: 3,0)
	(axis cs: 3,0)--(axis cs:3,.5)--(axis cs: 4,.5)--(axis cs: 4,0)
	(axis cs: 4,0)--(axis cs:4,.4)--(axis cs: 5,.4)--(axis cs: 5,0);
\draw (axis cs: .95,1.75) node {\scriptsize $y=a(x)$};
\end{axis}
\node [right] at (myplot.right of origin) {\scriptsize $x$};
\node [above] at (myplot.above origin) {\scriptsize $y$};
\end{tikzpicture}
\\(a)\\
\begin{tikzpicture}[alt={Right-end rectangles (red) under same decreasing y=a(x) curve; underestimate area on 1 < x < 5.}]
\begin{axis}[width=1.16\marginparwidth,tick label style={font=\scriptsize},
axis y line=middle,axis x line=middle,name=myplot,xtick={1,2,3,4,5},
ytick={1,2},ymin=-.1,ymax=2.5,xmin=-.5,xmax=5.6]
\addplot [very thick,draw={\colorone}, smooth,domain=-.5:5.6,samples=20] ({x},{2/(x+1)});
\draw [very thick,draw={\colortwo}] 
	(axis cs: 1,0) -- (axis cs:1,.667) -- (axis cs: 2,.667) -- (axis cs: 2,0)
	(axis cs: 2,0) -- (axis cs:2,.5) -- (axis cs: 3,.5) -- (axis cs: 3,0)
	(axis cs: 3,0) -- (axis cs:3,.4) -- (axis cs: 4,.4) -- (axis cs: 4,0)
	(axis cs: 4,0) -- (axis cs:4,.333) -- (axis cs: 5,.333) -- (axis cs: 5,0);
\draw (axis cs: .95,1.75) node {\scriptsize $y=a(x)$};
\end{axis}
\node [right] at (myplot.right of origin) {\scriptsize $x$};
\node [above] at (myplot.above origin) {\scriptsize $y$};
\end{tikzpicture}
\\(b)}

Combining Equations \eqref{eq:integral_testa} and \eqref{eq:integral_testb}, we have
\begin{equation}
\sum_{n=1}^\infty a_n< a_1 +\int_1^\infty a(x)\dd x < a_1 + \sum_{n=1}^\infty a_n.\label{eq:integral_testc}
\end{equation}
From \autoeqref{eq:integral_testc} we can make the following two statements:
% todo Tim etal are these statements backwards?  should the prior equation be rearranged?
\begin{enumerate}
	\item If $\ds \sum_{n=1}^\infty a_n$ diverges, so does $\ds\int_1^\infty a(x)\dd x$ \\\hfill(because $\ds \sum_{n=1}^\infty a_n < a_1 +\int_1^\infty a(x)\dd x$)
	\item	If $\ds \sum_{n=1}^\infty a_n$ converges, so does $\ds\int_1^\infty a(x)\dd x$ \\\hfill(because $\ds \ds \int_1^\infty a(x)\dd x < \sum_{n=1}^\infty a_n$.)
\end{enumerate}
Therefore the series and integral either both converge or both diverge. \autoref{thm:series_behavior} allows us to extend this theorem to series where $a(n)$ is positive and decreasing on $[b,\infty)$ for some $b>1$.

\youtubeVideo{ObiRjUFHJHo}{Integral Test for Series: Why It Works}

\begin{example}[Using the Integral Test]\label{ex_itest1}%
Determine the convergence of $\ds\sum_{n=1}^\infty \frac{\ln n}{n^2}$.  (The terms of the sequence $\{a_n\} = \{(\ln n)/n^2\}$ and the n\textsuperscript{th} partial sums are given in \autoref{fig:itest1}.)
\solution
\autoref{fig:itest1}
\mtable{Plotting the sequence and series in \autoref{ex_itest1}.}{fig:itest1}{\begin{tikzpicture}[alt={Blue dots aₙ fall toward 0 while red dots Sₙ climb and level near 0.75 across 1≤n≤20.}]
\begin{axis}[width=\marginparwidth,tick label style={font=\scriptsize},
axis y line=middle,axis x line=middle,name=myplot,axis on top,xtick={2,4,...,20},
ymin=-.05,ymax=.85,xmin=-1,xmax=20.9,legend columns=-1,legend style={at={(.5,-.1)},anchor=north}]
\addplot [only marks,draw={\colorone},fill={\colorone},mark size={1.5pt},domain=1:10,samples=10]coordinates {(1.,0)(2.,0.1733)(3.,0.1221)(4.,0.08664)(5.,0.06438)(6.,0.04977)(7.,0.03971)(8.,0.03249)(9.,0.02713)(10.,0.02303)(11.,0.01982)(12.,0.01726)(13.,0.01518)(14.,0.01346)(15.,0.01204)(16.,0.01083)(17.,0.009804)(18.,0.008921)(19.,0.008156)(20.,0.007489)};
\addlegendentry{\scriptsize$a_n$\quad\mbox{}}
\addplot [only marks,draw={\colortwo},fill={\colortwo},mark size={1.5pt}] coordinates{(1.,0)(2.,0.1733)(3.,0.2954)(4.,0.382)(5.,0.4464)(6.,0.4961)(7.,0.5359)(8.,0.5684)(9.,0.5955)(10.,0.6185)(11.,0.6383)(12.,0.6556)(13.,0.6708)(14.,0.6842)(15.,0.6963)(16.,0.7071)(17.,0.7169)(18.,0.7258)(19.,0.734)(20.,0.7415)};
\addlegendentry{\scriptsize$S_n$}
\end{axis}
\node [right] at (myplot.right of origin) {\scriptsize $n$};
\node [above] at (myplot.above origin) {\scriptsize $y$};
\end{tikzpicture}} 
implies that $a(n) = (\ln n)/n^2$ is positive and decreasing on $[2,\infty)$. We can determine this analytically, too. We know $a(n)$ is positive as both $\ln n$ and $n^2$ are positive on $[2,\infty)$. To determine that $a(n)$ is decreasing, consider $a'(n) = (1-2\ln n)/n^3$, which is negative for $n\geq 2$. Since $a'(n)$ is negative, $a(n)$ is decreasing.

Applying the Integral Test, we test the convergence of $\ds \int_1^\infty \frac{\ln x}{x^2}\dd x$. Integrating this improper integral requires the use of Integration by Parts, with $u = \ln x$ and $dv = 1/x^2\dd x$. 
\begin{align*}
	\int_1^\infty \frac{\ln x}{x^2}\dd x
	&= \lim_{t\to\infty} \int_1^t \frac{\ln x}{x^2}\dd x\\
	&= \lim_{t\to\infty}\biggl( -\frac1x\ln x\Big|_1^t + \int_1^t\frac1{x^2}\dd x \biggr)\\
	&= \lim_{t\to\infty} \biggl( \Bigl[-\frac1x\ln x -\frac 1x\Bigr]_1^t\biggr)\\
	&= \lim_{t\to\infty}\biggl(1-\frac1t-\frac{\ln t}{t}\biggr).\quad \text{Apply L'Hôpital's Rule:}\\
	&= 1-0- \lim_{t\to\infty} \frac1t\\
	&=1
\end{align*}
Since $\ds \int_1^\infty \frac{\ln x}{x^2}\dd x$ converges, so does $\ds \sum_{n=1}^\infty \frac{\ln n}{n^2}$.
\end{example}

\subsection{\texorpdfstring{$p$}{p}-Series}

Another important type of series is the \emph{p-series}.

\begin{definition}[$p$-Series, General $p$-Series]\label{def:pseries}%
\index{series!p@$p$-series}\index{p@$p$-series}\mbox{}\\[-2\baselineskip]\parbox[t]{\linewidth}{%
\begin{enumerate}
\item	A \textbf{$p$-series} is a series of the form
\qquad$\ds\sum_{n=1}^\infty \frac{1}{n^p}$.
\item	A \textbf{general $p$-series} is a series of the form 
\qquad$\ds\sum_{n=1}^\infty \frac{1}{(an+b)^p}$,\medskip\\
where $a>0$, $b$ is a real number, and $an+b\neq 0$ for all $n$.
\end{enumerate}}
\end{definition}

Like geometric series, one of the nice things about $p$-series is that they have easy to determine convergence properties.

\begin{theorem}[Convergence of General $p$-Series]\label{thm:pseries}%
Assume $a$ and $b$ are real numbers and  $an+b\neq 0$ for all $n$.\\
A general $p$-series $\ds\sum_{n=1}^\infty \frac{1}{(an+b)^p}$ will converge if, and only if, $p>1$.
\index{series!p@$p$-series}\index{p@$p$-series}
\index{convergence!of p@of $p$-series}\index{divergence!of p@of $p$-series}
\end{theorem}

\begin{proof}
Consider the integral $\ds\int_1^\infty \frac1{(ax+b)^p}\dd x$; assuming $p\neq 1$,
\begin{align*}
	\int_1^\infty \frac1{(ax+b)^p}\dd x
	&= \lim_{t\to\infty} \int_1^t \frac1{(ax+b)^p}\dd x \\
	&= \lim_{t\to\infty} \frac{1}{a(1-p)}(ax+b)^{1-p}\Big|_1^t\\
	&= \lim_{t\to\infty} \frac{1}{a(1-p)}\bigl((at+b)^{1-p}-(a+b)^{1-p}\bigr).
\end{align*}
This limit converges if and only if, $p>1$. It is easy to show that the integral also diverges in the case of $p=1$. (This result is similar to the work preceding \autoref{idea:impint1}.)

Therefore $\ds \sum_{n=1}^\infty \frac 1{(an+b)^p}$ converges if, and only if, $p>1$.
\end{proof}

\begin{example}[Determining convergence of series]\label{ex_series6}%
Determine the convergence of the following series.
\begin{multicols}{3}
\begin{enumerate}\lxAddClass{columns2}
\item		$\ds\sum_{n=1}^\infty \frac{1}{n}$
\item		$\ds\sum_{n=1}^\infty \frac{1}{n^2}$
\item		$\ds\sum_{n=1}^\infty \frac{1}{\sqrt{n}}$
\item		$\ds\sum_{n=1}^\infty \frac{(-1)^n}{n}$
\item		$\ds\sum_{n=11}^\infty \frac{1}{(\frac12n-5)^3}$
\item		$\ds\sum_{n=1}^\infty \frac{1}{2^n}$
\end{enumerate}
\end{multicols}
\solution
\begin{enumerate}
\item	This is a $p$-series with $p=1$. By \autoref{thm:pseries}, this series diverges.

This series is a famous series, called the \emph{Harmonic Series}, so named because of its relationship to \emph{harmonics} in the study of music and sound. 

\item	This is a $p$-series with $p=2$. By \autoref{thm:pseries}, it converges. Note that the theorem does not give a formula by which we can determine \emph{what} the series converges to; we just know it converges. A famous, unexpected result is that this series converges to $\ds{\pi^2}/{6}$.

\item	This is a $p$-series with $p=1/2$; the theorem states that it diverges.

\item	This is not a $p$-series; the definition does not allow for alternating signs. Therefore we cannot apply \autoref{thm:pseries}. We will consider this series again in \autoref{sec:alt_series}.
(Another famous result states that this series, the \emph{Alternating Harmonic Series}, converges to $-\ln 2$.)

\item	This is a general $p$-series with $p=3$, therefore it converges.

\item	This is not a $p$-series, but a geometric series with $r=1/2$. It converges.
\end{enumerate}
\end{example}

In the next section we consider two more convergence tests, both comparison tests. That is, we determine the convergence of one series by  comparing it to another series with known convergence. 

\printexercises{exercises/08-03-exercises}

\input{text/08_Comparison_Tests}
\section{Alternating Series and Absolute Convergence}\label{sec:alt_series}

The series convergence tests we have used require that the underlying sequence $\{a_n\}$ be a positive sequence. (We can relax this with \autoref{thm:series_behavior} and state that there must be an $N>0$ such that $a_n>0$ for all $n>N$; that is, $\{a_n\}$ is positive for all but a finite number of values of $n$.)

In this section we explore series whose summation includes negative terms. We start with a very specific form of series, where the terms of the summation alternate between being positive and negative.

\begin{definition}[Alternating Series]\label{def:alt_series}%
Let $\{b_n\}$ be a positive sequence. An \textbf{alternating series} is a series of either the form
\index{series!alternating}
\[
\sum_{n=1}^\infty (-1)^nb_n\qquad \text{or}\qquad \sum_{n=1}^\infty (-1)^{n+1}b_n.
\]
\end{definition}

We want to think that an alternating sequence $\{a_n\}$ is related to a positive sequence $\{b_n\}$ by $a_n=(-1)^n b_n$.

Recall that the terms of Harmonic Series come from the Harmonic Sequence $\{b_n\} = \{\frac1n\}$. An important alternating series is the \textbf{Alternating Harmonic Series}:
\[
\sum_{n=1}^\infty (-1)^{n+1}\frac1n
= 1-\frac12+\frac13-\frac14+\frac15-\frac16+\dotsb
\]

Geometric Series can also be alternating series when $r<0$. For instance, if $r=-1/2$, the geometric series is
\[
\sum_{n=0}^\infty \left(\frac{-1}{2}\right)^n
= 1-\frac12+\frac14-\frac18+\frac1{16}-\frac1{32}+\dotsb
\]

\autoref{thm:geom_series} states that geometric series converge when $\abs r<1$ and gives the sum: $\ds \sum_{n=0}^\infty r^n = \frac1{1-r}$. When $r=-1/2$ as above, we find
\[
\sum_{n=0}^\infty\left(\frac{-1}{2}\right)^n=\frac1{1-(-1/2)}=\frac 1{3/2}=\frac23.
\]

A powerful convergence theorem exists for other alternating series that meet a few conditions.

\begin{theorem}[Alternating Series Test]\label{thm:alt_series_test}%
Let $\{b_n\}$ be a positive, decreasing sequence where $\ds\lim_{n\to\infty}b_n=0$. Then
\index{series!Alternating Series Test}\index{Alternating Series Test!for series}\index{convergence!Alternating Series Test}\index{divergence!Alternating Series Test}
\[
\sum_{n=1}^\infty (-1)^{n}b_n \qquad \text{and}\qquad \sum_{n=1}^\infty (-1)^{n+1}b_n
\]
converge.
\end{theorem}

The basic idea behind \autoref{thm:alt_series_test} is illustrated in \autoref{fig:alt_series_converge}. A positive, decreasing sequence $\{b_n\}$ is shown along with the partial sums
\[S_n = \sum_{i=1}^n(-1)^{i+1}b_i =b_1-b_2+b_3-b_4+\dotsb+(-1)^{n+1}b_n.\]
Because $\{b_n\}$ is decreasing, the amount by which $S_n$ bounces up and down decreases. Moreover, the odd terms of $S_n$ form a decreasing, bounded sequence, while the even terms of $S_n$ form an increasing, bounded sequence. Since bounded, monotonic sequences converge (see \autoref{thm:monotonic_converge}) and the terms of $\{b_n\}$ approach 0, we will show below that the odd and even terms of $S_n$ converge to the same common limit $L$, the sum of the series.

\mtable{Illustrating convergence with the Alternating Series Test.}{fig:alt_series_converge}{\begin{tikzpicture}[alt={The sequence bₙ is decreasing from 1 toward 0.  The partial sums Sₙ alternate on which side of the dashed line y=L they are.}]
\begin{axis}[width=1.16\marginparwidth,tick label style={font=\scriptsize},
axis y line=middle,axis x line=middle,name=myplot,axis on top,xtick={2,4,6,8,10},
ymin=-.25,ymax=1.25,xmin=-1,xmax=11,legend columns=-1,legend style={at={(.5,0)},anchor=north}]
\draw [thick, dashed] (axis cs: -.1,.614) node [left] {\scriptsize $L$} -- (axis cs: 10.25,.614);
\addplot [only marks,draw={\colorone},fill={\colorone},mark size={1.5pt},domain=1:10,samples=10] {2/(x+1)};
\addlegendentry{\scriptsize$b_n$\quad\mbox{}}
\addplot [only marks,draw={\colortwo},fill={\colortwo},mark size={1.5pt},domain=1:10,samples=10]coordinates {(1., 1.)(2., 0.333333)(3., 0.833333)(4., 0.433333)(5.,0.766667)(6., 0.480952)(7., 0.730952)(8., 0.50873)(9.,0.70873)(10., 0.526912)};
\addlegendentry{\scriptsize$S_n$}
\end{axis}
\node [right] at (myplot.right of origin) {\scriptsize $n$};
\node [above] at (myplot.above origin) {\scriptsize $y$};
\end{tikzpicture}}

\begin{proof}
Because $\{b_n\}$ is a decreasing sequence, we have $b_n-b_{n+1}\geq 0$. 
We will consider the even and odd partial sums separately. First consider the even partial sums.
\begin{align*}
S_2&=b_1-b_2\geq 0  &\text{since } b_2&\leq b_1\\
S_4&=b_1-b_2+b_3-b_4=S_2+b_3-b_4\geq S_2 &\text{since } b_3- b_4&\geq 0\\
S_6&=S_4+b_5-b_6\geq S_4 &\text{since } b_5- b_6&\geq 0\\
\vdots \\
S_{2n}&=S_{2n-2}+b_{2n-1}-b_{2n}\geq S_{2n-2}  &\text{since } b_{2n-1}- b_{2n}&\geq 0
\end{align*}
We now have
\[0\leq S_2\leq S_4\leq S_6\leq  \dotsb \leq S_{2n}\leq \dotsb\]
so $\{S_{2n}\}$ is an increasing sequence. 
But we can also write 
\begin{align*}
S_{2n}&=b_1-b_2+b_3-b_4+b_5-\dotsb -b_{2n-2}+b_{2n-1}-b_{2n}\\
&=b_1-(b_2-b_3)-(b_4-b_5)-\dotsb -(b_{2n-2}-b_{2n-1})-b_{2n}
\end{align*}
Each term in parentheses is positive and $b_{2n}$ is positive so we have $S_{2n}\leq b_1$ for all $n$. We now have the sequence of even partial sums, $\{S_{2n}\}$, is increasing and bounded above so by \autoref{thm:monotonic_converge}
 $\{S_{2n}\}$ converges. Since we know it converges, we will assume it's limit is $L$ or
\[\lim_{n\to \infty}S_{2n}=L\]
Next we determine the limit of the sequence of odd partial sums.
\begin{align*}
\lim_{n\to \infty}S_{2n+1}&=\lim_{n\to \infty}(S_{2n}+b_{2n+1})\\
&=\lim_{n\to \infty}S_{2n}+\lim_{n\to \infty}b_{2n+1}\\
&=L+0\\
&=L
\end{align*}
Both the even and odd partial sums converge to $L$ so we have $\ds \lim_{n\to \infty}S_n=L$, which means the series is convergent.
\end{proof}

\youtubeVideo{aOiZvfFAMW8}{Alternating Series --- Another Example 4}

\begin{example}[Applying the Alternating Series Test]\label{ex_alt1}%
Determine if the Alternating Series Test applies to each of the following series.
\[
 \text{1. }\sum_{n=1}^\infty (-1)^{n+1}\frac1n\qquad
 \text{2. }\sum_{n=2}^\infty (-1)^n\frac{\ln n}{n}\qquad
 \text{3. }\sum_{n=1}^\infty (-1)^{n+1}\frac{\abs{\sin n}}{n^2}
\]
\solution
\begin{enumerate}
	\item This is the Alternating Harmonic Series as seen previously. The underlying sequence is $\{b_n\} = \{1/n\}$, which is positive, decreasing, and approaches 0 as $n\to\infty$. Therefore we can apply the Alternating Series Test and conclude this series converges. 
	
	While the test does not state what the series converges to, we will see later that $\ds \sum_{n=1}^\infty (-1)^{n+1}\frac1n=\ln2.$
	
	\item		The underlying sequence is $\{b_n\} = \{(\ln n)/n\}$. This is positive for $n\geq 2$ and $\ds \lim_{n\to \infty} \frac{\ln n}{n}=\lim_{n\to \infty}\frac{1}{n}=0$ (use L'Hôpital's Rule). However, the sequence is not decreasing for all $n$. It is straightforward to compute $b_1\approx0.347$, $b_2\approx 0.366$, and $b_3\approx 0.347$: the sequence is increasing for at least the first 2 terms. 
	
	We do not immediately conclude that we cannot apply the Alternating Series Test. Rather, consider the long-term behavior of $\{b_n\}$. Treating $b_n=b(n)$ as a continuous function of $n$ defined on $[2,\infty)$, we can take its derivative:
	\[b\primeskip'(n) = \frac{1-\ln n}{n^2}.\]
	The derivative is negative for all $n\geq 3$ (actually, for all $n>e$), meaning $b(n)=b_n$ is decreasing on $[3,\infty)$. We can apply the Alternating Series Test to the series when we start with $n=3$ and conclude that $\ds \sum_{n=3}^\infty(-1)^n\frac{\ln n}{n}$ converges; adding the terms with $n=2$ does not change the convergence (i.e., we apply \autoref{thm:series_behavior}).
	
	The important lesson here is that as before, if a series fails to meet the criteria of the Alternating Series Test on only a finite number of terms, we can still apply the test.
	
	\item  The underlying sequence is $\{b_n\} = \{\abs{\sin n}/n^2\}$. This sequence is positive and approaches $0$ as $n\to\infty$. However, it is not a decreasing sequence; the value of $\abs{\sin n}$ oscillates between $0$ and $1$ as $n\to\infty$. We cannot remove a finite number of terms to make $\{b_n\}$ decreasing, therefore we cannot apply the Alternating Series Test.
	
	Keep in mind that this does not mean we conclude the series diverges; in fact, it does converge. We are just unable to conclude this based on \autoref{thm:alt_series_test}.
\end{enumerate}
\end{example}

One of the famous results of mathematics is that the Harmonic Series,\vspace{-.5\baselineskip} $\ds \sum_{n=1}^\infty \frac1n$ diverges, yet the Alternating Harmonic Series, $\ds \sum_{n=1}^\infty (-1)^{n+1}\frac1n$, converges. The notion that alternating the signs of the terms in a series can make a series converge leads us to the following definitions.\index{Alternating Harmonic Series}

\begin{definition}[Absolute and Conditional Convergence]\label{def:abs_converge}%
\index{convergence!absolute}\index{convergence!conditional}\index{series!absolute convergence}\index{series!conditional convergence}%
\mbox{}\\[-2\baselineskip]\parbox[t]{\linewidth}{%
\begin{enumerate}
	\item A series $\ds \sum_{n=1}^\infty a_n$ \textbf{converges absolutely} if $\ds \sum_{n=1}^\infty \abs{a_n}$ converges.
	\item A series $\ds \sum_{n=1}^\infty a_n$ \textbf{converges conditionally} if $\ds \sum_{n=1}^\infty a_n$ converges but $\ds \sum_{n=1}^\infty \abs{a_n}$ diverges.
\end{enumerate}}
\end{definition}

\mnote{\textbf{Note:} In \autoref{def:abs_converge}, $\ds \sum_{n=1}^\infty a_n$ is not necessarily an alternating series; it just may have some negative terms.}

Thus we say the Alternating Harmonic Series converges conditionally.

\begin{example}[Determining absolute and conditional convergence.]\label{ex_alt_series2}%
Determine if the following series converge absolutely, conditionally, or diverge.\\
\begin{minipage}{.5\textwidth}
 \begin{enumerate}
  \item $\ds\sum_{n=1}^\infty (-1)^n\frac{n+3}{n^2+2n+5}$
%  \item $\ds\sum_{n=1}^\infty (-1)^n\frac{n^2+2n+5}{2^n}$ % ratio test is next section
 \end{enumerate}
\end{minipage}%
\begin{minipage}{.5\textwidth}
 \begin{enumerate}\setcounter{enumi}{1}
  \item $\ds\sum_{n=3}^\infty (-1)^n\frac{3n-3}{5n-10}$
 \end{enumerate}
\end{minipage}
\solution
\begin{enumerate}
	\item We can show the series
	\[\sum_{n=1}^\infty \abs{(-1)^n\frac{n+3}{n^2+2n+5}}= \sum_{n=1}^\infty \frac{n+3}{n^2+2n+5}\]
	diverges using the Limit Comparison Test, comparing with $1/n$. 
	
	The sequence $\{\dfrac{n+3}{n^2+2n+5}\}$ is monotonically decreasing, so that the series $\ds \sum_{n=1}^\infty (-1)^n\frac{n+3}{n^2+2n+5}$ converges using the Alternating Series Test; we conclude it converges conditionally.
	
%	\item	We can show the series
%	\[\sum_{n=1}^\infty \abs{(-1)^n\frac{n^2+2n+5}{2^n}}=\sum_{n=1}^\infty \frac{n^2+2n+5}{2^n}\]
%	converges using the Ratio Test. 
%	
%	Therefore we conclude $\ds \sum_{n=1}^\infty (-1)^n\frac{n^2+2n+5}{2^n}$ converges absolutely.
	
	\item	The series
	\[\sum_{n=3}^\infty \abs{(-1)^n\frac{3n-3}{5n-10}} = \sum_{n=3}^\infty \frac{3n-3}{5n-10}\]
	diverges using the Test for Divergence, so it does not converge absolutely. 
	
	The series $\ds \sum_{n=3}^\infty (-1)^n\frac{3n-3}{5n-10}$ fails the conditions of the Alternating Series Test as $(3n-3)/(5n-10)$ does not approach $0$ as $n\to\infty$. We can state further that this series diverges; as $n\to\infty$, the series effectively adds and subtracts $3/5$ over and over. This causes the sequence of partial sums to oscillate and not converge.
	
	Therefore the series $\ds \sum_{n=1}^\infty (-1)^n\frac{3n-3}{5n-10}$ diverges.
\end{enumerate}
\end{example}

Knowing that a series converges absolutely allows us to make two important statements, given in the following theorem. The first is that absolute convergence is  ``stronger'' than regular convergence.\vspace{-.5\baselineskip}
That is, just because {\small$\ds \sum_{n=1}^\infty a_n$} converges, we cannot conclude that {\small$\ds \sum_{n=1}^\infty \abs{a_n}$} will converge, but knowing a series converges absolutely tells us that {\small$\ds \sum_{n=1}^\infty a_n$} will converge. 

\begin{theorem}[Absolute Convergence Theorem]\label{thm:abs_convergence}%
Let $\ds \sum_{n=1}^\infty a_n$ be a series that converges absolutely.
\index{convergence!absolute}\index{Absolute Convergence Theorem}\index{series!Absolute Convergence Theorem}\index{rearrangements of series}\index{series!rearrangements}
\begin{enumerate}
	\item $\ds \sum_{n=1}^\infty a_n$ converges.
	\item	Let $\{b_n\}$ be any rearrangement of the sequence $\{a_n\}$. Then 
	\[\sum_{n=1}^\infty b_n = \sum_{n=1}^\infty a_n.\]
\end{enumerate}
\end{theorem}

One reason this is important is that our convergence tests all require that the underlying sequence of terms be positive. By taking the absolute value of the terms of a series where not all terms are positive, we are often able to apply an appropriate test and determine absolute convergence. This, in turn, determines that the series we are given also converges.

The second statement relates to \textbf{rearrangements}\index{rearrangements of series}\index{series!rearrangements} of series. When dealing with a finite set of numbers, the sum of the numbers does not depend on the order which they are added. (So $1+2+3 = 3+1+2$.) One may be surprised to find out that when dealing with an infinite set of numbers, the same statement does not always hold true: some infinite lists of numbers may be rearranged in different orders to achieve different sums. The theorem states that the terms of an absolutely convergent series can be rearranged in any way without affecting the sum.

% todo Tim do we want to have an absolutely convergent example to refer back to?
% the former part 2 was removed because it used the ratio test
%In \autoref{ex_alt_series2}, we determined the series in part 2 converges absolutely. \autoref{thm:abs_convergence} tells us the series converges (which we could also determine using the Alternating Series Test).

The theorem states that rearranging the terms of an absolutely convergent series does not affect its sum. This implies that perhaps the sum of a conditionally convergent series can change based on the arrangement of terms. Indeed, it can. The Riemann Rearrangement Theorem (named after Bernhard Riemann) states that any conditionally convergent series can have its terms rearranged so that the sum is any desired value or infinity.

Before we consider an example, we state the following theorem that illustrates how the alternating structure of an alternating series is a powerful tool when approximating the sum of a convergent series.

\begin{theorem}[The Alternating Series Approximation Theorem]\label{thm:alt_series_approx}%
Let $\{b_n\}$ be a sequence that satisfies the hypotheses of the Alternating Series Test, and let $S_n$ and $L$ be the $n^{\text{th}}$ partial sum and sum, respectively, of either $\ds \sum_{n=1}^\infty (-1)^{n}b_n$ or $\ds \sum_{n=1}^\infty (-1)^{n+1}b_n$. Then
\index{series!alternating!Approximation Theorem}
\begin{enumerate}
	\item $\abs{S_n-L} < b_{n+1}$, and
	\item	$L$ is between $S_n$ and $S_{n+1}$.
\end{enumerate}
\end{theorem}

Part 1 of \autoref{thm:alt_series_approx} states that the $n^{\text{th}}$ partial sum of a convergent alternating series will be within $b_{n+1}$ of its total sum. Consider the alternating series we looked at before the statement of the theorem, $\ds \sum_{n=1}^\infty \frac{(-1)^{n+1}}{n^2}$. Since $b_{14} = 1/14^2 \approx 0.0051$, we know that $S_{13}$ is within $0.0051$ of the total sum. 

Moreover, Part 2 of the theorem states that since $S_{13} \approx 0.8252$ and $S_{14}\approx 0.8201$, we know the sum $L$ lies between $0.8201$ and $0.8252$. One use of this is the knowledge that $S_{14}$ is accurate to two places after the decimal.

Some alternating series converge slowly. In \autoref{ex_alt1} we determined the series $\ds\sum_{n=2}^\infty (-1)^{n+1}\frac{\ln n}{n}$ converged. With $n=1001$, we find $(\ln n)/n \approx 0.0069$, meaning that $S_{1000} \approx 0.1633$ is accurate to one, maybe two, places after the decimal. Since $S_{1001} \approx 0.1564$, we know the sum $L$ is $0.1564\leq L\leq0.1633$.

\begin{example}[Approximating the sums of convergent alternating series]\label{ex_alt_series_approx}%
Approximate the sum of the following series, accurate to within $0.001$.
\[
 \text{1. }\sum_{n=1}^\infty (-1)^{n+1}\frac{1}{n^3}\qquad
 \text{2. }\sum_{n=1}^\infty (-1)^{n+1}\frac{\ln n}{n}.
\]
\solution
\begin{enumerate}
	\item  Using \autoref{thm:alt_series_approx}, we want to find $n$ where $1/n^3 \le 0.001$:
	\begin{align*}
	\frac1{n^3} &\leq 0.001=\frac{1}{1000} \\
	n^3 &\geq 1000\\
	n &\geq \sqrt[3]{1000}\\
	n &\geq 10.
	\end{align*}
	Let $L$ be the sum of this series. By Part 1 of the theorem, $\abs{S_9-L}<b_{10} = 1/1000$. We can compute $S_9=0.902116$, which our theorem states is within $0.001$ of the total sum. 
	
	We can use Part 2 of the theorem to obtain an even more accurate result. As we know the 10\textsuperscript{th} term of the series is $-1/1000$, we can easily compute $S_{10} = 0.901116$. Part 2 of the theorem states that $L$ is between $S_9$ and $S_{10}$, so $0.901116 <L<0.902116$.
	
	\item		We want to find $n$ where $(\ln n)/n \le 0.001$. We start by solving $(\ln n)/n = 0.001$ for $n$. This cannot be solved algebraically, so we will use Newton's Method to approximate a solution. 
	
	Let $f(x) = \ln(x)/x-0.001$; we want to know where $f(x) = 0$. We make a guess that $x$ must be ``large,'' so our initial guess will be $x_1=1000$. Recall how Newton's Method works: given an approximate solution $x_n$, our next approximation $x_{n+1}$ is given by
	\[x_{n+1} = x_n - \frac{f(x_n)}{\fp(x_n)}.\]
	We find $\fp(x) = \bigl(1-\ln(x)\bigr)/x^2$. This gives
	\begin{align*}
	x_2 &= 1000 - \frac{\ln(1000)/1000-0.001}{\bigl(1-\ln(1000)\bigr)/1000^2} \\
			&= 2000.
	\end{align*}
	Using a computer, we find that Newton's Method seems to converge to a solution $x=9118.01$ after 8 iterations. Taking the next integer higher, we have $n=9119$, where $\ln(9119)/9119 =0.000999903<0.001$.
	
	Again using a computer, we find $S_{9118} = -0.160369$. Part 1 of the theorem states that this is within $0.001$ of the actual sum $L$. Already knowing the 9,119\textsuperscript{th} term, we can compute $S_{9119} = -0.159369$, meaning
	\[-0.160369 < L < -0.159369.\]
	
\end{enumerate}
Notice how the first series converged quite quickly, where we needed only 10 terms to reach the desired accuracy, whereas the second series took over 9,000 terms.
\end{example}

We now consider the Alternating Harmonic Series once more. We have stated that 
\[
\sum_{n=1}^\infty (-1)^{n+1}\frac1n
= 1-\frac12+\frac13-\frac14+\frac15-\frac16+\frac17\dotsb = \ln 2,
\]
(see \autoref{ex_alt1}). 

Consider the rearrangement where every positive term is followed by two negative terms:
\[
1-\frac12-\frac14+\frac13-\frac16-\frac18+\frac15-\frac1{10}-\frac1{12}\dotsb
\]
(Convince yourself that these are exactly the same numbers as appear in the Alternating Harmonic Series, just in a different order.) Now group some terms and simplify:
\begin{align*}
\left(1-\frac12\right)-\frac14+\left(\frac13-\frac16\right)-\frac18
+\left(\frac15-\frac1{10}\right)-\frac1{12}+\dotsb &= \\
\frac12-\frac14+\frac16-\frac18+\frac1{10}-\frac{1}{12}+\dotsb &= \\
\frac12\left(1-\frac12+\frac13-\frac14+\frac15-\frac16+\dotsb\right) & = \frac12\ln 2.
\end{align*}

By rearranging the terms of the series, we have arrived at a different sum. (One could \emph{try} to argue that the Alternating Harmonic Series does not actually converge to $\ln 2$, because rearranging the terms of the series \emph{shouldn't} change the sum. However, the Alternating Series Test proves this series converges to $L$, for some number $L$, and if the rearrangement does not change the sum, then $L = L/2$, implying $L=0$. But the Alternating Series Approximation Theorem quickly shows that $L>0$. The only conclusion is that the rearrangement \emph{did} change the sum.) This is an incredible result.\bigskip


We mentioned earlier that the Integral Test did not work well with series containing factorial terms. The next section introduces the Ratio Test, which does handle such series well. We also introduce the Root Test, which is good for series where each term is raised to a power.

\printexercises{exercises/08-05-exercises}

\section{Ratio and Root Tests}\label{sec:ratio_root_tests}

\autoref{thm:series_nth_term} states that if a series $\ds\sum_{n=1}^\infty a_n$ converges, then $\ds \lim_{n\to \infty} a_n =0$. That is, the terms of $\{a_n\}$ must get very small. Not only must the terms approach 0, they must approach 0 ``fast enough'': while $\ds \lim_{n\to\infty}1/n=0$, the Harmonic Series $\ds\sum_{n=1}^\infty \frac1n$ diverges as the terms of $\{1/n\}$ do not approach 0 ``fast enough.''

The comparison tests of \autoref{sec:comp_tests} determine convergence by comparing terms of a series to terms of another series whose convergence is known. This section introduces the Ratio and Root Tests, which determine convergence by analyzing the terms of a series to see if they approach 0 ``fast enough.''

\subsection{Ratio Test}

\begin{theorem}[Ratio Test]\label{thm:ratio_test}%
Let $\{a_n\}$ be a sequence where $\ds \lim_{n\to\infty}\abs{\frac{a_{n+1}}{a_n}} = L$.%
\index{series!Ratio Comparison Test}\index{Ratio Comparison Test!for series}\index{convergence!Ratio Comparison Test}\index{divergence!Ratio Comparison Test}%
\begin{enumerate}
	\item If $L<1$, then $\ds\sum_{n=1}^\infty a_n$ converges.
	\item If $L>1$ or $L=\infty$, then $\ds\sum_{n=1}^\infty a_n$ diverges.
	\item If $L=1$, the Ratio Test is inconclusive.
\end{enumerate}
\end{theorem}

%\mnote{\textbf{Note:} \autoref{thm:series_behavior} allows us to apply the Ratio Test to series where $\{a_n\}$ is positive for all but a finite number of terms.}

The principle of the Ratio Test is this: if $\ds\lim_{n\to\infty}\abs{\frac{a_{n+1}}{a_n}} = L<1$, then for large $n$, each term of $\{a_n\}$ is significantly smaller than its previous term which is enough to ensure convergence.  A full proof can be found at \url{http://tutorial.math.lamar.edu/Classes/CalcII/RatioTest.aspx}.

\youtubeVideo{iy8mhbZTY7g}{Using the Ratio Test to Determine if a Series Converges \#1}

\begin{example}[Applying the Ratio Test]\label{ex_ratio1}%
Use the Ratio Test to determine the convergence of the following series:
\[
 \text{1. }\sum_{n=1}^\infty \frac{2^n}{n!}\qquad\qquad
 \text{2. }\sum_{n=1}^\infty \frac{3^n}{n^3}\qquad\qquad
 \text{3. }\sum_{n=1}^\infty \frac{1}{n^2+1}.
\]
\solution
\begin{enumerate}
	\item $\ds \sum_{n=1}^\infty \frac{2^n}{n!}$:\\[-3\baselineskip]\parbox[t]{\linewidth}{%
	\begin{align*}
		\lim_{n\to\infty}\frac{2^{n+1}/(n+1)!}{2^n/n!}
		&= \lim_{n\to\infty} \frac{2^{n+1}n!}{2^n(n+1)!}\\
		&= \lim_{n\to\infty} \frac{2}{n+1}\\
		&=0.
	\end{align*}}
	Since the limit is $0<1$, by the Ratio Test $\ds\sum_{n=1}^\infty \frac{2^n}{n!}$ converges.
	
	\item	$\ds\sum_{n=1}^\infty \frac{3^n}{n^3}$:\\[-3\baselineskip]\parbox[t]{\linewidth}{%
	\begin{align*}
		\lim_{n\to\infty} \frac{3^{n+1}/(n+1)^3}{3^n/n^3}
		&= \lim_{n\to\infty}\frac{3^{n+1}n^3}{3^n(n+1)^3}\\
		&= \lim_{n\to\infty} \frac{3n^3}{(n+1)^3}\\
		&= 3.
	\end{align*}}
	Since the limit is $3>1$, by the Ratio Test $\ds\sum_{n=1}^\infty \frac{3^n}{n^3}$ diverges.
	
	\item  $\ds\sum_{n=1}^\infty \frac{1}{n^2+1}$:\\[-3\baselineskip]\parbox[t]{\linewidth}{%
	\begin{align*}
		\qquad\lim_{n\to\infty} \frac{1/\bigl((n+1)^2+1\bigr)}{1/(n^2+1)}
		&= \lim_{n\to\infty} \frac{n^2+1}{(n+1)^2+1}\\
		&= 1.
	\end{align*}}
	Since the limit is 1, the Ratio Test is inconclusive. We can easily show this series converges using the Direct or Limit Comparison Tests, with each comparing to the series $\ds \sum_{n=1}^\infty \frac{1}{n^2}$.
\end{enumerate}
\end{example}

The Ratio Test is not effective when the terms of a series \emph{only} contain algebraic functions (e.g., polynomials). It is most effective when the terms contain some factorials or exponentials. The previous example also reinforces our developing intuition: factorials dominate exponentials, which dominate algebraic functions, which dominate logarithmic functions. In Part 1 of the example, the factorial in the denominator dominated the exponential in the numerator, causing the series to converge. In Part 2, the exponential in the numerator dominated the algebraic function in the denominator, causing the series to diverge.

While we have used factorials in previous sections, we have not explored them closely and one is likely to not yet have a strong intuitive sense for how they behave. The following example gives more practice with factorials.

\begin{example}[Applying the Ratio Test]\label{ex_ratio2}%
Determine the convergence of $\ds\sum_{n=1}^\infty \frac{n!n!}{(2n)!}$.
\solution
Before we begin, be sure to note the difference between $(2n)!$ and $2n!$. When $n=4$, the former is $8!=8\cdot7\cdot\dotsm\cdot 2\cdot1=40{,}320$, whereas the latter is $2(4\cdot3\cdot2\cdot1) = 48$.

Applying the Ratio Test:
\begin{align*}
	\lim_{n\to\infty} \frac{(n+1)!(n+1)!/\bigl(2(n+1)\bigr)!}{n!n!/(2n)!}
	&= \lim_{n\to\infty}\frac{(n+1)!(n+1)!(2n)!}{n!n!(2n+2)!}\\
\intertext{Noting that \intertextmath{(2n+2)! = (2n+2)\cdot(2n+1)\cdot(2n)!}, we have}
	&= \lim_{n\to\infty}\frac{(n+1)(n+1)}{(2n+2)(2n+1)}\\
	&= 1/4.
\end{align*}
Since the limit is $1/4<1$, by the Ratio Test we conclude $\ds \sum_{n=1}^\infty \frac{n!n!}{(2n)!}$ converges.
\end{example}

\subsection{Root Test}

The final test we introduce is the Root Test, which works particularly well on series where each term is raised to a power, and does not work well with terms containing factorials. %(but not those that include factorials). 

\mnote{\textbf{Note:} We won't go into the proof, but the idea (as with the Ratio Test) is that the series is behaving enough like the geometric series $\sum_{n=0}^\infty L^n$ that we can determine convergence.}

\begin{theorem}[Root Test]\label{thm:root_test}%
Let $\{a_n\}$ be a sequence %such that there is an $N\geq 1$ where for all $n\geq N$, $a_n\geq 0$, 
where $\ds \lim_{n\to \infty}\abs{a_n}^{1/n} = L$.
\index{series!Root Comparison Test}\index{Root Comparison Test!for series}\index{convergence!Root Comparison Test}\index{divergence!Root Comparison Test}
	\begin{enumerate}
		\item If $L<1$, then $\ds\sum_{n=1}^\infty a_n$ converges.
		\item If $L>1$ or $L=\infty$, then $\ds\sum_{n=1}^\infty a_n$ diverges.
		\item If $L=1$, the Root Test is inconclusive.
	\end{enumerate}
\end{theorem}

%\mnote{\textbf{Note:} \autoref{thm:series_behavior} allows us to apply the Root Test to series where $\{a_n\}$ is positive for all but a finite number of terms.}

\mnote{\textbf{Note:} We can use L'Hôpital's Rule to show that $n^{1/n}$ approaches 1.  If someone insists on using the Root Test with factorials, it can be useful to know that $(n!)^{1/n}$ approaches infinity.}

\begin{example}[Applying the Root Test]\label{ex_root1}%
Determine the convergence of the following series using the Root Test:
\[
 \text{1.}\quad\sum_{n=1}^\infty \left(\frac{3n+1}{5n-2}\right)^n\qquad\qquad
 \text{2.}\quad\sum_{n=2}^\infty\frac{n^4}{(\ln n)^n}\qquad\qquad
 \text{3.}\quad\sum_{n=1}^\infty \frac{2^n}{n^2}.
\]
\solution
\begin{enumerate}
	\item	$\ds\lim_{n\to\infty} \left(\left(\frac{3n+1}{5n-2}\right)^n\right)^{1/n} = \lim_{n\to\infty} \frac{3n+1}{5n-2} = \frac 35.$ 
	
	Since the limit is less than 1, we conclude the series converges. Note: it is difficult to apply the Ratio Test to this series.
	
	\item	$\ds\lim_{n\to\infty} \left(\frac{n^4}{(\ln n)^n}\right)^{1/n} = \lim_{n\to\infty} \frac {\bigl(n^{1/n}\bigr)^4}{\ln n}  $. 
	
	As $n$ grows, the numerator approaches 1 (apply L'Hôpital's Rule) and the denominator grows to infinity.  Thus
	\[\lim_{n\to\infty} \frac{\bigl(n^{1/n}\bigr)^4}{\ln n} = 0.\]
	Since the limit is less than 1, we conclude the series converges.
	
	\item	$\ds \lim_{n\to\infty} \left(\frac{2^n}{n^2}\right)^{1/n} = \lim_{n\to\infty} \frac{2}{\bigl(n^{1/n}\bigr)^2} = 2$. 
	
	Since this is greater than 1, we conclude the series diverges.
\end{enumerate}
\end{example}

We end here our study of tests to determine convergence. The next section of this text provides strategies for testing series, while the \hyperref[tab_series_tests]{back of the book} contains a table summarizing the tests that one may find useful. 

While series are worthy of study in and of themselves, our ultimate goal within calculus is the study of Power Series, which we will consider in \autoref{sec:power_series}. We will use power series to create functions where the output is the result of an infinite summation. %A special type of power series is something called a Taylor Series; in the context of Taylor Series we will finally show the Alternating Harmonic Series converges to $\ln 2$.

\printexercises{exercises/08-04-exercises}

\input{text/08_Series_Strategies}
\section{Power Series}\label{sec:power_series}

So far, our study of series has examined the question of ``Is the sum of these infinite terms finite?,'' i.e., ``Does the series converge?'' We now approach series from a different perspective: as a function. Given a value of $x$, we evaluate $f(x)$ by finding the sum of a particular series that depends on $x$ (assuming the series converges). We start this new approach to series with a definition.

\begin{definition}[Power Series]\label{def:power_series}%
Let $\{a_n\}$ be a sequence, let $x$ be a variable, and let $c$ be a real number.
\index{power series}\index{series!power}
	\begin{enumerate}
		\item The \textbf{power series in $x$} is the series
		\[\sum_{n=0}^\infty a_nx^n = a_0+a_1x+a_2x^2+a_3x^3+\dotsb\]
		
		\item The \textbf{power series in $x$ centered at $c$} is the series
		\[\sum_{n=0}^\infty a_n(x-c)^n = a_0+a_1(x-c)+a_2(x-c)^2+a_3(x-c)^3+\dotsb\]
	\end{enumerate}
\end{definition}

\begin{example}[Examples of power series]\label{ex_ps1}%
Write out the first five terms of the following power series:
\[
 \text{1. }\sum_{n=0}^\infty x^n \qquad\qquad
 \text{2. }\sum_{n=1}^\infty (-1)^{n+1}\frac{(x+1)^n}n\qquad\qquad
 \text{3. }\sum_{n=0}^\infty (-1)^{n+1} \frac{(x-\pi)^{2n}}{(2n)!}.
\]
\solution
\begin{enumerate}
	\item One of the conventions we adopt is that $x^0=1$ regardless of the value of $x$. Therefore
	\[\sum_{n=0}^\infty x^n = 1+x+x^2+x^3+x^4+\dotsb\]
	This is a geometric series in $x$.
	
	\item	This series is centered at $c=-1$. Note how this series starts with $n=1$. We could rewrite this series starting at $n=0$ with the understanding that $a_0=0$, and hence the first term is $0$.
	\begin{multline*}
	\sum_{n=1}^\infty (-1)^{n+1}\frac{(x+1)^n}n \\
	= (x+1) - \frac{(x+1)^2}{2} + \frac{(x+1)^3}{3} - \frac{(x+1)^4}{4}+\frac{(x+1)^5}{5}\dotsb
	\end{multline*}
	
	\item	This series is centered at $c=\pi$. Recall that $0!=1$.
	\begin{multline*}
	\sum_{n=0}^\infty (-1)^{n+1} \frac{(x-\pi)^{2n}}{(2n)!}\\
	= -1+\frac{(x-\pi)^2}{2} - \frac{(x-\pi)^4}{24}+ \frac{(x-\pi)^6}{6!}-\frac{(x-\pi)^8}{8!}\dotsb
	\end{multline*}
\end{enumerate}
\end{example}

We introduced power series as a type of function, where a value of $x$ is given and the sum of a series is returned. Of course, not every series converges. For instance, in part 1 of \autoref{ex_ps1}, we recognized the series $\ds \sum_{n=0}^\infty x^n$ as a geometric series in $x$. \autoref{thm:geom_series} states that this series converges only when $\abs x<1$. 

This raises the question: ``For what values of $x$ will a given power series converge?,'' which  leads us to a theorem and definition.

\begin{theorem}[Convergence of Power Series]\label{thm:radius_converge}%
Let a power series $\ds \sum_{n=0}^\infty a_n(x-c)^n$ be given. Then one of the following is true:
\index{convergence!of power series}\index{power series!convergence}\index{series!power}
\begin{enumerate}
	\item The series converges only at $x=c$.
	\item	There is an $R>0$ such that the series converges for all $x$ in \\	
	$(c-R,c+R)$ and diverges for all $x<c-R$ and $x>c+R$.
	\item	The series converges for all $x$.
\end{enumerate}
\end{theorem}

%A corollary to \autoref{thm:radius_converge} is this: given a series $\ds \sum_{n=0}^\infty a_n(x-c)^n$ for some real number $c$, one of the following is true:
	%\begin{enumerate}
		%\item The series converges only at $x=c$.
		%\item	There is an $R>0$ such that the series converges for all $x$ in $(c-R, c+R)$ and diverges for all $x<c-R$ and $x>c+R$.
		%\item	The series converges for all $x$.
	%\end{enumerate}
	
The value of $R$ is important when understanding a power series, hence it is given a name in the following definition. Also, note that part 2 of \autoref{thm:radius_converge} makes a statement about the interval $(c-R,c+R)$, but the not the endpoints of that interval. A series may or may not converge at these endpoints.
	
\begin{definition}[Radius and Interval of Convergence]\label{def:radius_converge}%
\mbox{}\\[-2\baselineskip]\index{convergence!radius of}\index{convergence!interval of}\index{radius of convergence}\index{interval of convergence}\index{series!radius of convergence}\index{series!interval of convergence}
\begin{enumerate}
	\item The number $R$ given in \autoref{thm:radius_converge} is the \textbf{radius of convergence} of a given series. When a series converges for only $x=c$, we say the radius of convergence is 0, i.e.,  $R=0$. When a series converges for all $x$, we say the series has an infinite radius of convergence, i.e., $R=\infty$.
	\item	The \textbf{interval of convergence} is the set of all values of $x$ for which the series converges.
\end{enumerate}
\end{definition}

To find the values of $x$ for which a given series converges, we will use the convergence tests we studied previously (especially the Ratio Test). However, the tests all required that the terms of a series be positive. The following theorem gives us a work-around to this problem.

\begin{theorem}[The Radius of Convergence of a Series and Absolute Convergence]\label{thm:abs_power}%
%\begin{itemize}
	%\item The series $\ds \sum_{n=0}^\infty a_nx^n$ and $\ds \sum_{n=0}^\infty \abs{a_nx^n}$ have the same radius of convergence $R$.
	%\item 
	The series $\ds \sum_{n=0}^\infty a_n(x-c)^n$ and $\ds \sum_{n=0}^\infty \abs{a_n(x-c)^n}$ have the same radius of convergence $R$.
%\end{itemize}
\end{theorem}

\autoref{thm:abs_power} allows us to find the radius of convergence $R$ of a series by applying the Ratio Test (or any applicable test) to the absolute value of the terms of the series. We practice this in the following example.

\youtubeVideo{01LzAU__J-0}{Power Series --- Finding the Interval of Convergence}

\begin{example}[Determining the radius and interval of convergence.]\label{ex_ps2}%
Find the radius and interval of convergence for each of the following series:
\[
 \text{1. }\sum_{n=0}^\infty \frac{x^n}{n!} \qquad
 \text{2. }\sum_{n=1}^\infty (-1)^{n+1}\frac{x^n}{n}\qquad
 \text{3. }\sum_{n=0}^\infty 2^n(x-3)^n\qquad
 \text{4. }\sum_{n=0}^\infty n!x^n
\]
\solution
\begin{enumerate}
	\item We apply the Ratio Test to the series $\ds \sum_{n=0}^\infty \abs{\frac{x^n}{n!}}$:
		\begin{align*}
		\lim_{n\to\infty} \frac{\abs{x^{n+1}/(n+1)!}}{\abs{x^n/n!}} &= \lim_{n\to\infty} \abs{\frac{x^{n+1}}{x^n}\cdot\frac{n!}{(n+1)!}}\\
			&= \lim_{n\to\infty} \abs{\frac x{n+1}}\\
			&= 0 \text{ for all } x.
		\end{align*}
		The Ratio Test shows us that regardless of the choice of $x$, the series converges. Therefore the radius of convergence is $R=\infty$, and the interval of convergence is $(-\infty,\infty)$.
		
	\item		We apply the Ratio Test to the series $\ds \sum_{n=1}^\infty \abs{(-1)^{n+1}\frac{x^n}{n}} = \sum_{n=1}^\infty \abs{\frac{x^n}{n}}$:
	\begin{align*}
	\lim_{n\to\infty} \frac{\abs{x^{n+1}/(n+1)}}{\abs{x^n/n}} &= \lim_{n\to\infty} \abs{\frac{x^{n+1}}{x^n}\cdot \frac{n}{n+1}} \\
			&= \lim_{n\to\infty} \abs x\frac{n}{n+1}\\
			&= \abs x.
	\end{align*}
	The Ratio Test states a series converges if the limit of $\abs{a_{n+1}/a_n} = L<1$. We found the limit above to be $\abs x$; therefore, the power series converges when $\abs x<1$, or when $x$ is in $(-1,1)$. Thus the radius of convergence is $R=1$.
	
	To determine the interval of convergence, we need to check the endpoints of $(-1,1)$. When $x=-1$, we have the opposite of the Harmonic Series:
	\begin{align*}
	\sum_{n=1}^\infty (-1)^{n+1}\frac{(-1)^n}{n}
	&= \sum_{n=1}^\infty \frac{(-1)^{2n+1}}{n}\\
	&= \sum_{n=1}^\infty \frac{-1}{n}\\
	&= -\infty.
	\end{align*}
	The series diverges when $x=-1$.
	
	When $x=1$, we have the series $\ds \sum_{n=1}^\infty (-1)^{n+1}\frac{(1)^n}{n}$, which is the Alternating Harmonic Series, which converges. Therefore the interval of convergence is $(-1,1]$.
	
	\item		We apply the Ratio Test to the series $\ds\sum_{n=0}^\infty \abs{2^n(x-3)^n}$:
	\begin{align*}
	\lim_{n\to\infty} \frac{\abs{2^{n+1}(x-3)^{n+1}}}{\abs{2^n(x-3)^n}} &= \lim_{n\to\infty} \abs{\frac{2^{n+1}}{2^n}\cdot\frac{(x-3)^{n+1}}{(x-3)^n}}\\
			&=\lim_{n\to\infty} \abs{2(x-3)}.
	\end{align*}
	
According to the Ratio Test, the series converges when $\abs{2(x-3)}<1 \implies \abs{x-3}< 1/2$. The series is centered at 3, and $x$ must be within $1/2$ of 3 in order for the series to converge. Therefore the radius of convergence is $R=1/2$, and we know that the series converges absolutely for all $x$ in $(3-1/2,3+1/2) = (2.5, 3.5)$.

We check for convergence at the endpoints to find the interval of convergence. When $x=2.5$, we have:
\begin{align*}
\sum_{n=0}^\infty 2^n(2.5-3)^n &= \sum_{n=0}^\infty 2^n(-1/2)^n \\
			&=\sum_{n=0}^\infty (-1)^n,
\end{align*}
which diverges. A similar process shows that the series also diverges at $x=3.5$. Therefore the interval of convergence is $(2.5, 3.5)$.

\item		We apply the Ratio Test to $\ds \sum_{n=0}^\infty \abs{n!x^n}$:
\begin{align*}
\lim_{n\to\infty} \frac{\abs{ (n+1)!x^{n+1}}}{\abs{n!x^n}} &= \lim_{n\to\infty} \abs{(n+1)x}\\
		&= \infty\ \text{ for all $x$, except $x=0$.}
\end{align*}

The Ratio Test shows that the series diverges for all $x$ except $x=0$. Therefore the radius of convergence is $R=0$.
\end{enumerate}
\end{example}

\subsection{Power Series as Functions}

We can use a power series to define a function:
\[f(x) = \sum_{n=0}^\infty a_nx^n\]
where the domain of $f$ is a subset of the interval of convergence of the power series. One can apply calculus techniques to such functions; in particular, we can find derivatives and antiderivatives. 

\begin{theorem}[Derivatives and Indefinite Integrals of Power Series Functions]\label{thm:calc_power_series}%
Let $\ds f(x) = \sum_{n=0}^\infty a_n(x-c)^n$ be a function defined by a power series, with radius of convergence $R$.
\index{series!power!derivatives and integrals}\index{integration!of power series}\index{derivative!power series}\index{power series!derivatives and integrals}
\begin{enumerate}
	\item $f(x)$ is continuous and differentiable on $(c-R,c+R)$.
	\item	$\ds \fp(x) = \sum_{n=1}^\infty a_nn(x-c)^{n-1}$, with radius of convergence $R$.
	\item	$\ds \int f(x)\dd x = C+\sum_{n=0}^\infty a_n\frac{(x-c)^{n+1}}{n+1}$, with radius of convergence $R$.
\end{enumerate}
\end{theorem}

A few notes about \autoref{thm:calc_power_series}:
\begin{enumerate}
	\item The theorem states that differentiation and integration do not change the radius of convergence. It does not state anything about the \emph{interval} of convergence. They are not always the same.
	\item	Notice how the summation for $\fp(x)$ starts with $n=1$. This is because the constant term $a_0$ of $f(x)$ goes to 0.
	\item	Differentiation and integration are simply calculated term-by-term using previous rules of integration and differentiation.
\end{enumerate}

\begin{example}[Derivatives and indefinite integrals of power series]\label{ex_ps3}%
Let $\ds f(x) = \sum_{n=0}^\infty x^n$. Find the following along with their respective intervals of convergence.
\[
 \text{1.}\quad\fp(x)
 \qquad\text{and}\qquad
 \text{2.}\quad\ds F(x) =\int f(x)\dd x
\]
\solution
We find the derivative and indefinite integral of $f(x)$, following \autoref{thm:calc_power_series}.

\begin{enumerate}
\item\hfill$\begin{aligned}[t]
f(x)&=1+x+\,x^2+\,x^3+\,x^4+\dotsb = \sum_{n=0}^\infty x^n\\
\fp(x) &= 0+1+2x+3x^2+4x^3+\dotsb=\sum_{n=1}^\infty nx^{n-1} 
\end{aligned}$\hfill\null

In \autoref{ex_ps1}, we recognized that $\ds \sum_{n=0}^\infty x^n$ is a geometric series in $x$. We know that such a geometric series converges when $\abs x<1$; that is, the interval of convergence is $(-1,1)$.

To determine the interval of convergence of $\fp(x)$, we consider the endpoints of $(-1,1)$.
When $x=-1$ we have
\[\fp(-1) = \sum_{n=1}^\infty n(-1)^{n-1}\]
which diverges by the Test for Divergence
and when $x=1$ we have
\[\fp(1) = \sum_{n=1}^\infty n\]
which also diverges by the Test for Divergence. Therefore, the interval of convergence of $\fp(x)$ is $(-1,1)$. 

\item\hfill$\begin{aligned}[t]
 f(x)&=\phantom{C+{}}1+\,x+\,x^2+\,x^3+\dotsb = \sum_{n=0}^\infty x^n\\
 F(x) = \int f(x)\dd x &= C+ x+\frac{x^2}{2}+\frac{x^3}3+\frac{x^4}4+\dotsb \\
 &= C+\sum_{n=0}^\infty \frac{x^{n+1}}{n+1}=C+\sum_{n=1}^\infty \frac{x^{n}}{n}  
\end{aligned}$\hfill\null

To find the interval of convergence of $F(x)$, we again consider the endpoints of $(-1,1)$.
When $x=-1$ we have
\[F(-1) = C+\sum_{n=1}^\infty \frac{(-1)^{n}}{n}\]
The value of $C$ is irrelevant; notice that the rest of the series is an Alternating Series that whose terms converge to 0. By the Alternating Series Test, this series converges. (In fact, we can recognize that the terms of the series after $C$ are the opposite of the Alternating Harmonic Series. We can thus say that $F(-1) = C-\ln 2$.)
\[F(1) = C+\sum_{n=1}^\infty \frac{1}{n} \]
Notice that this summation is $C\ +$ the Harmonic Series, which diverges. Since $F$ converges for $x=-1$ and diverges for $x=1$, the interval of convergence of $F(x)$ is $[-1,1)$.
\end{enumerate}
\end{example}

The previous example showed how to take the derivative and indefinite integral of a power series without motivation for why we care about such operations. We may care for the sheer mathematical enjoyment ``that we can'', which is motivation enough for many. However, we would be remiss to not recognize that we can learn a great deal from taking derivatives and indefinite integrals.\bigskip

Recall that $\ds f(x) = \sum_{n=0}^\infty x^n$ in \autoref{ex_ps3} is a geometric series. According to \autoref{thm:geom_series}, this series converges to $1/(1-x)$ when $\abs x<1$. Thus we can say
\begin{equation}
f(x) = \sum_{n=0}^\infty x^n = \frac 1{1-x},\quad \text{ on }\quad (-1,1).
\label{eq:power_series_geometric}
\end{equation}

Integrating the power series, (as done in \autoref{ex_ps3},) we find
\begin{equation}
F(x)  = C_1+\sum_{n=0}^\infty \frac{x^{n+1}}{n+1},\label{eq:ps3a}
\end{equation}
while integrating the function $f(x) = 1/(1-x)$ gives
\begin{equation}
F(x)  = -\ln\abs{1-x} + C_2.\label{eq:ps3b}
\end{equation}

Equating Equations \eqref{eq:ps3a} and \eqref{eq:ps3b}, we have 
\[F(x) = C_1+\sum_{n=0}^\infty \frac{x^{n+1}}{n+1} = -\ln\abs{1-x} + C_2.\]
Letting $x=0$, we have $F(0) = C_1 = C_2$. This implies that we can drop the constants and conclude
\[\sum_{n=0}^\infty \frac{x^{n+1}}{n+1} = -\ln\abs{1-x}.\]
We established in \autoref{ex_ps3} that the series $\ds\sum_{n=0}^\infty \frac{x^{n+1}}{n+1}$ converges at $x=-1$; substituting $x=-1$ on both sides of the above equality gives
\[-1+\frac12-\frac13+\frac14-\frac15+\dotsb = -\ln 2.\]
On the left we have the opposite of the Alternating Harmonic Series; on the right, we have $-\ln 2$. We conclude that 
\[1-\frac12+\frac13-\frac14+\dotsb = \ln 2.\]
In \autoref{ex_alt1} of \autoref{sec:alt_series} we said the Alternating Harmonic Series converges to $\ln 2$, but did not show why this was the case. The work above shows how we conclude that the Alternating Harmonic Series Converges to $\ln 2$. \index{Alternating Harmonic Series}

We use this type of analysis in the next example.

\begin{example}[Analyzing power series functions]\label{ex_ps4}%
Let $\ds f(x) = \sum_{n=0}^\infty \frac{x^n}{n!}$. Find $\ds \fp(x)$ and $\ds \int f(x)\dd x$, and use these to analyze the behavior of $f(x)$.
\solution
We start by making two notes: first, in \autoref{ex_ps2}, we found the interval of convergence of this power series is $(-\infty,\infty)$. Second, we will find it useful later to have a  few terms of the series written out:
\begin{equation}
\sum_{n=0}^\infty \frac{x^n}{n!}
= 1 + x + \frac{x^2}2+\frac{x^3}{6} + \frac{x^4}{24} +\dotsb\label{eq:ps4}
\end{equation}

We now find the derivative:
\begin{align*}
\fp(x) &= \sum_{n=1}^\infty n\frac{x^{n-1}}{n!} \\
&=\sum_{n=1}^\infty \frac{x^{n-1}}{(n-1)!} = 1+x+\frac{x^2}{2!}+\dotsb. \\
\intertext{Since the series starts at \intertextmath{n=1} and each term refers to \intertextmath{(n-1)}, we can re-index the series starting with \intertextmath{n=0}:}
		&= \sum_{n=0}^\infty \frac{x^{n}}{n!}\\
		&= f(x).
\end{align*}
We found the derivative of $f(x)$ is $f(x)$. The only functions for which this is true are of the form $y=ce^x$ for some constant $c$. As $f(0) = 1$ (see \autoeqref{eq:ps4}), $c$ must be 1. Therefore we conclude that 
\[f(x) = \sum_{n=0}^\infty \frac{x^n}{n!} = e^x\]% \quad\text{for all $x$}.
for all $x$.

We can also find $\ds \int f(x)\dd x$:
\begin{align*}
\int f(x)\dd x &= C+\sum_{n=0}^\infty \frac{x^{n+1}}{n!(n+1)} \\
				&= C+ \sum_{n=0}^\infty \frac{x^{n+1}}{(n+1)!}
\end{align*}
We write out a few terms of this last series:
\[
C+ \sum_{n=0}^\infty \frac{x^{n+1}}{(n+1)!}
= C+ x+ \frac{x^2}2+\frac{x^3}{6}+\frac{x^4}{24}+\dotsb
\]
The integral of $f(x)$ differs from $f(x)$ only by a constant, again indicating that $f(x) = e^x$.
\end{example}


\autoref{ex_ps4} and the work following \autoref{ex_ps3} established relationships between a power series function and ``regular'' functions that we have dealt with in the past. In general, given a power series function, it is difficult (if not impossible) to express the function in terms of elementary functions. We chose examples where things worked out nicely.

%In this section's last example, we show how to solve a simple differential equation with a power series.
%
%\begin{example}[Solving a differential equation with a power series.]\label{ex_ps5}%
%Give the first 4 terms of the power series solution to $y\primeskip' = 2y$, where $y(0) = 1$.
%\solution
%The differential equation $y\primeskip' = 2y$ describes a function $y=f(x)$ where the derivative of $y$ is twice $y$ and $y(0)=1$. This is a rather simple differential equation; with a bit of thought one should realize that if $y=Ce^{2x}$, then $y\primeskip' = 2Ce^{2x}$, and hence $y\primeskip' = 2y$. By letting $C=1$ we satisfy the initial condition of $y(0)=1$.
%
%Let's ignore the fact that we already know the solution and find a power series function that satisfies the equation. The solution we seek will have the form
%\[f(x)  = \sum_{n=0}^\infty a_nx^n = a_0+a_1x+a_2x^2+a_3x^3+\dotsb\]
%for unknown coefficients $a_n$. We can find $\fp(x)$ using \autoref{thm:calc_power_series}:
%\[\fp(x) = \sum_{n=1}^\infty a_n\cdot n\cdot x^{n-1} = a_1+2a_2x+3a_3x^2+4a_4x^3\dotsb.\]
%Since $\fp(x) = 2f(x)$, we have
%\begin{align*}
%a_1+2a_2x+3a_3x^2+4a_4x^3\dotsb &= 2\bigl(a_0+a_1x+a_2x^2+a_3x^3+\dotsb\bigr)\\
%			&=2a_0+2a_1x+2a_2x^2+2a_3x^3+\dotsb
%\end{align*}
%The coefficients of like powers of $x$ must be equal, so we find that
%\[a_1 = 2a_0,\quad 2a_2 = 2a_1,\quad 3a_3 = 2a_2,\quad 4a_4 = 2a_3,\quad \text{etc.}\]
%The initial condition $y(0) = f(0) = 1$ indicates that $a_0 = 1$; with this, we can find the values of the other coefficients:
%\begin{align*}
%a_0 = 1 \text{ and } a_1=2a_0 &\Rightarrow a_1 = 2;\\
%a_1 = 2 \text{ and } 2a_2 = 2a_1 &\Rightarrow a_2=4/2 =2;\\
%a_2=2 \text{ and } 3a_3 = 2a_2 &\Rightarrow a_3=8/(2\cdot3)=4/3;\\
%a_3=4/3 \text{ and } 4a_4 = 2a_3 &\Rightarrow a_4 =16/(2\cdot3\cdot4)= 2/3. 
%\end{align*}
%Thus the first 5 terms of the power series solution to the differential equation $y\primeskip'=2y$ is 
%\[f(x) = 1+ 2x+2x^2 + \frac43x^3+\frac23x^4+\dotsb\]
%In \autoref{sec:taylor_series}, as we study Taylor Series, we will learn how to recognize this series as describing $y=e^{2x}$.
%\end{example}

\subsection{Representations of Functions with Power Series}

%Our last example illustrates that it
It can be difficult to recognize an elementary function by its power series expansion. It is far easier to start with a known function, expressed in terms of elementary functions, and represent it as a power series function. One may wonder why we would bother doing so, as the latter function probably seems more complicated.
%In the next two sections, we show both \emph{how} to do this and \emph{why} such a process can be beneficial.

%We opened the last section by saying that we were going to start thinking about applications of series and then promptly spent the section talking about convergence again.  It's now time to actually start with the applications of series. With this section  we will start talking about how to represent functions with power series.  The natural question of why we might want to do this will be answered later once we actually learn how to do this.

Let's start off with a series we already know how to do, although when we first ran across this series we didn't think of it as a power series nor did we acknowledge that it represented a function. Recall that the geometric series is
\[\sum_{n=0}^\infty ar^n = \frac a{1-r} \qquad \text{provided }\abs r<1.\]
We also know that if $\abs r\geq 1$ the series diverges. Now, if we take $a=1$ and $r=x$ this becomes,
\begin{equation}
\label{eq_power_eq}
\sum_{n=0}^\infty x^n=\frac1{1-x}\qquad\text{provided }\abs x<1
\end{equation}

Turning this around we can see that we can represent the function
\begin{equation}
\label{eq_power_func}
f(x) = \frac1{1-x}
\end{equation}
with the power series
\begin{equation}
\label{eq_power_series}
\sum_{n=0}^{\infty}x^n \qquad \text{provided }\abs x<1.
\end{equation}

This provision is important.  We can clearly plug any number other than $x=1$ into the function, however, we will only get a convergent power series if $\abs x<1$.  This means the equality in \autoeqref{eq_power_eq} will only hold if $\abs x<1$.  For any other value of $x$ the equality won't hold.  Note as well that we can also use this to acknowledge that the radius of convergence of this power series is $R=1$ and the interval of convergence is $\abs x<1$.

This idea of convergence is important here.  We will be representing many functions as power series and it will be important to recognize that the representations will often only be valid for a range of $x$'s and that there may be values of $x$ that we can plug into the function that we can't plug into the power series representation.

In this section we are going to concentrate on representing functions with power series where the function can be related back to a geometric series. In this way we will hopefully become familiar with some of the kinds of manipulations that we will sometimes need when working with power series. We will see in \autoref{sec:taylor_series} that this strategy is useful for integrating functions that don't have elementary antiderivatives.

%In this section we are going to concentrate on representing functions with power series where the function can be related back to \autoeqref{eq_power_func}.  

%In this way we will hopefully become familiar with some of the kinds of manipulations that we will sometimes need when working with power series. We now consider several examples.

\begin{example}[Finding a Power Series]\label{ex_power_cube}%
Find a power series representation for $g(x) = \dfrac1{1+x^3}$ and determine its interval of convergence.
\solution
We want to relate this function back to \autoeqref{eq_power_func}.  This is actually easier than it might look.  Recall that the $x$ in \autoeqref{eq_power_func} is simply a variable and can represent anything.  So, a quick rewrite of $g(x)$ gives,
\[g(x)=\frac1{1-(-x^3)}\]
and so the $-x^3$ holds the same place as the $x$ in \autoeqref{eq_power_func}.  Therefore, all we need to do is replace the $x$ in \autoeqref{eq_power_series} and we've got a power series representation for $g(x)$.  
\[g(x)=\sum_{n=0}^{\infty}\left( -x^3 \right)^n \qquad \text{provided }\abs{-x^3}<1\]

Notice that we replaced both the $x$ in the power series and in the interval of convergence. All we need to do now is a little simplification.
\[g(x)=\sum_{n=0}^{\infty}\left( -1 \right)^n x^{3n} \qquad\text{provided }\abs x<1\]
So, in this case the interval of convergence is the same as the original power series.  This usually won't happen.  More often than not the new interval of convergence will be different from the original interval of convergence.
\end{example}

\begin{example}[Finding a Power Series]\label{ex_power_numerator}%
Find a power series representation for $h(x)=\dfrac{2x^2}{1+x^3}$ and determine its interval of convergence.
\solution
This function is similar to the previous function, however the numerator is different.  Since \autoeqref{eq_power_func} doesn't have an $x$ in the numerator it appears that we can't relate this function back to that.  However, now that we've worked the first example this one is actually very simple since we can use the result of the answer from that example.  To see how to do this let's first rewrite the function a little. 
\[h(x)=2x^2\frac1{1+x^3}.\]
Now, from the first example we've already got a power series for the second term so let's use that to write the function as, 
\[
h(x)=2x^2\sum_{n=0}^\infty\left( -1 \right)^n x^{3n} \qquad\text{provided }\abs x<1
\]

Notice that the presence of $x$'s outside of the series will NOT affect its convergence and so the interval of convergence remains the same.  
The last step is to bring the coefficient into the series and we'll be done.  When we do this, make sure to combine the $x$'s as well.  We typically only want a single $x$ in a power series.
\[
h(x)=\sum_{n=0}^\infty 2\left( -1 \right)^n x^{3n+2} \qquad\text{provided }\abs x<1.
\]
As we saw in the previous example we can often use previous results to help us out.  This is an important idea to remember as it can often greatly simplify our work.
\end{example}

\begin{example}[Finding a Power Series]\label{ex_power_five}%
Find a power series representation for $f(x)=\dfrac x{5-x}$ and determine its interval of convergence.
\solution
So again, we have an $x$ in the numerator. As with the last example factor $x$ out and we have $f(x)=x\dfrac1{5-x}$. If we had a power series representation for $g(x)=\dfrac1{5-x}$ we could get a power series representation for $f(x)$. We need the number in the denominator to be a one so we rewrite the denominator.  
\[g(x)=\frac15\frac1{1-\frac x5}\]

Now all we need to do to get a power series representation is to replace the $x$ in \autoeqref{eq_power_series} with $\dfrac x5$.  Doing this gives
\[
g(x)=\frac15\sum_{n=0}^\infty\left(\frac x5\right)^n
\qquad\text{provided } \abs{\frac x5}<1.
\]

Now simplify the series.
\begin{align*}
	g(x)
	& = \frac15\sum_{n=0}^\infty\frac{x^n}{5^n} \\
	& = \sum_{n=0}^\infty\frac{x^n}{5^{n+1}} 
\end{align*}

The interval of convergence for this series is
\[
\abs{\frac x5}<1 \qquad \Rightarrow \qquad \frac15\abs x<1
\qquad \Rightarrow \qquad\abs x<5
\]

We now have a power series representation for $g(x)$ but we need to find a power series representation for the original function.  All we need to do for this is to multiply the power series representative for $g(x)$ by $x$ and we'll have it.  
\begin{align*}
	f(x)
	&= x \frac1{5-x} \\
	&= x \sum_{n=0}^\infty\frac{x^n}{5^{n+1}}\\ 
	& = \sum_{n=0}^\infty\frac{x^{n+1}}{5^{n+1}}
\end{align*}
The interval of convergence doesn't change and so it will be $\abs x<5$.
\end{example}

\begin{example}[Re-indexing a Power Series]\label{ex_power_reindex}%
Find a power series representation for $f(x)=\dfrac{1+x}{1-x}$.
\solution
We can start by writing this as
\[f(x)=(1+x)\sum_{n=0}^\infty x^n.\]
The problem with this representation is that the $1+x$ makes this not a power series --- we only want a single occurrence of $x$.  The proceed, we'll split this into two sums
\[
f(x)
=\sum_{n=0}^\infty x^n+x\sum_{n=0}^\infty x^n
=\sum_{n=0}^\infty x^n+\sum_{n=0}^\infty x^{n+1}.
\]
Because the second sum uses $n+1$, we can re-index it to get a sum involving $x^n$, which is our power series
\[
f(x)
=\sum_{n=0}^\infty x^n+x\sum_{n=0}^\infty x^n
=\sum_{n=0}^\infty x^n+\sum_{n=1}^\infty x^n
=1+2\sum_{n=1}^\infty x^n.
\]
\end{example}

We now consider several examples where differentiation and integration of power series from \autoref{thm:calc_power_series} are used to write the power series for a function. % in \autoref{sec:power_series}.

\begin{example}[Differentiating a Power Series]\label{ex_power_deriv}%
Find a power series representation for $g(x)=\dfrac1{(1-x)^2}$ and determine its radius of convergence.
\solution
We know that
\[\frac1{(1-x)^2} = \frac{\dd}{\dd x} \left( \frac{1}{1-x} \right).\]
Since we have a power series  representation for $\dfrac1{1-x}$, we can differentiate that power series to get a power series representation for $g(x)$.  
\begin{align*}
	g(x)
	& = \frac{1}{(1-x)^2} \\
    & = \frac{\dd}{\dd x} \left( \frac{1}{1-x} \right) \\
	& =  \frac{\dd}{\dd x} \left( \sum_{n=0}^{\infty} x^n \right)\\
    & =  \sum_{n=1}^{\infty} nx^{n-1}
\end{align*}

Since the original power series had a radius of convergence of $R=1$ the derivative, and hence $g(x)$, will also have a radius of convergence of $R=1$.
\end{example}

\begin{example}[Integrating a Power Series]\label{ex_power_int}%
Find a power series representation for $h(x)=\ln(5-x)$ and determine its radius of convergence.
\solution
In this case we need the fact that
\[\int\frac1{5-x}\dd x=-\ln(5-x).\]
Recall that we found a power series representation for $\dfrac1{5-x}$ in \autoref{ex_power_five}. We now have
\begin{align*}
	\ln(5-x)
	& =  -\ds\int \frac{1}{5-x}\dd x \\
    & =  -\ds\int \sum_{n=0}^{\infty} \frac{x^n}{5^{n+1}}\dd x \qquad\text{where $\abs x<5$}\\
	& =  C - \sum_{n=0}^{\infty} \frac{x^{n+1}}{(n+1)5^{n+1}} \qquad\text{where $\abs x<5$}
\end{align*}

We can find the constant of integration, $C$, by substituting in a value of $x$.  A good choice is $x=0$ as the series is usually easy to evaluate there.  
\begin{align*}
	\ln(5-0) & = C - \sum_{n=0}^{\infty} \frac{0^{n+1}}{(n+1)5^{n+1}} \\
	\ln (5-0) & =  C 
\end{align*}

So, the final answer is, 
\[\ln(5-x) =\ln(5) - \sum_{n=0}^{\infty} \frac{x^{n+1}}{(n+1)5^{n+1}},\]
and the radius of convergence is 5. Notice that $x=-5$ allows for convergence so the interval of convergence is $[-5,5)$.
\end{example}

% todo write a transition paragraph from this section to the next

\printexercises{exercises/08-06-exercises}

%\input{text/08_Series_Functions}
\section{Taylor Polynomials}\label{sec:taylor_poly}

Consider a function $y=f(x)$ and a point $\bigl(c,f(c)\bigr)$. The derivative, $\fp(c)$, gives the instantaneous rate of change of $f$ at $x=c$. Of all lines that pass through the point $\bigl(c,f(c)\bigr)$, the line that best approximates $f$ at this point is the tangent line; that is, the line whose slope (rate of change) is $\fp(c)$.

\mtable{Plotting $y=f(x)$ and a table of derivatives of $f$ evaluated at 0.}{fig:taypolyintroa}{\begin{tikzpicture}[alt={Blue curve f; red line p_1 is the tangent at x = 0.}]
\begin{axis}[width=1.16\marginparwidth,tick label style={font=\scriptsize},
axis y line=middle,axis x line=middle,name=myplot,axis on top,
ymin=-5.5,ymax=8.5,xmin=-4.2,xmax=4.4]
\addplot [draw={\colorone},domain=-4:4,smooth,thick,samples=50] {exp(x)*sin(deg(x))*cos(deg(x))+2};
\addplot [draw={\colortwo},domain=-4:4,thick] {x+2};
\draw (axis cs:-3.,3) node {\scriptsize $y=f(x)$};
\draw (axis cs:-2.75,-2.2) node {\scriptsize $y=p_1(x)$};
\end{axis}
\node [right] at (myplot.right of origin) {\scriptsize $x$};
\node [above] at (myplot.above origin) {\scriptsize $y$};
\end{tikzpicture}
\begin{align*}
f(0)&=2 & \fp''(0)&=-1\\
\fp(0)&=1 & f\,^{(4)}(0)&=-12 \\
\fpp(0)&=2 & f\,^{(5)}(0)&=-19\vspace{-.5\baselineskip}
\end{align*}}

In \autoref{fig:taypolyintroa}, we see a function $y=f(x)$ graphed. The table below the graph shows that $f(0)=2$ and $\fp(0) = 1$; therefore, the tangent line to $f$ at $x=0$ is $p_1(x) = 1(x-0)+2 = x+2$. The tangent line is also given in the figure. Note that ``near'' $x=0$, $p_1(x) \approx f(x)$; that is, the tangent line approximates $f$ well.

One shortcoming of this approximation is that the tangent line only matches the slope of $f$; it does not, for instance, match the concavity of $f$. We can find a polynomial, $p_2(x)$, that does match the concavity without much difficulty, though. The table in \autoref{fig:taypolyintroa} gives the following information:
\[f(0) = 2 \qquad \fp(0) = 1\qquad \fp'(0) = 2.\]
Therefore, we want our polynomial $p_2(x)$ to have these same properties. That is, we need
\[p_2(0) = 2 \qquad p_2'(0) = 1 \qquad p_2''(0) = 2.\]

This is simply an initial-value problem. We can solve this using the techniques first described in \autoref{sec:antider}. To keep $p_2(x)$ as simple as possible, we'll assume that not only  $p_2''(0)=2$, but that $p_2''(x)=2$. That is, the second derivative of $p_2$ is  constant.

\mtable{Plotting $f$, $p_2$, and $p_4$.}{fig:taypolyintrob}{\begin{tikzpicture}[alt={Blue f with two red polynomials: cubic p_2 and quartic p_4; higher-degree curve hugs f near 0 then departs.}]
\begin{axis}[width=1.16\marginparwidth,tick label style={font=\scriptsize},
axis y line=middle,axis x line=middle,name=myplot,axis on top,
ymin=-5.5,ymax=8.5,xmin=-4.2,xmax=4.4]
\addplot [draw={\colorone},domain=-4:4,smooth,thick,samples=50]
 {exp(x)*sin(deg(x))*cos(deg(x))+2};
\addplot [draw={\colortwo},domain=-4:4,thick,smooth] {x^2+x+2};
\addplot [draw={\colortwo!40},domain=-4:4,thick,smooth] {-x^4/2-x^3/6+x^2+x+2};
\draw (axis cs:-3.,3.5) node {\scriptsize $y=p_2(x)$};
\draw (axis cs:-3.1,-2.5) node {\scriptsize $y=p_4(x)$};
\end{axis}
\node [right] at (myplot.right of origin) {\scriptsize $x$};
\node [above] at (myplot.above origin) {\scriptsize $y$};
\end{tikzpicture}}

If $p_2''(x) = 2$, then $p_2'(x) = 2x+C$ for some constant $C$. Since we have determined that $p_2'(0) = 1$, we find that $C=1$ and so $p_2'(x) = 2x+1$. Finally, we can compute $p_2(x) = x^2+x+C$. Using our initial values, we know $p_2(0) = 2$ so $C=2.$ We conclude that $p_2(x) = x^2+x+2.$ This function is plotted with $f$ in \autoref{fig:taypolyintrob}.

We can repeat this approximation process by creating polynomials of higher degree that match more of the derivatives of $f$ at $x=0$. In general, a polynomial of degree $n$ can be created to match the first $n$ derivatives of $f$. \autoref{fig:taypolyintrob} also shows $p_4(x)= -x^4/2-x^3/6+x^2+x+2$, whose first four derivatives at 0 match those of $f$. (Using the table in \autoref{fig:taypolyintroa}, start with $p_4^{(4)}(x)=-12$ and solve the related initial-value problem.)

\mtable{Plotting $f$ and $p_{13}$.}{fig:taypolyintroc}{\begin{tikzpicture}[alt={Blue f alongside red quintic p_s; p₅ wiggles but tracks f closely on −1 < x < 4}]
\begin{axis}[width=1.16\marginparwidth,tick label style={font=\scriptsize},
axis y line=middle,axis x line=middle,name=myplot,axis on top,
ymin=-5.5,ymax=8.5,xmin=-4.2,xmax=4.4]
\addplot [draw={\colorone},domain=-4:4,smooth,thick,samples=50] {exp(x)*sin(deg(x))*cos(deg(x))+2};
\addplot [draw={\colortwo},domain=-4:4,thick,smooth] coordinates {(-3.52,-6.364)(-3.36,-2.334)(-3.2,-0.1706)(-3.04,0.9569)(-2.88,1.528)(-2.72,1.809)(-2.56,1.944)(-2.4,2.009)(-2.24,2.038)(-2.08,2.048)(-1.92,2.046)(-1.76,2.031)(-1.6,2.006)(-1.44,1.969)(-1.28,1.924)(-1.12,1.872)(-0.96,1.82)(-0.8,1.775)(-0.64,1.747)(-0.48,1.747)(-0.32,1.783)(-0.16,1.866)(0,2.)(0.16,2.185)(0.32,2.411)(0.48,2.662)(0.64,2.908)(0.8,3.112)(0.96,3.227)(1.12,3.202)(1.28,2.988)(1.44,2.546)(1.6,1.855)(1.76,0.9264)(1.92,-0.1926)(2.08,-1.406)(2.24,-2.565)(2.4,-3.474)(2.56,-3.891)(2.72,-3.542)(2.88,-2.133)(3.04,0.6461)(3.2,5.139)(3.36,11.79)};
% the degree 13 polynomial
\draw (axis cs:-2,-2.75) node {\scriptsize $y=p_{13}(x)$};
\end{axis}
\node [right] at (myplot.right of origin) {\scriptsize $x$};
\node [above] at (myplot.above origin) {\scriptsize $y$};
\end{tikzpicture}}

As we use more and more derivatives, our polynomial approximation to $f$ gets better and better. In this example, the interval on which the approximation is ``good'' gets bigger and bigger. \autoref{fig:taypolyintroc} shows $p_{13}(x)$; we can visually affirm that this polynomial approximates $f$ very well on $[-2,3]$. The polynomial $p_{13}(x)$ is fairly complicated:
\[\scriptstyle
\frac{16901x^{13}}{6227020800}+\frac{13x^{12}}{1209600}-\frac{1321x^{11}}{39916800}-\frac{779x^{10}}{1814400}-\frac{359x^9}{362880}+\frac{x^8}{240}+\frac{139x^7}{5040}+\frac{11 x^6}{360}-\frac{19x^5}{120}-\frac{x^4}{2}-\frac{x^3}{6}+x^2+x+2.
\]

The polynomials we have created are examples of \emph{Taylor polynomials}, named after the British mathematician Brook Taylor who made important discoveries about such functions. While we created the above Taylor polynomials by solving initial-value problems, it can be shown that Taylor polynomials follow a general pattern that makes their formation much more direct. This is described in the following definition.

{
\tcbset{grow to right by=10em}
\begin{definition}[Taylor Polynomial, Maclaurin Polynomial]\label{def:taypoly}%
Let $f$ be a function whose first $n$ derivatives exist at $x=c$.
\index{Taylor Polynomial!definition}\index{Maclaurin Polynomial!definition} \index{Maclaurin Polynomial|seealso{Taylor Polynomial}}
\begin{enumerate}
	\item	The \textbf{Taylor polynomial of degree $n$ of $f$ at $x=c$} is 
	\begin{align*}
	p_n(x)
	&= f(c) + \fp(c)(x-c) + \frac{\fpp(c)}{2!}(x-c)^2+\frac{\fp''(c)}{3!}(x-c)^3+\dotsb+\frac{f\,^{(n)}(c)}{n!}(x-c)^n \\
	&=\sum_{k=0}^n\frac{f\,^{(k)}(c)}{k!}(x-c)^k.
	\end{align*}
	\item	A special case of the Taylor polynomial is the Maclaurin polynomial, where $c=0$. That is, the \textbf{Maclaurin polynomial of degree $n$ of $f$} is 
	\begin{align*}
	p_n(x)
	&= f(0) + \fp(0)x + \frac{\fpp(0)}{2!}x^2+\frac{\fp''(0)}{3!}x^3+\dotsb+\frac{f\,^{(n)}(0)}{n!}x^n \\
	&=\sum_{k=0}^n\frac{f\,^{(k)}(0)}{k!}x^k.
	\end{align*}
\end{enumerate}
\end{definition}
}

\mnote{\textbf{Note:} The summations in this definition use the convention that $x^0=1$ even when $x=0$ and that $f^{(0)}=f$.  They also use the definition that $0!=1$.}
Generally, we order the terms of a polynomial to have decreasing degrees, and that is how we began this section.  This definition, and the rest of this chapter, reverses this order to reflect the greater importance of the lower degree terms in the polynomials that we will be finding.

\youtubeVideo{UINFWG0ErSA}{Taylor Polynomial to Approximate a Function, Ex 3}

We will practice creating Taylor and Maclaurin polynomials in the following examples.

\begin{example}[Finding and using Maclaurin polynomials]\label{ex_taypoly1}%
\mbox{}\\[-2\baselineskip]\parbox[t]{\linewidth}{%
\begin{enumerate}
	\item	Find the $n^{\text{th}}$ Maclaurin polynomial for $f(x) = e^x$.
	\item	Use $p_5(x)$ to approximate the value of $e$.
\end{enumerate}}\vspace{0pt}
\solution
\mtable{The derivatives of $f(x)=e^x$ evaluated at $x=0$.}{fig:taypoly1a}{%
\begin{align*}
% todo Tim latexml https://github.com/brucemiller/LaTeXML/issues/1763
f(x)&=e^x \quad\Rightarrow & f(0)&=1 \\
\fp(x)&=e^x \quad\Rightarrow & \fp(0)&=1 \\
\fpp(x)&=e^x \quad\Rightarrow & \fpp(0)&=1 \\
\vdots && \vdots & \\
f\,^{(n)}(x)&=e^x \quad\Rightarrow & f\,^{(n)}(0)&=1
%f(x)&=e^x & \Rightarrow && f(0)&=1 \\
%\fp(x)&=e^x & \Rightarrow && \fp(0)&=1 \\
%\fpp(x)&=e^x & \Rightarrow && \fpp(0)&=1 \\
%\vdots &&&& \vdots & \\
%f\,^{(n)}(x)&=e^x & \Rightarrow && f\,^{(n)}(0)&=1
\end{align*}}
\begin{enumerate}
\item We start with creating a table of the derivatives of $e^x$ evaluated at $x=0$. In this particular case, this is relatively simple, as shown in \autoref{fig:taypoly1a}. By the definition of the Maclaurin polynomial, we have 
\[
	p_n(x)
	=\sum_{k=0}^n\frac{f\,^{(k)}(0)}{k!}x^k
	=\sum_{k=0}^n\frac1{k!}x^k.
\]

\item	Using our answer from part 1, we have
\[p_5(x) = 1+x+\frac{1}{2}x^2+\frac{1}{6}x^3 + \frac{1}{24}x^4 + \frac{1}{120}x^5.\]
To approximate the value of $e$, note that $e = e^1 = f(1) \approx p_5(1).$ It is very straightforward to evaluate $p_5(1)$:
%
\mtable{A plot of $f(x)=e^x$ and its 5\textsuperscript{th} degree Maclaurin polynomial $p_5(x)$.}{fig:taypoly1b}{\begin{tikzpicture}[alt={Blue e^x curve with red 5th-degree Maclaurin p_s; the two align near x=0 then red trails below for x>2.}]
\begin{axis}[width=1.16\marginparwidth,tick label style={font=\scriptsize},
axis y line=middle,axis x line=middle,name=myplot,axis on top,
ymin=-3,ymax=11,xmin=-3.75,xmax=2.9]
\addplot [draw={\colorone},domain=-3.5:2.5,smooth,thick,samples=50] {exp(x)};
\addplot [draw={\colortwo},domain=-4:4,smooth,thick] {1+x+x^2/2+x^3/6+x^4/24+x^5/120};
\draw (axis cs:-3.,-2) node {\scriptsize $y=p_5(x)$};
\end{axis}
\node [right] at (myplot.right of origin) {\scriptsize $x$};
\node [above] at (myplot.above origin) {\scriptsize $y$};
\end{tikzpicture}}
%
\[p_5(1) = 1+1+\frac12+\frac16+\frac1{24}+\frac1{120} = \frac{163}{60} \approx 2.71667.\]
This is an error of about $0.0016$, or $0.06\%$ of the true value.

A plot of $f(x)=e^x$ and $p_5(x)$ is given in \autoref{fig:taypoly1b}.
\end{enumerate}
\end{example}

\begin{example}[Finding and using Taylor polynomials]\label{ex_taypoly2}%
\mbox{}\\[-2\baselineskip]\parbox[t]{\linewidth}{\begin{enumerate}
	\item	Find the $n^{\text{th}}$ Taylor polynomial of $y=\ln x$ at $x=1$.
	\item	Use $p_6(x)$ to approximate the value of $\ln 1.5$.
	\item	Use $p_6(x)$ to approximate the value of $\ln 2$. 
\end{enumerate}}\vspace{0pt}
\solution
\begin{enumerate}
\item	We begin by creating a table of derivatives of $\ln x$ evaluated at $x=1$. While this is not as straightforward as it was in the previous example, a pattern does emerge, as shown in \autoref{fig:taypoly2a}.
\mtable{Derivatives of $\ln x$ evaluated at $x=1$.}{fig:taypoly2a}{%
\begin{align*}
% todo Tim latexml https://github.com/brucemiller/LaTeXML/issues/1763
f(x) &= \makebox[3em]{$\ln x$} \quad \Rightarrow & f(1)&=\phantom-0\\
\fp(x) &= \makebox[3em]{$\phantom-1/x$} \quad \Rightarrow & \fp(1) &= \phantom-1\\
\fp'(x) &= \makebox[3em]{$-1/x^2$} \quad \Rightarrow & \fp'(1) &= -1\\
\fp''(x) &= \makebox[3em]{$\phantom-2/x^3$} \quad \Rightarrow & \fp''(1) &= \phantom-2\\
f\,^{(4)}(x) &= \makebox[3em]{$-6/x^4$} \quad \Rightarrow & f\,^{(4)}(1) &= -6\\
%f\,^{(5)}(x) &= \makebox[3em]{$24/x^5$} \quad\Rightarrow  &f\,^{(5)}(1) &= 24\\
\vdots && \vdots & \\
f\,^{(n)}(x) &= \makebox[3em]{}\quad \Rightarrow & f\,^{(n)}(1) &= \smallskip\\
\makebox[1pt]{$\dfrac{(-1)^{n+1}(n-1)!}{x^n}$} && \makebox[1pt]{$(-1)^{n+1}(n-1)!$} &
%f(x) &= \ln x & \Rightarrow && f(1)&=\phantom-0\\
%\fp(x) &= \phantom-1/x & \Rightarrow && \fp(1) &= \phantom-1\\
%\fp'(x) &= -1/x^2 & \Rightarrow && \fp'(1) &= -1\\
%\fp''(x) &= \phantom-2/x^3 & \Rightarrow && \fp''(1) &= \phantom-2\\
%f\,^{(4)}(x) &= -6/x^4 & \Rightarrow && f\,^{(4)}(1) &= -6\\
%\vdots &&&& \vdots & \\
%f\,^{(n)}(x) &= & \Rightarrow && f\,^{(n)}(1) &= \smallskip\\
%\makebox[1pt]{$\dfrac{(-1)^{n+1}(n-1)!}{x^n}$} &&&& \makebox[1pt]{$(-1)^{n+1}(n-1)!$} &
\end{align*}}

Using \autoref{def:taypoly}, we have
\[
	p_n(x)
	= \sum_{k=0}^n\frac{f\,^{(k)}(c)}{k!}(x-c)^k
	= \sum_{k=1}^n\frac{(-1)^{k+1}}k(x-1)^k.
\]

\item	We can compute $p_6(x)$ using our work above:
\[
p_6(x)
 = (x-1)-\frac12(x-1)^2+\frac13(x-1)^3-\frac14(x-1)^4+\frac15(x-1)^5-\frac16(x-1)^6.
\]
Since $p_6(x)$ approximates $\ln x$ well near $x=1$, we approximate $\ln 1.5 \approx p_6(1.5)$:
\begin{align*}
	p_6(1.5)
	&= (1.5-1)-\frac12(1.5-1)^2+\frac13(1.5-1)^3 \\
	&\qquad\qquad{}-\frac14(1.5-1)^4+\frac15(1.5-1)^5-\frac16(1.5-1)^6\\
	&=\frac{259}{640}\\
	&\approx 0.404688.
\end{align*}
%
\mtable{A plot of $y=\ln x$ and its 6\textsuperscript{th} degree Taylor polynomial at $x=1$.}{fig:taypoly2b}{\begin{tikzpicture}[alt={Blue ln x and red 6th-degree Taylor about x=1; close near x=1, red rises above blue by x=3.}]
\begin{axis}[width=1.16\marginparwidth,tick label style={font=\scriptsize},
axis y line=middle,axis x line=middle,name=myplot,axis on top,
ymin=-4.5,ymax=2.4,xmin=-.5,xmax=3.2]
\addplot [draw={\colorone},domain=0.01:3,smooth,thick,samples=50] {ln(x)};
\addplot [draw={\colortwo},domain=-.5:3.2,smooth,thick]
 {(x-1)-(x-1)^2/2+(x-1)^3/3-(x-1)^4/4+(x-1)^5/5-(x-1)^6/6};
\draw (axis cs:2.5,1.65) node {\scriptsize $y=\ln x$};
\draw (axis cs:2.2,-2.75) node {\scriptsize $y=p_{6}(x)$};
\end{axis}
\node [right] at (myplot.right of origin) {\scriptsize $x$};
\node [above] at (myplot.above origin) {\scriptsize $y$};
\end{tikzpicture}}
%
This is a good approximation as a calculator shows that $\ln 1.5 \approx 0.4055.$  %This is an error of $0.0008$, or $0.2\%$. % this is in next example
\autoref{fig:taypoly2b} plots $y=\ln x$ with $y=p_6(x)$. We can see that $\ln 1.5\approx p_6(1.5)$.

\item	
We approximate $\ln 2$ with $ p_6(2)$:
\begin{align*}
p_6(2) &= (2-1)-\frac12(2-1)^2+\frac13(2-1)^3 \\
			&\qquad\qquad{}-\frac14(2-1)^4+\frac15(2-1)^5-\frac16(2-1)^6\\
			&=	1-\frac12+\frac13-\frac14+\frac15-\frac16 \\
			&= \frac{37}{60}\\ 
			&\approx 0.616667.
\end{align*}
This approximation is not terribly impressive: a hand held calculator shows that $\ln 2 \approx 0.693147.$
% This is an error of $0.08$, or $11\%$. % this is in next example
The graph in \autoref{fig:taypoly2b} shows that $p_6(x)$ provides less accurate approximations of $\ln x$ as $x$ gets close to 0 or 2. 

% todo use 20th degree polynomial instead of coordinates in Figure 9.9.8 fig:taypoly2c
\mtable{A plot of $y=\ln x$ and its 20\textsuperscript{th} degree Taylor polynomial at $x=1$.}{fig:taypoly2c}{\begin{tikzpicture}[alt={Blue ln x with red 20th-degree Taylor about x=1; ln x tracked well to about x=2.3 before diverging upward.}]
\begin{axis}[width=1.16\marginparwidth,tick label style={font=\scriptsize},
axis y line=middle,axis x line=middle,name=myplot,axis on top,
ymin=-4.5,ymax=2.4,xmin=-.5,xmax=3.2]
\addplot [draw={\colorone},domain=0.01:3,smooth,thick,samples=50] {ln(x)};
\addplot [draw={\colortwo},smooth,thick] coordinates {
 (-0.108,-7.567)(-0.052,-4.958)(0.004,-3.519)(0.06,-2.671)(0.116,-2.13)
 (0.172,-1.756)(0.228,-1.478)(0.284,-1.259)(0.34,-1.079)(0.396,-0.9263)
 (0.452,-0.7941)(0.508,-0.6773)(0.564,-0.5727)(0.62,-0.478)(0.676,-0.3916)
 (0.732,-0.312)(0.788,-0.2383)(0.844,-0.1696)(0.9,-0.1054)(0.956,-0.045)
 (1.012,0.01193)(1.068,0.06579)(1.124,0.1169)(1.18,0.1655)(1.236,0.2119)
 (1.292,0.2562)(1.348,0.2986)(1.404,0.3393)(1.46,0.3784)(1.516,0.4161)
 (1.572,0.4523)(1.628,0.4874)(1.684,0.5212)(1.74,0.5538)(1.796,0.5853)
 (1.852,0.6154)(1.908,0.6427)(1.964,0.6635)(2.02,0.6665)(2.076,0.621)
 (2.132,0.4476)(2.188,-0.04877)(2.244,-1.326)(2.3,-4.421)};
% the 20th degree Taylor approximation
\draw (axis cs:2.5,1.65) node {\scriptsize $y=\ln x$};
\draw (axis cs:1.7,-2.75) node {\scriptsize $y=p_{20}(x)$};
\end{axis}
\node [right] at (myplot.right of origin) {\scriptsize $x$};
\node [above] at (myplot.above origin) {\scriptsize $y$};
\end{tikzpicture}}

Surprisingly enough, even the 20\textsuperscript{th} degree Taylor polynomial fails to approximate $\ln x$ for $x>2$, as shown in \autoref{fig:taypoly2c}. We'll soon discuss why this is.
\end{enumerate}
\end{example}

Taylor polynomials are used to approximate functions $f(x)$ in mainly two situations:
\begin{enumerate}
	\item	When $f(x)$ is known, but perhaps ``hard'' to compute directly. For instance, we can define $y=\cos x$ as either the ratio of sides of a right triangle (``adjacent over hypotenuse'') or with the unit circle. However, neither of these provides a convenient way of computing $\cos 2$. A Taylor polynomial of sufficiently high degree can provide a reasonable method of computing such values using only operations usually hard-wired into a computer ($+$, $-$, $\times$ and $\div$).
	
	\item	When $f(x)$ is not known, but information about its derivatives is known. This occurs more often than one might think, especially in the study of differential equations.
\end{enumerate}

\mnote{\textbf{Note:} Even though Taylor polynomials \emph{could} be used in calculators and computers to calculate values of trigonometric functions, in practice they generally aren't. Other more efficient and accurate methods have been developed, such as the CORDIC algorithm.}
	
In both situations, a critical piece of information to have is ``How good is my approximation?'' If we use a Taylor polynomial to compute $\cos 2$, how do we know how accurate the approximation is? 

We had the same problem with Numerical Integration. \autoref{thm:numerical_error} provided bounds on the error when using, say, Simpson's Rule to approximate a definite integral. These bounds allowed us to determine that, for example, using $10$ subintervals provided an approximation within $\pm .01$ of the exact value. The following theorem gives similar bounds for Taylor (and hence Maclaurin) polynomials.

{\tcbset{grow to right by=.5em}
\begin{theorem}[Taylor's Theorem]\label{thm:taylorthm}%
\mbox{}\\[-2\baselineskip]\index{Taylor Polynomial!Taylor's Theorem}\index{Taylor's Theorem}
\begin{enumerate}
	\item	Let $f$ be a function whose $(n+1)^{\text{th}}$ derivative exists on an open interval $I$ and let $c$ be in $I$. Then, for each $x$ in $I$, there exists $z_x$ between $x$ and $c$ such that\vspace{-.3\baselineskip}
\[R_n(x) = f(x) - \sum_{k=0}^n\frac{f\,^{(k)}(c)}{k!}(x-c)^k = \frac{f\,^{(n+1)}(z_x)}{(n+1)!}(x-c)^{n+1}.\]
%\[f(x) = \sum_{k=0}^n\frac{f\,^{(k)}(c)}{k!}(x-c)^k+R_n(x),\]
%where $\ds R_n(x) = \frac{f\,^{(n+1)}(z_x)}{(n+1)!}(x-c)^{n+1}.$
	\item	$\abs{R_n(x)}\leq\dfrac{\max_z\abs{\,f\,^{(n+1)}(z)}}{(n+1)!}\abs{x-c}^{n+1}$, where $z$ is between $x$ and $c$.
\end{enumerate}
\end{theorem}}

The first part of Taylor's Theorem states that $f(x) = p_n(x) + R_n(x)$, where $p_n(x)$ is the $n^{\text{th}}$ order Taylor polynomial and $R_n(x)$ is the remainder, or error, in the Taylor approximation. The second part gives bounds on how big that error can be. If the $(n+1)^{\text{th}}$ derivative is large, the error may be large; if $x$ is far from $c$, the error may also be large. However, the $(n+1)!$ term in the denominator tends to ensure that the error gets smaller as $n$ increases.

The following example computes error estimates for the approximations of $\ln 1.5$ and $\ln 2$ made in \autoref{ex_taypoly2}.

\begin{example}[Finding error bounds of a Taylor polynomial]\label{ex_taypoly3}%
Use \autoref{thm:taylorthm} to find error bounds when approximating $\ln 1.5$ and $\ln 2$ with $p_6(x)$, the Taylor polynomial of degree 6 of $f(x)=\ln x$ at $x=1$, as calculated in \autoref{ex_taypoly2}.
\solution
\begin{enumerate}
\item	We start with the approximation of $\ln 1.5$ with $p_6(1.5)$.
% Taylor's Theorem references an open interval $I$ that contains both $x$ and $c$. The smaller the interval we use the better; it will give us a more accurate (and smaller) approximation of the error. We let $I = (0.9,1.6)$, as this interval contains both $c=1$ and $x=1.5$. 
%
Taylor's Theorem references $\max\abs{f\,^{(n+1)}(z)}$. In our situation, this is asking ``How big can the 7\textsuperscript{th} derivative of $y=\ln x$ be on the interval $[1,1.5]$?'' The seventh derivative is $y = 6!/x^7$. The largest absolute value it attains on $I$ is 720. Thus we can bound the error as:
\begin{align*}
	\abs{R_6(1.5)}
	&\leq \frac{\max\abs{f\,^{(7)}(z)}}{7!}\abs{1.5-1}^7\\
	&\leq \frac{720}{5040}\cdot\frac1{2^7}\\
	&\approx 0.001.
\end{align*}
We computed $p_6(1.5) = 0.404688$; using a calculator, we find $\ln 1.5 \approx 0.405465$, so the actual error is about $0.000778$ (or $0.2\%$), which is less than our bound of $0.001$. This affirms Taylor's Theorem; the theorem states that our approximation would be within about one thousandth of the actual value, whereas the approximation was actually closer.

	\item	%We again find an interval $I$ that contains both $c=1$ and $x=2$; we choose $I = (0.9,2.1)$.
	The maximum value of the seventh derivative of $f$ on $[1,2]$ %this interval
	is again 720 (as the largest values come at $x=1$). Thus 
\begin{align*}
	\abs{R_6(2)}
	&\leq \frac{\max\abs{f\,^{(7)}(z)}}{7!}\abs{2-1}^7\\
	&\leq \frac{720}{5040}\cdot1^7\\
	&\approx0.15.
\end{align*}
This bound is not as nearly as good as before. Using the degree 6 Taylor polynomial at $x =1$ will bring us within 0.15 of the correct answer. As $p_6(2)\approx 0.61667$, our error estimate guarantees that the actual value of $\ln 2$ is somewhere between $0.46$ and $0.76$. These bounds are not particularly useful.

In reality, our approximation was only off by about $0.07$ (or $11\%$). However, we are approximating ostensibly because we do not know the real answer. In order to be assured that we have a good approximation, we would have to resort to using a polynomial of higher degree.
\end{enumerate}
\end{example}

We practice again. This time, we use Taylor's theorem to find $n$ that guarantees our approximation is within a certain amount.

\clearpage

\mtable[yshift=-0.3\pageheight]{A table of the derivatives of $f(x)=\cos x$ evaluated at $x=0$.}{fig:taypoly4a}{\begin{align*}
f(x) &= \phantom-\cos x & \Rightarrow && f(0) &= \phantom-1\\
\fp(x) &= -\sin x & \Rightarrow && \fp(0) &= \phantom-0\\
\fp'(x) &= -\cos x & \Rightarrow && \fp'(0) &= -1\\
\fp''(x) &= \phantom-\sin x & \Rightarrow && \fp''(0) &= \phantom-0\\
f\,^{(4)}(x) &= \phantom-\cos x & \Rightarrow && f\,^{(4)}(0) &= \phantom-1\\
f\,^{(5)}(x) &= -\sin x & \Rightarrow && f\,^{(5)}(0) &= \phantom-0\\
f\,^{(6)}(x) &= -\cos x & \Rightarrow && f\,^{(6)}(0) &= -1\\
f\,^{(7)}(x) &= \phantom-\sin x & \Rightarrow && f\,^{(7)}(0) &= \phantom-0\\
f\,^{(8)}(x) &= \phantom-\cos x & \Rightarrow && f\,^{(8)}(0) &= \phantom-1\\
f\,^{(9)}(x) &= -\sin x & \Rightarrow && f\,^{(9)}(0) &= \phantom-0
\end{align*}}%
%
\mtable{A graph of $f(x)= \cos x$ and its degree 8 Maclaurin polynomial.}{fig:taypoly4b}{\begin{tikzpicture}[alt={Blue cos x; red 6-term Maclaurin matches near 0, then alternates above/below the curve out to ±π.}]
\begin{axis}[width=1.16\marginparwidth,tick label style={font=\scriptsize},
axis y line=middle,axis x line=middle,name=myplot,axis on top,
xtick={-5,-4,-3,-2,-1,1,2,3,4,5},ymin=-1.1,ymax=1.5,xmin=-5.5,xmax=5.5,legend style={at={(.5,0)},anchor=north,draw=none}]
\addplot [draw={\colorone},domain=-5:5,smooth,thick,samples=50] {cos(deg(x))};
\legend{\scriptsize$f(x)=\cos x$}
\addplot [draw={\colortwo},domain=-4:4,smooth,thick] {1-x^2/2+x^4/24-x^6/720+x^8/40320};
\draw (axis cs:-3.,1) node {\scriptsize $y=p_8(x)$};
\end{axis}
\node [right] at (myplot.right of origin) {\scriptsize $x$};
\node [above] at (myplot.above origin) {\scriptsize $y$};
\end{tikzpicture}}%

\vspace{-\topsep}\vspace{-\partopsep}%

\begin{example}[Finding sufficiently accurate Taylor polynomials]\label{ex_taypoly4}%
Find $n$ such that the $n^{\text{th}}$ Taylor polynomial of $f(x)=\cos x$ at $x=0$ approximates $\cos 2$ to within $0.001$ of the actual answer. What is $p_n(2)$?
\solution
Following Taylor's theorem, we need bounds on the size of the derivatives of $f(x)=\cos x$. In the case of this trigonometric function, this is easy. All derivatives of cosine are $\pm \sin x$ or $\pm \cos x$. In all cases, these functions are never greater than 1 in absolute value. We want the error to be less than $0.001$. To find the appropriate $n$, consider the following inequalities:
\begin{align*}
\frac{\max\abs{f\,^{(n+1)}(z)}}{(n+1)!}\abs{2-0}^{n+1} &\leq 0.001 \\
\frac1{(n+1)!}\cdot2^{n+1} &\leq 0.001
\end{align*}
We find an $n$ that satisfies this last inequality with trial-and-error. When $n=8$, we have $\ds \frac{2^{8+1}}{(8+1)!} \approx 0.0014$; when $n=9$, we have $\ds \frac{2^{9+1}}{(9+1)!} \approx 0.000282 <0.001$. Thus we want to approximate $\cos 2$ with $p_9(2)$.\bigskip

We now set out to compute $p_9(x)$. We again need a table of the derivatives of $f(x)=\cos x$ evaluated at $x=0$. A table of these values is given in \autoref{fig:taypoly4a}. Notice how the derivatives, evaluated at $x=0$, follow a certain pattern. All the odd powers of $x$ in the Taylor polynomial will disappear as their coefficient is 0. While our error bounds state that we need $p_9(x)$, our work shows that this will be the same as $p_8(x)$.

Since we are forming our polynomial at $x=0$, we are creating a Maclaurin polynomial, and:\vspace{-.3\baselineskip}
\[
	p_8(x) = \sum_{k=0}^8\frac{f\,^{(k)}(0)}{k!}x^k
	=  1-\frac{1}{2!}x^2+\frac{1}{4!}x^4-\frac{1}{6!}x^6+\frac{1}{8!}x^8
\]

We finally approximate $\cos 2$:
\[\cos 2 \approx p_8(2) = -\frac{131}{315} \approx -0.41587.\]
Our error bound guarantees that this approximation is within $0.001$ of the correct answer. Technology shows us that our approximation is actually within about $0.0003$ (or $0.07\%$) of the correct answer.

\autoref{fig:taypoly4b} shows a graph of $y=p_8(x)$ and $y=\cos x$. Note how well the two functions agree on about $(-\pi,\pi)$.
\end{example}

\begin{example}[Finding and using Taylor polynomials]\label{ex_taypoly5}%
\mbox{}\\[-2\baselineskip]\parbox[t]{\linewidth}{\begin{enumerate}
	\item	Find the degree 4 Taylor polynomial, $p_4(x)$, for $f(x)=\sqrt{x}$ at $x=4.$
	\item	Use $p_4(x)$ to approximate $\sqrt{3}$.
	\item	Find bounds on the error when approximating $\sqrt{3}$ with $p_4(3)$.
\end{enumerate}}\vspace{0pt}
\solution
\begin{enumerate}
	\item	We begin by evaluating the derivatives of $f$ at $x=4$. This is done in \autoref{fig:taypoly5a}. These values allow us to form the Taylor polynomial $p_4(x)$:
	%
\mtable{A table of the derivatives of $f(x)=\sqrt{x}$ evaluated at $x=4$.}{fig:taypoly5a}{%
\begin{align*}
f(x) &= \sqrt{x} & \Rightarrow && f(4) &= 2\\
\ds\fp(x) &= \frac{1}{2\sqrt{x}} & \Rightarrow && \ds\fp(4) &= \frac{1}{4}\\
\ds\fp'(x) &= \frac{-1}{4x^{3/2}} & \Rightarrow && \ds\fp'(4) &= \frac{-1}{32}\\
\ds\fp''(x) &= \frac3{8x^{5/2}} & \Rightarrow && \ds\fp''(4) &= \frac{3}{256}\\
\ds f\,^{(4)}(x) &= \frac{-15}{16x^{7/2}} & \Rightarrow
&& \ds f\,^{(4)}(4) &= \frac{-15}{2048}\vspace{-.5\baselineskip}
\end{align*}}
%
\begin{multline*}
p_4(x) = \\
2 + \frac14(x-4) +\frac{-1/32}{2!}(x-4)^2+\frac{3/256}{3!}(x-4)^3+\frac{-15/2048}{4!}(x-4)^4.
\end{multline*}

	\item	As $p_4(x) \approx \sqrt{x}$ near $x=4$, we approximate $\sqrt{3}$ with $p_4(3) = 1.73212$.

	\item	%To find a bound on the error, we need an open interval that contains $x=3$ and $x=4$. We set $I = (2.9,4.1)$.
	The largest value the fifth derivative of $f(x)=\sqrt{x}$ takes on $[3,4]$ is when $x=3$, at about $0.0234$. Thus
\[\abs{R_4(3)} \leq \frac{0.0234}{5!}\abs{3-4}^5 \approx 0.00019.\]
%
\mtable{A graph of $f(x)=\sqrt{x}$ and its degree 4 Taylor polynomial at $x=4$.}{fig:taypoly5b}{\begin{tikzpicture}[alt={Blue √x (x > 0); red 4-term Taylor at 4 tracks up to = 6, then climbs more steeply.}]
\begin{axis}[width=1.16\marginparwidth,tick label style={font=\scriptsize},
axis y line=middle,axis x line=middle,name=myplot,axis on top,
ymin=-.1,ymax=3.5,xmin=-1,xmax=11,legend style={at={(1,0.5)}}]
\addplot [draw={\colorone},domain=0:10,smooth,thick,samples=50] {sqrt(x)};
\addlegendentry{\scriptsize$y=\sqrt x$}
\addplot [draw={\colortwo},smooth,thick,domain=0:10] {2+(x-4)/4-(x-4)^2/64+3*(x-4)^3/1536-15*(x-4)^4/49152};
\addlegendentry{\scriptsize$y=p_4(x)$}
\end{axis}
\node [right] at (myplot.right of origin) {\scriptsize $x$};
\node [above] at (myplot.above origin) {\scriptsize $y$};
\end{tikzpicture}}%
%
This shows our approximation is accurate to at least the first 2 places after the decimal. It turns out that our approximation has an error of $0.00007$, or $0.004\%$. A graph of $f(x)=\sqrt x$ and $p_4(x)$ is given in \autoref{fig:taypoly5b}. Note how the two functions are nearly indistinguishable on $(2,7)$.
\end{enumerate}
\end{example}

%Our final example gives a brief introduction to using Taylor polynomials to solve differential equations.
%
%\begin{example}[Approximating an unknown function]\label{ex_taypoly6}%
%A function $y=f(x)$ is unknown save for the following two facts.
%\begin{enumerate}
%\item		$y(0) = f(0) = 1$, and
%\item		$y\primeskip'= y^2$
%\end{enumerate}
%(This second fact says that amazingly, the derivative of the function is actually the function squared!)
%
%Find the degree 3 Maclaurin polynomial $p_3(x)$ of $y=f(x)$.
%\solution
%One might initially think that not enough information is given to find $p_3(x)$. However, note how the second fact above actually lets us know what $y\primeskip'(0)$ is:
%\[y\primeskip' = y^2 \Rightarrow y\primeskip'(0) = y^2(0).\]
%Since $y(0) = 1$, we conclude that $y\primeskip'(0) = 1$.
%
%Now we find information about $y\primeskip''$. Starting with $y\primeskip'=y^2$, take derivatives of both sides, \emph{with respect to $x$}. That means we must use implicit differentiation.
%\begin{align*}
%y\primeskip' &= y^2\\
%\frac{\dd}{\dd x}\bigl(y\primeskip'\bigr) &= \frac{\dd}{\dd x}\bigl(y^2\bigr)\\
%y\primeskip'' &= 2y\cdot y\primeskip'.
%\intertext{Now evaluate both sides at \intertextmath{x=0}:}
%y\primeskip''(0) &= 2y(0)\cdot y\primeskip'(0)\\
%y\primeskip''(0) &= 2
%\end{align*}
%We repeat this once more to find $y\primeskip'''(0)$. We again use implicit differentiation; this time the Product Rule is also required.
%\begin{align*}
%\frac{\dd}{\dd x}\bigl(y\primeskip''\bigr) &= \frac{\dd}{\dd x} \bigl(2yy\primeskip'\bigr)\\
%y\primeskip''' &= 2y\primeskip'\cdot y\primeskip' + 2y\cdot y\primeskip''.
%\intertext{Now evaluate both sides at \intertextmath{x=0}:}
%y\primeskip'''(0) &= 2y\primeskip'(0)^2 + 2y(0)y\primeskip''(0)\\
%y\primeskip'''(0) &=	2+4=6
%\end{align*}
%In summary, we have:
%\[y(0) = 1 \qquad y\primeskip'(0) = 1  \qquad y\primeskip''(0) = 2 \qquad y\primeskip'''(0) = 6.\]
%We can now form $p_3(x)$:
%\begin{align*}
%p_3(x) &= 1 + x + \frac{2}{2!}x^2 + \frac{6}{3!}x^3\\
%				&= 1+x+x^2+x^3.
%\end{align*}
%\mtable{A graph of $y=-1/(x-1)$ and $y=p_3(x)$ from \autoref{ex_taypoly6}.}{fig:taypoly6}{\begin{tikzpicture}[alt={Blue curve y=-1/(x-1) starts at (-1,0.5), passes through (0,1), and approaches upward its vertical asymptote x=1.  Red curve y=p_3(x) starts near (-1,0), traces the blue curve for -0.5<x<0.5 and ends not as steep near (0.8,3).}]
%\begin{axis}[width=1.16\marginparwidth,tick label style={font=\scriptsize},
%axis y line=middle,axis x line=middle,name=myplot,axis on top,
%ymin=-.1,ymax=3.5,xmin=-1.1,xmax=1.1,legend style={at={(1.1,0.1)},anchor=south east}]
%\addplot [draw={\colorone},domain=-1:.7,smooth,thick,samples=50] {-1/(x-1)};
%\addlegendentry{\scriptsize $\displaystyle y= \frac{1}{1-x}$}
%\addplot [draw={\colortwo},smooth,thick] coordinates {(-1.,0)(-0.96,0.07686)(-0.92,0.1477)(-0.88,0.2129)(-0.84,0.2729)(-0.8,
%0.328)(-0.76,0.3786)(-0.72,0.4252)(-0.68,0.468)(-0.64,0.5075)(-0.6,0.
%544)(-0.56,0.578)(-0.52,0.6098)(-0.48,0.6398)(-0.44,0.6684)(-0.4,0.
%696)(-0.36,0.7229)(-0.32,0.7496)(-0.28,0.7764)(-0.24,0.8038)(-0.2,0.
%832)(-0.16,0.8615)(-0.12,0.8927)(-0.08,0.9259)(-0.04,0.9615)(0,1.)(0.
%04,1.042)(0.08,1.087)(0.12,1.136)(0.16,1.19)(0.2,1.248)(0.24,1.311)(0.
%28,1.38)(0.32,1.455)(0.36,1.536)(0.4,1.624)(0.44,1.719)(0.48,1.821)(0.
%52,1.931)(0.56,2.049)(0.6,2.176)(0.64,2.312)(0.68,2.457)(0.72,2.612)(
%0.76,2.777)(0.8,2.952)(0.84,3.138)(0.88,3.336)(0.92,3.545)(0.96,3.766)(1.,4.)};
%\addlegendentry{\scriptsize $y= p_3(x)$}
%\end{axis}
%\node [right] at (myplot.right of origin) {\scriptsize $x$};
%\node [above] at (myplot.above origin) {\scriptsize $y$};
%\end{tikzpicture}}
%It turns out that the differential equation we started with, $y\primeskip'=y^2$, where $y(0)=1$, can be solved without too much difficulty: $\ds y = \frac{1}{1-x}$. \autoref{fig:taypoly6} shows this function plotted with $p_3(x)$. Note how similar they are near $x=0$.
%\end{example}
%
%It is beyond the scope of this text to pursue error analysis when using Taylor polynomials to approximate solutions to differential equations. This topic is often broached in introductory Differential Equations courses and usually covered in depth in Numerical Analysis courses. Such an analysis is very important; one needs to know how good their approximation is. We explored this example simply to demonstrate the usefulness of Taylor polynomials.

\bigskip

Most of this chapter has been devoted to the study of infinite series. This section has stepped aside from this study, focusing instead on finite summation of terms. In the next section, we will combine power series and Taylor polynomials into \textbf{Taylor Series}, where we represent a function with an infinite series.

\printexercises{exercises/08-07-exercises}

% todo for exercises 21-24, reach back to lin approx / Newton’s method and redo the appropriate problems with Taylor series

\section{Taylor Series}\label{sec:taylor_series}

In \autoref{sec:power_series}, we showed how certain functions can be represented by a power series function. In \autoref{sec:taylor_poly}, we showed how we can approximate functions with polynomials, given that enough derivative information is available. In this section we combine these concepts: if a function $f(x)$ is infinitely differentiable, we show how to represent it with a power series function.

\begin{definition}[Taylor and Maclaurin Series]\label{def:taylor_series}%
Let $f(x)$ have derivatives of all orders at $x=c$.
\index{Taylor Series!definition}\index{Maclaurin Series!definition}\index{Maclaurin Series|seealso{Taylor Series}}\index{series!Taylor}\index{series!Maclaurin}
\begin{enumerate}
	\item The \textbf{Taylor Series of $f(x)$, centered at $c$} is
	\[\sum_{n=0}^\infty \frac{f\,^{(n)}(c)}{n!}(x-c)^n.\]
	\item	Setting $c=0$ gives the \textbf{Maclaurin Series of $f(x)$}:
	\[\sum_{n=0}^\infty \frac{f\,^{(n)}(0)}{n!}x^n.\]
\end{enumerate}
\end{definition}

\youtubeVideo{Os8OtXFBLkY}{Taylor and Maclaurin Series --- Example 2}

The difference between a Taylor polynomial and a Taylor series is the former is a polynomial, containing only a finite number of terms, whereas the latter is a series, a summation of an infinite set of terms. When creating the Taylor polynomial of degree $n$ for a function $f(x)$ at $x=c$, we needed to evaluate $f$, and the first $n$ derivatives of $f$, at $x=c$. When creating the Taylor series of $f$, we need to find a pattern that describes the $n^{\text{th}}$ derivative of $f$ at $x=c$. We demonstrate this in the next two examples.

\begin{example}[The Maclaurin series of $f(x) = \cos x$]\label{ex_ts1}%
Find the Maclaurin series of $f(x)=\cos x$.
\solution
In \autoref{ex_taypoly4} we found the $8^{\text{th}}$ degree Maclaurin polynomial of $\cos x$. In doing so, we created the table shown in \autoref{fig:ts1}.
%
\mtable{A table of the derivatives of $f(x)=\cos x$ evaluated at $x=0$.}{fig:ts1}{%
\begin{align*}
f(x) &= \phantom-\cos x &\Rightarrow & &f(0) &= \phantom-1\\
\fp(x) &= -\sin x &\Rightarrow & &\fp(0) &= \phantom-0\\
\fp'(x) &= -\cos x &\Rightarrow & &\fp'(0) &= -1\\
\fp''(x) &= \phantom-\sin x &\Rightarrow & &\fp''(0) &= \phantom-0\\
f\,^{(4)}(x) &= \phantom-\cos x &\Rightarrow & &f\,^{(4)}(0) &= \phantom-1\\
f\,^{(5)}(x) &= -\sin x &\Rightarrow & &f\,^{(5)}(0) &= \phantom-0\\
f\,^{(6)}(x) &= -\cos x &\Rightarrow & &f\,^{(6)}(0) &= -1\\
f\,^{(7)}(x) &= \phantom-\sin x &\Rightarrow & &f\,^{(7)}(0) &= \phantom-0\\
f\,^{(8)}(x) &= \phantom-\cos x &\Rightarrow & &f\,^{(8)}(0) &= \phantom-1\\
f\,^{(9)}(x) &= -\sin x &\Rightarrow & &f\,^{(9)}(0) &= \phantom-0
\end{align*}}
%
Notice how $f\,^{(n)}(0)=0$ when $n$ is odd,  $f\,^{(n)}(0)=1$ when $n$ is divisible by $4$, and $f\,^{(n)}(0)=-1$ when $n$ is even but not divisible by 4. Thus the Maclaurin series of $\cos x$ is
\[1-\frac{x^2}2+\frac{x^4}{4!}-\frac{x^6}{6!}+\frac{x^8}{8!} - \dotsb\]
We can go further and write this as a summation. Since we only need the terms where the power of $x$ is even, we write the power series in terms of $x^{2n}$:
\[\sum_{n=0}^\infty (-1)^{n}\frac{x^{2n}}{(2n)!}.\]
\end{example}

\begin{example}[The Taylor series of $f(x)=\ln x$ at $x=1$]\label{ex_ts2}%
Find the Taylor series of $f(x) = \ln x$ centered at $x=1$.
\solution
\autoref{fig:ts2} shows the $n^{\text{th}}$ derivative of $\ln x$ evaluated at $x=1$ for $n=0,\dotsc,5$, along with an expression for the $n^{\text{th}}$ term:
\[f\,^{(n)}(1) = (-1)^{n+1}(n-1)!\quad \text{for $n\geq 1$.}\]
Remember that this is what distinguishes Taylor series from Taylor polynomials; we are very interested in finding a pattern for the $n^{\text{th}}$ term, not just finding a finite set of coefficients for a polynomial.
\mtable{Derivatives of $\ln x$ evaluated at $x=1$.}{fig:ts2}{%
\begin{align*}
% todo Tim latexml https://github.com/brucemiller/LaTeXML/issues/1763
f(x) &= \makebox[3em]{$\ln x$} \quad \Rightarrow & f(1)&=\phantom-0\\
\fp(x) &= \makebox[3em]{$\phantom-1/x$} \quad \Rightarrow & \fp(1) &= \phantom-1\\
\fp'(x) &= \makebox[3em]{$-1/x^2$} \quad \Rightarrow & \fp'(1) &= -1\\
\fp''(x) &= \makebox[3em]{$\phantom-2/x^3$} \quad \Rightarrow & \fp''(1) &= \phantom-2\\
f\,^{(4)}(x) &= \makebox[3em]{$-6/x^4$} \quad \Rightarrow & f\,^{(4)}(1) &= -6\\
f\,^{(5)}(x) &= \makebox[3em]{$24/x^5$} \quad\Rightarrow  &f\,^{(5)}(1) &= 24\\
\vdots && \vdots & \\
f\,^{(n)}(x) &= \makebox[3em]{}\quad \Rightarrow & f\,^{(n)}(1) &= \smallskip\\
\makebox[1pt]{$\dfrac{(-1)^{n+1}(n-1)!}{x^n}$} && \makebox[1pt]{$(-1)^{n+1}(n-1)!$} &
%f(x) &= \ln x &\Rightarrow & &f(1) &= \phantom-0\\
%\fp(x) &= \phantom-1/x &\Rightarrow & &\fp(1) &= \phantom-1\\
%\fp'(x) &= -1/x^2 &\Rightarrow & &\fp'(1) &= -1\\
%\fp''(x) &= \phantom-2/x^3 &\Rightarrow & &\fp''(1) &= \phantom-2\\
%f\,^{(4)}(x) &= -6/x^4 &\Rightarrow & &f\,^{(4)}(1) &= -6\\
%f\,^{(5)}(x) &= 24/x^5 &\Rightarrow & &f\,^{(5)}(1) &= 24\\
%\ \vdots & & & & \ \vdots & \\
%f\,^{(n)}(x) &=  &\Rightarrow & & f\,^{(n)}(1) &= \smallskip\\
%\makebox[1pt]{$\dfrac{(-1)^{n+1}(n-1)!}{x^n}$} &&&& \makebox[1pt]{$(-1)^{n+1}(n-1)!$} &
\end{align*}}
Since $f(1) = \ln 1 = 0$, we skip the first term and start the summation with $n=1$, giving the Taylor series for $\ln x$, centered at $x=1$, as 
\[\sum_{n=1}^\infty (-1)^{n+1}(n-1)!\frac{1}{n!}(x-1)^n = \sum_{n=1}^\infty (-1)^{n+1}\frac{(x-1)^n}{n}.\]
\end{example}

It is important to note that \autoref{def:taylor_series} defines a Taylor series given a function $f(x)$; however, we \emph{cannot} yet state that $f(x)$ \emph{is equal} to its Taylor series. We will find that ``most of the time'' they are equal, but we need to consider the conditions that allow us to conclude this.

\autoref{thm:taylorthm} states that the error between a function and its $n^{\text{th}}$-degree Taylor polynomial is $R_n(x)$, where
\[\abs{R_n(x)}\leq \frac{\max\abs{\,f\,^{(n+1)}(z)}}{(n+1)!}\abs{x-c}^{n+1}.\]

If $R_n(x)$ goes to 0 for each $x$ in an interval $I$ as $n$ approaches infinity, we conclude that the function is equal to its Taylor series expansion.

\begin{theorem}[Function and Taylor Series Equality]\label{thm:function_series_equality}%
Let $f(x)$ have derivatives of all orders at $x=c$, let $R_n(x)$ be as stated in \autoref{thm:taylorthm}, and let $I$ be an interval on which the Taylor series of $f(x)$ converges. 
If $\ds\lim_{n\to\infty} R_n(x) = 0$ for all $x$ in $I$, then 
\index{Taylor Series!equality with generating function}
\[f(x) = \sum_{n=0}^\infty \frac{f\,^{(n)}(c)}{n!}(x-c)^n\ \text{ on $I$.}\]
\end{theorem}

We demonstrate the use of this theorem in an example.

\begin{example}[Establishing equality of a function and its Taylor series]\label{ex_ts3}%
Show that, for all $x$, $f(x) = \cos x$ is equal to its Maclaurin series as found in \autoref{ex_ts1}.
\solution
Given a value $x$, the magnitude of the error term $R_n(x)$ is bounded by
\[\abs{R_n(x)}\leq \frac{\max\abs{\,f\,^{(n+1)}(z)}}{(n+1)!}\abs{x}^{n+1}.\]
Since all derivatives of $\cos x$ are $\pm \sin x$ or $\pm\cos x$, whose magnitudes are bounded by $1$, we can state
\[\abs{R_n(x)}\leq \frac{1}{(n+1)!}\abs{x}^{n+1}\]
which implies
\begin{equation*}
 -\frac{\abs{x}^{n+1}}{(n+1)!} \leq R_n(x) \leq\frac{\abs{x}^{n+1}}{(n+1)!}.%\label{eq:coseqtaylor}
\end{equation*}
For any $x$, $\ds\lim_{n\to\infty} \frac{x^{n+1}}{(n+1)!} = 0$. Applying the Squeeze Theorem to our last inequality%\autoeqref{eq:coseqtaylor}
, we conclude that $\ds \lim_{n\to\infty} R_n(x) = 0$ for all $x$, and hence
\[\cos x = \sum_{n=0}^\infty (-1)^{n}\frac{x^{2n}}{(2n)!}\quad \text{for all $x$}.\]
\end{example}

It is natural to assume that a function is  equal to its Taylor series on the series' interval of convergence, but this is not necessarily the case. In order to properly establish equality, one must use \autoref{thm:function_series_equality}. This is a bit disappointing, as we developed beautiful techniques for determining the interval of convergence of a power series, and proving that $R_n(x)\to 0$ can be cumbersome as it deals with high order derivatives of the function.

There is good news. A function $f(x)$ that is equal to its Taylor series, centered at any point the domain of $f(x)$, is said to be an \textbf{analytic function},\index{analytic function} and most, if not all, functions that we encounter within this course are analytic functions. Generally speaking, any function that one creates with elementary functions (polynomials, exponentials, trigonometric functions, etc.) that is not piecewise defined is probably analytic. For most functions, we assume the function is equal to its Taylor series on the series' interval of convergence and only use \autoref{thm:function_series_equality} when we suspect something may not work as expected.  The converse is also true: if a function is equal to \emph{some} power series on an interval, then that power series is the Taylor series of the function.

We develop the Taylor series for one more important function, then give a table of the Taylor series for a number of common functions.\index{Binomial Series}\index{series!Binomial}

\begin{example}[The Binomial Series]\label{ex_ts4}%
Find the Maclaurin series of $f(x) = (1+x)^k$, $k\neq 0$.
\solution
When $k$ is a positive integer, the Maclaurin series is finite. For instance, when $k=4$, we have 
\[f(x) = (1+x)^4 = 1+4x+6x^2+4x^3+x^4.\]
The coefficients of $x$ when $k$ is a positive integer are known as the \emph{binomial coefficients}, giving the series we are developing its name.

When $k=1/2$, we have $f(x) = \sqrt{1+x}$. Knowing a series representation of this function would give a useful way of approximating $\sqrt{1.3}$, for instance.

To develop the Maclaurin series for $f(x) = (1+x)^k$ for any value of $k\neq0$, we consider the derivatives of $f$ evaluated at $x=0$:
{%
\small%
\begin{align*}
f(x) &= (1+x)^k & f(0) &= 1\\
\fp(x) &= k(1+x)^{k-1} & \fp(0) &=k\\
\fp'(x) &= k(k-1)(1+x)^{k-2} & \fp'(0) &=k(k-1)\\
\fp''(x) &= k(k-1)(k-2)(1+x)^{k-3} & \fp''(0) &=k(k-1)(k-2)\\
&\vdots & &\vdots\\
f\,^{(n)}(x) &= k(k-1)\dotsm\bigl(k-(n-1)\bigr)(1+x)^{k-n}
\mkern-3\thickmuskip&&\vdots \\
&& f\,^{(n)}(0) &=k(k-1)\dotsm\bigl(k-(n-1)\bigr)
\end{align*}}%
Thus the Maclaurin series for $f(x) = (1+x)^k$ is
\[
1+ kx + \frac{k(k-1)}{2!}x^2 + \frac{k(k-1)(k-2)}{3!}x^3 + \dotsb + \frac{k(k-1)\dotsm\bigl(k-(n-1)\bigr)}{n!}x^n+\dotsb
\]

It is important to determine the interval of convergence of this series. With 
\[a_n = \frac{k(k-1)\dotsm\bigl(k-(n-1)\bigr)}{n!}x^n,\]
we apply the Ratio Test:
\begin{align*}
	\lim_{n\to\infty}\frac{\abs{a_{n+1}}}{\abs{a_n}}
	&=\lim_{n\to\infty}\frac{\abs{\frac{k(k-1)\dotsm(k-n)}{(n+1)!}x^{n+1}}}{\abs{\frac{k(k-1)\dotsm\bigl(k-(n-1)\bigr)}{n!}x^n}}\\
		&=\lim_{n\to\infty} \abs{\frac{k-n}{n+1}x}\\
		&= \abs x.
\end{align*}

The series converges absolutely when the limit of the Ratio Test is less than 1; therefore, we have absolute convergence when $\abs x<1$. 

While outside the scope of this text, the interval of convergence depends on the value of $k$. When $k>0$, the interval of convergence is $[-1,1]$. When $-1<k<0$, the interval of convergence is $(-1,1]$. If $k\leq -1$, the interval of convergence is $(-1,1)$.
%When $x=1$, we can apply the Alternating Series Test and find the series converges. When $x=-1$, it can be shown (with some difficulty) that the series also converges. Therefore the interval of convergence is $[-1,1]$. We can apply \autoref{thm:function_series_equality} to prove equality between $f(x)$ and the series (or apply the discussion following the theorem). 
\end{example}

We learned that Taylor polynomials offer a way of approximating a ``difficult to compute'' function with a polynomial. Taylor series offer a way of exactly representing a function with a series. One probably can see the use of a good approximation; is there any use of representing a function exactly as a series?

While we appreciate the mathematical beauty of Taylor series (which is reason enough to study them), there are practical uses as well. They provide a valuable tool for solving a variety of problems, including problems relating to integration and differential equations.

In \autoref{idea:common_taylor} (on the following page) we give  a table of the Maclaurin series of a number of common functions. We then give a theorem about the ``algebra of power series,'' that is, how we can combine power series to create power series of new functions. This allows us to find the Taylor series of functions like $f(x) = e^x\cos x$ by knowing the Taylor series of $e^x$ and $\cos x$.

Before we investigate combining functions, consider the Taylor series for the arctangent function (see \autoref{idea:common_taylor}). Knowing that $\tan^{-1}(1) = \pi/4$, we can use this series to approximate the value of $\pi$:
\begin{align*}
\frac{\pi}4 &= \tan^{-1}(1) = 1-\frac13+\frac15-\frac17+\frac19-\dotsb\\
\pi &= 4\left(1-\frac13+\frac15-\frac17+\frac19-\dotsb\right)
\end{align*} 
Unfortunately, this particular expansion of $\pi$ converges very slowly. The first 
100 terms approximate $\pi$ as $3.13159$, which is not particularly good.

\clearpage

{
\tcbset{grow to right by=12em}
\begin{keyidea}[Important Maclaurin Series Expansions]\label{idea:common_taylor}%
%
\noindent
\tagpdfsetup{table/header-rows={1}}
\begin{tabular}{llc}
\textbf{Function and Series} & \textbf{First Few Terms} & \parbox{55pt}{\centering\textbf{Interval of}\\\textbf{Convergence}} \\
$\ds e^x = \sum_{n=0}^\infty \frac{x^n}{n!}$ & $\ds 1+ x+\frac{x^2}{2!} + \frac{x^3}{3!}+\dotsb$ & $(-\infty,\infty)$\medskip\\ % Ex {ex_taypoly1} (ish) 9.9.1
$\ds \sin x = \sum_{n=0}^\infty (-1)^n\frac{x^{2n+1}}{(2n+1)!}$ & $\ds x-\frac{x^3}{3!}+\frac{x^5}{5!} - \frac{x^7}{7!}+\dotsb$ & $(-\infty,\infty)$\medskip\\
$\ds \cos x = \sum_{n=0}^\infty (-1)^n\frac{x^{2n}}{(2n)!}$ & $\ds 1-\frac{x^2}{2!}+\frac{x^4}{4!} - \frac{x^6}{6!} +\dotsb$ & $(-\infty,\infty)$\smallskip\\ % ex {ex_ts1} 9.10.1
$\ds \ln(x+1) = \sum_{n=1}^\infty(-1)^{n+1}\frac{x^n}{n}$ & $\ds x-\frac{x^2}{2}+\frac{x^3}{3}-\dotsb$& $(-1,1]$\smallskip\\ % Ex 9.10.2 (ish)
$\ds \frac{1}{1-x} = \sum_{n=0}^\infty x^n$ &$\ds 1+x+x^2+x^3+\dotsb$& $(-1,1)$\\ % equ {eq:power_series_geometric} 9.8.1
\small$\ds (1+x)^k=\sum_{n=0}^\infty \frac{k(k-1)\dotsm\bigl(k-(n-1)\bigr)}{n!}x^n$ \normalsize& $\ds 1+kx+\frac{k(k-1)}{2!}x^2 + \dotsb$ & {$\begin{cases}(-1,1)&\phantom{-}k\le-1\\{}(-1,1]&-1<k<0\\{}[-1,1]&\phantom{-}0<k\end{cases}$}\\ % ex {ex_ts4} 9.10.4
$\ds \tan^{-1}x = \sum_{n=0}^\infty (-1)^n\frac{x^{2n+1}}{2n+1}$ & $\ds x-\frac{x^3}{3}+\frac{x^5}{5}-\frac{x^7}{7}+\dotsb$ & $[-1,1]$
\end{tabular}\index{Taylor Series!common series}
\end{keyidea}
}

% todo Tim should we mention where in the text we derive the ``Important Maclaurin Series Expansions''
% todo Tim Do we want to derive the Maclaurin series for sin and arctan?

\begin{theorem}[Algebra of Power Series]\label{thm:series_alg}%
Let $\ds f(x) = \sum_{n=0}^\infty a_nx^n$ and $\ds g(x) = \sum_{n=0}^\infty b_nx^n$ converge absolutely for $\abs x<R$, and let $h(x)$ be continuous.
\index{power series!algebra of} 
\begin{enumerate}
	\item	$\ds f(x)\pm g(x) = \sum_{n=0}^\infty (a_n\pm b_n)x^n$ \quad for $\abs x<R$.
	\item	$\ds f(x)g(x) = \left(\sum_{n=0}^\infty a_nx^n\right)\left(\sum_{n=0}^\infty b_nx^n\right) =\\
	\mbox{}\qquad
	\sum_{n=0}^\infty\bigl(a_0b_n+a_1b_{n-1}+\dotsb+ a_nb_0\bigr)x^n$ for $\abs x<R$.
	\item	$\ds f\bigl(h(x)\bigr) = \sum_{n=0}^\infty a_n\bigl(h(x)\bigr)^n$ \quad for $\abs{h(x)}<R$.
\end{enumerate}
\end{theorem}

\begin{example}[Combining Taylor series]\label{ex_ts5}%
Write out the first 3 terms of the Maclaurin Series for $f(x) = e^x\cos x$ using \autoref{idea:common_taylor} and \autoref{thm:series_alg}.
\solution
\autoref{idea:common_taylor} informs us that 
\[e^x = 1+x+\frac{x^2}{2!}+\frac{x^3}{3!}+\dotsb\qquad \text{and}\qquad \cos x = 1-\frac{x^2}{2!}+\frac{x^4}{4!}+\dotsb.\]
Applying \autoref{thm:series_alg}, we find that
\begin{align*}
e^x\cos x &= \left(1+x+\frac{x^2}{2!}+\frac{x^3}{3!}+\dotsb\right)\left(1-\frac{x^2}{2!}+\frac{x^4}{4!}+\dotsb\right).\\
\intertext{Distribute the right hand expression across the left:}
	&= 1\left(1-\frac{x^2}{2!}+\frac{x^4}{4!}+\dotsb\right)
	+x\left(1-\frac{x^2}{2!}+\frac{x^4}{4!}+\dotsb\right)\\
	&\phantom{=}+\frac{x^2}{2!}\left(1-\frac{x^2}{2!}+\frac{x^4}{4!}+\dotsb\right)
	+\frac{x^3}{3!}\left(1-\frac{x^2}{2!}+\frac{x^4}{4!}+\dotsb\right)\\
	&\phantom{=}+ \frac{x^4}{4!}\left(1-\frac{x^2}{2!}+\frac{x^4}{4!}+\dotsb\right)+\dotsb\\
	\intertext{Distribute again and collect like terms.}
	&= 1 + x -\frac{x^3}{3}-\frac{x^4}{6} - \frac{x^5}{30}+\frac{x^7}{630}+\dotsb
\end{align*}
While this process is a bit tedious, it is much faster than evaluating all the necessary derivatives of $e^x\cos x$ and computing the Taylor series directly.

Because the series for $e^x$ and $\cos x$ both converge on $(-\infty,\infty)$, so does the series expansion for $e^x\cos x$.
\end{example}

\begin{example}[Creating new Taylor series]\label{ex_ts6}%
Use \autoref{thm:series_alg} to create the Taylor series for $y=\sin(x^2)$ centered at $x=0$ and a series for $y=\ln (\sqrt{x})$ centered at $c=1$.
\solution
Given that 
\[\sin x = \sum_{n=0}^\infty (-1)^n\frac{x^{2n+1}}{(2n+1)!} = x-\frac{x^3}{3!}+\frac{x^5}{5!} -\frac{x^7}{7!}+\dotsb,\]
we simply substitute $x^2$ for $x$ in the series, giving
\[
\sin (x^2) = \sum_{n=0}^\infty (-1)^n\frac{(x^2)^{2n+1}}{(2n+1)!} = \sum_{n=0}^\infty (-1)^n\frac{x^{4n+2}}{(2n+1)!} = x^2-\frac{x^6}{3!}+\frac{x^{10}}{5!} -\frac{x^{14}}{7!}\dotsb.
\]
Since the Taylor series for $\sin x$ has an infinite radius of convergence, so does the Taylor series for $\sin(x^2)$.\bigskip

The Taylor expansion for $\ln(x+1)$ given in \autoref{idea:common_taylor} is centered at $x=0$, so we can center the series for $\ln (\sqrt{x})$ at $x=1$.
With 
\[\ln x = \sum_{n=1}^\infty(-1)^{n+1}\frac{(x-1)^n}{n} = (x-1)- \frac{(x-1)^2}{2} +\frac{(x-1)^3}{3}-\dotsb,\]
we substitute $\sqrt{x}$ for $x$ to obtain
\[\ln (\sqrt{x}) = \sum_{n=1}^\infty(-1)^{n+1}\frac{(\sqrt{x}-1)^n}{n} = (\sqrt{x}-1)- \frac{(\sqrt{x}-1)^2}{2} +\frac{(\sqrt{x}-1)^3}{3}-\dotsb.\]
While this is not strictly a power series because of the $\sqrt x$, it is a series that allows us to study the function $\ln(\sqrt{x})$. Since the interval of convergence of $\ln x$ is $(0,2]$, and the range of $\sqrt{x}$ on $(0,4]$ is $(0,2]$, the interval of convergence of this series expansion of $\ln(\sqrt{x})$ is $(0,4]$.
\end{example}

\mnote[yshift=12\baselineskip]{\textbf{Note:} In \autoref{ex_ts6}, one could create a series for $\ln(\sqrt{x})$ by simply recognizing that $\ln(\sqrt{x}) = \ln (x^{1/2}) = 1/2\ln x$, and hence multiplying the Taylor series for $\ln x$ by $1/2$. This example was chosen to demonstrate other aspects of series, such as the fact that the interval of convergence changes.}

\begin{example}[Using Taylor series to evaluate definite integrals]\label{ex_ts7}%
Use the Taylor series of $e^{-x^2}$ to evaluate $\ds \int_0^1e^{-x^2}\dd x$.
\solution
We learned, when studying Numerical Integration, that $e^{-x^2}$ does not have an antiderivative expressible in terms of elementary functions. This means any definite integral of this function must have its value approximated, and not computed exactly.

We can quickly write out the Taylor series for $e^{-x^2}$ using the Taylor series of $e^x$:\vspace{-.5\baselineskip}
\begin{align*}
e^x &= \sum_{n=0}^\infty \frac{x^n}{n!} = 1+x+\frac{x^2}{2!}+\frac{x^3}{3!}+\dotsb\\
\intertext{and so}
e^{-x^2} &= \sum_{n=0}^\infty \frac{(-x^2)^n}{n!} \\
				&= \sum_{n=0}^\infty (-1)^n\frac{x^{2n}}{n!}\\
				&= 1-x^2+\frac{x^4}{2!}-\frac{x^6}{3!}+\dotsb.
\end{align*}
We use \autoref{thm:calc_power_series} to integrate:
\[\int e^{-x^2}\dd x = C + x - \frac{x^3}{3}+\frac{x^5}{5\cdot2!}-\frac{x^7}{7\cdot3!}+\dotsb +(-1)^n\frac{x^{2n+1}}{(2n+1)n!}+\dotsb\]
This \emph{is} the antiderivative of $e^{-x^2}$; while we can write it out as a series, we cannot write it out in terms of elementary functions. We can evaluate the definite integral $\ds \int_0^1e^{-x^2}\dd x$ using this antiderivative; substituting 1 and 0 for $x$ and subtracting gives
\[\int_0^1e^{-x^2}\dd x = 1-\frac{1}{3}+\frac{1}{5\cdot 2!}-\frac{1}{7\cdot3!} + \frac{1}{9\cdot4!}-\dotsb.\]
Summing the 5 terms shown above gives the approximation of $0.74749.$ Since this is an alternating series, we can use the Alternating Series Approximation Theorem, (\autoref{thm:alt_series_approx}), to determine how accurate this approximation is. The next term of the series is $ 1/(11\cdot5!) \approx 0.00075758$. Thus we know our approximation is within $0.00075758$ of the actual value of the integral. This is arguably much less work than using Simpson's Rule to approximate the value of the integral.
\end{example}

Another advantage to using Taylor series instead of Simpson's Rule is for making subsequent approximations.  We found in \autoref{ex_num7} that the error in using Simpson's Rule for $\ds\int_0^1 e^{-x^2}\dd x$ with four intervals was $0.00026$.  If we wanted to decrease that error, we would need to use more intervals, essentially starting the problem over.  Using a Taylor series, if we wanted a more accurate approximation, we can just subtract the next term $1/(11\cdot5!)$ to get an approximation of $0.7467$, with an error of at most $1/(13\cdot6!)\approx0.0001$.\bigskip

%\begin{example}[Using Taylor series to solve differential equations]\label{ex_ts8}%
%Solve the differential equation $y\primeskip'=2y$ in terms of a power series, and use the theory of Taylor series to recognize the solution in terms of an elementary function.
%\solution
%We found the first 5 terms of the power series solution to this differential equation in \autoref{ex_ps5} in \autoref{sec:power_series}. These are:
%\[a_0=1,\quad a_1 = 2,\quad a_2 = \frac42=2,\quad a_3=\frac{8}{2\cdot3}=\frac43,\quad a_4=\frac{16}{2\cdot3\cdot4} = \frac23.\]
%We include the ``unsimplified'' expressions for the coefficients found in \autoref{ex_ps5} as we are looking for a pattern. It can be shown that $a_n = 2^n/n!$. Thus the solution, written as a power series, is
%\[y = \sum_{n=0}^\infty \frac{2^n}{n!}x^n = \sum_{n=0}^\infty \frac{(2x)^n}{n!}.\]
%Using \autoref{idea:common_taylor} and \autoref{thm:series_alg}, we recognize $f(x) = e^{2x}$:
%\[e^x = \sum_{n=0}^\infty \frac{x^n}{n!} \qquad \Rightarrow \qquad e^{2x} = \sum_{n=0}^\infty \frac{(2x)^n}{n!}.\]
%\end{example}

Finding a pattern in the coefficients that match the series expansion of a known function, such as those shown in \autoref{idea:common_taylor}, can be difficult. What if the coefficients %in the previous example were
are given in their reduced form; how could we still recover the function?% $y=e^{2x}$?

Suppose that all we know is that 
\[a_0=1,\quad a_1=2,\quad a_2=2,\quad a_3=\frac43,\quad a_4=\frac23.\]
\autoref{def:taylor_series} states that each term of the Taylor expansion of a function includes an $n!$. This allows us to say that
\[a_2=2=\frac{b_2}{2!},\quad a_3 = \frac43=\frac{b_3}{3!},\quad \text{and}\quad a_4 = \frac23=\frac{b_4}{4!}\]
for some values $b_2$, $b_3$ and $b_4$.
Solving for these values, we see that $b_2=4$, $b_3 = 8$ and $b_4=16$. That is, we are recovering the pattern $b_n=2^n$, %we had previously seen,
allowing us to write 
\begin{align*}
	f(x) = \sum_{n=0}^\infty a_nx^n
	&= \sum_{n=0}^\infty \frac{b_n}{n!}x^n \\
	&= 1+2x+ \frac{4}{2!}x^2 + \frac{8}{3!}x^3+\frac{16}{4!}x^4 + \dotsb
\end{align*}
From here it is easier to recognize that the series is describing an exponential function.

%There are simpler, more direct ways of solving the differential equation $y\primeskip' = 2y$. We applied power series techniques to this equation to demonstrate its utility, and went on to show how \emph{sometimes} we are able to recover the solution in terms of elementary functions using the theory of Taylor series. Most differential equations faced in real scientific and engineering situations are much more complicated than this one, but power series can offer a valuable tool in finding, or at least approximating, the solution.\bigskip

This chapter introduced sequences, which are ordered lists of numbers, followed by series, wherein we add up the terms of a sequence. We quickly saw that such sums do not always add up to ``infinity,'' but rather converge. We studied tests for convergence, then ended the chapter with a formal way of defining functions based on series. Such ``series-defined functions'' are a valuable tool in solving a number of different problems throughout science and engineering.

Coming in the next chapters are new ways of defining curves in the plane apart from using functions of the form $y=f(x)$. Curves created by these new methods can be beautiful, useful, and important. 

% todo have the exercises find the interval of convergence where appropriate

\printexercises{exercises/08-08-exercises}



\apexchapter{Curves in the Plane}{chapter:planar_curves}
%\apexchapter[text/09_Conic_Sections]{Curves in the Plane}{chapter:planar_curves}

We have explored functions of the form $y=f(x)$ closely throughout this text. We have explored their limits, their derivatives and their antiderivatives; we have learned to identify key features of their graphs, such as relative maxima and minima, inflection points and asymptotes; we have found equations of their tangent lines, the areas between portions of their graphs and the $x$-axis, and the volumes of solids generated by revolving portions of their graphs about a horizontal or vertical axis.

Despite all this, the graphs created by functions of the form $y=f(x)$ are limited. Since each $x$-value can correspond to only 1 $y$-value, common shapes like circles cannot be fully described by a function in this form.  Fittingly, the ``vertical line test''  excludes vertical lines from being functions of $x$, even though these lines are important in mathematics.

In this chapter we'll explore new ways of drawing curves in the plane. We'll still work within the framework of functions, as an input will still only correspond to one output. However, our new techniques of drawing curves will render the vertical line test pointless, and allow us to create important --- and beautiful --- new curves. Once these curves are defined, we'll apply the concepts of calculus to them, continuing to find equations of tangent lines and the areas of enclosed regions.

One aspect that we'll be interested in is ``how long is this curve?'' Before we explore that idea for these new ways to draw curves, we'll start by exploring how long a curve is when we've gotten it from a regular $y=f(x)$ function.

\inputprereq{text/09_Conic_Sections}
\section{Arc Length and Surface Area}\label{sec:arc_length}

In previous sections we have used integration to answer the following questions:
\begin{enumerate}
	\item	Given a region, what is its area?
	\item	Given a solid, what is its volume?
\end{enumerate}

In this section, we address a related question: Given a curve, what is its length? This is often referred to as \textbf{arc length}. 

\mtable{Graphing $y=\sin x$ on $[0,\pi]$ and approximating the curve with line segments.}{fig:arcintro}{%
\begin{tikzpicture}[alt={Blue y = sin x arch on 0 < x < π, peak at (π/2 , 1); ticks at π/4, π/2, 3π/4.}]
 \begin{axis}[width=\marginparwidth,%height=.46\marginparwidth,
   tick label style={font=\scriptsize},axis equal,
   axis y line=middle,axis x line=middle,name=myplot,axis on top,
   xtick={.79,1.57,2.36,3.14},
   xticklabels={$\frac{\pi}4$,$\frac{\pi}2$,$\frac{3\pi}{4}$,$\pi$},
   ymin=-.2,ymax=1.25,
   xmin=-.1,xmax=3.5]
  \addplot [draw={\colorone},smooth,thick,domain=0:3.14] {sin(deg(x))};
 \end{axis}
 \node [right] at (myplot.right of origin) {\scriptsize $x$};
 \node [above] at (myplot.above origin) {\scriptsize $y$};
\end{tikzpicture}
\\(a)\\
\begin{tikzpicture}[alt={Same sin curve; red polyline through five black nodes approximates it on 0 < x < π.}]
 \begin{axis}[width=\marginparwidth,%height=.46\marginparwidth,
   tick label style={font=\scriptsize},
   axis y line=middle,axis x line=middle,name=myplot,axis on top,%
   xtick={.79,1.57,2.36,3.14},
   xticklabels={$\frac{\pi}4$,$\frac{\pi}2$,$\frac{3\pi}{4}$,$\pi$},
   ytick={.71,1},%
   yticklabels={$\frac1{\sqrt2}$,$1$},axis equal,
   ymin=-.2,ymax=1.25,
   xmin=-.1,xmax=3.5]
  \addplot [draw={\colorone},smooth,thick,domain=0:3.14] {sin(deg(x))};
  \draw [draw={\colortwo},thick] (axis cs:0,0) -- (axis cs:.79,.71) -- (axis cs: 1.57,1) -- (axis cs:2.36,.71) -- (axis cs: 3.14,0);
  \filldraw (axis cs:0,0) circle (1pt) (axis cs:.79,.71) circle (1pt) (axis cs: 1.57,1) circle (1pt) (axis cs:2.36,.71) circle (1pt) (axis cs: 3.14,0)circle (1pt);
 \end{axis}
 \node [right] at (myplot.right of origin) {\scriptsize $x$};
 \node [above] at (myplot.above origin) {\scriptsize $y$};
\end{tikzpicture}
\\(b)}

Consider the graph of $y=\sin x$ on $[0,\pi]$ given in \autoref{fig:arcintro} (a). How long is this curve? That is, if we were to use a piece of string to exactly match the shape of this curve, how long would the string be?

As we have done in the past, we start by approximating; later, we will refine our answer using limits to get an exact solution.

The length of straight-line segments is easy to compute using the Distance Formula. We can approximate the length of the given curve by approximating the curve with straight lines and measuring their lengths. 

In \autoref{fig:arcintro} (b), the curve $y=\sin x$ has been approximated with 4 line segments (the interval $[0,\pi]$ has been divided into 4 equal length subintervals). It is clear that these four line segments approximate $y=\sin x$ very well on the first and last subinterval, though not so well in the middle. Regardless, the sum of the lengths of the line segments is $3.79$, so we approximate the arc length of $y=\sin x$ on $[0,\pi]$ to be $3.79$. 

%Using 6 evenly spaced subintervals is not too much more work and gives a better approximation. The six lines are shown in \autoref{fig:arcintro} (c) and provide an arc length approximation of $3.81$. 

In general,  we can approximate the arc length of $y=f(x)$ on $[a,b]$ in the following manner. Let $a=x_0 < x_1 < \dotso < x_{n-1}< x_n=b$ be a partition of $[a,b]$ into $n$ subintervals. Let $\Delta x_i$ represent the length of the $i^{\text{th}}$ subinterval $[x_{i-1},x_i]$.

\autoref{fig:arcintro2} zooms in on the $i^{\text{th}}$ subinterval where $y=f(x)$ is approximated by a straight line segment. The dashed lines show that we can view this line segment as the hypotenuse of a right triangle whose sides have length $\Delta x_i$ and $\Delta y_i$. Using the Pythagorean Theorem, the length of this line segment is $\ds \sqrt{(\Delta x_i)^2 + (\Delta y_i)^2}.$ Summing over all subintervals gives an arc length approximation
%
\mtable{Zooming in on the $i^{\text{th}}$ subinterval $[x_{i-1},x_i]$ of a partition of $[a,b]$.}{fig:arcintro2}{\begin{tikzpicture}[alt={Zoom on x_{i−1} < x < x_i: blue arc, red chord, labeled delta x horizontal and delta y vertical.}]
 \begin{axis}[width=\marginparwidth,tick label style={font=\scriptsize},
   axis y line=middle,axis x line=middle,name=myplot,axis on top,
   xtick=\empty,extra x ticks={.3,1.37},extra x tick labels={$x_{i-1}$,$x_i$},
   ytick=\empty,extra y ticks={.48,1},extra y tick labels={$y_{i-1}$,$y_i$},
   ymin=-.2,ymax=1.25,xmin=-.1,xmax=1.6]
  \addplot [draw={\colorone},smooth,thick,domain=0.3:1.37] {sin(deg(x+.2))};
  \draw (axis cs: .25,1) -- (axis cs:.1,1) 
        (axis cs: .25,.48) -- (axis cs:.1,.48)
        (axis cs: .175,.48)
         -- node [pos=.5,fill=white] {\scriptsize $\Delta y_i$} (axis cs:.175,1);
  \draw (axis cs: .3,.25) -- (axis cs: .3,.4)
        (axis cs: 1.37,.25) -- (axis cs: 1.37,.4)
        (axis cs: .3,.325)
         -- node [pos=.5,fill=white] {\scriptsize $\Delta x_i$} (axis cs: 1.37,.325);
  \draw [draw={\colortwo},thick] (axis cs: .3,.48) -- (axis cs:1.37,1);
  \draw [dashed,thin] (axis cs: .3,.48) -| (axis cs:1.37,1);
  \filldraw (axis cs: .3,.48) circle (1pt) (axis cs:1.37,1) circle (1pt);
 \end{axis}
 \node [right] at (myplot.right of origin) {\scriptsize $x$};
 \node [above] at (myplot.above origin) {\scriptsize $y$};
\end{tikzpicture}}
%
\[L \approx \sum_{i=1}^n \sqrt{(\Delta x_i)^2 + (\Delta y_i)^2}.\]

As it is written, this is \emph{not} a Riemann Sum. While we could conclude that taking a limit as the subinterval length goes to zero gives the exact arc length, we would not be able to compute the answer with a definite integral. We need first to do a little algebra.

In the above expression factor out a $\Delta x_i^2$ term:
\begin{align*}
	\sum_{i=1}^n \sqrt{(\Delta x_i)^2 + (\Delta y_i)^2}
	&= \sum_{i=1}^n \sqrt{(\Delta x_i)^2\left(1 + \frac{(\Delta y_i)^2}{(\Delta x_i)^2}\right)}.
\end{align*}
Now pull the $(\Delta x_i)^2$ term out of the square root:
\[L\approx\sum_{i=1}^n\sqrt{1 + \frac{(\Delta y_i)^2}{(\Delta x_i)^2}}\ \Delta x_i.\]
This is nearly a Riemann Sum. Consider the $(\Delta y_i)^2/(\Delta x_i)^2$ term. The expression $\Delta y_i/\Delta x_i$ measures the ``change in $y$/change in $x$,'' that is, the ``rise over run'' of $f$ on the $i^{\text{th}}$ subinterval. The Mean Value Theorem of Differentiation (\autoref{thm:mvt}) states that there is a $c_i$ in the $i^{\text{th}}$ subinterval where $\fp(c_i) = \Delta y_i/\Delta x_i$. Thus we can rewrite our above expression as:
\[L\approx\sum_{i=1}^n\sqrt{1+[\fp(c_i)]^2}\ \Delta x_i.\]
This \emph{is} a Riemann Sum. As long as \fp\ is continuous on $[a,b]$, we can invoke \autoref{thm:riemannSum} and conclude
\[L=\int_a^b\sqrt{1+[\fp(x)]^2}\dd x.\]

\begin{keyidea}[Arc Length]\label{idea:arclength}%
Let $f$ be differentiable on an open interval containing $[a,b]$, where $\fp$ is also continuous on $[a,b]$. Then the arc length of $f$ from $x=a$ to $x=b$ is
\index{integration!arc length}\index{arc length}
\[L = \int_a^b \sqrt{1+[\fp(x)]^2}\dd x.\]
\end{keyidea}

\youtubeVideo{PwmCZAWeRNE}{Arc Length}

As the integrand contains a square root, it is often difficult to use the formula in \autoref{idea:arclength} to find the length exactly. When exact answers are difficult to come by, we resort to using numerical methods of approximating definite integrals. The following examples will demonstrate this.

\begin{example}[Finding arc length]\label{ex_arc1}%
Find the arc length of $f(x) = x^{3/2}$ from $x=0$ to $x=4$.
%
\mtable{A graph of $f(x) = x^{3/2}$ from \autoref{ex_arc1}.}{fig:arc1}{\begin{tikzpicture}[alt={Blue y = x^(3/2) curve on 0 < x < 4 rises steeply, reaching 8 at x = 4.}]
\begin{axis}[width=1.16\marginparwidth,tick label style={font=\scriptsize},
axis y line=middle,axis x line=middle,name=myplot,axis on top,
ymin=-.2,ymax=8.5,xmin=-.1,xmax=4.5,]
\addplot [draw={\colorone},smooth,thick,domain=0:4,samples=50] {x^(3/2)};
%\addplot [draw={\colorone},smooth,thick,domain=0:4,samples=50] coordinates {(0,0)(0.1,0.03162)(0.2,0.08944)(0.3,0.1643)(0.4,0.253)(0.5,0.3536)(0.6,0.4648)(0.7,0.5857)(0.8,0.7155)(0.9,0.8538)(1.,1.)(1.3,1.482)(1.6,2.024)(1.9,2.619)(2.2,3.263)(2.5,3.953)(2.8,4.685)(3.1,5.458)(3.4,6.269)(3.7,7.117)(4.,8.)};
\end{axis}
\node [right] at (myplot.right of origin) {\scriptsize $x$};
\node [above] at (myplot.above origin) {\scriptsize $y$};
\end{tikzpicture}}
%
\solution
A graph of $f$ is given in \autoref{fig:arc1}.
We begin by finding $\fp(x)= \frac32x^{1/2}$. Using the formula, we find the arc length $L$ as
\begin{align*}
	L
	&=	\int_0^4 \sqrt{1+\left(\frac32x^{1/2}\right)^2}\dd x \\
	&=	\int_0^4 \sqrt{1+\frac94x}\ \dd x \\
	&= 	\int_0^4 \left(1+\frac94x\right)^{1/2}\dd x \\
	&=  \frac23\cdot\frac49\left(1+\frac94x\right)^{3/2}\Big|_0^4 \\
	&=	\frac{8}{27}\left(10^{3/2}-1\right)\text{units}. % \approx 9.07 
\end{align*}
\end{example}

\begin{example}[Finding arc length]\label{ex_arc2}%
Find the arc length of $\ds f(x) =\frac18x^2-\ln x$ from $x=1$ to $x=2$.
\solution
A graph of $f$ is given in \autoref{fig:arc2}; the portion of the curve measured in this problem is in bold.
\mtable{A graph of $f(x) =\frac18x^2-\ln x$ from \autoref{ex_arc2}.}{fig:arc2}{\begin{tikzpicture}[alt={Blue y = x²/8 − ln x on 0 < x < 3 dives from 1, bottoms near (2, -0.2), then turns upward.}]
\begin{axis}[width=\marginparwidth,tick label style={font=\scriptsize},
axis y line=middle,axis x line=middle,name=myplot,axis on top,
ymin=-.26,ymax=1.2,xmin=-.1,xmax=3.1]
\addplot [draw={\colorone},smooth,thin,domain=.2:3] {x*x/8-ln(x)};
\addplot [draw={\colorone},smooth,very thick,domain=1:2] {x*x/8-ln(x)};
\end{axis}
\node [right] at (myplot.right of origin) {\scriptsize $x$};
\node [above] at (myplot.above origin) {\scriptsize $y$};
\end{tikzpicture}}
This function was chosen specifically because the resulting integral can be evaluated exactly. We begin by finding $\fp(x) = x/4-1/x$. The arc length is 
\begin{align*}
	L
	&=  \int_1^2 \sqrt{1+ \left(\frac x4-\frac1x\right)^2}\dd x \\
	&= 	\int_1^2 \sqrt{ 1 + \frac{x^2}{16} -\frac12 + \frac1{x^2} } \dd x \\
	&=	\int_1^2 \sqrt{ \frac{x^2}{16} +\frac12 + \frac1{x^2} } \dd x \\
	&=	\int_1^2	\sqrt{ \left(\frac x4 + \frac1x\right)^2 }\dd x \\
	&= \int_1^2 \left(\frac x4 + \frac1x\right) \dd x \\
	&=  \left(\frac{x^2}8 + \ln x\right)\Bigg|_1^2\\
	&=	\frac38+\ln 2\ \text{units}. % \approx 1.07
\end{align*}
\end{example}

The previous examples found the arc length exactly through careful choice of the functions. In general, exact answers are much more difficult to come by and numerical approximations are necessary.

\begin{example}[Approximating arc length numerically]\label{ex_arc3}%
Find the length of the sine curve from $x=0$ to $x=\pi$.
\solution
This is somewhat of a mathematical curiosity; \autoref{ex_ftc4} showed that the area under one ``hump'' of the sine curve is 2 square units; now we are measuring its arc length.

\mtable{A table of values of $y=\sqrt{1+\cos^2x}$ to evaluate a definite integral in \autoref{ex_arc3}.}{fig:arc3}{%
\[\begin{array}{cc}
x & \sqrt{1+\cos^2x} \\ \midrule
 0 & \sqrt{2} \\
 \pi/4 & \sqrt{3/2} \\
 \pi/2 & 1 \\
 3 \pi/4 & \sqrt{3/2} \\
 \pi  & \sqrt{2} \\
\end{array}\]}

The setup is straightforward: $f(x) = \sin x$ and $\fp(x) = \cos x$. Thus 
\[L = \int_0^\pi \sqrt{1+\cos^2x}\dd x.\]
This integral \emph{cannot} be evaluated in terms of elementary functions so we will approximate it with Simpson's Method with $n=4$. \autoref{fig:arc3} gives the integrand evaluated at 5 evenly spaced points in $[0,\pi]$. Simpson's Rule then states that 
\begin{align*}
	\int_0^\pi \sqrt{1+\cos^2x}\dd x
	&\approx	\frac{\pi-0}{4\cdot 3}
	\left(\sqrt{2}+4\sqrt{3/2}+2(1)+4\sqrt{3/2}+\sqrt{2}\right) \\
	&\approx3.82918.
\end{align*}
Using a computer with $n=100$ the approximation is $L\approx 3.8202$; our approximation with $n=4$ is quite good.  Our approximation of $3.79$ from the beginning of this section isn't as close.
\end{example}

\subsection{Surface Area of Solids of Revolution}

\mtable{Establishing the formula for surface area.}{fig:surface_intro}{%
\myincludeasythree{
3Droll=126.0060482236849,
3Dortho=0.004875946324318647,
3Dc2c=0.44462037086486816 0.39410829544067383 0.8043577671051025,
3Dcoo=67.21983337402344 3.5258262157440186 -29.566415786743164,
3Droo=150.00000160608386}{\marginparwidth}{Blue arc y = f(x) between x = a and x = b; short red chord shows one sample slice. Revolving that slice about the x-axis gives a purple band on a blue tubular surface.}{figures/figarc4_b_3D}
% todo Tim should this be 3D at all?
\\(a)\\
\myincludeasythree{
3Droll=126.0060482236849,
3Dortho=0.004875946324318647,
3Dc2c=0.44462037086486816 0.39410829544067383 0.8043577671051025,
3Dcoo=67.21983337402344 3.5258262157440186 -29.566415786743164,
3Droo=150.00000160608386}{\marginparwidth}{Blue band on blue tube: frustum formed when the red chord of the slice is revolved about the x-axis}{figures/figarc4_3D}
\\(b)}

We have already seen how a curve $y=f(x)$ on $[a,b]$ can be revolved around an axis to form a solid. Instead of computing its volume, we now consider its surface area.

We begin as we have in the previous sections: we partition the interval $[a,b]$ with $n$ subintervals, where the $i\,^{\text{th}}$ subinterval is $[x_i,x_{i+1}]$. On each subinterval, we can approximate the curve $y=f(x)$ with a straight line that connects $f(x_i)$ and $f(x_{i+1})$ as shown in \autoref{fig:surface_intro}(a). Revolving this line segment about the $x$-axis creates part of a cone (called a \emph{frustum} of a cone) as shown in \autoref{fig:surface_intro}(b). The surface area of a frustum of a cone is
\[A=2\pi r_{\text{avg}} L,\]
where $r_{\text{avg}}$ is the average of $R_1$ and $R_2$.  The length is given by $L$; we use the material just covered by arc length to state that
\[L\approx \sqrt{1+[\fp(c_i)]^2}\Delta x_i\]
for some $c_i$ in the $i^{\text{th}}$ subinterval. The radii are just the function evaluated at the endpoints of the interval: $f(x_{i-1})$ and $f(x_i)$. Thus the surface area of this sample frustum of the cone is approximately 
\[2\pi\frac{f(x_{i-1})+f(x_i)}2\sqrt{1+[\fp(c_i)]^2}\Delta x_i.\]

Since $f$ is a continuous function, the Intermediate Value Theorem states there is some $d_i$ in $[x_{i-1},x_i]$ such that $f(d_i)=\dfrac{f(x_{i-1})+f(x_i)}2$; we can use this to rewrite the above equation as
\[2\pi f(d_i)\sqrt{1+[\fp(c_i)]^2}\Delta x_i.\]
Summing over all the subintervals we get the total surface area to be approximately 
\[\text{Surface Area}\approx \sum_{i=1}^n 2\pi f(d_i)\sqrt{1+[\fp(c_i)]^2}\Delta x_i,\]
which is almost a Riemann Sum (we would need $d_i=c_i$ to remove the ``almost''). Taking the limit as the subinterval lengths go to zero gives us the exact surface area, given in the upcoming Key Idea.

If instead we revolve $y=f(x)$ about the $y$-axis, the radii of the resulting frustum are $x_{i-1}$ and $x_i$; their average value is simply the midpoint of the interval. In the limit, this midpoint is just $x$. This gives the second part of \autoref{idea:surface_area}.

% todo Tim does this (and prev KI) only need f' continuous on (a,b)?
\begin{keyidea}[Surface Area of a Solid of Revolution]\label{idea:surface_area}%
Let $f$ be differentiable on an open interval containing $[a,b]$ where $\fp$ is also continuous on $[a,b]$. \index{integration!surface area}\index{surface area!solid of revolution}
\begin{enumerate}
	\item	The surface area of the solid formed by revolving the graph of $y=f(x)$, where $f(x)\geq0$, about the $x$-axis is
	\[\text{Surface Area} = 2\pi\int_a^b f(x)\sqrt{1+[\fp(x)]^2}\dd x.\]
	\item	The surface area of the solid formed by revolving the graph of $y=f(x)$ about the $y$-axis, where $a,b\geq0$, is
	\[\text{Surface Area} = 2\pi\int_a^b x\sqrt{1+[\fp(x)]^2}\dd x.\]
\end{enumerate}
\end{keyidea}

\begin{example}[Finding surface area of a solid of revolution]\label{ex_sa1}%
Find the surface area of the solid formed by revolving $y=\sin x$ on $[0,\pi]$ around the $x$-axis, as shown in \autoref{fig:sa1}.
\solution
The setup turns out to be easier than the resulting integral. Using \autoref{idea:surface_area}, we have the surface area $SA$ is:
\mtable{Revolving $y=\sin x$ on $[0,\pi]$ about the $x$-axis.}{fig:sa1}{%
\myincludeasythree{
3Droll=121.00092658103458,
3Dortho=0.0048407199792563915,
3Dc2c=0.2899017035961151 0.17925968766212463 0.9401186108589172,
3Dcoo=72.58285522460938 2.891160011291504 -26.545438766479492,
3Droo=149.99998703790175}{\marginparwidth}{Blue sinus-shaped surface: y = sin x (0 < x < pi) revolved about the x-axis, axes shown.}{figures/figsa1_3D}}
\begin{align*}
	\text{SA}
	&=	2\pi\int_0^\pi \sin x\sqrt{1+\cos^2x}\dd x \\
	&=	-2\pi\int_1^{-1}\sqrt{1+u^2}\dd u \qquad\qquad \text{substitute $u=\cos x$} \\
	&=	2\pi\int_{-\pi/4}^{\pi/4}\sec^3\theta\dd\theta \qquad\qquad \text{substitute $u=\tan\theta$} \\
	&=  \left.\pi\left(\sec\theta\tan\theta+\ln\abs{\sec\theta+\tan\theta}\right)\right|_{-\pi/4}^{\pi/4} \qquad \text{by \autoref{ex_trigint6}} \\
	&=  \pi\left(\sqrt2+\ln\left(\sqrt2+1\right)
	-\left(-\sqrt2+\ln\left(\sqrt2-1\right)\right)\right) \\
	&=	\pi\left(2\sqrt2+\ln\left(\frac{\sqrt2+1}{\sqrt2-1}\right)\right) \\
	&=	2\pi\left(\sqrt2+\ln\left(\sqrt2+1\right)\right)\ \text{units}^2 \qquad\qquad\text{rationalize the denominator.}
\end{align*}

It is interesting to see that the surface area of a solid, whose shape is defined by a trigonometric function, involves both a square root and a natural logarithm.
\end{example}

\mtable{The solids used in \autoref{ex_sa2}.}{fig:sa2}{%
\myincludeasythree{
3Droll=138.26595091591264,
3Dortho=0.0048407199792563915,
3Dc2c=0.16831636428833008 0.19470007717609406 0.966313362121582,
3Dcoo=81.32470703125 3.403531074523926 -28.20997428894043,
3Droo=149.9999886313645}{\marginparwidth}{Blue horn-like surface: y = x^2 (0 < x < 1) rotated about x-axis, tip at (0,0,0) flaring to x = 1.}{figures/figsa2a_3D}
\\(a)\\[15pt]
\myincludeasythree{
3Droll=138.26595091591264,
3Dortho=0.0048407199792563915,
3Dc2c=0.16831637918949127 0.19470007717609406 0.966313362121582,
3Dcoo=0.013853734359145164 79.42704010009766 -29.364774703979492,
3Droo=149.99998895240248}{\marginparwidth}{Blue parabolic bowl: y = x^2 (0 < x < 1) revolved about the y-axis; vertex at origin, rim near y = 1.}{figures/figsa2b_3D}\\
(b)}

\begin{example}[Finding surface area of a solid of revolution]\label{ex_sa2}%
Find the surface area of the solid formed by revolving the curve $y=x^2$ on $[0,1]$ about:\\
\null\hfill 1.\quad the $x$-axis\hfill2.\quad the $y$-axis.\hfill\null
\solution
\begin{enumerate}
	\item		The solid formed by revolving $y=x^2$ around the $x$-axis is graphed in \autoref{fig:sa2}(a). Like the integral in \autoref{ex_sa1}, this integral is easier to setup than to actually integrate. While it is possible to use a trigonometric substitution to evaluate this integral, it is significantly more difficult than a solution employing the hyperbolic sine:
\begin{align*}
	SA &= 2\pi\int_0^1 x^2\sqrt{1+(2x)^2}\dd x. \\
%	&= 2\pi\int_0^{\tan^{-1}2} \frac18\tan^2\theta\sec^3\theta\ d\theta\qquad\text{substitute $x=\frac12\tan\theta$} \\
%	&= \frac\pi4\int_0^{\tan^{-1}2}\sec^5\theta-\sec^3\theta\ d\theta \\
	&= \left.\frac{\pi}{32}\left(2(8x^3+x)\sqrt{1+4x^2}-\sinh^{-1}(2x)\right)\right|_0^1\\
	&=\frac{\pi}{32}\left(18\sqrt{5}-\sinh^{-1}2\right)\ \text{units}^2.
\end{align*}

	\item	 Since we are revolving around the $y$-axis, the ``radius'' of the solid is not $f(x)$ but rather $x$. Thus the integral to compute the surface area is:
\begin{align*}
	SA &= 2\pi\int_0^1x\sqrt{1+(2x)^2}\dd x \\
	&=	\frac{\pi}4\int_1^5 \sqrt{u}\dd u \qquad\text{substitute $u=1+4x^2$} \\
	&= \left.\frac{\pi}{4}\frac23 u^{3/2}\right|_1^5\\
	&= \frac{\pi}6\left(5\sqrt{5}-1\right)\ \text{units}^2.
\end{align*}
 The solid formed by revolving $y=x^2$ about the $y$-axis is graphed in \autoref{fig:sa2} (b).
\end{enumerate}
\end{example}

Our final example is a famous mathematical ``paradox.''

\mtable{A graph of Gabriel's Horn.}{fig:gabriel}{%
\myincludeasythree{
3Droll=107.09033109095482,
3Dortho=0.0048407199792563915,
3Dc2c=0.4647676646709442 0.1565970778465271 0.8714748024940491,
3Dcoo=72.58285522460938 2.8911631107330322 -26.545440673828125,
3Droo=149.99999776925813}{\marginparwidth}{Infinite blue trumpet (Gabriel’s Horn): surface from rotating y = 1 / x, 1 < x < infty, about the x-axis.}{figures/figgabriel_3D}}
\begin{example}[The surface area and volume of Gabriel's Horn]\label{ex_gabriel}%
Consider the solid formed by revolving $y=1/x$ about the $x$-axis on $[1,\infty)$. Find the volume and surface area of this solid. (This shape, as graphed in \autoref{fig:gabriel}, is known as ``Gabriel's Horn'' since it looks like a very long horn that only a supernatural person, such as an angel, could play.)\index{Gabriel's Horn}
\solution
To compute the volume it is natural to use the Disk Method. We have:
\begin{align*}
V &= \pi\int_1^\infty \frac{1}{x^2}\dd x \\
	&= \lim_{b\to\infty}\pi\int_1^b\frac{1}{x^2}\dd x \\
	&=	\lim_{b\to\infty} \left.\pi\left(\frac{-1}{x}\right)\right|_1^b \\
	&= \lim_{b\to\infty} \pi\left(1-\frac1b\right) \\
	&= \pi \ \text{units}^3.
\end{align*}
Gabriel's Horn has a finite volume of $\pi$ cubic units. Since we have already seen that regions with infinite length can have a finite area, this is not too difficult to accept.

We now consider its surface area. The integral is straightforward to setup:
\begin{align*}
SA &= 2\pi\int_1^\infty \frac{1}{x}\sqrt{1+1/x^4}\dd x.
\intertext{Integrating this expression is not trivial. We can, however, compare it to other improper integrals. Since \intertextmath{1< \sqrt{1+1/x^4}} on \intertextmath{[1,\infty)}, we can state that}
2\pi\int_1^\infty \frac{1}{x}\dd x &<2\pi\int_1^\infty \frac{1}{x}\sqrt{1+1/x^4}\dd x .
\end{align*}
By \autoref{idea:impint1}, the improper integral on the left diverges. Since the integral on the right is larger, we conclude it also diverges, meaning Gabriel's Horn has infinite surface area.

Hence the ``paradox'': we can fill Gabriel's Horn with a finite amount of paint, but since it has infinite surface area, we can never paint it.

Somehow this paradox is striking when we think about it in terms of volume and area. However, we have seen a similar paradox before, as referenced above. We know that the area under the curve $y=1/x^2$ on $[1,\infty)$ is finite, yet the shape has an infinite perimeter. Strange things can occur when we deal with the infinite.
\end{example}

%A standard equation from physics is ``Work = force $\times$ distance'', when the force applied is constant. In the next section we learn how to compute work when the force applied is variable.

\printexercises{exercises/07-04-exercises}

\section{Parametric Equations}\label{sec:param_eqs}

We are familiar with sketching  shapes, such as parabolas, by following this basic procedure:
\[
\begin{gathered}\text{Choose}\\x\end{gathered}
\longrightarrow
\begin{gathered}
\text{Use a function}\\f\text{ to find }y\\\bigl(y=f(x)\bigr)
\end{gathered}
\longrightarrow
\begin{gathered}\text{Plot point}\\(x,y)\end{gathered}
\]

In the rectangular coordinate system, the \textbf{rectangular equation} $y=f(x)$ works well for some shapes like a parabola with a vertical axis of symmetry, but in Precalculus and the review of conic sections in \autoref{sec:conic_sections}, we encountered several shapes that could not be sketched in this manner. (To plot an ellipse using the above procedure, we need to plot the ``top'' and ``bottom'' separately.)\index{curve!rectangular equation}

In this section we introduce a new sketching procedure:

\ifbool{latexml}{
\[
\begin{gathered}\text{Choose}\\t\end{gathered}\quad
\begin{gathered}
\rotatebox{-20}{\scalebox{2}{$\nearrow$}}\\
\rotatebox{20}{\scalebox{2}{$\searrow$}}
\end{gathered}\quad
{\addtolength{\jot}{-.5ex}
\begin{gathered}
\text{Use a function}\\f\text{ to find }x\\\bigl(x=f(t)\bigr)\\~\\
\text{Use a function}\\g\text{ to find }y\\\bigl(y=g(t)\bigr)
\end{gathered}
}\quad
\begin{gathered}
\rotatebox{20}{\scalebox{2}{$\searrow$}}\\
\rotatebox{-20}{\scalebox{2}{$\nearrow$}}
\end{gathered}\quad
\begin{gathered}\text{Plot point}\\(x,y)\end{gathered}
\]
}{
\begin{center}
\begin{tikzpicture}[alt={Pick t to compute x = f(t) and y = g(t) to plot (x,y).},>=latex]
\draw (0,0) node (A) [align=center] {Choose \\$t$} 
      (3,1) node[align=center] (B1) {Use a function\\ $f$ to find $x$\\$\bigl(x=f(t)\bigr)$}
			(3,-1) node[align=center] (B2) {Use a function\\ $g$ to find $y$\\$\bigl(y=g(t)\bigr)$}
			(6.25,0) node [align=center] (C) {Plot point \\ $(x,y)$};
\draw [->](A) --(B1);
\draw [->](A) --(B2);
\draw [->](B1) -- (C);
\draw [->](B2) -- (C);
\end{tikzpicture}
\end{center}
}

Here, $x$ and $y$ are found separately but then plotted together. This leads us to a definition.

\begin{definition}[Parametric Equations and Curves]\label{def:param_eq}%%
Let $f$ and $g$ be continuous functions on an interval $I$. The \textbf{graph} of the \textbf{parametric equations} $x=f(t)$ and $y=g(t)$ is the set of all points $\bigl(x,y\bigr) = \bigl(f(t),g(t)\bigr)$ in the Cartesian plane, as the \textbf{parameter} $t$ varies over $I$. A \textbf{curve} is a graph along with the parametric equations that define it.
\index{curve!parametrically defined}\index{parametric equations!definition}
\end{definition}

\youtubeVideo{9kKZHQtYP7g}{Parametric Equations --- Some basic questions}

This is a formal definition of the word \emph{curve}. When a curve lies in a plane (such as the Cartesian plane), it is often referred to as a \textbf{plane curve}. Examples will help us understand the concepts introduced in the definition.

\mtable{A table of values of the parametric equations in \autoref{ex_pareq1} along with a sketch of their graph.}{fig:pareq1}{%
\tagpdfsetup{table/header-rows={1}}
\begin{tabular}{r @{\hspace{2em}} rr}
$t$ & $x$ & $y$ \\ \midrule
$-2$ & \phantom{$-$}4 & $-1$ \\
$-1$ & 1 & 0 \\
0 & 0 & 1 \\
1 & 1 & 2 \\
2 & 4 & 3
\end{tabular}\\(a)\\
\begin{tikzpicture}[alt={Blue parametric curve x=t²−4, y=t with arrows showing direction as −2≤t≤2; key points labelled when t is integral.}]
\begin{axis}[width=1.16\marginparwidth,tick label style={font=\scriptsize},
axis y line=middle,axis x line=middle,name=myplot,
ymin=-2.2,ymax=4.2,xmin=-.9,xmax=5.2]
\addplot [thick,draw={\colorone}, smooth,domain=-2:2,] ({x^2},{x+1});
\draw [->,>=latex] (axis cs:3.9204,2.98) -- (axis cs:3.96,2.99);
\draw [->,>=latex] (axis cs:2.25,-.5) -- (axis cs:2.22,-.49);
\filldraw (axis cs:4,-1) circle (1.5pt) node [below] {\scriptsize $t=-2$};
\filldraw (axis cs:1,0) circle (1.5pt) node [shift={(0pt,-10pt)}]
 {\scriptsize $t=-1$};
\filldraw (axis cs:0,1) circle (1.5pt) node [left] {\scriptsize $t=0$};
\filldraw (axis cs:1,2) circle (1.5pt) node [above] {\scriptsize $t=1$};
\filldraw (axis cs:4,3) circle (1.5pt) node [above ] {\scriptsize $t=2$};
\end{axis}
\node [right] at (myplot.right of origin) {\scriptsize $x$};
\node [above] at (myplot.above origin) {\scriptsize $y$};
\end{tikzpicture}
\\(b)}

\begin{example}[Plotting parametric functions]\label{ex_pareq1}%%
Plot the graph of the parametric equations $x=t^2$, $y=t+1$ for $t$ in $[-2,2]$.
\solution
We plot the graphs of parametric equations in much the same manner as we plotted graphs of functions like $y=f(x)$: we make a table of values, plot points, then connect these points with a ``reasonable'' looking curve. \autoref{fig:pareq1}(a) shows such a table of values; note how we have 3 columns.

The points $(x,y)$ from the table are plotted in \autoref{fig:pareq1}(b). The points have been connected with a smooth curve. Each point has been labeled with its corresponding $t$-value. These values, along with the two arrows along the curve, are used to indicate the \textbf{orientation} of the graph. This information describes the \textbf{path} of a particle traveling along the curve.
\end{example}

We often use the letter $t$ as the parameter as we often regard $t$ as representing \emph{time}. Certainly there are many contexts in which the parameter is not time, but it can be helpful to think in terms of time as one makes sense of parametric plots and their orientation (for instance, ``At time $t=0$ the position is $(1,2)$ and at time $t=3$ the position is $(5,1)$.'').

\begin{example}[Plotting parametric functions]\label{ex_pareq2}%
%
\mtable{A table of values of the parametric equations in \autoref{ex_pareq2} along with a sketch of their graph.}{fig:pareq2}{%
\tagpdfsetup{table/header-rows={1}}
\begin{tabular}{c @{\hspace{2em}} cc}
$t$ & $x$ & $y$ \\ \hline
0 & 1 & 2 \\
$\pi/4$ & $1/2$ & $1+\sqrt{2}/2$ \\
$\pi/2$ & 0 & 1\\
$3\pi/4$ & $1/2$ & $1-\sqrt{2}/2$ \\
$\pi$ & 1 & 0
\end{tabular}
\\(a)\\
\begin{tikzpicture}[alt={Blue curve sweeps downward on 0.5 ≤ x ≤ 1.5; dots mark t = 0, π/4, π/2, 3π/4.}]
\begin{axis}[width=1.16\marginparwidth,tick label style={font=\scriptsize},
axis y line=middle,axis x line=middle,name=myplot,
ymin=-.1,ymax=2.1,xmin=-.1,xmax=1.6]
\addplot [thick,draw={\colorone}, smooth,domain=0:180,] ({(cos(x))^2},{cos(x)+1});
\draw [->,>=latex] (axis cs:0.998271, 0.00086485) -- (axis cs:0.999534, 0.000233112);
\draw [->,>=latex] (axis cs:0.770151, 1.87758) -- (axis cs:0.75311, 1.86782);
\filldraw (axis cs:1,2) circle (1.5pt) node [right] {\scriptsize $t=0$};
\filldraw (axis cs:.5,1.71) circle (1.5pt) node [below right]
 {\scriptsize $t=\pi/4$};
\filldraw (axis cs:0,1) circle (1.5pt) node [right] {\scriptsize $t=\pi/2$};
\filldraw (axis cs:.5,.292) circle (1.5pt) node [above right]
 {\scriptsize $t=3\pi/4$};
\filldraw (axis cs:1,0) circle (1.5pt) node [above right] {\scriptsize $t=\pi$};
\end{axis}
\node [right] at (myplot.right of origin) {\scriptsize $x$};
\node [above] at (myplot.above origin) {\scriptsize $y$};
\end{tikzpicture}
\\(b)}
%
Sketch the graph of the parametric equations $x=\cos^2t$, $y=\cos t+1$ for $t$ in $[0,\pi]$.
\solution
We again start by making a table of values in \autoref{fig:pareq2}(a), then plot the points $(x,y)$ on the Cartesian plane in \autoref{fig:pareq2}(b).

The curves in Examples \ref{ex_pareq1} and \ref{ex_pareq2} are portions of the same parabola $(y-1)^2+x=1$. While the \emph{parabola} is the same, the \emph{curves} are different. In \autoref{ex_pareq1}, if we let $t$ vary over all real numbers, we'd obtain the entire parabola. In this example, letting $t$ vary over all real numbers would still produce the same graph; this portion of the parabola would be traced, and re-traced, infinitely often. The orientation shown in \autoref{fig:pareq2} shows the orientation on $[0,\pi]$, but this orientation is reversed on $[\pi,2\pi]$.
\end{example}

\subsection{Converting between rectangular and parametric equations}

It is sometimes useful to transform rectangular form equations (i.e., $y=f(x)$) into parametric form equations, and vice-versa. Converting from rectangular to parametric can be very simple: given $y=f(x)$, the parametric equations $x=t$, $y=f(t)$ produce the same graph. As an example, given $y=x^2-x-6$, the parametric equations $x=t, y=t^2-t-6$ produce the same parabola. However, other parameterizations can be used. The following example demonstrates one possible alternative.

\mtable{The equation $f(x)=x^2-x-6$ with different parameterizations.}{fig_parab_param}{\begin{tikzpicture}[alt={Blue parabola segment y=x²−x−6, traced left-to-right as −3 ≤ t ≤ 3 with points marked at t = −3, −1, 1, 3.}]
 \begin{axis}[width=1.16\marginparwidth,tick label style={font=\scriptsize},
   axis y line=middle,axis x line=middle,name=myplot,
   ymin=-7.2,ymax=5.5,xmin=-3.75,xmax=4.5]
  \addplot [thick,draw={\colorone}, smooth,domain=-1.5:4.5,] ({x-1},{x^2-3*x-4});
  \draw [->,>=latex] (axis cs:-0.82, -4.51) -- (axis cs:-0.813, -4.53);
  \draw [->,>=latex] (axis cs:2.42, -2.57) -- (axis cs:2.55, -2.13);
  \filldraw (axis cs:-2,0) circle (1.5pt) node [above left] {\scriptsize $t=-1$};
  \filldraw (axis cs:0,-6) circle (1.5pt) node [below left] {\scriptsize $t=1$};
  \filldraw (axis cs:2,-4) circle (1.5pt) node [right] {\scriptsize $t=3$};
  \filldraw (axis cs:3,0) circle (1.5pt) node [above right] {\scriptsize $t=4$};
 \end{axis}
 \node [right] at (myplot.right of origin) {\scriptsize $x$};
 \node [above] at (myplot.above origin) {\scriptsize $y$};
\end{tikzpicture}
\\(a)\\
\begin{tikzpicture}[alt={Blue parabola, now traced right-to-left as 3 ≥ t ≥ −3 (i.e. −3 ≤ t ≤ 3 but reverse direction); labelled t = 3, 1, −1, −3.}]
 \begin{axis}[width=1.16\marginparwidth,tick label style={font=\scriptsize},
   axis y line=middle,axis x line=middle,name=myplot,
   ymin=-7.2,ymax=5.5,xmin=-3.75,xmax=4.5]
  \addplot [thick,draw={\colorone}, smooth,domain=-.5:5.5,] ({3-x},{x^2-5*x});
  \draw [->,>=latex] (axis cs:-1.20, -3.43) -- (axis cs:-1.21, -3.4);
  \draw [->,>=latex] (axis cs:2.55, -2.13) -- (axis cs:2.42, -2.57);
  \filldraw (axis cs:-2,0) circle (1.5pt) node [above left] {\scriptsize $t=5$};
  \filldraw (axis cs:0,-6) circle (1.5pt) node [below left] {\scriptsize $t=3$};
  \filldraw (axis cs:2,-4) circle (1.5pt) node [right] {\scriptsize $t=1$};
  \filldraw (axis cs:3,0) circle (1.5pt) node [above right] {\scriptsize $t=0$};
 \end{axis}
 \node [right] at (myplot.right of origin) {\scriptsize $x$};
 \node [above] at (myplot.above origin) {\scriptsize $y$};
\end{tikzpicture}
\\(b)\\
\begin{tikzpicture}[alt={Blue parabola, traced left-to-right on −3 ≤ t ≤ 3 but twice as fast (spacing of t-labels shows speed).}]
 \begin{axis}[width=1.16\marginparwidth,tick label style={font=\scriptsize},
   axis y line=middle,axis x line=middle,name=myplot,
   ymin=-7.2,ymax=5.5,xmin=-3.75,xmax=4.5]
  \addplot [thick,draw={\colorone}, smooth,domain=-5.75:5.75,] ({(x+1)/2},{.25*x^2-6.25});
  \draw [->,>=latex] (axis cs:-0.82, -4.51) -- (axis cs:-0.813, -4.53);
  \draw [->,>=latex] (axis cs:2.42, -2.57) -- (axis cs:2.55, -2.13);
  \filldraw (axis cs:-2,0) circle (1.5pt) node [above left] {\scriptsize $t=-5$};
  \filldraw (axis cs:0,-6) circle (1.5pt) node [below left] {\scriptsize $t=-1$};
  \filldraw (axis cs:2,-4) circle (1.5pt) node [right] {\scriptsize $t=3$};
  \filldraw (axis cs:3,0) circle (1.5pt) node [above right] {\scriptsize $t=5$};
 \end{axis}
 \node [right] at (myplot.right of origin) {\scriptsize $x$};
 \node [above] at (myplot.above origin) {\scriptsize $y$};
\end{tikzpicture}
\\(c)}

\begin{example}[Converting from rectangular to parametric]\label{ex_pareq5}%
Find parametric equations for $f(x)=x^2-x-6$.
\solution
Solution 1: For any choice for $x$ we can determine the corresponding $y$ by substitution. If we choose $x=t-1$ then $y=(t-1)^2-(t-1)-6=t^2-3t-4$. Thus $f(x)$ can be represented by the parametric equations
\[x=t-1 \qquad y=t^2-3t-4.\]
On the graph of this parameterization (\autoref{fig_parab_param}(a)) the points have been labeled with the corresponding $t$-values and arrows indicate the path of a particle traveling on this curve. The particle would move from the upper left, down to the vertex at $(.5,-6.25)$ and then up to the right.

Solution 2: If we choose $x=3-t$ then $y=(3-t)^2-(3-t)-6=t^2-5t$. Thus $f(x)$ can also be represented by the parametric equations
\[x=3-t \qquad y=t^2-5t.\]
On the graph of this parameterization (\autoref{fig_parab_param}(b)) the points have been labeled with the corresponding $t-$values and arrows indicate the path of a particle traveling on this curve. The particle would move down from the upper right, to the vertex at $(.5,-6.25)$ and then up to the left.

Solution 3: We can also parameterize any $y=f(x)$ by setting $t=\frac{\dd y}{\dd x}$. That is, $t=a$ corresponds to the point on the graph whose tangent line has a slope $a$. Computing $\frac{\dd y}{\dd x}$, $\fp(x) = 2x-1$ we set $t=2x-1$. Solving for $x$ we find $x=\frac{t+1}{2}$ and by substitution $y=\frac{1}{4}t^2 - \frac{25}{4}$. Thus $f(x)$ can be represented by the parametric equations
\[x=\frac{t+1}{2} \qquad y=\frac{1}{4}t^2 - \frac{25}{4}.\]

The graph of this parameterization is shown in \autoref{fig_parab_param}(c). To find the point where the tangent line has a slope of $0$, we set $t=0$. This gives us the point $(.5, -6.25)$ which is the vertex of $f(x)$.
\end{example}

\begin{example}[Converting from rectangular to parametric]\label{ex_par_circle}%
Find parametric equations for the circle $x^2+y^2=4$.
\solution
We will present three different approaches:

\textbf{Solution 1:} Consider the equivalent equation $\biggl(\dfrac{x}{2}\biggr)^2+ \biggl(\dfrac{y}{2}\biggr)^2=1$ and the Pythagorean Identity, $\sin^2t+\cos^2 t=1$. We set $\cos t=\frac{x}{2}$ and $\sin t=\frac{y}{2}$, which gives $x=2\cos t$ and $y=2\sin t$. To trace the circle once, we must have $0\leq t \leq 2\pi$. Note that when $t=0$ a particle tracing the curve would be at the point $(2,0)$ and would move in a counterclockwise direction.

\textbf{Solution 2:} Another parameterization of the same circle would be $x=2\sin t$ and $y=2\cos t$ for $0\leq t \leq 2\pi$. When $t=0$ a particle would be at the point $(0,2)$ and would move in a clockwise direction.

\textbf{Solution 3:} We could let $x=-2\sin t$ and $y=2\cos t$ for $0\leq t \leq 2\pi$. Also note that we could use $x=2\cos 2t$ and $y=2\sin 2t$ for $0\leq t \leq \pi$.
\end{example}

As we have shown in the previous examples, there are many different ways to parameterize any given curve. We sometimes choose the parameter to accurately model physical behavior.

\begin{example}[Converting from rectangular to parametric]\label{ex_para_ellipse}%
Find a parameterization that traces the ellipse $\dfrac{(x-2)^2}{9}+\dfrac{(y+3)^2}{4}=1$ starting at the point $(-1,-3)$ in a clockwise direction.
\solution
Applying the Pythagorean Identity, $\cos^2t+\sin^2t=1$, we set $\cos^2 t=\dfrac{(x-2)^2}{9}$ and $\sin^2 t=\dfrac{(y+3)^2}{4}$. Solving these equations for $x$ and $y$ we set $x=-3\cos t+2$ and $y=2\sin t-3$  for $0\leq t\leq 2\pi$.
\end{example}

\begin{example}[Converting from rectangular to parametric]\label{ex_para_hyper}%
Find a parameterization for the hyperbola $\dfrac{(x-2)^2}{9}-\dfrac{(y-3)^2}{4}=1$.
\solution
We use the form of the Pythagorean Identity $\sec^2t-\tan^2t=1$. We let  $\ds \sec^2 t=\frac{(x-2)^2}{9}$ and $\tan^2 t=\dfrac{(y-3)^2}{4}$. Solving these equations for $x$ and $y$ we have $x=3\sec t +2$ and $y=2\tan t +3$ for $0\leq t\leq 2\pi$ and $t\neq\frac\pi2,\frac{3\pi}2$.
\end{example}

\begin{example}[Converting from rectangular to parametric]\label{ex_pareq6}%
An object is fired from a height of 0ft and lands 6 seconds later, 192ft away. Assuming ideal projectile motion, the height, in feet, of the object can be described by $h(x) = -x^2/64+3x$, where $x$ is the distance in feet from the initial location. (Thus $h(0) = h(192) = 0$ft.) Find parametric equations $x=f(t)$, $y=g(t)$ for the path of the projectile where $x$ is the horizontal distance the object has traveled at time $t$ (in seconds) and $y$ is the height at time $t$.
\solution
Physics tells us that the horizontal motion of the projectile is linear; that is, the horizontal speed of the projectile is constant. Since the object travels 192ft in 6s, we deduce that the object is moving horizontally at a rate of 32ft/s, giving the equation $x=32t$. As $y=-x^2/64+3x$, we find $y= -16t^2+96t$. We can quickly verify that $y\primeskip''=-32$ft/s$^2$, the acceleration due to gravity, and that the projectile reaches its maximum at $t=3$, halfway along its path.

\mtable{Graphing projectile motion in \autoref{ex_pareq6}.}{fig:pareq6}{\begin{tikzpicture}[alt={Blue projectile arc from (0, 0) to ground at x=192; points marked at t=1 and peak t=2; param laws x=32t, y=–16t²+96t for 0 ≤ t ≤ 6.}]
\begin{axis}[width=\marginparwidth,tick label style={font=\scriptsize},
axis y line=middle,axis x line=middle,name=myplot,
ymin=0,ymax=160,xmin=0,xmax=210]
\addplot [thick, draw={\colorone},smooth,domain=0:6] ({32*x},{-16*x^2+96*x});
\draw [->,>=latex] (axis cs:32,80) -- (axis cs:35.2, 86.24);
\filldraw (axis cs:64, 128) circle (1.5pt) node [above left] {\scriptsize $t=2$};
%\draw (axis cs: 100,50) node [align=left]
% {\scriptsize $x=32t$\\ \scriptsize $y=-16t^2+96t$};
\draw (axis cs: 100,50) node [align=left]
 {\scriptsize $\begin{aligned}x&=32t\\y&=-16t^2+96t\end{aligned}$};
\end{axis}
\node [right] at (myplot.right of origin) {\scriptsize $x$};
\node [above] at (myplot.above origin) {\scriptsize $y$};
\end{tikzpicture}}

These parametric equations make certain determinations about the object's location easy: 2 seconds into the flight the object is at the point $\bigl(x(2),y(2)\bigr) = \bigl(64,128\bigr)$. That is, it has traveled horizontally 64ft and is at a height of 128ft, as shown in \autoref{fig:pareq6}.
\end{example}

It is  sometimes necessary to convert given parametric equations into rectangular form. This can be decidedly more difficult, as some ``simple'' looking parametric equations can have very ``complicated'' rectangular equations. This conversion is often referred to as ``eliminating the parameter,'' as we are looking for a relationship between $x$ and $y$ that does not involve the parameter $t$.

\begin{example}[Eliminating the parameter]\label{ex_pareq7}%
Find a rectangular equation for the curve described by
\[x= \frac{1}{t^2+1}\quad \text{and}\quad y=\frac{t^2}{t^2+1}.\]
\solution
There is not a set way to eliminate a parameter. One method is to solve for $t$ in one equation and then substitute that value in the second. We use that technique here, then show a second, simpler method.

Starting with $x= 1/(t^2+1)$, solve for $t$: $ t = \pm\sqrt{1/x-1}$. Substitute this value for $t$ in the equation for $y$:
{\allowdisplaybreaks
\begin{align*}
 y &= \frac{t^2}{t^2 +1} \\
		&= \frac{1/x-1}{1/x-1+1} \\
		&= \frac{1/x - 1}{1/x} \\
		&= \left(\frac1x-1\right)\cdot x \\
		&= 1-x.
\end{align*}}

\mtable{Graphing parametric and rectangular equations for a graph in \autoref{ex_pareq7}.}{fig:pareq7}{\begin{tikzpicture}[alt={Blue line y=1−x; thick segment from (1, 0) to (0, 1) traced by x = 1/(t²+1), y=t²/(t²+1) on −1≤t≤1.}]
\begin{axis}[width=1.16\marginparwidth,tick label style={font=\scriptsize},
axis y line=middle,axis x line=middle,name=myplot,
ymin=-1.1,ymax=2.1,xmin=-2.1,xmax=2.1,axis equal]
\addplot [ draw={\colorone},smooth,domain=-2:2] {1-x};
\addplot [draw={\colorone},ultra thick, smooth,domain=0:1] {1-x};
\draw [fill=white] (axis cs:0,1) circle (1.5pt);
\filldraw [fill=black] (axis cs:1,0) circle (1.5pt);
\draw (axis cs:1,0) node [align=left,shift={(5pt,40pt)}]
 {\scriptsize $x=\dfrac{1}{t^2+1}$};
\draw (axis cs:1,0) node [align=left,shift={(5pt,20pt)}]
 {\scriptsize $y=\dfrac{t^2}{t^2+1}$};
%\filldraw [fill=black] (axis cs:1,0) circle (1.5pt)
% node [align=left,shift={(5pt,30pt)}]
% {\scriptsize $x=\dfrac{1}{t^2+1}$\\[2pt] \scriptsize $y=\dfrac{t^2}{t^2+1}$};
\draw (axis cs: -1,1.4) node {\scriptsize $y=1-x$};
\end{axis}
\node [right] at (myplot.right of origin) {\scriptsize $x$};
\node [above] at (myplot.above origin) {\scriptsize $y$};
\end{tikzpicture}}

Thus $y=1-x$. One may have recognized this earlier by manipulating the equation for $y$:
\[y = \frac{t^2}{t^2+1} = 1-\frac{1}{t^2+1} = 1-x.\]
This is a shortcut that is very specific to this problem; sometimes shortcuts exist and are worth looking for.

We should be careful to limit the domain of the function $y=1-x$. The parametric equations limit $x$ to values in $(0,1]$, thus to produce the same graph we should limit the domain of $y=1-x$ to the same. 

The graphs of these functions are given in \autoref{fig:pareq7}. The portion of the graph defined by the parametric equations is given in a thick line; the graph defined by $y=1-x$ with unrestricted domain is given in a thin line.
\end{example}

\begin{example}[Eliminating the parameter]\label{ex_pareq8}%
Eliminate the parameter in $x=4\cos t+3$, $y= 2\sin t+1$
\solution
We should not try to solve for $t$ in this situation as the resulting algebra/trig would be messy. Rather, we solve for $\cos t$ and $\sin t$ in each equation, respectively. This gives
\[\cos t = \frac{x-3}{4} \quad \text{and}\quad \sin t=\frac{y-1}{2}.\]
The Pythagorean Theorem gives $\cos^2t+\sin^2t=1$, so:
%
\mtable{Graphing the parametric equations $x=4\cos t+3$, $y=2\sin t+1$ in \autoref{ex_pareq8}.}{fig:pareq8}{\begin{tikzpicture}[alt={Counter-clockwise ellipse centred (3, 1) with semi-axes 4 and 2; path x = 4 cos t + 3, y = 2 sin t + 1 for 0 ≤ t < 2pi.}]
\begin{axis}[width=1.16\marginparwidth,tick label style={font=\scriptsize},
axis y line=middle,axis x line=middle,name=myplot,minor x tick num=1,
minor y tick num=1,ymin=-3.1,ymax=5.1,xmin=-1.1,xmax=8.5]
\addplot [thick,draw={\colorone},smooth,domain=0:360,samples=60]
 ({4*cos(x)+3},{2*sin(x)+1});
\end{axis}
\node [right] at (myplot.right of origin) {\scriptsize $x$};
\node [above] at (myplot.above origin) {\scriptsize $y$};
\end{tikzpicture}}
%
\begin{align*}
\cos^2t+\sin^2t &=1 \\
\left(\frac{x-3}{4}\right)^2 +\left(\frac{y-1}{2}\right)^2 &=1\\
\frac{(x-3)^2}{16}+\frac{(y-1)^2}{4} &=1
\end{align*}
This final equation should look familiar --- it is the equation of an ellipse. \autoref{fig:pareq8} plots the parametric equations, demonstrating that the graph is indeed of an ellipse with a horizontal major axis and center at $(3,1)$.
\end{example}

%The Pythagorean Theorem can also be used to identify parametric equations for hyperbolas. We give the parametric equations for ellipses and hyperbolas in the following Key Idea.

%\begin{keyidea}[Parametric Equations for Ellipses]\label{idea:par_ellipse}%
%The parametric equations 
%\[x=a\cos t+h, \quad y=b\sin t+k\]
%define an ellipse with horizontal axis of length $2a$ and vertical axis of length $2b$, centered at $(h,k)$.
%\index{ellipse!parametric equations}%\\
%\end{keyidea}
%
%\begin{keyidea}[Parametric Equations for Hyperbolas]\label{idea:par_hyperbola}%
%The parametric equations
%\[x= a\tan t+h,\quad y=\pm b\sec t+k\]
%define a hyperbola with vertical transverse axis centered at $(h,k)$, and 
%\[x=\pm a\sec t+h, \quad y=b\tan t + k\]
%defines a hyperbola with horizontal transverse axis. Each has asymptotes at $y=\pm b/a(x-h)+k$.
%\index{hyperbola!parametric equations}
%\end{keyidea}

%\begin{keyidea}[Parametric Equations of Ellipses and Hyperbolas]\label{idea:par_ellipse_hyperbola}%
%\begin{itemize}
%	\item The parametric equations 
%\[x=a\cos t+h, \quad y=b\sin t+k\]
%define an ellipse with horizontal axis of length $2a$ and vertical axis of length $2b$, centered at $(h,k)$.
%\index{ellipse!parametric equations}
%	\item The parametric equations
%\[x= a\tan t+h,\quad y=\pm b\sec t+k\]
%define a hyperbola with vertical transverse axis centered at $(h,k)$, and 
%\[x=\pm a\sec t+h, \quad y=b\tan t + k\]
%defines a hyperbola with horizontal transverse axis. Each has asymptotes at $y=\pm b/a(x-h)+k$.
%\index{hyperbola!parametric equations}
%\end{itemize}
%\end{keyidea}
%%%%% End with this??
%%%%%

\subsection{Graphs of Parametric Equations}

These examples begin to illustrate the powerful nature of parametric equations. Their graphs are far more diverse than the graphs of functions produced by ``$y=f(x)$'' functions.

%\textbf{Technology Note:} Most graphing utilities can graph functions given in parametric form. Often the word ``parametric'' is abbreviated as ``PAR'' or ``PARAM'' in the  options. The user usually needs to determine the graphing window (i.e, the minimum and maximum $x$- and $y$-values), along with the values of $t$ that are to be plotted. The user is often prompted to give a $t$ minimum, a $t$ maximum, and a ``$t$-step'' or ``$\Delta t$.'' Graphing utilities effectively plot parametric functions just as we've shown here: they plots lots of points. A smaller $t$-step plots more points, making for a smoother graph (but may take longer). In \autoref{fig:pareq1}, the $t$-step is 1; in \autoref{fig:pareq2}, the $t$-step is $\pi/4$.\\

One nice feature of parametric equations is that their graphs are easy to shift. While this is not too difficult in the ``$y=f(x)$'' context, the resulting function can look rather messy. (Plus, to shift to the right by two, we replace $x$ with $x-2$, which is counterintuitive.) The following example demonstrates this.

\begin{example}[Shifting the graph of parametric functions]\label{ex_pareq3}%
Sketch the graph of the parametric equations $x=t^2+t$, $y=t^2-t$.  Find new parametric equations that shift this graph to the right 3 units and down 2.
%
\mtable{Illustrating how to shift graphs in \autoref{ex_pareq3}.}{fig:pareq3}{%
\begin{tikzpicture}[alt={Blue curve x=t²+t, y=t²−t on −3≤t≤3; arrow shows motion toward larger x,y, starting near (0, 0).}]
\begin{axis}[width=1.16\marginparwidth,tick label style={font=\scriptsize},
axis y line=middle,axis x line=middle,name=myplot,xtick={2,4,6,8,10},
ymin=-2.6,ymax=6.5,xmin=-.6,xmax=10.5]
\addplot [thick,draw={\colorone}, smooth,domain=-3:3,] ({x^2+x},{x^2-x});
\draw [->,>=latex] (axis cs:0.96, 4.16) -- (axis cs:0.9381, 4.1181);
\draw [->,>=latex] (axis cs:6,2) -- (axis cs:6.05,2.03);
\draw (axis cs: 4,5) node [align=left]
 {\scriptsize $\begin{aligned}x&=t^2+t\\y&=t^2-t\end{aligned}$};
% {\scriptsize $x=t^2+t$\\ \scriptsize $y=t^2-t$};
\end{axis}
\node [right] at (myplot.right of origin) {\scriptsize $x$};
\node [above] at (myplot.above origin) {\scriptsize $y$};
\end{tikzpicture}
\\(a) \\
\begin{tikzpicture}[alt={Blue curve x=t²+t, y=t²−t on −3≤t≤3; shifted right 3 and down 2; equations x=t²+t+3, y=t²−t−2.}]
\begin{axis}[width=1.16\marginparwidth,tick label style={font=\scriptsize},
axis y line=middle,axis x line=middle,name=myplot,xtick={2,4,6,8,10},
ymin=-2.6,ymax=6.5,xmin=-.6,xmax=10.5]
\addplot [thick,draw={\colorone}, smooth,domain=-3:3,] ({x^2+x+3},{x^2-x-2});
\draw [->,>=latex] (axis cs:3.96, 2.16) -- (axis cs:3.9381, 2.1181);
\draw [->,>=latex] (axis cs:9,0) -- (axis cs:9.05,0.03);
\draw (axis cs: 2.5,5) node [align=left]
 {\scriptsize $\begin{aligned}x&=t^2+t+3\\y&=t^2-t-2\end{aligned}$};
% {\scriptsize $x=t^2+t+3$\\ \scriptsize $y=t^2-t-2$};
\end{axis}
\node [right] at (myplot.right of origin) {\scriptsize $x$};
\node [above] at (myplot.above origin) {\scriptsize $y$};
\end{tikzpicture}
\\(b)}
%
\solution
We see the graph in \autoref{fig:pareq3}(a). It is a parabola with an axis of symmetry along the line $y=x$; the vertex is at $(0,0)$. It should be noted that finding the vertex is not a trivial matter and not something you will be asked to do in this text.

In order to shift the graph to the right 3 units, we need to increase the $x$-value by 3 for every point. The straightforward way to accomplish this is simply to add 3 to the function defining $x$: $x = t^2+t+3$. To shift the graph down by 2 units, we wish to decrease each $y$-value by 2, so we subtract 2 from the function defining $y$: $y = t^2-t-2$. Thus our parametric equations for the shifted graph are $x=t^2+t+3$, $y=t^2-t-2$. This is graphed in \autoref{fig:pareq3} (b). Notice how the vertex is now at $(3,-2)$.
\end{example}

Because the $x$- and $y$-values of a graph are determined independently, the graphs of parametric functions often possess features not seen on ``$y=f(x)$'' type graphs. The next example demonstrates how such graphs can arrive at the same point more than once.

\begin{example}[Graphs that cross themselves]\label{ex_pareq4}%
Plot the parametric functions $x=t^3-5t^2+3t+11$ and $y=t^2-2t+3$ and determine the $t$-values where the graph crosses itself.
\solution
Using the methods developed in this section, we again plot points and graph the parametric equations as shown in \autoref{fig:pareq4}. It appears that the graph crosses itself at the point $(2,6)$, but we'll need to analytically determine this.

We are looking for two different values, say, $s$ and $t$, where $x(s) = x(t)$ and $y(s) = y(t)$. That is, the $x$-values are the same precisely when the $y$-values are the same. This gives us a system of 2 equations with 2 unknowns:
%
\mtable{A graph of the parametric equations from \autoref{ex_pareq4}.}{fig:pareq4}{\begin{tikzpicture}[alt={Looping blue path x = t³ − 5t² + 3 t + 11, y = t² − 2 t + 3 on −2 ≤ t ≤ 6; arrows show forward direction crossing itself near x ≈ 1.}]
\begin{axis}[width=1.16\marginparwidth,tick label style={font=\scriptsize},
axis y line=middle,axis x line=middle,name=myplot,minor x tick num=4,
minor y tick num=4,ymin=-.9,ymax=16,xmin=-6,xmax=16]
\addplot [thick, draw={\colorone},smooth,domain=-2:5,samples=50]
 ({x^3-5*x^2+3*x+11},{x^2-2*x+3});
\draw [->,>=latex] (axis cs:-4.62287, 7.5225) -- (axis cs:-4.4041, 7.4756);
\draw [->,>=latex] (axis cs:7,11) -- (axis cs:7.1107, 11.0601);
\draw (axis cs: 7,14) node [align=left]
 {\scriptsize $\begin{aligned}x&=t^3-5t^2+3t+11\\y&=t^2-2t+3\end{aligned}$};
\end{axis}
\node [right] at (myplot.right of origin) {\scriptsize $x$};
\node [above] at (myplot.above origin) {\scriptsize $y$};
\end{tikzpicture}}
%
\begin{align*}
s^3-5s^2+3s+11 &= t^3-5t^2+3t+11 \\
s^2-2s+3 &= t^2-2t+3
\end{align*}

Solving this system is not trivial but involves only algebra. Using the quadratic formula, one can solve for $t$ in the second equation and find that $\ds t = 1\pm \sqrt{s^2-2s+1}$. This can be substituted into the first equation, revealing that the graph crosses itself at $t=-1$ and $t=3$. We confirm our result by computing $x(-1) = x(3)=2$ and $y(-1) = y(3) = 6$.
\end{example}

We now present a small gallery of ``interesting'' and ``famous'' curves along with parametric equations that produce them.\bigskip\vfill

\noindent\flushinner{%
\newcommand{\scalefactor}{1.16}
\tagpdfsetup{table/header-rows={2,4}}
 \begin{tabular}{ c c c }
\begin{tikzpicture}[alt={Astroid: Symmetric four-cusped star touching axes at ±1; given by x = cos³t and y = sin³t.}]
\begin{axis}[width=\scalefactor\marginparwidth,tick label style={font=\scriptsize},
axis y line=middle,axis x line=middle,name=myplot,xtick={-1,1},ytick={-1,1},
ymin=-1.1,ymax=1.1,xmin=-1.1,xmax=1.1,axis equal]
\addplot [thick,draw={\colorone}, smooth,domain=0:360,samples=360]
 ({(cos(x))^3},{(sin(x))^3});
\end{axis}
\node [right] at (myplot.right of origin) {\scriptsize $x$};
\node [above] at (myplot.above origin) {\scriptsize $y$};
\end{tikzpicture}
  &
  \begin{tikzpicture}[alt={Cycloid – One arch of a rolling-wheel path from (0, 0) to (2π, 0), crest at (π, 2 r); x = r(t–sin t), y = r(1–cos t).}]
   \begin{axis}[width=\scalefactor\marginparwidth,tick label style={font=\scriptsize},
                axis y line=middle,axis x line=middle,name=myplot,
                xtick={3.14,6.28},xticklabels={$\pi$,$2\pi$},
                ymin=-.2,ymax=4.2,xmin=-.1,xmax=6.3,axis equal image]
    \addplot [thick,draw={\colorone}, smooth,domain=0:6.28,samples=20]
     ({x-sin(deg(x))},{1-cos(deg(x))});
   \end{axis}
   \node [right] at (myplot.right of origin) {\scriptsize $x$};
   \node [above] at (myplot.above origin) {\scriptsize $y$};
  \end{tikzpicture} &
  \begin{tikzpicture}[alt={Witch of Agnesi – Smooth bell curve centred at the origin, peak near y = 2a, tails flatten for x > 2.}]
   \begin{axis}[width=\scalefactor\marginparwidth,tick label style={font=\scriptsize},
                axis y line=middle,axis x line=middle,name=myplot,
                %xtick={-1,1},ytick={-1,1},
                ymin=-.1,ymax=4.1,xmin=-3.1,xmax=3.1,axis equal image]
    \addplot [thick,draw={\colorone}, smooth,domain=-1.55:1.55,samples=20]
     ({2*x},{2/(1+x^2)});
   \end{axis}
   \node [right] at (myplot.right of origin) {\scriptsize $x$};
   \node [above] at (myplot.above origin) {\scriptsize $y$};
  \end{tikzpicture} \\
  \parbox{150pt}{\centering Astroid\\$x=\cos^3 t$\\$y=\sin^3t$} &
  \parbox{150pt}{\centering Cycloid\\$x=r(t-\sin t)$\\$y=r(1-\cos t)$} &
  \parbox{150pt}{\centering Witch of Agnesi\\$x=2at$\\$y=2a/(1+t^2)$} \\
  \\
\begin{tikzpicture}[alt={Hypotrochoid: Five pointed star looping around origin inside a box ±5; x = 2 cos t+5 cos (2 t/3), y = 2 sin t–5 sin (2 t/3).}]
\begin{axis}[width=\scalefactor\marginparwidth,tick label style={font=\scriptsize},
axis y line=middle,axis x line=middle,name=myplot,
ymin=-7.1,ymax=7.1,xmin=-7.1,xmax=7.5,axis equal]
\addplot [thick, draw={\colorone},smooth,domain=0:1440,samples=120]
 ({2*cos(x)+5*cos(2*x/3)},{2*sin(x)-5*sin(2*x/3)});
\end{axis}
\node [right] at (myplot.right of origin) {\scriptsize $x$};
\node [above] at (myplot.above origin) {\scriptsize $y$};
\end{tikzpicture}
  &
\begin{tikzpicture}[alt={Epicycloid: Three rounded outward petals meeting at origin, reaching about (±5, 0); x = 4 cos t–cos 4 t, y = 4 sin t–sin 4 t.}]
\begin{axis}[width=\scalefactor\marginparwidth,tick label style={font=\scriptsize},
axis y line=middle,axis x line=middle,name=myplot,
ymin=-7.1,ymax=7.1,xmin=-7.1,xmax=7.5,axis equal]
\addplot [thick,draw={\colorone}, smooth,domain=0:720,samples=240]
 ({4*cos(x)-cos(4*x)},{4*sin(x)-sin(4*x)});
\end{axis}
\node [right] at (myplot.right of origin) {\scriptsize $x$};
\node [above] at (myplot.above origin) {\scriptsize $y$};
\end{tikzpicture}
  &
  \begin{tikzpicture}[alt={Folium of Descartes – Leaf-shaped loop crossing itself at origin then curving rightward; x = 3 a t /(1+t³), y = 3 a t²/(1+t³).}]
   \begin{axis}[width=\scalefactor\marginparwidth,tick label style={font=\scriptsize},
                axis y line=middle,axis x line=middle,name=myplot,
                ymin=-2.1,ymax=2.1,xmin=-2.1,xmax=2.1,axis equal]
    \addplot [thick,draw={\colorone}, smooth,domain=-40:130,samples=20]
     ({3*sin(x)*cos(x)^2/((sin(x))^3+(cos(x))^3)},{3*sin(x)^2*cos(x)/((sin(x))^3+(cos(x))^3)});
   \end{axis}
   \node [right] at (myplot.right of origin) {\scriptsize $x$};
   \node [above] at (myplot.above origin) {\scriptsize $y$};
  \end{tikzpicture} \\
  \parbox{150pt}{\centering Hypotrochoid\\$x=2\cos(t)+5\cos(2t/3)$\\$y=2\sin(t)-5\sin(2t/3)$} &
  \parbox{150pt}{\centering Epicycloid\\$x=4\cos(t)-\cos(4t)$\\$y=4\sin(t)-\sin(4t)$} &
  \parbox{150pt}{\centering Folium of Descartes\\$x=3at/(1+t^3)$\\$y=3at^2/(1+t^3)$}
 \end{tabular}}
\bigskip

One might note a feature shared by three of these graphs: ``sharp corners,'' or \textbf{cusps}. We have seen graphs with cusps before and determined that such functions are not differentiable at these points. This leads us to a definition.

\begin{definition}[Smooth]\label{def:smooth}%
A curve $C$ defined by $x=f(t)$, $y=g(t)$ is \textbf{smooth} on an interval $I$ if $\fp$ and $g\primeskip'$ are continuous on $I$ and not simultaneously 0 (except possibly at the endpoints of $I$). A curve is \textbf{piecewise smooth} on $I$ if $I$ can be partitioned into subintervals where $C$ is smooth on each subinterval.
\index{curve!smooth}\index{smooth curve}\index{cusp}
\end{definition}

Consider the astroid, given by $x=\cos^3t$, $y=\sin^3t$. Taking derivatives, we have:
\[
x\primeskip' = -3\cos^2t\sin t\quad \text{and}\quad y\primeskip' = 3\sin^2t\cos t.
\]
It is clear that each is 0 when $t=0,\ \pi/2,\ \pi,\dotsc$. Thus the astroid is not smooth at these points, corresponding to the cusps seen in the figure. However, by restricting the domain of the astroid to all reals except $t = \frac{k\pi}{2}$ for $k \in\mathbb{Z}$ we have a piecewise smooth curve.

We demonstrate this once more.

\begin{example}[Determine where a curve is not smooth]\label{ex_pareq9}%
Let a curve $C$ be defined by the parametric equations $x=t^3-12t+17$ and $y=t^2-4t+8$. Determine the points, if any, where it is not smooth.
\solution
We begin by taking derivatives.
%
\mtable{Graphing the curve in \autoref{ex_pareq9}; note it is not smooth at $(1,4)$.}{fig:pareq9}{\begin{tikzpicture}[alt={Two thin blue line segments form a slight wedge; join at (1, 4) is a corner where the param curve isn’t smooth.}]
\begin{axis}[width=1.16\marginparwidth,tick label style={font=\scriptsize},
axis y line=middle,axis x line=middle,name=myplot,minor x tick num=4,
minor y tick num=1,ymin=-1,ymax=9,xmin=-1,xmax=11]
\addplot [thick,draw={\colorone}, smooth,domain=-2:3.5,samples=40]
 ({x^3-12*x+17},{x^2-4*x+8)});
\end{axis}
\node [right] at (myplot.right of origin) {\scriptsize $x$};
\node [above] at (myplot.above origin) {\scriptsize $y$};
\end{tikzpicture}}
%
\[x\primeskip' = 3t^2-12,\quad y\primeskip' = 2t-4.\]
We set each equal to 0:
\begin{align*}
 x\primeskip' =0 &\Rightarrow 3t^2-12=0 \Rightarrow t=\pm 2\\
 y\primeskip' =0 &\Rightarrow 2t-4 = 0 \Rightarrow t=2
\end{align*}
We consider only the value of $t=2$ since both $x'$ and $y'$ must be 0. Thus $C$ is not smooth at $t=2$, corresponding to the point $(1,4)$. The curve is graphed in \autoref{fig:pareq9}, illustrating the cusp at $(1,4)$.
\end{example}

If a curve is not smooth at $t=t_0$, it means that $x\primeskip'(t_0)=y\primeskip'(t_0)=0$ as defined. This, in turn, means that rate of change of $x$ (and $y$) is 0; that is, at that instant, neither $x$ nor $y$ is changing. If the parametric equations describe the path of some object, this means the object is at rest at $t_0$. An object at rest can make a ``sharp'' change in direction, whereas moving objects tend to change direction in a ``smooth'' fashion.

% cycloid courtesy of Paul Dawkins
\begin{example}[The Cycloid]\label{ex_cycloid}%
A well-known parametric curve is the cycloid.  Fix $r$, and let $x=r(t-\sin t)$, $y=r(1-\cos t)$.  This represents the path traced out by a point on a wheel of radius $r$ as it starts rolling to the right.  We can think of $t$ as the angle through which the point has rotated.
 
\noindent\begin{minipage}[t]{\linewidth}\noindent%
\captionsetup{type=figure}%
 \centering
  \begin{tikzpicture}[alt={Row of rolling circles with red cycloid track showing two arches and cusps beneath successive wheels.}]
   \begin{axis}[axis lines=none,xmin=-1,xmax=4*pi+1,ymin=-.5,ymax=3,axis equal,
     width=\textwidth,height=\ifbool{latexml}{.3}{}\textwidth]
    \foreach \x in { 0, 2*pi/3, 4*pi/3, 2*pi, 8*pi/3, 10*pi/3, 4*pi} {
     \edef\temp{%
      \noexpand\draw [draw={\colorone}] (axis cs:\x,1) circle (1);
      \noexpand\draw [draw={\colorone},fill={\colorone}] (axis cs:\x,1)circle (1.5pt)
       -- (axis cs:{\x-sin(deg(\x))},{1-cos(deg(\x))}) circle (1.5pt);
     }\temp
    }
    \draw (axis cs:0,2.5) node {$t=0$};
    \draw (axis cs:2*pi,2.5) node {$t=2\pi$};
    \draw (axis cs:4*pi,2.5) node {$t=4\pi$};
    \foreach \t in { 2, 4, 8, 10} {
     \edef\temp{%
      \noexpand\draw (axis cs:\t*pi/3,2.5) node {$t=\frac{\t\pi}3$};
     }\temp
    }
    \addplot[draw={\colortwo},smooth,thick,domain=0:4*pi]
     ({\x-sin(deg(\x))},{1-cos(deg(\x))});
   \end{axis}
  \end{tikzpicture}
 \caption{A cycloid traced through two revolutions.}
 \label{fig_cycloid}
\end{minipage}
 
\autoref{fig_cycloid} shows a cycloid sketched out with the wheel shown at various places.  The dot on the rim is the point on the wheel that we're using to trace out the curve.
 
From this sketch we can see that one arch of the cycloid is traced out in the range $0\le t\le2\pi$.  This makes sense when you consider that the point will be back on the ground after it has rotated through an angle of $2\pi$.
\end{example}

One should be careful to note that a ``sharp corner'' does not have to occur when a curve is not smooth. For instance, one can verify that $x=t^3$, $y=t^6$ produce the familiar $y=x^2$ parabola. However, in this parameterization, the curve is not smooth. A particle traveling along the parabola according to the given parametric equations comes to rest at $t=0$, though no sharp point is created.\bigskip

Our previous experience with cusps taught us that a function was not differentiable at a cusp. This can lead us to wonder about derivatives in the context of parametric equations and the application of other calculus concepts. Given a curve defined parametrically, how do we find the slopes of tangent lines? Can we determine concavity? We explore these concepts and more in the next section.

\printexercises{exercises/09-02-exercises}

\section{Calculus and Parametric Equations}\label{sec:par_calc}

The previous section defined curves based on parametric equations. In this section we'll employ the techniques of calculus to study these curves.

We are still interested in lines tangent to points on a curve. They describe how the $y$-values are changing with respect to the $x$-values, they are useful in making approximations, and they indicate instantaneous direction of travel.

The slope of the tangent line is still $\frac{\dd y}{\dd x}$, and the Chain Rule allows us to calculate this in the context of parametric equations. If $x=f(t)$ and $y=g(t)$, the Chain Rule states that
\[\frac{\dd y}{\dd t} = \frac{\dd y}{\dd x}\cdot\frac{\dd x}{\dd t}.\]
Solving for $\frac{\dd y}{\dd x}$, we get 
\[\frac{\dd y}{\dd x} = \frac{\dd y/\dd t}{\dd x/\dd t} = \frac{g\primeskip'(t)}{\fp(t)},\]
provided that $\fp(t)\neq 0$. This is important so we label it a Key Idea.

\begin{keyidea}[Finding $\frac{\dd y}{\dd x}$ with Parametric Equations.]\label{idea:dydxpar}%
Let $x=f(t)$ and $y=g(t)$, where $f$ and $g$ are differentiable on some open interval $I$ and $\fp(t)\neq 0$ on $I$. Then \index{parametric equations!finding $\frac{\dd y}{\dd x}$}\index{derivative!parametric equations}
\[
\frac{\dd y}{\dd x}=\frac{\dd y/\dd t}{\dd x/\dd t}=\frac{g\primeskip'(t)}{\fp(t)}.
\]
\end{keyidea}

We use this to define the tangent line.

\begin{definition}[Tangent Lines]\label{def:tangent_par}%
Let a curve $C$ be parameterized by $x=f(t)$ and $y=g(t)$, where $f$ and $g$ are differentiable functions on some interval $I$ containing $t=t_0$. The 
\textbf{tangent line} to $C$ at $t=t_0$ is the line
\[\fp(t_0)(y-g(t_0))=g'(t_0)(x-f(t_0)).\]
\index{tangent line}\index{parametric equations!tangent line}
\end{definition}

Notice that the tangent line goes through the point $(f(t_0),g(t_0))$.  It is possible for parametric curves to have horizontal and vertical tangents. As expected a horizontal tangent occurs whenever $\frac{\dd y}{\dd x} = 0$ or when $\frac{\dd y}{\dd t} = 0$ (provided $\frac{\dd x}{\dd t} \neq 0$). Similarly, a vertical tangent occurs whenever $\frac{\dd y}{\dd x}$ is undefined or when $\frac{\dd x}{\dd t} = 0$ (provided $\frac{\dd y}{\dd t} \neq 0$).  When $\frac{\dd y}{\dd x}$ is defined, however, the tangent line has slope $\frac{g'(t_0)}{\fp(t_0)}$.

\begin{definition}[Normal Lines]\label{def:normal_par}%
The \textbf{normal line} to a curve $C$ at a point $P$ is the line through $P$ and perpendicular to the tangent line at $P$. This has equation
\[g'(t_0)(y-g(t_0))=-\fp(t_0)(x-f(t_0)).\]
%For $t=t_0$ the normal line is the line through $(f(t_0), g(t_0))$ with slope $m=-\frac{\fp(t_0)}{g'(t_0)}$, provided $g'(t_0) \neq 0$.
\index{normal line}\index{parametric equations!normal line}
\end{definition}

As with the tangent line we note that it is possible for a normal line to be vertical or horizontal. A horizontal normal line occurs whenever $\frac{\dd y}{\dd x}$ is undefined or when $\frac{\dd x}{\dd t} = 0$ (provided $\frac{\dd y}{\dd t} \neq 0$). Similarly, a vertical normal line occurs whenever $\frac{\dd y}{\dd x} = 0$ or when $\frac{\dd y}{\dd t} = 0$ (provided $\frac{\dd x}{\dd t} \neq 0$). In other words, if the curve $C$ has a vertical tangent  at $(f(t_0), g(t_0))$ the normal line will be horizontal and if the tangent is horizontal the normal line will be a vertical line.

\youtubeVideo{k5QnaGVk1JI}{Derivatives of Parametric Functions}

\begin{example}[Tangent and Normal Lines to Curves]\label{ex_parcalc1}%
Let $x=5t^2-6t+4$ and $y=t^2+6t-1$, and let $C$ be the curve defined by these equations.
\begin{enumerate}
	\item	Find the equations of the tangent and normal lines to $C$ at $t=3$.
	\item	Find where $C$ has vertical and horizontal tangent lines.
\end{enumerate}
\solution
\begin{enumerate}
	\item We start by computing $\fp(t) = 10t-6$ and $g\primeskip'(t) =2t+6$. Thus \[\frac{\dd y}{\dd x} = \frac{2t+6}{10t-6}.\]
	Make note of something that might seem unusual: $\frac{\dd y}{\dd x}$ is a function of $t$, not $x$. Just as points on the curve are found in terms of $t$, so are the slopes of the tangent lines.

	The point on $C$ at $t=3$ is $(31,26)$. The slope of the tangent line is $m=1/2$ and the slope of the normal line is $m=-2$. Thus,
	\begin{itemize}
		\item the equation of the tangent line is $\ds y=\frac12(x-31)+26$, and
		\item	the equation of the normal line is $\ds y=-2(x-31)+26$.
	\end{itemize}
	This is illustrated in \autoref{fig:parcalc1}.

\mtable{Graphing tangent and normal lines in \autoref{ex_parcalc1}.}{fig:parcalc1}{\begin{tikzpicture}[alt={Blue parametric curve with two black dots at left and right; red line touching at left dot, steeper red line through left dot indicates normal.}]
\begin{axis}[width=\marginparwidth,tick label style={font=\scriptsize},
axis y line=middle,axis x line=middle,name=myplot,minor x tick num=1,
minor y tick num=1,ymin=-20,ymax=55,xmin=-5,xmax=90]
\addplot [draw={\colorone},thick, smooth,domain=-3.5:5,samples=30]
 ({5*x^2-6*x+4},{x^2+6*x-1});
\addplot [draw={\colortwo},thick,domain=-5:90] ({x},{.5*(x-31)+26});
\addplot[draw={\colortwo},thick,domain=-20:55] ({x},{-2*(x-31)+26});
\draw [>=latex,->] (axis cs: 12.14,-5.204) -- (axis cs: 12,-5.16);
\filldraw (axis cs: 67,-10) circle (1.5pt)
 (axis cs: 2.2,2.96) circle (1.5pt);
\end{axis}
\node [right] at (myplot.right of origin) {\scriptsize $x$};
\node [above] at (myplot.above origin) {\scriptsize $y$};
\end{tikzpicture}}

	\item	To find where $C$ has a horizontal tangent line, we set $\frac{\dd y}{\dd x}=0$ and solve for $t$. In this case, this amounts to setting $g\primeskip'(t)=0$ and solving for $t$ (and making sure that $\fp(t)\neq 0$). 
\[g\primeskip'(t)=0 \quad \Rightarrow \quad 2t+6=0 \quad \Rightarrow \quad t=-3.\]
	The point on $C$ corresponding to $t=-3$ is $(67,-10)$; the tangent line at that point is horizontal (hence with equation $y=-10$).
		
	To find where $C$ has a vertical tangent line, we find where it has a horizontal normal line, and set $-\frac{\fp(t)}{g\primeskip'(t)}=0$. This amounts to setting $\fp(t)=0$ and solving for $t$ (and making sure that $g\primeskip'(t)\neq 0$). 
	\[\fp(t)=0 \quad \Rightarrow \quad 10t-6=0 \quad \Rightarrow \quad t=0.6.\]
	The point on $C$ corresponding to $t=0.6$ is $(2.2,2.96)$. The tangent line at that point is $x=2.2$.
	
	The points where the tangent lines are vertical and horizontal are indicated on the graph in \autoref{fig:parcalc1}.
\end{enumerate}
\end{example}

\begin{example}[Tangent and Normal Lines to a Circle]\label{ex_parcalc2}%
\mbox{}\\[-2\baselineskip]\parbox[t]{\linewidth}{\begin{enumerate}
	\item	Find where the unit circle, defined by $x=\cos t$ and $y=\sin t$ on $[0,2\pi]$, has vertical and horizontal tangent lines. 
	\item	 Find the equation of the normal line at $t=t_0$.
\end{enumerate}}\vspace{0pt}
\solution
\begin{enumerate}
	\item We compute the derivative following \autoref{idea:dydxpar}:
	\[\frac{\dd y}{\dd x} = \frac{g\primeskip'(t)}{\fp(t)} = -\frac{\cos t}{\sin t}.\]
	The derivative is $0$ when $\cos t= 0$; that is, when $t=\pi/2,\ 3\pi/2$. These are the points $(0,1)$ and $(0,-1)$ on the circle.

	The normal line is horizontal (and hence, the tangent line is vertical) when $\sin t=0$; that is, when $t= 0,\ \pi,\ 2\pi$, corresponding to the points $(-1,0)$ and $(0,1)$ on the circle. These results should make intuitive sense.
	\item	The slope of the normal line at $t=t_0$ is $\ds m=\frac{\sin t_0}{\cos t_0} = \tan t_0$. This normal line goes through the point $(\cos t_0,\sin t_0)$, giving the line
%
\mtable{Illustrating how a circle's normal lines pass through its center.}{fig:parcalc2}{\begin{tikzpicture}[alt={Blue circle centered at origin; red radius-like normal line meets rim at upper right and passes through center.}]
\begin{axis}[width=1.16\marginparwidth,tick label style={font=\scriptsize},
axis y line=middle,axis x line=middle,name=myplot,ytick={-1,1},
ymin=-1.1,ymax=1.1,xmin=-1.3,xmax=1.3,axis equal]
\draw[draw={\colorone},thick,smooth](axis cs:0,0)circle(1);
\addplot [draw={\colortwo},thick,domain=-1:1] {2*x};
\filldraw (axis cs: .447,.894) circle (1.5pt);
\end{axis}
\node [right] at (myplot.right of origin) {\scriptsize $x$};
\node [above] at (myplot.above origin) {\scriptsize $y$};
\end{tikzpicture}}%
%
\begin{align*}
	y &=\frac{\sin t_0}{\cos t_0}(x-\cos t_0) + \sin t_0\\	
	&= (\tan t_0)x,
\end{align*}
as long as $\cos t_0\neq 0$. It is an important fact to recognize that the normal lines to a circle pass through its center, as illustrated in \autoref{fig:parcalc2}. Stated in another way, any line that passes through the center of a circle intersects the circle at right angles.
\end{enumerate}
\end{example}

\begin{example}[Tangent lines when $\frac{\dd y}{\dd x}$ is not defined]\label{ex_parcalc3}%
Find the equation of the tangent line to the astroid $x=\cos^3 t$, $y=\sin^3t$ at $t=0$, shown in \autoref{fig:parcalc3}.
\solution
We start by finding $x\primeskip'(t)$ and $y\primeskip'(t)$:
%
\mtable{A graph of an astroid.}{fig:parcalc3}{\begin{tikzpicture}[alt={Blue four-cusp astroid touching axes at ±1.}]
\begin{axis}[width=1.16\marginparwidth,tick label style={font=\scriptsize},
axis y line=middle,axis x line=middle,name=myplot,xtick={-1,1},ytick={-1,1},
ymin=-1.1,ymax=1.1,xmin=-1.1,xmax=1.1,axis equal]
\addplot [thick,draw={\colorone}, smooth,domain=0:360,samples=360]
 ({(cos(x))^3},{(sin(x))^3});
\end{axis}
\node [right] at (myplot.right of origin) {\scriptsize $x$};
\node [above] at (myplot.above origin) {\scriptsize $y$};
\end{tikzpicture}}
%
\[x\primeskip'(t) = -3\sin t\cos^2t, \qquad y\primeskip'(t) = 3\cos t\sin^2t.\]
Note that both of these are 0 at $t=0$; the curve is not smooth at $t=0$ forming a cusp on the graph. Evaluating $\frac{\dd y}{\dd x}$ at this point returns the indeterminate form of ``0/0''. 
%$\frac{\dd y}{\dd x}$:
%\[\frac{\dd y}{\dd x} = \frac{-3\sin t\cos^2t}{3\cos t\sin^2t} = -\frac{\sin t}{\cos t},\] as long as $\cos t\neq 0$ and $\sin t\neq 0$. When $t=0$, it is tempting to declare that
%\[\frac{\dd y}{\dd x} = -\frac{\sin 0}{\cos 0} = 0,\]
%but this overlooks the fact that we divided earlier with the stipulation that $\sin t\neq 0$. In fact, the graph of the curve has a cusp at $t=0$, as both $x\primeskip'=0$ and $y\primeskip'=0$. 

We can, however, examine the slopes of tangent lines near $t=0$, and take the limit as $t\to 0$. 
\begin{align*}
	\lim_{t\to0} \frac{y\primeskip'(t)}{x\primeskip'(t)}
	&=\lim_{t\to0} \frac{3\cos t\sin^2t}{-3\sin t\cos^2t}
	\quad \text{\small (We can reduce as $t\neq 0$.)}\\
	&= \lim_{t\to0}\left(-\frac{\sin t}{\cos t}\right)\\
	&= 0.
\end{align*}
We have accomplished something significant. When the derivative $\frac{\dd y}{\dd x}$ returns an indeterminate form at $t=t_0$, we can define its value by setting it to be $\ds \lim_{t\to t_0}\frac{\dd y}{\dd x}$, if that limit exists. This allows us to find slopes of tangent lines at cusps, which can be very beneficial. 

We found the slope of the tangent line at $t=0$ to be 0; therefore the tangent line is $y=0$, the $x$-axis.
\end{example}

\subsection{Concavity}

We continue to analyze curves in the plane by considering their concavity; that is, we are interested in $\frac{\dd^2y}{\dd x^2}$, ``the second derivative of $y$ with respect to $x$.'' To find this, we need to find the derivative of $\frac{\dd y}{\dd x}$ with respect to $x$; that is,
\[\frac{\dd^2y}{\dd x^2}=\frac{\dd}{\dd x}\left[\frac{\dd y}{\dd x}\right],\]
but recall that $\frac{\dd y}{\dd x}$ is a function of $t$, not $x$, making this computation not straightforward. \index{concavity}\index{parametric equations!concavity}

To make the upcoming notation a bit simpler, let $h(t) = \frac{\dd y}{\dd x}$. We want $\frac{\dd}{\dd x}[h(t)]$; that is, we want $\frac{\dd h}{\dd x}$. We again appeal to the Chain Rule. Note:
\[
\frac{\dd h}{\dd t} = \frac{\dd h}{\dd x}\cdot\frac{\dd x}{\dd t}
\quad \Rightarrow \quad
\frac{\dd h}{\dd x} = \frac{\dd h/\dd t}{\dd x/\dd t}.
\]

In words, to find $\dfrac{\dd^2y}{\dd x^2}$, we first take the derivative of $\dfrac{\dd y}{\dd x}$ \emph{with respect to $t$}, then divide by $x\primeskip'(t)$. We restate this as a Key Idea.
% should this be in terms of f and g?

\begin{keyidea}[Finding $\frac{\dd^2y}{\dd x^2}$ with Parametric Equations]\label{idea:second_der_par}%
Let $x=f(t)$ and $y=g(t)$ be twice differentiable functions on an open interval $I$, where $\fp(t)\neq 0$ on $I$. Then 
\index{parametric equations!finding $\frac{\dd^2y}{\dd x^2}$}
\[
\frac{\dd^2y}{\dd x^2}
\quad = \quad
\frac{\dfrac{\dd}{\dd t}\left[\dfrac{\dd y}{\dd x}\right]}{\dfrac{\dd x}{\dd t}}
\quad=\quad
\frac{\dfrac{\dd}{\dd t}\left[\dfrac{\dd y}{\dd x}\right]}{\fp(t)}.
\] 
\end{keyidea}

Examples will help us understand this Key Idea.

\begin{example}[Concavity of Plane Curves]\label{ex_parcalc4}%
Let $x=5t^2-6t+4$ and $y=t^2+6t-1$ as in \autoref{ex_parcalc1}. Determine the $t$-intervals on which the graph is concave up/down.
\solution
Concavity is determined by the second derivative of $y$ with respect to $x$, $\frac{\dd^2y}{\dd x^2}$, so we compute that here following \autoref{idea:second_der_par}.

In \autoref{ex_parcalc1}, we found $\ds\frac{\dd y}{\dd x} = \frac{2t+6}{10t-6}$ and $\fp(t) = 10t-6$. So:
\begin{align*}
	\frac{\dd^2y}{\dd x^2}
	&= \frac{\frac{\dd}{\dd t}\left[\frac{2t+6}{10t-6}\right]}{10t-6} \\
	&= \frac{-\frac{72}{(10t-6)^2}}{10t-6}\\
	&= -\frac{72}{(10t-6)^3} \\&= -\frac{9}{(5t-3)^3}
\end{align*}

\mtable{Graphing the parametric equations in \autoref{ex_parcalc4} to demonstrate concavity.}{fig:parcalc4}{\begin{tikzpicture}[alt={Blue curve from origin arcs right; t < 3/5 concave-up, then t > 3/5 concave-down.}]
\begin{axis}[width=\marginparwidth,tick label style={font=\scriptsize},
axis y line=middle,axis x line=middle,name=myplot,minor x tick num=1,
minor y tick num=1,ymin=-20,ymax=55,xmin=-5,xmax=90]
\addplot [draw={\colorone},thick, smooth,domain=-3.5:5,samples=30]
 ({5*x^2-6*x+4},{x^2+6*x-1});
\draw [>=latex,->] (axis cs: 12.14,-5.204) -- (axis cs: 12,-5.16);
\filldraw (axis cs: 2.2,2.96) circle (1.5pt);
\draw (axis cs: 31,26) node [above,rotate=26.56] {\scriptsize $t>3/5$; concave down};
\draw (axis cs: 50.25,-9.75) node [below] {\scriptsize $t<3/5$; concave up};
\end{axis}
\node [right] at (myplot.right of origin) {\scriptsize $x$};
\node [above] at (myplot.above origin) {\scriptsize $y$};
\end{tikzpicture}}

The graph of the parametric functions is concave up when $\frac{\dd^2y}{\dd x^2} > 0$ and concave down when $\frac{\dd^2y}{\dd x^2} <0$. We determine the intervals when the second derivative is greater/less than 0 by first finding when it is 0 or undefined.

As the numerator of $\ds -\frac{9}{(5t-3)^3}$ is never 0, $\frac{\dd^2y}{\dd x^2} \neq 0$ for all $t$. It is undefined when $5t-3=0$; that is, when $t= 3/5$. Following the work established in \autoref{sec:concavity}, we look at values of $t$ greater or less than $3/5$ on a number line:
\begin{center}
\includecodegraphics{figures/matrices/paramcalc}
\end{center}

Reviewing \autoref{ex_parcalc1}, we see that when $t=3/5=0.6$, the graph of the parametric equations has a vertical tangent line. This point is also a point of inflection for the graph, illustrated in \autoref{fig:parcalc4}.
\end{example}

\begin{example}[Concavity of Plane Curves]\label{ex_parcalc5}%
Find the points of inflection of the graph of the parametric equations $x=\sqrt{t}$, $y=\sin t$, for $0\leq t\leq 16$.
\solution
We need to compute $\frac{\dd y}{\dd x}$ and $\frac{\dd^2y}{\dd x^2}$. 
\begin{gather*}
\frac{\dd y}{\dd x} = \frac{y\primeskip'(t)}{x\primeskip'(t)}
= \frac{\cos t}{1/(2\sqrt{t})} = 2\sqrt{t}\cos t.\\
\frac{\dd^2y}{\dd x^2} = \frac{\frac{\dd}{\dd t}\left[\frac{\dd y}{\dd x}\right]}{x\primeskip'(t)} = \frac{\cos t/\sqrt{t}-2\sqrt{t}\sin t}{1/(2\sqrt{t})}=2\cos t-4t\sin t.
\end{gather*}
The possible points of inflection are found by setting $\frac{\dd^2y}{\dd x^2}=0$. This is not trivial, as equations that mix polynomials and trigonometric functions generally do not have ``nice'' solutions. 

\mtable{In (a), a graph of $\frac{\dd^2y}{\dd x^2}$, showing where it is approximately 0. In (b), graph of the parametric equations in \autoref{ex_parcalc5} along with the points of inflection.}{fig:parcalc5}{%
\begin{tikzpicture}[alt={Plot of y double prime versus t: blue wave crosses zero near t about one, five, nine, thirteen.}]
\begin{axis}[width=\marginparwidth,tick label style={font=\scriptsize},
axis y line=middle,axis x line=middle,name=myplot,minor x tick num=4,
minor y tick num=4,ymin=-59,ymax=59,xmin=-.5,xmax=16.9]
\addplot [draw={\colorone},thick, smooth,domain=0:16,samples=60]
 {2*cos(deg(x))-4*x*sin(deg(x))};
\draw  (axis cs: 5,-40) node {\scriptsize $y=2\cos t-4t\sin t$};
\end{axis}
\node [right] at (myplot.right of origin) {\scriptsize $t$};
\node [above] at (myplot.above origin) {\scriptsize $y$};
\end{tikzpicture}
\\(a)\\
\begin{tikzpicture}[alt={Corresponding x–y curve: blue oscillation left to right, black dots mark inflection points near x one, two, three, four.}]
\begin{axis}[width=\marginparwidth,tick label style={font=\scriptsize},
axis y line=middle,axis x line=middle,name=myplot,xtick={3,1,2,4},
ymin=-1.1,ymax=1.1,xmin=-.5,xmax=4.5]
\addplot [draw={\colorone},thick, smooth,domain=0:3,samples=100]
 ({sqrt(x)},{sin(deg(x))});
\addplot [draw={\colorone},thick, smooth,domain=3:18,samples=60]
 ({sqrt(x)},{sin(deg(x))});
\filldraw (axis cs: 0.808252, 0.607787) circle (1.5pt)
					(axis cs: 1.81447, -0.150147) circle (1.5pt)
					(axis cs: 2.52223, 0.0783547) circle (1.5pt)
					(axis cs: 3.07855, -0.0526833) circle (1.5pt)
					(axis cs: 3.55049, 0.0396324) circle (1.5pt)
					(axis cs:3.96733, -0.0317508) circle (1.5pt);
\end{axis}
\node [right] at (myplot.right of origin) {\scriptsize $x$};
\node [above] at (myplot.above origin) {\scriptsize $y$};
\end{tikzpicture}
\\(b)}

In \autoref{fig:parcalc5}(a) we see a plot of the second derivative. It shows that it has zeros at approximately $t=0.5,\ 3.5,\ 6.5,\ 9.5,\ 12.5$ and $16$. These approximations are not very good, made only by looking at the graph. Newton's Method provides more accurate approximations. Accurate to 2 decimal places, we have:
\[t=0.65,\ 3.29,\ 6.36,\ 9.48,\ 12.61\ \text{and}\ 15.74.\]
The corresponding points have been plotted on the graph of the parametric equations in \autoref{fig:parcalc5}(b). Note how most occur near the $x$-axis, but not exactly on the axis.
\end{example}


\subsection{Area with Parametric Equations}

We will now find a formula for determining the area under a parametric curve given by the parametric equations
\[x=f(t)\qquad y=g(t).\]
We will also need to further add in the assumption that the curve is traced out exactly once as t increases from $\alpha$ to $\beta$.
%
%We will do this in much the same way that we found the first derivative.  We will
First, recall how to find the area under $y=F(x)$ on $a\le x\le b$:
\[A=\int_a^b F(x)\dd x.\]
Now think of the parametric equation $x=f(t)$ as a substitution in the integral, assuming that $a=f(\alpha)$ and $b=f(\beta)$ for the purposes of this formula.  (There is actually no reason to assume that this will always be the case and so we'll give a corresponding formula later if it's the opposite case ($b=f(\alpha)$ and $a=f(\beta)$).)

In order to substitute, we'll need $\dd x=\fp(t)\dd t$. Plugging this into the area formula above and making sure to change the limits to their corresponding $t$ values gives us
\[A=\int_\alpha^\beta F(f(t))\fp(t)\dd t.\]
Since we don't know what $F(x)$ is, we'll use the fact that
\[y=F(x)=F(f(t))=g(t)\]
and arrive at the formula that we want.

\begin{keyidea}[Area Under a Parametric Curve]\label{ki_param_area}%
The area under the parametric curve given by $x=f(t)$, $y=g(t)$, for $f(\alpha)=a<x<b=f(\beta)$ is
\[A=\int_\alpha^\beta g(t)\fp(t)\dd t.\]
\end{keyidea}

On the other hand, if we should happen to have $b=f(\alpha)$ and $a=f(\beta)$, then the formula would be
\[A=\int_\beta^\alpha g(t)\fp(t)\dd t.\]

Let's work an example.

\begin{example}[Finding the area under a parametric curve]\label{eg_cycloid_area}%
Determine the area under the cycloid given by the parametric equations
\[x=6(\theta-\sin\theta)\qquad y=6(1-\cos\theta)\qquad 0\le\theta\le2\pi.\]
\solution
First, notice that we've switched the parameter to $\theta$ for this problem.  This is to make sure that we don't get too locked into always having $t$ as the parameter.

Now, we could graph this to verify that the curve is traced out exactly once for the given range if we wanted to.
 
There really isn't too much to this example other than plugging the parametric equations into the formula.  We'll first need the derivative of the parametric equation for $x$ however.
\[\frac{\dd x}{\dd\theta}=6(1-\cos\theta).\]
The area is then
\begin{align*}
A
&=\int_0^{2\pi}36(1-\cos\theta)^2\dd\theta\\
&=36\int_0^{2\pi}1-2\cos\theta+\cos^2\theta\dd\theta\\
&=36\int_0^{2\pi}\frac32-2\cos\theta+\frac12\cos(2\theta)\dd\theta\\
&=36\left[\frac32\theta-2\sin\theta+\frac14\sin(2\theta)\right]_0^{2\pi}\\
&=108\pi.
\end{align*}
\end{example}


\subsection{Arc Length}

We continue our study of the features of the graphs of parametric equations by computing their arc length.

Recall in \autoref{sec:arc_length} we found the arc length of the graph of a function, from $x=a$ to $x=b$, to be
\[L = \int_a^b\sqrt{1+\left(\frac{\dd y}{\dd x}\right)^2}\dd x.\]
We can use this equation and convert it to the parametric equation context. Letting $x=f(t)$ and $y=g(t)$, we know that $\frac{\dd y}{\dd x} = g\primeskip'(t)/\fp(t)$. Suppose that $\fp(t)>0$, and calculate the differential of $x$:
\[\dd x = \fp(t)\dd t \qquad \Rightarrow \qquad \dd t = \frac{1}{\fp(t)}\cdot \dd x.\]
Starting with the arc length formula above, consider:
\begin{align*}
	L &= \int_a^b\sqrt{1+\left(\frac{\dd y}{\dd x}\right)^2}\dd x\\
	&= \int_a^b \sqrt{1+\frac{[g\primeskip'(t)]^2}{[\fp(t)]^2}}\dd x \\
%		\intertext{Factor out the \intertextmath{\fp(t)^2}:}
	&= \int_a^b \sqrt{[\fp(t)]^2+[g\primeskip'(t)]^2}\cdot\underbrace{\frac1{\fp(t)}\dd x}_{=\dd t} \qquad\text{Factor out the $[\fp(t)]^2$} \\
	&= \int_{t_1}^{t_2} \sqrt{[\fp(t)]^2+[g\primeskip'(t)]^2}\dd t.\\
\end{align*}
Note the new bounds (no longer ``$x$'' bounds, but ``$t$'' bounds). They are found by finding $t_1$ and $t_2$ such that $a= f(t_1)$ and $b=f(t_2)$. This formula holds even when $\fp$ isn't positive and we restate it as a theorem.

\begin{theorem}[Arc Length of Parametric Curves]\label{thm:arc_length_parametric}%
Let $x=f(t)$ and $y=g(t)$ be parametric equations with $\fp$ and $g\primeskip'$ continuous on some open interval $I$ containing $t_1$ and $t_2$ on which the graph traces itself only once. The arc length of the graph, from $t=t_1$ to $t=t_2$, is
\index{parametric equations!arc length}\index{arc length}
\[L = \int_{t_1}^{t_2} \sqrt{[\fp(t)]^2+[g\primeskip'(t)]^2}\dd t.\]
\end{theorem}

As before, these integrals are often not easy to compute. We start with a simple example, then give  another where we approximate the solution.

\begin{example}[Arc Length of a Circle]\label{ex_parcalc6}%
Find the arc length of the circle parameterized by $x=3\cos t$, $y=3\sin t$ on $[0,3\pi/2]$. 
\solution
By direct application of \autoref{thm:arc_length_parametric}, we have
\begin{align*}
L &= \int_0^{3\pi/2} \sqrt{(-3\sin t)^2 +(3\cos t)^2} \dd t.\\
\intertext{Apply the Pythagorean Theorem.}
	&= \int_0^{3\pi/2} 3 \dd t\\
	&= 3t\Big|_0^{3\pi/2} = 9\pi/2.
	\end{align*}
	
This should make sense; we know from geometry that the circumference of a circle with radius 3 is $6\pi$; since we are finding the arc length of $3/4$ of a circle, the arc length is $3/4\cdot 6\pi = 9\pi/2$.
\end{example}

\mtable{A graph of the parametric equations in \autoref{ex_parcalc7}, where the arc length of the teardrop is calculated.}{fig:parcalc7}{\begin{tikzpicture}[alt={Blue teardrop loop crossing the origin, with long arms stretching up left and right before dipping into a rounded tip.}]
\begin{axis}[width=1.16\marginparwidth,tick label style={font=\scriptsize},
axis y line=middle,axis x line=middle,name=myplot,xtick={1,-1},ytick={-1,1},
ymin=-1.1,ymax=1.1,xmin=-1.1,xmax=1.1]
\addplot [draw={\colorone},thick, smooth,domain=-1.5:1.5,samples=40] ({x^3-x},{x^2-1});
\end{axis}
\node [right] at (myplot.right of origin) {\scriptsize $t$};
\node [above] at (myplot.above origin) {\scriptsize $y$};
\end{tikzpicture}}

\begin{example}[Arc Length of a Parametric Curve]\label{ex_parcalc7}%
The graph of the parametric equations $x=t(t^2-1)$, $y=t^2-1$ crosses itself as shown in \autoref{fig:parcalc7}, forming a ``teardrop.'' Find the arc length of the teardrop.
\solution
We can see by the parameterizations of $x$ and $y$ that when $t=\pm 1$, $x=0$ and $y=0$. This means we'll integrate from $t=-1$ to $t=1$. Applying \autoref{thm:arc_length_parametric}, we have
\begin{align*}
L 	&= \int_{-1}^1\sqrt{(3t^2-1)^2+(2t)^2}\dd t\\
		&=	\int_{-1}^1 \sqrt{9t^4-2t^2+1} \dd t.
\end{align*}
Unfortunately, the integrand does not have an antiderivative expressible by elementary functions. We turn to numerical integration to approximate its value. Using 4 subintervals, Simpson's Rule approximates the value of the integral as $2.65051$. Using a computer, more subintervals are easy to employ, and $n=20$ gives a value of $2.71559$. Increasing $n$ shows that this value is stable and a good approximation of the actual value.
\end{example}

% todo Tim Rogowski has a subsection on speed at this point.  Stewart doesn't.

\subsection{Surface Area of a Solid of Revolution}

Related to the formula for finding arc length is the formula for finding surface area. We can adapt the formula found in \autoref{idea:surface_area} from \autoref{sec:arc_length} in a similar way as done to produce the formula for arc length done before.

\begin{keyidea}[Surface Area of a Solid of Revolution]\label{idea:surface_area_parametric}%
Consider the graph of the parametric equations $x=f(t)$ and $y=g(t)$, where $\fp$ and $g\primeskip'$ are continuous on an open interval $I$ containing $t_1$ and $t_2$ on which the graph does not cross itself.\index{surface area!solid of revolution}\index{integration!surface area}\index{parametric equations!surface area}
\begin{enumerate}
	\item	The surface area of the solid formed by revolving the graph about the $x$-axis is (where $g(t)\geq0$ on $[t_1,t_2]$):
\[
\text{Surface Area}
= 2\pi\int_{t_1}^{t_2} g(t)\sqrt{[\fp(t)]^2+[g\primeskip'(t)]^2}\dd t.
\]
	
	\item	The surface area of the solid formed by revolving the graph about the $y$-axis is (where $f(t)\geq0$ on $[t_1,t_2]$):
\[
\text{Surface Area}
= 2\pi\int_{t_1}^{t_2} f(t)\sqrt{[\fp(t)]^2+[g\primeskip'(t)]^2}\dd t.
\]
\end{enumerate}
\end{keyidea}

\begin{example}[Surface Area of a Solid of Revolution]\label{ex_parcalc8}%
%
\mtable{Rotating a teardrop shape about the $x$-axis in \autoref{ex_parcalc8}.}{fig:parcalc8}{
\myincludeasythree{%deactivate=pageinvisible,
3Droll=0,
3Dortho=0.004,
3Dc2c=0.7469545602798462 0.6216799020767212 0.23573943972587585,
3Dcoo=0 0.000 0,
3Droo=120.71813597868993}{\marginparwidth}{Black teardrop curve in xy-plane; light-blue surface shows its sweep when spun 360° around the vertical x-axis.}{figures/figparcalc8_3D}
%{scale=1.5,trim=5mm 5mm 5mm 5mm,clip}{0in}{figures/figparcalc8}
}
%
Consider the teardrop shape formed by the parametric equations $x=t(t^2-1)$, $y=t^2-1$ as seen in \autoref{ex_parcalc7}. Find the surface area if this shape is rotated about the $x$-axis, as shown in \autoref{fig:parcalc8}.
\solution
The teardrop shape is formed between $t=-1$ and $t=1$. Using \autoref{idea:surface_area_parametric}, we see we need for $g(t)\geq 0$ on $[-1,1]$, and this is not the case. To fix this, we simplify replace $g(t)$ with $-g(t)$, which flips the whole graph about the $x$-axis (and does not change the surface area of the resulting solid). The surface area is: 
\begin{align*}
\text{Area}\ S &= 2\pi\int_{-1}^1 (1-t^2)\sqrt{(3t^2-1)^2+(2t)^2}\dd t\\
		&=	2\pi\int_{-1}^1 (1-t^2)\sqrt{9t^4-2t^2+1} \dd t.
		\end{align*}
Once again we arrive at an integral that we cannot compute in terms of elementary functions. Using Simpson's Rule with $n=20$, we find the area to be $S=9.44$. Using larger values of $n$ shows this is accurate to 2 places after the decimal.
\end{example}

After defining a new way of creating curves in the plane, in this section we have applied calculus techniques to the parametric equations defining these curves to study their properties. In the next section, we define another way of forming curves in the plane. To do so, we create a new coordinate system, called \emph{polar coordinates}, that identifies points in the plane in a manner different than from measuring distances from the $y$- and $x$- axes.

\printexercises{exercises/09-03-exercises}

\section{Introduction to Polar Coordinates}\label{sec:polar}

We are generally introduced to the idea of graphing curves by relating $x$-values to $y$-values through a function $f$. That is, we set $y=f(x)$, and plot lots of point pairs $(x,y)$ to get a good notion of how the curve looks. This method is useful but has limitations, not least of which is that curves that ``fail the vertical line test'' cannot be graphed without using multiple functions.

The previous two sections introduced and studied a new way of plotting points in the $x,y$-plane. Using parametric equations, $x$ and $y$ values are computed independently and then plotted together. This method allows us to graph an extraordinary range of curves. This section introduces yet another way to plot points in the plane: using \textbf{polar coordinates}.

\subsection{Polar Coordinates}

\mtable{Illustrating polar coordinates.}{fig:polar_intro1}{%
\begin{tikzpicture}[alt={Origin, initial ray on x-axis labelled 0; second ray makes angle θ (counter-clockwise) marked on arc to point (r, θ).},>=latex]
	\draw[thick,->] (0,0) node [below] {$O$} -- (3,0) node [below] {initial ray} ;
	\filldraw (0,0) circle (1.5pt);
	\filldraw [rotate=55] (2,0) circle (1.5pt);
	\draw [thick,rotate=55] (0,0)-- node [rotate=55,pos=.5,above] {$r$} (2,0) node [above] {$(r,\theta)$};
	\draw [->] (.75,0) arc(0:55:.75); 
	\draw [rotate=27.5] (1,0) node {$\theta$};
\end{tikzpicture}}

Start with a point $O$ in the plane called the \textbf{pole} (we will always identify this point with the origin). From the pole, draw a ray, called the \textbf{initial ray} (we will always draw this ray horizontally, identifying it with the positive $x$-axis). A point $P$ in the plane is determined by the distance $r$ that $P$ is from $O$, and the angle $\theta$ formed between the initial ray and the segment $\overline{OP}$ (measured counter-clockwise). We record the distance and angle as an ordered pair $(r,\theta)$.%To avoid confusion with rectangular coordinates, we will denote polar coordinates with the letter $P$, as in $P(r,\theta)$. This is illustrated in \autoref{fig:polar_intro1}
\index{polar coordinates}\index{polar!coordinates}

\youtubeVideo{r0fv9V9GHdo}{Polar Coordinates --- The Basics}

Practice will make this process more clear.

\begin{example}[Plotting Polar Coordinates]\label{ex_polar1}%
Plot the following polar coordinates:\index{polar coordinates!plotting points}
\[A(1,\pi/4)\quad B(1.5,\pi)\quad C(2,-\pi/3)\quad D(-1,\pi/4)\]
\solution
To aid in the drawing, a polar grid is provided
%
\parbox[t][0pt]{0pt}{%
\tikzset{gridshift/.style={}}
\iflatexml\else
 \checkoddpage%
 \ifbool{oddpage}{%
  \tikzset{gridshift/.append style={shift=(current page.south),shift={+(.5in,1.7in)}}}
 }{%
  \tikzset{gridshift/.append style={shift=(current page.south),shift={+(.5in,1.7in)}}}
 }%
\fi
\begin{tikzpicture}[alt={Three concentric circles (r = 1, 2, 3); points: A at 30 deg (π / 6), B at 210 deg (7π / 6), C at 315 deg (7π / 4), D at 240 deg (4π / 3).},remember picture,overlay,scale=.8,gridshift]
 \foreach \x in {0,30,45,60,90,120,135,150,180,210,225,240,270,300,315,330}
  {\draw [dashed,gray] (0,0) -- (\x:3.1);}
 \draw[thick,->,>=stealth] (0,0) node [below] {$O$} -- (3.5,0) ;
 \filldraw (0,0) circle (1.5pt);
 \foreach \x in {1,2,3} {%
  \draw (0,0) circle (\x cm);
  \draw (\x,0) node [below right] {\x};
 }
\end{tikzpicture}}%
\ifbool{latexml}{here.\\}{at the bottom of this page.}%
%
\mtable{Plotting polar points in \autoref{ex_polar1}.}{fig:polar1}{\begin{tikzpicture}[alt={Concentric polar grid: circles at radii 1, 2, 3 with dashed 30 deg spokes and a bold initial ray pointing right from O.},scale=.75,>=latex]
	\draw[thick,->] (0,0) node [below] {$O$} -- (3.5,0) ;
	\filldraw (0,0) circle (1.5pt);
	\foreach \x in {1,2,3} {%
		\draw (0,0) circle (\x cm);
		\draw (\x,0) node [below right] {\x};
	}
\filldraw [rotate=45] (1,0) circle (1.5pt) node [above right] {$A$};
\filldraw [rotate=180] (1.5,0) circle (1.5pt) node [above] {$B$};
\filldraw [rotate=-60] (2,0) circle (1.5pt) node [below right] {$C$};
\filldraw [rotate=45] (-1,0) circle (1.5pt) node [below] {$D$};
\end{tikzpicture}}
%
To place the point $A$, go out 1 unit along the initial ray (putting you on the inner circle shown on the grid), then rotate counter-clockwise $\pi/4$ radians (or $45^\circ$).  Alternately, one can consider the rotation first: think about the ray from $O$ that forms an angle of $\pi/4$ with the initial ray, then move out 1 unit along this ray (again placing you on the inner circle of the grid).

To plot $B$, go out $1.5$ units along the initial ray and rotate $\pi$ radians ($180^\circ$). 

To plot $C$, go out 2 units along the initial ray then rotate \emph{clockwise} $\pi/3$ radians, as the angle given is negative.

To plot $D$, move along the initial ray ``$-1$'' units --- in other words, ``back up'' 1 unit, then rotate counter-clockwise by $\pi/4$. The results are given in \autoref{fig:polar1}.
\end{example}

Consider the following two points: $A(1,\pi)$ and $B(-1,0)$. To locate $A$, go out 1 unit on the initial ray then rotate $\pi$ radians; to locate $B$, go out $-1$ units on the initial ray and don't rotate. One should see that $A$ and $B$ are located at the same point in the plane. We can also consider $C(1,3\pi)$, or $D(1,-\pi)$; all four of these points share the same location. 

This ability to identify a point in the plane with multiple polar coordinates is both a ``blessing'' and a ``curse.'' We will see that it is beneficial as we can plot beautiful functions that intersect themselves (much like we saw with parametric functions). The unfortunate part of this is that it can be difficult to determine when this happens. We'll explore this more later in this section.

\subsection{Polar to Rectangular Conversion}

\mtable{Converting between rectangular and polar coordinates.}{fig:polar_intro2}{\begin{tikzpicture}[alt={Right triangle showing rectangular–polar link: legs x, y, hypotenuse r; angle at origin labelled θ},scale=.75,>=latex]
	\draw [thick] (0,0) -- node [below,pos=.5] {$x$} (4,0)
	 -- node [right,pos=.5] {$y$} (4,2)
	 -- node [above,rotate=26.57,pos=.5] {$r$} (0,0);
	\draw [->,thick] (1.5,0) arc (0:26:1.5);
	\draw [rotate=13] (1.8,0) node {$\theta$};
	\filldraw (0,0) circle (1.5pt) node [below] {$O$}
			(4,2) circle (1.5pt) node [above] {$P$};
\end{tikzpicture}}

It is useful to recognize both the rectangular (or, Cartesian) coordinates of a point in the plane and its polar coordinates. \autoref{fig:polar_intro2} shows a point $P$ in the plane with rectangular coordinates $(x,y)$ and polar coordinates $(r,\theta)$. Using trigonometry, we can make the identities given in the following Key Idea.

\begin{keyidea}[Converting Between Rectangular and Polar Coordinates]\label{idea:polarconvert}%
Given the polar point $P(r,\theta)$, the rectangular coordinates are determined by \[x=r\cos \theta\qquad y=r\sin \theta.\]

Given the rectangular coordinates $(x,y)$, the polar coordinates are determined by
\[r^2=x^2+y^2\qquad \tan \theta = \frac yx.\]
\end{keyidea}

\begin{example}[Converting Between Polar and Rectangular Coordinates]\label{ex_polar2}%
\mbox{}\\[-2\baselineskip]\parbox[t]{\linewidth}{%
\begin{enumerate}
\item		Convert the polar coordinates $A(2,2\pi/3)$ and $B(-1,5\pi/4)$ to rectangular coordinates.
\item		Convert the rectangular coordinates $(1,2)$ and $(-1,1)$ to polar coordinates.
\end{enumerate}}
\solution
\begin{enumerate}
	\item \begin{enumerate}
		\item 
	We start with $A(2,2\pi/3)$. Using \autoref{idea:polarconvert}, we have 
	\[x= 2\cos (2\pi/3) = -1\qquad y = 2\sin (2\pi/3) = \sqrt{3}.\]
	So the rectangular coordinates are $(-1,\sqrt{3}) \approx (-1,1.732)$.
	
	\item The polar point $B(-1,5\pi/4)$ is converted to rectangular with:
	\[x=-1\cos (5\pi/4) = \sqrt{2}/2\qquad y= -1\sin (5\pi/4) = \sqrt{2}/2.\]
	So the rectangular coordinates are $(\sqrt{2}/2,\sqrt{2}/2) \approx (0.707,0.707)$.
	\end{enumerate}
	These points are plotted in \autoref{fig:polar2} (a). The rectangular coordinate system is drawn lightly under the polar coordinate system so that the relationship between the two can be seen.
	
\mtable{Plotting rectangular and polar points in \autoref{ex_polar2}.}{fig:polar2}{%
\begin{tikzpicture}[alt={Polar grid with bold circles r = 1 and 2; A at 120 deg (2pi / 3) on r = 2, B at 225 deg (5pi / 4) on r = 1.},>=latex]
 \draw [thin,gray] (-2.5,-2.5) grid (2.5,2.5);
 \draw [thick,->] (0,0) node [below] {$O$} -- (2.5,0);
 \foreach \x in {1,2}{%
  \draw [very thick] (0,0) circle (\x cm);
 }
 \filldraw (0,0) circle (1.5pt);
 \filldraw [rotate=120] (2,0) circle (1.5pt) node [ left] {\scriptsize $A(2,\frac{2\pi}{3})$};
 \filldraw [rotate=225] (-1,0) circle (1.5pt) node [shift={(0pt,12pt)}] {\scriptsize$B(-1,\frac{5\pi}{4})$};
\end{tikzpicture}
\smallskip\\(a)\bigskip\\
\begin{tikzpicture}[alt={Rectangular grid over light polar circles; dashed lines show angle 135 deg (3pi / 4) to locate (1, 2) and (–1, 1).},>=latex]
\draw [very thick] (-2.5,-2.5) grid (2.5,2.5);
\draw [thin,->] (0,0) node [below left,black] {\scriptsize$(0,0)$} -- (2.5,0);
\foreach \x in {1,2} {%
 \draw [thin] (0,0) circle (\x cm);
}
\filldraw (0,0) circle (1.5pt);
\filldraw [] (1,2) circle (1.5pt) node [above left] {\scriptsize $(1,2)$};
\filldraw [] (-1,1) circle (1.5pt) node [above right] {\scriptsize$(-1,1)$};
\draw [gray,dashed] (0,0) -- ($(0,0)!2.2!(-1,1)$);
\draw [gray,->] (.5,0) arc (0:135:.5);
\draw [rotate=35] (.65,0) node [] {\scriptsize$\frac{3\pi}{4}$};
\draw [gray,dashed] (0,0) -- ($(0,0)!1.2!(1,2)$);
\draw [gray,->] (.5,0) arc (0:135:.5);
\draw [gray,->] (-.5,0) arc (180:135:.5);
\draw [rotate=35] (.65,0) node [] {\scriptsize$\frac{3\pi}{4}$};
\draw [rotate=-30] (-.65,0) node [shift={(-3pt,0pt)}] {\scriptsize$\frac{-\pi}{4}$};
\draw [gray,->] (1.5,0) arc (0:63.4:1.5);
\draw [rotate=28] (1.75,0) node [] {\scriptsize$1.11$};
\end{tikzpicture}
\smallskip\\(b)}

	\item \begin{enumerate}
		\item To convert the rectangular point $(1,2)$ to polar coordinates, we use the Key Idea to form the following two equations:
		\[1^2+2^2 = r^2 \qquad \tan \theta = \frac{2}{1}.\]
		The first equation tells us that $r=\sqrt{5}$. Using the inverse tangent function, we find
		\[\tan \theta = 2 \quad \Rightarrow \quad \theta = \tan^{-1} 2 \approx 1.11\text{ radians}\approx 63.43^\circ.\]
		Thus polar coordinates of $(1,2)$ are $(\sqrt{5},1.11)$.
		\item		To convert $(-1,1)$ to polar coordinates, we form the equations 
		\[(-1)^2+1^2=r^2 \qquad \tan \theta = \frac{1}{-1}.\]
		Thus $r=\sqrt{2}$. We need to be careful in computing $\theta$: using the inverse tangent function, we have
		\[\tan\theta = -1 \quad \Rightarrow \quad \theta = \tan^{-1}(-1) = -\pi/4.\]
		This is not the angle we desire. The range of $\tan^{-1}x $ is $(-\pi/2,\pi/2)$; that is, it returns angles that lie in the $1^\text{st}$ and $4^\text{th}$ quadrants. To find locations in the $2^\text{nd}$  and $3^\text{rd}$ quadrants, add $\pi$ to the result of $\tan^{-1}x$. So  $\pi+(-\pi/4)$ puts the angle at $3\pi/4$. Thus the polar point is $(\sqrt{2},3\pi/4)$.
		
		An alternate method is to use the angle $\theta$ given by arctangent, but change the sign of $r$. Thus we could also refer to $(-1,1)$ as\\ $(-\sqrt{2},-\pi/4)$.
	\end{enumerate}
These points are plotted in \autoref{fig:polar2} (b). The polar system is drawn lightly under the rectangular grid with rays to demonstrate the angles used.
\end{enumerate}
\end{example}

\subsection{Polar Functions and Polar Graphs}

Defining a new coordinate system allows us to create a new kind of function, a \textbf{polar function.} Rectangular coordinates lent themselves well to creating functions that related $x$ and $y$, such as $y=x^2$. Polar coordinates allow us to create functions that relate $r$ and $\theta$. Normally these functions look like $r=f(\theta)$, although we can create functions of the form $\theta = f(r)$. The following examples introduce us to this concept.\index{polar!functions}\index{polar!functions!graphing}

\begin{example}[Introduction to Graphing Polar Functions]\label{ex_polar3}%
Describe the graphs of the following polar functions.
\begin{enumerate}
	\item $r = 1.5$
	\item $\theta = \pi/4 $
\end{enumerate}
%
%\noindent$1.\ r = 1.5 \qquad 2.\ \theta = \pi/4 $
\solution
\begin{enumerate}
\item		The equation $r=1.5$ describes all points that are 1.5 units from the pole; as the angle is not specified, any $\theta$ is allowable. All points 1.5 units from the pole describes a circle of radius 1.5.

We can consider the rectangular equivalent of this equation; using $r^2=x^2+y^2$, we see that $1.5^2=x^2+y^2$, which we recognize as the equation of a circle centered at $(0,0)$ with radius 1.5. This is sketched in \autoref{fig:polar3}.

\mtable{Plotting standard polar plots.}{fig:polar3}{\begin{tikzpicture}[alt={Concentric polar rings r = 0.5, 1 (bold), 1.5 with a thick ray at θ = π/4 crossing r = 1.5.},>=latex]
 \draw [thick,->] (0,0) node [below] {$O$} -- (2.5,0);
 \foreach \x in {1,2}
  {\draw [thin,gray] (0,0) circle (\x cm);
   \draw (\x,0) node [shift={(3pt,-4pt)}] {\scriptsize\x};
  }
 \filldraw (0,0) circle (1.5pt);
 \draw [very thick] (0,0) circle (1.5);
 \draw[rotate=90] (1.7,0) node {\scriptsize $r=1.5$};
 \draw [thick] (-2,-2)
  -- node [pos=1,left] {\scriptsize $\theta = \frac{\pi}4$} (2,2);
\end{tikzpicture}}

\item		The equation $\theta = \pi/4$ describes all points such that the line through them and the pole make an angle of $\pi/4$ with the initial ray. As the radius $r$ is not specified, it can be any value (even negative). Thus $\theta = \pi/4$ describes the line through the pole that makes an angle of $\pi/4 = 45^\circ$ with the initial ray.

We can again consider the rectangular equivalent of this equation. Combine $\tan \theta =y/x$ and $\theta =\pi/4$:
\[\tan(\pi/4) = y/x \quad \Rightarrow \quad x\tan(\pi/4) = y \quad \Rightarrow \quad y = x.\] 
This graph is also plotted in \autoref{fig:polar3}.
\end{enumerate}
\end{example}

The basic rectangular equations of the form $x=h$ and $y=k$ create vertical and horizontal lines, respectively; the basic polar equations $r= h$ and $\theta =\alpha$ create circles and lines through the pole, respectively. With this as a foundation, we can create more complicated polar functions of the form $r=f(\theta)$. The input is an angle; the output is a length, how far in the direction of the angle to go out.

We sketch these functions much like we sketch rectangular and parametric functions: we plot lots of points and ``connect the dots'' with curves. We demonstrate this in the following example.

\begin{example}[Sketching Polar Functions]\label{ex_polar4}%
Sketch
%
\mtable[-1in]{Graph of the polar function in \autoref{ex_polar4} by plotting points.}{fig:polar4}{%
\tagpdfsetup{table/header-rows={1}}
\begin{tabular}{c}
 {$\begin{array}{cc}
	\theta & r=1+\cos\theta \\ \midrule
	0 & 2 \\
	\pi/6 & 1+\sqrt{3}/2 \\
	\pi/4 & 1+1/\sqrt2 \\
	\pi/3 & 3/2 \\
	\pi/2 & 1 \\
	2\pi/3 & 1/2 \\
	3\pi/4 & 1-1/\sqrt2 \\
	5\pi/6 & 1-\sqrt3/2 \\
	\pi & 0 \\
	7\pi/6 & 1-\sqrt3/2 \\
	5\pi/4 & 1-1/\sqrt2 \\
	4\pi/3 & 1/2 \\
	3\pi/2 & 1 \\
	5\pi/3 & 3/2 \\
	7 \pi /4 & 1+1/\sqrt2 \\
	11\pi/6 & 1+\sqrt{3}/2 \\
 \end{array}$} \\ \\
 \begin{tikzpicture}[alt={Polar grid; right-facing cardioid r = 1 + cos θ touches r = 0 at π and peaks at 2 on initial ray.},scale=.95]
  \draw [dashed,gray] (-2.1,0) -- (0,0);
  \draw[thick,->,>=stealth] (0,0) node [below] {$O$} -- (2.5,0) ;
  \filldraw (0,0) circle (1.5pt);
  \foreach \x in {1,2}
   {\draw[gray] (0,0) circle (\x cm);
    \draw (\x,0) node [shift={(3pt,-4pt)}] {\scriptsize\x};
   }
  \foreach \x in {30,45,60,90,120,135,150}
   {\draw [rotate=\x,dashed,gray] (-2.3,0) -- (2.3,0);}
  \draw [thick,draw={\colorone},domain=0:360,samples=60,smooth] plot
   ({cos(\x)*(1+cos(\x))},{sin(\x)*(1+cos(\x))});
  \foreach \x in {0,30,45,60, 90,120,135,150, 180,210,225,240, 270,300,315,330}
   {\filldraw ({\x}:{1+cos(\x)}) circle (1pt); }
  \foreach \x/\y in {45/{\pi/4}, 90/{\pi/2}, 135/{3\pi/4}, 180/{\pi}, 225/{5\pi/4}, 270/{3\pi/2}, 315/{7\pi/4} }
   {\draw({\x}:{2.3}) node{$\scriptstyle\y$}; }
\end{tikzpicture}
\end{tabular}}
%
the polar function $r=1+\cos \theta$ on $[0,2\pi]$ by plotting points.
\solution
A common question when sketching curves by plotting selected points is ``Which points should I plot?'' With rectangular equations, we often chose ``easy'' values --- integers, then added more if needed. When plotting polar equations, start with the ``common'' angles --- multiples of $\pi/6$ and $\pi/4$. \autoref{fig:polar4} gives a table of just a few values of $\theta$ in $[0,\pi]$. 

Consider the point $(2,0)$ determined by the first line of the table. The angle is 0 radians --- we do not rotate from the initial ray -- then we go out 2 units from the pole. When $\theta=\pi/6$, $r = 1+\sqrt{3}/2$; so rotate by $\pi/6$ radians and go out $1+\sqrt{3}/2$ units.
\end{example}

%\paragraph{Technology Note:} Plotting functions in this way can be tedious, just as it was with rectangular functions. To obtain very accurate graphs, technology is a great aid. Most graphing calculators can plot polar functions; in the menu, set the plotting mode to something like \texttt{polar} or \texttt{POL}, depending on one's calculator. As with plotting parametric functions, the viewing ``window'' no longer determines the $x$-values that are plotted, so additional information needs to be provided. Often with the ``window'' settings are the settings for  the beginning and ending $\theta$ values (often called \texttt{$\theta_{\text{min}}$} and \texttt{$\theta_{\text{max}}$}) as well as the \texttt{$\theta_{\text{step}}$} --- that is, how far apart the $\theta$ values are spaced. The smaller the \texttt{$\theta_{\text{step}}$} value, the more accurate the graph (which also increases plotting time). Using technology, we graphed the polar function $r=1+\cos \theta$ from \autoref{ex_polar4} in \autoref{fig:polar4b}.
%
%\mtable{Using  technology to graph a polar function.}{fig:polar4b}{\begin{tikzpicture}[alt={Four-petal polar rose; dots 1–17 trace one full loop through symmetric petals about the origin.},scale=1]
%	\draw [dashed,gray] (-2.1,0) -- (0,0);
%	\draw[thick,->,>=stealth] (0,0) node [below] {$O$} -- (2.5,0) ;
%	\filldraw (0,0) circle (1.5pt);
%	\foreach \x in {1,2} {
%		\draw (0,0) circle (\x cm);
%		\draw (\x,0) node [shift={(3pt,-4pt)}] {\scriptsize\x};
%	}
%	\foreach \x in {30,45,60,90,120,135,150} {
%		\draw [rotate=\x,dashed,gray] (-2.3,0) -- (2.3,0);
%	}
%	\draw [thick,draw={\colorone},domain=0:360,samples=60,smooth]
%	 plot ({cos(\x)*(1+cos(\x))},{sin(\x)*(1+cos(\x))});
%\end{tikzpicture}}

\begin{example}[Sketching Polar Functions]\label{ex_polar5}%
Sketch the polar function $r=\cos (2\theta)$ on $[0,2\pi]$ by plotting points.
\solution
We start by making a table of $\cos (2\theta)$ evaluated at common angles $\theta$, as shown in \autoref{fig:polar5table}. These points are then plotted in \autoref{fig:polar5}. This particular graph ``moves'' around quite a bit and one can easily forget which points should be connected to each other. To help us with this, we numbered each point in the table and on the graph. 

\begin{center}
\captionsetup{type=figure}%
$\begin{array}{ccl @{\hspace{4em}} ccl}\toprule
\text{Pt.} & \theta & \cos (2\theta) & \text{Pt.} & \theta & \cos (2\theta) \\
\iflatexml
\midrule
\else
\cmidrule(r{4em}){1-3}\cmidrule(l{-1em}){4-6}
\fi
 1 & 0 & \phantom{-}1 & 10 & 7 \pi /6 & \phantom{-}0.5 \\
 2 & \pi /6 & \phantom{-}0.5 & 11 & 5 \pi /4 & \phantom{-}0 \\
 3 & \pi /4 & \phantom{-}0 & 12 & 4 \pi /3 & -0.5 \\
 4 & \pi /3 & -0.5 & 13 & 3 \pi /2 & -1 \\
 5 & \pi /2 & -1 & 14 & 5 \pi /3 & -0.5 \\
 6 & 2 \pi /3 & -0.5 & 15 & 7 \pi /4 & \phantom{-}0 \\
 7 & 3 \pi /4 & \phantom{-}0 & 16 & 11 \pi /6 & \phantom{-}0.5 \\
 8 & 5 \pi /6 & \phantom{-}0.5 & 17 & 2 \pi & \phantom{-}1 \\
 9 & \pi  & \phantom{-}1 \\\bottomrule
\end{array}$
\caption{Tables of points for plotting a polar curve.}
\label{fig:polar5table}
\end{center}

\mtable[-2in]{Polar plots from \autoref{ex_polar5}.}{fig:polar5}{%
 \begin{tikzpicture}[alt={Four-petal polar rose; dots 1–17 trace one full loop through symmetric petals about the origin.},scale=1]
  \draw [dashed,gray] (-2.1,0) -- (0,0);
  \draw[thick,->,>=latex] (0,0) node [below] {} -- (2.5,0) ;
  \filldraw (0,0) circle (1.5pt);
  \foreach \x in {1,2}
   {\draw[gray] (0,0) circle (\x cm);}
  \foreach \x in {30,45,60,90,120,135,150}
   {\draw [rotate=\x,dashed,gray] (-2.3,0) -- (2.3,0);}
%  \draw [thick,draw={\colorone}] plot coordinates {(2.,0)(0.866,0.5)(0,0)(-0.5,-0.866)(0,-2.)(0.5,-0.866)(0,0)(-0.866,0.5)(-2.,0)(-0.866,-0.5)(0,0)(0.5,0.866)(0,2.)(-0.5,0.866)(0,0)(0.866,-0.5)(2,0)};
  \foreach \x/\y/\z/\w in
   {2./0/1/{above right}, 0.866/0.5/2/above, 0/0/3/above, -0.5/-0.866/4/above,
    0/-2./5/above, 0.5/-0.866/6/above, 0/0/7/left, -0.866/0.5/8/above,
    -2./0/9/above left, -0.866/-0.5/10/above, 0/0/11/right, 0.5/0.866/12/above,
    0/2./13/above, -0.5/0.866/14/above, 0/0/15/below, 0.866/-0.5/16/above,
    2/0/17/{below right}}
   {\filldraw (\x,\y) circle (1.5pt) node [\w] {\scriptsize \z};}       
  \draw [thick,draw={\colorone},domain=0:360,samples=60,smooth]
   plot ({\x}:{2*cos(2*\x)});
\end{tikzpicture}}

%Using more points (and the aid of technology) a smoother plot can be made as shown in \autoref{fig:polar5} (b).
This plot is an example of a \emph{rose curve}.
\end{example}

It is sometimes desirable to refer to a graph via a polar equation, and other times by a rectangular equation. Therefore it is necessary to be able to convert between polar and rectangular functions, which we practice in the following example. We will make frequent use of the identities found in \autoref{idea:polarconvert}.

\begin{example}[Converting between rectangular and polar equations.]\label{ex_polar6}%
\begin{minipage}[t]{.5\linewidth}
Convert from rectangular to polar.
\begin{enumerate}
	\item $y=x^2$
	\item $xy = 1$
\end{enumerate}
\end{minipage}%
\begin{minipage}[t]{.5\linewidth}
Convert from polar to rectangular.
\begin{enumerate}\addtocounter{enumi}{2}
	\item $\ds r=\frac{2}{\sin \theta-\cos\theta}$
	\item $r=2\cos \theta$
\end{enumerate}
\end{minipage}
\solution
\begin{enumerate}
	\item Replace $y$ with $r\sin\theta$ and replace $x$ with $r\cos\theta$, giving:
	\begin{align*}
	y &=x^2\\
	r\sin\theta &= r^2\cos^2\theta\\
	\frac{\sin\theta}{\cos^2\theta} &= r
	\end{align*}
	We have found that $r=\sin\theta/\cos^2\theta = \tan\theta\sec\theta$. The domain of this polar function is $(-\pi/2,\pi/2)$; plot a few points to see how the familiar parabola is traced out by the polar equation.
	
	\item		We again replace $x$ and $y$ using the standard identities and work to solve for $r$:
	\begin{align*}
	xy &= 1 \\
	r\cos\theta\cdot r\sin\theta & = 1\\
	r^2 & = \frac{1}{\cos\theta\sin\theta}\\
	r & = \frac{1}{\sqrt{\cos\theta\sin\theta}}\\
	\end{align*}
%
\mtable{Graphing $xy=1$ from \autoref{ex_polar6}.}{fig:polar6}{\begin{tikzpicture}[alt={Rectangular plot of xy = 1; hyperbola branches in quadrants I and III approach both coordinate axes.}]
\begin{axis}[width=1.16\marginparwidth,tick label style={font=\scriptsize},
axis y line=middle,axis x line=middle,name=myplot,minor x tick num=4,
minor y tick num=4,ymin=-5.3,ymax=5.3,xmin=-5.3,xmax=5.3,axis equal]
\addplot [draw={\colorone},thick, smooth,domain=-5:-.1,samples=30] {1/x};
\addplot [draw={\colorone},thick, smooth,domain=.1:5,samples=30] {1/x};
\end{axis}
\node [right] at (myplot.right of origin) {\scriptsize $x$};
\node [above] at (myplot.above origin) {\scriptsize $y$};
\end{tikzpicture}}%
%
	This function is valid only when the product of $\cos\theta\sin\theta$ is positive. This occurs in the first and third quadrants, meaning the domain of this polar function is $(0,\pi/2) \cup (\pi,3\pi/2)$.
	
	We can rewrite the original rectangular equation $xy=1$ as $y=1/x$. This is graphed in \autoref{fig:polar6}; note how it only exists in the first and third quadrants.
		
	\item		There is no set way to convert from polar to rectangular; in general, we look to form the products $r\cos \theta$ and $r\sin\theta$, and then replace these with $x$ and $y$, respectively. We start in this problem by multiplying both sides by $\sin\theta-\cos\theta$:
	\begin{align*}
	r &= \frac{2}{\sin\theta-\cos\theta} \\
	r(\sin\theta-\cos\theta) &= 2\\
	r\sin\theta-r\cos\theta &= 2. \qquad \text{Now replace with $y$ and $x$:}\\
	y-x &= 2\\
	y &= x+2.
	\end{align*}
	The original polar equation, $r=2/(\sin\theta-\cos\theta)$ does not easily reveal that its graph is simply a line. However, our conversion shows that it is. The upcoming gallery of polar curves gives the general equations of lines in polar form.

	\item		By multiplying both sides by $r$, we obtain both an $r^2$ term and an $r\cos\theta$ term, which we replace with $x^2+y^2$ and $x$, respectively. 
	\begin{align*}
	r &=2\cos\theta \\
	r^2 &= 2r\cos\theta \\
	x^2+y^2 &= 2x. 
	\intertext{We recognize this as a circle; by completing the square we can find its radius and center.}
	x^2-2x+y^2 &= 0 \\
	(x-1)^2 + y^2 &=1.
	\end{align*}
	The circle is centered at $(1,0)$ and has radius 1. The upcoming gallery of polar curves gives the equations of \emph{some} circles in polar form; circles with arbitrary centers have a complicated polar equation that we do not consider here.
\end{enumerate}
\end{example}

Some curves have very simple polar equations but rather complicated rectangular ones. For instance, the equation $r=1+\cos\theta$ describes a \emph{cardioid} (a shape important to the sensitivity of microphones, among other things; one is graphed in the gallery in the Limaçon section). Its rectangular form is not nearly as simple; it is the implicit equation
$x^4+y^4+2x^2y^2-2xy^2-2x^3-y^2=0$. The conversion is not ``hard,'' but takes several steps, and is left as an exercise.

\subsection{Gallery of Polar Curves}

There are a number of basic and ``classic'' polar curves, famous for their beauty and/or applicability to the sciences. \index{polar!function!gallery of graphs} This section ends with a small gallery of some of these graphs. We encourage the reader to understand how these graphs are formed, and to investigate with technology other types of polar functions.

\newlength{\gallerywidth}
%\setlength{\gallerywidth}{.25\textwidth}
% no.  the textwidth doesn't change until \exercisegeometry in 50 lines
\setlength{\gallerywidth}{(0pt+\marginparwidth+\textwidth)/4}

\noindent
\flushinner{%
\noindent\hrulefill\bigskip\\
\tagpdfsetup{table/header-rows={1,2}}
\begin{tabular}{@{}p{\gallerywidth}@{}p{\gallerywidth}@{}p{\gallerywidth}@{}p{\gallerywidth}@{}}
%
\multicolumn{4}{l}{\bfseries\large Lines}\smallskip\\
\textbf{Through the origin:} & \textbf{Horizontal line:} & \textbf{Vertical line:} & \textbf{Not through origin:} \smallskip\\
$\theta = \alpha$ & $r=a\csc\theta$ & $r=a\sec\theta$ & $\ds r=\frac{b}{\sin\theta-m\cos\theta}$ \medskip\\
\begin{tikzpicture}[alt={Slanted blue ray through origin making angle α with the polar axis.},scale=.9,>=stealth]
	\draw [<->,] (-2.1,0) -- (2.1,0);
	\draw [<->,] (0,-2.1) -- (0,2.1);
	\draw [thick,draw={\colorone},rotate=50] (-2.1,0) -- (2.1,0);
	\draw [->] (.5,0) arc (0:50:.5);
	\draw [rotate=25] (.7,0) node {\scriptsize $\alpha$};
\end{tikzpicture}		
&
\begin{tikzpicture}[alt={Horizontal blue line at y = a drawn from r = a csc θ.},scale=.9,>=stealth]
	\draw [<->,] (-2.1,0) -- (2.1,0);
	\draw [<->,] (0,-2.1) -- (0,2.1);
	\draw [thick,draw={\colorone}] (-2,.6) -- (2,.6);
	\draw (-.2,.3) node {\scriptsize $a\left\{\rule[-.23cm]{0pt}{.23cm}\right.$};
\end{tikzpicture}		
&
\begin{tikzpicture}[alt={Vertical blue line at x = a drawn from r = a sec θ.},scale=.9,>=stealth]
	\draw [<->,] (-2.1,0) -- (2.1,0);
	\draw [<->,] (0,-2.1) -- (0,2.1);
	\draw [thick,draw={\colorone}] (.6,-2) -- (.6,2);
	\draw (.3,.2) node {\scriptsize $\overbrace{\makebox[.55cm]{}}$};
	\draw (.3,.4) node {\scriptsize $a$};
\end{tikzpicture}		
&
\begin{tikzpicture}[alt={Blue line with slope m intercept b, given by r = b / (sin θ − m cos θ).},scale=.9,>=stealth]
	\draw [<->,] (-2.1,0) -- (2.1,0);
	\draw [<->,] (0,-2.1) -- (0,2.1);
	\draw [thick,draw={\colorone},shift={(0,.6)},rotate=50] (-2.1,0) -- (1.7,0)
	 node [pos=.82,rotate=50,black,shift={(0,-5pt)}] {\scriptsize slope $=m$};
	\draw (.2,.3) node {\scriptsize $\left.\rule[-.23cm]{0pt}{.23cm}\right\}b$};
\end{tikzpicture}		
%
\end{tabular}}

\exercisegeometry

\thispagestyle{empty}%
\enlargethispage{\baselineskip}% otherwise, the table gets pushed to a new page for some reason
%
\noindent%\flushinner{%
\tagpdfsetup{table/header-rows={1,5,10,15}}%
\begin{tabular}{@{}p{\gallerywidth}@{}p{\gallerywidth}@{}p{\gallerywidth}@{}p{\gallerywidth}@{}}
%
\multicolumn{3}{l}{\bfseries\large Circles} & \textbf{\large Spiral} \\
\textbf{Centered on origin:} & $(x-\frac a2)^2+y^2=\frac{a^2}4$ & $x^2+(y-\frac a2)^2=\frac{a^2}4$ & \textbf{Archimedean spiral}\\
$r=a$ & $r=a\cos \theta$ & $r=a\sin\theta$ & $r=\theta$\\
\begin{tikzpicture}[alt={Blue circle centered at pole, radius a, equation r = a.},scale=.9,>=stealth]
	\draw [<->,] (-2.1,0) -- (2.1,0);
	\draw [<->,] (0,-2.1) -- (0,2.1);
	\draw [thick,draw={\colorone}] (0,0) circle (.9);
	\draw (.45,.1) node {\scriptsize $\overbrace{\makebox[.8cm]{}}$};
	\draw (.45,.3) node {\scriptsize $a$};
\end{tikzpicture}		
&
\begin{tikzpicture}[alt={Blue circle center (a/2, 0), radius a/2, equation r = a cos θ.},scale=.9,>=stealth]
	\draw [<->,] (-2.1,0) -- (2.1,0);
	\draw [<->,] (0,-2.1) -- (0,2.1);
	\draw [thick,draw={\colorone}] (.9,0) circle (.9);
	\draw (.9,.1) node {\scriptsize $\overbrace{\makebox[1.7cm]{}}$};
	\draw (.9,.3) node {\scriptsize $a$};
\end{tikzpicture}		
&
\begin{tikzpicture}[alt={Blue circle center (0, a/2), radius a/2, equation r = a sin θ.},scale=.9,>=stealth]
	\draw [<->,] (-2.1,0) -- (2.1,0);
	\draw [<->,] (0,-2.1) -- (0,2.1);
	\draw [thick,draw={\colorone}] (0,.9) circle (.9);
	\draw (-.2,.9) node {\scriptsize $a\left\{\rule[-.8cm]{0cm}{0.8cm}\right.$};
\end{tikzpicture}		
&
\begin{tikzpicture}[alt={Blue Archimedean spiral r = θ making several outward turns from the pole.},scale=.9,>=stealth]
	\draw [<->,] (-2.1,0) -- (2.1,0);
	\draw [<->,] (0,-2.1) -- (0,2.1);
	\draw [thick,draw={\colorone},domain=0:18.85,samples=100,smooth] plot
	 ({cos(\x r)*(\x/9.5)},{sin(\x r)*(\x/9.5)});
\end{tikzpicture}		
\bigskip\\
%
\multicolumn{4}{l}{\bfseries\large Limaçons}\\
\multicolumn{4}{l}{Symmetric about $x$-axis: $r=a\pm b\cos\theta$; \qquad
Symmetric about $y$-axis:  $r=a\pm b\sin \theta$; \qquad $a,b>0$}\\
\textbf{With inner loop:} & \textbf{Cardioid:} &
\textbf{Dimpled:} & \textbf{Convex:} \smallskip\\
$\dfrac ab < 1$ & $\dfrac ab=1$ & $1<\dfrac ab <2$ & $\dfrac ab>2$ \\
\begin{tikzpicture}[alt={Inner-loop limacon, symmetric about x-axis; small interior loop inside larger curve when a/b < 1.},scale=.9,>=stealth]
	\draw [<->,] (-2.1,0) -- (2.1,0);
	\draw [<->,] (0,-2.1) -- (0,2.1);
	\draw [thick,draw={\colorone},domain=0:360,samples=60,smooth] plot
	 ({cos(\x)*(.6+1.2*cos(\x))},{sin(\x)*(.6+1.2*cos(\x))});
\end{tikzpicture}		
&
\begin{tikzpicture}[alt={Cardioid-shaped limacon, symmetric about x-axis; single heart-like loop touching the origin when a/b = 1.},scale=.9,>=stealth]
	\draw [<->,] (-2.1,0) -- (2.1,0);
	\draw [<->,] (0,-2.1) -- (0,2.1);
	\draw [thick,draw={\colorone},domain=0:360,samples=60,smooth] plot
	 ({cos(\x)*(.9+.9*cos(\x))},{sin(\x)*(.9+.9*cos(\x))});
\end{tikzpicture}		
&
\begin{tikzpicture}[alt={Dimpled limacon, symmetric about y-axis; slight inward dent on each side for 1 < a/b < 2.},scale=.9,>=stealth]
	\draw [<->,] (-2.1,0) -- (2.1,0);
	\draw [<->,] (0,-2.1) -- (0,2.1);
	\draw [thick,draw={\colorone},domain=0:360,samples=60,smooth] plot
	 ({cos(\x)*(1+.8*cos(\x))},{sin(\x)*(1+.8*cos(\x))});
\end{tikzpicture}		
&
\begin{tikzpicture}[alt={Convex limacon, symmetric about y-axis; smooth oval without dimples when  /b > 2.},scale=.9,>=stealth]
	\draw [<->,] (-2.1,0) -- (2.1,0);
	\draw [<->,] (0,-2.1) -- (0,2.1);
	\draw [thick,draw={\colorone},domain=0:360,samples=60,smooth] plot
	 ({cos(\x)*(1.3+.6*cos(\x))},{sin(\x)*(1.3+.6*cos(\x))});
\end{tikzpicture}		
\bigskip\\
%
\multicolumn{4}{l}{\bfseries\large Rose Curves}\\
\multicolumn{4}{l}{Symmetric about $x$-axis: $r=a \cos(n\theta)$; \qquad
Symmetric about $y$-axis:  $r=a\sin(n\theta)$}\\
\multicolumn{4}{l}{Curve contains $2n$ petals when $n$ is even and $n$ petals when $n$ is odd.}\\
$r=a\cos (2\theta)$ & $r=a\sin(2\theta)$ &
$r=a\cos (3\theta)$ & $r=a\sin (3\theta)$ \smallskip\\
\begin{tikzpicture}[alt={Four-petal rose; long petals point along both axes, symmetric shape traced by r = a cos 2θ.},scale=.9,>=stealth]
	\draw [<->,] (-2.1,0) -- (2.1,0);
	\draw [<->,] (0,-2.1) -- (0,2.1);
	\draw [thick,draw={\colorone},domain=0:360,samples=60,smooth] plot
	 ({cos(\x)*(1.9*cos(2*\x))},{sin(\x)*(1.9*cos(2*\x))});
\end{tikzpicture}		
&
\begin{tikzpicture}[alt={Four-petal rose; petals lie between axes at 45 degree, symmetric curve from r = a sin 2θ.},scale=.9,>=stealth]
	\draw [<->,] (-2.1,0) -- (2.1,0);
	\draw [<->,] (0,-2.1) -- (0,2.1);
	\draw [thick,draw={\colorone},domain=0:360,samples=60,smooth] plot
	 ({cos(\x)*(1.9*sin(2*\x))},{sin(\x)*(1.9*sin(2*\x))});
\end{tikzpicture}		
&
\begin{tikzpicture}[alt={Three-petal rose; right-hand lobe on initial ray, other two every 120 degree, given by r = a cos 3θ.},scale=.9,>=stealth]
	\draw [<->,] (-2.1,0) -- (2.1,0);
	\draw [<->,] (0,-2.1) -- (0,2.1);
	\draw [thick,draw={\colorone},domain=0:360,samples=60,smooth] plot
	 ({cos(\x)*(1.9*cos(3*\x))},{sin(\x)*(1.9*cos(3*\x))});
\end{tikzpicture}		
&
\begin{tikzpicture}[alt={Three-petal rose; top lobe on positive y-axis, others every 120 deg, from r = a sin 3θ.},scale=.9,>=stealth]
	\draw [<->,] (-2.1,0) -- (2.1,0);
	\draw [<->,] (0,-2.1) -- (0,2.1);
	\draw [thick,draw={\colorone},domain=0:360,samples=60,smooth] plot
	 ({cos(\x)*(1.9*sin(3*\x))},{sin(\x)*(1.9*sin(3*\x))});
\end{tikzpicture}		
\bigskip\\
%
\multicolumn{4}{l}{\bfseries\large Special Curves} \\
\multicolumn{2}{l}{\bfseries Rose curves} & \textbf{Lemniscate:} & \textbf{Eight Curve:} \smallskip\\
$r=a\sin (\theta/5)$ & $r=a\sin(2\theta/5)$ &
$r^2=a^2\cos (2\theta)$ & $r^2=a^2\sec^4\theta\cos (2\theta)$ \medskip\\
\begin{tikzpicture}[alt={Five-arm spiral-rose: small central loop then larger coils, all from r = a sin (θ / 5).},scale=.9,>=stealth]
	\draw [<->,] (-2.1,0) -- (2.1,0);
	\draw [<->,] (0,-2.1) -- (0,2.1);
	\draw [thick,draw={\colorone},domain=0:900,samples=100,smooth] plot
	 ({cos(\x)*(1.9*sin(\x/5))},{sin(\x)*(1.9*sin(\x/5))});
\end{tikzpicture}		
&
\begin{tikzpicture}[alt={Nested rose with ten slender lobes arranged in two rings, produced by r = a sin (2θ / 5).},scale=.9,>=stealth]
	\draw [<->,] (-2.1,0) -- (2.1,0);
	\draw [<->,] (0,-2.1) -- (0,2.1);
	\draw [thick,draw={\colorone},domain=0:1800,samples=200,smooth] plot
	 ({cos(\x)*(1.9*cos(2*\x/5))},{sin(\x)*(1.9*cos(2*\x/5))});
\end{tikzpicture}		
&
\begin{tikzpicture}[alt={Horizontal figure-eight centered at origin, each lobe length √2.},scale=.9,>=stealth]
	\draw [<->,] (-2.1,0) -- (2.1,0);
	\draw [<->,] (0,-2.1) -- (0,2.1);
	\draw [thick,draw={\colorone},domain=-45:45,samples=60,smooth] plot
	 ({cos(\x)*(1.9*sqrt(cos(2*\x)))},{sin(\x)*(1.9*sqrt(cos(2*\x)))});
	\draw [thick,draw={\colorone},domain=-45:45,samples=60,smooth] plot
	 ({-cos(\x)*(1.9*sqrt(cos(2*\x)))},{-sin(\x)*(1.9*sqrt(cos(2*\x)))});
\end{tikzpicture}		
&
\begin{tikzpicture}[alt={Slim vertical figure-eight; wider loops meet at origin, axes of symmetry both ways.},scale=.9,>=stealth]
	\draw [<->,] (-2.1,0) -- (2.1,0);
	\draw [<->,] (0,-2.1) -- (0,2.1);
	\draw [thick,draw={\colorone},domain=-45:45,samples=60,smooth] plot
	 ({cos(\x)*1.9*sec(\x)*sec(\x)*sqrt(cos(2*\x))},
	  {sin(\x)*1.9*sec(\x)*sec(\x)*sqrt(cos(2*\x))});
	\draw [thick,draw={\colorone},domain=-45:45,samples=60,smooth] plot
	 ({-cos(\x)*1.9*sec(\x)*sec(\x)*sqrt(cos(2*\x))},
	  {-sin(\x)*1.9*sec(\x)*sec(\x)*sqrt(cos(2*\x))});
\end{tikzpicture}
%
\end{tabular}%}

\restoregeometry

Earlier we discussed how each point in the plane does not have a unique representation in polar form. This can be a ``good'' thing, as it allows for the beautiful and interesting curves seen in the preceding gallery. However, it can also be a ``bad'' thing, as it can be difficult to determine where two curves intersect.

\begin{example}[Finding points of intersection with polar curves]\label{ex_polar7}%
Determine where the graphs of the polar equations $r=1+3\cos\theta$ and $r=\cos \theta$ intersect.
\solution
As technology is generally readily available, it is usually a good idea to start with a graph. We have graphed the two functions in \autoref{fig:polar7}(a); to better discern the intersection points, part (b) of the figure zooms in around the origin.
\mtable{Graphs to help determine the points of intersection of the polar functions given in \autoref{ex_polar7}.}{fig:polar7}{%
\begin{tikzpicture}[alt={Blue two-turn spiral; red inner loop; both start at origin and cross initial ray.}]
\begin{axis}[width=1.16\marginparwidth,tick label style={font=\scriptsize},
axis y line=middle,axis x line=middle,name=myplot,minor x tick num=1,
minor y tick num=1,ymin=-2.6,ymax=2.6,xmin=-1.5,xmax=4.5,axis equal]
\addplot [draw={\colorone},thick, smooth,domain=0:360,samples=60]
 ({cos(x)*(1+3*cos(x))},{sin(x)*(1+3*cos(x))});
\draw[draw={\colortwo},thick](axis cs:.5,0)circle(.5);
\end{axis}
\node [right] at (myplot.right of origin) {\scriptsize $0$};
\node [above] at (myplot.above origin) {\scriptsize $\pi/2$};
\end{tikzpicture}
\\(a)\\
\begin{tikzpicture}[alt={Zoom view: red right-curving cardioid segment and blue concave left curve meet at origin and intersect twice more.}]
\begin{axis}[width=1.16\marginparwidth,tick label style={font=\scriptsize},
axis y line=middle,axis x line=middle,name=myplot,minor x tick num=1,
ymin=-.6,ymax=.6,xmin=-.72,xmax=.72,axis equal]
\addplot [draw={\colorone},thick, smooth,domain=0:360,samples=60]
 ({cos(x)*(1+3*cos(x))},{sin(x)*(1+3*cos(x))});
\draw[draw={\colortwo},thick](axis cs:.5,0)circle(.5);
\end{axis}
\node [right] at (myplot.right of origin) {\scriptsize $0$};
\node [above] at (myplot.above origin) {\scriptsize $\pi/2$};
\end{tikzpicture}
\\(b)}
We start by setting the two functions equal to each other and solving for $\theta$:
\begin{align*}
1+3\cos\theta &= \cos \theta \\
2\cos\theta &= -1\\
\cos\theta&= -\frac12\\
\theta &= \frac{2\pi}{3}, \frac{4\pi}{3}.
\end{align*}
(There are, of course, infinite solutions to the equation $\cos\theta=-1/2$; as the limaçon is traced out once on $[0,2\pi]$, we restrict our solutions to this interval.) 

We need to analyze this solution. When $\theta = 2\pi/3$ we obtain the point of intersection that lies in the 4\textsuperscript{th} quadrant. When $\theta = 4\pi/3$, we get the point of intersection that lies in the 1\textsuperscript{st} quadrant. There is more to say about this second intersection point, however. The circle defined by $r=\cos\theta$ is traced out once on $[0,\pi]$, meaning that this point of intersection occurs while tracing out the circle a second time. It seems strange to pass by the point once and then recognize it as a point of intersection only when arriving there a ``second time.'' The first time the circle arrives at this point is when $\theta = \pi/3$.
It is key to understand that these two points are the same: $(\cos \pi/3,\pi/3)$ and $(\cos 4\pi/3,4\pi/3)$. 

To summarize what we have done so far, we have found two points of intersection: when $\theta=2\pi/3$ and when $\theta=4\pi/3$. When referencing the circle $r=\cos \theta$, the latter point is better referenced as when $\theta=\pi/3$.

There is yet another point of intersection: the pole (or, the origin). We did not recognize this intersection point using our work above as each graph arrives at the pole at a different $\theta$ value.

A graph intersects the pole when $r=0$. Considering the circle $r=\cos\theta$, $r=0$ when $\theta = \pi/2$ (and odd multiples thereof, as the circle is repeatedly traced). The limaçon intersects the pole when $1+3\cos\theta =0$; this occurs when $\cos \theta = -1/3$, or for $\theta = \cos^{-1}(-1/3)$. This is a nonstandard angle, approximately $\theta = 1.9106\text{ radians} \approx 109.47^\circ$. The limaçon intersects the pole twice in $[0,2\pi]$; the other angle at which the limaçon is at the pole is the reflection of the first angle across the $x$-axis. That is, $\theta = 4.3726 \approx 250.53^\circ$.
\end{example}

If all one is concerned with is the $(x,y)$ coordinates at which the graphs intersect, much of the above work is extraneous. We know they intersect at $(0,0)$; we might not care at what $\theta$ value. Likewise, using $\theta =2\pi/3$ and $\theta=4\pi/3$ can give us the needed rectangular coordinates. However, in the next section we apply calculus concepts to polar functions. When computing the area of a region bounded by polar curves, understanding the nuances of the points of intersection becomes important.

\printexercises{exercises/09-04-exercises}

\section{Calculus and Polar Functions} \label{sec:polarcalc}

The previous section defined polar coordinates, leading to polar functions. We investigated plotting these functions and solving a fundamental question about their graphs, namely, where do two polar graphs intersect?

We now turn our attention to answering other questions, whose solutions require the use of calculus. A basis for much of what is done in this section is the ability to turn a polar function $r=f(\theta)$ into a set of parametric equations. Using the identities $x=r\cos \theta$ and $y=r\sin \theta$, we can create the parametric equations $x=f(\theta)\cos\theta$, $y=f(\theta)\sin\theta$ and apply the concepts of \autoref{sec:par_calc}.

\subsection{Polar Functions and \texorpdfstring{$\dfrac{\dd y}{\dd x}$}{dy/dx}}

We are interested in the lines tangent to a given graph, regardless of whether that graph is produced by rectangular, parametric, or polar equations. In each of these contexts, the slope of the tangent line is $\frac{\dd y}{\dd x}$. Given $r=f(\theta)$, we are generally \emph{not} concerned with $r\,'=\fp(\theta)$; that describes how fast $r$ changes with respect to $\theta$. Instead, we will use $x=f(\theta)\cos\theta$, $y=f(\theta)\sin\theta$ to compute $\frac{\dd y}{\dd x}$. 

Using \autoref{idea:dydxpar} we have
\[\frac{\dd y}{\dd x} = \frac{\dd y}{\dd\theta}\Big/\frac{\dd x}{\dd\theta}.\]
Each of the two derivatives on the right hand side of the equality requires the use of the Product Rule. We state the important result as a Key Idea.

\begin{keyidea}[Finding $\frac{\dd y}{\dd x}$ with Polar Functions]\label{idea:dydxpol}%
Let $r=f(\theta)$ be a polar function. With $x=f(\theta)\cos\theta$ and $y=f(\theta)\sin\theta$,
\index{polar!functions!finding $\frac{\dd y}{\dd x}$}\index{tangent line}
\[
 \frac{\dd y}{\dd x}
 =\frac{\frac{\dd y}{\dd\theta}}{\frac{\dd x}{\dd\theta}}
 =\frac{\fp(\theta)\sin\theta+f(\theta)\cos\theta}{\fp(\theta)\cos\theta-f(\theta)\sin\theta}.
\]
\end{keyidea}

\youtubeVideo{QTa9OZ4iGPo}{The Slope of Tangent Lines to Polar Curves}

% at some point, someone said `` answer should be ``-2'' not ``-1''. ''  whatever that means
\begin{example}[Finding $\frac{\dd y}{\dd x}$ with polar functions.]\label{ex_polcalc1}%
Consider the limaçon $r=1+2\sin\theta$ on $[0,2\pi]$.
\begin{enumerate}
	\item	Find the rectangular equations of the tangent and normal lines to the graph at $\theta=\pi/4$.
	\item	Find where the graph has vertical and horizontal tangent lines.
\end{enumerate}
\solution
\begin{enumerate}
	\item We start by computing $\frac{\dd y}{\dd x}$. With $\fp(\theta) = 2\cos\theta$, we have
	\begin{align*}
	\frac{\dd y}{\dd x} &= \frac{2\cos\theta\sin\theta + \cos\theta(1+2\sin\theta)}{2\cos^2\theta-\sin\theta(1+2\sin\theta)}\\
	&= \frac{\cos\theta(4\sin\theta+1)}{2(\cos^2\theta-\sin^2\theta)-\sin\theta}.
	\end{align*}
	When $\theta=\pi/4$, $\frac{\dd y}{\dd x}=-2\sqrt{2}-1$ (this requires a bit of simplification). In rectangular coordinates, the point on the graph at $\theta=\pi/4$ is $(1+\sqrt{2}/2,1+\sqrt{2}/2)$. Thus the rectangular equation of the line tangent to the limaçon at $\theta=\pi/4$ is 
\[y=(-2\sqrt{2}-1)\bigl(x-(1+\sqrt{2}/2)\bigr)+1+\sqrt{2}/2 \approx  -3.83 x+8.24.\]
The limaçon and the tangent line are graphed in \autoref{fig:polcalc1}. 
	
	The normal line has the opposite-reciprocal slope as the tangent line, so its equation is 
\[y \approx \frac{1}{3.83}x+1.26.\]
	
	\item		To find the horizontal lines of tangency, we find where $\frac{\dd y}{\dd x}=0$ (when the denominator does not equal 0); thus we find where the numerator of our equation for $\frac{\dd y}{\dd x}$ is 0.
	\[\cos\theta(4\sin\theta+1)=0\quad \Rightarrow \quad \cos\theta=0 \quad \text{or}\quad 4\sin\theta+1=0.\]
	On $[0,2\pi]$, $\cos\theta=0$ when $\theta=\pi/2,\ 3\pi/2$. 

Setting $4\sin\theta+1=0$ gives $\theta=\sin^{-1}(-1/4)\approx -0.2527 = -14.48^\circ$. We want the results in $[0,2\pi]$; we also recognize there are two solutions, one in the 3\textsuperscript{rd} quadrant and one in the 4\textsuperscript{th}. Using reference angles, we have our two solutions as $\theta =3.39$ and $6.03$ radians. The four points we obtained where the limaçon has a horizontal tangent line are given in \autoref{fig:polcalc1} with black-filled dots.\bigskip

To find the vertical lines of tangency, we determine where $\frac{\dd y}{\dd x}$ is undefined by setting the denominator of $\frac{\dd y}{\dd x}=0$ (when the numerator does not equal 0).
\begin{align*}
2(\cos^2\theta -\sin^2\theta)-\sin\theta &= 0 .
\intertext{Convert the $\cos^2\theta$ term to $1-\sin^2\theta$:}
2(1-\sin^2\theta-\sin^2\theta)-\sin\theta &= 0\\
4\sin^2\theta + \sin\theta -2 &= 0.
\intertext{Recognize this as a quadratic in the variable $\sin\theta$. Using the quadratic formula, we have}
\sin\theta &= \frac{-1\pm\sqrt{33}}{8}.
\end{align*}
%
\mtable{The limaçon in \autoref{ex_polcalc1} with its tangent line at $\theta=\pi/4$ and points of vertical and horizontal tangency.}{fig:polcalc1}{\begin{tikzpicture}[alt={Blue limacon with inner loop; red tangent at θ = π/4; filled/open dots show four vertical or horizontal tangent points.}]
\begin{axis}[width=1.16\marginparwidth,tick label style={font=\scriptsize},
axis y line=middle,axis x line=middle,name=myplot,minor x tick num=1,
ymin=-.35,ymax=3.2,xmin=-2.1,xmax=2.1]
\addplot [draw={\colorone},thick, smooth,domain=0:360,samples=60]
 ({cos(x)*(1+2*sin(x))},{sin(x)*(1+2*sin(x))});
\addplot [draw={\colortwo},thick, smooth,domain=0.1:2,samples=2] {-3.82843*x+8.24264};
\filldraw (axis cs: 0,3) circle (1.5pt)
		(axis cs: 0,1) circle (1.5pt)
		(axis cs: -.493,-0.125) circle (1.5pt)
		(axis cs: .483,-0.125) circle (1.5pt);
\filldraw [fill=white,draw=black,thick] (axis cs: 1.76011, 1.31048) circle (1.5pt)
		(axis cs: -1.76011, 1.31048) circle (1.5pt)
		(axis cs: 0.368977, 0.572639) circle (1.5pt)
		(axis cs: -0.369009, 0.578743) circle (1.5pt);
\end{axis}
\node [right] at (myplot.right of origin) {\scriptsize $0$};
\node [above] at (myplot.above origin) {\scriptsize $\pi/2$};
\end{tikzpicture}}%
%
We solve $\sin\theta = \frac{-1+\sqrt{33}}8$ and $\sin\theta = \frac{-1-\sqrt{33}}8$:
\begin{align*}
\sin\theta &=\frac{-1+\sqrt{33}}8 & \sin\theta &= \frac{-1-\sqrt{33}}{8}\\
\theta &= \sin^{-1}\left(\frac{-1+\sqrt{33}}8\right) & \theta &= \sin^{-1}\left(\frac{-1-\sqrt{33}}8\right)\\
\theta &\approx 0.6349 & \theta &\approx -1.0030
\end{align*}
In each of the solutions above, we only get one of the possible two solutions as $\sin^{-1}x$ only returns solutions in $[-\pi/2,\pi/2]$, the 4\textsuperscript{th} and 1\textsuperscript{st} quadrants. Again using reference angles, we have:
\[
\sin\theta = \frac{-1+\sqrt{33}}8 \quad \Rightarrow \quad
\theta \approx 0.6349,\ 2.5067 \text{ radians}
\]
and 
\[
\sin\theta = \frac{-1-\sqrt{33}}8 \quad \Rightarrow \quad
\theta \approx 4.1446,\ 5.2802 \text{ radians.}
\]
These points are also shown in \autoref{fig:polcalc1} with white-filled dots.
\end{enumerate}
\end{example}

When the graph of the polar function $r=f(\theta)$ intersects the pole, it means that $f(\alpha) = 0$ for some angle $\alpha$. Making this substitution in the formula for $\dfrac{\dd y}{\dd x}$ given in \autoref{idea:dydxpol} we see
\[
 \frac{\dd y}{\dd x}
 =\frac{\fp(\alpha)\sin\alpha+f(\alpha)\cos\alpha}{\fp(\alpha)\cos\alpha-f(\alpha)\sin\alpha}
 =\frac{\sin\alpha}{\cos\alpha}
 =\tan\alpha.
\]
%
%When the graph of the polar function $r=f(\theta)$ intersects the pole, it means that $f(\alpha)=0$ for some angle $\alpha$. Thus the formula for $\frac{\dd y}{\dd x}$ in such instances is very simple, reducing simply to \[\frac{\dd y}{\dd x} = \tan \alpha.\]
This equation makes an interesting point. It tells us the slope of the tangent line at the pole is $\tan \alpha$; some of our previous work (see, for instance, \autoref{ex_polar3}) shows us that the line through the pole with slope $\tan \alpha$ has polar equation $\theta=\alpha$. Thus when a polar graph touches the pole at $\theta=\alpha$, the equation of the tangent line at the pole is $\theta=\alpha$.

\begin{example}[Finding tangent lines at the pole]\label{ex_polcalc2}%
Let $r=1+2\sin\theta$, a limaçon. Find the equations of the lines tangent to the graph at the pole.
%
\mtable{Graphing the tangent lines at the pole in \autoref{ex_polcalc2}.}{fig:polcalc2}{\begin{tikzpicture}[alt={Blue looped curve through origin; four symmetric red lines through the pole mark its tangent directions.}]
\begin{axis}[width=1.16\marginparwidth,tick label style={font=\scriptsize},
axis y line=middle,axis x line=middle,name=myplot,
ymin=-.5,ymax=1.37,xmin=-1.1,xmax=1.1]
\addplot [draw={\colorone},thick, smooth,domain=0:360,samples=60]
 ({cos(x)*(1+2*sin(x))},{sin(x)*(1+2*sin(x))});
\addplot [draw={\colortwo},thick, smooth,domain=-1:1,samples=2] {x/sqrt(3)};
\addplot [draw={\colortwo},thick, smooth,domain=-1:1,samples=2] {-x/sqrt(3)};
\end{axis}
\node [right] at (myplot.right of origin) {\scriptsize $0$};
\node [above] at (myplot.above origin) {\scriptsize $\pi/2$};
\end{tikzpicture}}
\solution
We need to know when $r=0$. 
\begin{align*}
1+2\sin\theta &= 0\\
\sin\theta &= -1/2\\
\theta &= \frac{7\pi}{6},\ \frac{11\pi}6.
\end{align*}
Thus the equations of the tangent lines, in polar coordinates, are $\theta = 7\pi/6$ and $\theta = 11\pi/6$. In rectangular form, the tangent lines are $y=\tan(7\pi/6)x=\frac x{\sqrt3}$ and $y=\tan(11\pi/6)x=-\frac x{\sqrt3}$. The full limaçon can be seen in \autoref{fig:polcalc1}; we zoom in on the tangent lines in \autoref{fig:polcalc2}.
\end{example}

\subsection{Area}

When using rectangular coordinates, the equations $x=h$ and $y=k$ defined vertical and horizontal lines, respectively, and combinations of these lines create rectangles (hence the name ``rectangular coordinates''). It is then somewhat natural to use rectangles to approximate area as we did when learning about the definite integral.\index{polar!functions!area}

\mnote{\textbf{Note:} Recall that the area of a sector of a circle with radius $r$ subtended by an angle $\theta$ is $A = \frac12\theta r^2$.
\[
\begin{tikzpicture}[alt={Circular sector with radius r and angle θ; accompanying note states its area A = θ r² /2.},x=30pt,y=30pt,thick]
			\draw (2,0) arc (0:50:2) -- (0,0);
			\draw [] (0,0) -- (2,0) node [pos=.5,below] {$r$};
			\draw [fill=black] (0,0) circle (1pt);
			%\draw (1.95,1.0) node {$s$};
			\draw (0,0) node [shift={(15pt,8pt)}] {$\theta$};
			\end{tikzpicture}
\]}

When using polar coordinates, the equations $\theta=\alpha$ and $r=c$ form lines through the origin and circles centered at the origin, respectively, and combinations of these curves form sectors of circles. It is then somewhat natural to calculate the area of regions defined by polar functions by first approximating with sectors of circles. 

Consider \autoref{fig:polararea} (a) where a region defined by $r=f(\theta)$ on $[\alpha,\beta]$ is given. (Note how the ``sides'' of the region are the lines $\theta=\alpha$ and $\theta=\beta$, whereas in rectangular coordinates the ``sides'' of regions were often the vertical lines $x=a$ and $x=b$.)

Partition the interval $[\alpha,\beta]$ into $n$ equally spaced subintervals as $\alpha=\theta_0<\theta_1<\dotso<\theta_n=\beta$. The radian length of each subinterval is $\Delta\theta = (\beta-\alpha)/n$, representing a small change in angle. The area of the region defined by the $i^{\text{th}}$ subinterval $[\theta_{i-1},\theta_i]$ can be approximated with a sector of a circle with radius $f(c_i)$, for some $c_i$ in $[\theta_{i-1},\theta_i]$. The area of this sector is $\frac12[f(c_i)]^2\Delta\theta$. This is shown in part (b) of the figure, where $[\alpha,\beta]$ has been divided into 4 subintervals. We approximate the area of the whole region by summing the areas of all sectors:
%
\mtable{Computing the area of a polar region.}{fig:polararea}{%
\begin{tikzpicture}[alt={Blue r = f(θ) with dashed spokes; shaded wedge from θ = α to β.}]
\begin{axis}[width=1.16\marginparwidth,tick label style={font=\scriptsize},
axis y line=middle,axis x line=middle,name=myplot,
ymin=-.1,ymax=1.1,xmin=-.1,xmax=1.1,axis equal]
\addplot
 [draw={\colorone},fill={\coloronefill},area style, smooth,domain=18:72,samples=30]
% (axis cs:x:{1+.05*cos(9*x)})
 ({cos(x)*(1+.05*cos(9*x))},{sin(x)*(1+.05*cos(9*x))})
  -- (axis cs:0,0) -- cycle;
\addplot [draw={\colorone},thick, smooth,domain=0:90,samples=30]
 ({cos(x)*(1+.05*cos(9*x))},{sin(x)*(1+.05*cos(9*x))});
\draw [thick,draw={\colorone},] (axis cs:0,0) -- (axis cs: 0.905831, 0.294322)
 node [pos=.7,below,rotate=18,black] {\scriptsize $\theta=\alpha$};
\draw [thick,draw={\colorone},] (axis cs:0,0) -- (axis cs:0.313792, 0.965751)
 node [pos=.7,above,rotate=72,black] {\scriptsize $\theta=\beta$};
\draw [thick,draw={\colorone},dashed]
 (axis cs:0,0) -- (axis cs: 0.862592, 0.528597)
  (axis cs:0,0) -- (axis cs: 0.732107, 0.732107)
   (axis cs:0,0) -- (axis cs: 0.497095, 0.811186);
\draw (axis cs:.8,.85) node {\scriptsize $r=f(\theta)$};
\end{axis}
\node [right] at (myplot.right of origin) {\scriptsize $0$};
\node [above] at (myplot.above origin) {\scriptsize $\pi/2$};
\end{tikzpicture}
\\(a)\\
\begin{tikzpicture}[alt={Red r = f(θ) with dashed spokes; shaded wedge from θ = α to β}]
\begin{axis}[width=1.16\marginparwidth,tick label style={font=\scriptsize},
axis y line=middle,axis x line=middle,name=myplot,
ymin=-.1,ymax=1.1,xmin=-.1,xmax=1.1,axis equal]
\addplot
 [draw={\coloronefill},fill={\coloronefill},area style, smooth,domain=18:72,samples=30]
 ({cos(x)*(1+.05*cos(9*x))},{sin(x)*(1+.05*cos(9*x))}) -- (axis cs:0,0) -- cycle;
\draw [thick,draw={\colorone},] (axis cs:0,0) -- (axis cs:0.313792, 0.965751)
 node [pos=.7,above,rotate=72,black] {\scriptsize $\theta=\beta$};
\addplot [draw={\colortwo},fill={\colortwofill},thick, smooth,domain=18:31.5,samples=30]
 ({.96*cos(x)},{.96*sin(x)}) -- (axis cs:0,0) -- cycle;
\addplot [draw={\colortwo},fill={\colortwofill},thick, smooth,domain=31.5:45,samples=30]
 ({1.05*cos(x)},{1.05*sin(x)}) -- (axis cs:0,0) -- cycle;
\addplot [draw={\colortwo},fill={\colortwofill},thick, smooth,domain=45:58.5,samples=30]
 ({1.0*cos(x)},{1*sin(x)}) -- (axis cs:0,0) -- cycle;
\addplot [draw={\colortwo},fill={\colortwofill},thick, smooth,domain=58.5:72,samples=30]
 ({.96*cos(x)},{.96*sin(x)}) -- (axis cs:0,0) -- cycle;
\draw (axis cs:.8,.85) node {\scriptsize $r=f(\theta)$};
\draw [thick,draw={\colortwo},] (axis cs:0,0) -- (axis cs: 0.905831, 0.294322)
 node [pos=.7,below,rotate=18,black] {\scriptsize $\theta=\alpha$};
\addplot [draw={\colorone},thick, smooth,domain=0:90,samples=30]
 ({cos(x)*(1+.05*cos(9*x))},{sin(x)*(1+.05*cos(9*x))});
\end{axis}
\node [right] at (myplot.right of origin) {\scriptsize $0$};
\node [above] at (myplot.above origin) {\scriptsize $\pi/2$};
\end{tikzpicture}
\\(b)}%
%
\[\text{Area} \approx \sum_{i=1}^n \frac12[f(c_i)]^2\Delta\theta.\]
This is a Riemann sum. By taking the limit of the sum as $n\to\infty$, we find the exact area of the region in the form of a definite integral.

\begin{theorem}[Area of a Polar Region]\label{thm:polar_area}%
Let $f$ be continuous and non-negative on $[\alpha,\beta]$, where $0\leq \beta-\alpha\leq 2\pi$. The area  $A$ of the region bounded by the curve $r=f(\theta)$ and the lines $\theta=\alpha$ and $\theta=\beta$ is 
\[
A=\frac12\int_\alpha^\beta[f(\theta)]^2 \dd\theta
=\frac12\int_\alpha^\beta r^{\,2} \dd\theta
\]
\end{theorem}

The theorem states that $0\leq \beta-\alpha\leq 2\pi$. This ensures that region does not overlap itself, which would give a result that does not correspond directly to the area.

\begin{example}[Area of a polar region]\label{ex_polcalc3}%
Find the area of the circle defined by $r=\cos \theta$. (Recall this circle has radius $1/2$.)
\solution
This is a direct application of \autoref{thm:polar_area}. The circle is traced out on $[0,\pi]$, leading to the integral
%
\mnote{\textbf{Note:} \autoref{ex_polcalc3} requires the use of the integral $\ds\int \cos^2\theta\dd\theta$. This is handled well by using the half angle formula as found in the back of this text. Due to the nature of the area formula, integrating $\cos^2\theta$ and $\sin^2\theta$ is required often. We offer here these indefinite integrals as a time-saving measure.
\begin{align*}
	\int\cos^2\theta\dd\theta &= \frac12\theta+\frac14\sin(2\theta)+C\\
	\int\sin^2\theta\dd\theta &= \frac12\theta-\frac14\sin(2\theta)+C
\end{align*}}
%
\begin{align*}
	\text{Area} &= \frac12\int_0^\pi \cos^2\theta\dd\theta \\
	&= \frac12\int_0^\pi \frac{1+\cos(2\theta)}{2}\dd\theta\\
	&= \frac14\bigl(\theta +\frac12\sin(2\theta)\bigr)\Bigg|_0^\pi\\
	&= \frac\pi4.
\end{align*}
Of course, we already knew the area of a circle with radius $1/2$. We did this example to demonstrate that the area formula is correct.
\end{example}

\begin{example}[Area of a polar region]\label{ex_polcalc4}%
Find the area of the cardioid $r=1+\cos\theta$ bound between $\theta=\pi/6$ and $\theta=\pi/3$, as shown in \autoref{fig:polcalc4}.
\solution
This is again a direct application of \autoref{thm:polar_area}.
%
\mtable{Finding the area of the shaded region of a cardioid in \autoref{ex_polcalc4}.}{fig:polcalc4}{\begin{tikzpicture}[alt={Cardioid sector: blue curve r = f(θ) with shaded slice between θ = 0 and π/6, bounded by two radii.}]
\begin{axis}[width=1.16\marginparwidth,tick label style={font=\scriptsize},
axis y line=middle,axis x line=middle,name=myplot,ytick={1},
ymin=-.7,ymax=1.5,xmin=-.5,xmax=2.1]
\addplot
 [draw={\coloronefill},fill={\coloronefill},area style,smooth,domain=30:60,samples=30]
 ({cos(x)*(1+cos(x))},{sin(x)*(1+cos(x))}) -- (axis cs:0,0) -- cycle;
\addplot [draw={\colorone},thick, smooth,domain=0:360,samples=60]
 ({cos(x)*(1+cos(x))},{sin(x)*(1+cos(x))});
\draw [thick,draw={\colorone}] (axis cs:0,0) -- (axis cs: 1.61603, 0.933013)
 node [pos=.6,below,rotate=30,black] {\scriptsize $\theta=\pi/6$};
\draw [thick,draw={\colorone}] (axis cs:0,0) -- (axis cs:0.75, 1.29904)
 node [pos=.6,above,rotate=60,black] {\scriptsize $\theta=\pi/3$};
\end{axis}
\node [right] at (myplot.right of origin) {\scriptsize $0$};
\node [above] at (myplot.above origin) {\scriptsize $\pi/2$};
\end{tikzpicture}}
%
\begin{align*}
	\text{Area}
	&= \frac12\int_{\pi/6}^{\pi/3} (1+\cos\theta)^2\dd\theta\\
	&= \frac12\int_{\pi/6}^{\pi/3} (1+2\cos\theta+\cos^2\theta)\dd\theta\\
	&= \frac12\left[\theta+2\sin\theta+\frac12\theta
	+\frac14\sin(2\theta)\right]_{\pi/6}^{\pi/3} \\
	&= \frac18\bigl(\pi+4\sqrt{3}-4\bigr).
\end{align*}
\end{example}

\subsubsection{Area Between Curves}

Our study of area in the context of rectangular functions led naturally to finding area bounded between curves. We consider the same in the context of polar functions. \index{polar!functions!area between curves}

\mtable{Illustrating area bound between two polar curves.}{fig:polarea3}{\begin{tikzpicture}[alt={Blue outer r_2 and red inner r_1 enclose shaded region between θ = α and β, inner arc from r_1 to pole.}]
\begin{axis}[width=1.16\marginparwidth,tick label style={font=\scriptsize},
axis y line=middle,axis x line=middle,name=myplot,
ymin=-.1,ymax=1.1,xmin=-.1,xmax=1.1]
\addplot [draw={\colortwo}, thick,smooth,domain=0:90,samples=30]
 ({cos(x)*(.7+.05*sin(6*x))},{sin(x)*(.7+.05*sin(6*x))});
\draw (axis cs:.8,.85) node {\scriptsize $r_2=f_2(\theta)$}
      (axis cs:.2,.8) node {\scriptsize $r_1=f_1(\theta)$};
\addplot [draw={\colorone},thick, smooth,domain=0:90,samples=30]
 ({cos(x)*(1+.05*cos(9*x))},{sin(x)*(1+.05*cos(9*x))});
\addplot[draw={\colorone},fill={\coloronefill},thick, smooth,domain=30:60,samples=30]
 ({cos(x)*(1+.05*cos(9*x))},{sin(x)*(1+.05*cos(9*x))})--(axis cs:0,0)--cycle;
\addplot[draw={\colortwo},fill=white,thick,smooth,domain=30:60,samples=30]
 ({cos(x)*(.7+.05*sin(6*x))},{sin(x)*(.7+.05*sin(6*x))})--(axis cs:0,0)--cycle;
\draw [thick,draw={\colorone},] (axis cs:0,0) -- (axis cs: 0.866025, 0.5)
 node [pos=.4,below,rotate=30,black] {\scriptsize $\theta=\alpha$};
\draw [thick,draw={\colorone},] (axis cs:0,0) -- (axis cs:0.475, 0.822724)
 node [pos=.4,above,rotate=60,black] {\scriptsize $\theta=\beta$};
\end{axis}
\node [right] at (myplot.right of origin) {\scriptsize $0$};
\node [above] at (myplot.above origin) {\scriptsize $\pi/2$};
\end{tikzpicture}}

Consider the shaded region shown in \autoref{fig:polarea3}. We can find the area of this region by computing the area bounded by $r_2=f_2(\theta)$ and subtracting the area bounded by $r_1=f_1(\theta)$ on $[\alpha,\beta]$. Thus
\[
\text{Area}
\ = \ \frac12\int_\alpha^\beta r_2^{\,2}\dd\theta
- \frac12\int_\alpha^\beta r_1^{\,2}\dd\theta
= \frac12\int_\alpha^\beta \bigl(r_2^{\,2}-r_1^{\,2}\bigr)\dd\theta.
\]

\begin{keyidea}[Area Between Polar Curves]\label{idea:area_between_polar}%
The area $A$ of the region bounded by $r_1=f_1(\theta)$ and $r_2=f_2(\theta)$, $\theta=\alpha$ and $\theta=\beta$, where $0\le f_1(\theta)\leq f_2(\theta)$ on $[\alpha,\beta]$, is
\[
A = \frac{1}{2} \int_\alpha^\beta [f_2(\theta)]^2 - [f_1(\theta)]^2\dd\theta
= \frac12\int_\alpha^\beta \bigl(r_2^{\,2}-r_1^{\,2}\bigr)\dd\theta.
\]
\end{keyidea}

\mtable{Finding the area between polar curves in \autoref{ex_polcalc5}.}{fig:polcalc5}{\begin{tikzpicture}[alt={Blue circle r = 1 intersects red cardioid r = 1 − cos θ; right‐hand lens (−π/2 ≤ θ ≤ π/2) is shaded.}]
\begin{axis}[width=1.16\marginparwidth,tick label style={font=\scriptsize},
axis y line=middle,axis x line=middle,name=myplot,axis on top,
ymin=-1.6,ymax=1.6,xmin=-.6,xmax=3.2]
\draw [fill={\coloronefill},draw={\colorone},thick, smooth,domain=0:180,samples=60]  (axis cs:1.5,0)circle(1.5);
% fill=white to cover up what we don't want
\addplot [fill=white,draw={\colortwo}, thick,smooth,domain=0:360,samples=60]
 ({cos(x)*(1+cos(x))},{sin(x)*(1+cos(x))});
% redraw the circle for the part that was covered
\draw [draw={\colorone},thick, smooth,domain=0:180,samples=60]
 (axis cs:1.5,0)circle(1.5);
\end{axis}
\node [right] at (myplot.right of origin) {\scriptsize $0$};
\node [above] at (myplot.above origin) {\scriptsize $\pi/2$};
\end{tikzpicture}}

\begin{example}[Area between polar curves]\label{ex_polcalc5}%
Find the area bounded between the curves $r=1+\cos \theta$ and $r=3\cos\theta$, as shown in \autoref{fig:polcalc5}.
\solution
We need to find the points of intersection between these two functions. Setting them equal to each other, we find:
\begin{align*}
1+\cos\theta &= 3\cos \theta \\
 \cos\theta &=1/2\\
\theta &= \pm \pi/3
\end{align*}
Thus we integrate $\frac12\bigl((3\cos\theta)^2-(1+\cos\theta)^2\bigr)$ on $[-\pi/3,\pi/3]$.
\begin{align*}
	\text{Area}
	&= \frac12\int_{-\pi/3}^{\pi/3} \bigl((3\cos\theta)^2-(1+\cos\theta)^2\bigr)\dd\theta\\
	&= \frac12\int_{-\pi/3}^{\pi/3} \bigl( 8\cos^2\theta-2\cos\theta-1\bigr)\dd\theta \\
	&= \frac12\bigl(2\sin(2\theta) - 2\sin\theta+3\theta\bigr)\Bigg|_{-\pi/3}^{\pi/3}\\
	&= \pi.
\end{align*}
Amazingly enough, the area between these curves has a ``nice'' value.
\end{example}

\begin{example}[Area defined by polar curves]\label{ex_polcalc6}%
Find the area bounded between the polar curves $r=1$ and one petal of $r=2\cos(2\theta)$ where $y>0$, as shown in \autoref{fig:polcalc6}(a).
\solution
%
\mtable{Graphing the region bounded by the functions in \autoref{ex_polcalc6}.}{fig:polcalc6}{%
\begin{tikzpicture}[alt={Blue rose r = sin 4 θ overlaps red three‐leaf r = sin 2 θ; intersection rays at θ = ± π/6, ± π/2.}]
\begin{axis}[width=\marginparwidth,tick label style={font=\scriptsize},
axis y line=middle,axis x line=middle,name=myplot,axis on top,
ymin=-1.15,ymax=1.15,xmin=-.5,xmax=2.2]
\addplot[draw={\coloronefill},fill={\coloronefill},area style,domain=30:45]
 ({cos(x)*(2*cos(2*x))},{sin(x)*(2*cos(2*x))})--(axis cs:{sqrt(3)/2},0)--cycle;
\addplot[draw={\coloronefill},fill={\coloronefill},area style,domain=0:30]
 ({cos(x)},{sin(x)})--(axis cs:0,0)--cycle;
\addplot [draw={\colorone},thick, smooth,domain=0:360,samples=90]
 ({cos(x)*(2*cos(2*x))},{sin(x)*(2*cos(2*x))});
\draw[draw={\colortwo},thick](axis cs:0,0)circle(1);
\end{axis}
\node [right] at (myplot.right of origin) {\scriptsize $0$};
\node [above] at (myplot.above origin) {\scriptsize $\pi/2$};
\end{tikzpicture}
\\(a)\\[10pt]
\begin{tikzpicture}[alt={Close-up: sector α ≤ θ ≤ β between the two curves is filled to illustrate the integrand region.}]
\begin{axis}[width=\marginparwidth,tick label style={font=\scriptsize},
axis y line=middle,axis x line=middle,name=myplot,axis on top,
ymin=-.1,ymax=1,xmin=-.1,xmax=1.2]
\addplot[draw={\colorone},fill={\coloronefill},area style,domain=30:45]
 ({cos(x)*(2*cos(2*x))},{sin(x)*(2*cos(2*x))})--(axis cs:{sqrt(3)/2},0)--cycle;
\addplot[draw={\colortwo},fill={\colortwofill},area style,domain=0:30]
 ({cos(x)},{sin(x)})--(axis cs:0,0)--cycle;
\addplot [draw={\colorone},thick, smooth,domain=0:360,samples=90]
 ({cos(x)*(2*cos(2*x))},{sin(x)*(2*cos(2*x))});
\draw[draw={\colortwo},thick](axis cs:0,0)circle(1);
\draw [thick,dashed] (axis cs: 0,0) -- (axis cs:.866,.5);
\end{axis}
\node [right] at (myplot.right of origin) {\scriptsize $0$};
\node [above] at (myplot.above origin) {\scriptsize $\pi/2$};
\end{tikzpicture}
\\(b)}
%
We need to find the point of intersection between the two curves. Setting the two functions equal to each other, we have
\[
2\cos(2\theta) = 1 \quad \Rightarrow \quad \cos(2\theta) = \frac12
\quad \Rightarrow \quad 2\theta = \pi/3\quad \Rightarrow \quad \theta=\pi/6.
\]
In part (b) of the figure, we zoom in on the region and note that it is not really bounded \emph{between} two polar curves, but rather \emph{by} two polar curves, along with $\theta=0$. The dashed line breaks the region into its component parts. Below the dashed line, the region is defined by $r=1$, $\theta=0$ and $\theta = \pi/6$. (Note: the dashed line lies on the line $\theta=\pi/6$.) Above the dashed line the region is bounded by $r=2\cos(2\theta)$ and $\theta =\pi/6$. Since we have two separate regions, we find the area using two separate integrals.

Call the area below the dashed line $A_1$ and the area above the dashed line $A_2$. They are determined by the following integrals:
\[
A_1 = \frac12\int_0^{\pi/6} (1)^2\dd\theta\qquad
A_2 = \frac12\int_{\pi/6}^{\pi/4} \bigl(2\cos(2\theta)\bigr)^2\dd\theta.
\]
(The upper bound of the integral computing $A_2$ is $\pi/4$ as $r=2\cos(2\theta)$ is at the pole when $\theta=\pi/4$.)

We omit the integration details and let the reader verify that $A_1 = \pi/12$ and $A_2 = \pi/12-\sqrt{3}/8$; the total area is $A = \pi/6-\sqrt{3}/8$.
\end{example}

\subsection{Arc Length}

As we have already considered the arc length of curves defined by rectangular and parametric equations, we now consider it in the context of polar equations. Recall that the arc length $L$ of the graph defined by the parametric equations $x=f(t)$, $y=g(t)$ on $[a,b]$ is
\index{arc length}\index{polar!function!arc length}
\begin{equation}
L = \int_a^b \sqrt{[\fp(t)]^2 + [g\primeskip'(t)]^2}\dd t = \int_a^b \sqrt{[x\primeskip'(t)]^2+[y\primeskip'(t)]^2}\dd t.
\label{eq:polar_arclength}
\end{equation}

Now consider the polar function $r=f(\theta)$. We again use the identities $x=f(\theta)\cos\theta$ and $y=f(\theta)\sin\theta$ to create parametric equations based on the polar function. We compute $x\primeskip'(\theta)$ and $y\primeskip'(\theta)$ as done before when computing $\frac{\dd y}{\dd x}$, then apply \autoeqref{eq:polar_arclength}.

The expression $[x\primeskip'(\theta)]^2+[y\primeskip'(\theta)]^2$ can be simplified a great deal; we leave this as an exercise and state that
\[
[x\primeskip'(\theta)]^2+[y\primeskip'(\theta)]^2 = [\fp(\theta)]^2+[f(\theta)]^2.
\]
This leads us to the  arc length formula.

\begin{keyidea}[Arc Length of Polar Curves]\label{idea:polar_arclength}%
Let  $r=f(\theta)$ be a polar function with $\fp$ continuous on an open interval $I$ containing $[\alpha,\beta]$, on which the graph traces itself only once. The arc length $L$ of the graph on $[\alpha,\beta]$ is
\[L = \int_\alpha^\beta \sqrt{[\fp(\theta)]^2+[f(\theta)]^2}\dd\theta = \int_\alpha^\beta\sqrt{(r\,')^2+ r^2}\dd\theta.\]
\end{keyidea}

\begin{example}[Arc Length of Polar Curves]\label{ex_cardioid_length}%
Find the arc length of the cardioid $r=1+\cos \theta$.
\solution
With $r=1+\cos \theta$, we have $r' = -\sin \theta$. The cardioid is traced out once on $[0,2\pi]$, giving us our bounds of integration. Applying \autoref{idea:polar_arclength} we have
\begin{align*}
	L
	& = \int_0^{2\pi} \sqrt{(-\sin \theta)^2  + (1 + \cos \theta)^2}\dd\theta \\
	& = \int_0^{2\pi} \sqrt{\sin^2 \theta + (1 + 2\cos \theta + \cos^2 \theta)}\dd\theta\\
	& = \int_0^{2\pi} \sqrt{2 + 2\cos \theta}\dd\theta\\
	& = \int_0^{2\pi} \sqrt{2 + 2\cos \theta}\ \frac{\sqrt {2-2\cos \theta}}{\sqrt {2-2\cos \theta}}\dd\theta\\
	& = \int_0^{2\pi} \frac{\sqrt{4 - 4\cos^2 \theta}}{\sqrt {2-2\cos \theta}}\dd\theta\\
	& = 2\int_0^{2\pi} \frac{\sqrt{1 - \cos^2 \theta}}{\sqrt {2-2\cos \theta}}\dd\theta\\
	& = 2\int_0^{2\pi} \frac{\abs{\sin\theta}}{\sqrt {2-2\cos \theta}}\dd\theta
\end{align*}
Since the $\sin \theta > 0$ on $[0, \pi]$ and $\sin \theta < 0$ on $[\pi, 2\pi]$ we separate the integral into two parts
\[
2\int_0^{\pi} \frac{\sin \theta}{\sqrt {2-2\cos \theta}}\dd\theta - 2\int_{\pi}^{2\pi} \frac{\sin \theta}{\sqrt {2-2\cos \theta}}\dd\theta
\]
Using the symmetry of the cardioid and $u$-substitution ($u = 2 - 2\cos \theta$) we simplify the integration to
\begin{align*}
L
& = 4\int_0^{\pi} \frac{\sin \theta}{\sqrt {2-2\cos \theta}}\dd\theta\\
& = 2\int_0^4 \frac{1}{\sqrt u}\dd u\\
& = 4  u^{1/2}\biggr|_0^4 = 8.
\end{align*}
\end{example}

\begin{example}[Arc length of a limaçon]\label{ex_polcalc7}%
Find the arc length of the limaçon $r=1+2\sin\theta$.
\solution
With $r=1+2\sin\theta$, we have $r\,' = 2\cos\theta$. The limaçon is traced out once on $[0,2\pi]$, giving us our bounds of integration. Applying \autoref{idea:polar_arclength}, we have
\begin{align*}
	L
	&= \int_0^{2\pi} \sqrt{(2\cos\theta)^2+(1+2\sin\theta)^2}\dd\theta \\
	&=	\int_0^{2\pi} \sqrt{4\cos^2\theta+4\sin^2\theta +4\sin\theta+1}\dd\theta\\
	&=	\int_0^{2\pi} \sqrt{4\sin\theta+5}\dd\theta\\
	&\approx 13.3649.
\end{align*}
%
\mtable{The limaçon in \autoref{ex_polcalc7} whose arc length is measured.}{fig:polcalc7}{\begin{tikzpicture}[alt={Limacon r = 1 + 2 cos θ plotted; small inner loop (θ around π) and outer loop (whole 0 ≤ θ < 2 π) }]
\begin{axis}[width=1.16\marginparwidth,tick label style={font=\scriptsize},
axis y line=middle,axis x line=middle,name=myplot,axis on top,
ymin=-.5,ymax=3.2,xmin=-2.22,xmax=2.22]
\addplot [draw={\colorone},thick, smooth,domain=0:360,samples=90]
 ({cos(x)*(1+2*sin(x))},{sin(x)*(1+2*sin(x))});
\end{axis}
\node [right] at (myplot.right of origin) {\scriptsize $0$};
\node [above] at (myplot.above origin) {\scriptsize $\pi/2$};
\end{tikzpicture}}
%
The final integral cannot be solved in terms of elementary functions, so we resorted to a numerical approximation. (Simpson's Rule, with $n=4$, approximates the value with $13.0608$. Using $n=22$ gives the value above, which is accurate to 4 places after the decimal.)
\end{example}

\subsection{Surface Area}

The formula for arc length leads us to a formula for surface area. The following Key Idea is based on \autoref{idea:surface_area_parametric}.

\begin{keyidea}[Surface Area of a Solid of Revolution]\label{idea:surface_area_polar}%
Consider the graph of the polar equation $r=f(\theta)$, where $\fp$ is continuous on an open interval containing $[\alpha,\beta]$ on which the graph does not cross itself.
\index{surface area!solid of revolution}\index{polar!function!surface area}\index{integration!surface area}
\begin{enumerate}
	\item The surface area of the solid formed by revolving the graph about the initial ray ($\theta=0$) is:
	\[\text{Surface Area} = 2\pi\int_\alpha^\beta f(\theta)\sin\theta\sqrt{[\fp(\theta)]^2+[f(\theta)]^2}\dd\theta.\]
	\item The surface area of the solid formed by revolving the graph about the line $\theta=\pi/2$ is:
	\[\text{Surface Area} = 2\pi\int_\alpha^\beta f(\theta)\cos\theta\sqrt{[\fp(\theta)]^2+[f(\theta)]^2}\dd\theta.\]
\end{enumerate}
\end{keyidea}

\begin{example}[Surface area determined by a polar curve]\label{ex_polcalc8}%
Find the surface area formed by revolving one petal of the rose curve $r=\cos(2\theta)$ about its central axis (see \autoref{fig:polcalc8}).
\solution
\mtable{Finding the surface area of a rose-curve petal that is revolved around its central axis.}{fig:polcalc8}{%
\begin{tikzpicture}[alt={Red and blue roses r = cos 2 θ and r = sin 2 θ share four petals, alternating every π/4.}]
\begin{axis}[width=1.16\marginparwidth,tick label style={font=\scriptsize},
axis y line=middle,axis x line=middle,name=myplot,axis on top,xtick={-1,1},
ytick={-1,1},ymin=-1.1,ymax=1.1,xmin=-1.3,xmax=1.3]
\addplot [draw={\colorone},thick, smooth,domain=0:45,samples=30]
 ({cos(x)*(cos(2*x))},{sin(x)*(cos(2*x))});
\addplot [draw={\colortwo},thick, smooth,domain=45:360,samples=60]
 ({cos(x)*(cos(2*x))},{sin(x)*(cos(2*x))});
\end{axis}
\node [right] at (myplot.right of origin) {\scriptsize $0$};
\node [above] at (myplot.above origin) {\scriptsize $\pi/2$};
\end{tikzpicture}
\\(a)\\[10pt]
\myincludeasythree{
3Droll=0,
3Dortho=0.009,
3Dc2c=0.41893383860588074 -0.887881875038147 0.1901581883430481,
3Dcoo=65.85308837890625 6.341495513916016 17.98039436340332,
3Droo=85}{\marginparwidth}{One blue petal (0 ≤ θ ≤ pi/2) revolved about its axis, forming a translucent spindle-shaped surface for surface-area computation.}{figures/figpolcalc8a_3D}
\\(b)}
We choose, as implied by the figure, to revolve the portion of the curve that lies on $[0,\pi/4]$ about the initial ray. Using \autoref{idea:surface_area_polar} and the fact that $\fp(\theta) = -2\sin(2\theta)$, we have
\begin{align*}
\text{Surface Area}
&= 2\pi\int_0^{\pi/4} \cos(2\theta)\sin(\theta)
\sqrt{\bigl(-2\sin(2\theta)\bigr)^2+\bigl(\cos(2\theta)\bigr)^2}\dd\theta \\
&\approx 1.36707.
\end{align*}
The integral is another that cannot be evaluated in terms of elementary functions. Simpson's Rule, with $n=4$, approximates the value at $1.36751$.%; with $n=10$, the value is accurate to 4 decimal places.
\end{example}


This chapter has been about curves in the plane. While there is great mathematics to be discovered in the two dimensions of a plane, we live in a three dimensional world and hence we should also look to do mathematics in 3D --- that is, in \emph{space}. The next chapter begins our exploration into space by introducing the topic of \emph{vectors}, which are incredibly useful and powerful mathematical objects.

\printexercises{exercises/09-05-exercises}



\cleardoublepage

\index{$"!$|see {factorial}} % 08 sequences
